\documentclass[11pt,a4paper]{article}
\usepackage[margin=27mm]{geometry}
\usepackage{fontspec}
\setmainfont{Latin Modern Roman}
\usepackage{amsmath,amssymb,amsthm,mathtools}
\usepackage{microtype,booktabs,longtable,array}
\usepackage{xcolor}
\usepackage[unicode,colorlinks=true,linkcolor=blue!45!black,citecolor=blue!45!black,urlcolor=blue!45!black]{hyperref}
\usepackage{xurl}
\usepackage[authoryear,round]{natbib}
\setlength{\emergencystretch}{3em}
\renewcommand{\thesection}{\Alph{section}}
\newcommand{\Fp}{\mathbf F_p}
\newcommand{\Z}{\mathbf Z}
\newcommand{\Q}{\mathbf Q}
\newcommand{\F}{\mathbf F}
\newcommand{\G}{\mathcal G}
\newcommand{\cl}{\mathrm{cl}}
\newcommand{\pr}{\mathrm{pr}}
\newcommand{\et}{\text{ét}}
\newcommand{\Gal}{\operatorname{Gal}}
\newcommand{\Hom}{\operatorname{Hom}}
\newcommand{\Aut}{\operatorname{Aut}}
\newcommand{\Out}{\operatorname{Out}}
\newcommand{\im}{\operatorname{im}}
\newcommand{\Spec}{\operatorname{Spec}}
\newcommand{\Pic}{\operatorname{Pic}}
\newcommand{\Inn}{\operatorname{Inn}}
\newcommand{\crr}{\mathrm{c}r}
\providecommand{\mathscr}{\mathcal}
\newenvironment{claim}[3]{\par\medskip\noindent\hypertarget{#1}{\textbf{#1. #2.}}\ \textbf{#3.}\par\nopagebreak\smallskip\begingroup\itshape}{\endgroup\par\smallskip}
\newcommand{\statementref}[1]{\hyperlink{#1}{#1}}
\title{ANABELIAN\_A1\_SUCCESSOR\\Point sections, cuspidal exactness, and arithmetic continuations}
\author{AI-produced by GPT-6 Astra\\OpenAI, ChatGPT Pro mode}
\date{18 November 2026}

\newtheorem{lemma}{Lemma}
\newtheorem{remark}[lemma]{Remark}

\tracinglostchars=3
\setlength{\tabcolsep}{5.8pt}

\begin{document}
\maketitle
\begin{abstract}
We continue the prime-to-characteristic point-separation argument with
explicit dependency separation, direct algebraic configuration exactness,
and independently implemented finite certificates. We retain every earlier
numbered assertion and distinguish checked proofs from work still awaiting
completion. The further arithmetic questions are treated with the same
distinction. No priority claim or different-model-family verification is made.
\end{abstract}
\noindent\textbf{Overall task status: INCOMPLETE.}
Every open item remains an open item. Individual proved results are stated
with their hypotheses and their check records. The exact available model
identification is GPT-6 Astra (OpenAI), ChatGPT Pro mode; a separate internal
reasoning-effort setting is not exposed. Prior external readings are recorded
as prior events, not as source openings or computations in this writing.
\tableofcontents
\part{Checked results}
\section{The point-separation argument}
\subsection{Conventions, conclusion, and the logical chain}
We retain the whole arithmetic quotient. If $k=\mathbf F_q$ has
characteristic $p$, $C/k$ is smooth, proper and geometrically connected,
and $\Sigma$ is a set of primes different from $p$, our notation is
\[
 \Delta_C^\Sigma=\pi_1(C_{\bar k})^\Sigma,\qquad
 1\longrightarrow\Delta_C^\Sigma\longrightarrow
 \Pi_C^\Sigma\longrightarrow G_k\longrightarrow1.
\]
For $\Sigma=\{\ell:\ell\ne p\}$ we omit the superscript. All
homomorphisms of profinite groups are continuous. Point sections are
classes under conjugation by the geometric kernel, not by a finite
approximation to that kernel. We reserve $E_C$ for the image on closed
points of the exceptional set of Sa\"idi--Tamagawa.

\begin{claim}{A.1}{Proved}{The main theorem, including every constant extension}
Let $C/k$ be a smooth proper geometrically connected curve of genus at
least two over a finite field of characteristic $p$. The map
\[
 C(k)\longrightarrow
 \{s:G_k\to\Pi_C\mid \operatorname{pr}\circ s=\operatorname{id}\}/\Delta_C
\]
is injective. The same assertion holds for $C_L/L$ for every finite
extension $L/k$. Thus AS-A.3 (AR.3), with precisely these hypotheses,
is correct. This answers cote 149, scan 36, for proper hyperbolic curves.
\end{claim}
\begin{proof}
Fix $r\ne p$. A.45 extends every finite $r$-group fibre torsor in
the labelled product family, and A.46 proves tame right exactness.
The finite proper-cover and core argument in A.13 therefore gives
the relative configuration sequence. A.47 proves survival of full
$r$-adic cusp inertia algebraically; A.48 constructs the diagonal
boundary group and transports the actual cusp between the two
arithmetic fibres whenever the proper sections are geometrically
conjugate. It retains the entire $G_k$ and is over the identity of
$\Pi_C$.

A.14--A.15 and A.49 recognize the transported noncuspidal sections
as actual points. B.68 identifies their abelian differences with
the positive Kummer classes, so A.21 puts all potentially changed
labels in a proper support field. A.42--A.43 construct and evaluate
the central difference cocycle. A.20's Kummer cover and Weil-bound
argument then force $P=Q$. This chain uses neither lifting to
characteristic zero nor finite exceptional support. Alternatively,
A.16 and A.44 give the finite-exception endgame A.19. The
topological proof A.2--A.5 is retained as a separately checked
route, rather than a dependency of the direct route.

Every argument applies over any finite constant extension $L/k$.
Thus the conclusion is the asserted all-extension injectivity in
the entire arithmetic quotient, under conjugation by its entire
geometric subgroup.
\end{proof}
\paragraph{Check record.} Also checked in this writing, 2026-11-18: GPT-6 Astra (OpenAI, ChatGPT Pro mode), root; separate effort setting not exposed. Checked the complete direct characteristic-p chain A.45–A.48, A.13, A.14–A.15, A.49, B.68, A.42–A.43, A.20–A.21, including full arithmetic quotient and whole-geometric-subgroup conjugacy.



The attached register identifies AS-A.87--A.90 with AG.39--AG.42,
and AS-A.97--A.98 with AR.8--AR.9. We use the working labels only
when discussing that earlier manuscript. Our statements are numbered
independently by Part. The proof audit and the computational checks
are AI-produced; we claim no priority.



\subsection{The relative configuration sequence, with all maps specified}

Throughout this subsection, a finite $\Sigma$-group means a finite group
whose order has all prime divisors in $\Sigma$. Completion means the
inverse limit of all such quotients. A prime $r\in\Sigma$ is fixed.
For a complex smooth proper curve $C$ of genus $g\geq2$, put
\[
 \Lambda=\pi_1(C-P),\quad \Gamma=\pi_1((C\times C)-\operatorname{diag}),
 \quad B=\pi_1(C),\quad \Delta=B^{\Sigma}.
\]
The second projection is used, so that $\Lambda$ is its discrete normal
fibre subgroup. This exact discrete sequence is the
Fadell--Neuwirth sequence; the precise surface hypotheses and the
vanishing of $\pi_2(C)$ are given in
\citet[Section 2.1]{Bellingeri2004}. Superscripts $\Sigma$ below denote
pro-$\Sigma$ completion, not a quotient of the arithmetic Galois group.

\begin{claim}{A.2}{Proved}{The free-fibre monodromy lemma}
Let $1\to L\to H\to B_0\to1$ be an exact sequence of discrete groups,
where $L$ is free of finite rank at least two. If conjugation by $H$
acts through an $r$-group on $L/[L,L]L^r$, the natural continuous map
$L^{(r)}\to H^{(r)}$ is injective. Every prescribed finite $r$-quotient
of $L$ is consequently dominated by the restriction of a finite
$r$-quotient of $H$.
\end{claim}
\begin{proof}
Write $L_r=L^{(r)}$. Its iterated Frattini subgroups are open and
characteristic, and they are cofinal: their images eventually vanish in
each finite $r$-quotient. For a finite $r$-group $Q$, the kernel of
$\Aut(Q)\to\Aut(Q/\Phi(Q))$ is an $r$-group. Indeed, it acts freely
on the $|\Phi(Q)|^d$ ordered lifts of a fixed basis of $Q/\Phi(Q)$,
and every such lift generates $Q$. Passing to characteristic finite
quotients proves that the kernel of
\[
 \Aut(L_r)\longrightarrow\Aut(L_r/\Phi(L_r))
\]
is pro-$r$. Thus the closure of the conjugation image of $H$ in the
profinite group $\Aut(L_r)$ is pro-$r$. The action extends continuously
to $H^{(r)}$.

For completeness, $L_r$ has trivial centre. The Magnus map sends its
free generators to $1+X_i$ in the units of the noncommutative formal
power-series algebra $\F_r\langle\!\langle X_1,\ldots,X_d\rangle\!\rangle$.
It is injective: for each finite $r$-quotient $Q$ of $L_r$, substitution
$X_i\mapsto q_i-1$ is defined because the augmentation ideal of
$\F_r[Q]$ is nilpotent, and it recovers the quotient map to $Q$.
Nilpotence follows by induction on $|Q|$, using a central subgroup of
order $r$ and $(z-1)^r=0$. A power series commuting with both $X_1$
and $X_2$ is scalar. To see this, compare homogeneous terms: a
homogeneous polynomial commuting with $X_1$ has only the monomial
$X_1^n$, by comparing initial and terminal letters successively; then
commutation with $X_2$ eliminates every positive-degree term. A central
Magnus unit of augmentation one is therefore one.

The composite $L_r\to H^{(r)}\to\Aut(L_r)$ is the inner action,
so it is injective. A continuous injective map from a compact group
into a Hausdorff profinite group is a closed embedding. Given an open
normal subgroup $W$ of $L_r$, there is an open normal subgroup $N$ of
$H^{(r)}$ with $N\cap L_r\subset W$: use the induced topology and
then take the finite normal core of an open neighbourhood. The quotient
$H^{(r)}/N$ gives the stated domination.
\end{proof}
\paragraph{Check record.} Also checked in this writing, 2026-11-18: GPT-6 Astra Pro; Pro mode; separate effort setting not exposed; ANABELIAN\_A1\_SUCCESSOR/configuration. Blind proof saved before reading the prior proof or audit; SHA256 d435494f\allowbreak{}67de82bd\allowbreak{}06f08a04\allowbreak{}cc519501\allowbreak{}e739a7ab\allowbreak{}f0f88dc3\allowbreak{}b8b42ef5\allowbreak{}c96c6ac9. | Also checked in this writing, 2026-11-18: GPT-6 Astra Pro; Pro mode; separate effort setting not exposed; ANABELIAN\_A1\_SUCCESSOR/configuration. Every source line1–140 compared after blind save; both arguments valid; finite automorphism-centre distinction resolved in S1\_comparison.tex.


\begin{claim}{A.3}{Proved}{Exactness for a relative single-prime cusp kernel}
Set $F=\Lambda^\Sigma$, $E=\Gamma^\Sigma$, and
\[
 \begin{split}
 K_F&=\ker(F\to\Delta),\qquad
 K_E=\ker(E\to\Delta\times\Delta),\\
 F_{\crr}&=F/\ker(K_F\to K_F^{(r)}),\qquad
 E_{\crr}=E/\ker(K_E\to K_E^{(r)}).
 \end{split}
\]
The induced continuous sequence
\[
 1\longrightarrow F_{\crr}\longrightarrow E_{\crr}
       \xrightarrow{\pr_2}\Delta\longrightarrow1
\]
is exact. Every finite quotient of $F_{\crr}$ is dominated on the
fibre by a finite quotient of $E_{\crr}$ which has a normal $r$-subgroup
and a quotient $A\times A$, for a finite $\Sigma$-quotient $A$ of
$\Delta$.
\end{claim}
\begin{proof}
Let $q:F_{\crr}\twoheadrightarrow Q$ be continuous with $Q$ finite.
The image $R$ of $K_F^{(r)}$ is a normal $r$-subgroup, and
$A=Q/R$ is a quotient of $\Delta$. Let $a:C'\to C$ be the connected
Galois proper cover corresponding to this quotient. Put
\[
 Y=(C'\times C')-(a\times a)^{-1}(\operatorname{diag} C).
\]
Its removed divisor is the disjoint union of the graphs of the deck
transformations of $a$, so $Y\to C'$ has individually labelled
punctures. Choose a point of $C'$ above $P$. The fibre is
$V'=C'-a^{-1}(P)$. The surface fibration gives a discrete exact sequence
\[
 1\to L=\pi_1(V')\to H=\pi_1(Y)\to\pi_1(C')\to1.
\]
Here $L$ is free of rank $2g(C')+|A|-1\geq2$. In
$H_1(V',\F_r)$ the boundary classes span a subspace $W$, and there is
an exact sequence
\[
 0\to W\to H_1(V',\F_r)\to H_1(C',\F_r)\to0.
\]
Monodromy acts trivially on $W$, because the boundary sections are
labelled and oriented. It acts trivially on the quotient, because the
compactified family is the product $C'\times C'\to C'$. Choosing any
vector-space splitting shows that its image is contained in the additive
$r$-group $\Hom_{\F_r}(H_1(C',\F_r),W)$. This does not require
$r\nmid|A|$. By \statementref{A.2}, $L^{(r)}$ embeds in $H^{(r)}$.

The subgroup $H$ is normal in $\Gamma$, with quotient $A\times A$,
and $\Lambda\cap H=L$. The restriction of $q$ to $L$ takes values in
$R$. Choose a finite $r$-quotient of $H$, with kernel $J$, whose
restriction dominates this quotient of $L$. Thus
$J\cap L\subset\ker(q)$. Take
\[
 J^*=\bigcap_{\gamma\in\Gamma/H}\gamma J\gamma^{-1}.
\]
This is a finite intersection since $J$ is normal in $H$. The group
$H/J^*$ embeds in a finite product of $r$-groups. Consequently
$\Gamma/J^*$ is a finite $\Sigma$-group, an extension of $A\times A$
by an $r$-group, and $J^*\cap\Lambda\subset\ker(q)$. This last
containment remains true on completed fibre groups, as is seen in the
joint finite image of the two quotient maps.

The quotient of $\Gamma$ factors through $E_{\crr}$. Its image of
$K_E$ lies in $H/J^*$ because the $A\times A$ quotient comes from
the proper product; it therefore kills $\ker(K_E\to K_E^{(r)})$.
Conversely, the fibre map $F\to E_{\crr}$ has pro-$r$ image on
$K_F$, so it factors through $F_{\crr}$. Detection of every finite
quotient proves that this induced map is injective.

Finally, right exactness of pro-$\Sigma$ completion follows directly
from its universal property: quotienting $E$ by the closed image of
$F$ represents precisely the finite $\Sigma$-quotients of $B$, and
hence gives $\Delta$. The subgroup killed to form $E_{\crr}$ is
inside $K_E$, which is inside this fibre kernel. Thus the image of
$F_{\crr}$ is the whole kernel of $\pr_2$, and $\pr_2$ is onto.
\end{proof}
\paragraph{Check record.} Also checked in this writing, 2026-11-18: GPT-6 Astra Pro; Pro mode; separate effort setting not exposed; ANABELIAN\_A1\_SUCCESSOR/configuration. Blind proof saved before reading the prior proof or audit; SHA256 d435494f\allowbreak{}67de82bd\allowbreak{}06f08a04\allowbreak{}cc519501\allowbreak{}e739a7ab\allowbreak{}f0f88dc3\allowbreak{}b8b42ef5\allowbreak{}c96c6ac9. | Also checked in this writing, 2026-11-18: GPT-6 Astra Pro; Pro mode; separate effort setting not exposed; ANABELIAN\_A1\_SUCCESSOR/configuration. Every source line1–140 compared after blind save; both arguments valid; finite automorphism-centre distinction resolved in S1\_comparison.tex.


\begin{claim}{A.4}{Proved}{Specialization and the arithmetic fibre}
Let $C$ be smooth proper geometrically connected of genus at least two
over an algebraically closed field of characteristic $p$, and let
$\Sigma$ be any set of primes not containing $p$, with $r\in\Sigma$.
The preceding relative exact sequence holds for geometric etale groups.
If $C/k$ is defined over a finite field and $P\in C(k)$, its arithmetic
cuspidally pro-$r$ group is the inverse image of a chosen point section
$s_P(G_k)$ under the second projection of the arithmetic configuration
quotient. The comparison is over the identity of the whole $G_k$.
\end{claim}
\begin{proof}
We record the exact specialization input. SGA~1, Expose XIII,
Corollary~2.9, assumes a smooth proper morphism of finite presentation
$X\to S$, the complement $U$ of a relative normal-crossings divisor,
and the set of primes different from every residual characteristic of
$S$. For a locally constant constructible sheaf of finite groups with
these primes and tame along the boundary, it gives locally constant
$R^1$ and cohomological properness through degree one. Section~2.10
constructs the specialization homomorphism by strict localization and
a path; its ambiguity is inner conjugation. Applied to constant finite
groups, Corollary~2.9 identifies their torsors on the two geometric
fibres, so the resulting maximal prime-to-$p$ specialization map is an
isomorphism. It is functorial in the pair. Lemma~2.11 identifies the
images of inertia groups under this specialization. These statements
are in the opened text at pages 288--290, corresponding to original
margin pages 389--392 of \cite{LA_SGA1}.

Here a suitable lift exists. SGA~1, Expose III, Theorem~7.3 assumes
$H^2(C,T_C)=H^2(C,\mathcal O_C)=0$ for a smooth projective special
fibre over the residue field of a complete local ring and gives a
smooth projective lift; Corollary~7.4 states the curve case. Both
vanishings hold for a curve. A chosen point lifts successively by
formal smoothness, and the compatible points give a section over the
complete ring. Equivalently, the marked deformation obstruction lies
in $H^2(C,T_C(-P))=0$. Take the Witt ring of a perfect algebraically
closed field of definition. The pairs
\[
 (\mathcal C,\mathcal P),\qquad
 (\mathcal C\times\mathcal C,\operatorname{diag}\mathcal C)
\]
are smooth proper pairs with smooth relative boundary. Their fibre
inclusion is transverse with multiplicity one. Functorial
specialization therefore identifies the groups and maps used in
\statementref{A.3}; the characteristic relative kernels and their
pro-$r$ quotients are identified as well.

For general algebraically closed fields, first descend the pointed
curve to a finitely generated field over $\F_p$, and take its
algebraic closure there. This is countable. Its Witt fraction field has
cardinality at most that of $\mathbb C$, so an algebraic closure
embeds into $\mathbb C$. Invariance of tame prime-to-$p$ groups of
smooth proper pairs under algebraically closed field extension, also
obtained from the same proper-pair comparison of finite torsors,
permits this descent and the return to the original field. Taking a
further maximal pro-$\Sigma$ quotient transfers the result for every
$\Sigma$ under consideration.

For the arithmetic assertion, write $\Pi_C$ for the geometrically
pro-$\Sigma$ group over the unchanged $G_k$. The configuration group
maps onto $\Pi_C\times_{G_k}\Pi_C$. Its kernel relative to this
product is geometric; replacing it by its pro-$r$ quotient gives a
group $\mathcal E_{\crr}$ with geometric subgroup $E_{\crr}$.
The actual fibre $C-P$ induces a map of extensions
\[
 \Pi_{C-P,\crr}\longrightarrow
 \mathcal E_{\crr}\times_{\Pi_C,s_P}G_k.
\]
On geometric kernels it is the isomorphism of
\statementref{A.3}, and on arithmetic quotients it is the identity
of $G_k$. Injectivity and surjectivity follow by the elementary
comparison of short exact sequences. Since $s_P$ is injective, the
right-hand group identifies with $\pr_2^{-1}(s_P(G_k))$.
Paths give exactly geometric inner conjugacy, not a change of the
arithmetic quotient.
\end{proof}
\paragraph{Check record.} Also checked in this writing, 2026-11-18: GPT-6 Astra Pro; Pro mode; separate effort setting not exposed; ANABELIAN\_A1\_SUCCESSOR/configuration. Specialization and arithmetic fibre proved with opened SGA III7.3–7.4 and XIII2.9–2.11; hypotheses verified; preserves whole G\_k.


\begin{claim}{A.5}{Proved}{The common diagonal boundary and cusp transport}
In the arithmetic situation of \statementref{A.4}, the image
$\mathcal D$ of a fixed diagonal boundary component fits into a
continuous exact sequence
\[
 1\to\Z_r(1)\to\mathcal D\to\Pi_C\to1,
\]
and maps diagonally to $\Pi_C\times_{G_k}\Pi_C$. With paths chosen
through this boundary, the actual cusp subgroup of $C-P$ is
$\mathcal D\cap\pr_2^{-1}(s_P(G_k))$. Consequently, geometrically
conjugate proper point sections give an isomorphism of their actual
cuspidally pro-$r$ punctured groups over the identity of $\Pi_C$,
carrying the cusp subgroup to a geometric conjugate of the other cusp.
\end{claim}
\begin{proof}
Put $X=C\times C$, $D=\operatorname{diag} C$, and
$N=\mathcal O_X(D)|_D$. For an integer $n$ prime to $p$, the boundary
of $\sqrt[n]{(X,D)}$ is the gerbe of $n$-th roots of $N$ on $D$.
Here is the tame comparison being used, rather than a configuration
fibre exactness assertion. SGA~1 XIII, Proposition~5.2, assumes a
regular local scheme, a normal-crossings divisor with local equations
$f_i$, and a finite cover tame along that divisor. If $n_i$ is the
least common multiple of the relevant inertia orders, these integers
are invertible and the cover extends uniquely to a finite etale cover
after adjoining $T_i$ with $T_i^{n_i}=f_i$. The conclusion is over the
entire Kummer chart; the proof invokes purity. This is the opened
statement at pages 315--316, original margin page 429, of \cite{LA_SGA1}.
Thus every finite tame cover of $X-D$ extends to a sufficiently high
root stack. Restriction to its boundary gerbe defines compatible maps
from the boundary fundamental groups to the configuration group.

A finite etale cover of the boundary gerbe with trivial stabilizer
action descends to $D$, and covers of $D$ pull back fully faithfully.
Etale locally the gerbe is $B\mu_n$; consequently the kernel of its
map to $\pi_1(D)$ is the image of $\mu_n$. An injection of this
finite stabilizer has not been assumed. Taking the inverse limit over
root orders, and then the geometrically pro-$\Sigma$ quotient, gives
a right exact sequence
\[
 \widehat\Z^{\Sigma}(1)\longrightarrow\mathcal T
                 \longrightarrow\Pi_C\longrightarrow1.
\]
Exactness follows also directly from the same trivial-stabilizer
descent criterion. The transition maps of the root gerbes are
surjective on fundamental groups: their maps on the base are the
identity and the maps of stabilizer images are onto. Compact inverse
limits preserve these right exact sequences. Right exactness of
pro-$\Sigma$ completion shows that discarding the geometric base
kernel introduces no further kernel over $\Pi_C$.

Let $\mathcal D$ be the image of $\mathcal T$ in
$\mathcal E_{\crr}$. Its map to the proper product is the injective
diagonal on $\Pi_C$, so its kernel over $\Pi_C$ is precisely the
image of inertia, a quotient of $\Z_r(1)$. It is the full
$\Z_r(1)$. Indeed, the geometric group of the fibre $C-P$ maps
onto the free pro-$r$ group on $2g$ generators. Its boundary loop is
the product of their commutators, up to orientation. Sending one pair
to $1+e_{12},1+e_{23}$ in the upper-unitriangular group over $\Z_r$
and the others to one sends that boundary loop to $1+e_{13}$, of
infinite $r$-adic order. The fibre injection of
\statementref{A.3} retains this inertia in the total quotient.
Transversality identifies it with diagonal inertia with multiplicity
one. This argument includes $r=2$.

The local transverse Kummer chart sends the actual cusp decomposition
group into $\mathcal D$ over $s_P(G_k)$, identically on inertia and
on the arithmetic residue group. Comparing the two short exact
sequences therefore identifies it with the whole inverse image of
$s_P(G_k)$ in $\mathcal D$. The paths for $P$ and $Q$ may both be
chosen in the same connected boundary tower, so this uses the same
$\mathcal D$ without assuming any prior separation of points.

Suppose $\gamma\in\Delta$ conjugates $s_P(G_k)$ to $s_Q(G_k)$.
Choose a geometric $d\in\mathcal D$ over $\gamma$. Conjugation by
$d$ identifies the two fibre subgroups and their cusp subgroups; on
the first proper coordinate it induces $\Inn(\gamma)$. The first
projection from the geometric part of the $Q$-fibre onto $\Delta$
is surjective. Choose there an $e$ mapping to $\gamma^{-1}$. Then
$\Inn(e)\Inn(d)$ is over the identity of $\Pi_C$ and $G_k$, and
carries the first cusp to a geometric conjugate of the second. No
lifting of a point section to a section of the configuration fibration
has been assumed.
\end{proof}
\paragraph{Check record.} Also checked in this writing, 2026-11-18: GPT-6 Astra Pro; Pro mode; separate effort setting not exposed; ANABELIAN\_A1\_SUCCESSOR/configuration. Boundary/transport proof verified; root-stack descent distinguished from SGA chart theorem. Alternative direct all-Sigma proof supplied as A.48.


\subsection{Finite positive and negative tests in genus two}

\begin{claim}{A.6}{Proved}{A mixed-prime finite quotient detecting the relative fibre}
For genus two, $p=2$, $r=3$, and
$A=(\Z/3)^4\times\Z/5$, there is a quotient of the fibre of
order $1215$, with central cuspidal subgroup of order three and proper
quotient $A$, which embeds in a quotient of the configuration group of
order $492075$ with central kernel of order three over $A\times A$.
The defining surface braid relations and the fibre injection have been
checked by the retained program.
\end{claim}
\begin{proof}
Let $V=\F_3^4$ have standard alternating form
$J(v,w)=v_1w_2-v_2w_1+v_3w_4-v_4w_3$. On
\[
 G=(V\oplus V)\times\F_3\times\F_5^2
\]
define multiplication by adding the noncentral coordinates and adding
$2J(v_1+v_2,w_1+w_2)$ to the central coordinate. This is associative
because the added function is bilinear; $2$ is $1/2$ in $\F_3$.
Its order is $3^9\cdot25=492075$. Factor interchange preserves the
law. Form $G\rtimes\langle\tau\rangle$, where $\tau$ interchanges
the two $V$ coordinates and the two $\F_5$ coordinates.

Use \citet[Theorem 1.2]{Bellingeri2004} with $n=g=2$. Write its first
strand generators as $a,b,c,d$, put
$t=[a,b^{-1}][c,d^{-1}]$, and denote the interchange generator by
$\sigma$. Send $a,b,c,d$ to the four standard basis vectors in the
first $V$, giving $a$ also the coordinate $(1,0)$ in $\F_5^2$;
all central coordinates are zero. Send $\sigma$ to the product of
central coordinate $2$ with $\tau$. The eleven relations of the
presentation become identities: the four same-generator relations
use zero self-pairing, the four different-handle relations use zero
cross-handle pairing, the two same-handle mixed relations use the fact that
$[a,b]$ and $[c,d]$ both have central coordinate $1$ in $\F_3$,
and the surface relation
uses $-2=1$ in $\F_3$, giving $t=\sigma^2$.
The code checks each relation directly in this finite group.

The pure braid image lies in $G$ and is all of $G$. The first and
second strand images give eight independent $\F_3$ coordinates;
$a^3$ and its $\sigma$-conjugate give the two $\F_5$ coordinates,
and $[a,b]$ gives the central coordinate. The first fibre image is
\[
 Q=V\times\F_3\times\F_5\hookrightarrow G,
 \qquad (v,z,f)\longmapsto((v,0),z,(f,0)).
\]
This is injective by its coordinates and has order $3^5\cdot5=1215$.
The normal closure of $t$ is the central $\F_3$, in both groups.
Thus these finite maps factor through the respective cuspidally
pro-three groups over the full pro-odd proper groups, and their
proper quotients are exactly $A$ and $A\times A$. The program also
enumerates the $1215$ fibre elements and checks closure under its
specified generators. By \statementref{A.4}, the same geometric
finite test applies to the genus-two curve $y^2+y=x^5$ over
$\overline{\F}_2$.
\end{proof}
\paragraph{Check record.} Also checked in this writing, 2026-11-18: GPT-6 Astra Pro; Pro mode; separate effort setting not exposed; ANABELIAN\_A1\_SUCCESSOR/configuration. Finite group law, orders, commutator signs and source braid conventions checked mathematically. Computational checks are separately attributed to finite\_checks agent.


\begin{claim}{A.7}{Proved}{An explicit finite obstruction to full pro-odd exactness}
Let $F_4=\langle a,b,c,d\rangle$ be the first fibre in the genus-two
pure surface braid group $P_2$. Map $F_4$ onto the Heisenberg group
$H=\operatorname{UT}_3(\F_3)$ by sending $a,b$ to its standard
generators and $c,d$ to one. Let $K$ be the kernel and
$N=[K,K]K^5$. Then the finite odd quotient $F_4/N$, of order
$27\cdot5^{82}$, is not dominated on the fibre by any finite odd
quotient of $P_2$. In particular $F_4^{(2')}\to P_2^{(2')}$ is not
injective. The conclusion applies to the geometric groups of
$y^2+y=x^5$ in characteristic two by specialization.
\end{claim}
\begin{proof}
This is a finite certificate, supplemented by the proof that it
obstructs every possible finite odd total quotient. In Bellingeri's
presentation used above, set
$C=\sigma c\sigma^{-1}$ and $D=\sigma d\sigma^{-1}$.
Conjugation by them fixes $a,b$. On the other generators it is
\[
 \begin{array}{ll}
 \varphi_C(c)=tct^{-1},&
 \varphi_C(d)=tdt^{-1}(tc^{-1}tct^{-1}),\\
 \varphi_D(c)=ct^{-1},&
 \varphi_D(d)=tdt^{-1}.
 \end{array}                                                    \tag{CF.1}
\]
These are the actual conjugation automorphisms, not guessed
endomorphisms. For example the same-generator relation and its
$\sigma$-conjugate give
$CcC^{-1}=tct^{-1}$ and
$\varphi_C(t)=tc^{-1}tct^{-1}$. The mixed relation gives
$\varphi_C(dt^{-1})=tdt^{-1}$. Its conjugate gives
$\varphi_D(c)=ct^{-1}$; the analogous same-generator relation gives
the formula for $d$. The different-handle relations give the asserted
fixing of $a,b$. Normality of the discrete fibre follows from the
Fadell--Neuwirth sequence. The code independently verifies these
rewritten relations and the images of $t$ by free reduction.

The Heisenberg group is represented by triples with product
\[
 (x,y,z)(u,v,w)=(x+u,y+v,z+w+xv)\quad\text{over }\F_3.
\]
Both $\varphi_C^3$ and $\varphi_D^3$ preserve the quotient map to
$H$, as direct evaluation of all four generator images verifies.
Thus they normalize $K$, and they normalize $N$. A spanning tree in
the Cayley graph of this quotient gives a free Schreier basis of $K$
of rank $1+27(4-1)=82$. Its homology coordinates are computed by
recording oriented non-tree edges; the program verifies that its
basis-coordinate matrix is the identity.

Let $z=[a,b]\in H$. In the integral homology of $K$, write $c_s$
and $d_s$ for the lifted $c,d$ loops at vertex $s$. They are non-tree
edges, since $c,d$ map to one. On the independent vectors
\[
                 x=c_1-c_z,\qquad y=d_1-d_z
\]
the complete integral Schreier calculation gives
\[
 M_C=\begin{pmatrix}1&3\\0&1\end{pmatrix},\qquad
 M_D=\begin{pmatrix}1&0\\-3&1\end{pmatrix}.                    \tag{CF.2}
\]
Here $M_C$ is the action of $C^3$ and $M_D$ that of $D^3$.
The program exports the entire two $82$ by $82$ integral matrices
and checks all four equalities defining (CF.2), so this is not a
numerical approximation or a guessed invariant subspace. Hence
\[
 M_CM_D=\begin{pmatrix}-8&3\\-3&1\end{pmatrix},\qquad
 (M_CM_D)(x-y)=-(x-y)\pmod5.
\]
The class $x-y$ in $K/N$ is nonzero, having four nonzero Schreier
coordinates.

Suppose a finite odd quotient $P_2\to T$ with kernel $J$ dominated
$F_4/N$. Then $J\cap F_4\subset N$. Let
$H_0=\langle F_4,C^3,D^3\rangle$. It normalizes $K$ and $N$.
The quotient of the image of $K$ by the image of $N$ in $T$ is
exactly $K/N$, since $J\cap K\subset N$. The conjugation action of
$H_0$ on it therefore factors through the finite odd group
$\im(H_0\to T)$. In particular $C^3D^3$ must act with odd order.
Its action on the nonzero line generated by $x-y$ has order two,
a contradiction. No invariance of $K$ under the entire $P_2$ was
needed. Finally, an injective map between the two compact pro-odd
groups would induce the full topology on the fibre and would supply
a dominating finite total quotient for this specified finite fibre
quotient. Thus it cannot be injective.

The retained files are
\path{code/configuration_finite_tests.py}, its exact output, and
\path{code/configuration_finite_certificates.json}; the latter
contains the presentation inputs, the free words, the Schreier basis,
the integral matrices, and the sparse vector. The factorized arithmetic size is $27\cdot5^{82}$; its exact decimal
value is also retained in the program output.
\end{proof}
\paragraph{Check record.} Also checked in this writing, 2026-11-18: GPT-6 Astra Pro; Pro mode; separate effort setting not exposed; ANABELIAN\_A1\_SUCCESSOR/configuration. Normalizer obstruction checked: finite odd total quotient would act with odd order on K/N, contradicting the certified minus-one eigenline. Computation attributed separately. | Also checked in this writing, 2026-11-18: GPT-6 Astra Pro; finite\_checks agent; effort not separately exposed. Read full source proof, verified normalizer argument and finite odd obstruction; independently recomputed full integral matrices; characteristic-p specialization entrusted to configuration reader, not separately read here.


\subsection{Which sets of primes preserve the full configuration fibre?}

\begin{claim}{A.8}{Proved}{Exactness would pass to proper covers}
Suppose $C/\mathbb C$ has genus at least two and the full pro-$\Sigma$
fibre map for $(C\times C)-\operatorname{diag}$ is injective.
For every connected finite Galois etale $\Sigma$-cover $C'\to C$,
the analogous full fibre map for $C'$ is injective as well.
\end{claim}
\begin{proof}
First, if a subgroup $H$ of a discrete group $G$ contains a normal
subgroup of finite $\Sigma$-index, it has its full pro-$\Sigma$
topology induced from $G$. Given a normal $\Sigma$-index subgroup
of $H$, intersect with the fixed normal subgroup and take the
$G$-core. The quotient of the latter normal subgroup embeds in a
finite product of $\Sigma$-groups, and the resulting quotient of
$G$ is an extension of $\Sigma$-groups. It is therefore a
$\Sigma$-group. This proves the topology assertion and the
identification of the corresponding completed open subgroup.

Let $a:C'\to C$ have group $A$, and set
$Y=(C'\times C')-\bigcup_{\alpha\in A}\operatorname{graph}(\alpha)$.
It is a normal $A\times A$ cover of the original configuration
space. Choose $P'\in C'$ above $P$. Its fibre is
$V'=C'-a^{-1}(P)$, the corresponding proper-cover subgroup of
$C-P$. The preceding topology assertion and the assumed injection
identify $(\pi_1V')^\Sigma$ with a closed normal subgroup of
$(\pi_1Y)^\Sigma$; the quotient is $(\pi_1C')^\Sigma$.

Let $N$ be the closed normal subgroup of this fibre generated by
inertias at the points of $a^{-1}(P)$ other than $P'$. Monodromy
preserves these individual puncture labels. Thus the discrete total
group normalizes $N$, and so does its completion, since the
normalizer of a closed subgroup is closed. Filling the nonidentity
graphs in $Y$ kills exactly the closed total normal subgroup generated
by these same meridians. This elementary assertion follows by putting
a transverse disk through each smooth divisor and applying van
Kampen; every meridian along that divisor is a total conjugate of
its meridian in the chosen fibre. Since $N$ is already total-normal,
that total filling kernel equals $N$.

It follows that the quotient of the completed total group is
$(\pi_1((C'\times C')-\operatorname{diag}))^\Sigma$, and that its
injected fibre is $(\pi_1(C'-P'))^\Sigma$. Here we used right
exactness of pro-$\Sigma$ completion for both fillings; it follows
from the finite-quotient universal property exactly as in
\statementref{A.3}. This proves the assertion without assuming
that arbitrary completion commutes with kernels.
\end{proof}
\paragraph{Check record.} Also checked in this writing, 2026-11-18: GPT-6 Astra Pro; Pro mode; separate effort setting not exposed; ANABELIAN\_A1\_SUCCESSOR/configuration. Open subgroup pro-Sigma topology and filling-meridian argument checked; normal cores and total normalization justified.


\begin{claim}{A.9}{Proved}{A finite symplectic monodromy obstruction}
Let $C/\mathbb C$ have genus $h\geq2$, and let $r,\ell$ be distinct
primes. There is a connected degree-$r^3$ cover of $C-P$, with
$r$-group Galois group and cusp inertia of order $r$, whose compact
homology modulo $\ell$ has a subspace of dimension $2h-2$ on which
a subgroup of the configuration group preserving the cover acts
onto $\operatorname{Sp}_{2h-2}(\F_\ell)$. Consequently, if $r,\ell\in\Sigma$
and the full pro-$\Sigma$ configuration fibre were injective, every
prime divisor of $|\operatorname{Sp}_{2h-2}(\F_\ell)|$ would belong
to $\Sigma$.
\end{claim}
\begin{proof}
Choose a one-holed torus containing $P$ and a complementary connected
subsurface $M$ of genus $h-1$ with one boundary component. Map the
free group of $C-P$ onto $\operatorname{UT}_3(\F_r)$ by sending
the first two handle generators to its standard generators and all
generators supported on $M$ to one. Its boundary commutator is a
nontrivial central element $z$ of order $r$ (up to inversion).
Let $D^\circ\to C-P$ be the resulting regular cover and $D$ its
compactification. The restriction over $M$ is the disjoint union of
$r^3$ copies $M_s$, labelled by deck elements $s$. For a fixed $s$,
put
\[
 W=\{u_s-u_{zs}:u\in H_1(M,\F_\ell)\}\subset H_1(D,\F_\ell).
\]
The map defining $W$ is injective. Indeed, the usual symplectic handle
cycles in the disjoint subsurfaces $M_s$ have block-diagonal
intersection pairings with their respective dual cycles, so all their
classes are independent in $D$. This argument also works when
$\ell=2$, where the restricted pairing on $W$ itself can be
 degenerate. We use its identified vector space, not a false claim
that this restricted pairing is nondegenerate.

We give the topological calculation of the action explicitly. Let
$\alpha$ be an oriented simple nonseparating curve in $M$, joined to
$P$ by an embedded arc avoiding a fixed first-strand basepoint.
Pushing $P$ once round this based simple loop is the product of two
opposite Dehn twists on its left and right push-offs, denoted
$\alpha_0,\alpha_1$. To verify the identity, identify an annular
neighbourhood of the based loop with angular coordinate $\theta$ and
transverse coordinate $t$, and move the marked point by
$\theta\mapsto\theta+2\pi u$ during the isotopy parameter $u$.
At its two boundary collars the resulting map is one positive and
one negative twist. It is the identity outside this neighbourhood.
The two push-offs and the puncture bound a pair of pants. Their
cover monodromies are respectively $1,z$ (possibly after replacing
$z$ by $z^{-1}$); centrality makes the joining path immaterial.

The $r$-th power of each twist lifts and fixes the chosen sheet
outside the twisting annuli. On an annular cover of degree $n$,
where $n\mid r$, a twist of exponent $r$ lifts to exponent $r/n$:
in coordinates $\theta_{\rm base}=n\theta_{\rm cover}$ the lift is
$\theta_{\rm cover}\mapsto\theta_{\rm cover}+2\pi(r/n)f(t)$.
It is the identity on the boundary circles, so it glues to the
identity outside all annuli. This proves existence of the lift,
its sheet convention, and preservation of the cover without an
unproved assertion about outer actions.

The lifts of $\alpha_0$ have degree one and are the curves
$\alpha_s$ in the separate copies of $M$. Their twists have
exponent $r$. The lifts of $\alpha_1$ have degree $r$, one for
each central $\langle z\rangle$-orbit, and their twists have
exponent one. Each such long lift has homology class
\[
                        \sum_{j=0}^{r-1}[\alpha_{z^j s}]
                                                               \tag{CF.3}
\]
up to its common orientation. For this equality, lift the pair of
pants: each connected component is a degree-$r$ cyclic cover with
$r$ short boundaries, one long boundary, and one boundary above the
puncture. After compactification the last boundary bounds a disk.
The oriented boundary relation gives exactly (CF.3), with no extra
homology component.

The intersection of (CF.3) with $u_s-u_{zs}$ is zero: the two
coefficients are equal and cancel. Thus the long-lift twists act
trivially on $W$. The short-lift twists send
\[
 u_s-u_{zs}\longmapsto
 u_s-u_{zs}\ \mathbin{\pm}\ r\langle u,\alpha\rangle
                          (\alpha_s-\alpha_{zs}).              \tag{CF.4}
\]
The whole lifted point-push therefore preserves $W$ and induces the
usual symplectic transvection of exponent $\pm r$ on its identified
space $H_1(M,\F_\ell)$. All these lifts have the same sheet labels,
because they fix every point above the chosen basepoint outside the
annuli.

Since $r$ is invertible modulo $\ell$, powers give the transvection
of any coefficient. Every primitive integral homology class of $M$
is represented by a simple nonseparating curve: Euclidean handle
slides reduce a primitive vector to a handle basis vector, and
reversing the slides supplies the curve. Every nonzero vector modulo
$\ell$ has such a primitive integral lift (adjust one coordinate
by multiples of $\ell$ if necessary). We therefore obtain all
symplectic transvections on $H_1(M,\F_\ell)$.

These generate its full symplectic group. Here is a linear-algebra
proof, including characteristic two. A transvection of suitable
coefficient about $v-u$ sends $u$ to $v$ whenever
$\langle u,v\rangle\ne0$. If the pairing is zero, an intermediate
vector with both pairings nonzero exists, since two proper
hyperplanes do not exhaust the vector space. Thus the transvection
group is transitive on nonzero vectors. After fixing the first vector
$e$ of a symplectic basis, write its partner as $f+w+ae$, with $w$
in the symplectic complement. The product
$T_{e+w}T_w^{-1}T_e^{-1}$ fixes $e$ and removes $w$, with signs
chosen according to $\langle e,f\rangle=1$; a power of $T_e$
then removes $ae$. Induction in the symplectic complement fixes the
whole basis. This proves generation of $\operatorname{Sp}_{2h-2}$.

To obtain the last assertion, suppose the completed fibre is a
closed subgroup of the pro-$\Sigma$ configuration group. The
$r^3$-index subgroup defining $D^\circ$ has its full pro-$\Sigma$
topology there by the normal-core argument. Its quotient
$H_1(D,\F_\ell)$ is finite. The chosen point-push powers normalize
both the covering subgroup and the closed kernel defining this
homology quotient. Their closures still do so, because normalizers
of closed subgroups are closed. Conjugation on a finite quotient is
continuous, as one checks on finitely many coset representatives.
It therefore has finite $\Sigma$-image. Restricting to $W$ gives
a quotient of that image, but the preceding calculation makes this
quotient the full $\operatorname{Sp}_{2h-2}(\F_\ell)$. Its prime
divisors must lie in $\Sigma$.
\end{proof}
\paragraph{Check record.} Also checked in this writing, 2026-11-18: GPT-6 Astra Pro; Pro mode; separate effort setting not exposed; ANABELIAN\_A1\_SUCCESSOR/configuration. Lifted point-push transvections and long-boundary cancellation checked, including ell=2 and r=2 cases; no nondegenerate restricted pairing assumed.


\begin{claim}{A.10}{Proved}{Classification of prime sets for the full fibre sequence}
For a smooth proper complex curve of genus at least two, the full
pro-$\Sigma$ configuration fibre sequence is exact if and only if
$\Sigma$ is empty, consists of one prime, or is the set of all
primes. If $C$ instead has characteristic $p$ and
$\Sigma\subseteq\{\ell:\ell\ne p\}$, it is exact if and only if
$|\Sigma|\leq1$. In particular the full prime-to-$p$ fibre map is
noninjective for every $p$ and every such curve, as claimed in R4.
The relative cuspidally pro-$r$ sequence of
\statementref{A.3} remains exact for every $r\in\Sigma$.
\end{claim}
\begin{proof}
The empty case is trivial. For $\Sigma=\{r\}$, monodromy on
$H_1(C-P,\F_r)=H_1(C,\F_r)$ is trivial, because the compactified
family is the product and there is only one boundary class, which
vanishes in abelian homology. Apply \statementref{A.2} and right
exactness. For the set of all primes, we use
\citet[Proposition 2.2]{LA_MT2008}, with a characteristic-zero
hyperbolic curve and the projection from its second to its first
configuration space: the proposition gives the full profinite fibre
exact sequence. Its prime-set hypothesis allows the total set (or a
singleton); it does not assert the sequence for a proper multi-prime
set.

Suppose now that $\Sigma$ is proper and contains distinct primes
$r,\ell$, and choose $p\notin\Sigma$. If the fibre sequence were
exact for $C$, it would be exact for every connected cyclic proper
cover of degree $r^m$, by \statementref{A.8}. Such covers exist by
the surjection $H_1(C,\Z)\twoheadrightarrow\Z/r^m$ and have genus
$h=1+r^m(g-1)$, hence arbitrarily large genus. Let $d$ be the
multiplicative order of $\ell$ modulo $p$ (for $p=2$, take $d=1$).
Choose $m$ such that $h-1\geq d$. Then $p$ divides
$\ell^{2d}-1$, and therefore divides
\[
 |\operatorname{Sp}_{2h-2}(\F_\ell)|
       =\ell^{(h-1)^2}\prod_{i=1}^{h-1}(\ell^{2i}-1).
\]
For reference, this order formula follows inductively by choosing a
nonzero first vector, its partner of pairing one, and a symplectic
basis of the complement; the successive factors are
$(\ell^{2i}-1)\ell^{2i-1}$. The divisibility contradicts
\statementref{A.9}, since $p\notin\Sigma$. Thus left injectivity
fails. Middle and right exactness still hold by the universal
finite-quotient argument, so this identifies the precise lost part
of the sequence.

For characteristic $p$, the functorial proper-pair specialization in
\statementref{A.4} identifies the pro-$\Sigma$ fibre map with the
complex one. As $p\notin\Sigma$, a nonempty such set is always
proper, giving exactly the stated classification. The mechanism also
explains the relative repair: a single retained cuspidal prime makes
all finite covering monodromy kernels used in
\statementref{A.3} $r$-groups; keeping two independent cusp primes
permits the symplectic finite actions above and their omitted-prime
divisors.
\end{proof}
\paragraph{Check record.} Also checked in this writing, 2026-11-18: GPT-6 Astra Pro; Pro mode; separate effort setting not exposed; ANABELIAN\_A1\_SUCCESSOR/configuration. Prime-set classification checked; MT2008 PT-formation restriction verified from opened source; characteristic-p scope excludes p.


\begin{claim}{A.11}{Proved}{The labelled several-section extension}
Let $C_0,B/\mathbb C$ be smooth proper curves, with $B$ of genus at
least one and $C_0$ of genus at least two, and let
$f_1,\ldots,f_m:B\to C_0$ have pairwise disjoint graphs. For
$Y=(C_0\times B)-\bigcup_i\operatorname{graph}(f_i)$, the relative
cuspidally pro-$r$ fibre sequence over pro-$\Sigma$ proper groups
is exact for every $r\in\Sigma$. The same assertion holds for simultaneously liftable
smooth proper pairs in characteristic prime to every member of
$\Sigma$ by the preceding specialization. In particular it applies
to the diagonal families after finite proper covers used in
cuspidally abelian reconstruction.
\end{claim}
\begin{proof}
For a finite fibre quotient, take its proper quotient $A$ and the
associated Galois cover $C_1\to C_0$. Pull this cover back by every
$f_i$, and take a common finite Galois $\Sigma$-cover $B_1\to B$
which trivializes all these $A$-torsors. Such a cover is given by
the intersection of their kernels; its group is a subgroup of a
finite product of copies of $A$. Now the inverse boundary in
$C_1\times B_1$ consists of disjoint labelled sections. Its compact
family is a product, its base is aspherical, and the homology action
modulo $r$ is the same shear group used in
\statementref{A.3}. Apply \statementref{A.2} and take the
normal core in $\pi_1(Y)$. The resulting total finite quotient is
an $r$-group by $A\times A_B$, where $A_B$ is the chosen base-cover
group; this is a $\Sigma$-group. Every finite fibre quotient is
detected, and the middle and right exactness proof is unchanged.
For the specific characteristic-$p$ diagonal-cover families, lift
the proper etale covers together with the smooth curve (finite
prime-to-$p$ etale covers lift uniquely on the proper-pair comparison)
and specialize. For a general characteristic-$p$ family of graphs,
the same conclusion requires a simultaneous lift of the indicated
smooth proper pair; the statement is understood with that hypothesis,
not as an assertion that arbitrary morphisms of curves lift.
\end{proof}
\paragraph{Check record.} Also checked in this writing, 2026-11-18: GPT-6 Astra Pro; Pro mode; separate effort setting not exposed; ANABELIAN\_A1\_SUCCESSOR/configuration. Labelled-section cover/shear/core proof checked. m>=1 free-fibre case explicit; m=0 is the product case. Direct positive-characteristic constant-family version now follows from A.45–A.46.


\begin{claim}{A.12}{Proved}{A finite two-prime test beyond the cofinite case}
For a genus-three curve, the full pro-$\{2,3\}$ configuration fibre
map is noninjective. A particular obstructing fibre quotient has order
$8\cdot3^{41}$; an exact computation of the preserving braid action
on a four-dimensional homology subspace gives
$\operatorname{Sp}_4(\F_3)$, of order $51840$, and an explicit
acting element of order five.
\end{claim}
\begin{proof}
Let $F_6$ have handle generators $a,b,c,d,e,f$, and let $K$ be the
kernel of its quotient onto $\operatorname{UT}_3(\F_2)$ obtained
from $a,b$ alone. Then $K$ is free of rank $1+8(6-1)=41$.
Use $N=[K,K]K^3$. The actual point-push lifts
$C,D,E,F$ are the $\sigma$-conjugates of $c,d,e,f$ in Bellingeri's
genus-three presentation. The program
\path{code/configuration_symplectic_crosscheck.py} derives their
conjugation words by the same relations as (CF.1), squares the words,
and checks preservation of the quotient to $\operatorname{UT}_3(\F_2)$.
It computes their action on the full integral Schreier homology and
on the submodule spanned by the four differences of the $c,d,e,f$
loops at the identity and at the nontrivial central vertex.
The five actions of $C^2,D^2,E^2,F^2,(CE)^2$ are exactly the
symplectic transvections of coefficient two, including the
transvection in direction $c+e$. Exhaustive multiplication of the
resulting four-by-four matrices modulo three gives $51840$ distinct
matrices. It also checks that
\[
                    (CE)^2 F^2 D^2 D^2 C^2
\]
acts on this subspace by a matrix of exact order five.
The matrices, group order, and word are retained in the JSON
certificate and exact output alongside the program. Any finite
$\{2,3\}$ total quotient dominating $F_6/N$ would make this action
factor through a finite $\{2,3\}$ group, contradicting that order.
The normalizer argument is exactly the one in
\statementref{A.7}. This finite computation verifies an instance
of the higher-dimensional mechanism in \statementref{A.9}; the
general proof of that statement does not rely on the computation.
\end{proof}
\paragraph{Check record.} Also checked in this writing, 2026-11-18: GPT-6 Astra Pro; Pro mode; separate effort setting not exposed; ANABELIAN\_A1\_SUCCESSOR/configuration. Finite symplectic obstruction logic checked; numerical certificates independently rerun by finite\_checks agent, not claimed as this checker computation. | Also checked in this writing, 2026-11-18: GPT-6 Astra Pro; finite\_checks agent; effort not separately exposed. Read full source proof and finite normalizer contradiction; original fullSchreier program rerun; independent transvection enumeration; geometric point-push derivation not independently reproved here.


\subsection{Direct characteristic-positive proof, including proper wild covers}

\begin{claim}{A.13}{Proved}{Relative cuspidal exactness for arbitrary prime sets containing a non-residue prime}
Let \(C\) be a smooth proper geometrically connected curve of genus at
least two over an algebraically closed field of characteristic \(p>0\).
Let \(\Sigma\) be any set of primes and choose
\(r\in\Sigma\setminus\{p\}\). In the definitions of A.3, now formed
with maximal pro-\(\Sigma\) groups even when \(p\in\Sigma\), the
relative cuspidally pro-\(r\) configuration sequence is exact. The
boundary comparison and arithmetic transport of A.5 remain valid.
More generally, the finite-detection proof works for a constant proper
compactification with disjoint labelled sections in characteristic
\(p\), without a simultaneous lifting hypothesis.
\end{claim}
\begin{proof}
Let $C/k$ be smooth proper of genus at least two over an algebraically
closed field of characteristic $p>0$. Let $\Sigma$ be arbitrary and
$r\in\Sigma\setminus\{p\}$. Start with a finite quotient $Q$ of
$F_{\mathrm{cr}}$, let $R$ be its normal $r$-subgroup given by the
image of $K_F$, and let $A=Q/R$. The finite quotient $A$ of
$\Delta=\pi_1(C)^\Sigma$ defines a connected finite \'etale proper
Galois cover $C'\to C$. There is no restriction on the $p$-part of
$|A|$. Pull back the two coordinates to obtain
\[
 V=(C'\times C')-\coprod_{a\in A}\Gamma_a.
\]
This is a connected finite \'etale Galois cover of $U$ with group
$A\times A$; its second projection has constant compactification and
the disjoint labelled sections $y\mapsto(ay,y)$. Its fibre is the
cover of $C-P$ corresponding to the kernel of its proper quotient $A$.
The restriction of $Q$ to that fibre has $r$-group image. A.45
dominates it by a finite $r$-quotient of $H=\pi_1(V)$, with kernel
$J$. Write $G=\pi_1(U)$ and
$J_0=\bigcap_{g\in G/H}gJg^{-1}$. Since $J$ is normal in $H$,
this finite intersection is normal in $G$, and $H/J_0$ embeds
in the product of the conjugate finite $r$-quotients $H/gJg^{-1}$.
Consequently $T=G/J_0$ fits into
an exact sequence
\[
 1\to R_T\to T\to A\times A\to1,
 \qquad R_T\text{ a finite }r\text{-group}.
\]
Every prime divisor of $|T|$ belongs to $\Sigma$, so $G\to T$
factors through $E=G^\Sigma$. Its map to $A\times A$ factors
through the proper product $\Delta\times\Delta$; hence the image
of $K_E$ lies in $R_T=H/J_0$, and the map factors through
$E_{\mathrm{cr}}$. If a fibre element maps trivially to $T$,
its proper image in $A$ is trivial, so it belongs to the fibre
of $V\to C'$; there it is killed by the chosen quotient of $H$
and therefore by $Q$. Thus the fibre restriction dominates $Q$.
This proves injectivity of
$F_{\mathrm{cr}}\to E_{\mathrm{cr}}$ for every such $\Sigma$.

We spell out the middle exactness, which cannot be inferred from an
unrestricted homotopy sequence for an open morphism in characteristic
$p$. Every inertia group along the diagonal has trivial image in the
proper product. Its image in every finite quotient of
$E_{\mathrm{cr}}$ consequently lies in the image of $K_E$, an
$r$-group. Since $r\ne p$, the corresponding cover is tame along the
diagonal. A.46 applies. If $T$ is such a finite quotient and $I$ its
fibre image, $T/I$ comes from the second proper curve. It is a finite
$\Sigma$-group, so this factorization is through $\Delta$.
Thus every element of $\ker(E_{\mathrm{cr}}\to\Delta)$ has image in
$I$ in every finite quotient, and hence belongs to the closed image
of $F_{\mathrm{cr}}$. The reverse inclusion is immediate. Projection
to $\Delta$ is onto, and we obtain
\[
 1\to F_{\mathrm{cr}}\to E_{\mathrm{cr}}
 \overset{\operatorname{pr}_2}{\longrightarrow}\Delta\to1.
\]
This argument proves both the originally asserted relative exactness
and the full finite-domination assertion with quotient $A\times A$.
It uses neither Stix's theorem nor a lift to characteristic zero.
Stix's Proposition 2.7 remains a verified alternative source for the
single-prime relative fibre step, not a dependency of this proof.

For the constant-compactification variant, choose the proper fibre
cover for a prescribed finite fibre quotient, and pull it back along
each of the finitely many section maps. A common proper Galois base
cover trivializing those torsors has group a subgroup of a finite
product of their $\Sigma$-deck groups. On the resulting product the
inverse boundary is disjoint and labelled. A.45 and the same finite
core produce an $r$-group by the product of the two proper deck
groups, and A.46 gives middle and right exactness. Finally A.48
below proves the boundary and arithmetic transport clause of A.13.
That proof derives its own geometric fibre input from A.45--A.46,
so this reference creates no circular dependency.


\end{proof}
\paragraph{Check record.} Also checked in this writing, 2026-11-18: GPT-6 Astra Pro; Pro mode; separate effort setting not exposed; ANABELIAN\_A1\_SUCCESSOR/configuration. Omitted characteristic-p middle/right exactness justification identified and repaired by A.46; independent direct finite-r detection and transport supplied in A.45–A.48. Statement remains Proved with expanded proof. | Also checked in this writing, 2026-11-18: GPT-6 Astra Pro; Pro mode; separate effort setting not exposed; ANABELIAN\_A1\_SUCCESSOR/configuration. Direct characteristic-p proof through nonabelian central torsor induction and tame normalization; p in Sigma allowed, no curve lifting or excluded SGA specialization input.



\subsection{The cuspidal criterion and its arithmetic inputs}

We retain the whole arithmetic Galois quotient throughout this subsection.

The manuscript question addressed is cote~149, scan~36: can the covering condition force the equality of the two points?

\begin{claim}{A.14}{Proved}{The same-prime quotient lemma}

Let $F$ be profinite, $K\triangleleft F$ closed, $r$ prime, and
$N=\ker(K\to K^{(r)})$. For every open $I\le F$ containing $N$, the
continuous quotient map induces a natural isomorphism
$I^{(r)}\simeq(I/N)^{(r)}$. No restriction on $[F:I]$ is imposed.

\end{claim}

\begin{proof}

Let F be a profinite group, K a closed normal subgroup, r a prime, and N the kernel of the maximal pro-r quotient \(K{\to}K^{(r)}\). Because N is characteristic in K, it is normal in F. Let \(I{\le}F\) be open and contain N. We prove that the continuous epimorphism \(I{\to}I/N\) induces an isomorphism \(I^{(r)}\to(I/N)^{(r)}\).

Take an arbitrary continuous quotient I/J which is a finite r-group, where J is open and normal in I. Put \(I_0=\operatorname{core}_F(I)\); it is open normal in F, contained in I, and contains N. Put \(J_0=J{\cap}I_0\). Then \(J_0\) is open normal in \(I_0\) and \(I_0/J_0\hookrightarrow I/J\), so \(I_0/J_0\) is an r-group. Take \(M=\operatorname{core}_F(J_0)\). There are only finitely many F-conjugates of \(J_0\), because \(J_0\) is open. Since \(I_0\) is normal in F, \(I_0/M\) embeds in the product of the quotients \(I_0/{}^fJ_0\), each isomorphic to \(I_0/J_0\). Thus \(I_0/M\) is an r-group.

The image of K in \(F/I_0\) is an r-group: \(N{\subset}I_0\), so this image is a finite quotient of K/N, a pro-r group. The image of K in F/M is an extension of that image by a subgroup of \(I_0/M\). It too is an r-group. By the universal defining property of N, N is killed by \(K{\to}F/M\). Therefore \(N\subset M\subset J_0\subset J\).

Every finite r-quotient of I consequently kills N. This is exactly the asserted isomorphism of maximal pro-r quotients, with its natural map, not merely an abstract isomorphism. No finite-generation assumption, normality of I, or r-power-index assumption on I was used.


\end{proof}
\paragraph{Check record.} Also checked in this writing, 2026-11-18: GPT-6 Astra (OpenAI, ChatGPT Pro mode), root; separate effort setting not exposed. Read the full finite-core argument for the same-prime topology and the integral Frobenius/cusp-count comparison. No mixed-prime full-fibre injection is used.


\begin{claim}{A.15}{Proved}{The full compactified pro-$r$ group at every open level}
Let $C/k$ be a smooth proper geometrically connected curve over a
finite field of characteristic $p$, let $P\in C(k)$, and let
$r\in\Sigma\setminus\{p\}$, with no other restriction on $\Sigma$.
Put $U=C-P$, $F=\pi_1(U_{\bar k})^\Sigma$,
$\Delta=\pi_1(C_{\bar k})^\Sigma$, $K=\ker(F\to\Delta)$, and
$N=\ker(K\to K^{(r)})$. Let $G_P$ be the arithmetic relative
cuspidally pro-$r$ quotient, with geometric kernel $F/N$ and full
arithmetic quotient $G_k$.

For every open geometric subgroup $I/N\le F/N$, let $V$ be its
actual connected geometric cover of $U_{\bar k}$ and let $\bar V$
be its smooth compactification. If $T_I$ is the closed normal
subgroup of $I/N$ generated by all its cusp inertia images, then
the actual covering and compactification maps induce an isomorphism
\[
             ((I/N)/T_I)^{(r)}\simeq\pi_1(\bar V)^{(r)}.
\]
This isomorphism is compatible with every arithmetic outer action
preserving the cover. Consequently, a cusp-preserving isomorphism
of such arithmetic relative groups over $G_k$ preserves the
integral Frobenius polynomial and the point counts of every pair
of corresponding open covers, over every finite constant extension.
In particular these assertions allow $p\in\Sigma$.
\end{claim}

\begin{proof}
Write $\Gamma=\pi_1(U_{\bar k})$ before taking a prime-set quotient.
We first justify the two geometric facts needed at an open level.

\paragraph{Actual covers, tameness, and the induced $\Sigma$ topology.}
The inverse image $\widetilde I$ of $I$ in $\Gamma$ is open,
and hence defines an actual connected finite \emph{\'etale} cover
$V\to U_{\bar k}$. Its Galois closure has group a finite quotient
of $F/N$. The image of every cusp inertia subgroup lies in the
image of $K/N$, a pro-$r$ group. Its finite inertia groups
therefore have $r$-power order and are tame because $r\ne p$.
Thus the cover extends by normalization to a finite tame morphism
of smooth proper curves, and its boundary is the finite set
$\bar V-V$. The same reasoning applies to arithmetic covers after
a finite extension of constants. This tameness assertion places
no restriction on the $p$-part of their proper deck groups.

The natural map $\widetilde I^\Sigma\to I$ is an isomorphism.
Here is the full normal-core argument, including $p\in\Sigma$.
The subgroup $B=\operatorname{core}_{\Gamma}(\widetilde I)$ is
open normal, and $\Gamma/B$ is a finite $\Sigma$-group, since
$\widetilde I$ is pulled back from the pro-$\Sigma$ group $F$.
Given any open normal $J\triangleleft\widetilde I$ with
$\widetilde I/J$ a finite $\Sigma$-group, put $J'=J\cap B$ and
\[
                 J_0=\bigcap_{g\in\Gamma/B}gJ'g^{-1}.
\]
Then $J'$ is normal in $B$, $B/J'$ is a finite $\Sigma$-group,
and $J_0$ is open normal in $\Gamma$. The group $B/J_0$ embeds
in the product of the conjugates of $B/J'$. Hence both $B/J_0$
and $\Gamma/J_0$ are finite $\Sigma$-groups, and $J_0\subseteq J$.
Every finite $\Sigma$-quotient of $\widetilde I$ is therefore
detected by a finite $\Sigma$-quotient of $\Gamma$. This proves
the asserted induced topology and the isomorphism. It uses only
that finite $\Sigma$-groups are closed under subgroups, finite
products, and extensions.

\paragraph{Killing the actual boundary at the retained prime.}
Let $T\triangleleft I$ be the closed normal subgroup generated by
the images of the actual inertia subgroups of $\widetilde I$ at
the points of $\bar V-V$. A finite \emph{\'etale} cover of $V$
extends \emph{\'etale} over $\bar V$ exactly when each of those
inertia groups acts trivially. Indeed its normalization is finite
over the smooth curve, and the criterion is the unramified
criterion at each discrete valuation ring. The universal property
of pro-$\Sigma$ completion therefore gives
\[
                       I/T\simeq\pi_1(\bar V)^\Sigma.
\]
This is right exactness of completion, and uses no left-exactness
claim for completion.

The same-prime quotient lemma A.14 applies to $N\subset I$:
every continuous homomorphism from $I$ to a finite $r$-group
kills $N$. In particular the finite $r$-quotients of $I$ killing
$T$ are exactly those killing the closed product $TN$.
The image of $T$ in $I/N$ is precisely $T_I$. Therefore
\[
 \begin{aligned}
 ((I/N)/T_I)^{(r)}
   &\simeq (I/TN)^{(r)}\\
   &\simeq (I/T)^{(r)}\\
   &\simeq (\pi_1(\bar V)^\Sigma)^{(r)}
    =\pi_1(\bar V)^{(r)}.
 \end{aligned}
\]
In the final equality we used $r\in\Sigma$. Each map is induced
by an actual geometric quotient and is consequently natural.
There is no assumption that $N\subseteq T$, or that $[F:I]$
is an $r$-power.

\paragraph{Arithmetic compatibility and Frobenius.}
Let $L/k$ be finite and let $H\le(G_P)_L$ be an open subgroup
whose arithmetic projection is the whole $G_L$. Its actual
connected cover $V_H\to U_L$ is geometrically connected: the
image of its arithmetic fundamental group in $G_L$ is all of
$G_L$, so its field of constants is $L$. Its geometric subgroup
is an $I/N$ as above. Arithmetic conjugation by $H$ permutes
the actual boundary inertia groups and preserves every kernel
in the displayed construction. The isomorphism is thus equivariant
for the outer $G_L$-action on the geometric compactification group.
It identifies, in particular,
\[
 H^1(\bar V_{H,\bar L},\mathbf Z_r)
   =\operatorname{Hom}_{\mathrm{cont}}
        (\pi_1(\bar V_{H,\bar L})^{(r)},\mathbf Z_r)
\]
with its actual continuous Frobenius action.

An isomorphism of the original arithmetic relative groups over
$G_k$, preserving the cusp decomposition conjugacy class, identifies
the groups belonging to corresponding open subgroups $H$.
It also identifies their systems of cusp inertia. Their compactified
$r$-adic $H^1$ representations are therefore isomorphic over
$G_L$. For the geometric Frobenius $\operatorname{Fr}_L$, the
smooth proper curve trace formula and integrality theorem give
\[
 P_{\bar V_H}(t)=\det(1-t\operatorname{Fr}_L\mid
       H^1(\bar V_{H,\bar L},\mathbf Q_r))\in\mathbf Z[t],
 \qquad
 \#\bar V_H(L_m)=|L|^m+1-
       \operatorname{Tr}(\operatorname{Fr}_L^m\mid H^1),
\]
where $L_m/L$ is the degree-$m$ extension. These are
\citet[formula (1.5.1) and Theorem (1.6)]{WeilCurves}; their
hypotheses apply to every smooth proper geometrically connected
compactification, including genus zero or one. Equality of the
$r$-adic characteristic polynomials implies equality of their
integer polynomials because $\mathbf Z[t]\hookrightarrow\mathbf Q_r[t]$.
The compactified point counts consequently agree for every $m$.

\paragraph{Recovering the deleted points.}
Choose one actual cusp decomposition subgroup $D_P$ of $(G_P)_L$.
The boundary points of $\bar V_H$ are indexed by the finite double
coset set $H\backslash(G_P)_L/D_P$. For a representative $g$, the
residue degree of the corresponding boundary point is
\[
       [G_L:\operatorname{pr}(H\cap gD_Pg^{-1})].
\]
This follows directly from the decomposition-group action on the
finite fibre of the cover: the stabilizer is the displayed
intersection and its arithmetic projection is the Galois group
of the residue field. It remains valid in the relative quotient,
because the cover factors through that quotient and the actual
decomposition action has exactly the same image. The isomorphism
over $G_L$ identifies the double cosets and these degrees.
A deleted point of degree $d$ contributes $d$ points over $L_m$
if $d\mid m$, and zero otherwise. Subtracting the equal boundary
counts from the compactified counts proves
$\#V_H(L_m)=\#V_{H'}(L_m)$ for all corresponding covers and all $m$.
For an open subgroup with a smaller arithmetic projection, begin
over its actual finite field of constants; this covers every open
level claimed in the statement.
\end{proof}
\paragraph{Check record.} Also checked in this writing, 2026-11-18: GPT-6 Astra (OpenAI, ChatGPT Pro mode), root; separate effort setting not exposed. Read the full finite-core argument for the same-prime topology and the integral Frobenius/cusp-count comparison. No mixed-prime full-fibre injection is used. | Also checked in this writing, 2026-11-18: GPT-6 Astra (OpenAI, ChatGPT Pro mode), root; separate effort setting not exposed. Read the complete replacement proofs for arbitrary retained prime sets and arbitrary perfect fields, respectively, including the normal-core topology, primitive orientation class and tame-boundary restriction. | Also checked in this writing, 2026-11-18: GPT-6 Astra Pro; Pro mode; separate effort setting not exposed; ANABELIAN\_A1\_SUCCESSOR/configuration. Full replacement for arbitrary Sigma containing r different from p. Explicit core gives induced Sigma topology; A14 kills N in every finite r quotient; compactified full pro-r group and arithmetic action, Frobenius polynomial and boundary double-coset point counts recovered. Checked in this writing; parent review requested.


\begin{claim}{A.16}{Known}{The exact uniform exceptional-set input}

For a proper hyperbolic curve $C/k$ over a finite field of
characteristic $p$, and $\Sigma=\mathrm{Primes}-\{p\}$, there is a
finite set $E_C\subset C^{\cl}$ such that outside its inverse image
a decomposition subgroup has a unique geometric point label, even
up to commensurability, over every finite constant extension at once.
The numbering of Definitions~1.3 and 1.5, Remark~1.2, and
Proposition~1.8(vi) is the same in the attached arXiv text and the
journal version of \cite{LA_ST2009}.

\end{claim}

\begin{proof}

Let k be finite of characteristic p, X/k a smooth proper geometrically connected curve, \(U=X-S\) hyperbolic, and \({\Sigma}\) a set of primes containing at least one prime different from p.~Let \({\Pi}_U\) be the arithmetic group retaining the whole \(G_k\) and the maximal geometric pro-\({\Sigma}\) quotient of the tame group. Its universal pro-cover compactification is denoted by \(\widetilde X_U\).

The attached ST2009 Definition 1.3(iii) calls a point a of a set map f exceptional when \(\#f^{-1}(f(a))>1\). Definition 1.5 applies this to the map \(\overline D\) from closed points upstairs to commensurability classes of decomposition subgroups. Definition 1.10 puts \(E_U\) equal to its image in \(X^{\rm cl}\).

Proposition 1.8(iii) states injectivity of \(\overline D\) over a fixed geometric point; (iv) gives finite fibres over one closed point and injectivity over one rational point; (v) says the exceptional set is \({\Pi}_U\)-stable; under the additional hypothesis that \({\Sigma}\) is cofinite, (vi) states that \(E_{\widetilde U}/\Pi_U\) is finite. This quotient is \(E_U\). The same numbering and hypotheses occur in the opened journal PDF.

Remark 1.2 says that the commensurability maps are unchanged upon passing to any open subgroup. More explicitly, if H is open in \({\Pi}_U\) and \(D_1,D_2\) are closed subgroups, then \(D_1\) and \(D_2\) are commensurate if and only if \(D_1{\cap}H\) and \(D_2{\cap}H\) are commensurate. One direction is immediate. For the other, \(D_i{\cap}H\) has finite index in \(D_i\); the common intersection has finite index in both. Universal pro-cover points are the same on replacing U by its H-cover. In particular, when H is the preimage of \(G_L\) for a finite constant extension L/k, the exceptional points upstairs are exactly the same points and the exceptional image on \(X_L\) is the finite inverse image of \(E_U\).

For \({\Sigma}=\mathrm{Primes}-\{p\}\), all hypotheses are met when C is proper of genus at least two: U=C is hyperbolic and \({\Sigma}\) is nonempty and cofinite. Consequently there is one finite set \(E_C\subset C^{\rm cl}\) whose geometric inverse image excludes all exceptional labels over all finite extensions. If R is a closed point outside \(E_C\), a decomposition subgroup at R cannot be commensurate, even after conjugation, with a subgroup at another geometric point. Conjugation moves the lift and preserves the underlying closed-point label. This is exactly the use in AS-AG.40 Step 3.

\textbf{Check of their proof and its hypotheses.}

Their proof first replaces U by an open-subgroup cover whose compactification has genus \textgreater1 and is nonhyperelliptic. Commensurability invariance just proved makes finiteness upstairs descend under this finite cover. For a nonhyperelliptic compactification X, the difference morphism \(\delta:X\times X\to J_X\), \((P,Q)\mapsto[P-Q]\), is injective on geometric points away from the diagonal. Indeed, equality of two differences gives \(P+Q'\sim P'+Q\); distinct linearly equivalent effective degree-two divisors would supply a degree-two map to the projective line, contrary to nonhyperellipticity. Equality of the divisors then gives the equal ordered pair, unless both differences vanish.

The image W of \({\delta}\) contains no translate of a positive-dimensional abelian subvariety. A two-dimensional such translate would be W itself, hence an abelian variety containing zero and generating \(J_X\). Then \(J_X\) would have dimension two, contrary to nonhyperellipticity. A one-dimensional translate would, by the quasi-finite bijection off zero, admit a finite radicial cover mapping nontrivially to X. Its normalization has genus one (purely inseparable extensions of the elliptic function field preserve genus), so it cannot map nontrivially to a genus\textgreater1 curve; after factoring off inseparability, the separable part contradicts Riemann--Hurwitz.

An exceptional pair has zero abelian Kummer difference after a finite constant extension. Therefore its nonzero difference class lies in \(J_X(\bar k)\{\Sigma'\}\), the torsion supported on the complementary primes. Here \({\Sigma}'\) is finite. The opened proof applies Boxall's finite-prime torsion intersection theorem to W and concludes that its intersection with that subgroup is finite. Injectivity of \({\delta}\) away from the diagonal then gives finitely many exceptional geometric pairs on X. For the prime-to-p choice the complementary prime set is exactly \{p\}, with no missing restriction to odd p, p-rank, ordinarity, or nonhyperellipticity of the original curve. The cover reduction is what handles genus two and hyperelliptic original curves.

We have checked the derivation and the hypotheses of the stated ST theorem. We have not independently reproved Boxall's theorem or the nonhyperelliptic-cover existence result cited within ST's proof.


\end{proof}
\paragraph{Check record.} Also checked in this writing, 2026-11-18: GPT-6 Astra (OpenAI, ChatGPT Pro mode), root; separate effort setting not exposed. Opened supplied Saidi--Tamagawa source, Proposition 1.8(vi), definitions, and finite support argument; cofinite hypothesis retained. Known result, not independently reproved.


\begin{claim}{A.17}{Proved}{A rational point section is self-normalizing}

Let $C/L$ be smooth proper geometrically connected of genus at least
two over a finite field, and let $\Sigma$ contain a prime
$r\ne\operatorname{char}L$. If $D\subset\Pi_C^\Sigma$
is a continuous section subgroup mapping isomorphically onto $G_L$,
then $N_{\Pi_C^\Sigma}(D)=D$.

\end{claim}

\begin{proof}

The exact ST2009 Lemma 1.7 hypotheses are: \(Z{\subset}{\Pi}_U\) is closed, pr(Z) is open in \(G_k\), and \(pr|_Z\) is injective. Its conclusion is that pr embeds the commensurator in \(G_k\); commensurator, normalizer, and centralizer coincide and contain Z with finite index. If Z is the section of an L-rational point and we work in \({\Pi}_{U_L}\), its projection is the whole \(G_L\). Therefore the conclusion specializes to

\[N_{\Pi_{U_L}}(Z)=Z_{\Pi_{U_L}}(Z)=Z.\]

This does not assert self-normalization of every closed-point subgroup in the original \({\Pi}_U\) without enlarging constants. It is exactly sufficient for AS-AG.40.

For clarity, the key proof of the lemma can be re-derived. A geometric commensurator \({\sigma}\) centralizes the open intersection of Z and \({\sigma}Z{\sigma}^{-1}\), because both project injectively to the same abelian arithmetic group. If \({\sigma}{\ne}1\), choose an open characteristic \(N{\trianglelefteq}{\Delta}_U\) not containing \({\sigma}\). The open subgroup \(\bar H=\langle N,\sigma\rangle\) has a nonzero image of \({\sigma}\) in its abelianization, detected by the cyclic quotient \(\bar H/N\). The open part of Z centralizing \({\sigma}\) stabilizes \(\bar H\), so this abelian image is fixed by the Frobenius action on the corresponding genuine cover. The abelian geometric group is a product of prime-adic Tate modules of the generalized Jacobian (proper Jacobian for proper covers); none has a nonzero Frobenius-fixed element because its eigenvalues have nonzero Weil weights, and its Tate module is torsion-free. This contradiction proves that the geometric part of the commensurator is trivial. Its injection into the procyclic arithmetic group gives the rest.

The opened source used in the first argument is
\cite[Lemma 1.7]{LA_ST2009}; it has no cofiniteness assumption on \(\Sigma\).


\end{proof}
\paragraph{Check record.} Also checked in this writing, 2026-11-18: GPT-6 Astra (OpenAI, ChatGPT Pro mode), root; separate effort setting not exposed. Checked the self-normalizer proof, both compactness steps in point recognition, and commensurability-to-finite-constant-extension equivalence. Opened ST2009 Lemma 1.7 and Proposition 1.8(iii).


\begin{claim}{A.18}{Proved}{Point recognition in the whole cuspidal quotient}

Under the hypotheses of AS-AG.40, a cusp-preserving isomorphism
$\beta:G_P\simeq G_Q$ over $\mathrm{id}_{\Pi_C}$ carries every
noncuspidal point section, after its finite residue-field extension,
to an actual noncuspidal point section. Outside $E_C\cup\{P,Q\}$ the
target point is the same point of $C$. This conclusion holds in the
whole quotient, and not merely at a finite level.

\end{claim}

\begin{proof}

Suppose \(\beta:G_P\xrightarrow\sim G_Q\) is the isomorphism in AS-AG.40, over \(\operatorname{id}_{{\Pi}_C}\), preserving the cusp decomposition conjugacy class. For a finite extension L/k, restrict both groups to the preimage of \(G_L\); \({\beta}\) remains an isomorphism over \(\operatorname{id}_{{\Pi}_{C_L}}\) and carries the base-changed cusps. Choose \(R{\in}(C-P)(L)\) and a decomposition subgroup \(D_R^P\). Let \(S={\beta}(D_R^P)\), a closed subgroup mapping isomorphically onto \(G_L\).

Let \(\mathcal H\) be the directed set of all open subgroups \(H{\le}(G_Q)_L\) containing S, ordered by reverse inclusion. All these subgroups have full arithmetic projection, so their connected covers \(X_H\) of \(C_L-Q\) are geometrically connected over L. The source cover for \({\beta}^{-1}(H)\) has a rational point above R because it contains \(D_R^P\). By Statement \statementref{A.15}'s point-count equality, \(X_H(L)\) is a nonempty finite set.

For \(H'{\subset}H\) the covering map gives a map \(X_{H'}(L){\to}X_H(L)\). The inverse limit of these sets is nonempty even though no individual transition map is asserted surjective. To prove this, take their product, a compact space. Compatibility along each transition defines a closed condition. Any finite set of compatibility conditions is satisfied by choosing a point on the cover for the intersection of the finitely many subgroups involved and projecting it to the relevant covers. Thus the closed conditions have the finite-intersection property, and compactness supplies a compatible family \((x_H)\).

The component at \(H=(G_Q)_L\) is a point \(T{\in}(C-Q)(L)\). Choose a geometric lift \(\widetilde T\) over T in the universal pro-cover with every projection equal to \(x_H\). This is possible by the same compactness argument in the profinite fibre above T. Its decomposition subgroup \(D_{\widetilde T}\) maps isomorphically to \(G_L\): all these covers are étale at T and the residue field upstairs is \(\bar L\). For every H, rationality of \(x_H\) says that \(D_{\widetilde T}{\cap}H\) surjects onto \(G_L\). Since \(D_{\widetilde T}\) itself projects injectively, this implies \(D_{\widetilde T}{\subset}H\). Therefore

\[
D_{\widetilde T}\subset\bigcap_{H\supset S}H=S.
\]

Both sides project isomorphically onto \(G_L\), so they are equal. Thus \({\beta}\) carries a genuine noncuspidal point decomposition subgroup to the decomposition subgroup of a genuine noncuspidal point of the target, in the whole profinite quotient.

If R lies outside \(E_C\) and P,Q, its proper image \(D_R\) in \({\Pi}_{C_L}\) equals the proper image \(D_T\). The finite-exception uniqueness from Statement \statementref{A.16} forces T=R. This works for every closed R of C outside the one fixed finite set, upon taking L=k(R). It does not assume an a priori bijection of rational point sets, and does not identify one finite quotient with the inverse limit.


\end{proof}
\paragraph{Check record.} Also checked in this writing, 2026-11-18: GPT-6 Astra (OpenAI, ChatGPT Pro mode), root; separate effort setting not exposed. Checked the self-normalizer proof, both compactness steps in point recognition, and commensurability-to-finite-constant-extension equivalence. Opened ST2009 Lemma 1.7 and Proposition 1.8(iii).


\begin{claim}{A.19}{Proved}{The cuspidal criterion with finite exceptional support}
Let $k$ be finite of characteristic $p$, let $C/k$ be smooth proper
geometrically connected of genus at least two, and let $P,Q\in C(k)$.
For $r\ne p$, a cusp-preserving isomorphism $G_P\simeq G_Q$ of the
arithmetic relative cuspidally pro-$r$ quotients over
$\operatorname{id}_{\Pi_C}$ forces $P=Q$.
\end{claim}
\begin{proof}
A.49 carries every noncuspidal point section to an actual point
section in the whole target quotient. A.16 identifies its label with
the original point outside one finite exceptional image, uniformly
over all finite residue fields. Thus the hypothesis of A.44 holds.
A.42 constructs the full central extensions and their difference
cocycle, A.43 evaluates it at each unchanged label, and A.44 gives
the contradiction. Each of those arguments is written out separately
below. This proof uses A.16. The alternative proof A.20--A.21 cites
A.42--A.43 and A.49 directly and has no such dependency.
\end{proof}
\paragraph{Check record.} Also checked in this writing, 2026-11-18: GPT-6 Astra (OpenAI, ChatGPT Pro mode), root; separate effort setting not exposed. Rebuilt the central extension, residue cocycle, same-label evaluation and two separate endgames. Dependency correction uses A.42--A.44 and A.49; support-field proof uses B.68, not A.16 or Boxall.


\begin{claim}{A.20}{Proved}{Proper-field support suffices for the criterion}

Let $\Sigma$ be any set of primes containing $r\ne p$.
Use the full proper geometric pro-$\Sigma$ group in the preceding
criterion. It suffices that the noncuspidal point transport induced
by $\beta$ have unchanged labels outside $C(F)$, for one proper
algebraic subfield $k\subset F\subsetneq\bar k$. In particular,
an exceptional image supported over such an $F$ suffices, even if it
is infinite. The annihilating integer for $[P-Q]$ need not be
coprime to $r$.

\end{claim}

\begin{proof}

This strengthening also applies to the \(J_C\)-large prime sets of \cite{LA_ST2017}.

Fix any prime set \({\Sigma}\) containing a prime \(r\ne p\), and use the full proper geometric pro-\({\Sigma}\) group in the definitions of \(G_P,G_Q\). The same-prime lemma, compactification comparison, central extension construction, and self-normalizer argument remain valid. Assume \({\beta}\) is a cusp-preserving isomorphism over \(\operatorname{id}_{{\Pi}_C^{\Sigma}}\). Suppose that there is a proper algebraic subfield \(F{\subsetneq}\bar k\) containing k such that, for every closed R whose geometric points are not F-rational, the point transported by \({\beta}\) has the same label R. Then P=Q.

We use A.42 for the full central extensions and the residue cocycle, A.43 for same-label evaluation, and A.49 for point recognition. None of these uses A.16. The final Kummer argument is as follows. Let n be the exact order of \([P-Q]\) in \(J_C(k)\), and choose \(f{\in}k(C)^{\times}\) with \(\operatorname{div}(f)=n(P-Q)\). This order exists because \(J_C(k)\) is finite. The argument for the central extensions uses neither that n is p-primary nor that n is coprime to r: the difference \({\kappa}_r(f)-n{\chi}\) has zero residues, hence is a constant class \({\kappa}_r(a)\), and \({\chi}\) evaluates to zero at every point with unchanged label. Thus \({\kappa}_r(f(R)/a)=0\) for every such R.

The function f/a is not an r-th power in \(\bar k(C)^{\times}\). If \(r\nmid n\) its valuations at P,Q already rule this out. If \(r\mid n\) and \(f/a=b^r\), then \(\operatorname{div}(b)=(n/r)(P-Q)\), contradicting the exact order of \([P-Q]\). The geometric divisor class has the same order as its k-rational class, since \(J_C(k){\to}J_C(\bar k)\) is injective. Equivalently, a rational principal divisor descends by Hilbert 90.

Let D be the smooth compactification of the Kummer cover \(z^r=f/a\). It is geometrically connected, since r is prime and f/a is not a geometric r-th power; \(r{\ne}p\) makes it generically separable. Take a finite L/k containing \({\mu}_r\). Because F is a proper subfield of \(\bar k\), some prime ell has finite exponent in the supernatural degree {[}F:k{]}. To see this directly, the nonzero closed subgroup \(G_F{\subset}\widehat{\mathbf Z}\) has a nonzero ell-adic projection for some ell, and a nonzero closed subgroup of \(\mathbf{Z}_\ell\) is open. Put \(L_m\) equal to the degree-\(\ell^m\) extension of L. Then

\[F_0=F\cap\bigcup_m L_m\]

is a finite extension of k: the degrees outside ell are bounded by {[}L:k{]}, and the ell-part of {[}F:k{]} is bounded. Thus all excluded F-rational points lying in \(C(L_m)\) belong to the fixed finite set \(C(F_0)\). Write \(B=\#C(F_0)+2\), including P,Q.

For every remaining \(R{\in}C(L_m)\), the zero r-adic Kummer class implies that f(R)/a is an r-th power in \(L_m\). Since \({\mu}_r{\subset}L_m\) and f has no zero or pole away from P,Q, the fibre of D above R consists of exactly r \(L_m\)-rational points. Hence

\[\#D(L_m)\ge r\bigl(\#C(L_m)-B\bigr).\]

Writing \(Q_m=\#L_m\) and applying the Weil bounds to the geometrically connected proper curves C and D gives

\[(r-1)(Q_m+1)\le (2g_D+2r g_C)\sqrt{Q_m}+rB.\]

This is impossible as \(m{\to}{\infty}\). Thus P=Q. The auxiliary ell is allowed to equal p or r; the extension is an extension of finite constants and no ramification issue occurs. Only the Kummer degree r must be different from p.

In particular, if \(E_C^{\Sigma}\) is supported over a proper algebraic subfield, the needed same-label assertion follows from compactness recognition and the definition of E. Thus, once the generalized configuration transport is established, almost \({\Sigma}\)-separated in the sense of ST2017 implies \({\Sigma}\)-separated, i.e.~\(E_C^{\Sigma}=\varnothing\). This conclusion uses no finiteness assertion about \(E_C\) itself.

The finite-exception case is a special case: finitely many closed-point residue fields generate one finite F, which is proper in \(\bar k\). In fact the original roots-of-unity proof also works directly with arbitrary n, so no coprimality condition was hidden there.


\end{proof}
\paragraph{Check record.} Also checked in this writing, 2026-11-18: GPT-6 Astra (OpenAI, ChatGPT Pro mode), root; separate effort setting not exposed. Rebuilt the central extension, residue cocycle, same-label evaluation and two separate endgames. Dependency correction uses A.42--A.44 and A.49; support-field proof uses B.68, not A.16 or Boxall.


\begin{claim}{A.21}{Proved}{An independent elementary support-field proof}

For a proper hyperbolic curve $C/k$ over a finite field of
characteristic $p$, put
\[
B_C=\{R\in C(\bar k):\exists S\ne R,
       [R-S]\in J_C(\bar k)[p^\infty]\}.
\]
There is a proper algebraic subfield $F\subsetneq\bar k$ containing
$k$ for which $B_C\subset C(F)$. Hence the cuspidal criterion admits a proof from A.42--A.43 and A.49, using
\statementref{A.20} and this support field in place of the
exceptional-set finiteness theorem.

\end{claim}

\begin{proof}

We can remove the Boxall finiteness input entirely from an alternative proof of the main point-separation theorem. Define

\[B_C=\{R\in C(\bar k):\exists S\ne R,\ [R-S]\in J_C(\bar k)[p^\infty]\}.\]

We prove that \(B_C\) is supported over one proper algebraic subfield of \(\bar k\). This is a statement about divisor classes, not a replacement of group collisions by a claimed converse.

\textbf{The fibres of the difference map.}

For every nonzero \(t{\in}J_C(\bar k)\), the set

\[\{(R,S)\in C(\bar k)^2:R\ne S,[R-S]=t\}\]

has at most two elements. Indeed, for two such pairs, R+S' is linearly equivalent to R'+S. If these effective degree-two divisors are equal, the nonzero-difference condition forces the ordered pairs to be equal. Otherwise they define a degree-two pencil on C. There is at most one such pencil when \(g{\ge}2\): two distinct degree-two maps would give a birational map of C to a curve of bidegree (2,2) in \(P^1{\times}P^1\), whose arithmetic genus is one, impossible. The degree-two pencil is separable even in characteristic two, because a purely inseparable degree-two map to \(P^1\) would identify C with a Frobenius twist of a rational curve, again impossible in positive genus. Denote its involution by iota. Equality of the two fibre divisors then forces (R',S')=(iota(S),iota(R)). This is the sole possible second pair.

\textbf{A nontrivial arithmetic subgroup fixing all p-primary torsion.}

Let f be the p-rank of \(J_C\). Then \(T_pJ_C\cong Z_p^f\) and its finite quotients identify with the geometric p-power torsion groups. Choose an odd prime \(\ell{\ne}p\) not dividing \(|\operatorname{GL}_f(\mathbf{F}_p)|\), with \(|\operatorname{GL}_0(\mathbf{F}_p)|=1\). The kernel of \(\operatorname{GL}_f(\mathbf{Z}_p){\to}\operatorname{GL}_f(\mathbf{F}_p)\) is pro-p, as its congruence filtration has elementary abelian p-group successive quotients. The ell-adic factor \(H=\mathbf{Z}_\ell{\subset}G_k\) has trivial image in \(\operatorname{GL}_f(\mathbf{F}_p)\), since the latter order is prime to ell; its image in the pro-p congruence kernel is also trivial. Thus H fixes all of \(J_C(\bar k)[p^{\infty}]\).

If \(t=[R-S]\) is a nonzero p-primary difference, H acts on the at-most-two-element fibre over t. A pro-ell group with ell odd cannot act nontrivially on a set of cardinality at most two. Hence H fixes R and S individually. Therefore

\[B_C\subset C(F),\qquad F=\bar k^H\subsetneq\bar k.\]

If f=0, the p-primary torsion group itself is zero, so \(B_C\) is empty because the Abel map of a positive-genus curve is injective.

\textbf{Application to the conditional criterion.}

From the compactness proof, \({\beta}\) carries the point R to an actual target point S over the same finite field. Because \({\beta}\) is over \(\operatorname{id}_{{\Pi}_C}\), their full proper point sections have equal geometric conjugacy class. Passing to the Jacobian Kummer map proves \([R-S]\) is p-primary. This is only the necessary implication from group collision to divisor-class torsion, which has been proved directly in B.68. Therefore R outside \(B_C\) must have S=R.

Statement \statementref{A.20} now applies with the explicit proper field F above and proves P=Q. This independently establishes the conditional cuspidal criterion without ST2009's exceptional-set finiteness, Boxall's theorem, or an assumption that a p-primary divisor difference produces a group collision. It retains only the exact cusp isomorphism as input, the full same-prime compactification lemma, the compactness recognition, the elementary degree-two-pencil bound, standard curve cohomology, finite-field Kummer theory, and the Weil bound.

Together with A.45--A.48 and the finite-core proof in A.13, this gives the direct route to A.1, without an exceptional-set finiteness input.


\end{proof}
\paragraph{Check record.} Also checked in this writing, 2026-11-18: GPT-6 Astra (OpenAI, ChatGPT Pro mode), root; separate effort setting not exposed. Rebuilt the central extension, residue cocycle, same-label evaluation and two separate endgames. Dependency correction uses A.42--A.44 and A.49; support-field proof uses B.68, not A.16 or Boxall.


\begin{claim}{A.22}{Proved}{Exceptional emptiness and all-extension point separation}

For a proper hyperbolic $C/k$ over a finite field of characteristic
$p$ and every prime set $\Sigma$ containing some $r\ne p$,
$E_C^\Sigma=\varnothing$ is equivalent to injectivity of the
geometric pro-$\Sigma$ point-section map after every finite
constant extension of $k$.

\end{claim}

\begin{proof}

For every prime set \({\Sigma}\) containing some prime \(r\ne p\), let \({\Pi}_C^{\Sigma}\) retain the full arithmetic quotient. Assume C is proper hyperbolic over finite k. Then \(E_C^{\Sigma}\) is empty if and only if the point-section map is injective after every finite constant extension.

If distinct L-rational points have geometrically conjugate sections, choose their lifts so that the subgroups are equal. The universal pro-cover for \(C_L\) is the same universal pro-cover with the open arithmetic group \({\Pi}_{C_L}^{\Sigma}\). Equal decomposition groups are commensurate. The lifts remain distinct because their geometric images are distinct. Thus the original exceptional set is nonempty.

Conversely, take two distinct upstairs points with commensurate decomposition groups \(D_1,D_2\). Proposition 1.8(iii), whose proof only needs Lemma 1.7 and no cofiniteness, shows that their geometric images are distinct. The intersection \(D_1{\cap}D_2\) is open in both. Its arithmetic projection is open in \(G_k\) because the \(D_i\) project injectively onto open subgroups. Choose a finite L/k containing both point residue fields and with \(G_L\) contained in that projection. Since each \(D_i\) projects injectively, their restrictions over \(G_L\) are precisely the same inverse image of \(G_L\) in \(D_1{\cap}D_2\). They give equal sections of distinct L-points. This is a collision in the whole \({\Pi}_{C_L}^{\Sigma}\).

ST2017 Proposition 2.11(i), read in full, states \(E_{C'}=f^{-1}(E_C)\) for any étale \({\Sigma}\)-covering f, including constant extensions. This is compatible with the equivalence. Its ramified statement is an inclusion with the ramification locus, so must not be promoted to equality. Proposition 2.12, under \(J_C\)-largeness, only confines exceptional points to a proper algebraic subfield, finite of degree at most two over the \({\Sigma}'\)-torsion field (degree one for nonhyperelliptic C); it does not exhibit a nonempty exceptional set or contradict emptiness for the full prime-to-p set.


\end{proof}
\paragraph{Check record.} Also checked in this writing, 2026-11-18: GPT-6 Astra (OpenAI, ChatGPT Pro mode), root; separate effort setting not exposed. Checked the self-normalizer proof, both compactness steps in point recognition, and commensurability-to-finite-constant-extension equivalence. Opened ST2009 Lemma 1.7 and Proposition 1.8(iii).


\begin{claim}{A.23}{Proved}{Large prime sets and the finite-or-infinite dichotomy}

Let $\Sigma$ contain a prime $r\ne p$, allowing $p\in\Sigma$. The
transport theorem A.48 applies, and
$E_C^\Sigma=\varnothing$ whenever $\Sigma$ is $J_C$-large in the
sense of \cite{LA_ST2017}; this includes every cofinite $\Sigma$.
More generally, almost $\Sigma$-separation implies
$E_C^\Sigma=\varnothing$, and a finite $E_C^\Sigma$ is necessarily
empty. Thus the exceptional image is either empty or infinite.
This dichotomy alone does not classify every prime set.

\end{claim}

\begin{proof}

For any prime set \({\Sigma}\) containing some \(r\ne p\), put

\[T_{\Sigma'}J_C=\prod_{s\in\Sigma'}T_sJ_C,\qquad H=\ker(G_k\to\operatorname{Aut}(T_{\Sigma'}J_C)).\]

Here the p-adic Tate module means the étale p-adic module, of rank equal to the p-rank. Under the ST2017 definition, \({\Sigma}\) is \(J_C\)-large exactly when H is nontrivial. Every \({\Sigma}'\)-primary geometric torsion point is fixed by H. The fibre argument of the difference-map argument in \statementref{A.21} shows that 2H fixes each ordered pair realizing a nonzero \({\Sigma}'\)-primary difference. Since multiplication by two is injective on the closed subgroup \(H{\subset}\widehat{\mathbf Z}\), 2H is still nonzero. Thus

\[B_C^\Sigma\subset C(\bar k^{2H}),\qquad \bar k^{2H}\ne\bar k.\]

This recovers and makes explicit the support mechanism of ST2017 Proposition 2.12: for a nonhyperelliptic curve the fibre is a singleton, so H itself suffices; for a hyperelliptic curve one may need the index-at-most-two subgroup 2H.

The Kummer assertion used here includes the characteristic prime when
$p\in\Sigma$. Over a finite field $L$, multiplication by $p^n$ is
surjective on $J_C(\bar L)$, with kernel $J_C[p^n](\bar L)$.
It gives an injective map
\[
 J_C(L)/p^nJ_C(L)\longrightarrow
 H^1(G_L,J_C[p^n](\bar L)).
\]
These are the same classes as those read by the abelian geometric
pro-$p$ fundamental group. Indeed, factor $[p^n]$ through the
quotient by the connected component $J_C[p^n]^0$: the first
isogeny is radicial and bijective on geometric points, and the
remaining isogeny is finite \emph{\'etale}, with geometric kernel
$J_C[p^n](\bar L)$. Its lifting cocycle therefore equals the
multiplication-$p^n$ Kummer cocycle. The Abel--Jacobi map identifies
the abelian curve group with the Jacobian group at this prime as
well, by \cite[Proposition 9.1]{MilneJac}. Passing to the inverse
limit is legitimate for these finite groups and detects the entire
$p$-primary part of the finite group $J_C(L)$. Thus equality of
point sections kills the components at every prime in $\Sigma$,
including $p$ when it is retained; the remaining difference has
precisely complementary-prime torsion. No ordinary-Jacobian
hypothesis is required.

In the generalized criterion, compactness carries each R to S with \([R-S]\) \({\Sigma}'\)-primary by Kummer theory at all primes in \({\Sigma}\). Hence outside \(B_C^{\Sigma}\) the labels agree. Statement \statementref{A.20} then proves the generalized conditional criterion for every \(J_C\)-large \({\Sigma}\). Since A.48 proves the configuration transport for every such \({\Sigma}\), it follows that \(E_C^{\Sigma}=\varnothing\) for every \(J_C\)-large \({\Sigma}\) and every proper hyperbolic C over a finite field. Every cofinite \({\Sigma}\) is \(J_C\)-large: choose an auxiliary odd ell outside the finite set \({\Sigma}'\) and outside the prime divisors of all finite groups \(\operatorname{GL}_{\operatorname{rank}T_s}(\mathbf F_s)\), \(s{\in}{\Sigma}'\); then \(\mathbf{Z}_\ell\) acts trivially on the complementary Tate module.

Even without \(J_C\)-largeness, finiteness of \(E_C^{\Sigma}\) supplies finite-field support and Statement \statementref{A.20} applies. Therefore, by the generalized transport theorem A.48, a finite exceptional set must be empty. Every \(E_C^{\Sigma}\) is consequently either empty or infinite. This is a dichotomy, not a classification of which alternative occurs for every arbitrary \({\Sigma}\). The singleton theorem in Statement \statementref{A.26} supplies explicit infinite cases; the \(J_C\)-large theorem supplies empty cases. We have not proved that these exhaust all possible prime sets.


\end{proof}
\paragraph{Check record.} Also checked in this writing, 2026-11-18: GPT-6 Astra (OpenAI, ChatGPT Pro mode), root; separate effort setting not exposed. Read prime-set proofs and their limits, including the degree-two-pencil bound, singleton Frobenius kernel, and Frattini obstruction. Finite data are checked in separate computation events. A.28 remains Open.


\begin{claim}{A.24}{Proved}{Genus-one negative control}

The elliptic curve $E/\mathbf F_2:y^2+xy=x^3+1$ has distinct
rational points $O$ and $T=(0,1)$ whose sections collide in the
entire geometrically prime-to-two arithmetic group. Its exceptional
image is all of $E^{\cl}$. The finite-exception hypothesis used for
proper hyperbolic curves fails here.

\end{claim}

\begin{proof}

Take \(E/\mathbf{F}_2\) given by \(y^2+xy=x^3+1\), with origin O at infinity. Its affine \(\mathbf{F}_2\)-points are (0,1),(1,0),(1,1), so it has four rational points. It is nonsingular: a potential affine singularity must have x=0 from \(\mathbf{F}_y=x\), but at (0,1) the derivative \(\mathbf{F}_x=y+x^2\) is 1; the standard projective infinity point is also nonsingular. The inverse of (x,y) is (x,y+x). Hence T=(0,1) is a nonzero point of exact order two.

The entire geometric fundamental group of an elliptic curve is abelian. Its prime-to-two quotient is the product of the odd Tate modules. Using O as splitting, the class of the point section of T is exactly its Kummer class in that product. The two-primary point T is divisible by every odd integer in the finite group \(E(\mathbf{F}_2)\); consequently every odd-adic Kummer component vanishes. Thus O,T collide in the entire geometrically prime-to-two group, not just in an abelianized approximation of a larger nonabelian group.

For every \(x\in E(\bar F_2)\), the same argument over a finite field containing x gives a collision between x and x+T. These are distinct. Thus the exceptional image is all of \(E^{cl}\); it is infinite. The exact failed AS-AG.40 hypothesis is the finite-exception statement in Step 3. ST2009 Proposition 1.8 assumes hyperbolicity of U; the proper genus-one curve is not hyperbolic. The rest of the Kummer argument does not repair that missing input.


\end{proof}
\paragraph{Check record.} Also checked in this writing, 2026-11-18: GPT-6 Astra (OpenAI, ChatGPT Pro mode), root; separate effort setting not exposed. Read prime-set proofs and their limits, including the degree-two-pencil bound, singleton Frobenius kernel, and Frattini obstruction. Finite data are checked in separate computation events. A.28 remains Open.


\begin{claim}{A.25}{Proved}{The abelianized negative control in genus two}

The smooth proper genus-two curve
$C/\mathbf F_2:y^2+x(x+1)y=x^5+1$ has distinct rational points
$P=(0,1)$ and $Q=(1,0)$ with $[P-Q]$ of exact order two. They
collide in the entire arithmetic extension by the geometric
abelian prime-to-two quotient. Both the cusp-residue and the
compactified-cover cohomology arguments cease to apply to that
abelianized replacement.

\end{claim}

\begin{proof}

Take \(C/\mathbf{F}_2\) given by

\[y^2+x(x+1)y=x^5+1.\]

At finite ramification points x=0,1 there are respectively P=(0,1),Q=(1,0), and \(\mathbf{F}_x=y+x^4\) is 1 at both. The infinity chart \(u=1/x,v=y/x^3\) is

\[v^2+(u+u^2)v=u+u^6,\]

whose derivative in u is 1 at (0,0). This is a smooth genus-two hyperelliptic curve, with three simple Artin--Schreier branch points \((0,1,{\infty})\). For an elementary genus check, put z=y/{[}x(x+1){]} in the function field. The right-hand Artin--Schreier function is \((x^5+1)/(x^2(x+1)^2)\). At 0, subtract the Artin--Schreier coboundary \((1/x)^2+1/x\); the remaining pole has order one. At 1 the original pole has order one, and at infinity it has order one. The degree-two separable map therefore has different contribution two at each of three points, giving \(2g-2=-4+6=2\).

The rational function x/(x+1) has divisor \(2P-2Q\). The class \([P-Q]\) is nonzero, because a function with divisor \(P-Q\) would give a degree-one map to the projective line, impossible in positive genus. Hence its exact order is two. Both point sections become equal up to geometric conjugation in the entire extension by \(\Delta_C^{\rm ab,(2')}\), by finite-field Kummer theory. The program gives \(N_1=N_2=3\) and \(\#J_C(\mathbf{F}_2)=4\), consistently.

The full AS-AG.40 input does not survive replacing all groups by their geometric abelianizations. At a once-punctured curve the sole puncture inertia is zero in the geometric abelianization, since the puncture word is a product of commutators. Thus there is no central cusp kernel \(\mathbf{Z}_r(1)\) in that abelianized punctured group and no cocycle with residue +1 at that cusp. Moreover, even the cohomology-retention statement needed for point recognition fails for general open subgroups of the proper abelian quotient: an index-r geometric subgroup has r-adic abelian rank 2g, while its actual degree-r étale curve cover has genus \(1+r(g-1)\), hence r-adic rank \(2+2r(g-1)>2g\). Therefore the relevant failure occurs before using finite exceptions or Kummer evaluations; AG.39.15 must be applied to the full punctured pro-prime-to-p group and its actual cuspidal kernel, not to an abelianized replacement.


\end{proof}
\paragraph{Check record.} Also checked in this writing, 2026-11-18: GPT-6 Astra (OpenAI, ChatGPT Pro mode), root; separate effort setting not exposed. Read prime-set proofs and their limits, including the degree-two-pencil bound, singleton Frobenius kernel, and Frattini obstruction. Finite data are checked in separate computation events. A.28 remains Open.


\begin{claim}{A.26}{Proved}{Empty and singleton prime sets have infinitely many exceptions}

For every proper hyperbolic curve $C/k$ over a finite field and every
prime $r$, including $r=\operatorname{char}k$, the exceptional
image $E_C^{\{r\}}$ is infinite. Here the full arithmetic quotient
$G_k$ is retained, and for $r=\operatorname{char}k$ we use the
same definition by commensurable point decomposition groups.
For $\Sigma=\varnothing$, the group is $\Pi_C^\varnothing=G_k$;
all rational point sections are equal and every closed point is
exceptional. Thus retaining at most one geometric prime does not
separate points over every finite constant extension.

\end{claim}

\begin{proof}

\textbf{Elementary automorphism lemma.} Let D be a finitely generated pro-r group, with Frattini quotient \(D/\Phi(D)\cong F_r^d\). Then the kernel of \(\operatorname{Aut}(D){\to}\operatorname{GL}_d(\mathbf{F}_r)\) is pro-r.

Here is a proof avoiding an unopened automorphism theorem. For a finite r-group H, every maximal subgroup has index r, so a tuple whose images generate \(H/{\Phi}(H)\) generates H. Fix one basis of \(H/{\Phi}(H)\), and consider the set of all ordered lifts of that basis to H. It has cardinality \(|{\Phi}(H)|^d\), an r-power. The kernel of \(\operatorname{Aut}(H){\to}\operatorname{Aut}(H/{\Phi}(H))\) acts freely on this set: fixing a generating tuple fixes H pointwise. Therefore the kernel's order divides this r-power, so it is an r-group. A finitely generated pro-r group has a cofinal system of open characteristic subgroups (intersect all subgroups of bounded index). Its finite quotients are r-groups; the image of \({\Phi}(D)\) in each is the finite quotient's Frattini subgroup. Passing to the inverse limit of the preceding finite kernels proves the assertion.

Now assume C/k proper hyperbolic has a rational point O, write \({\Delta}_r\) for its geometric pro-r group, and use \(s_O\) to express \({\Pi}_C^{\{r\}}={\Delta}_r{\rtimes}G_k\). Its Frattini rank is \(d=2g\) when \(r\ne p\), and \(d=f_C\), the geometric p-rank of the Jacobian, when \(r=p\). The latter equality follows from the Artin--Schreier comparison and the all-prime Abel--Jacobi comparison in \cite[Proposition 9.1 and Lemma 9.4]{MilneJac}. In particular it is finite. For a pro-r group, finitely many lifts of a basis of its Frattini quotient generate the group: otherwise their closed generated subgroup lies in a maximal open subgroup of index r, contradicting that they generate the quotient. Thus the automorphism lemma applies at the characteristic prime as well. If \(f_C=0\), the pro-p group itself is trivial. Choose a prime q different from r and not dividing \(|\operatorname{GL}_d(\mathbf{F}_r)|\), with \(\operatorname{GL}_0\) interpreted as the trivial group. The q-adic direct factor \(Z_q\subset\widehat Z=G_k\) has trivial image under \(G_k{\to}\operatorname{Aut}({\Delta}_r)\): its image in \(\operatorname{GL}_d(\mathbf{F}_r)\) is a q-group, hence trivial, and its remaining image in the pro-r kernel is both pro-q and pro-r, hence trivial. Consequently \(s_O(\mathbf{Z}_q)\) is central in the entire arithmetic group.

For every closed point x of degree divisible by q, choose \(h{\in}\mathbf{Z}_q\) with q-adic coordinate 1. This h acts nontrivially on a geometric point over x, because its image modulo deg(x) is nonzero. Let \(z=s_O(h)\), and choose any lift \(\tilde{x}\) in the universal pro-cover. Then \(z\tilde{x}\) differs from \(\tilde{x}\) already on the geometric curve, whereas

\[D_{z\tilde{x}}=zD_{\tilde{x}}z^{-1}=D_{\tilde{x}}.\]

Thus x is exceptional. There are infinitely many such closed points: for sufficiently large j the Weil bound implies \(\#C(k_{q^j})>\#C(k_{q^{j-1}})\), so there is a point of exact degree \(q^j\). The equality of the displayed decomposition groups, restricted to a finite field defining both geometric points, gives equal whole section subgroups for two distinct rational points. No finite-level approximation is involved.

If C(k) is empty, first take a finite constant extension with a point. Equality of two decomposition groups in this open arithmetic subgroup makes their decomposition groups in the original arithmetic group commensurable: each original group contains that restricted group with finite index. The geometric points are still distinct. Infinitely many closed points over the extension map to infinitely many closed points over k because the extension is finite. This proves descent of infiniteness directly, without invoking ST2009's additional prime-set hypotheses. Therefore \(E_C^{\{r\}}\) is infinite for every proper hyperbolic C over every finite field and every prime r.

When \(\Sigma=\varnothing\), the geometric quotient is trivial. A geometric point over a closed point of degree n has decomposition group \(G_{k_n}\subset G_k\). Any two such open subgroups of the procyclic group are commensurable. Since the curve has infinitely many geometric points, every closed point is exceptional, and all rational points over any fixed finite extension give the same identity section of that extension's whole arithmetic group.

This proof gives actual all-level collisions after finite extensions. It also shows exactly why ST2009's cofiniteness hypothesis cannot be omitted from AS-AG.40's finite-exception argument. We make no inference that arbitrary finite sets with more than one prime have virtually pro-\({\Sigma}\) automorphism groups; that would require a separate theorem and is not supplied here.


\end{proof}
\paragraph{Check record.} Also checked in this writing, 2026-11-18: GPT-6 Astra (OpenAI, ChatGPT Pro mode), root; separate effort setting not exposed. Read prime-set proofs and their limits, including the degree-two-pencil bound, singleton Frobenius kernel, and Frattini obstruction. Finite data are checked in separate computation events. A.28 remains Open. | Also checked in this writing, 2026-11-18: GPT-6 Astra Pro; finite\_checks agent; effort not separately exposed. Read and checked elementary finite-Frattini automorphism lemma and central Zq collision argument, including descent by commensurability; characteristic-p H1 identification cited but not separately sourced by this reader.


\begin{claim}{A.27}{Proved}{A completely specified pro-three collision}

Let $L=\mathbf F_2[t]/(t^7+t+1)$ and $C/L:y^2+y=x^5$.
The distinct $L$-points
$P=(t,t^6)$ and $Q=(t^2,t^6+t^5)$ have conjugate sections in
the entire arithmetic extension by the geometric pro-three group.
The assertion is not merely a finite-quotient collision.

\end{claim}

\begin{proof}

Let \(L=F_2[t]/(t^7+t+1)\), and take \(C/L:y^2+y=x^5\) with \(O={\infty}\). It is smooth of genus two. The program verifies the field identity for every nonzero element of the 128-element ring, so irreducibility of the chosen modulus is checked rather than assumed. Put

\[P=(t,t^6),\qquad Q=(t^2,t^6+t^5).\]

They satisfy the equation and Q is the \(\mathbf{F}_2\)-Frobenius image of P. The orbit has seven distinct points; the code gives their complete coordinates. Work in the geometrically pro-three quotient over \(\mathbf{F}_2\). We have \(|GL_4(F_3)|=24\,261\,120\), which is not divisible by seven. The preceding argument gives a central image of \(\mathbf{Z}_7{\subset}G_{\mathbf{F}_2}\). Its element with 7-adic coordinate 1 has image 1 in \(\operatorname{Gal}(L/\mathbf{F}_2)\), so sends a lift of P to a lift of Q while leaving the decomposition group unchanged. Intersect that common decomposition subgroup with the preimage of \(G_L\). This gives equal section subgroups for the two L-rational points in the entire geometrically pro-three arithmetic group over L.

The program finds \(\#C(L)=129\), of which 126 points have exact degree seven over \(\mathbf{F}_2\), making 18 exceptional degree-seven closed points. Infinitely many other degrees are handled by the preceding theorem; no bounded computation is being used to prove infiniteness.


\end{proof}
\paragraph{Check record.} Also checked in this writing, 2026-11-18: GPT-6 Astra (OpenAI, ChatGPT Pro mode), root; separate effort setting not exposed. Read prime-set proofs and their limits, including the degree-two-pencil bound, singleton Frobenius kernel, and Frattini obstruction. Finite data are checked in separate computation events. A.28 remains Open. | Also checked in this writing, 2026-11-18: GPT-6 Astra Pro; finite\_checks agent; effort not separately exposed. Read and checked complete r3 central-Z7 proof using A26 argument; explicit finite field and point orbit independently implemented and checked.



\noindent\emph{Statement A.28 and its produced argument are retained in Part II.}
\subsection{A concrete obstruction to the finite-Frattini extension}

\begin{claim}{A.29}{Proved}{Mixed-prime surface groups have infinite Frattini quotient}
Let $\Sigma$ contain two distinct primes $r,s$, and let $\Delta$
be the pro-$\Sigma$ completion of a surface group of genus at least
two. Its Frattini quotient $\Delta/\Phi(\Delta)$ is infinite,
where $\Phi$ is the intersection of the maximal open subgroups.
In particular this applies to the geometric pro-$\Sigma$ group
of a proper curve in characteristic outside $\Sigma$.

For $\Sigma=\{2,3\}$ there is, moreover, an explicit two-generated
group $G$ of order $576$, with $\Phi(G)=1$, and an automorphism of
order seven acting trivially on $G^{\mathrm{ab}}$.
Neither assertion is a counterexample to mixed-prime point separation.
They show why the finite-Frattini and abelianization arguments used
for a single prime do not extend by the same reasoning.
\end{claim}
\begin{proof}
For each positive integer $n$, put $d_n=\operatorname{ord}_{s^n}(r)$.
Choose an element $\zeta_n$ of exact order $s^n$ in
$\mathbf F_{r^{d_n}}^\times$, and form
\[
 G_n=(\mathbf F_{r^{d_n}},+)\rtimes\langle\zeta_n\rangle,
 \qquad (a,i)(b,j)=(a+\zeta_n^i b,i+j).
\]
The action on the additive normal subgroup is irreducible over
$\mathbf F_r$, because $\mathbf F_r[\zeta_n]=\mathbf F_{r^{d_n}}$.
It is faithful. Consequently the complement $H_n=\langle\zeta_n\rangle$
is a maximal subgroup: a subgroup containing it intersects the
additive group in an invariant subspace. Its core is trivial,
since a nonidentity scalar fails to fix some additive element.
Thus $\Phi(G_n)\subset\bigcap_{g\in G_n}gH_ng^{-1}=1$.
The elements $(1,0)$ and $(0,1)$ generate $G_n$, since the scalar
orbit of $1$ spans the whole additive group.

The surface group surjects onto the free group on two generators
by sending two of its $a$-generators to free generators and every
$b$-generator, and all remaining $a$-generators, to one. Its
pro-$\Sigma$ completion therefore surjects onto each $G_n$.
The image of a Frattini subgroup under a surjection is contained
in the target Frattini subgroup: pull back every maximal subgroup.
Hence $\Delta/\Phi(\Delta)$ surjects onto every $G_n$.
Their orders are unbounded, since they are divisible by $s^n$;
the quotient is infinite. The characteristic-$p$ interpretation
uses the geometric specialization isomorphism already checked in
A.4, with all primes of $\Sigma$ invertible.

For the finite example use
\[
 \mathbf F_{64}=\mathbf F_2[t]/(t^6+t+1),\quad
 \zeta=t^7,\quad\lambda=t^9,
 \qquad G=(\mathbf F_{64},+)\rtimes\langle\zeta\rangle.
\]
Here $\zeta$ has order nine and generates the degree-six field;
$\lambda$ has order seven. Multiplication by $\lambda$ on the
additive factor commutes with multiplication by $\zeta$, so
\[
                    (a,i)\longmapsto(\lambda a,i)
\]
is an automorphism of order seven. Since $\zeta-1$ is invertible,
the commutators with $\zeta$ fill the additive factor. Thus
$G^{\mathrm{ab}}=\mathbf Z/9$, on which this automorphism is
the identity. The program \path{code/frattini_two_prime.py}
checks the field, the root orders, all $576^2$ automorphism
compatibilities, the two-generator orbit, the commutator image,
and the trivial complement core. This supplies the exact finite
counterexample to replacing the Frattini quotient by abelianization
in the attempted extension of the single-prime argument.
\end{proof}
\paragraph{Check record.} Also checked in this writing, 2026-11-18: GPT-6 Astra (OpenAI, ChatGPT Pro mode), root; separate effort setting not exposed. Read prime-set proofs and their limits, including the degree-two-pencil bound, singleton Frobenius kernel, and Frattini obstruction. Finite data are checked in separate computation events. A.28 remains Open. | Also checked in this writing, 2026-11-18: GPT-6 Astra Pro; finite\_checks agent; effort not separately exposed. Read and checked full infinite-Frattini proof via irreducible semidirect quotients, plus complete finite576group verification; geometric specialization cited and checked by configuration reader.



\subsection{The covering condition and finitely generated fields}

\begin{claim}{A.30}{Proved}{Prime-to-characteristic Kummer kernels}
Let $K$ be a field finitely generated over $\mathbf F_p$, and let $G/K$
be a semiabelian variety, an extension of an abelian variety by a torus.
For the prime-to-$p$ Kummer map on rational points, the kernel is
\[
 \ker\left(G(K)\longrightarrow
       \prod_{\ell\ne p}H^1(G_K,T_\ell G)\right)=G(K)[p^\infty].
\]
For a field finitely generated over $\mathbf Q$, the analogous map
using every prime has zero kernel. No finite-generation assertion
about $K^\times$ is needed or made.
\end{claim}
\begin{proof}
For each integer $n$ invertible in $K$, the finite Kummer sequence
injects $G(K)/nG(K)$ into $H^1(G_K,G[n])$. Compatibility in $n$
identifies the displayed kernel with
$\bigcap_{(n,p)=1}nG(K)$; in characteristic zero use all positive $n$.
Passing to a finite separable extension splitting the torus preserves
any equality $x=ny$. We may therefore establish the required
implication there; the assertion that a rational point is killed by
a power of $p$ descends through the injective inclusion on points.

For a split torus, the divisor sequence on a normal proper integral
model shows that $K^\times$ modulo its constant-field units embeds in
the free abelian group on prime divisors. Its image is free abelian.
The constants are finite in positive characteristic. In characteristic
zero they form a number field; the unit theorem and the principal
fractional ideal group show again that its multiplicative group is
finite torsion plus a free abelian group. Thus $T(K)$ is free abelian
plus finite torsion. The image of $G(K)$ in the abelian quotient $A(K)$
is finitely generated. Here we use the Lang--N\'eron theorem
\cite[Theorem 2.1]{Conrad2006}: for a finitely generated regular
extension of a constant field, rational points on an abelian variety
modulo the points of its trace form a finitely generated group. The
constant-field trace has finite rational points over a finite field,
and finitely generated points over a number field.

An extension of a free abelian group plus finite torsion by a finitely
generated abelian group is itself a free abelian group plus a finitely
generated group. Indeed choose finitely many lifts of generators and
relations of the quotient. The finitely many relation values involve
only finitely many elements of a basis of the free kernel; all remaining
basis elements split off. In such a group the intersection of all
prime-to-$p$ multiples is exactly its $p$-primary torsion; the
intersection of all multiples is zero. This proves the assertion.
The converse for a $p$-primary point is immediate because every integer
prime to $p$ acts invertibly on its finite cyclic subgroup.
\end{proof}
\paragraph{Check record.} Also checked in this writing, 2026-11-18: GPT-6 Astra (OpenAI, ChatGPT Pro mode), root; separate effort setting not exposed. Read Kummer-kernel proof, cover criterion in both directions with uniqueness/descent, specialization and all hyperbolic types. Corrected the finite-generation wording for tori; nonreduced modulus excluded from the semiabelian kernel statement.


\begin{claim}{A.31}{Proved}{The covering condition is equivalent to collision}
Let $C/k$ be as in \statementref{A.1}, and let $P,Q\in C(k)$.
Their sections are geometrically conjugate if and only if the
following condition holds. For every finite extension $L/k$ and every
connected finite \emph{geometrically} prime-to-$p$ Galois cover
$D\to C_L$ with a specified $P'\in D(L)$ over $P$, there is
$Q'\in D(L)$ above $Q$ such that
$[Q'-P']\in J_D(L)[p^\infty]$. When this condition holds, the lift
$Q'$ is unique. This is the condition of cote 149, scan 36, with
\emph{$p$-torsion} understood as $p$-primary torsion.
\end{claim}
\begin{proof}
Equality of section classes identifies the Galois actions on the
prime-to-$p$ fibre functor. A fixed point $P'$ in the fibre therefore
gives a fixed point $Q'$ with equal lifted section classes. For a
geometric Galois cover of order prime to $p$, the geometric open
subgroup has its full prime-to-$p$ topology: a core is obtained by
intersecting finitely many conjugate subgroups, and the resulting
quotient embeds in a product of groups of prime-to-$p$ order.
Thus the lifted conjugacy is in the correct full arithmetic quotient
of $D$. Abelianization and \statementref{A.30} give the torsion.

We justify uniqueness independently of the theorem being proved.
Over an algebraic closure suppose two such lifts differ by a
nonidentity deck transformation $a$ of order $m$ prime to $p$.
Then $n(aR-R)$ is principal for some power $n$ of $p$. Choose $f$
with that divisor. Its norm under $\langle a\rangle$ is a constant;
scale $f$ to make the norm one. Hilbert 90 writes
$f=a(u)/u$, with a convention for the action on divisors.
Consequently $\operatorname{div}(u)-nR$ is $a$-invariant (replacing
$a$ by its inverse changes no conclusion). Every orbit of this free
action has length $m$, so the degree of an invariant divisor is
divisible by $m$. Its degree is $-n$, a contradiction.

Conversely use a cofinal system of open characteristic normal
subgroups $N$ of $\Delta_C$. They are stable under the action defined
by $s_P$, and the subgroup $N\rtimes s_P(G_k)$ defines a connected
pointed cover $D_N/k$. Its geometric deck group need not be constant.
After a finite constant extension making it constant, the assumed
condition supplies a lift $Q_N$. Uniqueness makes this lift invariant
under the original $G_k$, hence rational over $k$. Norm maps on
Jacobians send the distinguished difference class to one at every
coarser level; uniqueness gives compatibility. Compatible rational
lifts determine a $G_k$-fixed element of the profinite path torsor
from $P$ to $Q$, equivalently a geometric conjugacy of their sections.
This proves the converse in the whole group. By \statementref{A.1}
the condition in fact forces $P=Q$.
\end{proof}
\paragraph{Check record.} Also checked in this writing, 2026-11-18: GPT-6 Astra (OpenAI, ChatGPT Pro mode), root; separate effort setting not exposed. Read Kummer-kernel proof, cover criterion in both directions with uniqueness/descent, specialization and all hyperbolic types. Corrected the finite-generation wording for tori; nonreduced modulus excluded from the semiabelian kernel statement.


\begin{claim}{A.32}{Proved}{Specialization from finite type to a finite field}
If proper hyperbolic curves over all finite fields of characteristic
$p$ have injective geometrically prime-to-$p$ point-section maps,
then so do proper hyperbolic curves over every field $K$ finitely
generated over $\mathbf F_p$.
\end{claim}
\begin{proof}
Suppose $P\ne Q$ collide on $C/K$. Spread the smooth proper curve,
the two sections, and its constant genus over a smooth connected
finite-type scheme $S/\mathbf F_q$ with function field $K$, shrinking
$S$ so that the sections are disjoint. Choose a closed point $s$,
let $R=\mathcal O^h_{S,s}$ and $L=\operatorname{Frac}(R)$. Pullback
along $G_L\to G_K$ preserves the proposed geometric conjugacy.

Smooth proper specialization identifies the geometric prime-to-$p$
fundamental groups of the generic and special fibres over this
henselian base. We use \cite[Expos\'e X, Corollary 3.9, and
Expos\'e XIII, Corollaries 2.9--2.10]{LA_SGA1}, with a smooth proper
curve and no boundary or wild inertia. The comparison is equivariant
for the arithmetic map
$G_L\to\pi_1(\operatorname{Spec}R)=G_{k(s)}$ and gives a map of
arithmetic extensions. This arithmetic map is surjective: a connected
finite \emph{\'etale} cover of the normal henselian domain $R$ has
connected generic fibre, so every finite quotient of $G_{k(s)}$
is already seen by $G_L$.

The specialization map carries the sections supplied by $P,Q$ to
those supplied by $P_s,Q_s$ after this surjection. A single geometric
conjugator therefore conjugates their values for every element of
$G_{k(s)}$. Since $P_s\ne Q_s$, this contradicts the finite-field
hypothesis. This explains the direction of specialization and the
arithmetic surjectivity needed in the reduction.
\end{proof}
\paragraph{Check record.} Also checked in this writing, 2026-11-18: GPT-6 Astra (OpenAI, ChatGPT Pro mode), root; separate effort setting not exposed. Read Kummer-kernel proof, cover criterion in both directions with uniqueness/descent, specialization and all hyperbolic types. Corrected the finite-generation wording for tori; nonreduced modulus excluded from the semiabelian kernel statement.


\begin{claim}{A.33}{Proved}{Every hyperbolic curve over a finitely generated field}
Let $K$ be finitely generated over $\mathbf F_p$, and let $X/K$ be
a smooth geometrically connected hyperbolic curve. Its map from
$X(K)$ to geometric conjugacy classes of sections of its arithmetic
geometrically prime-to-$p$ group is injective. For $K$ finitely
generated over $\mathbf Q$, the same assertion holds for the full
geometric fundamental group. The positive-characteristic assertion
is equivalent to \statementref{A.1}.
\end{claim}
\begin{proof}
First make a finite algebraic extension over which the smooth
geometric compactification and its distinct boundary points descend.
This is possible because their equations and morphisms use finitely
many algebraic coefficients. The purely inseparable part induces
universal homeomorphisms and equivalences on finite \emph{\'{e}tale}
covers, and identifies the absolute Galois groups. The separable part
restricts both sections along the same Galois homomorphism. Either
part therefore preserves a collision, and distinct points remain
distinct. The enlarged field is still finitely generated. This step
handles imperfect fields and boundary residue fields that cannot be
split by a separable extension alone.

For compactified genus at least two, the quotient induced by
$X\hookrightarrow\overline X$ carries a collision to one on the
proper curve. Statements \statementref{A.1} and
\statementref{A.32} exclude it.

If the compactification has genus zero, make a finite separable
extension splitting three boundary points, and choose their
coordinates to be $0,1,\infty$. The coordinate is a morphism
$X\to\mathbf G_m$ injective on geometric points. A collision
has equal Kummer coordinates. By \statementref{A.30} their ratio
is $p$-primary torsion in $K^\times$, hence one in characteristic $p$.

For compactified genus one, a putative rational point $P$ makes
the compactification an elliptic curve $E$, with origin $P$ and
nonempty boundary $B$. A second colliding point $Q$ has class of
exact order $p^a$ in $E(K)$, by \statementref{A.30}. Choose $m$
prime to $p$ with $m^2>p^a$. Pull back the punctured curve by
$[m]:E\to E$, and make a finite separable extension splitting the
boundary and deck group. The lifted collision supplies $P'=0,Q'$
in $D=E-[m]^{-1}B$. Its class in the generalized Jacobian $G_D$ is
$p$-primary. Its projection to $E$ has exact order $p^a$, because
$[m]$ is invertible on geometric $p$-primary torsion. The torus
kernel of $G_D\to E$ has no nonzero geometric $p$-primary points,
so the generalized class itself has exact order $p^a$.
The definition of the generalized Jacobian with reduced modulus
then supplies a function $f$ with
\[
 \operatorname{div}(f)=p^a(Q'-P'),\qquad
 f(T)=1\quad(T\in[m]^{-1}B).
\]
But $f-1$ has at least $m^2$ distinct zeros and pole divisor of
degree $p^a$, impossible. The boundary points are distinct because
$[m]$ is \emph{\'etale}, and none is $P'$ or $Q'$.

In characteristic zero use the Abel map into the generalized
Jacobian and the all-prime case of \statementref{A.30}. Its
injectivity on points follows directly: equality would give a
function with divisor $Q-P$ and value one on the reduced boundary.
For positive compactified genus a degree-one function is impossible;
for genus zero hyperbolicity supplies at least three boundary points,
while a degree-one function takes a specified value only once.
Finally finite fields are finitely generated and proper genus at
least two is hyperbolic, proving the equivalence. This answers
cote 149, scans 35--36, including the affine cases.
\end{proof}
\paragraph{Check record.} Also checked in this writing, 2026-11-18: GPT-6 Astra (OpenAI, ChatGPT Pro mode), root; separate effort setting not exposed. Read Kummer-kernel proof, cover criterion in both directions with uniqueness/descent, specialization and all hyperbolic types. Corrected the finite-generation wording for tori; nonreduced modulus excluded from the semiabelian kernel statement.




\subsection{The configuration statements in the cited literature}

\begin{claim}{A.34}{Proved}{A10, A14}
Mochizuki's original author Proposition~1.8(ii), equivalently
Proposition~1.6(ii) in the separate LaTeX version, asserts the full
prime-to-$p$ punctured-fibre identification when the ground field has
characteristic $p$ and $\Sigma$ consists of all primes different from $p$.
Consequently the noninjectivity established in the configuration audit
contradicts this specialization of the assertion itself. The February
2014 Comments do not amend this assertion.
\end{claim}
\begin{proof}
The conventions preceding original Proposition~1.6 define the geometric
configuration group by taking the maximal pro-$\Sigma$ quotient of the
tame fundamental group and then the maximal cuspidally pro-$\Sigma^\dagger$
quotient over the proper product. For the present finite-field choice,
$\Sigma^\dagger=\Sigma$. A closed subgroup of a pro-$\Sigma$ group is
already pro-$\Sigma$, so the latter quotient makes no change. The same
observation applies to the punctured fibre. Original Proposition~1.8(ii)
then identifies the inverse image of a rational-point decomposition
group under a configuration projection with that punctured group.
It is therefore not a statement about a singleton prime or about only
the completion of the fibre kernel.

The proof cites Stix's monodromy paper. Its relative completion retains
the entire base fundamental group and completes the fibre kernel.
The arXiv version's \S2.3 and Proposition~2.7 give that construction.
They do not justify additionally completing the base. The four corrections
in the 2014 Comments concern descended forms in decomposition recognition,
targets of diagonal covers, a Kummer injection, and infinitude of a
cyclotomic image, respectively. None replaces the fibre identification.
These readings establish the precise conflict; they do not establish a
verdict about every theorem of the paper.
\citep[original pp.~11,14--15]{LA_Mochizuki2007}
\citep{LA_Mochizuki2014,LA_Stix2005}
\end{proof}
\paragraph{Check record.} Also checked in this writing, 2026-11-18: GPT-6 Astra Pro; Pro mode; separate effort setting not exposed; ANABELIAN\_A1\_SUCCESSOR/configuration. Also checked in this writing: Read the complete carried proof and both original/LaTeX Proposition1.8(ii)/1.6(ii), plus the complete2014 Comments. The full mixed-prime fibre assertion is literal; the extra cuspidal Sigma completion is redundant when Sigma=Sigma-dagger. Its failure follows from A10. No verdict about all other theorems is inferred.


\paragraph{Version control.}
The following numbering distinctions matter. In Sa\"idi--Tamagawa 2009,
the supplied arXiv Remark~3.13 is journal Remark~3.14, and the supplied
arXiv Theorem~2.7 is journal Theorem~2.6. The former is the explicitly
conditional configuration-space route to elimination of exceptional
points; journal Remark~3.13 concerns a different extension of the Isom
theorem. Mochizuki's original Theorem~3.10 is Theorem~3.1 in the LaTeX
version. Statements below specify the original numbering when necessary.
\citep{LA_ST2009,LA_Mochizuki2007}

\begin{claim}{A.35}{Proved}{The tame central, abelian, and fixed-prime quotient repair}
Let $k$ be a perfect field, let $C/k$ be a proper smooth
geometrically connected curve of genus at least two, and let
$\Sigma^\dagger=\{r\in\Sigma:r\text{ is invertible in }k\}$
be nonempty. Thus $\Sigma^\dagger=\Sigma$ in characteristic
zero and $\Sigma^\dagger=\Sigma\setminus\{p\}$ in
characteristic $p>0$.
Let $F\to E\to\Delta_2\to1$ be the geometric right-exact configuration
sequence, let $F\twoheadrightarrow\Delta_1$ be the proper fibre quotient,
and let $E\twoheadrightarrow\Delta_1\times\Delta_2$ be the proper product
quotient. Write $K_F,K_E$ for their cuspidal kernels. Suppose the relative
groups are formed from the \emph{tame} boundary groups, with Mochizuki's
maximal cuspidally pro-$\Sigma^\dagger$ convention. The proper groups
may retain $p$ in $\Sigma$. The relative cuspidally pro-$r$ fibre
theorem of Part~A, including A.13, holds for every
$r\in\Sigma^\dagger$.
Then
\[
 K_F^{(r)}\xrightarrow{\sim}K_E^{(r)}\quad(r\in\Sigma^\dagger),
 \qquad K_F^{\mathrm{ab}}\xrightarrow{\sim}K_E^{\mathrm{ab}}.
\]
Consequently the configuration fibre sequence is left exact after
cuspidal abelianization, as well as after the specified fixed-prime
cuspidal completion. The one-puncture geometrically cuspidally central
fibre identification also holds. For a $k$-rational point, the
corresponding arithmetic inverse-image identifications hold over
the unchanged quotient $G_k$.
\end{claim}

\begin{proof}
All geometric groups in the proof are formed over $\bar k$.
No finite-field argument is used. We first supply the uniform
geometric fibre input, then prove the central assertion separately.

\paragraph{The relative fibre input in every characteristic.}
Fix $r\in\Sigma^\dagger$ and a finite quotient $Q$ of the
relative cuspidally pro-$r$ fibre group. Let $R$ be the image of
its cuspidal kernel and put $A=Q/R$. The proper finite
$\Sigma$-quotient $A$ defines a connected \emph{\'etale} Galois
cover $C'\to C_{\bar k}$. The cover
\[
 V=(C'\times C')-\coprod_{a\in A}\Gamma_a
       \longrightarrow U=(C_{\bar k}\times C_{\bar k})-\operatorname{diag}
\]
has group $A\times A$. Its projection to $C'$ has constant proper
compactification and nonempty disjoint labelled boundary sections.
The base genus is at least two. The fibre restriction of $Q$ has
$r$-group image. A.45 applies, also in characteristic zero, and
dominates it by a finite $r$-quotient of $H=\pi_1(V)$.
If $J$ is the kernel of that quotient, the intersection of its
conjugates over $\pi_1(U)/H$ is normal in $\pi_1(U)$, and its
quotient in $H$ embeds in a finite product of conjugate $r$-groups.
The resulting total quotient is therefore an $r$-group by
$A\times A$. It is a $\Sigma$-group, its cuspidal kernel has
$r$-group image, and its fibre restriction dominates $Q$.
Thus it factors through the relative cuspidally pro-$r$ total
group and proves the fibre injection. Every finite quotient of
that total group is tame along the diagonal, since its diagonal
inertia has $r$-power image. A.46 proves middle and right exactness.
This also verifies the right exactness required before passing to
cuspidal quotients: apply A.46 to all the finite tame quotients
defining $E$.

The same argument works for a constant proper compactification
with finitely many disjoint labelled sections. For a prescribed
proper fibre cover, pull it back along the section maps and choose
a common proper Galois base cover trivializing these finitely many
torsors. Its group is a subgroup of a finite product of their
proper $\Sigma$-groups. A.45 then applies to the labelled inverse
boundary; the finite core has an $r$-group kernel over the product
of the two proper deck groups. A.46 again gives right exactness.
This is the finite-proper-pullback form needed in the applications.

\paragraph{Cuspidal kernels and abelianization.}
Right exactness gives a surjection $K_F\twoheadrightarrow K_E$:
an element of $K_E$ lifts to $F$, and its trivial first proper
projection forces the lift into $K_F$. In the relative pro-$r$
fibre isomorphism just proved, taking the kernels over $\Delta_1$
identifies $K_F^{(r)}$ with $K_E^{(r)}$.

There is no cuspidal $p$-primary abelian factor in characteristic
$p$. Each cuspidal kernel is topologically generated by the
ambient conjugates of tame inertia, and a homomorphism to a
finite $p$-group kills every such subgroup. Thus
$K_F^{(p)}=K_E^{(p)}=1$ even when the proper groups retain $p$.
This would not be true for full wild cusp kernels; it is precisely
where the tame-boundary convention is used. An abelian pro-$\Sigma$
group is the product of its primary factors, and the abelian
pro-$r$ quotient is the abelianization of the maximal pro-$r$
quotient. The already proved isomorphisms for all
$r\in\Sigma^\dagger$ consequently give
$K_F^{\mathrm{ab}}\simeq K_E^{\mathrm{ab}}$.
These are the actual maps, equivariant for the geometric proper
fibre action. Comparing extensions over $\Delta_1$ proves the
cuspidally abelian fibre identification.

\paragraph{The full primitive orientation module on the fibre.}
Write $M=\widehat{\mathbf Z}^{\Sigma^\dagger}(1)$, with the local
convention $z^{1/n}\mapsto\zeta z^{1/n}$ for actual tame inertia.
For one puncture $P$, the coinvariants of $K_F^{\mathrm{ab}}$
under the whole geometric fibre are a quotient $M_F$ of $M$:
the cusp kernel is normally generated by the one inertia subgroup,
and all its conjugates agree in these coinvariants. We prove that
no quotient of $M$ has occurred.

Fix $r\in\Sigma^\dagger$ and $n=r^a$. Algebraic Kummer theory
and the Picard group of the proper curve give
\[
 H^2(C_{\bar k},\mu_n)\simeq\operatorname{Pic}(C_{\bar k})/n
       \simeq\mathbf Z/n,
 \qquad c_1(\mathcal O_C(P))\longmapsto1.
\]
The exact source is \citet[Propositions 14.1--14.2]{MilneEtale};
the algebraic Jacobian multiplication theorem used there is
\citet[Theorem 8.2 and Remark 8.4]{CF_MilneAV1986}.
This also proves the comparison
\[
 H^2_{\mathrm{cont}}(\Delta_1,\mu_n)
           \simeq H^2(C_{\bar k},\mu_n)
\]
by an entirely algebraic finite-cover tower argument. At every
finite proper $\Sigma$-cover, a degree-one cohomology class is
killed by its cyclic $r$-power torsor. A degree-two class is killed
by a tower of connected degree-$r$ \emph{\'etale} covers of
total degree divisible by $n$; each exists from nonzero
$\operatorname{Pic}^0[r]$, because the genus remains positive.
The Galois closures of these towers over the given finite
$\Sigma$-cover are still $\Sigma$-covers by the finite-core
product argument. Cohomology above degree two is zero. The
Cartan--Leray tower sequence
\citep[Theorem 14.9 and Remark 14.10]{MilneEtale} therefore
gives the comparison; no complex surface-group presentation is
involved.

The gerbe of $n$-th roots of $\mathcal O_C(P)$, with its
canonical trivialization on $C-P$, now gives the corresponding
extension of $\Delta_1$ by $\mu_n$ and a homomorphism from the
punctured group. Choose its band identification by the actual
local Kummer convention above. The geometric extension class
is $\varepsilon c_1(\mathcal O_C(P))$, with the same convention
sign $\varepsilon=\pm1$ at every compatible level; in particular
it is primitive. Its inertia map is the identity modulo $n$.
The homomorphism is surjective and factors through the
geometrically cuspidally central quotient. Thus $M_F$ retains
the full quotient $\mu_{r^a}$ for every $a$ and every
$r\in\Sigma^\dagger$. A closed subgroup of the product of the
procyclic primary factors is the product of its primary
subgroups, so a kernel killed at all these finite levels is zero.
Consequently $M_F=M$.

\paragraph{The second proper factor adds no central relation.}
It remains to verify this assertion rather than infer it from
the order of inertia before centralization. Let
$X=C_{\bar k}\times C_{\bar k}$ and let $D$ be its diagonal.
For each $n$ supported on $\Sigma^\dagger$, form the local
root charts $z^n=t$ of the smooth divisor $D$. SGA~1 XIII
Proposition 5.2 extends each tame cover with inertia dividing
$n$ to these charts, uniquely up to unique isomorphism.
Their descent isomorphisms therefore extend as well. The
boundary is the root gerbe of $\mathcal O_X(D)|_D$.
Finite covers of that gerbe with trivial stabilizer action
descend to $D$, and its fundamental group thus maps onto
$\pi_1(D)$. Passing to the compatible inverse limit and to
the retained proper pro-$\Sigma$ quotient gives a boundary
image in $E$ mapping onto the diagonal $\Delta_1=\Delta_2$.
This surjectivity is right exactness of the finite-cover
quotient; no finite gerbe stabilizer is assumed to inject.

Conjugation by this boundary image preserves the inertia
orientation. Over $\bar k$ the cyclotomic action is trivial,
so it fixes the image of every actual inertia element in
$M_F$. The first proper factor already acts trivially by the
definition of $M_F$. An element with proper coordinates
$(1,\delta)$ is the product of an element with coordinates
$(\delta^{-1},1)$ from the fibre and a boundary element with
coordinates $(\delta,\delta)$, up to an element of the
cuspidal kernel, whose action on its abelianization is trivial.
It follows that the second proper factor also acts trivially
on $M_F$. Taking its coinvariants therefore adds no relation.
This proves the geometrically cuspidally central fibre
identification with the full module $M$.

\paragraph{Arithmetic groups over the perfect ground field.}
All maps just used come from the actual covers, compactifications,
local Kummer charts, and their natural quotient maps. They are
compatible with the $G_k$-action. The action on surviving inertia
is cyclotomic, as is evident in the same local chart. For a
$k$-rational point, the actual punctured arithmetic group maps
to the inverse image of its proper point section in the relative
arithmetic configuration group. On the geometric kernels this
is the fixed-prime, abelian, or central isomorphism just proved,
and on the arithmetic quotients it is the identity of the full
$G_k$. The elementary isomorphism criterion for two short exact
sequences makes the middle map an isomorphism. No quotient of
$G_k$ is taken. When finitely many closed punctures must first
be labelled, a finite separable splitting extension exists
because $k$ is perfect; the natural constructions commute with
its descent action, so the identifications descend. This proves
the stated field scope and the required finite-pullback form.
\end{proof}
\paragraph{Check record.} Also checked in this writing, 2026-11-18: GPT-6 Astra (OpenAI, ChatGPT Pro mode), root; separate effort setting not exposed. Read the complete replacement proofs for arbitrary retained prime sets and arbitrary perfect fields, respectively, including the normal-core topology, primitive orientation class and tame-boundary restriction. | Also checked in this writing, 2026-11-18: GPT-6 Astra Pro; Pro mode; separate effort setting not exposed; ANABELIAN\_A1\_SUCCESSOR/configuration. Full replacement over arbitrary perfect fields. A45, A46 give uniform relative fibre exactness; Picard/Kummer cycle class and an etale Sigma-cover tower prove primitive central inertia directly; diagonal boundary lifts show the second proper factor adds no relation. No finite-field-only A42 premise. Checked in this writing; parent review requested.


\paragraph{Exact repair boundaries.}
This proves a replacement for the false intermediate assertion over
every perfect ground field with these tame cuspidal conventions,
including proper $\Sigma$-groups which retain the positive
characteristic. It covers all fields of the original paper:
its \S0 defines a nonarchimedean local field to mean a finite
extension of $\mathbb Q_p$, so its local convention
$\Sigma^\dagger=\Sigma$ is exactly the characteristic-zero
case here. It does not turn the original
full sequence into an exact sequence. The following dependency loci were
checked in the original author text. Proposition~1.8(iii),(iv) uses
1.8(ii) for central quotients. The paragraph preceding Proposition~1.12
(original p.~20) cites 1.8(ii),(iii) for surjectivity of the symmetrized
fundamental section. Proposition~1.14 builds cuspidally abelian quotients
as inverse limits. Before Proposition~3.8 (original p.~63), a punctured
fibre is inserted in the configuration group by 1.8(ii); Proposition~3.9(v)
identifies the fixed-$l$ cuspidal quotient, which enters Theorem~3.10.

At these loci one may use the central or abelian fibre identification
when the subsequent construction is central or abelian, and the relative
pro-$l$ identification when it is cuspidally pro-$l$. This specifies the
replacement actually proved here. Constructions on finite proper covers
retain their labelled boundary components. Here is the needed extension
of the relative proof. For a product compactification $C_0\times B$
with finitely many disjoint section graphs, choose the finite Galois
$A$-cover $C_1\to C_0$ required by a finite fibre quotient. Take a
finite Galois $\Sigma$-cover of $B$ dominating its pullbacks along
every section; their finite intersection and core give such a cover.
All lifted punctures are then labelled, the compactification remains
a product, and mod-$r$ monodromy is a shear between proper homology
and the boundary module. The same relative pro-$r$ and normal-core
argument applies, with finite proper deck group $A\times A_B$.
When this group's order is divisible by $p$, the direct algebraic
argument of \statementref{A.13} uses Stix's relative pro-$r$
completion over the full base; it requires no equivariant lift of the
proper cover to characteristic zero.
This proves the required compatibility for these finite proper covers.
For naturality under a finite proper pullback not itself visible
as a $\Sigma$-cover, apply the construction to that curve with
its own proper pro-$\Sigma$ group: naturality of the restriction
maps and the diagonal cycle class gives the comparison square.
One does not pretend that this cover is an open subgroup of the
original pro-$\Sigma$ group. Any conclusion retaining the entire
mixed-prime punctured kernel still loses its asserted fibre
identification. The dependency census below records that loss
separately from the repaired quotient constructions.
\citep[Propositions~1.8,1.12,1.14,3.9 and Theorem~3.10]{LA_Mochizuki2007}

\begin{claim}{A.36}{Proved}{A10}
Mochizuki--Tamagawa's configuration exactness theorem and Wakabayashi's
configuration theorem do not assert the mixed prime-to-$p$ exactness
refuted in Part~A.
\end{claim}
\begin{proof}
Mochizuki--Tamagawa Definition~1.1 restricts the formation to a
PT-formation, namely primary or total. Proposition~2.2 additionally
requires invertibility on the connected base and the stated formation
conditions on the cyclotomic and compactified-fibre mod-$l$ monodromy
images. For classes of all finite $\Sigma$-groups, primary means one
prime; total means all finite groups. The proof uses that the kernel of
the automorphism action of a finitely generated pro-$l$ group on its
mod-$l$ abelianization is pro-$l$. A set of several primes omitting the
characteristic is not one of these two formations. The 2015 Comments
do not enlarge the formation scope. Wakabayashi fixes one invertible
prime throughout. Thus neither theorem supplies the disputed
mixed-prime completion step.
\citep[Definition~1.1 and Proposition~2.2]{LA_MT2008}
\citep{LA_MT2015,LA_Wakabayashi2016}
\end{proof}
\paragraph{Check record.} Also checked in this writing, 2026-11-18: GPT-6 Astra Pro; Pro mode; separate effort setting not exposed; ANABELIAN\_A1\_SUCCESSOR/configuration. Also checked in this writing: Read complete carried proof; MT2008 Definition1.1 and Proposition2.2 with proof, complete2015 Comments, and Wakabayashi formation conventions. Source-scope comparison is correct. A10 is contextual, not a premise of the formation distinction.


\begin{claim}{A.37}{Proved}{A11}
Fix a prime $l$ invertible in the two finite fields. Wakabayashi's
Theorem~A, proved in the more general logarithmic form as Theorem~3.9,
starts with a Frobenius-preserving isomorphism of the geometrically
pro-$l$ groups of once-punctured curves, carrying the specified new
cusp decomposition groups to one another. It extends this isomorphism
to the two-point configuration groups, compatibly with the base-curve
quotient and switching, up to the stated inner ambiguity. Corollary~3.10
transports punctures at other points with corresponding base-curve
decomposition groups. Theorem~C, namely Theorem~4.4, reconstructs the
higher configuration groups from the geometrically pro-$l$ curve group.
These are reconstruction results in the singleton-prime category.
\end{claim}
\begin{proof}
These are the hypotheses and conclusions of the supplied source. Its
proof uses the inverse image of a point decomposition group under the
configuration projection, so the fibre method itself is already present.
Remark~3.10.1 explicitly permits nonuniqueness of the corresponding
point. The theorem retains the geometrically pro-$l$ group of the unpunctured
base curve. In the proper case this is a pro-$l$ proper group; it is
not the full prime-to-$p$ proper group with only a pro-$r$ cuspidal kernel. No operation
in the theorem reconstructs mixed finite quotients from separate
pro-$l$ quotients. Nor may the cofinite-prime exceptional-set theorem
be applied to a singleton set without verifying its hypothesis. Thus
these results do not conflict with AG.40--AG.42 and, with the stated
Sa\"idi--Tamagawa finiteness input alone, do not supply their universal
point-separation conclusion.
\citep[Theorems~3.9,4.4 and Corollary~3.10]{LA_Wakabayashi2016}
\end{proof}
\paragraph{Check record.} Also checked in this writing, 2026-11-18: GPT-6 Astra Pro; Pro mode; separate effort setting not exposed; ANABELIAN\_A1\_SUCCESSOR/configuration. Also checked in this writing: Read complete carried proof and Wakabayashi Proposition1.4, Theorem3.9, Corollary3.10/Remark3.10.1, Theorem4.4 and proof. Corrected proper quotient to base-curve quotient: the source also permits affine base curves. Scope remains singleton, and matching-point existence does not assert uniqueness.


\paragraph{Sa\"idi--Tamagawa dependencies.}
The 2009 configuration route in journal Remark~3.14 is expressly
conditional. Its failure does not by itself invalidate that paper's
Isom theorem. Theorem~2.2 reconstructs the configuration cuspidally
abelian group; Corollary~2.3 restricts the labelled puncture comparison
to the complements of exceptional sets; journal Theorem~2.6 gives
fixed-$l$ cuspidal reconstruction. The 2017 text repeats these as
Theorem~3.2, Corollary~3.3 and Theorem~3.6. Thus their imported
configuration machinery depends on the central, abelian, and fixed-$l$
consequences just distinguished, not on a final theorem asserting
full mixed-prime configuration exactness.

There is a further scope distinction. Original Mochizuki Theorem~3.10
allows arbitrary matching rational points, without a separation
hypothesis. Its proof begins with configuration cuspidally abelian
reconstruction from Theorem~1.16(iii), chooses matching splittings in
the diagonal decomposition groups, and restricts to the corresponding
cuspidally abelian fibres. Thus the reference to an outside-exceptional-set
corollary in Sa\"idi--Tamagawa is not on its own proof that the original
arbitrary-point theorem has no such argument. The needed repair is the
fibre input at the stated quotient level. Our main proof uses its own
relative fibre and transport constructions and does not assume the
published reconstruction theorem in order to prove that repair.
\citep{LA_ST2009,LA_ST2017,LA_Mochizuki2007}

\subsection{Proper support fields and the later exceptional-set theorem}

\begin{claim}{A.38}{Proved}{A10, A14}
In Sa\"idi--Tamagawa 2017, Definition~2.6 calls $\Sigma$ $J_C$-large
when the omitted-prime representation
$\rho_{J_C,\Sigma'}:G_k\to\Aut(T_{\Sigma'}J_C)$ is noninjective.
Definition~2.9 calls $C$ almost $\Sigma$-separated when its exceptional
closed points are supported over some proper algebraic subfield
$F\subsetneq\bar k$. Proposition~2.12 proves this latter conclusion
under $J_C$-largeness. It does not in general assert that the
exceptional set is finite.
\end{claim}
\begin{proof}
Put $F_0=\bar k^{\ker\rho_{J_C,\Sigma'}}$. The proof uses the difference
map
\[
 \delta:C\times C\longrightarrow J_C,\qquad (P,Q)\longmapsto[P-Q].
\]
Off the diagonal its geometric fibres have cardinality at most two,
or at most one when $C$ is not hyperelliptic. Indeed, two different
pairs in a fibre give $P+Q'\sim P'+Q$. A nonconstant degree-two linear
system is the hyperelliptic system; in that case the only possible
second pair is $(\iota Q,\iota P)$, with $\iota$ the hyperelliptic
involution. This argument is the predecessor explicitly cited from
2009 Claim~1.9(i), where the curve has first been made
nonhyperelliptic.

An exceptional pair has distinct images on $C_{\bar k}$ and its
difference lies in $J_C(\bar k)\{\Sigma'\}\subset J_C(F_0)$.
Taking the subgroup of squares of $G_{F_0}$ kills the possible
two-element permutation. Hence both coordinates are defined over a
finite extension of $F_0$, still proper in $\bar k$. This is exactly
the argument of Proposition~2.12. Proposition~2.13 then chooses a prime
$l$ for which the corresponding $\Z_l$ tower has finite intersection
with every fixed finite extension of that support field. It may choose
$l\in\Sigma\cup\{p\}$ under largeness, and $l\in\Sigma$ if the
Jacobian also has positive $p$-rank. Proposition~2.11(i) gives
$E_{C'}=f^{-1}(E_C)$ for finite \'etale $\Sigma$-covers and transports
almost separation through them.
\citep[Definitions~2.6,2.9 and Propositions~2.11--2.13]{LA_ST2017}
\end{proof}
\paragraph{Check record.} Also checked in this writing, 2026-11-18: GPT-6 Astra Pro; Pro mode; separate effort setting not exposed; ANABELIAN\_A1\_SUCCESSOR/configuration. Also checked in this writing: Read complete proof and ST2017 Definitions2.6/2.9, Propositions2.11–2.13 with proofs. The difference-fibre bound and index-at-most-two subgroup give proper support; finite support is not asserted. Proposition2.11(i) is equality for etale Sigma covers; ramified2.11(ii) is only inclusion with ramification.


\begin{claim}{A.39}{Proved}{A10, A14}
The separation argument extends from finite exceptional support to
almost $\Sigma$-separation whenever $\Sigma\setminus\{p\}$ is nonempty,
including $p\in\Sigma$. We use the relative fixed-$r$ fibre and
transport statements of Part~A in that $\Sigma$-category, as extended
by \statementref{A.13}. Consequently the $J_C$-large
almost-separated cases of Sa\"idi--Tamagawa become separated.
\end{claim}
\begin{proof}
Statement~\statementref{A.23}, in its unconditional form in this
paper, proves that almost $\Sigma$-separation implies
$E_C^\Sigma=\varnothing$ for every $\Sigma$ containing some $r\ne p$.
Its hypotheses permit $p\in\Sigma$, and its proof incorporates the
verified transport theorem. It therefore applies to the first assertion
exactly as stated. Sa\"idi--Tamagawa's Proposition~2.12 says that
$J_C$-largeness implies almost $\Sigma$-separation. Applying the first
assertion to that conclusion proves the last assertion.
The support-field condition is not replaced by a finite exceptional-set
condition at either step.
\citep[Definition~2.9 and Proposition~2.12]{LA_ST2017}
\end{proof}
\paragraph{Check record.} Also checked in this writing, 2026-11-18: GPT-6 Astra Pro; Pro mode; separate effort setting not exposed; ANABELIAN\_A1\_SUCCESSOR/configuration. Also checked in this writing: Read original full endgame and replaced its duplicate argument by the precise unconditional A23 corollary. ST2017 Proposition2.12 supplies J-large implies almost-separated. No finite-exception theorem is used.


\paragraph{The elementary full prime-to-$p$ specialization.}
For the main theorem one can avoid the earlier Boxall-based finiteness
input entirely. Let $f_p$ be the $p$-rank of $J_C$. Choose an odd prime
$l\ne p$ not dividing $|\mathrm{GL}_{f_p}(\F_p)|$, with the evident
convention when $f_p=0$. The subgroup $\Z_l\subset G_k$ acts trivially
on $T_pJ_C$, since the congruence kernel in its automorphism group is
pro-$p$, and therefore acts trivially on all geometric $p$-primary
torsion. The preceding two-point difference-fibre bound then puts
every distinct pair with $p$-primary difference over
$\bar k^{\Z_l}$. Along the $l$-power constant-field tower these bad
labels are uniformly bounded. This is a direct special case of the
known support-field method, combined with the new criterion as proved
in Part~A; the difference-fibre lemma itself is not new.

\subsection{What is and is not supplied by the opened sources}

\begin{claim}{A.40}{Proved}{A14}
The following comparison records exact predecessors and the additional
work required in this audit. It is not a priority claim.
\end{claim}
\begin{center}\small
\begin{longtable}{p{0.17\linewidth}p{0.37\linewidth}p{0.38\linewidth}}
\toprule
Step&Opened predecessor&Additional issue resolved or retained\\\midrule
AG.40, point recognition&
Mochizuki original Remark~1.18.2 and its 2014 correction; the conclusion
is ``Any Frobenius-preserving isomorphism $\alpha$ is quasi-point-theoretic.''&
One must prove that the relative quotient still recovers the compactified
$r$-adic cohomology of every visible cover; the finite-intersection
argument must allow all relevant descended forms.\\
AG.40, Kummer evaluation&
Mochizuki Proposition~2.1(ii),(iii): residue exactness and evaluation of
a function's Kummer class at a point. Sa\"idi--Tamagawa repeats it.&
The over-identity cusp isomorphism yields a difference cocycle with
residues $+1,-1$; combining it with transported point labels gives the
contradiction. No identical full criterion was located.\\
AG.41, exactness&
Stix's relative completion; M--T Proposition~2.2 in PT formations;
Hoshi--Sawada--Tsujimura Definition~1.6(ii) and Corollary~1.14;&
Neither states the mixed proper-base, fixed-$r$ cuspidal quotient used
here. The direct proof and its finite-quotient checks remain essential.\\
AG.42, fibre transport&
Mochizuki original Remark~1.7.1 at the central level; Wakabayashi
Corollary~3.10 for a fixed prime; Sun Definition~3.1 and Corollary~3.8.&
The general pullback trick is known. Its use in the relative mixed-base
quotient also requires the local diagonal-boundary lift and compatibility.\\
Specialization and inertia&
SGA1 XIII Corollary~2.9, \S2.10, Lemma~2.11 and Proposition~5.2.&
The hypotheses require a smooth proper pair with relative normal crossings
and invertible coefficient primes; the local Kummer extension exists on
the whole regular chart, by purity.\\
Orbifold completion&
Boggi--Zalesskii Definition~1.1 and Lemma~3.1, with a torsion-free normal
subgroup and quotient in the extension-closed formation.&
The prime-to-$p$ formation meets these hypotheses when the chosen
orbifold quotient has prime-to-$p$ order; no $p$-acyclicity is required.\\
\bottomrule
\end{longtable}
\end{center}
The 2024 Hoshi--Sawada--Tsujimura result requires good log-smooth
pairs and a proper geometrically connected morphism, smooth away
from the base boundary, with smooth relative boundary there, and
log smooth as a morphism of pairs. Definition~1.6(ii) takes the
prime-to-characteristic quotient of the \emph{kernel} of the map
between full tame fundamental groups. Corollary~1.14 identifies
the geometric prime-to-characteristic fibre when it is centre-free.
The base remains the full tame group. This is a closer recent
predecessor for relative completion, and supplies no simultaneous
prime-to-$p$ completion of that base. Its proof in \S1 was checked
at these precise statements.
\citep[Definitions~1.1,1.3,1.6 and Corollary~1.14]{LA_HST2024}

The October 2026 Hoshi--Tsujimura preprint also has a point-section
injectivity theorem: Theorem~A(i), proved as Theorem~1.8, concerns
affine varieties over a characteristic-$p$ field admitting a Galois
extension of degree $p$. It uses the full \'etale fundamental group
and Artin--Schreier evaluations. It does not assert the proper-curve
prime-to-$p$ injectivity at issue here.
\citep[Theorem~A(i) and Theorem~1.8]{LA_HT2026}

The sources were opened in full or at the identified proofs, not merely
through search snippets. Searches included each named paper, both Sun
preprints, and combinations of prime-to-$p$ point separation, exceptional
sets, cuspidalization, and configuration exactness. No independently
published universal prime-to-$p$ point-separation proof was found in
that bounded search. This does not prove novelty or exclude an unknown
predecessor. \citep{LA_Mochizuki2007,LA_Mochizuki2014,LA_ST2009,LA_ST2017,LA_MT2008,LA_Stix2005,LA_Wakabayashi2016,LA_Sun2025,LA_SGA1,LA_BoggiZalesskii2019,LA_HST2024,LA_HT2026}


\paragraph{Check record.} Also checked in this writing, 2026-11-18: GPT-6 Astra Pro; Pro mode; separate effort setting not exposed; ANABELIAN\_A1\_SUCCESSOR/configuration. Also checked in this writing: Read the full predecessor comparison and opened its cited loci. Added HMM2022 Theorem2.5 and HSS2026 Theorem5.6/8.4 scope records. This is a bibliographic comparison, not a new mathematical theorem or a priority claim. Source-proof scope and the bounded search are recorded separately.
\subsection{A finite dependency census for the disputed exactness}

\begin{claim}{A.41}{Proved}{Which imported fibre arguments require replacement}
For the five papers specified in A10, the dependencies on the
configuration fibre assertion are as listed below. At every
listed central, abelian, or fixed-prime use, A.35 supplies the
same quotient statement with its maps and boundary data. The
literal mixed-prime fibre insertion and any argument requiring
that full insertion are lost. This is a census of the disputed
dependency, not a reproof of unrelated theorems in those papers.
\end{claim}
\begin{center}\small
\begin{longtable}{p{0.23\linewidth}p{0.35\linewidth}p{0.34\linewidth}}
\toprule
Source locus&Use of the fibre input&Verified replacement or exact loss\\\midrule
Mochizuki 2007, Proposition~1.8(ii)&
Identifies the full punctured group with a configuration inverse
image; concludes slimness of the full configuration group.&
The full mixed-prime identification is false in the specialization
of A.34. Only its central, abelian and fixed-prime versions are
proved here. Full mixed-group slimness is not inferred from those
versions or from the invalid full-fibre proof.\\
Proposition~1.8(iii),(iv)&
Specializes the fundamental section, proves surjectivity from
boundary inertia, and identifies central quotients.&
The diagonal cycle class and its point restrictions give the
same maps. Primitive boundary inertia gives surjectivity and
the universal central quotient, by A.35. The separate slimness
claim for curve groups in (iii) is proved there from curve-group
results, not from the disputed configuration insertion.\\
Paragraph before Proposition~1.12; Propositions~1.12, 1.14--1.16, 
1.19&
Surjectivity of symmetrized fundamental sections; inverse limits
of central quotients; reconstruction and switching of the
cuspidally abelian configuration group.&
Apply the repaired central maps on each finite proper pullback
and take their compatible inverse limits. The needed comparison
and rigidity are at the abelian level.\\
Paragraph before Proposition~3.8; Propositions~3.8--3.9;
Theorem~3.10&
Writes the full punctured fibre as a subgroup to define diagonal
actions, then uses its fixed-$l$ Lie quotients and finite proper
pullbacks.&
Replace the ambient group and the fibre from the outset by their
relative cuspidally pro-$l$ quotients. A.35 and the retained
boundary give those actions and inclusions. Every Lie object
used here factors through this quotient. The displayed full
mixed-group inclusion on original p.~63 remains invalid.\\
Mochizuki--Tamagawa 2008, Proposition~2.2 and its subsequent
configuration sequences&
Uses its own PT-formation homotopy theorem. The citation to
Mochizuki 2007 in Remark~4.8.2 concerns a proposed later
application.&
No proof in this dependency chain invokes the false mixed
formation. Its singleton and total hypotheses are retained.\\
Sa\"idi--Tamagawa 2009, Theorem~2.2, Corollary~2.3,
Proposition~2.4, Theorem~2.6&
Imports Mochizuki's abelian reconstruction, marked restrictions,
central quotient, and fixed-$l$ reconstruction, respectively.&
These are exactly the repaired quotient inputs. The conditional
full-group proposal in journal Remark~3.14 is lost; it is not a
premise of their Isom proof.\\
Sa\"idi--Tamagawa 2017, Theorem~3.2, Corollary~3.3,
Proposition~3.4, Theorem~3.6&
Reuses the preceding four constructions, with its own exceptional
sets and largeness hypotheses.&
The same four quotient repairs apply. Its support-field and
function-reconstruction arguments have independent hypotheses,
which remain as in C.5--7.\\
Wakabayashi 2016, Proposition~1.4; Theorems~3.9,4.4&
The fibre theorem comes from M--T Proposition~2.2 or Stix in
the single-prime category. The proper two-point case of 4.4
cites Mochizuki's fixed-$l$ theorem.&
No full mixed-prime fibre assertion is needed in the category
used by these proofs. Its imported fixed-$l$ theorem is in the
repaired scope, and also in the singleton case where the
homotopy theorem is valid.\\
\bottomrule
\end{longtable}
\end{center}
\begin{proof}
The full-text occurrence census for original Proposition~1.8(ii)
has three direct loci: its use in the proof of (iii), the
surjectivity paragraph before Proposition~1.12, and the fibre
insertion immediately before Proposition~3.8. The other explicit
citations to Proposition~1.8 use its proper-group slimness in
(i), its curve or central assertions in (iii),(iv), or the
definitions of the punctured groups. We checked the mathematical
content of these loci, not just the citation strings.

For the first two loci, the source's required map sends each
boundary generator to its corresponding orientation-module
generator in the extension defined by the diagonal class.
A.35 constructs this map without a full-fibre injection. Since
the central cuspidal kernel is generated by these images, the
map is surjective and has the universal central quotient kernel.
On a finite proper cover, conjugate boundary components give the
separate factors of the symmetrized extension. The labelled-cover
argument in A.35 proves compatibility with their deck actions.
Taking inverse limits identifies the full cuspidal abelian
module, hence gives the identifications in Proposition~1.14.
The proof of Theorem~1.16(iii) reconstructs these extensions and
their transition maps, and its uniqueness argument uses
Proposition~1.15 at this abelian level. The induced geometric
switching map exists before taking any quotient; its needed
abstract characterization in Proposition~1.19 likewise uses
Proposition~1.15, not a full mixed-group fibre theorem.

At the third locus the paper has $S=S_* =\{x_*\}$, hence
$S_0=\varnothing$. Its ensuing objects are constructed from the
pro-$l$ groups of the punctured finite proper covers, their
weight quotients and their Lie realizations. Consequently each
action in Propositions~3.8--3.9 factors through the relative
cuspidally pro-$l$ fibre with its proper pro-$\Sigma$ quotient.
Replace the displayed full configuration group by that relative
quotient. A.35 identifies its fibre with the relative
punctured group and identifies the image of the diagonal
decomposition group with its retained tame boundary. Conjugation
therefore supplies the same continuous diagonal actions on every
pro-$l$ object. Naturality on finite proper pullbacks supplies
their transition maps. No weight filtration, trace map or
Frobenius splitting in the subsequent proof changes. The
slimness used at the end of Proposition~3.9(v) is explicitly
slimness of the \emph{singleton-$l$} curve and configuration
groups, where M--T Proposition~2.2 applies. Thus the fixed-$l$
reconstruction of Theorem~3.10 has the repaired input, without
restoring the original full mixed-group insertion.

The two Sa\"idi--Tamagawa papers cite precisely these central,
abelian and fixed-$l$ results at the four listed theorem loci.
Their subsequent use in Green compatibility and multiplicative
function reconstruction only uses these quotients; C.5--7
give the separate consequences of removing exceptional sets.
For M--T, the full text restricts configuration exactness to
PT formations in Proposition~2.2. Its sole invocation of the
cuspidalization paper in Remark~4.8.2 is motivational, not a
proof input. Finally Wakabayashi fixes one invertible prime
throughout and gives the exact source of its fibre comparison
in Proposition~1.4; its references to Mochizuki's weight
filtration and fixed-$l$ reconstruction stay in that category.
These checks prove the stated dependency classification.
\citep[Propositions~1.8,1.12,1.14--1.16,1.19,3.8--3.9;
Theorem~3.10]{LA_Mochizuki2007}
\citep[Proposition~2.2, Remark~4.8.2]{LA_MT2008}
\citep[\S2, Remark~3.14]{LA_ST2009}
\citep[\S3]{LA_ST2017}
\citep[Proposition~1.4, Theorems~3.9,4.4]{LA_Wakabayashi2016}
\end{proof}
\paragraph{Check record.} Also checked in this writing, 2026-11-18: GPT-6 Astra Pro; Pro mode; separate effort setting not exposed; ANABELIAN\_A1\_SUCCESSOR/configuration. Also checked in this writing: Read complete carried census/proof and all three direct original1.8(ii) citation loci. Read original1.12, 1.14–1.16, 1.19, 3.8–3.9 and3.10 at their full proof loci for this fibre-input audit. Source central/abelian/fixed-l constructions factor through A35 repairs, including finite proper pullbacks. A34 supplies the exact full-fibre loss. The long Lie-algebra and unrelated global Isom theorems are not independently reproved by this audit.




\subsection{The three logically separate parts of the cuspidal criterion}

The following statements separate the construction, its evaluation,
and the exceptional-set input. This separation is essential to the
two different proofs of A.1. In each statement $k$ is finite of
characteristic $p$, $C/k$ is smooth proper geometrically connected
of genus at least two, and $r\ne p$ belongs to the retained prime
set $\Sigma$. The group $\Pi_C^\Sigma$ retains the whole $G_k$.
For $A\in C(k)$, $G_A$ is the arithmetic maximal cuspidally
pro-$r$ quotient of $C-A$ relative to $\Pi_C^\Sigma$.

\begin{claim}{A.42}{Proved}{Central extensions and the difference cocycle}
Suppose $P,Q\in C(k)$ and
$\beta:G_P\xrightarrow{\sim}G_Q$ is a continuous isomorphism over
$\operatorname{id}_{\Pi_C^\Sigma}$ preserving the cusp decomposition
conjugacy class. Put $M=\mathbf Z_r(1)$, with its arithmetic
cyclotomic action and its actual local inertia convention:
$z^{1/r^a}\mapsto\zeta z^{1/r^a}$ represents $\zeta\in\mu_{r^a}$.
The geometrically cuspidally central quotients
fit into exact sequences
\[
 1\to M\to E_A\to\Pi_C^\Sigma\to1\qquad(A=P,Q).
\]
The isomorphism induced by $\beta$ is the identity on $M$.
On the arithmetic geometrically pro-$\Sigma$ group $\Gamma$ of
$C-\{P,Q\}$ there is a continuous $M$-valued crossed homomorphism
$\chi$ with residues $(1,-1)$ at $(P,Q)$. If
$\operatorname{div}(f)=n(P-Q)$ for $f\in k(C)^\times$ and $n>0$,
then there is $a\in k^\times$ such that
\[
                  \kappa_r(f)-n[\chi]=\kappa_r(a)
                  \quad\text{in }H^1(\Gamma,M).           \tag{A.42.1}
\]
No assertion about exceptional points or transported point labels
is used in this construction.
\end{claim}
\begin{proof}
Quotient the cuspidal kernel by its commutators with the entire
geometric subgroup. The remaining kernel is a quotient of $M$:
the cusp kernel is normally generated by inertia at the single
puncture, and all its conjugates now agree. Its arithmetic action
is cyclotomic, as is the action on tame inertia.

We check that no further quotient of $M$ has occurred by an
algebraic argument. For each $a\ge1$, the degree isomorphism is
\[
 H^2(C_{\bar k},\mu_{r^a})\simeq\mathbf Z/r^a,
                    \qquad c_1(\mathcal O_C(A))\longmapsto1.
\]
This is the smooth proper curve case of
\citet[Proposition 14.2 and Theorem 24.1(a)]{MilneEtale}.
The comparison with $H^2(\Delta_C^\Sigma,\mu_{r^a})$ follows
without lifting the curve: in the tower of finite geometric
$\Sigma$-covers, a degree-one cohomology class is killed by its
finite cyclic torsor, and a degree-two class is killed by an
\emph{\'etale} cover of degree divisible by $r^a$. Such covers
exist because a positive-genus Jacobian has nonzero prime-to-$p$
Tate module and the Abel map identifies its abelian covers.
The same argument applies after each cover; cohomology above
degree two is zero. The Hochschild--Serre tower argument in
\citet[Theorem 14.9 and Remark 14.10]{MilneEtale} therefore gives
the comparison. The arithmetic comparison follows by applying
Hochschild--Serre for the unchanged $G_k$ to both sides.

Consider the gerbe of $r^a$-th roots of $\mathcal O_C(A)$.
Its canonical trivialization on $C-A$ gives a homomorphism from
the punctured group to the corresponding extension of
$\Pi_C^\Sigma$ by $\mu_{r^a}$. Fix the band identification once,
at every compatible level, by requiring that the local Kummer
coordinate $w^{r^a}=z$ identify the image of actual inertia with
$w\mapsto\zeta w$, hence with $\zeta$. Its extension class is
$\eta_{A,a}=\varepsilon c_1(\mathcal O_C(A))$, where the single
sign $\varepsilon\in\{1,-1\}$ depends only on the convention
relating gerbes and group-extension factor sets. It is independent
of $A$ and of $a$: the local root chart and its band identification
are the same and commute with the power transition maps. This
fixes the local residue at $+1$ without requiring a convention for
the sign of a spectral-sequence transgression. The homomorphism
is surjective:
the punctured group surjects onto the proper group, and its
inertia surjects onto the kernel. It factors through the
geometrically central cuspidal quotient, since its geometric
kernel is central. Thus that quotient retains inertia modulo
every $r^a$. A closed subgroup of $\mathbf Z_r$ is either zero
or $r^b\mathbf Z_r$; it follows that its inertia kernel is zero.
This proves the exact sequence with full kernel $M$ and also
identifies the geometric extension class with the primitive
orientation class $\eta_{A,a}$ normalized by that actual inertia.
For any two rational points $P,Q$, the degree isomorphism gives
$\eta_{P,a}=\eta_{Q,a}$, and both are generators. All these arguments work in characteristic
$p$, including $p=2$ and $r=2$ when $r\ne p$.

The induced map $\bar\beta:E_P\to E_Q$ acts on the procyclic
kernel by a unit $u\in\mathbf Z_r^\times$. Naturality of
extension classes over the identity of the geometric base gives
$u\eta_{P,a}=\eta_{Q,a}=\eta_{P,a}$. Since this common class
is primitive, $u\equiv1$
modulo every $r^a$, and $u=1$.

On the continuous fibre product $E_P\times_{\Pi_C^\Sigma}E_Q$
put $\chi(e_P,e_Q)=\bar\beta(e_P)e_Q^{-1}\in M$, using additive
notation on $M$. Direct multiplication gives
$\chi(ee')=\chi(e)+e\chi(e')$. Pull back along the two natural
punctured-group maps from $\Gamma$. At $P$ the second map kills
inertia and the first has residue $1$; at $Q$ the first kills
inertia and the inverse second map has residue $-1$.
The Kummer cocycle of $f$ has residues $(n,-n)$, so their
difference vanishes on both inertia groups. The coefficient
module is trivial under the geometric kernel; the cocycle
therefore vanishes on their closed normal closure. That closure
is the kernel of the proper quotient. The difference descends
to $H^1(\Pi_C^\Sigma,M)$.

Finally inflation is an isomorphism
\[
                  H^1(G_k,M)\xrightarrow{\sim}
                  H^1(\Pi_C^\Sigma,M).
\]
Indeed the possible extra term is
$H^1(\Delta_C^\Sigma,M)^{G_k}$. This is a submodule of
$\operatorname{Hom}(T_rJ_C,\mathbf Z_r(1))$, whose Frobenius
eigenvalues have complex absolute value $\sqrt{\#k}$; none is
$1$. The module is torsion-free, so its invariant submodule
vanishes. Finite-field Kummer theory identifies the left side
with the $r$-primary completion of the finite group $k^\times$.
It is represented by some $a\in k^\times$, proving (A.42.1).
\end{proof}
\paragraph{Check record.} Checked only in this writing, GPT-6 Astra Pro, 18 November 2026: complete proof written, root targeted reading and configuration-reader second reading; separate effort setting not exposed.


\begin{claim}{A.43}{Proved}{Evaluation at a point with unchanged label}
Keep the hypotheses and notation of A.42. Let $R$ be a closed
point different from $P,Q$, put $L=k(R)$, and suppose that
$\beta$ carries a point decomposition subgroup at $R$ in $G_P$
to a geometric conjugate of the point decomposition subgroup
at that same $R$ in $G_Q$. Then
\[
             [\chi]|_{s_R(G_L)}=0,
             \qquad \kappa_{r,L}(f(R)/a)=0.               \tag{A.43.1}
\]
There is no finiteness hypothesis on the set of points where
this same-label assumption holds.
\end{claim}
\begin{proof}
Choose the two point subgroups to have the same image $D_R$
in $\Pi_{C_L}^\Sigma$. A conjugator in the target has image
normalizing $D_R$. A.17, equivalently
\citet[Lemma 1.7]{LA_ST2009}, says its normalizer is exactly
$D_R$: the point subgroup projects isomorphically onto the
entire $G_L$. Multiply the conjugator on the appropriate side
by the unique element of the target point subgroup with
that same proper image. Its conjugation of the point subgroup
does not change, and its proper image becomes $1$.
In $E_Q$ it therefore belongs to $M$. Conjugating a splitting
by an element of $M$ changes its $M$-cocycle by a coboundary.
The difference cocycle consequently restricts to zero in
$H^1(G_L,M)$. Restrict (A.42.1) along the actual point section;
functorial Kummer theory evaluates $f$ at $R$, yielding the
second equality. The field extension is the residue field of
this individual point; no common field for all $R$ is required.
\end{proof}
\paragraph{Check record.} Checked only in this writing, GPT-6 Astra Pro, 18 November 2026: complete proof written, root targeted reading and configuration-reader second reading; separate effort setting not exposed.


\begin{claim}{A.44}{Proved}{The finite-exception endgame}
Keep the hypotheses of A.42. If all but finitely many closed
points of $C$ satisfy the same-label assumption of A.43, then
$P=Q$.
\end{claim}
\begin{proof}
Assume $P\ne Q$. Since $J_C(k)$ is finite, the class $[P-Q]$
has finite positive order $n$. A geometric function with divisor
$n(P-Q)$ descends to $k(C)$: its Galois ratios are constants,
and Hilbert 90 removes their scalar cocycle. Thus choose $f$
as in A.42. The finite excluded set has only finitely many
geometric points and hence finitely many values under $f/a$.
The finite nonconstant map $f/a:C_{\bar k}\to\mathbf P^1$
is surjective. There are infinitely many $r$-power roots of
unity in $\bar k$, because $r\ne p$. Choose a nonidentity
one, $\zeta$, outside the finite excluded values, and a
geometric point above it. Its residue field $L$ is finite,
and its underlying closed point satisfies A.43. But $f(R)/a$
is the nonzero $r$-primary element $\zeta$ of $L^\times$,
whose image in its $r$-primary completion is nonzero.
This contradicts (A.43.1). Hence $P=Q$.
\end{proof}
\paragraph{Check record.} Checked only in this writing, GPT-6 Astra Pro, 18 November 2026: complete proof written, root targeted reading and configuration-reader second reading; separate effort setting not exposed.


\paragraph{Check record and exact dependency separation.}
Checked only in this writing as separate statements: GPT-6 Astra
(OpenAI, ChatGPT Pro mode), 18 November 2026, direct proof
checking and algebraic rederivation. They replace the corresponding
three portions of A.19; those portions were also checked in this
writing after the earlier readings. A.42 depends on curve
cohomology, Abel covers and Kummer theory, and A.43 additionally
uses A.17. Neither uses A.16. A.44 is where finiteness of the
exceptional labels is used. In the alternative proof, A.20
uses A.42--A.43 and the Weil bounds, while A.21 supplies
proper-field support by the degree-two-pencil argument and
the complementary torsion representation. It uses neither
A.16 nor A.44. Thus the dependency edge A.20$\to$A.19
is replaced by A.20$\to$A.42,A.43; the edge through A.16
is removed rather than hidden in a citation to the earlier
unsplit criterion.

% New working identifiers A.45--A.48, allocated paper numbers A.45--A.48.
% GPT-6 Astra Pro, Pro mode; separate effort setting not exposed.
% This is a direct algebraic replacement for the characteristic-zero route.
\subsection{A direct algebraic configuration argument}

We give a proof in characteristic $p$ whose inputs contain neither lifting
of curves to characteristic zero nor the specialization statements of
SGA~1, Expos\'e XIII, 2.9--2.11. It also allows $p\in\Sigma$.
We retain the definitions of $F_{\mathrm{cr}}$ and $E_{\mathrm{cr}}$ in
A.3, now using geometric \emph{\'etale} fundamental groups, with
$r\in\Sigma$ and $r\ne p$.

\begin{claim}{A.45}{Proved}{Finite single-prime torsors in a constant compactification}
Let $k$ be algebraically closed of characteristic $p\geq0$, let $r$ be
a prime invertible in $k$, and let $C_0,S$ be smooth proper connected
$k$-curves, with $g(S)\geq1$. Let $m\geq1$, and let
$d_1,\ldots,d_m:S\to C_0\times S$ be pairwise disjoint sections.
Put $X=C_0\times S$, $D=\coprod_i d_i(S)$, $U=X-D$, and
$f:U\to S$. For a geometric point $s$ of $S$, every finite $r$-group
quotient of $\pi_1(U_s)$ is dominated by the restriction of a finite
$r$-group quotient of $\pi_1(U)$. Equivalently,
$\pi_1(U_s)^{(r)}\to\pi_1(U)^{(r)}$ is injective.
More precisely, every finite $r$-group torsor on $U_s$ extends to a
torsor with the same group on $U\times_S S'$ for some connected
finite \'etale cover $S'\to S$ whose Galois closure is an $r$-cover.
\end{claim}

\begin{proof}
Write $a=\mathbf F_r$ for the constant sheaf and fix an identification
$a(1)\simeq a$, possible over $k$. All cohomology in this proof is
\'etale cohomology unless continuous group cohomology is explicitly
indicated. We use the following algebraic inputs, in the exact forms
recorded in our source table: cohomological purity for a smooth pair of
codimension one (\citealp[Theorem 16.7]{MilneEtale}), proper base change
(ibid., Theorem 17.10), the cohomology of a proper smooth curve (ibid.,
Proposition 14.2), Leray (ibid., Theorem 12.7), and Cartan--Leray for
finite \'etale covers and their inverse limits (ibid., Theorem 14.9 and
Remark 14.10). None of these inputs is a characteristic-zero comparison.

\paragraph{The relative abelian calculation.}
Let $j:U\hookrightarrow X$ and $i:D\hookrightarrow X$.
Purity and localization give the distinguished triangle
\[
 i_*a_D(-1)[-2]\longrightarrow a_X\longrightarrow Rj_*a_U
 \longrightarrow i_*a_D(-1)[-1].
\]
Push it forward by the proper projection $\bar f:X\to S$.
Proper base change identifies $R^q\bar f_*a$ with the constant sheaf
$H^q(C_0,a)_S$. Since $D$ is a disjoint union of $m$ sections,
its higher direct images vanish in positive degrees. Thus
\[
 R^0f_*a=a_S,\qquad R^qf_*a=0\ (q\geq2),
\]
and, with $\mathcal M=R^1f_*a$,
\[
 0\to H^1(C_0,a)_S\to\mathcal M
 \to a(-1)_S^{\,m}\overset{\sum}{\longrightarrow}a(-1)_S\to0.       \tag{DC.1}
\]
The last map sends every labelled point class to its degree, namely one.
It is surjective even when $r\mid m$. In particular $\mathcal M$ is
a finite locally constant sheaf. The two terms in the resulting short
exact sequence
\[
 0\to H^1(C_0,a)_S\to\mathcal M
 \to\ker(a^m\overset{\sum}{\longrightarrow}a)_S(-1)\to0
\]
are constant. After a vector-space splitting at one geometric point,
monodromy is a subgroup of the additive group of homomorphisms from the
last vector space to the first. Its image is therefore an elementary
abelian $r$-group. No equivariant splitting and no division by $m$ is
used.

Here the identifications with the cohomology of $U_s$ are themselves
part of the argument, not an unproved base-change assertion for an open
morphism. Pulling the displayed localization triangle to a geometric
fibre gives its localization triangle: a section is transverse to the
fibre with intersection multiplicity one, so its local Kummer residue
map is unchanged. Proper base change for $X$ and $D$, followed by the
five lemma in (DC.1), identifies the stalk of $\mathcal M$ with
$H^1(U_s,a)$ and the remaining stalks with $H^q(U_s,a)$. The same
argument commutes with every finite \'etale base change $S'\to S$.
The local compatibility just used can also be checked in a relative
parameter $t$: the residue of a Kummer class $t$ is $1$, both before
and after restriction to the fibre. All schemes in this calculation
are smooth over $k$ and all torsion has order invertible in $k$.

\paragraph{Killing cohomology by $r$-covers of the base.}
If $\mathcal V$ is a finite locally constant $a$-sheaf on $S$ with
$r$-group monodromy image $B$, a class of $H^1(S,\mathcal V)$ is a
continuous crossed homomorphism $c:\pi_1(S)\to V$. The map
\[
 g\longmapsto(c(g),\rho(g))\in V\rtimes B
\]
is a continuous homomorphism to a finite $r$-group. Its kernel defines
a connected finite Galois $r$-cover on which both the sheaf and the
class become trivial.

Also, $H^2(S,a)$ is killed by pullback to any connected \'etale cover
of degree $r$. Such a cover exists because $g(S)\geq1$ and
$H^1(S,\mu_r)\simeq\operatorname{Pic}^0(S)[r]\ne0$.
Under the degree description
$H^2(S,\mu_r)\simeq\operatorname{Pic}(S)/r\simeq\mathbf Z/r$,
pullback multiplies degree by $r$, and is zero. These statements remain
true after every finite \'etale base change, since its genus is still
at least one. A tower of finite Galois $r$-covers is dominated by a
finite Galois $r$-cover of the original base: take a finite normal core,
whose kernel quotient embeds in a product of conjugate $r$-groups.

The first five terms of Leray show that every prescribed class of
$H^1(U_s,a)$ extends to $H^1(U_{S'},a)$ after such a base change.
Indeed, first kill the $r$-group monodromy of $\mathcal M$; the class
then becomes a section of $\mathcal M$. Its sole obstruction to lifting
is its transgression in $H^2(S',a)$, which the preceding degree-$r$
cover kills.

Moreover every class in $H^2(U,a)$ is killed after a finite $r$-cover
of the base. Since $R^2f_*a=0$ and $H^3(S,a)=0$, the Leray filtration
gives the exact sequence
\[
 H^2(S,a)\longrightarrow H^2(U,a)
 \longrightarrow H^1(S,\mathcal M)\longrightarrow0.              \tag{DC.2}
\]
Kill the image in the last term by its affine $r$-group representation
as above. The restricted class then comes from $H^2(S',a)$, and a
further degree-$r$ \'etale cover kills that term. The filtration and
these operations are functorial in the base change.

\paragraph{Central induction for nonabelian torsors.}
We prove the extension assertion by induction on the order of the finite
$r$-group $G$. The trivial group is immediate. Choose a central
subgroup $A\subseteq G$ of order $r$, and put $G_0=G/A$.
Inductively, the quotient of the prescribed fibre $G$-torsor extends to
a $G_0$-torsor on $U_{S_1}$, with a chosen identification on the fibre.
Write $h_0:\pi_1(U_{S_1})\to G_0$ for its continuous homomorphism.
The central extension $1\to A\to G\to G_0\to1$ defines a class
$\omega\in H^2(G_0,A)$, and its pullback
$h_0^*\omega\in H^2_{\mathrm{cont}}(\pi_1(U_{S_1}),A)$ is the
obstruction to a continuous lift to $G$. This assertion follows
directly by choosing a set-theoretic section $G_0\to G$: its factor
set is a two-cocycle, and changing a lift by a one-cochain removes the
defect precisely when that class vanishes. All cochains here are
continuous and have finite values.

For any connected scheme $Z$ used here there is a natural injection
\[
 H^2_{\mathrm{cont}}(\pi_1(Z),A)\hookrightarrow H^2(Z,A).         \tag{DC.3}
\]
To verify it, apply Cartan--Leray to the inverse system of all connected
finite \'etale Galois covers. The first cohomology of that universal
pro-cover is zero: every finite $A$-torsor occurring at a finite level
is trivialized at a higher level. The five-term sequence therefore
gives (DC.3). This argument also shows the natural identification of
continuous and \'etale $H^1$, and is functorial for our base changes.
In particular we have not assumed that the total space is an \'etale
$K(\pi,1)$.

By (DC.2), after a finite $r$-cover of the base the image of
$h_0^*\omega$ in \'etale $H^2$ vanishes. Injectivity in (DC.3) supplies
a continuous lift $h:\pi_1(U_{S_2})\to G$. Its restriction to the
fibre and the prescribed fibre lift have the same image in $G_0$.
Their pointwise ratio, because $A$ is central, is a homomorphism to
$A$, hence a class of $H^1(U_s,A)$. Extend that class after one more
$r$-cover of the base by the preceding $H^1$ argument, and multiply
$h$ by the resulting central character. The adjusted lift restricts
to the prescribed fibre homomorphism. Thus its associated $G$-torsor
is the required extension. If the original fibre torsor is connected,
the total torsor is connected as well, since each connected component
of a finite \'etale cover surjects onto the connected base space and
would otherwise disconnect its fibre.

Finally replace $S'\to S$ by a dominating Galois $r$-cover.
Its pullback gives a normal $r$-index subgroup $H'$ of $H=\pi_1(U)$.
The extended torsor is a finite $r$-quotient of $H'$, with kernel $J$.
The intersection of $\gamma J\gamma^{-1}$ over the finitely many
cosets $H/H'$ is normal in $H$; its quotient is an extension of an
$r$-group by an $r$-group. Restriction to the fibre dominates the
prescribed quotient. Finite quotients separate points in a pro-$r$
group, which proves injectivity.
\end{proof}
\paragraph{Check record.} Checked only in this writing, 18 November 2026, GPT-6 Astra (OpenAI), ChatGPT Pro mode; separate effort setting not exposed. Complete proof written and read by its named authoring reader; see CHECK\_EVENTS.tsv for additional second readings and computation scope.


\begin{claim}{A.46}{Proved}{Right exactness for tame configuration quotients}
In the situation of A.45, let $U\to S$ be the open projection.
For every connected finite Galois \'etale cover $W\to U$ tame along
$D$, with group $T$, the image $I\subseteq T$ of $\pi_1(U_s)$ is
normal, and $T/I$ is a quotient of $\pi_1(S)$ through the projection.
Consequently the sequence of tame fundamental groups is exact at its
middle and right terms. The same conclusion holds for any quotient of
$\pi_1(U)$ all of whose finite quotients are tame along $D$.
\end{claim}

\begin{proof}
Normalize $X$ in $W$, obtaining $\bar W\to X$. The normalization is
finite, since $X$ is excellent, and is proper over $S$.
SGA~1, Expos\'e XIII, Proposition 5.2, says that a tame cover of a
regular local scheme with normal-crossings boundary becomes \'etale
after adjoining roots of the boundary parameters of orders equal to
the relevant inertia least common multiples; the orders are invertible.
For our smooth relative divisor with disjoint sections, each strictly
local connected branch is therefore a Kummer chart $z^e=t$, where
$t$ is a relative parameter and $e$ is invertible. This chart is smooth
over $S$, with relative coordinate $z$; its boundary $z=0$ is \'etale
over $S$. Thus $\bar W\to S$ is smooth proper, and
$\bar W-W$ is finite \'etale over $S$.

The Stein factor $S_W$ of $\bar W\to S$ is finite \'etale over $S$.
For clarity, this is the usual proper-smooth connected-component
statement: the geometric fibres are reduced, and the algebra of
global functions on a proper smooth fibre is a product of copies of
its algebraically closed ground field; proper base change and the
Stein factorization identify these algebras in a finite \'etale family.
Equivalently, SGA~1, Expos\'e X, Proposition 1.2 applies to this proper
flat morphism with geometrically separable fibres.
Its fibre over $s$ is the set of connected components of $\bar W_s$.
Deleting finitely many points from a connected smooth curve does not
disconnect it, so this is also the set of components of $W_s$.

The $T$-action extends uniquely to the normalization and to $S_W$.
At $s$, its transitive action on components identifies that finite set
with $T/I$ as a homogeneous set. Since $S_W$ is connected, its base
monodromy is transitive and commutes with $T$. The automorphism group
of the transitive $T$-set $T/I$ is $N_T(I)/I$; it can be transitive
only when $N_T(I)=T$. Hence $I$ is normal. It follows that
$S_W\to S$ is the finite \'etale Galois cover with group $T/I$, and
the quotient map $\pi_1(U)\to T/I$ factors through $\pi_1(S)$.
Surjectivity to the base follows because the pullback of every
connected finite \'etale base cover remains connected. Taking all
finite quotients gives the assertion; images of profinite groups are
closed, so equality of their images in every finite quotient is
equality of the closed subgroups themselves.
\end{proof}
\paragraph{Check record.} Checked only in this writing, 18 November 2026, GPT-6 Astra (OpenAI), ChatGPT Pro mode; separate effort setting not exposed. Complete proof written and read by its named authoring reader; see CHECK\_EVENTS.tsv for additional second readings and computation scope.


\begin{claim}{A.47}{Proved}{Algebraic survival of the entire $r$-adic cusp inertia}
For $C/k$ smooth proper of genus at least one over an algebraically
closed field of characteristic $p\geq0$, $P\in C(k)$, and $r\ne p$,
the tame $r$-inertia $\mathbf Z_r(1)$ at $P$ injects into
$\pi_1(C-P)^{(r)}$. Consequently it injects into every
$F_{\mathrm{cr}}$ with $r\in\Sigma$.
\end{claim}

\begin{proof}
Choose a nonzero element of $\operatorname{Pic}^0(C)[r]$.
Its Kummer torsor is a connected cyclic \'etale cover $C'\to C$
of degree $r$. There are at least two distinct points $P_1,P_2$ above
$P$. For every $n\geq1$, divisibility of $\operatorname{Pic}^0(C')(k)$
under $r^n$ gives a degree-zero divisor class $[B]$ with
$r^n[B]=[P_1-P_2]$. Choosing a representing divisor, we obtain a
rational function $h$ with
\[
 \operatorname{div}(h)=P_1-P_2-r^nB.
\]
The extension $z^{r^n}=h$ has degree $r^n$ because the valuation of
$h$ at $P_1$ is congruent to $1$ modulo $r^n$. It has inertia of
order $r^n$ at $P_1$ and $P_2$, and is unramified elsewhere outside
the preimage of $P$: every other valuation is divisible by $r^n$,
and the residue fields are algebraically closed. Taking the compositum
of its conjugates under the cyclic deck group of $C'\to C$ produces
a Galois extension of $k(C)$ with $r$-group Galois group, ramified
only above $P$, whose inertia surjects onto a cyclic group of order
$r^n$.

These are finite \'etale $r$-covers of $C-P$ for every $n$.
The image of the local tame $\mathbf Z_r(1)$ therefore has unbounded
$r$-power finite quotients; its kernel, a closed subgroup of
$\mathbf Z_r$, must be zero. All these quotients factor through
$F_{\mathrm{cr}}$, since their images of $K_F$ are $r$-groups.
The Kummer and Picard facts used here are algebraic; see
\citet[Propositions 14.1--14.2]{MilneEtale}. In particular the proof
does not use the topological presentation of a punctured surface.
\end{proof}
\paragraph{Check record.} Checked only in this writing, 18 November 2026, GPT-6 Astra (OpenAI), ChatGPT Pro mode; separate effort setting not exposed. Complete proof written and read by its named authoring reader; see CHECK\_EVENTS.tsv for additional second readings and computation scope.


\begin{claim}{A.48}{Proved}{Diagonal boundary and arithmetic transport, with $p$ allowed in $\Sigma$}
Let $C/k$ be smooth proper geometrically connected of genus at least
two over a finite field of characteristic $p$, and let
$r\in\Sigma\setminus\{p\}$. Let $\Pi_C$ be the geometrically
pro-$\Sigma$ arithmetic fundamental group, retaining the whole $G_k$.
Let $\mathcal E_{\mathrm{cr}}$ be the arithmetic relative
cuspidally pro-$r$ configuration group. There is a diagonal boundary
subgroup $\mathcal D\subseteq\mathcal E_{\mathrm{cr}}$ with exact
sequence
\[
 1\to\mathbf Z_r(1)\to\mathcal D\to\Pi_C\to1.
\]
Its map to $\Pi_C\times_{G_k}\Pi_C$ is diagonal.
For each $P\in C(k)$, with compatible paths, the arithmetic group
$\Pi_{C-P,\mathrm{cr}}$ equals
$\operatorname{pr}_2^{-1}(s_P(G_k))$, and its actual cusp subgroup
is $\mathcal D\cap\operatorname{pr}_2^{-1}(s_P(G_k))$.
If the proper sections at $P$ and $Q$ are geometrically conjugate,
their punctured arithmetic groups are isomorphic over the identity of
$\Pi_C$ and $G_k$, with the first actual cusp taken to a geometric
conjugate of the second.
\end{claim}

\begin{proof}
We first obtain the geometric fibre input directly from A.45--A.46.
For a finite fibre quotient $Q$, divide by its cuspidal $r$-subgroup
to obtain the proper quotient $A$, and use the $A\times A$ deck-graph
cover $V$. A.45 gives a finite $r$-quotient of $\pi_1(V)$ detecting
the restriction of $Q$. The intersection of its kernel's conjugates
over $A\times A$ has quotient in $\pi_1(V)$ embedded in a finite
product of $r$-groups. Its total quotient is an $r$-group by
$A\times A$, factors through $E_{\mathrm{cr}}$, and dominates $Q$
on the fibre: an element killed in the total quotient first has
trivial proper $A$-image and is then killed by the chosen fibre
quotient. Thus the fibre map is injective. Its finite total quotients
are tame along the diagonal, since diagonal inertia has $r$-power
image; A.46 gives middle and right exactness. This establishes
the required geometric short exact sequence within the present
proof, independently of A.13.

Put $X=C\times_k C$, $D=\operatorname{diag}(C)$, and
$N=\mathcal O_X(D)|_D$. For $n=r^a$, the root stack
$\mathfrak X_n=\sqrt[n]{(X,D)}$ is defined locally, if $D=(t)$, by
\[
 \bigl[\operatorname{Spec}(R[z]/(z^n-t))/\mu_n\bigr].
\]
Its boundary $\mathfrak G_n$ is the gerbe of $n$-th roots of $N$ on
$D$. We use this language for descent of the Kummer charts; we do not
attribute a statement about stacks to SGA~1.
Specifically, SGA~1 XIII, Proposition 5.2, extends any finite tame
cover whose inertia orders divide $n$ to an \'etale cover of each
Kummer chart. Uniqueness of normalization extends the descent
isomorphisms, including their $\mu_n$-equivariance, from the dense open
complement. They glue to a finite \'etale cover of $\mathfrak X_n$.
Conversely, restriction of a finite \'etale cover of that stack to
$X-D$ has inertia of order dividing $n$. This proves precisely the
cover comparison used below, by finite \'etale descent.

A finite \'etale cover of $\mathfrak G_n$ on which its stabilizer
$\mu_n$ acts trivially descends to $D$: this is ordinary descent on
local trivializations $B\mu_n\times D_i$. Pullbacks of covers of $D$
are fully faithful. The corresponding fundamental group sequence is
therefore right exact:
\[
 \mu_n\longrightarrow\pi_1(\mathfrak G_n)
 \longrightarrow\pi_1(C)\longrightarrow1.                       \tag{DC.4}
\]
No injectivity of the finite stabilizer is assumed. The transition
maps for $r^{a+1}\to r^a$ are onto on the base and on the stabilizer
images, hence onto on the middle groups. Compact inverse limits of
(DC.4) give
\[
 \mathbf Z_r(1)\longrightarrow\mathcal T
 \longrightarrow\pi_1(C)\longrightarrow1.                      \tag{DC.5}
\]
Every finite quotient of $\mathcal E_{\mathrm{cr}}$ is tame with
$r$-power inertia, for the reason just given.
Restriction of its extended root-stack cover to the boundary produces
compatible maps from $\mathcal T$; thus there is a continuous map
$\mathcal T\to\mathcal E_{\mathrm{cr}}$ over the diagonal proper
map. Define $\mathcal D$ as its closed image.

Here is the precise quotient argument at the proper geometric base.
Write $\mathcal T_{\bar k}$ for the geometric kernel of $\mathcal T$,
and replace it by its maximal pro-$\Sigma$ quotient, leaving $G_k$
unchanged. The resulting arithmetic quotient $\mathcal T_\Sigma$
has a right exact sequence
\[
 \mathbf Z_r(1)\longrightarrow\mathcal T_\Sigma
 \longrightarrow\Pi_C\longrightarrow1.                        \tag{DC.6}
\]
Indeed, on geometric groups (DC.5) is an extension of $\pi_1(C_{\bar k})$
by the inertia image. After maximal pro-$\Sigma$ quotient, quotienting
out this image represents exactly the finite $\Sigma$-quotients of
$\pi_1(C_{\bar k})$; this is right exactness of completion, not an
assumption that it preserves arbitrary kernels. The map to
$\mathcal E_{\mathrm{cr}}$ factors through $\mathcal T_\Sigma$
because its geometric image is pro-$\Sigma$. Consequently
$\mathcal D\to\Pi_C$ is onto, with kernel exactly its inertia image,
a quotient of $\mathbf Z_r(1)$.

The fibre immersion $C-P\to X-D$, $x\mapsto(x,P)$, extends to a
map of pairs transverse to $D$ with multiplicity one. In the local
Kummer charts it takes the diagonal inertia to the actual cusp inertia
with that multiplicity. A.47 shows that the latter is the whole
$\mathbf Z_r(1)$ in $F_{\mathrm{cr}}$, and the geometric injection
just proved embeds this fibre group into $E_{\mathrm{cr}}$. Thus inertia in (DC.6) has
zero kernel in $\mathcal D$, proving the displayed short exact
sequence of the statement. This argument also works for $r=2$.

The actual arithmetic fibre induces a homomorphism
\[
 \Pi_{C-P,\mathrm{cr}}\longrightarrow
 \mathcal E_{\mathrm{cr}}\times_{\Pi_C,s_P}G_k.
\]
It is an isomorphism on geometric kernels by the exact sequence
proved above and the identity
on $G_k$; comparison of short exact sequences makes it an isomorphism.
Since $s_P$ is injective, the target is
$\operatorname{pr}_2^{-1}(s_P(G_k))$. The local punctured henselian
trait at $P$ gives the actual cusp decomposition subgroup. Its
inertia is $\mathbf Z_r(1)$ by A.47 and its residue quotient is the
whole $G_k$, retained injectively by the proper point section.
Its map into $\mathcal D$ is the identity on these inertia and
residue groups. The short exact sequence comparison therefore
identifies it with the full inverse image of $s_P(G_k)$ in
$\mathcal D$. Paths may be chosen in the same connected boundary
tower for all points under comparison, so this statement uses one
common $\mathcal D$ and does not assume point separation.

Finally suppose $\gamma\in\Delta$ satisfies
$\gamma s_P(g)\gamma^{-1}=s_Q(g)$ for every $g\in G_k$.
Choose a geometric element $d\in\mathcal D$ mapping to $\gamma$.
Conjugation by $d$ identifies the two fibre subgroups and their actual
cusps. On the first proper coordinate it induces
$\operatorname{Inn}(\gamma)$. The geometric subgroup of the $Q$-fibre
surjects onto $\Delta$ under that first coordinate, so choose in it
$e$ mapping to $\gamma^{-1}$. Then conjugation by $ed$ is over the
identity of $\Pi_C$ and $G_k$. It carries the first actual cusp to
$e$ times the second cusp times $e^{-1}$, a geometric conjugate.
All choices change the resulting identifications only by the geometric
inner conjugacies explicitly allowed in the statement.
\end{proof}
\paragraph{Check record.} Checked only in this writing, 18 November 2026, GPT-6 Astra (OpenAI), ChatGPT Pro mode; separate effort setting not exposed. Complete proof written and read by its named authoring reader; see CHECK\_EVENTS.tsv for additional second readings and computation scope.


\paragraph{Dependency conclusion for G1(ii).}
The geometric proof of A.13 is A.45 plus the finite proper-cover
and finite-core construction written out above, followed by A.46
for exactness at the middle term. Its boundary and arithmetic clause
uses A.48. The latter independently derives its geometric input
from A.45--A.46 and uses A.47 for inertia; it does not depend on
A.13. These four arguments are in characteristic $p$,
including when $p\in\Sigma$; they invoke no lifting theorem and none
of SGA~1 III 7.3--7.4 or XIII 2.9--2.11. SGA~1 XIII 5.2 is used for
local tame Kummer charts only. In contrast, the original Stix route
uses XIII 4.3--4.4, and the latter explicitly cites XIII 2.9 for its
cohomological hypotheses. It proves the desired geometry without an
actual lift, but does not by itself satisfy the stricter literal
source-exclusion condition. The present direct central-induction proof
removes that dependency rather than concealing it.

\subsection{Recognition without an exceptional-set hypothesis}

\begin{claim}{A.49}{Proved}{Point recognition before identifying its label}
Let $k$ be finite of characteristic $p$, let $C/k$ be smooth proper
geometrically connected of genus at least two, and let
$r\in\Sigma\setminus\{p\}$. Write $G_P,G_Q$ for the relative
cuspidally pro-$r$ arithmetic groups of $C-P,C-Q$ over
$\Pi_C^\Sigma$. An isomorphism $\beta:G_P\simeq G_Q$ over the
identity of $\Pi_C^\Sigma$ preserving the actual cusp decomposition
conjugacy class carries every point section of $C-P$, over its
finite residue field, to an actual point section of $C-Q$.
The source and target points have geometrically conjugate proper
sections. No finiteness or uniqueness assertion concerning
exceptional points is assumed.
\end{claim}
\begin{proof}
Let $L/k$ be finite and $R\in(C-P)(L)$. Restrict $\beta$ to the
preimages of $G_L$, and put $S=\beta(D_R)$ after choosing a lift
of the source point. The group $S$ projects isomorphically to
$G_L$. For each open subgroup $H\le(G_Q)_L$ containing $S$,
let $U_H\to C_L-Q$ be the connected cover it defines.
Its full arithmetic projection makes $U_H$ geometrically connected.
The source cover corresponding to $\beta^{-1}(H)$ has an
$L$-point above $R$. A.14 identifies the full pro-$r$ geometric
compactification group at this open level. A.15 therefore compares
the integral Frobenius polynomials of the two compactifications,
and the cusp-preserving isomorphism compares their boundary
permutations. Subtracting those boundary counts gives
$\#U_H(L)>0$. These are finite sets.

The inverse system $\{U_H(L)\}_{H\supset S}$ has a point.
Indeed, its product is compact. Every finite collection of
compatibility conditions is satisfied by taking a point on
the cover for the intersection of the finitely many subgroups
involved and projecting it. The closed compatibility conditions
thus have the finite-intersection property. This argument does
not require that the individual transition maps be surjective.

Let $(R_H)_H$ be a compatible family and let
$T\in(C-Q)(L)$ be its bottom point. The finite geometric fibres
over $T$ have compatible lifts of all the $R_H$. A second
compactness argument supplies one geometric lift $\widetilde T$
in the full pro-cover. At an open point there is no geometric
inertia, so its decomposition subgroup $D_{\widetilde T}$
projects isomorphically onto $G_L$. Since each $R_H$ is rational,
$D_{\widetilde T}\cap H$ surjects onto $G_L$, and injectivity
of the arithmetic projection on $D_{\widetilde T}$ implies
$D_{\widetilde T}\subseteq H$. Hence
\[
 D_{\widetilde T}\subseteq\bigcap_{H\supset S}H=S.
\]
The last equality holds for every closed subgroup of a profinite
group: an element outside it is excluded by an open subgroup
containing it, by compactness and finite quotients. Both groups
on the two sides project isomorphically onto $G_L$, so the
inclusion is equality. Projection to the proper group, over
which $\beta$ is the identity, proves the final assertion.
\end{proof}
\paragraph{Check record.} Checked only in this writing, GPT-6 Astra Pro, 18 November 2026: complete proof written, root targeted reading and configuration-reader second reading; separate effort setting not exposed.


\paragraph{Check record.}
Checked only in this writing as a separate statement: GPT-6 Astra
(OpenAI, ChatGPT Pro mode), 18 November 2026, direct proof check
and explicit separation of the compactness argument from the
exceptional-set argument. The same recognition argument inside
A.18 was also checked in this writing. The dependencies of this
statement are A.14--A.15, not A.16. A.18 is recovered by adding
the uniqueness statement outside the finite exceptional image
of A.16. A.20--A.21 instead use the present statement and the
Kummer implication proved in B.68.

\subsection{An independent finite separation certificate}

\begin{claim}{A.50}{Proved}{A common cyclic-nineteen witness for the two characteristic-two points}
Let
\[
 C/\mathbf F _2:\quad y^2+x(x+1)y=x^5+1,
 \qquad P=(0,1),\quad Q=(1,0),\quad O=\infty.
\]
Set $L=\mathbf F _8=\mathbf F _2[t]/(t^3+t+1)$ and $J=J_{C_L}$.
Then $\#J(L)=76$ and $H=J(L)[2^\infty]$ has order four.
The Lang map $F_8-1:J\longrightarrow J$ factors as
\[
 J\longrightarrow J/H\xrightarrow{\lambda}J,
 \qquad \deg\lambda=19.
\]
Pulling $\lambda$ back along $j_P:R\mapsto[R-P]$ gives a connected,
constant cyclic, finite \emph{\'etale} cover $D\to C_L$ of degree
nineteen, with the chosen point $(P,0)$ above $P$. Its fibre above $Q$
also splits completely. It has genus twenty and
\[
 \frac{\#J_D(L)}{\#J(L)}=12284193340665686881,
\]
an odd integer. In particular this cover is a class-number witness at
both $P$ and $Q$, independent of the proof of the point-section theorem.
The same construction based at any of the nine points of $C(L)$ is a
class-number witness. Its relative class number is the displayed integer
at $P,Q,O$ and is $81428104953590501$ at each of the other six points.
\end{claim}

\begin{proof}
We give the finite data, their interpretation, and the collision criterion.
First we check the geometric model. The partial derivative with respect
to $y$ is $x(x+1)$. At its zeros the derivative with respect to $x$
is $y+x^4$, which equals one at both points on the curve. Thus the
affine model is smooth. Dividing by $x^2(x+1)^2$ gives the
Artin--Schreier equation for $y/(x(x+1))$. Its poles at $1$ and
infinity are simple; at zero the pole of order two becomes simple
after subtracting $(1/x)^2+1/x$. Thus there are precisely three
simple poles. Their different exponents are two, and
Riemann--Hurwitz gives genus two. There is a unique ramified point
at infinity. The degree-two hyperelliptic pencil has canonical class
$2O$; equivalently the regular differential $dx/(x(x+1))$ has divisor
$2O$.
Write an element of $L$ as the integer whose binary digits are its
coefficients in $1,t,t^2$. In that convention the points, in the order
used below, are
\[
 (0,1),(1,0),(3,2),(3,4),(5,4),(5,6),(7,2),(7,6),O.
\]
Direct enumeration over $\mathbf F _2,\mathbf F _4,\mathbf F _8,
\mathbf F _{64}$ gives respectively $3,3,9,87$ points. The genus-two
identity
\[
 \#J(\mathbf F _q)
       =\tfrac12\bigl(\#C(\mathbf F _q)^2+
                           \#C(\mathbf F _{q^2})\bigr)-q
\]
follows from the first two Newton identities and the reciprocal Weil
polynomial. At $q=8$ it gives $76$. An independent enumeration of all
reduced Mumford representatives gives exactly seventy-six as well.
Here a representative is $(u,v)$ with $u$ monic of degree at most two,
$\deg v<\deg u$, and $u\mid v^2+(x^2+x)v+x^5+1$.
The program gives every representative and the whole addition table;
it checks all $76^3$ associativity triples and all inverses.

We specify the arithmetic sufficiently to reproduce that table.
For $D_i=(u_i,v_i)$ put
\[
 d=\gcd(u_1,u_2,v_1+v_2+h),\qquad
 s_1u_1+s_2u_2+s_3(v_1+v_2+h)=d,
 \quad h=x^2+x,
\]
and form
\[
 u=\frac{u_1u_2}{d^2},\qquad
 v=\frac{s_1u_1v_2+s_2u_2v_1+s_3(v_1v_2+f)}d\pmod u,
 \quad f=x^5+1.
\]
While $\deg u>2$, replace $u$ by the monic associate of
$(f+hv+v^2)/u$ and $v$ by $h+v$ modulo this new polynomial.
These formulas are the ideal product and reduction for the quadratic
algebra $L[x,y]/(y^2+hy-f)$: cancellation of the common ideal is
encoded by the displayed Bezout relation, and replacing the ideal by
its complementary factor in the norm gives the reduction formula.
All divisions have zero remainder in the retained execution.

Multiplication by four maps the seventy-six representatives onto a
cyclic group of order nineteen with kernel $H$. Choose the generator
specified in the JSON certificate; taking discrete logarithms after
multiplication by four gives a homomorphism
$\chi:J(L)\twoheadrightarrow\mathbf Z/19$. Its table is checked on
all $76^2$ pairs. On the displayed points its values on $[R-P]$ are
\[
       0, 0, 1, 18, 7, 12, 8, 11, 0.                 \tag{A.50.1}
\]
The same list holds with $Q$ or $O$ as base point. Direct Mumford
arithmetic checks that $[P-Q]$ is nonzero of exact order two.
Geometrically, $\operatorname{div}(x/(x+1))=2P-2Q$, since the two
points are ramification points of the degree-two $x$-map; a principal
divisor $P-Q$ would give a degree-one map to the projective line,
contradicting the genus. This also proves its exact order without
relying on the table.

The Lang map is separable, with constant kernel $J(L)$. Quotienting its
source by $H$ leaves a separable isogeny with constant kernel
$J(L)/H$ of order nineteen. The Abel map induces the isomorphism on
prime-to-characteristic geometric abelian fundamental groups established
in the Kummer compatibility argument. Its pullback is therefore a
geometrically connected cyclic cover. The point $P$ splits by
construction. Equation (A.50.1) shows that $Q$ splits too. Equivalently,
the Frobenius translation on its fibre is $\chi([Q-P])=0$.
Riemann--Hurwitz gives $g(D)=1+19(g(C)-1)=20$.

For completeness we compute the relative class number, not only its
parity. For a nontrivial character $\zeta^\chi$ of the cyclic deck
group, expand the unramified Euler product by effective divisors.
The resulting polynomial in genus two is
\[
 L_P(\chi,T)=1+
   \left(\sum_{R\in C(L)}\zeta^{\chi([R-P])}\right)T
       +8\zeta^{\chi([K_C-2P])}T^2.
\]
Indeed Riemann--Roch makes every coefficient above degree two vanish
and gives the displayed coefficient of degree two; the sum of a
nontrivial character over $J(L)$ is zero. The canonical class is
$2O=2P$, so the last root of unity is one. Therefore
\[
 L_P(\chi,1)=12+\zeta+\zeta^{18}+\zeta^7+
                      \zeta^{12}+\zeta^8+\zeta^{11}.
\]
Factoring the regular-representation Euler factor at each closed point
gives $Z(D,T)/Z(C,T)=\prod_{i=1}^{18}L_P(\chi^i,T)$. Evaluation at
one gives the asserted class-number ratio. The independent integer
resultant is
\[
 \operatorname{Res}_Z\left(\sum_{j=0}^{18}Z^j,
      12+Z+Z^{18}+Z^7+Z^{12}+Z^8+Z^{11}\right)
           =12284193340665686881.
\]
There is also a short parity check: $2$ has order eighteen modulo
nineteen, so the first polynomial is irreducible over $\mathbf F _2$;
the second polynomial reduces to six nonzero monomials and is not
its multiple. For the other base points, subtract the corresponding
entry of (A.50.1) from every exponent and use
$\chi([K_C-2R])=-2\chi([R-O])$. The same finite resultant calculation
gives the second integer in the statement. Each is odd.

We finish by proving why an odd relative class number excludes a
collision, without invoking A1. If the two sections of $C_L$ were
geometrically conjugate, choose an invariant geometric path between
their base points. Its lift to the pointed cyclic cover produces an
$L$-rational $Q'$ above $Q$ and geometrically conjugate sections at
$P',Q'$. The Kummer compatibility applied on $D$ implies
$T=[Q'-P']\in J_D(L)[2^\infty]$. The class is nonzero because the
Abel map of a positive-genus smooth proper curve is injective on
geometric points.

Let $B=(\ker(J_D\to J))^0$. Pullback followed by norm is
multiplication by nineteen, and $J\times B\to J_D$ is an isogeny.
Thus $\#B(L)$ is exactly the computed odd ratio. On the rational
two-primary group the averaging idempotent
$19^{-1}\sum_{g\in\mathbf Z/19}g$ has complementary image in
$B(L)[2^\infty]=0$. Hence every deck transformation fixes $T$.
For a generator $g$ this gives the equivalence of the distinct disjoint
effective divisors $gQ'+P'$ and $Q'+gP'$. They define a degree-two
pencil. Such a pencil on a curve of positive genus in characteristic
two is separable: a purely inseparable degree-two map to the projective
line would identify the square subfield with a rational function
field and force genus zero. Thus $D$ is hyperelliptic. Its unique
hyperelliptic involution is central in its automorphism group. The
order-nineteen deck transformation descends to an order-nineteen
automorphism of the projective line. A fixed point of this latter
automorphism has at most two geometric points above it, and an
odd-order permutation fixes both. This contradicts the free action
of the nonidentity deck transformation on an \emph{\'etale} cover.
All implications used in this finite test have now been proved.
\end{proof}
\paragraph{Check record.} Checked only in this writing; GPT-6 Astra Pro; 2026-11-18; complete finite group arithmetic and second plain-integer verifier, exact integer norms; effort not separately exposed.


\paragraph{Check record.}
Checked only in this writing: GPT-6 Astra Pro, finite-checks reader,
18 November 2026. We wrote independent finite-field and divisor-class
arithmetic, enumerated the full group, checked the full operation and
character tables, and computed integer resultants. A second program
recomputed every addition using plain integer polynomial arithmetic
and all nine norms by integer Bareiss determinants. The effort setting
is not separately exposed by the runtime. Inputs, outputs, group
tables and cover definitions are in
\path{code/s2_f8_witness.py} and \path{code/s2_f8_witness.json}.

\subsection{A finite nonnilpotent test of the relative configuration sequence}

\begin{claim}{A.51}{Proved}{An explicit nonnilpotent cusp quotient and a dominating configuration quotient}
Let $\Gamma=P_2(S_2)$ be the pure braid group on two points of a
closed oriented genus-two surface, and let
$\Lambda=\langle a,b,c,d\rangle$ be its first fibre. For $p=2$,
$r=5$, and $A=\mathbf Z/3$, the fibre has a quotient
\[
 Q=\mathbf F _5^2\rtimes\mathbf Z/3,
 \qquad \sigma|_{\mathbf F _5^2}=
 S=\begin{pmatrix}0&-1\\1&-1\end{pmatrix},
 \qquad |Q|=75,
\]
in which the normal cusp subgroup is $\mathbf F _5^2$ and is
noncentral. An explicitly specified finite quotient $T$ of $\Gamma$
has order $9\cdot5^{40}$ and a normal exponent-five subgroup of
order $5^{40}$ and nilpotency class three with quotient $A\times A$. Its
restriction to $\Lambda$ dominates $Q$. Consequently the two maps
factor through the respective cuspidally pro-five groups over their
pro-odd proper quotients. This is a nonnilpotent finite test of the
relative configuration sequence.
\end{claim}

\begin{proof}
We use Bellingeri's Theorem 1.2 \citep[Theorem~1.2]{Bellingeri2004} for a closed oriented surface of
genus $g\ge1$, specialized to $g=n=2$. This gives generators
$a,b,c,d,\sigma$ for the full braid group. Its four same-generator,
four different-handle, two mixed same-handle, and one surface
relations are transcribed in the input certificate. Put
$t=[a,b^{-1}][c,d^{-1}]=\sigma^2$ and
$A_i=\sigma x_i\sigma^{-1}$, where
$(x_1,x_2,x_3,x_4)=(a,b,c,d)$. Index-two Reidemeister--Schreier
rewriting with representatives $1,\sigma$ gives a presentation of
$\Gamma$ with the eight generators $x_i,A_i$. Rewriting each of
the eleven relations and its $\sigma$-conjugate gives twenty-two
relators. During the rewriting, a positive $\sigma$-edge from the
second coset contributes $t$, a negative $\sigma$-edge from the
first contributes $t^{-1}$, and the other two edges contribute
one. This completely specifies the presentation computation.

Let $K$ be the kernel of $\Lambda\to\mathbf Z/3$ given by
$a\mapsto1$ and $b,c,d\mapsto0$. The cover of the compact surface
is the degree-three cover $C'\to C$; its punctured fibre has three
punctures and free group $K$ of rank ten. Representatives
$1,a,a^2$ give a free basis by all positively oriented Cayley edges
except the first two $a$-edges. The program retains the ten resulting
free words. Define
\[
 P=K/\bigl(\gamma_3(K)K^5\bigr).
\]
This is the free class-two exponent-five group on ten generators,
of order $5^{55}$. A concrete model is
\[
 P=V\oplus\bigwedge^2V,
 \quad (v,z)(w,u)=(v+w,z+u+\tfrac12v\wedge w),
 \quad V=\mathbf F _5^{10}.
\]
Every automorphism used below is specified by two matrices $(M,B)$,
of sizes $10\times10$ and $45\times10$, respectively, acting by
\[
              (v,z)\longmapsto(Mv,(\bigwedge^2M)z+Bv).
                                                               \tag{A.51.1}
\]
The eight pairs of matrices in the certificate thus specify elements
of a concrete finite matrix group, not formal symbols.

For the first four generators, use the action on $K$ by conjugation
with $x_i$. For the second four, the point-push words are as follows:
write $u_i=t x_i^{-1}t x_i t^{-1}$ and $w_i=u_i^{-1}t$.
They send $x_i$ to $t x_i t^{-1}$. If $i$ is odd, the partner
$x_{i+1}$ goes to $t x_{i+1}t^{-1}u_i$; if $i$ is even,
$x_{i-1}$ goes to $x_{i-1}t^{-1}$. Earlier handles are fixed,
and later handles are conjugated by $w_i$. These formulas follow
by rewriting Bellingeri's same-handle and different-handle relations.
For an independently checkable safeguard, we do not rely on that
derivation alone: substituting the eight finite matrices into all
twenty-two relators gives the identity, exactly over $\mathbf F _5$.
Thus they define a homomorphism
$\rho:\Gamma\to\operatorname{Aut}(P)$ even independently of
any statement about automorphisms of a free group.

Append the two deck coordinates
$\delta:\Gamma\twoheadrightarrow(\mathbf Z/3)^2$, assigning
$(1,0)$ to $a$, $(0,1)$ to $A_1$, and zero to the other six
generators. Put $T=\operatorname{im}(\rho,\delta)$.
This defines a finite group explicitly as the subgroup generated
by eight specified matrices with their specified deck coordinates.
Let $H_0=\ker\delta$. Geometrically it is the fundamental group of
$C'\times C'$ minus the three deck graphs. Representatives
$a^iA_1^j$ for $0\le i,j<3$ give seventy-two Schreier generators
of $H_0$. Their linear parts are all $1+N_h$. If $W\subset V$
is the two-dimensional space spanned by the three cusp classes,
the certificate checks
\[
 N_h(V)\subset W,\qquad N_h(W)=0
 \quad\hbox{for all seventy-two generators}.
\]
It checks all $72^2$ identities $N_hN_{h'}=0$ as well; the span of
the $N_h$ has dimension four. Therefore the linear image of $H_0$
is an elementary abelian five-group. The kernel of
$\operatorname{Aut}(P)\to\operatorname{GL}(V)$ in (A.51.1) is
$\operatorname{Hom}(V,\bigwedge^2V)$, also a five-group.
It follows that $\rho(H_0)$ is a five-group. We now determine its
order, without enumerating all of its elements.

Here is the finite logarithm argument used for this purpose. If $I$
is an associative algebra over $\mathbf F _5$ with $I^5=0$, the
polynomials
\[
 \log(1+D)=D-\frac{D^2}{2}+\frac{D^3}{3}-\frac{D^4}{4},
 \qquad
 \exp(X)=1+X+\frac{X^2}{2}+\frac{X^3}{6}+\frac{X^4}{24}
\]
are inverse. Multiplication is transported to the BCH polynomial
through bracket degree four; its denominators divide $24$, so they
are invertible in the field. For matrices $g_i=\exp(X_i)$, their
generated group is exactly $\exp(L)$, where $L$ is the Lie algebra
generated by the $X_i$. To prove the reverse inclusion, which is the
nontrivial direction, first work in the free nilpotent Lie algebra
of class four on the symbols $X_i$. Give a bracket word its ordinary
bracket length. The logarithm of an iterated group commutator equals
its corresponding Lie bracket modulo larger bracket lengths.
Integer powers supply all five scalar multiples. Products add the
terms of the current smallest length modulo larger lengths.
Consequently every homogeneous linear combination of bracket length
$j$ can be realized modulo terms of length $j+1$. Starting at $j=1$
and correcting successively through $j=4$ realizes the exponential
of every Lie element exactly. Mapping this free algebra to $I$
proves the assertion for the matrices in question. In particular,
$|\langle g_i\rangle|=5^{\dim L}$. The same argument identifies
successive lower central factors on the group and Lie sides.

To check the hypothesis, put $E_h=\bigwedge^2(1+N_h)$. The
operator $E_h-1$ lowers the filtration
\[
 \bigwedge^2 V\supset W\wedge V\supset\bigwedge^2 W\supset0
\]
by one step. The operators $N_h$ lower $V\supset W\supset0$
by one step. The full augmentation matrices have blocks
$\left(\begin{smallmatrix}N_h&0\\ B_h&E_h-1\end{smallmatrix}\right)$.
A product of five such matrices is zero: a nonzero block path uses
at most one $N$ step, at most one $B$ step, and at most two $E-1$
steps. Thus the algebra $I$ they generate satisfies $I^5=0$.

The certificate computes the seventy-two logarithms, whose linear
span has dimension twenty, and closes their span under brackets.
The resulting Lie algebra has dimension forty. Its lower central
series has dimensions
\[
                   40,\quad24,\quad4,\quad0.
\]
The retained basis and all bracket reductions are exact over
$\mathbf F _5$. Hence $\rho(H_0)$ has order $5^{40}$, exponent
five, and nilpotency class exactly three. Since the appended deck
coordinates are surjective,
\[
             |T|=9\cdot5^{40}
                =81854523159563541412353515625.
\]
A second computation uses no logarithms. It constructs a common
invariant flag of dimensions $19,41,47,55$, makes all seventy-two
matrices unitriangular, and collects actual matrix commutators in
increasing triangular pivot order. It obtains forty generators of
order five. All pair commutators collect, as do all original
generators. Collection terminates because every commutator raises
the triangular pivot filtration. Distinct ordered products have
distinct first unequal pivot coordinates, so their $5^{40}$ normal
forms are distinct. This independently confirms the order.

This construction also realizes the finite-core step of the relative
completion proof. The subgroup
$J_0=\ker(H_0\to\operatorname{Aut}(P))$ is invariant under
$\Gamma$, since it is the intersection of $H_0$ with a homomorphism
kernel. Thus its core over the nine cosets is already $J_0$.
The subgroup $K$ is preserved by every generator, as the displayed
words preserve its abelian degree-three quotient. Its quotient $P$
is characteristic. Hence the core just used has precisely the
cover, shear, and normality interpretation required by that proof.

It remains to check domination of the nonnilpotent fibre quotient.
Define $q:\Lambda\to Q$ by
\[
 a\mapsto(0,\sigma),\qquad b\mapsto((1,0),1),
 \qquad c,d\mapsto1.
\]
The matrix $S$ has order three and no invariant line over
$\mathbf F _5$. The images generate all seventy-five elements of
$Q$. The nonzero image of $t$ and its $S$-translate span the
normal subgroup $\mathbf F _5^2$, so this is exactly the normal
cusp subgroup. In particular $Q$ is nonnilpotent. The restriction
of $q$ to $K$ factors through a rank-two linear map
$\lambda:V\twoheadrightarrow\mathbf F _5^2$; its matrix is
provided in the certificate.

Suppose $x\in\Lambda$ is killed by $(\rho,\delta)$. Its first
deck coordinate shows that $x\in K$. Its action on $P$ is the
inner action of its image. The centre of $P$ is exactly
$\bigwedge^2V$: if $v\ne0$, choose $w$ independent of $v$, so
$v\wedge w\ne0$. Equivalently the certificate computes rank ten
for $v\mapsto(w\mapsto v\wedge w)$. Thus $x$ has zero image in
$V$ and is killed by $\lambda$, hence by $q$. This proves the
required kernel inclusion and therefore domination.

Finally the kernel of the proper-product quotient maps into
$\rho(H_0)$, a five-group. The maps therefore kill the kernels of
the maximal pro-five quotients of the corresponding cusp groups.
This proves the asserted factorization. No nilpotent quotient of
the entire fibre has been substituted for $Q$.
\end{proof}
\paragraph{Check record.} Checked only in this writing; GPT-6 Astra Pro; 2026-11-18. All 22 relations, 72 Schreier generators, 80 independent tensor-algebra word evaluations, and the full kernel argument checked. Exact order checked by Lie closure and independent actual matrix collection; effort not separately exposed.


\paragraph{Check record.}
Checked only in this writing: GPT-6 Astra Pro, finite-checks reader,
18 November 2026, source inspection, explicit presentation rewriting,
matrix computation, and the complete kernel argument above.
The effort setting is not separately exposed by the runtime.
A second program independently recomputed all eighty generator images
in a truncated tensor algebra and verified all twenty-two relations
as full $55\times55$ matrices. The source, all words, all matrices,
all shear tests and the factor map
are retained in \path{code/s3_nonnilpotent.py} and its JSON output.
The two independent order methods and their full bases are in
\path{code/s3_group_order.py} and \path{code/s3_order_collect.py},
with their JSON certificates.

\subsection{Independent reruns of the four retained finite claims}

\begin{claim}{A.52}{Proved}{Agreement of the independent finite certificates}
The finite assertions in A.7, A.12, A.27, and A.29 have been checked
by a newly written program from their mathematical inputs, before
opening their original program sources, and by rerunning those original
programs. The two computations agree. More precisely:
\begin{enumerate}
\item The two integral rank-eighty-two homology matrices of A.7 are
integrally conjugate by a supplied matrix of determinant $\pm1$.
Their restricted matrices are
$\left(\begin{smallmatrix}1&3\\0&1\end{smallmatrix}\right)$ and
$\left(\begin{smallmatrix}1&0\\-3&1\end{smallmatrix}\right)$,
and their product has the claimed nonzero minus-one eigenvector
modulo five.
\item The five transvections in A.12 agree entry by entry, generate
$51840$ matrices, and the specified word has exact order five.
\item The finite-field, point, orbit, and group-order calculations
in A.27 agree: $129$ points over $\mathbf F _{128}$, $126$ of exact
degree seven, and $|\operatorname{GL}_4(\mathbf F _3)|=24261120$.
\item The group in A.29 has order $576$, the stated generators,
commutator subgroup of order $64$, trivial complement core, and the
specified order-seven automorphism. All $576^2$ product
compatibilities hold.
\end{enumerate}
This statement concerns the finite certificates. The profinite and
geometric consequences still require the separate proofs in the
corresponding statements.
\end{claim}

\begin{proof}
For A.7 we computed the induced action on the graph-cycle space of
the Heisenberg cover using a breadth-first tree that uses positive
generators only. Its eighty-two non-tree edges give coordinates.
Substitution of the point-push words in each closed lifted word gives
the integral action. The original program uses a different, signed
tree. We therefore evaluated every original basis word in our basis,
obtaining a matrix $U$ with determinant $\pm1$, and verified
$M_{
m new}U=UM_{
m old}$ for both complete matrices. This
explains the only coordinate differences; the mathematical data agree.

For A.12 the independent computation starts from the geometric
transvection formula on $\mathbf F _3^4$, with directions
$e_1,e_2,e_3,e_4,e_1+e_3$ and coefficient two. Exhaustive multiplication
gives $51840$ distinct matrices and five successive powers verify
the exact order of the indicated word. The original program was
separately rerun to check its derivation from full Schreier homology.
We do not describe the independent calculation as a second derivation
of all of its free-group words.

For A.27 we implemented polynomial-basis binary-field arithmetic,
verified $a^{127}=1$ for every nonzero element of the proposed
field, enumerated its curve points, and computed the full seven-point
Frobenius orbit. For A.29 we independently enumerated all elements
of $\mathbf F _{64}\rtimes\mathbf Z/9$, all products required by
the automorphism test, the generator orbit and the conjugate
complement intersection. The complete outputs and comparisons are
retained. All assertions are finite identities or exhaustive finite
enumerations, so their retained programs constitute reproducible
certificates for the assertions just listed.
\end{proof}
\paragraph{Check record.} Also checked in this writing; GPT-6 Astra Pro; 2026-11-18; independentcode, reruns, exact comparisons; effort not separately exposed.


\paragraph{Check record.}
Also checked in this writing, to the finite scope just stated:
GPT-6 Astra Pro, finite-checks reader, 18 November 2026. The blind
program hash, the precise scope of independence, the unmodified
rerun inputs, all outputs, and the comparison are recorded in
\path{S6_BLIND_PROTOCOL.json}, \path{code/independent_s6.py},
and \path{code/s6_comparison.json}.

\subsection{Why base-Jacobian support cannot settle finite prime sets}

\begin{claim}{A.53}{Proved}{Abelian collision candidates at every geometric point for a finite prime set}
Let $C/\mathbb F_q$ be smooth, proper, and geometrically connected
of genus $g\ge2$, and let $J$ be its Jacobian.  Let $\Sigma$ be
any finite set of primes, possibly empty and possibly containing
$p=\operatorname{char}\mathbb F_q$.  Put $k_n=\mathbb F_{q^n}$
and
\[
 M_n=\#J(k_n)_{\Sigma},
\]
where the subscript denotes the product of the primary subgroups
at the primes in $\Sigma$.  There is a constant $c=c(J,\Sigma)>0$
such that, for every $n\ge1$,
\[
                       M_n\le c n^{2g}.                 \tag{A.53.1}
\]
For every $P\in C(k_n)$, let
\[
 N_{n,P}=\#\{Q\in C(k_n):[Q-P]_{\Sigma}=0\}.
\]
Then, uniformly in $P$,
\[
 \left|N_{n,P}-\frac{q^n+1}{M_n}\right|
       \le 2\left(g-1+\frac1{M_n}\right)q^{n/2}.        \tag{A.53.2}
\]
In particular, for every $P\in C(\overline{\mathbb F}_q)$ there
are infinitely many distinct $Q$ for which $[Q-P]$ has order
supported outside $\Sigma$.  Consequently
\[
 B_C^{\Sigma}:=
 \{P\in C(\overline{\mathbb F}_q):
       \exists Q\ne P,\ [Q-P]\text{ is }\Sigma'\text{-primary}\}
       =C(\overline{\mathbb F}_q).                      \tag{A.53.3}
\]
Thus this base-Jacobian candidate set cannot be supported over a
proper algebraic subfield of $\overline{\mathbb F}_q$ for any
finite $\Sigma$, in particular for a set of two primes.
The assertion concerns abelian collision candidates.  It does
not assert a collision in the full pro-$\Sigma$ group or change
the Open status of A.28.
\end{claim}

\begin{proof}
\medskip\noindent\textbf{A uniform primary bound, including the characteristic prime.}
For a prime $\ell$, let $T_\ell J$ mean the Tate module of
geometric torsion points.  Its rank is $d_\ell=2g$ for
$\ell\ne p$ and the $p$-rank $d_p=f_J\le g$ for $\ell=p$.
At the characteristic prime only the \'etale part contributes:
\[
 J[p^a](\overline{\mathbb F}_q)\simeq
                  (\mathbb Z/p^a)^{f_J}.
\]
Let $A_\ell\in\operatorname{GL}_{d_\ell}(\mathbb Z_\ell)$
be the action of the $q$-power automorphism on this Tate module.
It is invertible also at $p$, since that automorphism is an
automorphism of $\overline{\mathbb F}_q$ and acts bijectively
on geometric torsion.  If $d_\ell=0$, all the corresponding
primary groups of rational points are trivial and we omit this
factor below.

For every $n\ge1$, $A_\ell^n-I$ is invertible over
$\mathbb Q_\ell$.  Indeed a nonzero fixed vector, scaled to
a primitive vector of $T_\ell J$, would give fixed torsion
points of arbitrarily large $\ell$-power order.  They would
all lie in the finite group $J(k_n)$, a contradiction.
Apply $A_\ell^n-I$ to the exact sequence
\[
 0\longrightarrow T_\ell J\longrightarrow
 T_\ell J\otimes\mathbb Q_\ell\longrightarrow
 J[\ell^\infty](\overline{\mathbb F}_q)\longrightarrow0.
\]
The resulting kernel-cokernel sequence gives
\[
 J(k_n)[\ell^\infty]\simeq
 T_\ell J/(A_\ell^n-I)T_\ell J,
 \qquad
 \#J(k_n)[\ell^\infty]
      =\ell^{v_\ell\det(A_\ell^n-I)}.                  \tag{A.53.4}
\]
These assertions use the \'etale $p$-adic module, not a
$2g$-dimensional $p$-adic module containing nonunit crystalline
eigenvalues.

Choose $e_\ell\ge1$ for which
\[
 B_\ell:=A_\ell^{e_\ell}\equiv I
 \pmod {\ell}\quad(\ell\text{ odd}),\qquad
 B_2\equiv I\pmod4.
\]
Such an exponent exists by finiteness of the corresponding
matrix group over $\mathbb Z/\ell$ or $\mathbb Z/4$.
Set $b_\ell=v_\ell\det(B_\ell-I)<\infty$.
We claim that, for every positive integer $m$,
\[
 v_\ell\det(B_\ell^m-I)
                =b_\ell+d_\ell v_\ell(m).              \tag{A.53.5}
\]
Here no diagonalization or semisimplicity is needed.  Write
$S_m(B)=I+B+\cdots+B^{m-1}$.  If $\ell\nmid u$, then
$S_u(B)\equiv uI\pmod\ell$ is invertible.  If $B\equiv I$
modulo $\ell$ for odd $\ell$, the binomial expansion gives
\[
 S_\ell(B)=\ell U,\qquad U\equiv I\pmod\ell.
\]
To check the divisibility of the remaining terms, write
$B=I+\ell X$: the coefficient of $(\ell X)^j$ in
$S_\ell(B)$ is $\binom\ell{j+1}$; for $1\le j\le\ell-2$
it is divisible by $\ell$, and the final term has valuation
at least $\ell-1\ge2$.  For $\ell=2$ use $B=I+4X$,
so $S_2(B)=2(I+2X)$.  The same congruence holds for every
$B^{\ell^i}$.  For $m=\ell^a u$ with $\ell\nmid u$,
\[
 S_m(B)=S_u(B^{\ell^a})
              \prod_{i=0}^{a-1}S_\ell(B^{\ell^i}).
\]
The first factor has unit determinant and every other factor
has determinant valuation $d_\ell$.  Since
$B^m-I=(B-I)S_m(B)$, this proves (A.53.5).

The inclusion $J(k_n)\subset J(k_{e_\ell n})$, followed
by (A.53.4)--(A.53.5), now gives a bound valid for every $n$,
without requiring $e_\ell\mid n$:
\[
 \#J(k_n)[\ell^\infty]
       \le \ell^{b_\ell}\ell^{d_\ell v_\ell(n)}.
\]
Because $\Sigma$ is finite, we can multiply these inequalities.
Using $d_\ell\le2g$ gives
\[
 M_n\le
 \left(\prod_{\substack{\ell\in\Sigma\\d_\ell>0}}
                           \ell^{b_\ell}\right)
 \prod_{\ell\in\Sigma}\ell^{2g v_\ell(n)}
       \le c n^{2g}.
\]
For an empty prime set take $c=1$.  This proves (A.53.1).

\medskip\noindent\textbf{The connected abelian cover that tests the retained part.}
Fix $n$ and $P\in C(k_n)$.  Let $\varphi$ denote the
$q$-power Frobenius endomorphism of $J$ and consider
\[
          \mathcal L_n=\varphi^n-I:J\longrightarrow J.
\]
This is a finite \'etale surjective homomorphism.  Indeed its
differential is $-I$, so it is \'etale and quasi-finite; its
proper image is open and closed in the connected target, hence
the whole target; properness then makes it finite.  Its kernel
is the constant finite group $J(k_n)$.  Thus it is a connected
abelian Galois cover over $k_n$, including any $p$-primary
part of its degree.

Take the intermediate cover obtained by quotienting its deck
group by $J(k_n)_{\Sigma'}$.  Its deck group is
$J(k_n)_\Sigma$ and its degree is $M_n$.  Pull it back along
the Abel map
\[
                 j_P:C_{k_n}\longrightarrow J_{k_n},
                      \qquad Q\longmapsto[Q-P],
\]
and denote the resulting cover by $D_{n,P}\to C_{k_n}$.
It is geometrically connected: over an algebraic closure, the
Abel map induces a surjection of fundamental groups onto that
of the Jacobian, indeed an isomorphism on abelianizations.
The exact all-prime statement is
\citet[Proposition~9.1]{MilneJac}; its characteristic-prime
argument is given in Lemmas~9.4--9.5 and the paragraph on
Witt vectors following them, so no prime-to-$p$ restriction
is being inserted here.  It follows that $D_{n,P}$ is a
smooth proper geometrically connected curve and, by the
unramified Hurwitz formula,
\[
                 g(D_{n,P})=1+M_n(g-1).                 \tag{A.53.6}
\]

For $Q\in C(k_n)$ choose $z\in J(\overline{\mathbb F}_q)$
with $\varphi^n z-z=[Q-P]$.  The arithmetic Frobenius
displacement of $z$ in this fibre of the full Lang cover is
exactly $[Q-P]\in J(k_n)$.  Consequently the displacement
in the intermediate cover is $[Q-P]_\Sigma$.
The fibre of $D_{n,P}$ over $Q$ has a rational point if and
only if this displacement is zero.  When it is zero, all
$M_n$ points of that fibre are rational, since the deck group
is constant over $k_n$.  Therefore
\[
                   \#D_{n,P}(k_n)=M_n N_{n,P}.          \tag{A.53.7}
\]
The Weil bound applied to the proper connected curve
$D_{n,P}$, followed by (A.53.6)--(A.53.7), proves (A.53.2).
This is the curve bound recalled in
\citet[Theorem~11.1 and Corollary~11.4]{MilneJac}.

\medskip\noindent\textbf{Every geometric point has infinitely many partners.}
For a fixed geometric point $P$, restrict $n$ to multiples
of a field of definition degree of $P$.  Equations
(A.53.1)--(A.53.2) give
\[
 N_{n,P}\ge \frac{q^n}{c n^{2g}}-2gq^{n/2},
\]
which tends to infinity along these multiples.  Thus the
union of these finite sets of partners is infinite, and
removing the one point $P$ does not change that conclusion.
This proves (A.53.3).

For clarity, $C(\overline{\mathbb F}_q)$ is not contained
in $C(F)$ for any proper algebraic subfield
$F\subsetneq\overline{\mathbb F}_q$.  Its finite constant
subfields are restricted in some prime-power degree: write
$F$ as the union of its finite subfields and choose a prime
$t$ for which the permitted $t$-adic valuations of the
extension degrees are bounded.  If no such prime existed,
the least-common-multiple property of these degrees would
give every finite extension and hence the entire algebraic
closure.  For all sufficiently large $a$, the Weil bound
gives a point of exact degree $t^a$ over the field of
definition: indeed
\[
 \#C(\mathbb F_{q^{t^a}})>
 \#C(\mathbb F_{q^{t^{a-1}}}),
\]
and all proper intermediate subfields containing $\mathbb F_q$
are contained in $\mathbb F_{q^{t^{a-1}}}$.  Such a point is
not $F$-rational.  If necessary replace $F$ by its compositum
with $\mathbb F_q$, which is still proper: adjoining a finite
field changes any bounded prime-adic degree by only a finite
amount.  Thus (A.53.3) is an actual obstruction to proper-field
support, rather than just a cardinality observation.

Finally the choice of $Q$ has been made separately for a fixed
curve and constant extension.  Repeating the argument on each
finite cover produces many candidates on that individual
cover, but neither compatibility with its maps to other
covers nor a fixed base partner is obtained.  A collision of
full point sections requires a compatible point in the entire
profinite path torsor.  No such passage is made in this proof.
\end{proof}
\paragraph{Check record.} Checked only in this writing, 18 November 2026, GPT-6 Astra (OpenAI), ChatGPT Pro mode; separate effort setting not exposed. Complete proof written and read by its named authoring reader; see CHECK\_EVENTS.tsv for additional second readings and computation scope.


\subsection{An exact all-cover criterion for the arithmetic outer kernel}

\begin{claim}{A.54}{Proved}{Auxiliary-prime Frobenius orders and the outer arithmetic kernel}
Let $C/k$ be a smooth proper geometrically connected curve of
genus at least two over $k=\mathbb F_q$, let $O\in C(k)$,
and let $\varnothing\ne\Sigma\subseteq\{\text{primes}\}\setminus\{p\}$,
where $p=\operatorname{char}k$.  Fix $r\in\Sigma$ and put
$\Delta=\pi_1(C_{\bar k})^\Sigma$.  The section at $O$
defines a based action
\[
 \alpha:G_k\longrightarrow\operatorname{Aut}(\Delta)
\]
and its associated outer action $\rho$.  Identify $G_k$ with
$\widehat{\mathbb Z}$ using the $q$-power arithmetic Frobenius.

Let $N$ range over all open normal subgroups of $\Delta$ stable
under $\alpha(G_k)$.  The subgroup $N\rtimes s_O(G_k)$ of
the arithmetic geometrically pro-$\Sigma$ group defines a
pointed cover $D_N\to C$ over $k$.  Write $\bar A_N$ for
Frobenius on
$H^1_{\mathrm{et}}((D_N)_{\bar k},\mathbb F_r)$.
For each omitted prime $t\notin\Sigma$, set
\[
 b_N(t)=v_t\bigl(\operatorname{ord}(\bar A_N)\bigr),
 \qquad B_t=\sup_N b_N(t)\in\mathbb N\cup\{\infty\}.
\]
Then
\[
 \ker(\rho)\cap\mathbb Z_t
       =\ker(\alpha)\cap\mathbb Z_t
       =t^{B_t}\mathbb Z_t,                             \tag{A.54.1}
\]
where $t^\infty\mathbb Z_t$ means $0$.  Moreover,
\[
 \ker(\rho)=\prod_{t\notin\Sigma}t^{B_t}\mathbb Z_t
       \ \subseteq\widehat{\mathbb Z},                 \tag{A.54.2}
\]
with zero components at the primes retained in $\Sigma$.
Thus the outer action is faithful if and only if $B_t=\infty$
for every omitted prime $t$.

If $B_t<\infty$ for even one omitted prime, the exceptional
image $E_C^\Sigma$ is infinite: all closed points of degrees
$t^a$ with $a>B_t$ are exceptional.  In particular empty
$E_C^\Sigma$ implies the stated unbounded-order condition
for every omitted prime.  No converse from outer faithfulness
to point-section injectivity is asserted.

If $C(k)$ is empty, the criterion applies after a finite constant
extension having a rational point.  Faithfulness of the original
outer action is equivalent to faithfulness on that open arithmetic
subgroup, and an infinite exceptional image over the extension
descends to an infinite exceptional image over $k$.
\end{claim}

\begin{proof}
\medskip\noindent\textbf{The pointed covers form a cofinal system.}
Every connected finite geometric cover descends to a finite
extension of the finite field: its equations and the finite
map use finitely many algebraic coefficients.  Accordingly an
open normal subgroup $N_0$ has a finite orbit under $G_k$.
The intersection of that orbit is open normal, stable under
$G_k$, and contained in $N_0$.  These intersections remain
pro-$\Sigma$ quotients, since their finite quotient embeds
in a finite product of finite $\Sigma$-groups.  The indicated
subgroups $N$ are therefore cofinal among open normal
subgroups of $\Delta$ and have intersection $1$.

The subgroup $N\rtimes s_O(G_k)$ defines a connected finite
\'etale cover with a specified rational lift $O_N$ of $O$.
Its geometric deck group is $G_N=\Delta/N$; it need not be
a constant deck group over $k$.  Its geometric fundamental
group with its full pro-$\Sigma$ topology is $N$.  To see
the topology assertion explicitly, let $\widetilde N$ be
the corresponding subgroup of $\pi_1(C_{\bar k})$, and let
$J$ be the kernel of a finite $\Sigma$-quotient of
$\widetilde N$.  Intersect its conjugates over the finite
group $\pi_1(C_{\bar k})/\widetilde N$.  The resulting
normal quotient of $\pi_1(C_{\bar k})$ is an extension
of a finite product subgroup of the conjugate
$\widetilde N/J$ by $\Delta/N$.  It is a finite
$\Sigma$-group and detects the prescribed quotient of
$\widetilde N$.  This proves
$\widetilde N^\Sigma\simeq N$.

In particular the $r$-adic abelianization and its cohomological
dual on $N$ are the full Tate module of $J_{D_N}$ and
$H^1_{\mathrm{et}}((D_N)_{\bar k},\mathbb Z_r)$.
This follows also from the Abel comparison
\citep[Proposition~9.1, Corollary~9.6]{MilneJac}.
The $G_k$-action on this cohomology is the action induced
by $\alpha$ on $N$.  These modules are free of rank
$2g(D_N)$ and reduce modulo $r$ to the displayed
$\mathbb F_r$-cohomology.

\medskip\noindent\textbf{Deck transformations act faithfully on rational cohomology.}
Put
\[
 V_N=H^1_{\mathrm{et}}((D_N)_{\bar k},\mathbb Q_r).
\]
The geometric deck representation
$\theta_N:G_N\to\operatorname{GL}(V_N)$ is faithful.
Indeed, if a nonidentity deck transformation $u$ acted
trivially, so would the cyclic group $H=\langle u\rangle$.
Its action is free, since $D_N\to C$ is \'etale.  The
quotient $D_N/H$ is therefore a smooth proper curve and
\[
 g(D_N/H)=1+\frac{g(D_N)-1}{|H|}<g(D_N).
\]
On the other hand finite \'etale descent and averaging over
$\mathbb Q_r$ give
\[
 H^1_{\mathrm{et}}((D_N/H)_{\bar k},\mathbb Q_r)
       \simeq V_N^H.
\]
For example, this follows from the finite-cover
Hochschild--Serre sequence
\citep[Theorem~14.9]{MilneEtale}; its positive-degree
finite-group cohomology is killed by $|H|$ and vanishes
after tensoring with $\mathbb Q_r$.  This remains valid
when $r$ divides $|H|$, since $|H|$ is invertible in
$\mathbb Q_r$.  Trivial action would imply
$2g(D_N/H)=2g(D_N)$, a contradiction.  This proves the
faithfulness assertion.

\medskip\noindent\textbf{All-cover cohomology detects the based action.}
Suppose $h\in G_k$ acts trivially on every $V_N$.
Functoriality of the deck action gives, for every $x\in\Delta$,
\[
 \theta_N(\alpha_h(x)\bmod N)
       =A_N(h)\theta_N(x\bmod N)A_N(h)^{-1}
       =\theta_N(x\bmod N),
\]
where $A_N(h)$ is the action on $V_N$.
One may use inverse pullback for the deck action so that
$\theta_N$ is a homomorphism; the displayed equivariance
then has the stated order.  Faithfulness of $\theta_N$
shows that $\alpha_h(x)x^{-1}\in N$ for every $N$.
Cofinality yields $\alpha_h(x)=x$ for every $x$.
Conversely the identity automorphism of $\Delta$ is the
identity on each $N$ and on its first cohomology.  Hence
\[
 \ker\alpha=
  \bigcap_N\ker\bigl(G_k\longrightarrow
           \operatorname{GL}(V_N)\bigr).                \tag{A.54.3}
\]

\medskip\noindent\textbf{The omitted-prime calculation.}
Fix $t\notin\Sigma$, so $t\ne r$.
The kernel of
\[
 \operatorname{GL}_d(\mathbb Z_r)
           \longrightarrow\operatorname{GL}_d(\mathbb F_r)
\]
is pro-$r$: its successive congruence quotients are
additive subgroups of matrix spaces over $\mathbb F_r$.
A pro-$t$ subgroup intersects it trivially.  Thus the
kernel on $\mathbb Z_t\subset G_k$ of the action on
integral or rational $r$-adic first cohomology is the
same as its kernel modulo $r$.
The finite cyclic image generated by $\bar A_N$ has a
$t$-Sylow subgroup of order $t^{b_N(t)}$.  The canonical
$\mathbb Z_t$ factor maps onto precisely that subgroup,
so its kernel is $t^{b_N(t)}\mathbb Z_t$.
Taking the intersection in (A.54.3) proves
\[
                  \ker\alpha\cap\mathbb Z_t
                       =t^{B_t}\mathbb Z_t.
\]

If $K_t=\ker\rho\cap\mathbb Z_t$, then
$\alpha(K_t)\subset\operatorname{Inn}(\Delta)$.
The latter is a pro-$\Sigma$ group, being the continuous
image of $\Delta$ under its inner action.  The former
is pro-$t$.  Since $t\notin\Sigma$, a profinite group
that is both pro-$t$ and pro-$\Sigma$ is trivial.
Consequently $K_t\subset\ker\alpha$, and the reverse
inclusion is automatic.  This proves (A.54.1) without
assuming centrelessness of $\Delta$.

\medskip\noindent\textbf{Every retained prime factor is already detected on the base.}
Let $\ell\in\Sigma$.  The action of $\mathbb Z_\ell$
on $T_\ell J_C$ is faithful.  Otherwise its image would
be finite, since a nonzero closed subgroup of
$\mathbb Z_\ell$ is open.  The closure of the cyclic
Frobenius subgroup in
$\operatorname{GL}_{2g}(\mathbb Z_\ell)$ has an open
pro-$\ell$ congruence subgroup, so finiteness of its
$\ell$-Sylow image would make the entire closure finite.
Some positive power of Frobenius would then act as the
identity on $T_\ell J_C$.  This would make the finite
group $J_C(\mathbb F_{q^m})$ contain $\ell$-primary
points of arbitrarily large order, a contradiction.
An inner automorphism of $\Delta$ acts trivially on
its abelianization, and hence on $T_\ell J_C$.
It follows that $\ker\rho\cap\mathbb Z_\ell=0$.

Every closed subgroup of $\widehat{\mathbb Z}$ is the
product of its intersections with the prime factors:
it is stable under multiplication by $\mathbb Z$ and,
by closure, by $\widehat{\mathbb Z}$, so it is stable
under the idempotents projecting to individual factors.
Finite partial sums and closure recover the full product.
Combining this fact with (A.54.1) and the retained-prime
calculation proves (A.54.2).

\medskip\noindent\textbf{A bounded auxiliary order gives actual all-level collisions.}
Suppose $B_t<\infty$.  Equation (A.54.1) says that
$\alpha$ kills $t^{B_t}\mathbb Z_t$.
Its image under $s_O$ is therefore central in
$\Delta\rtimes G_k$: it commutes with $\Delta$ by
triviality of the action and with $s_O(G_k)$ because
$G_k$ is abelian.
Choose $h\in\mathbb Z_t\subset\widehat{\mathbb Z}$
with $t$-adic valuation exactly $B_t$, and set $z=s_O(h)$.
For a closed point $x$ of degree $t^a$, $a>B_t$, this
$h$ acts nontrivially on the geometric points over $x$.
Thus for every lift $\widetilde x$ to the universal
pro-$\Sigma$ cover the geometric images of
$\widetilde x$ and $z\widetilde x$ are distinct, whereas
\[
 D_{z\widetilde x}
        =zD_{\widetilde x}z^{-1}=D_{\widetilde x}.
\]
These are actual decomposition subgroups in the whole
arithmetic group, not just images in one finite quotient.
Every such $x$ is exceptional.  For all sufficiently large
$a$, the Weil bound gives points of exact degree $t^a$,
by comparing the numbers of points over
$\mathbb F_{q^{t^a}}$ and
$\mathbb F_{q^{t^{a-1}}}$.
Hence the exceptional image is infinite.
\citep[Theorem~11.1]{MilneJac}

Finally a closed subgroup of $\widehat{\mathbb Z}$ has
zero intersection with an open subgroup only if it is
zero.  Faithfulness of $\rho$ can consequently be checked
after a finite constant extension.  Equal decomposition
subgroups in its open arithmetic group give commensurable
decomposition subgroups in the original group.  Infinitely
many closed points over a finite extension have infinitely
many images over $k$, since the morphism has finite fibres.
This proves the assertions concerning the absence of an
initial rational point.

The use of $p\notin\Sigma$ above is explicit: it ensures
that every retained prime has a base Tate module of
positive rank.  If $p\in\Sigma$ and the base $p$-rank
is zero, the retained-$p$ part of the argument does not
follow.  The omitted-prime identity (A.54.1) still follows
whenever one fixes an invertible $r\in\Sigma$ and uses
the same cofinal pointed-cover construction, but the
complete formula (A.54.2) is only asserted in the stated
prime-to-$p$ range.  Even in that range, unboundedness
of these arithmetic orders is not inferred merely from
the existence of automorphisms of abstract finite
$\Sigma$-groups, and faithful outer action is not
identified with point-section separation.
\end{proof}
\paragraph{Check record.} Checked only in this writing, 18 November 2026, GPT-6 Astra (OpenAI), ChatGPT Pro mode; separate effort setting not exposed. Complete proof written and read by its named authoring reader; see CHECK\_EVENTS.tsv for additional second readings and computation scope.


\begin{claim}{A.55}{Proved}{No finite prime set is Jacobian-large}
Let $C/\mathbb F_q$ be a smooth proper geometrically connected
curve of positive genus, with Jacobian $J$, and let $\Sigma$
be a finite set of primes, possibly containing the
characteristic.  Then
\[
 G_{\mathbb F_q}\longrightarrow
 \operatorname{Aut}\left(\prod_{\ell\notin\Sigma}T_\ell J\right)
                                                        \tag{A.55.1}
\]
is injective.  Here $T_pJ$ means the \'etale Tate module
of geometric $p$-power torsion points.
In Sa\"idi--Tamagawa's terminology, $\Sigma$ is therefore
not $J$-large.  In particular every finite set of two primes
lies outside the large-set emptiness theorem.  This does
not decide whether its full exceptional image is empty
or infinite.
\end{claim}

\begin{proof}
We first fix the exact definition.  For an abelian variety
$A/k$, \citet[Definition~2.6]{LA_ST2017} calls a prime
set $S$ \emph{$A$-small} when
$G_k\to\operatorname{Aut}(\prod_{\ell\in S}T_\ell A)$
is not injective, and calls $\Sigma$ \emph{$A$-large}
when its complement $\Sigma'$ is $A$-small.  Thus the
asserted injectivity is exactly the negation of
$J$-largeness, rather than a different notion of smallness.

Let $H$ be the kernel of (A.55.1), and suppose $H\ne0$.
Identify $G_k$ with $\widehat{\mathbb Z}$.  A closed
subgroup is stable under multiplication by
$\widehat{\mathbb Z}$, since it is stable under integer
multiplication and is closed.  Applying a prime-factor
idempotent to a nonzero element shows that there exist
a prime $t$ and $b\ge0$ such that
\[
                    t^b\mathbb Z_t\subset H.           \tag{A.55.2}
\]
There is no need to assume $t\in\Sigma$; in particular
this argument also covers a possible characteristic-prime
factor when the $p$-rank is zero.

Put $k_m=\mathbb F_{q^m}$.  For $a\ge b$ the arithmetic
subgroups of $k_{t^a}$ and $k_{t^b}$ are respectively
\[
 G_a=t^a\mathbb Z_t\times\prod_{u\ne t}\mathbb Z_u,
 \qquad
 G_b=t^b\mathbb Z_t\times\prod_{u\ne t}\mathbb Z_u.
\]
They have the same image on the complementary Tate module,
because $G_b=G_a+(t^b\mathbb Z_t)$ and the added subgroup
is killed by (A.55.2).
The action on all complementary-primary geometric torsion
is determined by these Tate modules:
\[
 J[\Sigma'^\infty](\bar k)
    =\bigoplus_{\ell\notin\Sigma}
          \bigl(T_\ell J\otimes\mathbb Q_\ell/\mathbb Z_\ell\bigr).
\]
At $p$ this equality concerns geometric points and hence
only the \'etale part, exactly as in (A.55.1).
The invariant subgroups are therefore equal:
\[
 J(k_{t^a})_{\Sigma'}=J(k_{t^b})_{\Sigma'}
                         \qquad(a\ge b).                \tag{A.55.3}
\]
This proves actual stabilization of these finite groups,
not just a bound for any one complementary prime.

Write $c_0=\#J(k_{t^b})_{\Sigma'}$.
The primary-growth argument in A.53 gives a constant $c$
such that
\[
 \#J(k_{t^a})_\Sigma\le c(t^a)^{2g}.
\]
That argument uses only the Tate modules of a
positive-dimensional abelian variety and therefore also
applies when the present curve has genus one; alternatively
one may repeat its determinant and geometric-series
calculation verbatim.
Combining this with (A.55.3) gives
\[
                   \#J(k_{t^a})\le c_0c(t^a)^{2g}.      \tag{A.55.4}
\]
On the other hand the finite-field Weil formula supplies
algebraic numbers $\alpha_1,\ldots,\alpha_{2g}$, each
of complex absolute value $q^{1/2}$, with
\[
 \#J(k_m)=\prod_{i=1}^{2g}(1-\alpha_i^m)
            \ge (q^{m/2}-1)^{2g}.                       \tag{A.55.5}
\]
For the inequality take complex absolute values of the
positive integer on the left and use the reverse triangle
inequality on every factor.  The exact primary source is
\citet[Theorem~19.1(b),(c)]{CF_MilneAV1986}.
At $m=t^a$ its exponential lower bound contradicts the
polynomial upper bound (A.55.4) as $a\to\infty$.
Therefore $H=0$, which proves the statement by the
definition of $J$-largeness.
\end{proof}
\paragraph{Check record.} Checked only in this writing, 18 November 2026, GPT-6 Astra (OpenAI), ChatGPT Pro mode; separate effort setting not exposed. Complete proof written and read by its named authoring reader; see CHECK\_EVENTS.tsv for additional second readings and computation scope.


\subsection{A limitation of abelian-prime towers in the two-prime problem}

\begin{claim}{A.56}{Proved}{Uniform auxiliary-prime bounds on abelian-$s$ then pro-$r$ towers}
Let $C/\mathbb F_q$ be smooth, proper and geometrically
connected of genus $g\ge2$, with a rational base point $O$.
Let $r,s$ be distinct primes, put $p=\operatorname{char}\mathbb F_q$,
and put $\Sigma=\{r,s\}$.  Either retained prime is allowed
to equal $p$.
Fix an omitted prime $t\notin\{r,s\}$.
Consider geometrically connected finite \'etale covers $D_N\to C$
with a specified rational lift of $O$, normal after geometric
base change, whose geometric deck group $G_N$ has a normal $r$-subgroup
$R$ with $G_N/R$ an abelian $s$-group.  There is a constant
$B=B(C,r,s,t)$, independent of $N$, such that
\[
 v_t\!\left(\operatorname{ord}
   \bigl(\operatorname{Frob}_q\mid
       H^1_{\mathrm{et}}((D_N)_{\bar k},\mathbb F_r)\bigr)
        \right)\le B.                                      \tag{A.56.1}
\]
Equivalently, the same uniform bound holds on every pointed
abelian $s$-cover of $C$ followed by an arbitrary tower with
geometric $r$-group Galois closure, whenever the resulting
cover and its base point descend to $\mathbb F_q$.

Here is an explicit bound.  Define
\[
 c_{r,s}=\sup_{n\ge1}
       \frac{\varphi(s^n)}{\operatorname{ord}_{s^n}(r)},
\]
and put
\[
 M=
 \begin{cases}
  \max\{2g,(2g-2)c_{r,s}\},&r\ne p,\\
  \max\{f_C,\lceil1+(g-1)\dfrac{s}{s-1}c_{r,s}\rceil\},&r=p,
 \end{cases}
\]
where $f_C$ is the $p$-rank of $J_C$,
and let
\[
 b_s(t)=v_t\!\left(\operatorname{ord}
       (\operatorname{Frob}_q\mid H^1_{\mathrm{et}}
                         (C_{\bar k},\mathbb F_s))\right).
\]
The integer $c_{r,s}$ is finite.  With
\[
 e_t(r)=
 \begin{cases}
  v_t(r^{t-1}-1),&t\text{ odd},\\
  v_2(r^2-1)-1,&t=2,
 \end{cases}
\]
one may take
\[
 B=b_s(t)+e_t(r)+v_t(s-1)+\lfloor\log_t M\rfloor.
                                                               \tag{A.56.2}
\]
Thus this entire restricted family of two-prime covers cannot
establish the unbounded-order condition in A.54, or its
retained-$p$ extension in A.57.  No bound
for all finite $\{r,s\}$-covers, and no conclusion about
the full exceptional image $E_C^{\{r,s\}}$, follows from
this statement.
\end{claim}

\begin{proof}
\medskip\noindent\textbf{The elementary order constants are uniform.}
For odd $s$, let $d=\operatorname{ord}_s(r)$ and
$h=v_s(r^d-1)$.  The binomial expansion, or multiplication
of geometric series, gives
\[
 \operatorname{ord}_{s^n}(r)=
  \begin{cases}
   d,&n\le h,\\
   d s^{n-h},&n>h.
  \end{cases}
\]
Consequently
\[
                 c_{r,s}=\frac{s-1}{d}s^{h-1}
                    \qquad(s\text{ odd}).               \tag{A.56.3}
\]
For $s=2$, the prime $r$ is odd.  If
$a=v_2(r-1)$ and $b=v_2(r+1)$, exactly one of $a,b$
equals $1$.  For $m\ge1$ the factorization of
$r^{2^m}-1$ gives
\[
                v_2(r^{2^m}-1)=a+b+m-1.
\]
Together with the valuation for an odd exponent, this gives
\[
                    c_{r,2}=2^{a+b-2}.                  \tag{A.56.4}
\]
These formulas also prove that $c_{r,s}$ is an integer.
For every $n$, the degree
$f_n=\operatorname{ord}_{s^n}(r)$ divides
$\varphi(s^n)=s^{n-1}(s-1)$, and therefore
\[
                      v_t(f_n)\le v_t(s-1).             \tag{A.56.5}
\]

We will use a second elementary estimate.  For any positive
integer $u$ and any $t\ne r$,
\[
                 v_t(r^u-1)\le e_t(r)+v_t(u).           \tag{A.56.6}
\]
For odd $t$, the valuation is zero unless
$\operatorname{ord}_t(r)$ divides $u$; in the latter case
the geometric-series valuation gives equality with
$v_t(r^{t-1}-1)+v_t(u)$, since
$\operatorname{ord}_t(r)$ divides $t-1$.
At $t=2$, an odd $u$ gives $v_2(r-1)$, which is at most
$e_2(r)$; an even $u$ gives
$v_2(r^2-1)-1+v_2(u)$ by the displayed factorization.
Thus (A.56.6) includes the prime $2$.

\medskip\noindent\textbf{All relevant abelian-cover multiplicities are bounded.}
Let $Y\to C$ be a pointed geometrically connected finite
\'etale cover with geometric deck group $A$ an abelian
$s$-group, stable under the arithmetic action determined
by $O$.  Put
\[
             W=H^1_{\mathrm{et}}(Y_{\bar k},\mathbb F_r).
\]
Since $r\nmid |A|$, the algebra $\mathbb F_r[A]$ is a
product of finite fields and $W$ is semisimple over this
algebra.  Its simple factors correspond to the
$r$-power orbits of characters
$\chi:A\to\overline{\mathbb F}_r^{\times}$.
If $\chi$ has order $s^n$, its coefficient field has
degree $f_n$ over $\mathbb F_r$.

Let $K=\ker\chi$.  The quotient $A/K$ is cyclic of order
$s^n$.  Let $K'/K$ be its unique subgroup of order $s$.
Finite \'etale descent and the vanishing of higher
cohomology of the finite $s$-groups with $\mathbb F_r$
coefficients identify
\[
 W^K\simeq H^1_{\mathrm{et}}((Y/K)_{\bar k},\mathbb F_r),
 \quad
 W^{K'}\simeq H^1_{\mathrm{et}}((Y/K')_{\bar k},\mathbb F_r).
\]
This is the finite-cover Hochschild--Serre argument of
\citet[Theorem~14.9]{MilneEtale}, here with group
order already invertible in the coefficient field.
The two quotient genera are
$1+s^n(g-1)$ and $1+s^{n-1}(g-1)$, respectively,
by the \'etale Riemann--Hurwitz formula, even if $s=p$.
The simple character factor belonging to $\chi$ occurs
only in $W^K/W^{K'}$: $K$ acts trivially on it and
$K'/K$ acts nontrivially.  If its multiplicity is $m_\chi$,
its dimension over $\mathbb F_r$ is $m_\chi f_n$.

If $r\ne p$, the usual cohomology dimension gives
\[
        \dim_{\mathbb F_r}(W^K/W^{K'})
                   =(2g-2)\varphi(s^n).                \tag{A.56.7}
\]
Consequently
\[
             m_\chi\le(2g-2)
                          \frac{\varphi(s^n)}{f_n}
                        \le(2g-2)c_{r,s}.              \tag{A.56.8}
\]
The multiplicity of the trivial character is then
$\dim W^A=2g$.

If $r=p$, the cohomology dimension is the $p$-rank,
not twice the genus.  For any proper smooth geometric
curve $Z$ this rank is at most $g(Z)$.
For completeness, the Artin--Schreier exact sequence
identifies $H^1_{\mathrm{et}}(Z,\mathbb F_p)$ with
\[
 \ker\bigl(F-1:H^1(Z,\mathcal O_Z)
                     \longrightarrow H^1(Z,\mathcal O_Z)\bigr),
\]
because $F-1$ is surjective on its algebraically closed
field of global constants.  Fixed vectors that are linearly
independent over $\mathbb F_p$ are linearly independent
over that algebraically closed field: in a shortest linear
relation, normalize one coefficient to one and subtract its
Frobenius transform.  Minimality forces every remaining
coefficient to satisfy $a^p=a$, contradicting the initial
$\mathbb F_p$-independence.  Their number is therefore at
most $\dim H^1(Z,\mathcal O_Z)=g(Z)$.
Apply this to $Y/K$.  We obtain
\[
 m_\chi f_n\le\dim W^K\le1+s^n(g-1),
 \qquad
 m_\chi\le1+(g-1)\frac{s}{s-1}c_{r,s}.              \tag{A.56.8p}
\]
Here $1/f_n\le1$ and
$s^n/f_n=(s/(s-1))\varphi(s^n)/f_n$.
The trivial multiplicity is $f_C$.

Thus in both cases every simple field factor of
$\mathbb F_r[A]$ acts on an isotypic component of
dimension at most the stated $M$ over that field.
Only quotient genera, finite-cover descent, and the
Artin--Schreier calculation in the characteristic case
have been used; no characteristic-zero comparison or
exact Chevalley--Weil multiplicity formula is needed.

\medskip\noindent\textbf{The arithmetic action becomes field-linear after a bounded power.}
Identify $G_k$ with $\widehat{\mathbb Z}$ using arithmetic
Frobenius, and consider its canonical factor $\mathbb Z_t$.
The action on the abelian $s$-deck group $A$ factors
through the action on $T_sJ_C$.  Set $d_s=2g$ if
$s\ne p$ and $d_s=f_C$ if $s=p$; the latter is the
rank of the \'etale Tate module.  The kernel of reduction
\[
 \operatorname{GL}_{d_s}(\mathbb Z_s)
             \longrightarrow\operatorname{GL}_{d_s}(\mathbb F_s)
\]
is pro-$s$, because its successive congruence quotients
are additive subgroups of matrix spaces over
$\mathbb F_s$.  The image of a pro-$t$ group intersects
this kernel trivially.  Therefore
\[
                 U=t^{b_s(t)}\mathbb Z_t
\]
acts trivially on $T_sJ_C$ and on every possible $A$.
The comparison between the $s$-adic abelianization and
$T_sJ_C$ is the usual Abel comparison
\citep[Proposition~9.1, Corollary~9.6]{MilneJac};
duality does not change the order of an automorphism.
This comparison includes the \'etale characteristic-primary
module.  If $s=p$ and $f_C=0$, every abelian $s$-cover is
trivial, the matrix group has rank zero, and $b_s(t)=0$;
the argument remains valid with these conventions.

Functoriality of geometric deck transformations shows
that the action of $U$ on $W$ commutes with the whole
image of $\mathbb F_r[A]$.  On a simple factor that image
is its full coefficient field $\mathbb F_{r^{f_n}}$.
Consequently this action is linear over that field,
and lies in $\operatorname{GL}_{m_\chi}
(\mathbb F_{r^{f_n}})$; it is not merely semilinear.
The trivial character gives
$\operatorname{GL}_{d_r}(\mathbb F_r)$, where $d_r=2g$
for $r\ne p$ and $d_r=f_C$ for $r=p$.

A $t$-element of $\operatorname{GL}_m(\mathbb F_{r^f})$
is semisimple because $t\ne r$.  Each eigenvalue belongs
to an extension of $\mathbb F_{r^f}$ of degree
$j\le m$, so the exponent of its $t$-part is at most
\[
                  \max_{1\le j\le m}v_t(r^{fj}-1).
\]
Using (A.56.5), (A.56.6), and $m\le M$, this is at most
\[
              L=e_t(r)+v_t(s-1)+\lfloor\log_t M\rfloor.
                                                               \tag{A.56.9}
\]
This bound also includes the trivial factor, where $f=1$.
The order of an element in a product of such factors
is the least common multiple of the component orders,
so its $t$-valuation is bounded by their maximum, not
by their sum.  It follows that
\[
                    t^{b_s(t)+L}\mathbb Z_t
\]
acts trivially on $W$, uniformly over all the covers $Y$.

\medskip\noindent\textbf{Further pro-$r$ towers preserve the same bound.}
Let $\Lambda_Y=\pi_1(Y_{\bar k},O_Y)^{(r)}$.
It is finitely generated: its mod-$r$ abelianization is
dual to the finite-dimensional space $W$, and the pro-$r$
Frattini generation criterion applies, also when $r=p$.
The kernel of
\[
 \operatorname{Aut}(\Lambda_Y)\longrightarrow
 \operatorname{Aut}(\Lambda_Y/\Phi(\Lambda_Y))
\]
is pro-$r$.  Here is the needed finite-level argument.
A finitely generated pro-$r$ group has a characteristic
open basis obtained by intersecting the kernels of all
maps to $r$-groups of bounded order.  For a finite
$r$-group $P$, the kernel of
$\operatorname{Aut}(P)\to\operatorname{Aut}(P/\Phi(P))$
is an $r$-group: filter $P$ by its lower exponent-$r$
central series; automorphisms trivial on its first
quotient are trivial on the generated associated graded,
and successive deeper congruence kernels embed in
additive groups of homomorphisms between these elementary
abelian $r$-quotients.  Taking the finite characteristic
quotients proves the asserted pro-$r$ kernel.

By the preceding step, the pro-$t$ group
$t^{b_s(t)+L}\mathbb Z_t$ acts trivially on
$H^1(Y_{\bar k},\mathbb F_r)$, hence on the Frattini
quotient of $\Lambda_Y$.  Its image in that pro-$r$
kernel is both pro-$t$ and pro-$r$ and is therefore
trivial.  It acts identically on the whole based group
$\Lambda_Y$.

For a connected finite geometric $r$-cover $D\to Y$
with $r$-group Galois closure, the geometric group
$\pi_1(D_{\bar k})^{(r)}$ is the corresponding open
subgroup of $\Lambda_Y$ with its full pro-$r$ topology.
Indeed, a finite $r$-quotient of that subgroup can be
dominated by its normal core; the resulting finite
quotient upstairs is an extension of an $r$-group by
an $r$-group.  Thus the identity action on $\Lambda_Y$
restricts to the identity on the full pro-$r$ group
of $D$, and in particular on
$H^1(D_{\bar k},\mathbb F_r)$.

For the covers in the first formulation of the claim,
the subgroup $R$ is the unique Sylow $r$-subgroup of
$G_N$, hence characteristic.  The intermediate quotient
$Y=D_N/R$ is accordingly stable under the arithmetic
action and has a rational lift of $O$.  It is an abelian
$s$-cover, and $D_N\to Y$ is an $r$-cover.  The preceding
argument applies to it.  Finally, restriction of a
finite cyclic Frobenius representation to the canonical
$\mathbb Z_t$ factor has image equal to the $t$-Sylow
subgroup of that cyclic image.  Its being killed by
$t^{b_s(t)+L}$ is exactly (A.56.1), with (A.56.2).

The unrestricted supremum in A.54, or in A.57 when $p$
is retained, ranges over all
normal two-prime covers.  Those deck groups need not
have the form just analyzed.  Thus the boundedness
proved here is an obstruction to a restricted testing
route, rather than boundedness of that unrestricted
supremum.
\end{proof}
\paragraph{Check record.} Checked only in this writing, 18 November 2026, GPT-6 Astra (OpenAI), ChatGPT Pro mode; separate effort setting not exposed. Complete proof written and read by its named authoring reader; see CHECK\_EVENTS.tsv for additional second readings and computation scope.


\subsection{Detecting a retained characteristic prime on a proper cover}

\begin{claim}{A.57}{Proved}{Positive $p$-rank in a fixed-prime tower and the retained-$p$ extension of A.54}
Let $C$ be a smooth proper geometrically connected curve of
genus $g\ge2$ over an algebraically closed field of
characteristic $p>0$, and fix any prime $\ell\ne p$.
There is a connected finite \'etale cyclic $\ell$-power
cover $Y\to C$ with positive $p$-rank.  This asserts
positive stable rank on at least one nontrivial character
orbit; it does not assert that the entire new Jacobian
factor is ordinary.

Consequently let $C/k$ be such a curve over a finite field,
and let $\Sigma$ be any prime set containing an invertible
prime $r\ne p$, with $p$ allowed in $\Sigma$.
Every retained prime factor of $G_k$ is detected by the
outer action on $\Delta=\pi_1(C_{\bar k})^\Sigma$.
In particular, when $p\in\Sigma$,
\[
          \ker\bigl(G_k\to\operatorname{Out}(\Delta)\bigr)
                         \cap\mathbb Z_p=0.             \tag{A.57.1}
\]
With a rational base point and the definitions of
$b_N(t)$ and $B_t$ from A.54, using any fixed invertible
$r\in\Sigma$, its complete formulas now hold in this
larger range:
\[
 \ker\rho\cap\mathbb Z_t=\ker\alpha\cap\mathbb Z_t
                =t^{B_t}\mathbb Z_t\quad(t\notin\Sigma),
 \qquad
 \ker\rho=\prod_{t\notin\Sigma}t^{B_t}\mathbb Z_t.       \tag{A.57.2}
\]
Thus a finite $B_t$ still gives the actual infinite
exceptional image described in A.54.  If no rational
base point is initially available, use a finite constant
extension as in that claim.  This removes its retained
$p$-rank-zero limitation; it does not prove unboundedness
of the omitted-prime orders or a converse from outer
faithfulness to point separation.
\end{claim}

\begin{proof}
\medskip\noindent\textbf{The exact theta-divisor input.}
Write $F$ for the absolute Frobenius of $C$ and set
\[
                  B=\operatorname{coker}
                      (\mathcal O_C\to F_*\mathcal O_C).
\]
This is the bundle of locally exact differentials, with
its scalar action through $p$-th powers.  It has rank
$p-1$ and degree $(p-1)(g-1)$.
Raynaud's theorem asserts that
\[
 \Theta_B=\{M\in\operatorname{Pic}^0(C):
                      H^0(C,B\otimes M)\ne0\}
\]
is the support of an effective Cartier divisor, in
particular a proper closed subset of the Jacobian.
The exact source is
\citet[\S4.1, Theorem~4.1.1; proof in \S4.2]{CF_Raynaud1982}.
Its relative-Frobenius formulation places $B$ on
$C^{(p)}$.  Over the perfect algebraically closed field,
the semilinear isomorphism $C^{(p)}\to C$ transports
that assertion to the absolute-Frobenius notation above.

For every nontrivial degree-zero line bundle $M$ of
order prime to $p$, projection formula gives the exact
sequence
\[
 0\longrightarrow M\longrightarrow F_*M^{\otimes p}
             \longrightarrow B\otimes M\longrightarrow0.
                                                               \tag{A.57.3}
\]
The first two bundles in cohomology have no global
sections, and their first cohomology groups both have
dimension $g-1$.  Hence
\[
 M\notin\Theta_B
 \quad\Longleftrightarrow\quad
 F:H^1(C,M)\longrightarrow H^1(C,M^{\otimes p})
                \text{ is a semilinear isomorphism}.    \tag{A.57.4}
\]
Replacing $M$ by $M^{\otimes p^j}$ uses the same bundle
$B$ and the same divisor $\Theta_B$ at every step.
The parameter permutation is multiplication by $p$ on
$\operatorname{Pic}^0(C)$, not arithmetic Frobenius on
its coefficient field.  Formula (A.57.3) is also the
tensorized exact sequence in Raynaud's proof; no
ordinarity assumption on $C$ has been imposed.

\medskip\noindent\textbf{A uniform torsion-intersection bound.}
Choose a symmetric very ample line bundle $H$ on
$J=J_C$, and let $D$ be the effective divisor supplied
above.  There is a constant
\[
                         K=(D\cdot H^{g-1})
\]
such that, for every integer $N$ prime to $p$,
\[
                       \#(D\cap J[N])\le K N^{2g-2}.
                                                               \tag{A.57.5}
\]
Here and below intersections with the torsion subgroup
mean sets of geometric points.
To prove (A.57.5), choose $g-1$ hyperplanes in the
embedding defined by $H$, each passing through the
origin, whose intersection with $[N](\operatorname{Supp}D)$
is zero-dimensional.  Such hyperplanes exist over the
infinite algebraically closed field: their common base
locus on $J$ is just the origin, so successive general
members lower the dimension of every positive-dimensional
component.  Their pullbacks under the finite map $[N]$
therefore meet $D$ properly in a zero-dimensional scheme
containing every point of $D\cap J[N]$.
Its intersection degree is
\[
 (D\cdot([N]^*H)^{g-1})
                 =N^{2g-2}(D\cdot H^{g-1}),
\]
because $[N]^*H$ is numerically equivalent to $H^{\otimes N^2}$.
Every geometric point contributes at least one to that
degree.  This proves (A.57.5), uniformly in $N$.
This is the hyperplane-cut form of the intersection calculation,
proved explicitly in
\citet[Lemma~3.9(ii),(iii)]{CF_Tamagawa2003} for arbitrary
prime-to-$p$ $N$.
The same source's Remark~3.11 identifies a missing finite-intersection
hypothesis in \citet[Lemma~4.3.5]{CF_Raynaud1982} and corrects
the numerical threshold of that paper's Theorem~4.3.1.
The argument above ensures the required proper intersections
by construction and does not use that numerical threshold.

\medskip\noindent\textbf{A full multiplication-$p$ orbit avoids the divisor.}
Put $N=\ell^n$ and
$h_N=\operatorname{ord}_{N}(p)$.
Multiplication by $p$ permutes $J[N]$, and every orbit
has length dividing $h_N$.  In particular
\[
                     h_N\le\varphi(N)\le N.
\]
This holds also for $\ell=2$: then $p$ is odd, and
$h_N$ divides $\varphi(2^n)$, with no lower bound on
the orbit size required.
The set of torsion points whose orbit meets $D$ is
\[
 \mathcal B_N=\bigcup_{j=0}^{h_N-1}
       [p]^{-j}(D\cap J[N]).
\]
Each permutation preserves cardinality, so
\[
              \#\mathcal B_N\le K N^{2g-1}.
\]
On the other hand $\#J[N]=N^{2g}$.  Choose $n$ so
large that $N^{2g}-K N^{2g-1}>1$.  There is then a
nonzero $M\in J[N]$ whose entire multiplication-$p$
orbit avoids $D$.  Its exact order is $\ell^m$ for
some $m\ge1$; it need not have exact order $\ell^n$,
and the proof does not require that stronger assertion.

Kummer theory associates to $M$ a connected cyclic
\'etale cover $Y\to C$ of degree $\ell^m$, with
\[
        \pi_*\mathcal O_Y
                  =\bigoplus_{i=0}^{\ell^m-1}M^{\otimes i},
\]
after choosing a trivialization of $M^{\otimes\ell^m}$.
Exact order gives connectedness, and $p\nmid\ell^m$
gives \'etaleness.  On first coherent cohomology,
absolute Frobenius sends the summand indexed by $i$
to that indexed by $pi$.  On the orbit containing
$i=1$ all these maps are isomorphisms by (A.57.4).
The resulting direct summand $V$ is Frobenius-stable,
nonzero, and Frobenius is bijective on it; each of its
summands has dimension $g-1>0$.

The Artin--Schreier sequence identifies
$H^1_{\mathrm{et}}(Y,\mathbb F_p)$ with the fixed
vectors of Frobenius on $H^1(Y,\mathcal O_Y)$.
Indeed, $x\mapsto x^p-x$ is surjective on the
algebraically closed ground field, so the preceding
$H^0$ term contributes no cokernel.
A bijective $p$-semilinear operator on a vector space
of positive dimension $d$ has $p^d$ fixed vectors.
For an elementary verification, in coordinates it is
$v\mapsto A v^{[p]}$ with $A$ invertible; the fixed
equations $v^{[p]}-A^{-1}v=0$ have leading monomials
$v_1^p,\ldots,v_d^p$, hence define a scheme of length
$p^d$, and their constant invertible Jacobian makes
that scheme \'etale.  Thus $V$ supplies a nonzero
class in $H^1_{\mathrm{et}}(Y,\mathbb F_p)$.
Equivalently $J_Y$ has positive $p$-rank.  This proves
the first assertion.  It uses one good orbit only;
no statement about all nontrivial character orbits is
being made.

\medskip\noindent\textbf{A pointed arithmetic level with positive $p$-rank.}
Return to the finite-field situation, initially with
$O\in C(k)$, and suppose $p\in\Sigma$.
Apply the first assertion with $\ell=r$ to obtain a
geometric cyclic $r$-power cover $Y\to C_{\bar k}$.
Its normal subgroup has a finite orbit under $G_k$,
because a finite cover descends to a finite field.
Intersect that orbit.  The resulting subgroup $N$
is open, normal, and $G_k$-stable, and its quotient
is an abelian $r$-group, being a subgroup of the
product of its conjugate cyclic quotients.
In particular this is a retained $\Sigma$-level.
The corresponding pointed arithmetic cover $D_N/k$
defined by $N\rtimes s_O(G_k)$ dominates $Y$ geometrically
and has a rational lift of $O$.

Its geometric degree over $Y$ is an $r$-power.  Pullback
and norm on Jacobians compose to multiplication by
that degree, which is a unit on the \'etale $p$-adic
Tate module.  Hence
\[
                         T_pJ_Y\hookrightarrow T_pJ_{D_N},
\]
and $J_{D_N}$ has positive $p$-rank.

\medskip\noindent\textbf{Inner arithmetic action cannot hide the characteristic factor.}
Frobenius on $T_pJ_{D_N}$ gives a representation of
$G_k$ in $\operatorname{GL}_f(\mathbb Z_p)$ with
$f>0$.  Its restriction to $\mathbb Z_p$ is faithful.
Indeed, a nonzero kernel in $\mathbb Z_p$ would be
open and its image finite.  The closure of the cyclic
Frobenius image in this matrix group has an open
pro-$p$ subgroup.  Finiteness of its $p$-Sylow image
would then make the whole closure finite.  Some
positive integral power of Frobenius would act as
the identity on $T_pJ_{D_N}$, making the finite group
$J_{D_N}(k')$ contain geometric $p$-power torsion of
arbitrarily large order for one finite field $k'$.
That is impossible.

Let $K_p=\ker\rho\cap\mathbb Z_p$.  For every
$h\in K_p$, the based action $\alpha_h$ on $\Delta$
is conjugation by some element of $\Delta$.
Its induced action on the abelianization of $N$
therefore belongs to the finite image of the deck
group $\Delta/N$.  The full induced $\Sigma$ topology
at this open level identifies its $p$-adic abelianization
with $T_pJ_{D_N}$, by the normal-core argument and
the all-prime Abel comparison
\citep[Proposition~9.1, Corollary~9.6]{MilneJac}.
Consequently $K_p$ has finite image on $T_pJ_{D_N}$.
If $K_p$ were nonzero it would be open in
$\mathbb Z_p$, contradicting the faithful action
just established.  Thus $K_p=0$.
The same conclusion without an initial rational point
follows after a finite constant extension: a nonzero
closed subgroup of $\mathbb Z_p$ has nonzero
intersection with every open arithmetic subgroup.

\medskip\noindent\textbf{Extension of the all-cover criterion.}
Every step of A.54 preceding its retained-prime
calculation remains valid with $p\in\Sigma$.
Specifically, the pointed normal levels are cofinal;
the finite-core argument supplies the full induced
$\Sigma$ topology; deck transformations act faithfully
on $H^1(-,\mathbb Q_r)$ because a nontrivial free
cyclic quotient lowers the genus; and averaging over
$\mathbb Q_r$ remains valid even when the deck order
is divisible by $r$ or $p$.  Finally, at an omitted
prime $t$, both the pro-$r$ congruence-kernel argument
and the comparison of pro-$t$ inner images with the
pro-$\Sigma$ group $\operatorname{Inn}(\Delta)$ apply
unchanged.
They give the first identity of (A.57.2).
Every retained invertible prime is detected on the
base Tate module as in A.54; (A.57.1) supplies the
remaining retained characteristic prime.  Decomposition
of a closed subgroup of $\widehat{\mathbb Z}$ into its
prime factors now gives the second identity of
(A.57.2).  The central arithmetic element and actual
decomposition-group collision argument of A.54 then
applies verbatim at any omitted prime with finite
$B_t$.  This completes the claimed extension.
\end{proof}
\paragraph{Check record.} Checked only in this writing, 18 November 2026, GPT-6 Astra (OpenAI), ChatGPT Pro mode; separate effort setting not exposed. Complete proof written and read by its named authoring reader; see CHECK\_EVENTS.tsv for additional second readings and computation scope.


\subsection{Automorphism orders in two-prime quotients}

\begin{claim}{A.58}{Proved}{Unbounded omitted-prime IA orders in two-generated two-prime groups}
Fix distinct primes $r,s$, and let $t\notin\{r,s\}$.
For every integer $b\ge1$ there is a two-generated finite
$\{r,s\}$-group $G_b$ and an automorphism $\beta_b$ such that:
\begin{enumerate}
\item $G_b$ has a normal elementary abelian $r$-subgroup
      $W_b$ and a metabelian $s$-group quotient;
      in particular, its derived length is at most three;
\item $\beta_b$ fixes that $s$-group complement pointwise
      and acts trivially on $G_b^{\mathrm{ab}}$;
\item both $\beta_b$ and its class in
      $\operatorname{Out}(G_b)$ have exact order $t^b$.
\end{enumerate}
Each $G_b$ is a quotient of the pro-$\{r,s\}$ completion
of every closed orientable surface group of genus at least two.
Consequently, if $p\notin\{r,s\}$, it is also a quotient
of the geometric pro-$\{r,s\}$ fundamental group of every
smooth proper genus-at-least-two curve in characteristic $p$.

Thus the number of generators, the two allowed group primes,
and even the derived-length bound three do not bound
omitted-prime orders in the kernel of
$\operatorname{Aut}(G)\to\operatorname{Aut}(G^{\mathrm{ab}})$.
This statement does not assert that the automorphisms
$\beta_b$ lift to an automorphism of a single geometric
pro-$\{r,s\}$ surface group, form a compatible inverse system,
or arise from arithmetic Frobenius of a fixed curve.
It does not determine the all-cover quantities $B_t$ of A.54.
\end{claim}

\begin{proof}
Put
\[
 f=\operatorname{ord}_s(r),\qquad
 F=\mathbb F_{r^f},\qquad q_0=r^f,
\]
and choose $\omega\in F^\times$ of exact order $s$.
When $s=2$, this means $f=1$ and $\omega=-1$.
For an integer $n\ge1$ to be chosen below, put $d=s^n$.
Let $V=F^d$, with column basis $e_1,\ldots,e_d$.
For $1\le i\le d$, let $u_i$ be the diagonal matrix
which multiplies $e_i$ by $\omega$ and fixes the other
basis vectors.  Let $v$ cyclically permute this basis.
These matrices generate the group
\[
 S=\langle u_1,v\rangle
    =\langle u_1,\ldots,u_d\rangle\rtimes\langle v\rangle
    \simeq (C_s)^d\rtimes C_d.                         \tag{A.58.1}
\]
Indeed, conjugation by $v$ cycles the $u_i$, the diagonal
subgroup is $(C_s)^d$, and no nonidentity power of the
cyclic permutation matrix is diagonal.  Thus $S$ is an
$s$-group of order $s^d d$, generated by two elements.
Both its diagonal subgroup and its cyclic quotient are
abelian, so $S$ is metabelian.

\medskip\noindent\textbf{The coefficient algebra is the full matrix algebra.}
Regard $V$ as an $\mathbb F_r[S]$-module and let
$\mathcal A\subseteq\operatorname{End}_{\mathbb F_r}(V)$
be the image of its group algebra.  All the displayed
matrices are $F$-linear, so
$\mathcal A\subseteq\operatorname{Mat}_d(F)$.
Writing $E_{ij}$ for the standard matrix units gives
\[
                 E_{ii}=(u_i-I)^{q_0-1}\in\mathcal A.  \tag{A.58.2}
\]
This equality holds because the only nonzero entry of
$u_i-I$ is $\omega-1$, whose $(q_0-1)$st power is one.
Also $u_iE_{ii}=\omega E_{ii}$.
Since $\mathbb F_r[\omega]=F$, polynomial expressions
in $u_i$ multiplied by $E_{ii}$ give $aE_{ii}$ for
every $a\in F$.  For any $i,j$, a suitable power $v^k$
sends $e_j$ to $e_i$, and hence
\[
                    E_{ii}v^kE_{jj}=E_{ij}.
\]
Multiplying by the diagonal scalar matrices just found
gives every $aE_{ij}$.  Consequently
\[
                     \mathcal A=\operatorname{Mat}_d(F).
                                                               \tag{A.58.3}
\]
In particular, there is no division by a scalar outside
the originally generated $\mathbb F_r$-algebra in this
argument.

\medskip\noindent\textbf{The semidirect product still has two generators.}
Set $m=d-1$, and regard
\[
                     W=V^m=\operatorname{Mat}_{d\times m}(F)
\]
as an additive group, with $S$ acting by left multiplication.
It is elementary abelian of order $r^{f d(d-1)}$.
Let
\[
                 w=(e_2,\ldots,e_d)\in W.
\]
The $\mathbb F_r[S]$-submodule generated by $w$ is all
of $W$: by (A.58.3), left multiplication by an arbitrary
$d\times d$ matrix prescribes independently the columns
$2,\ldots,d$, which are exactly the columns of its product
with $w$.

Form $G=W\rtimes S$, with multiplication
\[
             (a,h)(a',h')=(a+ha',hh').
\]
The element $u=u_1$ fixes $w$, since every column of $w$
has zero first coordinate.  Therefore the two elements
\[
                         x=(w,u),\qquad y=(0,v)
\]
satisfy
\[
                          x^s=(sw,1).                 \tag{A.58.4}
\]
Because $r\ne s$, multiplication by $s$ is invertible
on the elementary abelian $r$-group $W$.
Thus a suitable integral power of $x^s$ is $(w,1)$,
and $x(w,1)^{-1}=(0,u)$ belongs to $\langle x,y\rangle$.
Together with $y$, this generates the whole complement
$S$.  Conjugates of $(w,1)$ by that complement generate
its $\mathbb F_r[S]$-span, which is $W$.
It follows that
\[
                            G=\langle x,y\rangle.     \tag{A.58.5}
\]
Its order is $r^{f d(d-1)}s^d d$, so its only prime
divisors are $r,s$.  Since it is an abelian extension of
the metabelian group $S$, its derived length is at most three.

\medskip\noindent\textbf{Multiplicity automorphisms are IA.}
Right multiplication by any $T\in\operatorname{GL}_m(F)$
commutes with the left $S$-action on $W$.
Consequently
\[
              \beta_T:(a,h)\longmapsto(aT,h)           \tag{A.58.6}
\]
is an automorphism of $G$ fixing the complement $S$
pointwise.  The right-action convention reverses the order
of composition in $T$, which has no effect on injectivity
or on the order of an individual $\beta_T$.

For every $i$, the image $(u_i-I)W$ contains all matrices
supported on row $i$, because $\omega-1$ is a nonzero
element of $F$.  The sum of these subgroups is $W$.
These images are generated by commutators of elements of
$W$ with $u_i$, up to an immaterial sign or inverse.
Hence
\[
                             [W,S]=W.                 \tag{A.58.7}
\]
It follows directly from the semidirect product law that
$G^{\mathrm{ab}}\simeq S^{\mathrm{ab}}$.
Every $\beta_T$ acts as the identity on this quotient,
so (A.58.6) supplies IA automorphisms.

\medskip\noindent\textbf{Their omitted-prime orders are unbounded.}
Fix the requested exponent $b\ge1$, and put
\[
                       a=\operatorname{ord}_{t^b}(q_0).
\]
This is defined because $t\ne r$.
The multiplicative group of $\mathbb F_{q_0^a}$ contains
an element $\zeta$ of exact order $t^b$.
Multiplication by $\zeta$ on the $a$-dimensional
$F$-vector space $\mathbb F_{q_0^a}$ is an invertible
matrix of exact order $t^b$.
Choose $n$ so large that $d-1=s^n-1\ge a$, and extend
this matrix by the identity on the remaining coordinates.
The resulting $T\in\operatorname{GL}_{d-1}(F)$ has
exact order $t^b$.  Its action on $W$ is faithful,
so $\beta_T$ has that same order.

Finally, $\operatorname{Inn}(G)$ is a quotient of the
$\{r,s\}$-group $G$ and has no nontrivial $t$-subgroup.
The cyclic group $\langle\beta_T\rangle$, of order $t^b$,
therefore meets $\operatorname{Inn}(G)$ trivially.
The image of $\beta_T$ in $\operatorname{Out}(G)$ has
exactly the same order.  Take this $G$ and $\beta_T$
as $G_b$ and $\beta_b$.

\medskip\noindent\textbf{Geometric quotient realization and its limitation.}
A closed orientable surface group of genus $g\ge2$
has the presentation
\[
 \langle a_1,b_1,\ldots,a_g,b_g\mid
                   \prod_{i=1}^g[a_i,b_i]=1\rangle.
\]
Sending $a_1$ to $x$, $a_2$ to $y$, and every other
displayed generator to the identity respects the relation
and is surjective by (A.58.5).
Since $G$ is a $\{r,s\}$-group, this map factors through
the pro-$\{r,s\}$ completion.
For a proper smooth curve in characteristic
$p\notin\{r,s\}$, the geometric comparison and
specialization input explicitly established in the proof
of A.4 identifies this completion with the geometric
pro-$\{r,s\}$ group.  This proves the asserted geometric
quotient realization.  This last step uses the existing
proper-curve comparison; it is not part of the direct
characteristic-$p$ configuration proof of A.45--A.48.

Neither this surjection nor the matrix construction gives
a lift of $\beta_T$ to the geometric surface group.
Still less does it identify $\beta_T$ with Frobenius for
a fixed curve over a fixed finite field.
The claimed failure of a group-theoretic uniform bound
is already the exact finite statement proved above.
\end{proof}
\paragraph{Check record.} Checked only in this writing, 18 November 2026, GPT-6 Astra (OpenAI), ChatGPT Pro mode; separate effort setting not exposed. Complete proof written and read by its named authoring reader; see CHECK\_EVENTS.tsv for additional second readings and computation scope.


\subsection{A single arithmetic tower for the two-prime order test}

\begin{claim}{A.59}{Proved}{Cofinal alternating Jacobian covers and their Frobenius orders}
Let $C/k$ be a smooth proper geometrically connected curve of
genus $g\ge2$ over $k=\mathbb F_q$, with $O\in C(k)$.
Let $r,s$ be distinct primes with $r\ne\operatorname{char}k$;
$s$ is allowed to equal the characteristic.
Put $\Sigma=\{r,s\}$ and $\Delta=\pi_1(C_{\bar k},O)^\Sigma$.
For a profinite group $H$ and prime $\ell$, write
\[
 \Psi_\ell(H)=\overline{[H,H]H^\ell}
       =\ker\bigl(H\longrightarrow H^{\mathrm{ab}}/\ell\bigr).
\]
Starting with $N_0=\Delta$, define
\[
 N_{2j+1}=\Psi_r(N_{2j}),\qquad
 N_{2j+2}=\Psi_s(N_{2j+1})\quad(j\ge0).                \tag{A.59.1}
\]
These are open characteristic subgroups of $\Delta$,
cofinal among all its open subgroups.  They determine a
single tower of pointed geometrically connected finite
\'etale covers
\[
 \cdots\longrightarrow(D_2,O_2)\longrightarrow(D_1,O_1)
              \longrightarrow(D_0,O_0)=(C,O)           \tag{A.59.2}
\]
over the same field $k$, with each $O_n\in D_n(k)$.
At an $\ell$-step, the geometric deck group is the full
maximal elementary abelian $\ell$-quotient of
$\pi_1((D_n)_{\bar k})$.
Equivalently, the step is the pullback along
$x\mapsto[x-O_n]$ of the corresponding full elementary
abelian geometric cover of $J_{D_n}$, with its zero lift.

Fix an omitted prime $t\notin\{r,s\}$ and put
\[
 b_n(t)=v_t\!\left(\operatorname{ord}
    \bigl(\operatorname{Frob}_q\mid
       H^1_{\mathrm{et}}((D_n)_{\bar k},\mathbb F_r)\bigr)\right).
\]
Then $b_n(t)$ is nondecreasing, and the all-cover supremum
$B_t$ in A.54, with the extension of A.57 when needed, is
\[
                         B_t=\sup_{n\ge0}b_n(t).       \tag{A.59.3}
\]
In particular, a finite $B_t$ is equivalent to eventual
constancy of this integer sequence.

Each $b_n(t)$ is determined by the point-count Weil
polynomial $P_n(X)\in\mathbb Z[X]$ of $D_n/k$.
Factor its reduction into distinct monic irreducibles:
\[
                  P_n(X)\bmod r=\prod_h h(X)^{m_h}.
\]
None of these factors is $X$.  In the field
$F_h=\mathbb F_r[X]/(h)$, let $\xi_h$ be the class of $X$.
Then
\[
                b_n(t)=\max_h v_t
                         \bigl(\operatorname{ord}_{F_h^\times}
                                      (\xi_h)\bigr).   \tag{A.59.4}
\]
The polynomial $P_n$ is determined by
$\#D_n(\mathbb F_{q^m})$ for $1\le m\le2g(D_n)$.
Thus (A.59.3) is a test using the actual arithmetic action
on one fixed, recursively defined tower of curves.
It supplies no bound on a stage at which eventual
constancy can be recognized and no proof that the sequence
is unbounded.  In particular, it does not settle A12.
\end{claim}

\begin{proof}
\medskip\noindent\textbf{The exact finite-group input.}
The external group-theoretic input is
\citet[Theorem~III, p.~392]{CF_Burnside1904}:
a finite group whose order involves at most two distinct
primes is solvable.  Its hypotheses apply to every finite
quotient of $\Delta$ and to every subgroup of such a quotient.
The primary five-page paper, including its proof and
appended shorter argument, was opened.  No stronger
assertion about automorphism groups is being imported;
A.58 shows why such a stronger assertion would be false.

\medskip\noindent\textbf{Openness, characteristicity, and cofinality.}
The group $\Delta$ is finitely generated, as is the geometric
fundamental group of a proper curve.  Every open subgroup
is finitely generated, by the profinite Schreier argument:
coset representatives and the finitely many original
generators give finitely many generators of the open subgroup.
For a finitely generated profinite group $H$, its quotient
$H^{\mathrm{ab}}/\ell$ is a finite elementary abelian
$\ell$-group.  Therefore $\Psi_\ell(H)$ is open in $H$.
It is characteristic by its definition.  Induction proves
that each $N_n$ is open and characteristic in $\Delta$.

Let $a:\Delta\twoheadrightarrow G$ be any finite quotient,
and put $G_n=a(N_n)$.  Homomorphisms commute with the
formation of commutators and powers, and the target is
finite.  Hence
\[
 G_{2j+1}=[G_{2j},G_{2j}]G_{2j}^{\,r},\qquad
 G_{2j+2}=[G_{2j+1},G_{2j+1}]G_{2j+1}^{\,s}.          \tag{A.59.5}
\]
The descending sequence of subgroups $G_n$ of a finite
group is eventually constant, with some value $H$.
After both parities have occurred in the constant tail,
\[
                       H=[H,H]H^r=[H,H]H^s.
\]
The finite abelian group $H^{\mathrm{ab}}$ has only
$r$- and $s$-primary parts.  The first equality kills
its mod-$r$ quotient and the second its mod-$s$ quotient.
Thus $H^{\mathrm{ab}}=1$, so $H$ is perfect.
It is also solvable by the cited theorem, and therefore
$H=1$: its derived series is constant if it is perfect,
whereas solvability makes that series terminate at one.
Consequently $N_n\subseteq\ker a$ for some $n$.
The kernels of finite quotients form an open normal basis,
and every open subgroup contains its open normal core.
This proves the claimed cofinality and
$\bigcap_n N_n=1$.

\medskip\noindent\textbf{The tower exists over the fixed ground field.}
The section at $O$ gives a based $G_k$-action on $\Delta$.
Every $N_n$ is characteristic and hence stable under that
action.  In the geometrically pro-$\Sigma$ arithmetic
group, the subgroup
\[
                              N_n\rtimes s_O(G_k)
\]
therefore determines a connected finite \'etale cover
$D_n\to C$.  Its geometric subgroup is $N_n$, so it is
geometrically connected.  The identity coset is fixed by
$s_O(G_k)$, giving the specified rational point $O_n$.
The subgroup inclusions yield the transition maps in
(A.59.2), and their compatibility with the identity cosets
preserves the chosen points.

We check the induced prime-set topology, to ensure that
the successive quotients are the full geometric homology
quotients of these actual covers.  Let
$\widetilde N\subset\pi_1(C_{\bar k})$ be the preimage
of any open subgroup $N\subset\Delta$.
Given a finite $\Sigma$-quotient of $\widetilde N$ with
kernel $V$, let $N^0$ be the normal core of $N$ in $\Delta$
and let $\widetilde N^0$ be its preimage.
The core $U$ of $V$ in $\pi_1(C_{\bar k})$ is open and normal.
The quotient $\widetilde N^0/U$ embeds in the finite product
of the groups $\widetilde N^0/(\widetilde N^0\cap gVg^{-1})$
as $g$ runs through coset representatives for
$\pi_1(C_{\bar k})/\widetilde N$.
Each factor is a subgroup of a conjugate of
$\widetilde N/V$, hence is a $\Sigma$-group.
The group $\pi_1(C_{\bar k})/U$ is therefore an extension
of a $\Sigma$-subgroup of that product by the
$\Sigma$-group $\Delta/N^0$, and is itself a $\Sigma$-group.
Thus every finite $\Sigma$-quotient of the subgroup is
detected in the original maximal pro-$\Sigma$ quotient.
This gives
\[
                \pi_1((D_n)_{\bar k})^\Sigma=N_n
\]
with its full induced topology.  In particular,
$N_n/\Psi_\ell(N_n)$ is the full mod-$\ell$
abelianization of the actual geometric fundamental group,
for $\ell=r$ or $s$.

The Abel comparison
\citep[Proposition~9.1 and Corollary~9.6]{MilneJac}
identifies these abelian quotients with those of
$J_{D_n}$, including the \'etale characteristic-primary
quotient.  Pulling back the corresponding pointed cover
of the Jacobian therefore gives precisely the next
connected pointed cover in (A.59.2).
When $\ell$ is invertible, this is the pullback of
$[\ell]:J_{D_n}\to J_{D_n}$.
At $\ell=p$, the cover meant is the finite \'etale one
corresponding to $T_pJ_{D_n}/p$; the inseparable part of
$[p]$ is not part of that cover.  If the $p$-rank is zero,
that particular step has degree one, which is allowed.
The cofinality proof above still applies.

This construction uses the full geometric elementary
abelian quotient, with its arithmetic action.
It does not replace it by the group of rational points
of a Jacobian or by a rational class-field quotient.
The successive covers may have nonconstant deck groups
over $k$, while retaining the specified rational point.

\medskip\noindent\textbf{Comparison of omitted-prime orders.}
For every pointed normal level $D\to C$ considered in
A.54, its subgroup $N$ contains some $N_n$ by cofinality.
The resulting map $D_n\to D$ is defined over $k$ and is
finite \'etale.  Pullback induces a $G_k$-equivariant
injection
\[
 H^1_{\mathrm{et}}(D_{\bar k},\mathbb Q_r)
       \lhook\joinrel\longrightarrow
 H^1_{\mathrm{et}}((D_n)_{\bar k},\mathbb Q_r).         \tag{A.59.6}
\]
Indeed, trace composed with pullback is multiplication
by the positive degree.  That integer is invertible in
$\mathbb Q_r$, even when $r$ divides the degree.

For each of these curves, the kernel of
$\operatorname{GL}_d(\mathbb Z_r)\to
\operatorname{GL}_d(\mathbb F_r)$ is pro-$r$.
The image of the canonical factor $\mathbb Z_t\subset G_k$
meets that kernel trivially, because $t\ne r$.
Its order on rational $r$-adic cohomology therefore equals
its order after reduction modulo $r$; in both cases it is
the $t$-part of the finite mod-$r$ Frobenius order.
Equation (A.59.6) consequently gives
\[
                               b_D(t)\le b_n(t).
\]
Conversely, each $D_n$ is one of the allowed levels.
Taking suprema proves (A.59.3).
Applying the same injection to $D_{n+1}\to D_n$ proves
that $b_n(t)$ is nondecreasing.  A bounded nondecreasing
sequence of nonnegative integers is eventually constant,
which gives the stated equivalence.
The formula for the arithmetic kernel and the consequence
of a finite $B_t$ are precisely A.54 and A.57; no converse
from arithmetic faithfulness to point separation is added.

\medskip\noindent\textbf{Actual point counts determine the finite-stage orders.}
Write $h_n=g(D_n)$.
The Weil polynomial $P_n$ is monic of degree $2h_n$,
with constant coefficient $q^{h_n}$.
The traces of the $m$th power of its Frobenius are
\[
                         q^m+1-\#D_n(\mathbb F_{q^m}),
\]
so the first $2h_n$ point counts determine all its
coefficients by Newton's identities
\citep[Theorem~11.1]{MilneJac}.
It is the characteristic polynomial for the point-count
Frobenius on rational cohomology, or equivalently on the
covariant Tate module.
If the cohomological $G_k$-action is written using the
$q$-power arithmetic Frobenius, its eigenvalues are the
inverses of these; inversion does not change their finite
multiplicative orders.  Thus this convention difference
does not change the following calculation.

Because $r\ne\operatorname{char}k$, the constant term is
nonzero modulo $r$.  The eigenvalues of the finite matrix
are therefore nonzero, and the irreducible factors $h$
above all differ from $X$.
Over a finite splitting field, the matrix is a product
of commuting semisimple and unipotent parts.  The
unipotent part has $r$-power order, so contributes nothing
to its $t$-part.  The semisimple part has order the least
common multiple of the multiplicative orders of its
eigenvalues.  Roots of one irreducible $h$ are conjugate
over $\mathbb F_r$ and have the same order, namely that
of $\xi_h$ in $F_h^\times$.
The $t$-valuation of a least common multiple is the maximum
of the individual valuations.  This proves (A.59.4),
independently of the multiplicities $m_h$ or possible
nonsemisimplicity modulo $r$.

All entries in this finite-stage formula are attached to
the actual curves of (A.59.2) over the fixed $k$.
The proof gives no finite stopping rule for the infinite
sequence and no outcome for the remaining universal
classification.  It is a cofinal arithmetic test, not an
arithmetic realization of the arbitrary automorphisms of A.58.
\end{proof}
\paragraph{Check record.} Checked only in this writing, 18 November 2026, GPT-6 Astra (OpenAI), ChatGPT Pro mode; separate effort setting not exposed. Complete proof written and read by its named authoring reader; see CHECK\_EVENTS.tsv for additional second readings and computation scope.


\subsection{An actual increase of an omitted-prime Frobenius order}

\begin{claim}{A.60}{Proved}{The cyclic-nineteen cover introduces an arithmetic factor of seven}
Let $C/\mathbb F _8$ and its pointed constant cyclic-nineteen
cover $D\to C$ be those of A.50, based at $P=(0,1)$.
Thus $C$ is the smooth proper model of
$y^2+x(x+1)y=x^5+1$, and $D$ is the pullback of the
degree-nineteen Lang quotient of its Jacobian along
$R\mapsto[R-P]$.
Take $\Sigma=\{3,19\}$ and use $\mathbb F _3$ coefficients
in the all-cover order test.
Then the prime-to-three parts of the actual Frobenius
orders on $H^1_{\mathrm{et}}$ of $C$ and $D$ are respectively
$4$ and $56$.  In particular,
\[
\begin{array}{c|cc}
 & C & D\\\hline
 v_7(\operatorname{ord}\operatorname{Frob}_8)&0&1\\
 v_2(\operatorname{ord}\operatorname{Frob}_8)&2&3.
\end{array}                                                   \tag{A.60.1}
\]
These are orders of the actual arithmetic action on a
geometric $\Sigma$-cover of one fixed curve over one fixed
finite field.
Consequently the all-cover supremum cannot in general be
replaced by the order on the base Jacobian alone.

For the canonical tower of A.59, alternating first the
full mod-$3$ and then the full mod-$19$ quotient, its second
level dominates this $D$.  Hence
\[
                         b_0(7)=0,\qquad b_2(7)\ge1.          \tag{A.60.2}
\]
Moreover, every further pointed geometrically connected
finite \'etale cover $E\to D$ over $\mathbb F _8$ with
geometric $3$-group Galois closure has
\[
 v_7\!\left(\operatorname{ord}
   (\operatorname{Frob}_8\mid H^1_{\mathrm{et}}
                     (E_{\overline{\mathbb F}_8},\mathbb F _3))\right)=1.
                                                                  \tag{A.60.3}
\]
Thus continuing with only $3$-group covers cannot turn
this factor $7$ into a factor $7^2$.
No unboundedness of the full alternating tower and no
classification in A12 is asserted.
For this curve and these primes, the explicit bound in
A.56 on the entire class of normal $3$-group by abelian
$19$-group covers is exactly one at the omitted prime
seven; this cover attains that bound.
\end{claim}

\begin{proof}
\medskip\noindent\textbf{The actual cover and its character polynomials.}
A.50 proves that the cover is geometrically connected,
finite \'etale, constant cyclic of degree nineteen, and
pointed over $\mathbb F _8$.  Its genus is twenty.
That proof also supplies the exact nontrivial character
factors.  If $\zeta$ is a primitive nineteenth root of
unity, put
\[
 S=\zeta+\zeta^{18}+\zeta^7+\zeta^{12}
                           +\zeta^8+\zeta^{11}.
\]
The factor belonging to one nontrivial character is
\[
                       L_\chi(T)=1+(3+S)T+8T^2.             \tag{A.60.4}
\]
The full list of nontrivial character factors replaces
$\zeta$ by its powers $\zeta^a$, $1\le a\le18$; there
are eighteen factors in total, including the displayed one.
Their product is the relative zeta numerator.
The sign in (A.60.4) is the sign of the Euler-product
coefficient in A.50.  Thus its point-count Frobenius
polynomial is $X^2+(3+S)X+8$.

The six exponents form the subgroup
$H=\langle8\rangle$ of $(\mathbb Z/19)^\times$.
There are three Gaussian periods, associated with its
three cosets, and each occurs six times in the eighteen
characters.  Reduction using
$1+\zeta+\cdots+\zeta^{18}=0$ gives
\[
                             S^3+S^2-6S-7=0.                \tag{A.60.5}
\]
This is also a finite exact identity in
$\mathbb Z[Z]/(1+Z+\cdots+Z^{18})$.
The cubic $f(U)=U^3+U^2-6U-7$ has no rational root:
the only possible integral roots are $\pm1,\pm7$, and
none is a root.  Hence it is irreducible, so those three
periods are precisely its three conjugates.

Their product of reciprocal character factors is
\[
\begin{split}
 N(X)&=\operatorname{Res}_U
             \bigl(f(U),X^2+(3+U)X+8\bigr)\\
     &=X^6+8X^5+39X^4+135X^3+312X^2+512X+512.
\end{split}                                                   \tag{A.60.6}
\]
For a direct expansion, write $A=X^2+3X+8$; then
$N=A^3-XA^2-6X^2A+7X^3$.
The repeated occurrence of the periods therefore gives
the exact primitive polynomial $N(X)^6$.

A.50 counts $9$ and $87$ points on $C$ over
$\mathbb F _8$ and $\mathbb F _{64}$, respectively.
The first trace is zero and the second is $-22$.
Newton's identities and the genus-two reciprocal
symmetry give
\[
 P_C(X)=X^4+11X^2+64,\qquad
                     P_D(X)=P_C(X)N(X)^6.                    \tag{A.60.7}
\]
As a consistency check, $N(1)=1519$ and
$1519^6=12284193340665686881$, the independently proved
relative class number in A.50.

\medskip\noindent\textbf{The exact modular orders.}
Over $\mathbb F _3$,
\[
 \overline P_C=(X^2+1)^2,
 \qquad n(X):=\overline N=X^6-X^5-X-1.                       \tag{A.60.8}
\]
The roots of $X^2+1$ have exact order four.
We give explicit finite identities for $n$, rather than
inferring its root orders from its degree.
In the quotient by $n$, the class $x$ of $X$ satisfies
\[
\begin{split}
 x^8&=x^5+x^3-x^2-x+1,\\
 x^{28}&=-1,\qquad x^{56}=1.
\end{split}                                                   \tag{A.60.9}
\]
In particular $x^{729}=x$, because $729=13\cdot56+1$.
The following two remainders have the indicated inverses
modulo $n$:
\[
\begin{array}{c|c}
 X^9-X\pmod n & X^5-X^4+X^2-X-1\\
 X^{27}-X\pmod n & -X^5-X^3+X+1
\end{array}                                                   \tag{A.60.10}
\]
Here the remainders in the left column are, explicitly,
\[
 X^5+X^4-X^3-X^2+X+1,
              \qquad -X^5+X^4-X+1,
\]
and multiplying each by its right-column entry gives
remainder one modulo $n$.
Thus $n$ divides $X^{3^6}-X$ and is coprime to both
$X^{3^2}-X$ and $X^{3^3}-X$.
Every irreducible factor has degree dividing six, and
the latter two coprimality statements exclude degrees
one, two and three.  Since $n$ itself has degree six,
it is irreducible.

Consequently its quotient is a field.  Equations
(A.60.9) show that the order of $x$ divides $56$, but
does not divide $28$ or $8$.  These are the quotients
of $56$ by its two prime divisors.  Hence its order is
exactly $56=2^3\cdot7$.

For a finite invertible matrix over $\mathbb F _3$,
the unipotent part has $3$-power order, whereas the
semisimple order is the least common multiple of the
orders of its eigenvalues.  Therefore (A.60.7)--(A.60.9)
give prime-to-three orders four and
$\operatorname{lcm}(4,56)=56$.
Arithmetic and point-count Frobenius are inverse
conventions on cohomology, as explained in A.59;
inversion preserves these finite orders.
This proves (A.60.1) for the actual arithmetic action,
independently of any possible Jordan blocks.

\medskip\noindent\textbf{The canonical tower and the next obstruction.}
Let $\Delta$ be the geometric pro-$\{3,19\}$ group and
$a:\Delta\twoheadrightarrow C_{19}$ the quotient of
the pointed cover $D$.
The image of $N_1=\Psi_3(\Delta)$ is all of $C_{19}$,
because cubing is an automorphism of that cyclic group.
The image of $N_2=\Psi_{19}(N_1)$ is trivial.
Thus $N_2\subseteq\ker a$, and the second pointed
canonical level maps to $D$ over the same field.
The rational-cohomology trace injection and the pro-$3$
congruence argument of A.59 give (A.60.2).

For completeness, put
$\Lambda=\pi_1(D_{\overline{\mathbb F}_8})^{(3)}$.
It is a finitely generated pro-$3$ group, and
\[
 \ker\bigl(\operatorname{Aut}(\Lambda)\longrightarrow
              \operatorname{Aut}(\Lambda/\Phi(\Lambda))\bigr)
                        \quad\hbox{is pro-}3.                \tag{A.60.11}
\]
Indeed take the characteristic open subgroups obtained
by intersecting the finitely many kernels of quotients
of each bounded order.  On each finite $3$-group quotient,
automorphisms trivial on its Frattini quotient form a
$3$-group: they act freely on the ordered generating
lifts of one generating tuple, whose number is a power
of three.  Triviality on the Frattini quotient of
$\Lambda$ implies that property on every finite quotient.
Passing to the inverse limit proves (A.60.11).

The image of the canonical $\mathbb Z_7$ factor of
$G_{\mathbb F _8}$ consequently has order exactly seven
on $\Lambda$, since it has that order on its mod-$3$
abelianization by (A.60.1), and meets the kernel in
(A.60.11) trivially.
Each pointed $3$-group cover of $D$ has its full induced
pro-$3$ topology, by the same finite-core argument as
in A.59.  Its based arithmetic action is the restriction
of the one on $\Lambda$, so its mod-$3$ cohomology has
$7$-part of order at most seven.
Conversely rational $3$-adic pullback from $D$ is injective
by trace, and reduction does not change the finite
$7$-part.  Its order is therefore at least seven.
This proves (A.60.3).

The stronger restricted-tower bound follows directly
from A.56.  Here $g=2$,
$\operatorname{ord}_{19}(3)=18$ and
$v_{19}(3^{18}-1)=1$, so its constant $c_{3,19}$ is one
and $M=\max\{4,2\}=4$.
Also $e_7(3)=v_7(3^6-1)=1$ and $v_7(19-1)=0$.
Finally $b_{19}(7)=0$: the group
$\operatorname{GL}_4(\mathbb F _{19})$ has order prime
to seven, since $\operatorname{ord}_7(19)=6>4$.
Thus A.56 gives
\[
                    B=0+1+0+\lfloor\log_7 4\rfloor=1.       \tag{A.60.12}
\]
The actual cover $D$ has trivial normal $3$-subgroup
and abelian $19$-quotient, so it belongs to that class
and (A.60.1) proves sharpness.

The finite calculation has been independently reproduced
in \texttt{a50\_cover19\_order\_jump.py},
with its retained JSON and output.  That implementation
checks the cyclotomic relation and conjugate multiplicities,
computes the norm both as a resultant and as a multiplication
determinant, and checks irreducibility and the exact modular
order.  Its geometric premise is precisely A.50.
The calculation is finite and does not test an infinite
family of covers.
\end{proof}
\paragraph{Check record.} Checked only in this writing, 18 November 2026, GPT-6 Astra (OpenAI), ChatGPT Pro mode; separate effort setting not exposed. Complete proof written and read by its named authoring reader; see CHECK\_EVENTS.tsv for additional second readings and computation scope.


\paragraph{What an iteration would still have to prove.}
An orbit of length seven on an isotypic decomposition is
not by itself a factor of $7^2$ in Frobenius order.
If an invertible operator cyclically permutes $m$ nonzero
summands and its $m$th power restricts to $S$ on one
summand, then its order is $m\operatorname{ord}(S)$:
an identity power must first be divisible by $m$, and
the remaining condition is exactly an identity power
of $S$; the other restrictions are conjugate.
For $m=7$, a second factor of seven therefore requires
a factor of seven in this return action.
A.60 does not compute that return action on a further
alternating retained-prime cover.  Equation (A.60.3)
rules out obtaining the second factor by a further
pure pro-$3$ tower alone, and (A.60.12) requires any
successful construction to leave the entire class with
normal $3$-subgroup and abelian $19$-quotient.

\subsection{Where a second factor of seven could first occur}

\begin{claim}{A.61}{Proved}{The small character orbits cannot produce the next arithmetic factor}
Let $C/k$ be a smooth proper geometrically connected curve
of genus two over a finite field of characteristic different
from $3$ and $19$, with a rational base point.  Put
$\Delta=\pi_1(C_{\bar k})^{\{3,19\}}$ and
\[
 N_1=\Psi_3(\Delta),\qquad N_2=\Psi_{19}(N_1),
 \quad A=\Delta/N_1,\quad V=N_1/N_2,
 \quad G=\Delta/N_2,
\]
using the full mod-prime kernels of A.59.
Let $C_1,C_2$ be the associated pointed covers over $k$,
and let $\Gamma_7\simeq\mathbb Z_7$ be the canonical
pro-seven factor of $G_k$ acting through the chosen point.
Then $\Gamma_7$ acts trivially on $G$.

There is a splitting $G=V\rtimes A$, with
$A\simeq(\mathbb Z/3)^4$ and
$V\simeq(\mathbb Z/19)^{164}$.
Over $F=\mathbb Q_3(\mu_{19},\mu_3)$, the nontrivial
$V$-characters in an irreducible $G$-representation
form one $A$-orbit of length $d\in\{1,3,9,27,81\}$.
That representation has dimension $d$ and occurs in
$H^1(C_{2,\bar k},F)$ with multiplicity $2d$.
On this multiplicity space the order of the actual
$\Gamma_7$-action is at most
\[
                 7^{\,1+\lfloor\log_7(2d)\rfloor}.
                                                           \tag{A.61.1}
\]
The $V$-invariant part has trivial $\Gamma_7$-action.
Consequently
\[
 v_7\!\left(\operatorname{ord}
    (\operatorname{Frob}_k\mid H^1(C_{2,\bar k},\mathbb F_3))
          \right)\le3.                                   \tag{A.61.2}
\]
Any factor $7^2$ at this level must occur in an orbit
of length at least nine, and any factor $7^3$ must
occur in an orbit of length at least twenty-seven.
For the curve and cover of A.60, its canonical second
level therefore satisfies $1\le b_2(7)\le3$.

There is also a smaller-tower obstruction.  Start with
any pointed geometric cyclic-three cover $Y\to C$
over $k$, then take a pointed abelian nineteen-power
cover $Z\to Y$ over $k$, and finally any pointed
cover $W\to Z$ over $k$ with geometric three-group
Galois closure.  The abelian cover may have arbitrary
rank and exponent.  Then
\[
 v_7\!\left(\operatorname{ord}
    (\operatorname{Frob}_k\mid H^1(W_{\bar k},\mathbb F_3))
          \right)\le1.                                   \tag{A.61.3}
\]
Thus this extra cyclic-three step still cannot produce
the second factor of seven.  None of these bounds
asserts that a second factor actually occurs on the
full canonical tower or settles the unrestricted A12
criterion.
\end{claim}
\begin{proof}
\medskip\noindent\textbf{The geometric deck multiplicities.}
We first record precisely the deck-representation
calculation that will be used.  If $X\to C$ is a
connected finite etale Galois cover with group $H$
of order involving only $3$ and $19$, then for any
characteristic-zero splitting coefficient field $L$
of residue characteristic $3$ or $19$,
\[
 H^1(X_{\bar k},L)\simeq
                    L^{\oplus2}\oplus L[H]^{\oplus2}
                   \quad\text{as $L[H]$-modules}.          \tag{A.61.4}
\]
Only the geometric representation is compared in this
formula; no Frobenius is transported to characteristic
zero.  The proper prime-to-characteristic comparison
proved in A.4 lifts the finite cover and identifies
its geometric cohomology, with its deck action, with
that of a complex covering.  A cell structure on the
genus-two base has one vertex, four one-cells and one
two-cell.  The lifted cellular complex has modules
$L[H],L[H]^4,L[H]$ in degrees zero, one and two.
Both cohomology in degree zero and that in degree two
are the trivial representation, since the covering
is connected and the deck action preserves orientation.
The Euler identity in the representation ring gives
(A.61.4), and semisimplicity in characteristic zero
turns that identity into the displayed isomorphism.
In particular, every nontrivial absolutely irreducible
representation of degree $d$ has multiplicity $2d$;
the trivial representation has multiplicity four.

We will also use two elementary pro-prime facts,
already proved in A.59 and A.60.  For a finitely
generated pro-$r$ group, the kernel of the action of
its automorphism group on the mod-$r$ Frattini quotient
is pro-$r$.  For a free finite-rank $\mathbb Z_r$
lattice, the kernel of reduction of its automorphism
group modulo $r$ is pro-$r$.  It follows that a
pro-seven action embeds into the relevant finite
mod-$r$ image when $r=3$ or $19$.  These facts apply
to every curve and every pointed cover here, because
the full induced pro-prime topology is supplied by
the finite-core argument in A.59.

\medskip\noindent\textbf{The action on the first two deck groups is trivial.}
The mod-three Frattini quotient of the geometric
pro-three group of $C$ has rank four.  Since
$\operatorname{ord}_7(3)=6>4$, the group
$\operatorname{GL}_4(\mathbb F_3)$ has order prime to
seven.  Therefore $\Gamma_7$ acts trivially on the
entire based pro-three group of $C$.  It acts
trivially on its quotient $A$ and, by restriction,
on the pro-three group of each pointed three-group
cover of $C$.

For the full mod-three cover $C_1$, the group $A$
is elementary abelian of rank four.  The coefficient
field $\mathbb Q_{19}$ already contains $\mu_3$,
so every character of $A$ is defined over it.
Formula (A.61.4) for $C_1\to C$ decomposes its
$\mathbb Q_{19}$-cohomology into a trivial-character
space of dimension four and nontrivial-character
spaces of dimension two.  The actual $\Gamma_7$
action commutes with $A$, so it preserves these
spaces and acts linearly over $\mathbb Q_{19}$.
Its compact image preserves a lattice in each one.
Since $\operatorname{ord}_7(19)=6$, neither
$\operatorname{GL}_4(\mathbb F_{19})$ nor
$\operatorname{GL}_2(\mathbb F_{19})$ has order
divisible by seven.  Reduction is injective on a
pro-seven subgroup, so the action on all these
spaces is trivial.  It is then trivial on integral
and mod-nineteen cohomology and, by the Frattini
kernel fact, on the whole based pro-nineteen group
of $C_1$.

The arithmetic action on $G$ is consequently the
identity on both $A$ and $V$.  An automorphism of
an extension of $A$ by the abelian group $V$ that
is the identity on $A$ and on $V$ is described by
a crossed homomorphism $A\to V$: the difference
between the image of a chosen lift of $a$ and
that lift belongs to $V$, independently of the
chosen lift.  Such automorphisms form a subgroup
of the additive nineteen-group of all maps
$A\to V$.  A pro-seven group has no nontrivial
image there.  This proves that $\Gamma_7$ acts
trivially on $G$, not merely on its two factors.

The ranks follow from the full induced topology
and the unramified genus formula:
$|A|=3^4$, $g(C_1)=1+3^4=82$, and hence
$\dim_{\mathbb F_{19}}V=164$.
The extension of $A$ by $V$ splits because
$H^2(A,V)=0$: multiplication by $|A|$ is invertible
on $V$, and averaging contracts the positive-degree
finite-group cochain complex.  More explicitly,
for a normalized two-cocycle $c$, the one-cochain
$b(a)=|A|^{-1}\sum_{x\in A}c(a,x)$ satisfies
$\delta b=c$ by summing the cocycle identity.

\medskip\noindent\textbf{The multiplicity space is the only place for the new factor.}
Let $F=\mathbb Q_3(\mu_{19},\mu_3)$.
The extension obtained by adjoining $\mu_{19}$
is unramified of degree
$\operatorname{ord}_{19}(3)=18$, while adjoining
$\mu_3$ is totally ramified.  Thus the residue
field of $F$ is $\mathbb F_{3^{18}}$.

Every irreducible representation of $G=V\rtimes A$
over $F$ is described directly as follows.  Its
$V$-weights form a single $A$-orbit; otherwise the
sum of the weights in one orbit would be a proper
$G$-submodule.  Fix a weight $\lambda$ and its
stabilizer $A_\lambda$.  The character $\lambda$
extends to $V\rtimes A_\lambda$ by making it
trivial on $A_\lambda$.  Since $A_\lambda$ is
abelian and $F$ contains $\mu_3$, its irreducible
representations have degree one.  Inducing the
product of $\lambda$ with such a character gives
an irreducible representation of degree
$d=[A:A_\lambda]$; its distinct weight spaces
show irreducibility and describe every irreducible
representation.  This also proves that $F$ is a
splitting field here, with no assumption about a
Schur index.

For $\lambda\ne1$, formula (A.61.4) gives the
multiplicity $2d$.  Since $\Gamma_7$ centralizes
$G$, it acts as the identity on the fixed
irreducible factor and acts on its multiplicity
space through $\operatorname{GL}_{2d}(F)$.
The image is compact and preserves a lattice;
reduction into $\operatorname{GL}_{2d}
(\mathbb F_{3^{18}})$ is injective on this
pro-seven image.  For $j\ge1$,
\[
 v_7(3^{18j}-1)=1+v_7(j),
 \qquad v_7(3^{18}-1)=1.                                \tag{A.61.5}
\]
Indeed $\operatorname{ord}_7(3)=6$ and
$v_7(3^6-1)=1$, and the elementary geometric-series
valuation gives the displayed formula.

A finite seven-power-order matrix over
$\mathbb F_{3^{18}}$ is semisimple.  Each of its
eigenvalues lies in a field extension of degree
at most the matrix dimension.  Formula (A.61.5)
therefore bounds its order by
$7^{1+\lfloor\log_7(2d)\rfloor}$, proving
(A.61.1).  This bound uses the maximum of the
eigenvalue orders, not the seven-part of the
order of the entire general linear group.

The $V$-invariant subspace is
$H^1(C_{1,\bar k},F)$, by rational cohomological
descent, and its $\Gamma_7$-action is already
trivial from the pro-three restriction above.
The possible nontrivial orbit lengths and their
bounds are consequently
\[
\begin{array}{c|ccccc}
 d&1&3&9&27&81\\\hline
 1+\lfloor\log_7(2d)\rfloor&1&1&2&3&3.
\end{array}                                               \tag{A.61.6}
\]
The full integral cohomology of a proper smooth
curve is a free $\mathbb Z_3$-lattice, and its
mod-three cohomology is its reduction.  The
congruence kernel is pro-three, so the order of
the pro-seven image is unchanged upon this
reduction.  This proves (A.61.2) and both orbit
thresholds.  A.60 supplies the lower bound one
for its particular canonical second level.

\medskip\noindent\textbf{A cyclic-three first step remains insufficient.}
For a pointed cyclic-three cover $Y\to C$, the
same $\mathbb Q_{19}$ character argument, with
trivial multiplicity four and other multiplicities
two, proves that $\Gamma_7$ acts trivially on
the entire based pro-nineteen group of $Y$.
Given $Z\to Y$ as in the statement, take its
geometric normal core over $C$.  This core is
stable under the point-section arithmetic action
and hence gives a pointed cover over the same
field.  Its normal nineteen-group is a subgroup
of the product of the three conjugate abelian
nineteen-groups, so it is abelian.  Its quotient
is cyclic of order three.

The argument with crossed homomorphisms again
makes the $\Gamma_7$ action on this finite deck
group trivial.  Write its abelian normal group
as $V'$ and its exponent as $19^n$.  The group
splits over its cyclic-three quotient by the
same averaging argument.  Every irreducible
representation over
$F_n=\mathbb Q_3(\mu_{19^n},\mu_3)$ has degree
one or three by the weight-orbit argument.
Its nontrivial multiplicity in geometric
cohomology is therefore at most six.

We have
\[
 \operatorname{ord}_{19^n}(3)=18\cdot19^{n-1},
 \qquad v_7(3^{18\cdot19^{n-1}}-1)=1,                    \tag{A.61.7}
\]
because $v_{19}(3^{18}-1)=1$ and $7\nmid19$.
Thus the same finite-matrix bound on dimensions
at most six gives seven-part at most seven,
uniformly in $n$ and in the rank of $V'$.
The trivial $V'$-weight part is cohomology of
$Y$ and has no seven-part.  Rational pullback
and trace transfer this bound from the normal
core cover to $Z$.

Finally the Frattini automorphism kernel shows
that the pro-seven action on the full based
pro-three group of $Z$ has order at most seven.
Restriction to every pointed three-group cover,
using the full induced topology, preserves that
bound.  Its mod-three cohomology therefore
satisfies (A.61.3).
\end{proof}
\paragraph{Check record.} Checked only in this writing, 18 November 2026, GPT-6 Astra (OpenAI), ChatGPT Pro mode; separate effort setting not exposed. Complete proof written and read by its named authoring reader; see CHECK\_EVENTS.tsv for additional second readings and computation scope.


\paragraph{Check and scope.}
The complete proof was independently read by the root
and finite-checks readers on 18 November 2026; both
readings passed without a requested correction.
All readers are GPT-6 Astra Pro in Pro mode, with
no separately exposed effort setting; this is
same-model-family verification and claims no priority.
It gives an upper bound and localizes the only possible
new arithmetic factor to specified multiplicity spaces.
It does not compute those spaces' actual Frobenius
matrices and does not assert that the bound is attained.
The geometric comparison A.4 is used only for the
deck-representation formula; A.59 supplies the canonical
tower and finite-core topology; A.60 supplies its
specific lower bound and the repeated elementary
Frattini argument.  No arbitrary automorphism of a
finite group is substituted for arithmetic Frobenius.

\subsection*{Additional primary-source check for the two-prime gate}
\citet[Lemma~4.4(ii), Remark~4.5]{CF_MinTsuj2025}
treats hyperbolic genus-zero curves over finite fields with
rational cusps.  Its cyclotomic injectivity assertion for
the cusp-fixing Frobenius centralizer assumes that the
retained prime set is a singleton or has finite complement.
The remark leaves removal of that assumption unresolved.
The actual journal PDF and the complete lemma proof were
opened; the author's 2024 version numbers these passages
Lemma~2.4 and Remark~2.4.1.  This source therefore supplies
no general two-prime centralizer theorem for A.54--A.57,
and its affine genus-zero scope is kept distinct from the
proper-curve point-separation question.

For A.57, the input from
\citet[Theorem~4.1.1]{CF_Raynaud1982} is precisely the
existence of the proper theta divisor of the locally exact
differentials bundle.  Its entire proof in \S4.2 and the
intersection argument in Lemma~4.3.5 were read.  A.57
separately proves its fixed-prime-power orbit estimate and
the arithmetic consequence.  Raynaud's prime-degree
new-ordinariness theorem is not restated as a theorem for
an arbitrary fixed prime power.

\citet[Lemma~3.9(ii),(iii), Corollary~3.10,
Remark~3.11, pp.~77--79]{CF_Tamagawa2003} gives the precise
intersection-theoretic qualification.
The original Raynaud lemma needs the intersection of the
divisor with the pulled-back auxiliary curve to be finite.
Tamagawa supplies the hyperplane-cut proof for every
prime-to-$p$ integer $N$ and corrects the sufficient numerical
threshold in Raynaud's Theorem~4.3.1, for example to
$\ell+1\ge(p-1)3^{g-1}g!$.
The full statements and proofs of Lemma~3.9 and Corollary~3.10,
and the correction itself, were opened.
A.57 already proves the proper-intersection version and
does not use Raynaud's original threshold; its mathematical
argument is unchanged by this source correction.

The focused retry of the Stix journal PDF used the direct
Wiley PDF and ePDF routes, the legacy Hindawi download,
the OUP article-PDF route, and an older author archive.
They returned access errors or a missing page.  The
current author link still provides Proposition~2.7 in the
18 April 2005 text; the exact correspondence with the
journal citation Proposition~2.3 remains unverified.
No proof in the direct characteristic-$p$ route depends
on that unverified numbering correspondence.

\medskip\noindent\textbf{Fixed-field class towers are not all geometric covers.}
The actual publisher offprint of \citet[Section~2, pp.~168--170]{CF_SuthVoloch2019}
was opened and that section read in full.
Its recovery assertion from all iterated Hilbert-class-tower
$L$-functions is explicitly Conjecture~2.2.
The discussion on p.~169 distinguishes multiplication-by-two
geometric covers from that fixed-ground-field tower and notes
that their Prym factors need not descend individually.
The comment on p.~168 that solvable covers suffice for the
known full-group theorem is attributed to personal communication;
it is not a two-prime theorem proved in this source.
No result from this source is used as a premise in A.53--A.58.

\subsection{Independent re-proof and the configuration proof comparison}

The independent proof was written from the statement of A.3 and the
hypothesis of A.2 before either earlier proof or any earlier audit was
opened. Its preserved source is \path{work/configuration/blind_proof.tex},
with SHA-256
\begin{center}\footnotesize\ttfamily
d435494f67de82bd06f08a04cc519501e739a7abf0f88dc3b8b42ef5c96c6ac9.
\end{center}
The comparison then checked every line 1--140 of
\path{background/A1_checkpoint03/sections/A_configuration.tex}; the
line-level record is \path{work/configuration/S1_line_comparison.tsv}.
This is an independent re-proof by GPT-6 Astra Pro, followed by a
targeted line-by-line comparison. It is not a claim of review by a
different model family. The run is in Pro mode; a separate effort
setting is not exposed.

The proofs agree. The independent proof detects finite fibre quotients
inside finite automorphism groups rather than first passing to the
whole profinite automorphism group. To make this version completely
self-contained, we record its group-theoretic details.
Let $P$ be a nonabelian finitely generated free pro-$r$ group, and let
$U_n$ be a decreasing characteristic open basis. Put $P_n=P/U_n$ and
\[
 Z_n=\{x\in P:[x,P]\subseteq U_n\}.
\]
The finite group
$\ker(\operatorname{Aut}(P_n)\to\operatorname{Aut}(P_n/\Phi(P_n)))$
is an $r$-group: it acts freely on the ordered lifts of a basis of
$P_n/\Phi(P_n)$, a set of cardinality $|\Phi(P_n)|^d$.
Each lift generates $P_n$, since otherwise its generated subgroup is
contained in a maximal subgroup and its image cannot span the
Frattini quotient. Freeness of the action and the orbit formula give
the claim. Thus a group acting through an $r$-group on $P/\Phi(P)$
has $r$-group image in every $\operatorname{Aut}(P_n)$.

The centre of $P$ is trivial. For completeness the Magnus homomorphism
sends a free basis $x_i$ to $1+X_i$ in
$1+(X_1,\ldots,X_d)\mathbf F_r\langle\!\langle X_1,\ldots,X_d\rangle\!\rangle$.
Its target is pro-$r$ under degree truncation, so the map exists by
the universal property of $P$. For every finite $r$-quotient $Q$,
substitution $X_i\mapsto q_i-1$ recovers the quotient in
$\mathbf F_r[Q]^\times$. The augmentation ideal is nilpotent:
choose a central element $z$ of order $r$, apply induction to
$Q/\langle z\rangle$, and use $(z-1)^r=0$; if the quotient
augmentation ideal has exponent $m$, the original ideal has exponent
at most $mr$. Hence the substitution is defined on formal series,
and separation by finite quotients proves injectivity of Magnus.
A homogeneous polynomial commuting with $X_1$ has only the monomial
$X_1^n$: a monomial beginning in another letter would contribute to
$fX_1$ but not to $X_1f$; factor out $X_1$ and repeat.
Commutation with $X_2$ eliminates every term of positive degree.
A central Magnus unit therefore equals $1$.

Consequently $\bigcap_nZ_n=Z(P)=1$. Given an open normal $V\subset P$,
compactness implies $Z_n\subset V$ for some $n$. The finite
$r$-group image of $H$ in $\operatorname{Aut}(P_n)$ has kernel on
the fibre contained in $V$. This proves A.2 and its domination
assertion without assuming residual-$r$ properties of the base.
It also explains why using one prescribed finite fibre quotient
alone would be insufficient: its centre could be lost. The earlier
proof avoids that loss by using the whole centreless $P$, and is
equally valid. The earlier proof supplies a more explicit argument
for Magnus injectivity than the initially saved blind proof; the
argument just given incorporates and verifies that detail.

For A.3, the independent proof wrote the peripheral homology module
explicitly as
\[
 M_A=\mathbf F_r[A]/\mathbf F_r\Bigl(\sum_{a\in A}e_a\Bigr).
\]
The vector in the denominator is nonzero in every characteristic.
The monodromy is trivial on $M_A$ and on $H_1(C',\mathbf F_r)$,
and hence consists of shears. This verifies the earlier assertion
even when $r\mid|A|$; no averaging or equivariant splitting appears
in either proof. The earlier choice $A=Q/R$ is sufficient; the blind
proof permits any larger proper quotient, which gives the same
conclusion and can dominate several fibre quotients simultaneously.

Both proofs use the normal cover with deck group $A\times A$, the
identity $\Lambda\cap\pi_1(U_A)=\pi_1(C'-a^{-1}(P))$, and the finite
core over the cosets of that normal subgroup. The normal core has
$r$-group kernel over $A\times A$, because it embeds in a product of
conjugate $r$-group quotients. The image of $K_E$ lies in this kernel,
so the quotient factors through $E_{\mathrm{cr}}$. The simultaneous
finite image of the maps to $T$ and $Q$ verifies that fibre domination
persists after completion. Finally both proofs derive middle and
right exactness from the finite-quotient universal property of
pro-$\Sigma$ completion. There is no appeal to left exactness of a
completion functor. We found no error in A.2 or A.3.

\paragraph{Check record.}
A.2 and A.3 were also checked in this writing, by independent
re-proof before access to the earlier proofs and by subsequent
line-by-line comparison, on 18 November 2026. Their earlier checks
are retained in the inherited check ledger. The direct
characteristic-$p$ argument A.45--A.48 was checked only in this
writing; it is logically separate from this blind comparison.

\subsection{Opened sources for the configuration arguments}

The following records specify the statements actually used. A reference
to a source theorem is distinguished from the additional deduction
made here. All the attached sources in this table were opened in this
writing. The bibliography marks the checked source versions as verified;
the unresolved journal-number cross-reference is marked separately.

\paragraph{SGA 1, Expos\'e III, Theorem 7.3 and Corollary 7.4.}
Theorem 7.3 assumes a complete local ring $A$ with residue field $k$
and a smooth projective $k$-scheme $X_0$ satisfying
\[
 H^2(X_0,T_{X_0/k})=H^2(X_0,\mathcal O_{X_0})=0.
\]
It gives a smooth projective $A$-scheme reducing to $X_0$; the
following paragraph also treats a prescribed smooth truncated lift.
Corollary 7.4 states the smooth proper curve case. These are at
original margin pages 85--86 in the opened SMF source. They support
the curve lift in A.4, since a curve has no coherent cohomology in
degree two. The marked point is an additional elementary step:
formal smoothness lifts it through the nilpotent thickenings, and
properness and completeness give the compatible $A$-section. Neither
statement is used in A.45--A.48.

\paragraph{SGA 1, Expos\'e XIII, Corollary 2.9.}
The source assumes a smooth proper morphism of finite presentation
$X\to S$, its open complement $U$ of a relative normal-crossings
divisor, and the primes invertible at every residue characteristic of
$S$. A finite locally constant constructible sheaf of groups with
those primes, tame relative to the pair, has locally constant
constructible first direct image and is cohomologically proper through
degree one. This is a statement about finite torsors and their base
change. In the strictly henselian lifting situation of A.4, apply
it to every constant finite prime-to-$p$ group; the resulting torsor
comparison gives the prime-to-$p$ fundamental-group isomorphism.
This conclusion is a deduction from the quoted cohomological statement,
not a configuration exactness assertion in the source.

\paragraph{SGA 1, Expos\'e XIII, Section 2.10 and Lemma 2.11.}
Section 2.10 constructs specialization by strict localization of the
base and a chosen \'etale path. Changing that path changes the
homomorphism by an inner automorphism. Under the smooth proper-pair
hypotheses of Corollary 2.9 the prime-to-residue-characteristic
comparison is an isomorphism. Lemma 2.11 assumes a proper morphism of
finite presentation with relative normal-crossings boundary, a
specialization of geometric base points, and compatible boundary
points under generization; it says that the image of the first tame
inertia subgroup is an inertia subgroup at the second boundary point.
The opened text is at original margin pages 390--392. In A.4 the
point section and diagonal are smooth relative divisors, and the fibre
is transverse with multiplicity one, so the fundamental groups, maps,
and specified inertia are all transported. These specialization
statements are excluded from the direct proof A.45--A.48.

\paragraph{SGA 1, Expos\'e XIII, Proposition 4.3 and Examples 4.4.}
The notation in Section 4.0 is essential. If
$K=\ker(\pi_1(X)\to\pi_1(S))$, the middle group $\pi_1'(X)$ is
$\pi_1(X)/\ker(K\to K^L)$; it retains the full base group.
Proposition 4.3 assumes connected $S$, a morphism $f:X\to S$ that
is $0$-acyclic and locally $1$-aspherical for $L$, a section of $f$,
cohomological properness in degrees at most one for finite constant
$L$-groups, constructibility of the direct-image torsor stacks, and
the indicated constructibility of direct images of finite covers.
It concludes
\[
 1\to\pi_1(X_{\bar s})^L\to\pi_1'(X)
 \to\pi_1(S)\to1.
\]
Examples 4.4 verify these hypotheses for the proper smooth case and
for a smooth proper compactification with relative normal-crossings
boundary and a section, with primes invertible on the base. The
open-pair case explicitly cites XIII 2.9. Thus the source does not
justify applying $L$ to the whole total and base groups, and does
not allow $p$ in $L$ in characteristic $p$. These points are respected
by the alternative Stix argument. They are also why the direct
central-induction proof was supplied for the literal dependency
requirement in G1(ii).

\paragraph{SGA 1, Expos\'e XIII, Proposition 5.2 and Corollary 5.3.}
Proposition 5.2 assumes a regular local scheme $X$, a normal-crossings
divisor with equations $f_i$ belonging to a regular parameter system,
and a finite \'etale cover of its complement tame along that divisor.
Let $n_i$ be the least common multiple of the inertia orders appearing
above the generic point of its $i$-th component. Then the $n_i$ are
invertible, and after adjoining $T_i$ with $T_i^{n_i}=f_i$ the cover
extends uniquely, up to unique isomorphism, to a finite \'etale cover
of the entire chart. The proof uses normality for uniqueness and
purity for extension. Corollary 5.3 computes the strictly local tame
group as a product of prime-to-residue-characteristic procyclic
inertias. These statements occur at original margin pages 428--431.

No root stack appears in this source. The root-stack formulation in
A.48 is proved from these chart statements: extend the cover on
each chart, extend the descent isomorphisms by uniqueness, and descend
the resulting equivariant finite \'etale covers. The boundary is the
root gerbe of the normal line. Its stabilizer is not assumed to inject
into its fundamental group; injectivity of the surviving infinite
$r$-inertia is proved separately by A.47 and the configuration fibre
injection. The same local chart, with one relative parameter, proves
smoothness over the second curve in A.46.

\paragraph{SGA 1, Expos\'e X, Proposition 1.2.}
For a proper separable morphism to a locally noetherian base, its Stein
factor is finite \'etale. Here ``separable'' means flat with
geometrically reduced fibres (Definition 1.1). This is the statement
at original margin page 262, also identified in its editorial note
with EGA III 7.8.10(i). It applies to the normalization in A.46,
which the tame local charts show to be smooth proper over the base.
The result, not a nonexistent ``Corollary 1.3'', supplies the finite
\'etale family of connected components used in that proof.

\paragraph{Stix, the opened Proposition 2.7.}
In the attached arXiv version and the author-hosted text dated
18 April 2005, the hypotheses are a connected locally noetherian
scheme $S$, a smooth proper genus-$g$ curve with $n$ disjoint ordered
sections, $2-2g-n<0$, and a set $L$ of primes invertible on $S$.
For the complement $C^0$ of the sections, the source defines
$\pi_1^{(L)}(C^0)$ by completing only the kernel of
$\pi_1(C^0)\to\pi_1(S)$. Proposition 2.7 gives the exact sequence
\[
 1\to\pi_1(C^0_{\bar s})^L\to\pi_1^{(L)}(C^0)
 \to\pi_1(S)\to1
\]
and identifies it as the pullback of
$1\to\operatorname{Inn}(P)\to\operatorname{Aut}(P)
\to\operatorname{Out}(P)\to1$ through the base monodromy.
Its conjugation map restricts to the inner action of the fibre.
Taking $L=\{r\}$ and the constant compactified deck-graph family
meets every hypothesis, even if the proper deck group has a $p$-part.
The resulting monodromy on mod-$r$ fibre homology is an $r$-group of
shears, so the finite automorphism argument in A.2 supplies the
required total pro-$r$ detection. This validates the source input in
the earlier A.13 proof. The omitted tame middle-exactness step is
supplied by A.46.

The arXiv source and author text both number this statement 2.7.
The published journal has not yet been opened directly: OUP requests
returned HTTP 403. Mochizuki--Tamagawa Remark 2.2.1 cites a related
Stix assertion as Proposition 2.3, but we do not treat that
cross-reference as verification of a journal-number match.
The source version and theorem content used here are verified; that
number correspondence is explicitly unverified. This limitation does
not affect A.45--A.48, which do not use Stix's theorem.

\paragraph{Bellingeri, Theorem 1.2 and Section 2.1.}
Theorem 1.2 is for the braid group $B(n,F)$ of a closed orientable
surface of genus $g\geq1$. It uses generators
$\sigma_1,\ldots,\sigma_{n-1},a_1,\ldots,a_g,b_1,\ldots,b_g$,
commutator $[a,b]=aba^{-1}b^{-1}$, the Artin braid relations,
relations (R1)--(R4), and the surface relation (TR)
\[
 \prod_{j=1}^g[a_j,b_j^{-1}]
 =\sigma_1\cdots\sigma_{n-2}\sigma_{n-1}^2
   \sigma_{n-2}\cdots\sigma_1.
\]
For $n=2$, (R1) and the ordinary braid relations are empty. There are
$2g$ instances of (R2), $4\binom g2$ of (R3), $g$ of (R4), and
one (TR); for genus two these are eleven relations. A.6--A.7
use exactly this convention. Section 2.1 states the Fadell--Neuwirth
fibration and its exact homotopy sequence, and observes that the
higher homotopy groups vanish for an orientable surface other than
the sphere. For our $g\geq2$, this gives the discrete short exact
sequence. Interchanging the two coordinates changes the convention
of forgetting the last strand into the second projection used here.
The attached arXiv source calls the closed-surface theorem 1.2.
The journal metadata (\emph{Journal of Algebra} 274(2), 543--563,
DOI 10.1016/j.jalgebra.2003.12.009) was checked on the publisher's
primary page; no unobserved journal re-numbering is asserted.

\paragraph{Mochizuki--Tamagawa, Definition 1.1 and Proposition 2.2.}
Definition 1.1 makes a PT-formation either a primary full formation
(a singleton prime) or the total formation of all finite groups.
Proposition 2.2 assumes connected $S$, a PT-formation $\mathcal C$
invertible on $S$, and, for every $\ell\in\Sigma_{\mathcal C}$,
$\mathcal C$-group images of both the mod-$\ell$ cyclotomic character
and the mod-$\ell$ homology monodromy of the compactified proper
curve. It gives the length-one configuration fibre exact sequence
in (i), finite generation and slimness in (ii), the base exact
sequence in (iii), algebraically closed field-extension invariance
in (iv), and complete-DVR specialization in (v).
For A.10 with the total formation over $\mathbb C$, invertibility
holds, the base is a geometric point, and the monodromy conditions
are automatic. Thus (i) supplies precisely the full profinite
exactness used there. The proposition does not assert exactness for
any proper prime set with at least two primes. Its restriction is
therefore retained; the converse classification is proved by
A.8--A.9, not attributed to the source. The attached author PDF
was opened through its full proof, printed pages 11--14 (journal
pages 85--88 for the edition beginning at page 75).

\paragraph{The additional direct-proof cohomology source.}
We opened and downloaded J.~S.~Milne's \emph{Lectures on \'Etale
Cohomology}, version 2.21, dated 22 March 2013. The relevant numbered
results are Theorem 12.7 (Leray), Propositions 14.1--14.2 (prime-to-
characteristic Picard and curve cohomology), Theorem 14.9 and Remark
14.10 (Cartan--Leray and pro-covers), Theorem 16.7 (purity for a smooth
pair), and Theorem 17.10 (proper base change). Their hypotheses hold
for our smooth varieties over an algebraically closed field with
$r$ invertible. The open-family calculation and nonabelian induction
are proved here from these inputs. The Picard argument uses the
algebraic Jacobian proof, not the illustrative complex-analytic
discussion preceding it in the notes. The downloaded PDF has SHA-256
\begin{center}\footnotesize\ttfamily
ac4f122f371d38a44c58c296b7dbf88081d89d2de2334070bff3606771c01077.
\end{center}

\subsection{Source-proof audit of the direct characteristic-$p$ route}

This audit was carried out on 18 November 2026 by GPT-6 Astra Pro,
Pro mode; a separate reasoning-effort setting was not exposed.
It distinguishes a proof newly supplied here, a source statement
actually opened, and a foundational source theorem used without
claiming an independent proof of its entire theory. The direct
route has the following internal dependencies:
\[
 \mathrm{A.45},\mathrm{A.46},\mathrm{A.47}\Longrightarrow\mathrm{A.48},
 \qquad
 \mathrm{A.45},\mathrm{A.46},\mathrm{A.48}\Longrightarrow\mathrm{A.13}.
\]
The finite-core argument is written out independently inside A.13
and A.48. The former uses A.48 only for its boundary and arithmetic
clause, preserving the inherited statement without a cycle. Thus this
route does not import all of A.3 merely to use that elementary
argument. The replacement A.15 uses A.14 and the actual
compactification and trace-formula arguments written in its proof;
A.49 then uses A.14--A.15 without A.16.

\paragraph{The precise excluded chain.}
The original characteristic-zero route uses SGA~1 III 7.3--7.4
and XIII 2.9--2.11. The alternative citation to Stix's
Proposition 2.7 is valid for relative pro-$r$ completion and
does not itself construct a characteristic-zero lift. However,
Stix's Section 2.3 invokes SGA~1 XIII 4.3--4.4, and the open-pair
example in XIII 4.4 explicitly invokes XIII 2.9. The literal
G1(ii) requirement therefore cannot be met by relabelling that
alternative as a proof free of the excluded source input.
Neither Stix's theorem nor XIII 4.3--4.4 is used in A.45--A.48.

\paragraph{A.45: formal cohomology and the algebraic curve inputs.}
Theorem 12.7 of \citet{MilneEtale} is the Leray spectral sequence
for a morphism of schemes and an abelian \emph{\'etale} sheaf. Its
printed proof, pages 82--83, is the derived-functor spectral
sequence for direct image followed by global sections; direct
image preserves injectives. Theorem 14.9, pages 96--97, is the
Cartan--Leray sequence for a finite \emph{\'etale} Galois cover.
Its proof is the same formal construction, using descent and
group invariants. Remark 14.10 passes to continuous cochains on
an inverse tower. These statements contain no specialization or
characteristic-zero comparison. We use them to obtain the two
low-degree Leray sequences and the injection of continuous group
$H^2$ into \emph{\'etale} $H^2$; the latter is proved in A.45 by
the vanishing of $H^1$ on the universal finite-cover tower, not
by presuming the total space to be a $K(\pi,1)$.

Milne's Propositions 14.1--14.2, pages 89--91, provide Picard
divisibility, prime-to-characteristic torsion, and the cohomology
of a proper smooth curve. The proof of Proposition 14.1 first
illustrates the assertion over $\mathbf C$, then gives a separate
algebraic route through the Jacobian. We use the latter route.
To verify that this does not hide a lifting argument, we also
opened \citet[Theorem 8.2 and Remark 8.4]{CF_MilneAV1986} in the
author's corrected 2 January 2022 text, printed pages 13--14.
Theorem 8.2 proves algebraically that $[n]$ has degree $n^{2g}$
and is \emph{\'etale} when $n$ is invertible. Its proof uses a
symmetric ample line bundle, the algebraic theorem of the cube,
intersection numbers, and the differential $d[n]=n$ at the
origin. Remark 8.4 then determines the free rank-$2g$ torsion
module. No comparison with a complex torus is used in that
argument. The author's preface says that the numbering agrees
with the 1986 article; we opened the corresponding original
article PDF as well, but rely for the readable proof on the
corrected author text.

For completeness, the curve-cohomology calculation in Milne
Proposition 14.2 follows from Kummer and Theorem 13.7. We opened
the latter proof, pages 86--88: it uses the Weil-divisor sequence,
Hilbert 90, Tsen's theorem for a function field of one variable
over an algebraically closed field, and the strictly henselian
discrete-valuation version of the algebraic cohomological
dimension calculation. This is an algebraic input valid in
positive characteristic. In A.45, the degree description makes
the pullback on $H^2(S,\mathbf F_r)$ zero under a degree-$r$
\emph{\'etale} cover. Such a cover exists from a nonzero
$r$-torsion Picard class, and the argument remains valid after
each cover because the base genus remains positive.

\paragraph{A.45: purity and proper base change, with open base change derived.}
Milne Theorem 16.7, pages 111--112, is purity for a smooth pair
over the algebraically closed ground field with coefficients of
invertible order. The opened proof reduces \emph{\'etale} locally
to a standard coordinate pair by Lemma 16.8 and starts from the
one-dimensional local calculation of Proposition 14.3. The latter
is the Kummer calculation at a henselian discrete valuation ring.
We use only this smooth-pair theorem; no general absolute-purity
conjecture appearing elsewhere in the old notes is assumed.

Milne Theorem 17.10, pages 116--117, is proper base change for
torsion sheaves and arbitrary base change. It is applied only
to $C_0\times S\to S$ and the proper finite union of labelled
sections. Its proof reduces to the stalk theorem 17.7. The
following source discussion describes the algebraic proof through
the relative Picard functor for curves and induction on relative
dimension, and refers to the detailed proof in Artin's treatment.
We treat proper base change as an established foundational theorem;
we do not claim to have independently reproved that theorem.
The topological discussion at the start of Milne's Section 17 is
an analogy, not an input to the algebraic statement.

The conclusion for the nonproper map $f:U\to S$ is proved in
A.45 from the localization triangle, the proper product and the
proper labelled boundary. In particular the proof does not invoke
smooth base change for $f$, an open homotopy sequence, or the
cohomological properness result SGA~1 XIII 2.9. Fibre compatibility
of localization is checked at $t=0$ by the local Kummer residue
of $t$, which is one and remains one under transverse restriction.
The two constant pieces of $R^1f_*\mathbf F_r$ therefore give an
elementary $r$-group of shears. The sum of the labelled point
classes is a map with every coefficient one; its surjectivity
uses $m\ge1$, not invertibility of $m$.

The remaining nonabelian step is supplied in full in A.45.
For a central subgroup of order $r$, a group-extension factor set
produces the obstruction in continuous $H^2$. The Leray filtration
and finite affine $r$-group representations kill that obstruction
after an $r$-cover of the base. The difference between two lifts
is a central $H^1$ class, which extends after a further $r$-cover.
This induction does not call on any nonabelian base-change theorem.

\paragraph{A.46: checking the proof of the tame-chart source.}
SGA~1 XIII Proposition 5.2 and Corollary 5.3 were opened with
their proof, original margin pages 428--431. The source proof
of 5.2 uses regularity of the Kummer chart from 5.1, uniqueness
of normalization, purity away from codimension one, and the
discrete-valuation Abhyankar lemma SGA~1 X 3.6. We opened the
latter proof, original margin pages 279--280: the composite
inertia embeds in a product of cyclic tame inertia groups; if
one order is divisible by the other, its projection to the
larger cyclic group is an isomorphism. This is a finite cyclic
group argument. The specialization discussion following that
lemma is not used.

The purity cited by XIII 5.2 is SGA~2 XIV 1.11. We downloaded
and opened its actual statement and proof, together with SGA~2
X 3.4, to which it reduces. In the 2005 SMF edition these are
printed pages 169 and 94--97, respectively. The regular case
extends covers across a subset of codimension at least two.
The proof in dimension two extends the finite algebra, uses
depth and the Auslander--Buchsbaum formula to make it locally
free, and then the discriminant ideal excludes ramification
supported only at the closed point. Higher dimension follows
by completion and induction on a regular parameter, with
effective descent. These are algebraic local arguments, with
no characteristic-zero lift and no SGA~1 XIII 2.9--2.11 input.
In particular the relevant purity theorem is the homotopical
theorem XIV 1.11, not the separate cohomological theorem XIV 1.10
on the same page.

The other source input in A.46 is SGA~1 X Proposition 1.2,
original margin page 262, for the finite \emph{\'etale} Stein
factor of a proper separable morphism. Its proof reduces to
a complete local base and the coherent $H^0$ base-change and
flatness argument. Here the tame charts make the normalization
smooth proper over the curve, so its fibres are geometrically
reduced and every hypothesis is met. The source proof cites
EGA~III 7.8.10(i), not an open-family specialization theorem.
The commuting transitive-action argument proving normality of
the fibre image and the quotient from the base is supplied in
A.46 itself.

\paragraph{A.47 and A.48: inertia and boundary descent.}
A.47 uses only the algebraic Picard/Kummer inputs just checked.
It constructs branched cyclic covers of order $r^n$ after a proper
cyclic $r$-cover, and takes their conjugate compositum. Thus the
unbounded order of actual cusp inertia is established without a
surface presentation or lifting to characteristic zero.

For A.48, the language of root stacks is a description of
\emph{\'etale} descent from the charts in SGA~1 XIII 5.2. The
source itself contains no root-stack assertion. Unique extension
extends the transition isomorphisms and their equivariance;
ordinary finite \emph{\'etale} descent glues them. Descent from
the boundary gerbe for covers with trivial stabilizer action gives
the right exact boundary sequence. Its finite stabilizer is not
assumed injective. The inverse limit, right exactness of the
geometrically pro-$\Sigma$ quotient, and the actual transverse
fibre map are written out in A.48. Full infinite inertia follows
from A.47 and the geometric fibre injection proved at the start of
A.48 directly from A.45--A.46. The final two conjugations are
group-theoretic and preserve the full arithmetic quotient.

\paragraph{What the dependency audit establishes.}
The direct proofs have no edge to curve lifting, to the
characteristic-zero configuration theorem, to Stix's exactness
theorem, or to SGA~1 XIII 2.9--2.11. Their foundational inputs
remain algebraic purity, proper base change, Picard theory,
Kummer theory, descent, and formal spectral sequences. This is
a proof-dependency statement, not a claim that every foundational
theorem has been independently reproved here. The unresolved
published-number match for Stix's proposition is therefore a
bibliographic limitation of the alternative route and does not
affect the direct proof or its stated prime-set range.

\subsection{Source and predecessor audit completed in the successor run}

\paragraph{Scope of this reading.}
The configuration checker read the complete inherited statements and proofs
of A.34, A.36--A.41 and C.5--C.13 on 18 November 2026. The replacement
A.35 was checked separately. This was a targeted mathematical and source
dependency reading by GPT-6 Astra Pro in Pro mode; a separate effort setting
was not exposed. It was not an independent cross-model review. The source
statements and proof loci actually opened are recorded in
\texttt{configuration/source\_records.tsv}; individual current check events
are recorded in \texttt{configuration/CHECK\_EVENTS.tsv}. Download hashes
are retained in \texttt{configuration/literature\_download\_manifest.json}.

\paragraph{Literal source corrections.}
Three source descriptions required correction. In A.37 the quotient
retained by Wakabayashi is the group of the base curve, which need not be
proper. The singleton-prime formation and the source's distinction between
existence and uniqueness of a matching point are unchanged.
In C.12, Sun's defining morphism is smooth and induces a surjection on
fundamental groups; the definition does not require that morphism itself
to be surjective. In the C.9 table, Sun 2025 Theorem~B permits an arbitrary
characteristic-zero field, with its additional boundary and cyclotomic
hypotheses. Theorems~A and C impose the stated arithmetic-field hypotheses.
The corresponding complete replacement claim blocks are supplied.
\citep[Theorem~3.9, Corollary~3.10 and Remark~3.10.1]{LA_Wakabayashi2016}
\citep[Definition~3.1 and Theorems~A--C]{LA_Sun2025}

\paragraph{Precise reconstruction inputs.}
The disputed full fibre statement is original Mochizuki
Proposition~1.8(ii), or Proposition~1.6(ii) in the separate LaTeX text.
The original proof's three direct uses occur in the proof of (iii), before
Proposition~1.12, and before Proposition~3.8. The first two feed central
and abelian constructions; the third feeds fixed-$l$ Lie constructions.
A.35 supplies those quotient-level replacements, compatibly on finite
proper pullbacks. The full mixed-prime insertion, and a slimness inference
based only on that insertion, remain unavailable. This finite census does
not independently reprove the unrelated Lie-algebra or global Isom results.
The 2014 Comments were read in full and do not amend the disputed
statement.\citep[Propositions~1.8,1.12,1.14--1.16,1.19,3.8--3.9;
Theorem~3.10]{LA_Mochizuki2007}\citep{LA_Mochizuki2014}

The central, abelian, and fixed-$l$ inputs imported in ST2009
Theorem~2.2, Corollary~2.3, Proposition~2.4, and Theorem~2.6 reappear
as ST2017 Theorem~3.2, Corollary~3.3, Proposition~3.4, and
Theorem~3.6. Thus C.5 and C.7 record A.35 explicitly as their repair
input. The other hypotheses of the published reconstruction theorems
remain. In particular, eliminating exceptional point labels does not
eliminate complementary-primary Jacobian torsion or the corresponding
correction to the reconstruction of functions.
\citep[Section~2, Theorem~3.6, Theorem~3.9 and Corollary~3.10]{LA_ST2009}
\citep[Section~3 and Theorems~4.22--4.23]{LA_ST2017}

\paragraph{Numbering and an explicit verification limit.}
Both source texts were opened to check that ST2009 arXiv
Theorem~2.7 is journal Theorem~2.6, and arXiv Remark~3.13 is journal
Remark~3.14. The latter is the conditional full-configuration route;
it is not a premise of the published Isom proof. The source called
ST2017 throughout this paper is the 2017 preprint of the article
published in 2018. Its cited body numbers are taken from the opened
supplied source.
\citep[Theorem~2.6 and Remark~3.14]{LA_ST2009}\citep{LA_ST2017}

Stix's arXiv and author-hosted Proposition~2.7 and Section~2.3 were
opened and checked. The relative construction retains the full base
group. The exact correspondence with the journal numbering remains
unverified: the publisher full-text link returned a purchase page or
HTTP~403. Mochizuki and Mochizuki--Tamagawa cite Proposition~2.3;
that citation is recorded as their citation, not promoted to a direct
journal-text verification. No mathematical argument in the direct
characteristic-$p$ chain depends on this unresolved bibliographic detail.
\citep[Section~2.3 and Proposition~2.7, author version]{LA_Stix2005}

\paragraph{Later configuration sources.}
Hoshi--Minamide--Mochizuki's Theorem~2.5 assumes that $\Sigma$ is a
singleton or all primes and that the characteristic lies outside
$\Sigma$. Accordingly, the singleton case includes positive
characteristic, whereas the total case is characteristic zero. Its
geometric reconstruction of numerical invariants and generalized fibre
subgroups cites Mochizuki--Tamagawa Proposition~2.2 and retains that
formation restriction. It neither supplies a proper mixed-prime
formation nor addresses the finite-field arithmetic point-section map.
\citep[Proposition~1.2(iv) and Theorem~2.5]{CF_HMM2022}

The closer positive-characteristic relative predecessor is
Hoshi--Sawada--Tsujimura: Definition~1.6(ii) completes the kernel of
the map between the full tame groups, and Corollary~1.14 identifies
the prime-to-characteristic geometric fibre when it is centre-free.
The good-pair and good-morphism hypotheses are retained. Their base
group stays full tame. The complete target proof in Section~1 was
read; it does not justify simultaneous completion of that base.
\citep[Definitions~1.1,1.3,1.6 and Corollary~1.14]{LA_HST2024}

Hoshi--Tsujimura's 2026 point-section theorem uses affine varieties,
the full \'etale group, and an Artin--Schreier argument over a
characteristic-$p$ field admitting a degree-$p$ Galois extension.
Theorem~1.8 and its proof were read. Its hypotheses do not cover
proper curves in the geometric prime-to-$p$ quotient.
\citep[Theorem~A(i) and Theorem~1.8]{LA_HT2026}

Holzschuh--Schmidt--Stix 2026 gives a further characteristic-zero
predecessor. Theorem~5.6 identifies the full profinite \'etale homotopy
fibre of a geometric fibration over a normal qcqs characteristic-zero
base. Its proof uses Friedlander's fibre comparison. Theorem~8.4 has
a normal connected excellent characteristic-zero base with sub-$p$-adic
generic field and imposes the stated flatness, normality, generic-fibre
and surjective-smooth-locus hypotheses on the source. These results are
relevant to the full-group characteristic-zero setting of C.10--C.13.
They give no arbitrary-$\Sigma$ completion statement. Remark~5.7
refers the arbitrary-characteristic extension to a work still listed
as in preparation; no conclusion of that unprovided work is used here.
\citep[Definition~5.5, Theorem~5.6, Remark~5.7 and Theorem~8.4]{CF_HSS2026}

\paragraph{Nearest opened predecessors for the main chain.}
The following identifies the nearest verified source input, rather than
asserting that the full present statement is already in that source.
Elementary finite-group and finite-core arguments are proved in this
paper where no identical numbered predecessor was located.

\begin{center}\small
\begin{longtable}{p{0.13\linewidth}p{0.43\linewidth}p{0.34\linewidth}}
\toprule
Statements&Nearest opened source input&Additional argument here\\\midrule
A.2--A.3&
M--T Proposition~2.2 proof; Stix Proposition~2.7;
Bellingeri Section~2.1&
Finite-$r$ monodromy detection, deck-graph shear and finite core
for an arbitrary proper $\Sigma$ base.\\
A.4&SGA1 III~7.3--7.4; XIII~2.9--2.11&
Verification of the smooth proper pair and boundary maps.\\
A.5, A.13, A.48&
SGA1 XIII~5.2; original Mochizuki Remark~1.7.1;
Stix relative completion&
Actual diagonal boundary, retained inertia, arithmetic fibres
and transport in the specified mixed-base quotient.\\
A.6--A.7, A.9, A.12&
Bellingeri Theorem~1.2 and its surface conventions&
The stated finite obstruction and point-push calculations;
no identical obstruction theorem located.\\
A.8, A.14&
Finite-quotient universal properties&
Open-subgroup core and same-prime kernel arguments are given
in full; no source theorem is substituted for them.\\
A.10&
M--T Definition~1.1 and Proposition~2.2;
HMM Theorem~2.5&
Necessity of the prime-set restriction and the exact
noninjectivity obstruction.\\
A.11&
Stix Proposition~2.7 for marked families&
Labelled-cover shear and core construction.\\
A.15, A.18, A.49&
Mochizuki Remark~1.18.2 and its2014 correction;
curve cohomology and the Weil trace formula&
Retained compactification at every visible open level and
compactness for the entire section subgroup.\\
A.16--A.17&
ST2009 Proposition~1.8(vi), Definitions~1.3/1.5/1.10,
Remark~1.2 and Lemma~1.7&
Exact specialization to uniform labels and a full section
over the current finite constant field.\\
A.19, A.42--A.44&
Mochizuki Proposition~2.1(ii),(iii); degree-one cycle class&
Actual-inertia multiplier, difference cocycle, unchanged-label
evaluation, and the separate endgame.\\
A.20--A.21&
ST2009 Claim~1.9(i); ST2017 Propositions~2.12--2.13&
Proper support field suffices for the criterion, with
prime-power constants and no finite exceptional-set premise.\\
A.22--A.23&
ST2017 Proposition~1.3, Definitions~1.2/1.4/2.6/2.9,
Propositions~2.11--2.12&
All-extension equivalence and upgrade from proper support
to emptiness using the proved transport.\\
A.35, A.45--A.47&
Smooth-pair purity, proper base change, Picard/Kummer
theory; SGA1 XIII~5.2 and X~1.2&
Central torsor induction, tame Stein-factor exactness and
unbounded actual inertia, without the excluded lifting route.\\
\bottomrule
\end{longtable}
\end{center}
The source hypotheses for the foundational rows are written explicitly
in the configuration source and proof-dependency tables. They do not
silently invoke the open-family specialization assertions excluded by
G1(ii).\citep{LA_SGA1,MilneEtale,CF_MilneAV1986,Bellingeri2004,LA_MT2008,LA_Stix2005,LA_ST2009,LA_ST2017}

\paragraph{Bounded search and citation chasing.}
The web search service's system1 and system2 engines were used on
18 November 2026. The exact 30 queries, requested domain restrictions
and service names are in \texttt{configuration/S5\_search\_log.tsv}.
They include the exceptional set and cofinite primes; proper
prime-to-$p$ point-section injectivity; failure of configuration
exactness; the necessity of the primary-or-total condition; the four
named seed titles; and the Stix version-number question. Broad searches
returned many irrelevant items. Only opened primary texts support the
mathematical comparisons.

The bounded forward and backward citation chase checked ST2017's
citations of ST2009; HST2024's references to both Sa\"idi--Tamagawa
papers, M--T and HMM; Wakabayashi's use of M--T Proposition~2.2;
and HMM's actual M--T fibre-exactness citations. Primary author
publication lists were also consulted. The seven recorded routes are
in \texttt{configuration/S5\_citation\_chase.tsv}. This is not an
exhaustive bibliometric list of all citing papers, and no complete
citation-index export was obtained.

Within that bounded search and those opened passages, no independent
published universal prime-to-$p$ point-separation theorem, nor the
present full prime-set nonexactness classification, was located.
That negative search result is not a novelty or priority claim.
The conclusions supported by the audit are the explicit source scopes,
the retained hypotheses, the exact quotient repairs, and the separately
proved mathematical results of this paper.

\section{Independent existence and reconstruction questions}
\subsection{The exact relative criterion and the missing existence assertion}

We retain the whole arithmetic group throughout. A cover is
geometrically connected, and ``constant'' refers to its deck group
over the displayed finite field. A $P$-split cover has a chosen rational
point over $P$; for a constant Galois group this makes its entire fibre
above $P$ rational.

\begin{claim}{B.1}{Proved}{The relative class-number test, with its group hypothesis}
Let $C/L$ be proper of genus at least two over a finite field of
characteristic $p$, and let $D\to C$ be a nontrivial constant Galois
\'etale $P$-split cover of degree $d$ prime to $p$. Put
$B=(\ker(J_D\to J_C))^0$. Then
\[
 \frac{\#J_D(L)}{\#J_C(L)}=\#B(L)\in\mathbf Z_{>0}.
\]
If the group has a nonidentity element of odd order, the condition
$B(L)[p]=0$ directly excludes a point-section collision at $P$.
This includes every odd prime cyclic cover used in B1.
The conclusion also holds for the remaining prime-to-$p$ groups
by the independently proved A.1, but the odd-order stabilizer
argument by itself does not prove that broader group assertion.
\end{claim}
\begin{proof}
Pullback and norm satisfy $\pi_*\pi^*=[d]$. Hence
$J_C\times B\to J_D$ is an isogeny over $L$. The equality of
finite-field point counts under isogeny gives the asserted integer.
On $J_D(L)[p^\infty]$, multiplication by $d$ is invertible and
\[
 e=d^{-1}\pi^*\pi_*=d^{-1}\sum_{\delta\in G}\delta
\]
is an idempotent. Its complementary image lies in $B[p^\infty]$:
it is the image of the homomorphism $[d]-\pi^*\pi_*$, followed
by division by $d$ on this finite $p$-primary group. Thus
$B(L)[p]=0$ implies that all rational $p$-primary classes are
fixed by $G$.

Under a collision, A.31 produces the unique lift $Q'$ above the
distinct point $Q$ with $T=[Q'-P']\ne0$ rational and $p$-primary.
If a nonidentity deck transformation fixes $T$, the effective
divisors $\delta Q'+P'$ and $Q'+\delta P'$ are distinct and
disjoint, and are linearly equivalent. They define a degree-two
pencil, so $D$ is hyperelliptic. Its hyperelliptic involution is
central. An odd-prime-order deck transformation descends to an
automorphism of the quotient projective line of that same order.
It has a geometric fixed point there; its permutation of the fibre,
which has at most two points, is trivial. It therefore has a fixed
point on $D$, contrary to the free deck action. This proves the
direct criterion. A.1, proved from the configuration, boundary,
and cuspidal criteria in Part A, supplies the final broader conclusion.
The explicit even-order control in the norm-module discussion
explains why we have kept these two justifications distinct.
\end{proof}
\paragraph{Check record.} Also checked in this writing, 2026-11-18: GPT-6 Astra (OpenAI, ChatGPT Pro mode), root; separate effort setting not exposed. Read relative class-number proof, simultaneous-character formula in all genera, corrected pair bound and first-field fixed-prime reduction. Universal existence B.5 remains Open; no A.1 is used in the odd-degree pair exclusion. | Also checked in this writing, 2026-11-18: GPT-6 Astra Pro; finite\_checks agent; effort not separately exposed. Reproved relative-classnumber odd-deck criterion directly in A50, including char2 inseparabledegree2 exclusion; noA1 use.


\begin{claim}{B.2}{Proved}{Simultaneous character criterion in every genus}
Let $a:J_C(L)\twoheadrightarrow A$ be a nontrivial finite abelian
quotient of order prime to $p$, and use the Abel map based at $P$
to pull back its constant Lang quotient. Put
\[
 \Theta_{C,P,A}=\sum_{d=0}^{g-1}
       \sum_{E\in\operatorname{Div}^{+}_d(C)(L)}
                    [a([E-dP])]\quad\text{in }\mathbf F_p[A].
\]
Its relative class number is prime to $p$ if and only if
$\overline\psi(\Theta_{C,P,A})\ne0$ for every nontrivial
character $\overline\psi:A\to\overline{\mathbf F}_p^\times$.
For $A=\mathbf Z/\ell$ this is precisely B1's simultaneous
test; it is not enough to check one reduced character unless
all primitive characters form one Frobenius orbit.
\end{claim}
\begin{proof}
For a nontrivial complex character $\psi$ put
\[
 c_d(\psi)=\sum_{E\in\operatorname{Div}^{+}_d(C)(L)}
                 \psi(a([E-dP])),\qquad
 \kappa=[K_C-(2g-2)P].
\]
The Euler product, expanded by effective divisors, gives
$L_P(\psi,T)=\sum c_d(\psi)T^d$. A line bundle $M$ of
degree $d$ has $(q^{h^0(M)}-1)/(q-1)$ effective divisors in
its class. For $d>2g-2$, Riemann--Roch makes this number
independent of the class, so the nontrivial character sum
vanishes. For $g-1\le d\le2g-2$, apply Riemann--Roch to
$M$ and $K_C-M$; the class-independent term again vanishes.
This gives
\[
 c_d(\psi)=q^{d+1-g}\psi(a(\kappa))
                       c_{2g-2-d}(\psi^{-1}).
\]
In particular the coefficients with $d\ge g$ vanish modulo
every prime over $p$. This is the stated truncation at $g-1$.

On each closed fibre of the cover, its Frobenius translation
has the character eigenvalues $\psi(a_x)$ in the regular
representation. Factoring this permutation determinant and
multiplying the Euler factors proves
\[
 \frac{Z(D,T)}{Z(C,T)}=\prod_{\psi\ne1}L_P(\psi,T),
 \qquad
 \frac{\#J_D(L)}{\#J_C(L)}=
                       \prod_{\psi\ne1}L_P(\psi,1).
\]
Since $p\nmid|A|$, all characters remain distinct after
reduction in an algebraic closure of $\mathbf F_p$ and give
all reduced characters. The product is a unit exactly when
each factor is a unit. The preceding coefficient calculation
identifies the factors with the evaluations of $\Theta$.
If $A=\mathbf Z/\ell$, the rational factors are Galois
conjugates and their product is the cyclotomic norm. Over
$\mathbf F_p$, their orbits have length
$\operatorname{ord}_\ell(p)$, which proves the last assertion.
\end{proof}
\paragraph{Check record.} Also checked in this writing, 2026-11-18: GPT-6 Astra (OpenAI, ChatGPT Pro mode), root; separate effort setting not exposed. Read relative class-number proof, simultaneous-character formula in all genera, corrected pair bound and first-field fixed-prime reduction. Universal existence B.5 remains Open; no A.1 is used in the odd-degree pair exclusion. | Also checked in this writing, 2026-11-18: GPT-6 Astra Pro; finite\_checks agent; effort not separately exposed. Read Riemann–Roch character criterion; reproved genus2 exact polynomial in A50 and genus3 Cartier interpretation in B80, full general-genus proof not separately rewritten here.


\begin{claim}{B.3}{Proved}{The corrected pair estimate and its precise limitation}
For any quotient $A$ of $J_C(L)$ let $N=\#C(L)$ and let
$N_b$ be the occupancies of the Abel-image points in its cosets.
Then
\[
 \sum_bN_b(N_b-1)
 \le \frac{N^2}{|A|}
       +\frac{|A|-1}{|A|}(2g-2)^2q-N.                    \tag{B.CL.1}
\]
If all nonzero occupancies are multiples of $p$, the left side
is at least $(p-1)(N-N_0)$. In genus two, if $p$ generates
$(\mathbf Z/\ell)^\times$, a surjection with kernel
$H\subset J_C(L)$ is a class-number witness whenever
\[
                       N\ge2|H|,\qquad \ell>N.         \tag{B.CL.2}
\]
Neither $p>(2g-2)^2$ nor these estimates alone proves universal
existence of such a quotient. For $g>2$, occupancies of points
alone are not the character test in B.2.
\end{claim}
\begin{proof}
Character orthogonality gives
\[
 \sum_bN_b^2=\frac{1}{|A|}
       \left(N^2+\sum_{\psi\ne1}
          \left|\sum_{R\in C(L)}\psi(a([R-P]))\right|^2\right).
\]
Each nontrivial Lang character pulls back to a geometrically
nontrivial unramified rank-one local system: the Abel map
induces the isomorphism on geometric abelian fundamental groups.
Its $H^0$ and $H^2$ vanish, and its $H^1$ has dimension
$2g-2$ by the unramified Euler characteristic. It is a summand
of the cohomology of the finite \emph{\'etale} Lang pullback,
so its Frobenius eigenvalues have complex absolute value
$\sqrt q$, by the opened Weil theorem
\citep[Theorem~1.6 and formula~1.5.1]{WeilCurves}.
The trace formula bounds the inner sum by $(2g-2)\sqrt q$.
Subtracting $N$ proves (B.CL.1). For a positive occupancy
divisible by $p$, $N_b(N_b-1)\ge(p-1)N_b$; sum over $b\ne0$.

For $g=2$, B.2 reads $[0]+\sum_R[a([R-P])]$.
The primitive-root assumption makes $\Phi_\ell$ irreducible
over $\mathbf F_p$. Failure of the test therefore makes all
its coefficients equal. If $\ell>N$, an empty coset exists
and is nonzero, since $P$ occupies the zero coset. All nonzero
occupancies would consequently be divisible by $p$.
For a nonzero Jacobian class, at most two ordered pairs of
distinct geometric points have that difference: two different
representations give the unique degree-two pencil, and the
second pair must be $(\iota S,\iota R)$ when the first is
$(R,S)$. This proves the bound in all characteristics.
At most $2(|H|-1)$ points therefore share a coset with a
different point. Under (B.CL.2) at least two points occupy
singleton cosets, one of them nonzero. Its occupancy is not
divisible by $p$, a contradiction. Finally the unbounded
quotient-existence requirement and the effective-divisor terms
for $g>2$ have not been removed by this reasoning; the stated
limitations follow from the actual hypotheses just used.
\end{proof}
\paragraph{Check record.} Also checked in this writing, 2026-11-18: GPT-6 Astra (OpenAI, ChatGPT Pro mode), root; separate effort setting not exposed. Read relative class-number proof, simultaneous-character formula in all genera, corrected pair bound and first-field fixed-prime reduction. Universal existence B.5 remains Open; no A.1 is used in the odd-degree pair exclusion. | Also checked in this writing, 2026-11-18: GPT-6 Astra Pro; finite\_checks agent; effort not separately exposed. Read and checked genus2 pair bound and primitive-root quotient criterion; explicitS2 certificate avoids depending on mere numerical bound.


\begin{claim}{B.4}{Proved}{A finite all-constant-extension decision for each fixed prime}
Fix $(C,P)/\mathbf F_q$ and an odd prime $\ell\ne p$.
There is a finite procedure deciding whether B1 has an
order-$\ell$ witness over \emph{any} constant extension.
For each geometric nonzero character
$\chi:\pi_1(C_{\overline k})\to\mathbf F_\ell$, let $m$
be its least Frobenius period. Construct its unique $P$-split
constant cover over $k_m=\mathbf F_{q^m}$. It suffices to test
these finitely many covers at these first fields. If $A_\chi$
is Frobenius on their new geometric $p$-torsion, the exact
extension degrees preserving the witness are
\[
  n\ge1\quad\text{such that}\quad
                  \det(1-A_\chi^n)\ne0.                \tag{B.CL.3}
\]
A failed first-field test never becomes a witness at a later
field. A successful one remains successful under all
extensions of $p$-power degree.
\end{claim}
\begin{proof}
The character space is the finite vector space
$H^1_{\et}(C_{\overline k},\mathbf F_\ell)$ of dimension
$2g$, with its computable invertible Frobenius action.
Its nonzero elements, or their one-dimensional quotients,
are finite in number. An invariant character extends uniquely
to the semidirect arithmetic group by requiring it to be zero
on the chosen section $s_P$. This is exactly its geometrically
connected constant cyclic $P$-split cover. An extension permits
that same constant character exactly when its degree is a
multiple of $m$. Abelian class field theory, or the Lang
construction and the Abel isomorphism on geometric abelian
fundamental groups, identifies it with a quotient of
$J_C(k_m)$. These are the covers in B1, with the pointing fixed.

Pullback and norm split their $p$-primary Tate data, because
$p\ne\ell$. The quotient abelian variety's nonzero rational
$p$-torsion is therefore exactly the fixed space of $A_\chi$.
This proves (B.CL.3) and the first-field assertion: an existing
fixed vector remains fixed by every power. Conversely, a
constant cover over $k_{mn}$ with this geometric character and
split point is the base change of the unique split cover just
constructed, so no further twist has escaped the test.
For $n=p^a$, the polynomial identity
$X^{p^a}-1=(X-1)^{p^a}$ shows that $A_\chi^n-1$ is
invertible exactly when $A_\chi-1$ is invertible.

For an effective procedure, one may construct the finitely
many covers and compute their zeta numerators by point counts,
or compute the induced Frobenius on coherent cohomology and
its stable part by exact finite-field linear algebra. The
Artin--Schreier sequence identifies that stable part with
geometric \(p\)-torsion cohomology; nilpotent directions do
not affect $\det(1-F)$. Thus the procedure terminates for
every fixed $\ell$. It supplies no bound on the primes that
must be tested for the unrestricted existence assertion B1.
\end{proof}
\paragraph{Check record.} Also checked in this writing, 2026-11-18: GPT-6 Astra (OpenAI, ChatGPT Pro mode), root; separate effort setting not exposed. Read relative class-number proof, simultaneous-character formula in all genera, corrected pair bound and first-field fixed-prime reduction. Universal existence B.5 remains Open; no A.1 is used in the odd-degree pair exclusion. | Also checked in this writing, 2026-11-18: GPT-6 Astra Pro; finite\_checks agent; effort not separately exposed. Read fixedprime firstfield algorithm and permanence of failedfixedvectors; independently checked minimum-period deduction used B80.



\noindent\emph{Statement B.5 and its produced argument are retained in Part II.}
\subsection{The fixed norm module, full sections, and the free-orbit question}

We fix a finite field $k$ of characteristic $p$, a smooth proper
geometrically connected curve $C/k$ of genus at least two, and
$P\in C(k)$. All arithmetic groups below retain the whole $G_k$.
Put
\[
 \Omega=\pi_1(C_{\bar k},\bar P),\qquad
 \Delta=\Omega^{(p')},\qquad
 N=\ker(\Omega\longrightarrow\Delta),\qquad
 V=N^{\mathrm{ab},p}.
\]
Write $F$ both for arithmetic Frobenius through $s_P$ and for its
actions, and write $\alpha$ for its action on $\Delta$. For each
open normal $H\triangleleft\Delta$, choose the compatible pointed
geometric cover $D_H\to C_{\bar k}$. Its Jacobian norm maps define
\[
 M_P=\varprojlim_{H\triangleleft_o\Delta}
                 J_{D_H}(\bar k)[p^\infty].
\]
The action of $F$ on this entire inverse system permutes its
indices and satisfies $F(gm)=\alpha(g)F(m)$. We never assume that
it commutes with the geometric deck action.

\begin{claim}{B.6}{Proved}{The integral description of the fixed norm module}
There is a natural exact sequence of compact abelian pro-$p$ groups
\[
 0\longrightarrow V\xrightarrow{F-1}V
   \xrightarrow{\varepsilon}M_P^F\longrightarrow0.
\]
For every $F$-stable normal level $H$, projection gives a continuous
surjection
\[
 M_P^F\longrightarrow J_{D_H}(k)[p^\infty].
\]
In particular $M_P^F$ need not vanish. The first term in the exact
sequence is a Tate-module inverse limit; its Frobenius invariants
vanish, whereas its Frobenius coinvariants need not vanish.
\end{claim}
\begin{proof}
The $\alpha$-stable open normal subgroups form a cofinal directed
system. A finite pointed geometric cover descends to a finite
constant extension, so its subgroup has a finite $\alpha$-orbit;
intersecting that orbit gives a stable refinement. Finite
intersections give common refinements. For such $H$, the subgroup
generated by $H$ and $s_P(G_k)$ defines a model $D_H/k$ with
compatible point $P_H\in D_H(k)$. Put
\[
 U_H=\operatorname{pr}^{-1}(H)\subset\Omega,
 \qquad T_H=T_pJ_{D_H}.
\]
The pointed Abel map identifies $U_H^{\mathrm{ab},p}$ with $T_H$,
including the characteristic prime. The exact geometric comparison
is \cite[Proposition 9.1, Lemmas 9.4--9.5, and Corollary 9.6]{MilneJac}:
for a smooth complete curve over a separably closed field, the Abel
map identifies its full abelian fundamental group with that of its
Jacobian. The characteristic-$p$ part is explicitly included, using
Artin--Schreier cohomology and its Witt-vector extensions.

One can describe the Tate module in this comparison without treating
$[p^a]$ as an étale isogeny. Quotient the source of $[p^a]$ by
the connected component of its kernel. The resulting isogeny to the
Jacobian is étale, with geometric kernel
$J_{D_H}[p^a](\bar k)$. These covers are cofinal among the
connected étale abelian $p$-covers of the Jacobian. Their inverse
limit is $T_H$, and all constructions commute with $G_k$.
For $H'\subset H$, the map induced by
$U_{H'}\hookrightarrow U_H$ is the Jacobian norm on Tate modules:
the pointed Abel maps commute with the curve morphism and its
Jacobian norm. This is an inclusion-induced map, rather than a
transfer multiplied by the covering degree.

We next prove the $F$-equivariant identification
\[
                    V\simeq\varprojlim_H T_H.                 \tag{NM.1}
\]
Let $c:N\to\Q_p/\Z_p$ be a continuous character. Its image is
finite because $N$ is compact and the target is discrete. Choose
an open normal $O\triangleleft\Omega$ with
$N\cap O\subset\ker c$. The formula
\[
                  \widetilde c(no)=c(n)\quad(n\in N,o\in O)
\]
defines a continuous character on $NO$. It is well-defined since
two decompositions differ by $N\cap O$, and is additive since
$[O,N]\subset O\cap N$. The subgroup $NO$ is open, contains $N$,
and corresponds to an open normal subgroup of $\Delta$. A stable
refinement therefore extends $c$ to some $U_H$. Conversely, a
character of $U_H$ vanishing on $N$ factors through the
pro-prime-to-$p$ group $H$ and has $p$-primary image, so it is
zero. Restriction consequently identifies
\[
 \Hom_{\mathrm{cont}}(N,\Q_p/\Z_p)
 =\varinjlim_H\Hom_{\mathrm{cont}}(U_H,\Q_p/\Z_p).
\]
Pontryagin duality proves (NM.1), with its actual norm transition
maps and $F$-action.

At a stable level set $W_H=T_H\otimes_{\Z_p}\Q_p$. Its quotient
$W_H/T_H$ is the group of geometric $p$-primary points of the
Jacobian. The fixed subgroup $T_H^F$ is zero. Indeed an invariant
compatible $p$-division sequence would lie in the finite group
$J_{D_H}(k)$; multiplying sufficiently far down that sequence by
the exponent of its $p$-primary subgroup kills every coordinate.
Thus $F-1$ is injective on the finite free module $T_H$, and
invertible on $W_H$. There is an exact sequence
\[
 0\longrightarrow T_H\xrightarrow{F-1}T_H
 \xrightarrow{\varepsilon_H} J_{D_H}(k)[p^\infty]
 \longrightarrow0,\qquad
 \varepsilon_H(a)=(F-1)^{-1}a\bmod T_H.                       \tag{NM.2}
\]
For the image assertion, a coset $w+T_H$ is $F$-fixed exactly
when $(F-1)w\in T_H$. Its fixedness is exactly rationality over
$k$. Formula (NM.2) is also the finite-field Kummer comparison,
with cocycle convention $F(\widetilde x)-\widetilde x$.

All maps in (NM.2) commute with norms. Cofinality of the stable
levels gives
\[
 M_P^F=\varprojlim_H J_{D_H}(k)[p^\infty].                    \tag{NM.3}
\]
This supplies its compact pro-$p$ topology, without imposing a
bounded exponent on $M_P$.

For completeness, inverse-limit surjectivity requires no uniform
denominator. Given a compatible family $m_H$ in (NM.3), the sets
$\varepsilon_H^{-1}(m_H)\subset T_H$ are nonempty and compact.
Any finite collection of compatibility equations can be solved by
choosing a preimage at a common refinement and taking its norms.
The product of these compact fibres therefore has the finite
intersection property for all compatibility conditions. A compatible
lift exists. If a compatible family $a_H$ maps to zero, its unique
preimages $b_H$ under $F-1$ are compatible, since the norms commute
with $F$. This proves the kernel equality; left injectivity holds
at each level. Together with (NM.1), this proves the asserted
exact sequence.

Finally, for a transition of degree $d$ prime to $p$,
$\operatorname{Nm}\circ\pi^*=[d]$. Both maps are defined over
$k$. On the finite rational $p$-primary groups, multiplication by
$d$ is an automorphism. Thus norm has the section
$d^{-1}\pi^*$ and is surjective. An inverse system of nonempty
finite fibres with these compatible maps has surjective projection
to each coordinate, again by compactness. This proves the final
surjectivity assertion.
\end{proof}
\paragraph{Check record.} Also checked in this writing, 2026-11-18: GPT-6 Astra (OpenAI, ChatGPT Pro mode), root; separate effort setting not exposed. Read inverse norm/Tate-module identification, Frobenius cokernel, full-section construction, every-normal-quotient compactness criterion and examples. B.12 remains Open; nonzero fixed vectors alone do not give all-level free vectors.


\begin{claim}{B.7}{Proved}{All fixed classes have full arithmetic section realizations}
Every $m\in M_P^F$ is the compatible $p$-adic abelianized
difference of a continuous full étale arithmetic section from
$s_P$. Its projection to the geometrically prime-to-$p$ quotient
is exactly $s_P$. If $m\ne0$, the realizing section is not
geometrically conjugate to a point section. This conclusion uses
point separation proved by the cuspidal criterion and transport
above; it does not assert free deck orbits for $m$.
\end{claim}
\begin{proof}
Choose $v\in V$ with $\varepsilon(v)=m$ and choose $n\in N$
mapping to $v$. In the full arithmetic group write
$\mathfrak f=s_P(\operatorname{Frob}_k)$ and consider
$n\mathfrak f$. Every element $a$ of a profinite group defines a
continuous homomorphism $\widehat\Z\to\overline{\langle a\rangle}$:
the ordinary integer power map extends, because its composition
with each finite quotient has a finite period. Apply this to
$a=n\mathfrak f$. Its projection to $G_k\simeq\widehat\Z$ sends
the dense generator $1$ to $1$, so the composite is the identity.
It therefore gives a continuous section $s_n$. Since $n\in N$,
its prime-to-$p$ projection is exactly $s_P$.

At every stable pointed level, $n\in U_H$ and its image in
$U_H^{\mathrm{ab},p}=T_H$ is $v_H$. Hence the abelianized
section-difference cocycle takes the value $v_H$ at Frobenius.
Its class is $v_H\bmod(F-1)T_H$, and (NM.2) identifies this with
$m_H$. These identities commute with norm and prove realization
in the whole inverse limit.

Every section whose prime-to-$p$ projection is exactly $s_P$ has
this form: its value at Frobenius is uniquely $n\mathfrak f$ for
some $n\in N$, and that value determines the section. Conjugation
by $a\in N$ replaces $n$ by $a n F(a)^{-1}$, hence changes its
image in $V$ by $(1-F)\bar a$ and leaves $m$ unchanged. We obtain
a surjection from these $N$-conjugacy classes of full sections to
$V/(F-1)V=M_P^F$. No injectivity of this abelianization map is
being asserted.

Suppose $m\ne0$ and $s_n$ were geometrically conjugate to the
section of $Q\in C(k)$. Projection and point separation give
$Q=P$. If $a\in\Omega$ conjugates $s_n$ to $s_P$, its image in
$\Delta$ centralizes the common section $s_P(G_k)$. The geometric
centralizer is trivial by \statementref{A.17}. Thus $a\in N$,
and the preceding coboundary calculation would give $m=0$, a
contradiction. The dependencies on the point-separation proof are
the point-separation theorem A.1, using its direct support-field route. The point-to-section Kummer comparison is B.68.
\end{proof}
\paragraph{Check record.} Also checked in this writing, 2026-11-18: GPT-6 Astra (OpenAI, ChatGPT Pro mode), root; separate effort setting not exposed. Read inverse norm/Tate-module identification, Frobenius cokernel, full-section construction, every-normal-quotient compactness criterion and examples. B.12 remains Open; nonzero fixed vectors alone do not give all-level free vectors.


\begin{claim}{B.8}{Proved}{An exact finite criterion for free orbits at every normal level}
For an $F$-stable open normal $H\triangleleft\Delta$, let $S_H$
be the finite set of $a\in J_{D_H}(k)[p^\infty]$ such that, for
every normal $K\triangleleft\Delta$ containing $H$, the norm of
$a$ to $D_K$ has trivial stabilizer in $\Delta/K$. The quotient
$K/H$ and the curve $D_K$ need not be $F$-stable. Then the set of
fixed norm-module vectors with free orbits at every normal level
is naturally
\[
                         \varprojlim_H S_H.
\]
It is nonempty if and only if every $S_H$ is nonempty. Any such
vector is nonzero. Consequently, nonexistence of an all-level free
vector has a finite certificate of this specific form: one stable
pointed normal cover for which $S_H$ is empty.
\end{claim}
\begin{proof}
There are only finitely many intermediate normal subgroups between
$H$ and $\Delta$, so $S_H$ is a finite, explicitly specified set.
If $H'\subset H$ are stable, norm maps $S_{H'}$ into $S_H$:
its composite with each norm to $D_K$ is the norm from $D_{H'}$
to that same quotient. This proves preservation of all defining
stabilizer conditions, whether or not $K$ is stable.

An all-level free fixed vector projects to every $S_H$. Conversely,
a compatible family in these sets gives a vector of $M_P^F$ by
(NM.3). For an arbitrary normal level $K$, choose a stable
$H\subset K$; the definition of $S_H$ then proves freeness of
that coordinate. Thus the displayed inverse limit is exact.

If every $S_H$ is nonempty, impose the norm compatibility equations
on the compact product $\prod_H S_H$. A finite list is satisfied
by choosing an element at a common refinement and projecting to
the finitely many levels involved. Compactness supplies a point
of the inverse limit. This argument does not require the restricted
maps between the sets $S_H$ to be surjective. The converse is
immediate. There exists a nontrivial prime-to-$p$ geometric normal
cover, for the abelian geometric group surjects onto $\Z/r$ for
any $r\ne p$. At such a level the zero class has nontrivial
stabilizer, so an all-level free vector cannot be zero.

The criterion can also be written directly in terms of a section
representative from \statementref{B.7}. Let $v_H\in T_H$ be
its abelianized value at Frobenius and put
$w_H=(F-1)^{-1}v_H\in W_H$. For $K\supset H$ let
$\nu_{HK}:W_H\to T_K\otimes\Q_p$ be the rationalized norm.
For $g\in\Delta/K$ its action fixes the coordinate precisely when
\[
                         (g-1)\nu_{HK}(w_H)\in T_K.           \tag{NM.4}
\]
Indeed that coordinate is $\nu_{HK}(w_H)\bmod T_K$.
Formula (NM.4), tested for every $g\ne1$ and every intermediate
normal quotient, is equivalent to membership in $S_H$. It
respects the semilinear action: we have not commuted $g$ with
$F-1$, which would generally be incorrect.
\end{proof}
\paragraph{Check record.} Also checked in this writing, 2026-11-18: GPT-6 Astra (OpenAI, ChatGPT Pro mode), root; separate effort setting not exposed. Read inverse norm/Tate-module identification, Frobenius cokernel, full-section construction, every-normal-quotient compactness criterion and examples. B.12 remains Open; nonzero fixed vectors alone do not give all-level free vectors.


\begin{claim}{B.9}{Proved}{A relative class-number witness obstructs free vectors}
Suppose there are a finite extension $L/k$ and compatible pointed
étale prime-to-$p$ covers $Z\to Y\to C_L$, both geometrically
normal over $C$, such that $Z\to Y$ is nontrivial Galois and
\[
                           p\nmid\frac{\#J_Z(L)}{\#J_Y(L)}.
\]
Then $M_P^F$ contains no vector with free deck orbit at every
normal level. The geometric deck groups need not be constant over
$L$. This is a sufficient condition, not a claimed equivalence
with the existence of one cyclic witness over $C$.
\end{claim}
\begin{proof}
The chosen rational lifts identify these covers with levels in the
pointed tower whose descent is induced by a power of $s_P(F)$.
Thus the coordinate of any $m\in M_P^F$ on $Z$ belongs to
$J_Z(L)[p^\infty]$. Write $d=\deg(Z/Y)$ and
$\pi:Z\to Y$. Since $d$ is invertible on $p$-primary groups,
the identities $\operatorname{Nm}\pi^*=[d]$ give a direct-sum
decomposition of finite abelian groups
\[
 J_Z(L)[p^\infty]
 =\pi^*J_Y(L)[p^\infty]\oplus
       \ker\!\left(\operatorname{Nm}:J_Z(L)[p^\infty]
                                  \to J_Y(L)[p^\infty]\right).
\]
The order of the second summand is the $p$-part of the displayed
class-number ratio, so the hypothesis makes that summand zero.
Every coordinate is consequently a pullback from $Y$. All elements
of the nontrivial geometric group $\Gal(Z_{\bar k}/Y_{\bar k})$
fix that pullback, so it fails the required free-orbit condition
at the normal level $Z$. Norm and pullback are defined over $L$,
which is all that this argument needs; individual deck
transformations need not be defined over $L$.
\end{proof}
\paragraph{Check record.} Also checked in this writing, 2026-11-18: GPT-6 Astra (OpenAI, ChatGPT Pro mode), root; separate effort setting not exposed. Read inverse norm/Tate-module identification, Frobenius cokernel, full-section construction, every-normal-quotient compactness criterion and examples. B.12 remains Open; nonzero fixed vectors alone do not give all-level free vectors.


\begin{claim}{B.10}{Proved}{Nonzero fixed classes on an explicit genus-two curve}
For the smooth proper model over $\F_2$ of
\[
                     C:\quad y^2+y=x+\frac1x+\frac1{x+1},
\]
one has $J_C(\F_2)\simeq(\Z/2)^2$. For every rational base point
$P$, the group $M_P^F$ therefore has nonzero classes and these
have non-point-theoretic full arithmetic section realizations.
This assertion does not provide an all-level free vector.
\end{claim}
\begin{proof}
The three poles at $0,1,\infty$ are simple. The Artin--Schreier
extension is geometrically connected, since a function of the form
$b^2+b$ cannot have an odd pole order. It is étale away from these
poles. At a simple pole, $1/y$ is a local uniformizer upstairs and
the involution $y\mapsto y+1$ changes it by a quantity of valuation
two. Thus the local different exponent is two. Riemann--Hurwitz
gives $2g-2=2(-2)+3\cdot2=2$, so $g=2$. Each pole gives one
rational point, denoted $W_0,W_1,W_\infty$.

There are no other $\F_2$-points. Over
$\F_4=\F_2[t]/(t^2+t+1)$, the only additional possible
$x$-values are $t,t+1$. Their Artin--Schreier right sides have
absolute trace one, so neither gives a point. Hence $N_1=N_2=3$.
The genus-two Frobenius identities, or
\cite[Theorem 11.1 and Corollary 11.4]{MilneJac}, then give
\[
              P_C(T)=1-T^2+4T^4,\qquad \#J_C(\F_2)=P_C(1)=4.
\]
The divisor identities
\[
 \operatorname{div}(x)=2(W_0-W_\infty),\qquad
 \operatorname{div}(x+1)=2(W_1-W_\infty)
\]
give two distinct nonzero classes of order two. Nonzero and
distinct follow because the difference of two distinct points on
a positive-genus curve cannot be principal. Thus these classes
give all of $J_C(\F_2)\simeq(\Z/2)^2$.

The retained exact computation in
\path{code/criterion_norm_module.py} checks the field arithmetic,
both point counts, and the Frobenius polynomial. Now apply
\statementref{B.6} and \statementref{B.7}. All actual point
sections are distinct by the point-separation theorem; the new
sections constructed from a nonzero fixed class are different
from all of them.
\end{proof}
\paragraph{Check record.} Also checked in this writing, 2026-11-18: GPT-6 Astra (OpenAI, ChatGPT Pro mode), root; separate effort setting not exposed. Read inverse norm/Tate-module identification, Frobenius cokernel, full-section construction, every-normal-quotient compactness criterion and examples. B.12 remains Open; nonzero fixed vectors alone do not give all-level free vectors. | Also checked in this writing, 2026-11-18: GPT-6 Astra Pro; finite\_checks agent; effort not separately exposed. Read full proof; checked Artin–Schreier geometry, both point counts, the Jacobian order and the two independent rational two-torsion classes. Reran criterion\_norm\_module. Full arithmetic section consequences depend on B.6–7 and A1, checked by parent, not independently by this finite reader.


\begin{claim}{B.11}{False as stated}{Freeness at one normal level need not survive norm}
The implication from a free vector in a finite coprime deck module
to a free norm at every normal quotient is false, even for an
abelian odd group acting on an $\F_2$-vector space.
\end{claim}
\begin{proof}
Take $G=(\Z/3)^2$, $\F_4=\F_2[\omega]$ with
$\omega^2+\omega+1=0$, and $V=\F_4\oplus\F_4$. Let
\[
          (a,b)(u,v)=(\omega^a u,\omega^b v).
\]
The vector $(1,1)$ has trivial stabilizer. For the diagonal
subgroup $H=\langle(1,1)\rangle$, its norm is zero, since
$1+\omega+\omega^2=0$ in both summands. The quotient $G/H$ is
nontrivial and fixes this zero norm. This is an exact finite-module
counterexample to that inference; it is not claimed to be the
Jacobian module of a particular curve. The program just cited
enumerates the nine group elements, the orbit, and the norm.
It explains why the finite sets in \statementref{B.8} must test
every normal quotient of each chosen cover.
\end{proof}
\paragraph{Check record.} Also checked in this writing, 2026-11-18: GPT-6 Astra (OpenAI, ChatGPT Pro mode), root; separate effort setting not exposed. Read inverse norm/Tate-module identification, Frobenius cokernel, full-section construction, every-normal-quotient compactness criterion and examples. B.12 remains Open; nonzero fixed vectors alone do not give all-level free vectors. | Also checked in this writing, 2026-11-18: GPT-6 Astra Pro; finite\_checks agent; effort not separately exposed. Read and checked the full explicit F4² counterexample, all stabilizers and the zero diagonal norm; reran criterion\_norm\_module.



\noindent\emph{Statement B.12 and its produced argument are retained in Part II.}
\subsection{An exact class-number criterion at every finite abelian norm level}

\begin{claim}{B.13}{Proved}{Finite abelian norm obstructions, including nonconstant deck groups}
Let $D\to C$ be a pointed geometrically connected finite
\emph{\'etale} cover over $k=\mathbf F_q$, with abelian geometric
Galois group $A$ of order prime to $p=\operatorname{char}k$.
The chosen point of $D$ lies rationally above $P$.
Let $F$ act on $A$ by the pointed arithmetic descent, and put
\[
 S_D(k)=\{a\in J_D(k)[p^\infty]:
    \operatorname{Nm}_{D/D_H}(a)
    \text{ has trivial stabilizer in }A/H
                   \text{ for every }H\le A\},                 \tag{B.AB.1}
\]
where $D_H=D_{\overline k}/H$, with its pointed descent whenever
the subgroup is stable under a power of $F$.

For each subgroup $H$ for which $A/H$ is a nontrivial cyclic
prime-power group $\mathbf Z/\ell^a$, let $m_H$ be its orbit
length under $F$, and let $H^+/H$ be the unique subgroup of
order $\ell$ in $A/H$. Both quotient curves have their pointed
models over $k_{m_H}=\mathbf F_{q^{m_H}}$. Then
\[
 S_D(k)\ne\varnothing
 \quad\Longleftrightarrow\quad
 p\mid
 \frac{\#J_{D_H}(k_{m_H})}{\#J_{D_{H^+}}(k_{m_H})}
 \quad\text{for every such }H.                                \tag{B.AB.2}
\]
Thus an empty set at a finite abelian normal level is equivalent
to a relative class-number unit in one cyclic prime-power tower,
at that quotient's first pointed field of definition. The step
$D_H\to D_{H^+}$ has prime degree $\ell$; its lower curve need
not be $C$.
\end{claim}
\begin{proof}
We use the geometric $p$-primary torsion group
\[
                   T=J_D(\overline k)[p^\infty].
\]
For a subgroup $H\le A$, the element
\[
                   e_H=\frac1{|H|}\sum_{h\in H}h
                         \in\mathbf Z_p[A]
\]
is an idempotent. Pullback from $D_H$ is injective on
$p$-primary torsion, and pullback followed by norm is
$\sum_{h\in H}h=|H|e_H$. Its image is $e_HT$; this follows
also by applying norm and dividing by $|H|$. Consequently
the stabilizer in $A$ of $e_Ha$ is the inverse image of the
stabilizer of the norm coordinate in $A/H$. Freeness at this
quotient means exactly
\[
                         \operatorname{Stab}_A(e_Ha)=H.         \tag{B.AB.3}
\]

Here is the needed character decomposition, including its
integrality. Over a finite unramified extension of $\mathbf Z_p$
containing all roots of unity of exponent $\exp(A)$, the
idempotents of the characters $\chi$ diagonalize $A$.
For each subgroup $K$ with $A/K$ cyclic, sum these idempotents
over all characters whose kernel is exactly $K$, and denote
the result by $\varepsilon_K$. That sum is invariant under
the unramified Galois group and hence belongs to
$\mathbf Z_p[A]$. The $\varepsilon_K$ are pairwise orthogonal
and sum to one. Moreover
\[
       e_H=\sum_{K\supset H}\varepsilon_K.                     \tag{B.AB.4}
\]
Every nonzero element of $\varepsilon_KT$ has stabilizer
exactly $K$. To see this after the faithful unramified scalar
extension, if $g\notin K$, every corresponding value
$\chi(g)-1$ is a unit: a nontrivial root of unity of order
prime to $p$ remains nontrivial in the residue field. Thus
$g-1$ is invertible on that entire summand. This reasoning
holds for arbitrary finite $p$-power torsion, not only for
torsion killed by $p$.

If $A/H$ is cyclic of order $\ell^a$, its primitive packet is
particularly simple:
\[
                         \varepsilon_H=e_H-e_{H^+}.            \tag{B.AB.5}
\]
Every strict overgroup of $H$ contains $H^+$. Hence (B.AB.3)
forces $\varepsilon_Ha\ne0$: otherwise the whole norm
coordinate would be fixed by $H^+/H$. Conversely, suppose
$a$ has nonzero primitive component at \emph{every} subgroup
with cyclic prime-power quotient. For an arbitrary $K\le A$
and $g\notin K$, a prime-primary character of the finite
abelian group $A/K$ detects $g$. Its kernel lifts to a subgroup
$H\supset K$ with $A/H$ cyclic of prime-power order and
$g\notin H$. The nonzero component $\varepsilon_Ha$ survives
in $e_Ka$, and $g$ cannot fix it. Thus (B.AB.3) holds for
every $K$. We have proved the exact equivalence
\[
 a\text{ satisfies all norm-freeness conditions}
 \quad\Longleftrightarrow\quad
 \varepsilon_Ha\ne0\text{ for every cyclic prime-power }A/H.
                                                                  \tag{B.AB.6}
\]

We next impose Frobenius, without assuming that the deck group
is constant. Conjugation carries $\varepsilon_H$ to
$\varepsilon_{F(H)}$. If $a\in T^F=J_D(k)[p^\infty]$,
then $\varepsilon_Ha$ is fixed by $F^{m_H}$. Necessity in
(B.AB.6) therefore requires
\[
                         (\varepsilon_HT)^{F^{m_H}}\ne0.        \tag{B.AB.7}
\]
Conversely assume (B.AB.7) for every indicated $H$.
Choose one representative in each Frobenius orbit of these
subgroups and a nonzero $b_H$ in the displayed fixed group.
Then
\[
                  a_H=\sum_{i=0}^{m_H-1}F^i b_H
\]
is fixed by $F$. Its components in the distinct packets
$\varepsilon_{F^i(H)}T$ are all nonzero. Summing $a_H$ over
the finitely many subgroup orbits gives a rational torsion
class satisfying (B.AB.6). No division by $m_H$ occurs;
the proof is valid even when $p\mid m_H$.

Finally both $H$ and $H^+$ are stable under $F^{m_H}$, and
their quotient maps and rational lifts of $P$ descend to
$k_{m_H}$. Put
\[
 B_H=\bigl(\ker(J_{D_H}\longrightarrow J_{D_{H^+}})\bigr)^0.
\]
Pullback identifies $B_H[p^\infty]$ with the primitive packet
$\varepsilon_HT$, compatibly with $F^{m_H}$. Norm and
pullback split the $p$-primary groups because their degree
$\ell$ is prime to $p$. The isogeny
$J_{D_{H^+}}\times B_H\to J_{D_H}$ shows
\[
 \#B_H(k_{m_H})
       =\frac{\#J_{D_H}(k_{m_H})}{\#J_{D_{H^+}}(k_{m_H})}.
\]
Thus (B.AB.7) is exactly divisibility by $p$ of the indicated
integer. This proves (B.AB.2).
\end{proof}
\paragraph{Check record.} Also checked in this writing, 2026-11-18: GPT-6 Astra (OpenAI, ChatGPT Pro mode), root; separate effort setting not exposed. Read abelian primitive-packet decomposition including nonconstant Frobenius and p-divisible orbit lengths; solvable induction with all normal quotient idempotents; nonabelian finite-module limitation. Numerical checks are separate events.


\begin{claim}{B.14}{Proved}{The elementary-prime case and the stronger higher-power test}
If the group $A$ in \statementref{B.13} is constant, then all
$m_H=1$. If in addition $A=(\mathbf Z/\ell)^d$, then
$S_D(k)$ is nonempty exactly when every degree-$\ell$ quotient
$D_H\to C$ has relative class number divisible by $p$.
Thus at a constant elementary-abelian level, an obstruction
is exactly a prime-degree class-number witness among those
quotients.

For higher-exponent abelian groups, the exact criterion also
tests the primitive parts at cyclic prime-power quotients.
These tests are strictly stronger as finite module conditions
than testing only prime-degree quotients. This strictness does
not assert that a particular curve has no prime-degree witness
over all constant extensions.
\end{claim}
\begin{proof}
For elementary $A$, every nontrivial cyclic quotient has degree
$\ell$, and $H^+=A$. The relative factor in (B.AB.2) is
therefore exactly that of $D_H/C$. This proves the first
assertion and identifies the relationship with B1 at the
displayed finite field and level.

Here is a complete finite-module control for the higher-power
assertion. Take characteristic two and
\[
        A=\mathbf Z/9\times\mathbf Z/3,
                     \qquad M=\mathbf F_4^4.
\]
Let $\omega\in\mathbf F_4$ have order three. An element
$(a,b)$ acts diagonally on the four summands with scalars
\[
          \omega^a,\qquad\omega^b,\qquad
          \omega^{a+b},\qquad\omega^{a+2b}.                    \tag{B.AB.8}
\]
These are all four character kernels of prime-degree
quotients. A vector with all four coordinates nonzero has
free norm orbit at every such quotient. There are $3^4=81$
such vectors. Nevertheless the subgroup
\[
                   \{(0,0),(3,0),(6,0)\}
\]
fixes every element of $M$, so no vector has free orbit at
the full level. The three missing primitive packets have
cyclic quotient of order nine.

For comparison, choose a primitive ninth root
$\zeta\in\mathbf F_{64}$ and add three summands
$\mathbf F_{64}$ with characters
\[
                \zeta^{a+3jb},\qquad j=0,1,2.
\]
Now all seven cyclic prime-power kernels are present.
The vector with every coordinate one has free norm orbit
at all ten subgroups of $A$. In fact the exact number of
all-quotient free vectors in this module is
\[
                       3^4\cdot63^3=20\,253\,807.
\]
The packet criterion proves the formula; the accompanying
program checks every subgroup and every norm for the displayed
vector. It exhausts the $256$ vectors of the first module and
independently verifies its count $81$ and its absence of an
all-quotient free vector.

There is also a check of the nonconstant Frobenius step when
the orbit length is divisible by the characteristic. Set
$A=(\mathbf Z/3)^2$, retain (B.AB.8), and define
\[
 F(a,b)=(b,a),\qquad
             F(u,v,w,z)=(v,u,w,z^2).
\]
These actions satisfy $F(gm)=F(g)F(m)$. The fixed module
has $32$ vectors, of which exactly nine have free norm orbit
at every subgroup. The first two character kernels form an
orbit of length two, and the other two are stable. The sum
\[
                  (1,0,0,0)+F(1,0,0,0)=(1,1,0,0)
\]
is fixed and is nonzero in both moving packets, despite
working in characteristic two. Thus the Frobenius-orbit
argument in \statementref{B.13} does not hide an averaging
division by an inadmissible integer.

All finite-field operations, subgroup enumeration, norms,
stabilizers and counts are independently implemented in
\path{code/criterion_abelian_norm.py}; its complete results
are in the adjacent JSON certificate. These are finite
representation examples, with no claimed Jacobian realization.
The universal genus-two characteristic-two assertion in B2
still requires an empty level for every pointed curve. The
present theorem exactly reduces every \emph{abelian} candidate
level to the displayed relative class-number tests; it does
not settle the existence of such a level, or rule out a
stronger obstruction at a nonabelian level.
\end{proof}
\paragraph{Check record.} Also checked in this writing, 2026-11-18: GPT-6 Astra (OpenAI, ChatGPT Pro mode), root; separate effort setting not exposed. Read abelian primitive-packet decomposition including nonconstant Frobenius and p-divisible orbit lengths; solvable induction with all normal quotient idempotents; nonabelian finite-module limitation. Numerical checks are separate events. | Also checked in this writing, 2026-11-18: GPT-6 Astra Pro; finite\_checks agent; effort not separately exposed. Read full proof and all three finite module examples. Rechecked character-kernel packets, higher-exponent strictness, Frobenius compatibility and no averaging over its characteristic-divisible order. Both supplied abelian norm programs rerun.



\subsection{Solvable deck groups: a single common-stabilizer obstruction}

\begin{claim}{B.15}{Proved}{Simultaneous norm freeness for solvable coprime actions}
Let $G$ be a finite solvable group of order prime to $p$, acting
on a finite abelian $p$-group $M$. Let $F$ be an automorphism
of $M$, and let $\sigma\in\operatorname{Aut}(G)$ satisfy
$F(gm)=\sigma(g)F(m)$. For $K\triangleleft G$ put
\[
                       e_K=|K|^{-1}\sum_{g\in K}g.
\]
There exists $a\in M^F$ with
\[
                 \operatorname{Stab}_G(e_Ka)=K
                      \quad(K\triangleleft G)                 \tag{B.SV.1}
\]
if and only if
\[
        \bigcap_{b\in M^F}\operatorname{Stab}_G(e_Kb)=K
                      \quad(K\triangleleft G).                \tag{B.SV.2}
\]
Neither $M^F$ nor an individual norm image is assumed to be
$G$-stable. The order of $F$ or of $\sigma$ may be divisible
by $p$.
\end{claim}
\begin{proof}
The implication from (B.SV.1) to (B.SV.2) is immediate.
We prove the converse by induction on the derived length
of $G$. If $G=1$, choose $a=0$. Otherwise let $N$ be the
last nontrivial derived subgroup of $G$. It is nontrivial,
abelian, normal, and characteristic, so it is $\sigma$-stable.
There is a $G$- and $F$-stable direct sum
\[
                    M=e_NM\oplus W,\qquad W=(1-e_N)M.          \tag{B.SV.3}
\]
The first summand is a module for $G/N$. Its fixed group is
$e_N(M^F)$, since the idempotent commutes with $F$.
For every normal subgroup $K\supset N$, its common
stabilizer after $e_K$ is $K/N$ by (B.SV.2). The induction
hypothesis supplies $v\in(e_NM)^F$ such that
\[
                \operatorname{Stab}_G(e_Kv)=K
                        \quad(K\triangleleft G,\ N\subset K).
                                                                  \tag{B.SV.4}
\]

We need one simultaneous support lemma for the abelian group
$N$. For each subgroup $L\le N$ with $N/L$ cyclic, let
$\varepsilon_L^N\in\mathbf Z_p[N]$ be the sum of the
character idempotents with exact kernel $L$, including the
trivial character when $L=N$. As in \statementref{B.13},
these are integral orthogonal idempotents summing to one;
every nonzero element in their $L$-summand has stabilizer
exactly $L$ in $N$. The reason is that $\chi(n)-1$ is a
unit for every nontrivial prime-to-$p$ root $\chi(n)$.

All the $e_K$, for $K\triangleleft G$, are central in
$\mathbf Z_p[G]$, hence commute with each other and with
these $N$-idempotents. Form their simultaneous orthogonal
idempotents by taking all products
\[
 \varepsilon_L^N
       \prod_{K\in\mathcal S}e_K
       \prod_{K\notin\mathcal S}(1-e_K),                       \tag{B.SV.5}
\]
where $\mathcal S$ ranges through subsets of the finite set
of normal subgroups of $G$, and discard zero summands.
They decompose $W$. On each resulting summand, every $e_K$
is either zero or the identity, and every nonzero vector has
one fixed stabilizer $L$ in $N$.

The automorphism $F$ permutes these summands: it carries
$e_K$ to $e_{\sigma(K)}$ and $\varepsilon_L^N$ to
$\varepsilon_{\sigma(L)}^N$. Decompose $W$ further into
the direct sums belonging to these $F$-orbits. In each such
orbit whose fixed subgroup is nonzero choose one nonzero
fixed vector. Its projection to every summand in that orbit
is nonzero, since $F$ cyclically identifies all projections.
Sum the chosen vectors over the finitely many available
orbits. The resulting $w\in W^F$ has a nonzero component
in exactly every joint summand that occurs in any fixed
vector. Consequently, simultaneously for every normal $K$,
\[
 \operatorname{Stab}_N(e_Kw)
       =\bigcap_{b\in W^F}\operatorname{Stab}_N(e_Kb).           \tag{B.SV.6}
\]
Indeed $e_K$ retains a specified collection of the joint
summands; on each retained nonzero component the stabilizer
is its exact $N$-character kernel. Thus both sides are the
intersection of the same kernels. This argument uses no
averaging over $F$ and works for arbitrary $p$-power torsion.

The summand $e_NM$ is fixed by $N$, so (B.SV.2) and
(B.SV.3) identify the right side of (B.SV.6) with $N\cap K$.
Set $a=v+w$. If $g\in G$ fixes $e_Ka$, the $G$-stable
decomposition (B.SV.3) shows that it fixes $e_Kv$ and
$e_Kw$ separately. Since
\[
                         e_Kv=e_{KN}v,
\]
equation (B.SV.4) implies $g\in KN$. Write $g=nk$, with
$n\in N$ and $k\in K$. The group $K$ fixes every $e_K$
image. Hence $g$ fixes $e_Kw$ exactly when $n$ does.
Equation (B.SV.6) gives $n\in N\cap K$, so $g\in K$.
The reverse inclusion is automatic. This proves (B.SV.1)
for every normal $K$ and completes the induction.
\end{proof}
\paragraph{Check record.} Also checked in this writing, 2026-11-18: GPT-6 Astra (OpenAI, ChatGPT Pro mode), root; separate effort setting not exposed. Read abelian primitive-packet decomposition including nonconstant Frobenius and p-divisible orbit lengths; solvable induction with all normal quotient idempotents; nonabelian finite-module limitation. Numerical checks are separate events. | Also checked in this writing, 2026-11-18: GPT-6 Astra Pro; finite\_checks agent; effort not separately exposed. Read full induction line by line. Checked the characteristic last derived subgroup, integral exact-kernel idempotents, simultaneous norm idempotents, Frobenius support lemma without averaging, and final KN stabilizer argument. Finite solvable norm program rerun.


\begin{claim}{B.16}{Proved}{Constant solvable covers have exactly the relative-unit obstruction}
Let $D\to C$ be a constant pointed normal prime-to-$p$
cover over a finite field $k$, with solvable deck group $G$.
The finite all-normal-quotient norm-free set $S_D(k)$ is empty
if and only if there exist normal subgroups
$K\subsetneq L\triangleleft G$ such that
\[
                         p\nmid
                   \frac{\#J_{D_K}(k)}{\#J_{D_L}(k)}.          \tag{B.SV.7}
\]
In particular, for constant solvable deck groups, the relative
class-number-unit sufficient criterion is also necessary at
that finite level. The lower curve $D_L$ can differ from $C$,
and the degree of the step is not required to be prime.
\end{claim}
\begin{proof}
Apply \statementref{B.15} with
$M=J_D(k)[p^\infty]$ and $F=1$; the actual arithmetic
Frobenius has already been imposed by taking rational points.
This is a $G$-module because the deck group is constant.
The common stabilizer $L$ of $e_KM$ is a normal subgroup of
$G$, since $e_KM$ is $G$-stable, and contains $K$.

Pullback identifies $J_{D_K}(k)[p^\infty]$ with $e_KM$:
the averaging idempotent is defined over $k$, and pullback
and norm split because $|K|$ is prime to $p$.
If $L\supsetneq K$ is its common stabilizer, all of this
group is invariant under $L/K$. Pullback from $D_L$ is
therefore all of $J_{D_K}(k)[p^\infty]$, again by averaging.
The $p$-parts of the two Jacobian orders agree, giving
(B.SV.7). Conversely that relative unit makes every rational
$p$-primary class on $D_K$ a pullback from $D_L$, so its
stabilizer contains the nontrivial group $L/K$. No class can
satisfy all normal quotient conditions. The equivalence in
\statementref{B.15} proves that these are the only ways the
finite set can be empty.
\end{proof}
\paragraph{Check record.} Also checked in this writing, 2026-11-18: GPT-6 Astra (OpenAI, ChatGPT Pro mode), root; separate effort setting not exposed. Read abelian primitive-packet decomposition including nonconstant Frobenius and p-divisible orbit lengths; solvable induction with all normal quotient idempotents; nonabelian finite-module limitation. Numerical checks are separate events. | Also checked in this writing, 2026-11-18: GPT-6 Astra Pro; finite\_checks agent; effort not separately exposed. Read full proof; checked rational p-primary pullback/invariant identification, normality of the common stabilizer for constant deck group, and both directions of the relative-unit equivalence.


\begin{claim}{B.17}{Proved}{A single-element finite certificate in characteristic two}
For any proper hyperbolic pointed curve $(C,P)/k$ in
characteristic two, $M_P^F$ has no all-level free vector if
and only if there are a pointed $F$-stable normal cover
$D\to C$, a normal subgroup $K$ of its geometric deck group
$G$, and an element $g\in G\setminus K$ such that
\[
 g\,\operatorname{Nm}_{D/D_K}(a)
       =\operatorname{Nm}_{D/D_K}(a)
        \qquad\text{for every }a\in J_D(k)[2^\infty].          \tag{B.SV.8}
\]
The quotient $D_K$ need only be defined over its pointed
finite constant extension. Thus in characteristic two a
finite obstruction can always be witnessed by one quotient
and one element fixing every candidate norm coordinate.
It cannot require only a union of different stabilizers
covering the candidates.

This reformulation does not prove that such a finite
certificate exists for every genus-two curve. That universal
existence assertion in B2 remains open here.
\end{claim}
\begin{proof}
The geometric deck group of a prime-to-two cover has odd
order. The exact theorem of Feit--Thompson is that every
finite group of odd order is solvable; these are its only
hypotheses \citep[Chapter~I, \S1, p.~775, unnumbered main
theorem]{FeitThompson1963}. Hence \statementref{B.15} applies
to every finite level under consideration.

Choose $2^e$ killing the finite group $J_D(k)[2^\infty]$
and apply that theorem to
\[
                      M=J_D[2^e](\overline k).
\]
The geometric deck group acts on $M$; arithmetic Frobenius
normalizes that action and has fixed group exactly
$J_D(k)[2^\infty]$. Pullback identifies the stabilizer of
each norm coordinate with the stabilizer of its averaging
idempotent, so \statementref{B.15} says that $S_D(k)$
is empty precisely when (B.SV.8) holds for some $K,g$.
By the compact finite criterion \statementref{B.8}, absence
of an all-level free vector is equivalent to emptiness at
one stable finite level. This proves the stated equivalence.

The distinction between this statement and
\statementref{B.16} matters. With nonconstant deck group,
the Frobenius-fixed group need not be $G$-stable, so the
common stabilizer in (B.SV.8) need not be normal and cannot
automatically be interpreted as a relative class-number unit
over $k$. The semilinear common-stabilizer criterion is the
assertion proved in this generality.
\end{proof}
\paragraph{Check record.} Also checked in this writing, 2026-11-18: GPT-6 Astra (OpenAI, ChatGPT Pro mode), root; separate effort setting not exposed. Read abelian primitive-packet decomposition including nonconstant Frobenius and p-divisible orbit lengths; solvable induction with all normal quotient idempotents; nonabelian finite-module limitation. Numerical checks are separate events. | Also checked in this writing, 2026-11-18: GPT-6 Astra Pro; finite\_checks agent; effort not separately exposed. Read full proof; checked finite geometric torsion reduction with rational fixed group, the nonconstant quotient-field interpretation and use of B.15. Feit–Thompson and compact criterion B.8 not independently resourced by this reader.


\paragraph{Exhaustive finite checks of the induction.}
The program \path{code/criterion_solvable_norm.py} verifies the
theorem on a nonabelian solvable group with a Frobenius of order
divisible by the characteristic. Over $\mathbf F_3$, put
\[
 i=\begin{pmatrix}0&1\\2&0\end{pmatrix},\qquad
 j=\begin{pmatrix}1&1\\1&2\end{pmatrix},\qquad
 T=\begin{pmatrix}1&1\\0&1\end{pmatrix}.
\]
The first two generate $Q_8$, since $i^2=j^2=-1$ and
$ij=-ji$. The matrix $T$ has order three and normalizes this
group. Let $V=\mathbf F_3^2$ be the natural module and add
the three nontrivial one-dimensional characters of
$Q_8/\{\pm1\}$, giving $M=V\oplus\mathbf F_3^3$.
Define Frobenius by $T$ on $V$ and by the induced cyclic
permutation on the three character summands. It satisfies
the semilinear compatibility in \statementref{B.15}.

The program enumerates all six normal subgroups, all $243$
vectors, and every norm and stabilizer. With $F=1$, exactly
$64$ vectors satisfy all norm-free conditions. With the
displayed order-three Frobenius, exactly nine vectors are
fixed and four satisfy all conditions; every common-stabilizer
test is exact. On the natural module alone and on the three
character summands alone, the all-level set is empty, and
the program identifies a failed common-stabilizer condition
in each case, both with constant and nonconstant Frobenius.
These finite representation checks are recorded in
\path{code/criterion_solvable_norm.json}; no geometric
realization is asserted by the checks.


\subsection{A genuinely nonabelian finite obstruction mechanism}

\begin{claim}{B.18}{Proved}{A finite norm obstruction with no vanishing normal-step module}
The equivalence of \statementref{B.13} cannot be extended to
arbitrary finite groups using only the vanishing of relative
$p$-primary modules at normal quotient steps. There is a
faithful action of a finite nonabelian simple $7'$-group $G$
on a nonzero finite $7$-group $M$ such that the set of vectors
with free norm orbit at every normal quotient is empty,
although the relative module at every nontrivial normal step
is nonzero. This is an exact finite representation example,
not a claimed realization by a Jacobian or a counterexample
to universal B1.
\end{claim}
\begin{proof}
Let $M=\mathbf F_7^5$, and define $G\le\operatorname{GL}(M)$
by the three matrices
\[
 U=\operatorname{diag}(1,4,1,4,4),\qquad
 P=\begin{pmatrix}
 1&0&0&0&0\\
 0&1&0&0&0\\
 0&0&0&0&1\\
 0&0&1&0&0\\
 0&0&0&1&0
 \end{pmatrix},\qquad
 Q=\begin{pmatrix}
 2&4&4&4&4\\
 2&5&4&5&5\\
 2&4&5&5&5\\
 2&5&5&5&4\\
 2&5&5&4&5
 \end{pmatrix}.                                                \tag{B.NA.1}
\]
Here $U^3=P^3=Q^2=1$, and $UP\ne PU$.
The following finite certificate proves both the order and
simplicity needed in the argument, without assuming an
identification with a named finite classical group.

Starting with the identity, close the set of matrices under
right multiplication by $U,P,Q$. Exact arithmetic modulo
seven gives
\[
                        |G|=25920=2^6\cdot3^4\cdot5.
\]
Conjugation by the same generators partitions this complete
matrix list into twenty classes, with sizes
\[
\begin{gathered}
 1,40,40,45,240,270,360,360,480,540,720,720,1440,\\
                   2160,2160,2160,2880,2880,3240,5184.
\end{gathered}                                                   \tag{B.NA.2}
\]
For each nonidentity class representative $a$, impose the
equivalence relation $1\sim a$ on the matrix list and close
it under left and right multiplication by $U,P,Q$. This is
the group congruence whose identity class is the normal
closure of $a$. In all nineteen cases its identity class
has all $25920$ elements. Every nontrivial normal subgroup
contains one of these representatives up to conjugacy, so
$G$ is simple.

The data are fully reproducible. The generating program saves
every matrix in
\path{code/criterion_nonabelian_norm_group.json}. The independent
checker \path{code/criterion_nonabelian_norm_verify.py} reads
the three generators and regenerates the complete group using
a separate general matrix multiplication routine. It checks
the conjugacy partition and the normal closures just described.
Congruence closure uses a finite union-find calculation: after
a new pair is identified, its six images under left and right
generator multiplication are queued. This computes the least
invariant equivalence relation and therefore the asserted
normal closure. The checker also verifies that the common
fixed space of the three generators has dimension zero and
that $\sum_{g\in G}g$ is the zero operator on $M$.

The action is faithful by its definition as a matrix group,
and it is coprime because $7\nmid25920$. Nevertheless
\[
                         |M|=7^5=16807<|G|.
\]
No orbit in $M$ can therefore have size $|G|$. In particular
the norm-freeness condition already fails at the identity
subgroup, where the norm coordinate is the vector itself.

On the other hand, simplicity leaves only the normal
subgroups $1$ and $G$. The old module at the latter level is
$M^G=0$, and the relative module for the sole nontrivial
normal step is the entire nonzero module $M$. Its order is
$7^5$, not a unit. Thus there is no vanishing relative module
at any nontrivial normal step, yet the full finite norm-free
set is empty. The obstruction is the union of nontrivial
stabilizers, rather than one absent abelian character packet.

For an actual cover in B2, the separate task is to determine
whether its Frobenius-fixed Jacobian module has such a
configuration. This example proves that a search for
nonabelian finite obstructions has a mechanism not captured
by the relative class-number-unit sufficient criterion. It
does not supply that geometric realization or settle the
universal assertion.
\end{proof}
\paragraph{Check record.} Also checked in this writing, 2026-11-18: GPT-6 Astra (OpenAI, ChatGPT Pro mode), root; separate effort setting not exposed. Read abelian primitive-packet decomposition including nonconstant Frobenius and p-divisible orbit lengths; solvable induction with all normal quotient idempotents; nonabelian finite-module limitation. Numerical checks are separate events. | Also checked in this writing, 2026-11-18: GPT-6 Astra Pro; finite\_checks agent; effort not separately exposed. Read full proof; checked finite group simplicity certificate logic, conjugacy-class normal closures, pigeonhole free-orbit obstruction and unique nontrivial normal-step module. Supplied independent complete group verifier rerun.



\subsection{Point-isolating covers and exact constructions}

\begin{claim}{B.19}{Proved}{A pointed cover isolating one rational point}
For every $C/k$ as in \statementref{A.1} and every $P\in C(k)$,
there is a connected finite \emph{\'etale} cover $D\to C$ with
geometrically Galois group of order prime to $p$, and a specified
$P'\in D(k)$ over $P$, such that the image of $D(k)$ in $C(k)$ is
exactly $\{P\}$. The geometric deck group need not be constant over $k$.
\end{claim}
\begin{proof}
Split $\Pi_C=\Delta_C\rtimes s_P(G_k)$. For $Q\ne P$, the equations
expressing geometric conjugacy of $s_Q$ and $s_P$ have no solution
in the compact group $\Delta_C$, by \statementref{A.1}. If they had
a solution modulo every open characteristic subgroup $N$, the
closed solution sets in $\Delta_C$ would have the finite-intersection
property and yield a solution in $\Delta_C$. Thus some such $N_Q$
detects the failure. We can take characteristic subgroups cofinally:
a finitely generated profinite group has only finitely many open
subgroups of any prescribed index, and their finite intersection
is characteristic.

Intersect the finitely many $N_Q$, for $Q\in C(k)-\{P\}$, to obtain
$N$. The subgroup $N\rtimes s_P(G_k)$ defines the desired pointed
cover. A rational lift of $Q$ would give a fixed coset in its finite
fibre, equivalently a conjugacy of the two section subgroups modulo
$N$, contradicting the choice of $N_Q$. Its geometric deck group
is $\Delta_C/N$, of prime-to-$p$ order. This proof uses precisely
the relative exactness, boundary transport, and cuspidal criterion
audited in Part A; it gives no class-number unit assertion.
\end{proof}
\paragraph{Check record.} Also checked in this writing, 2026-11-18: GPT-6 Astra (OpenAI, ChatGPT Pro mode), root; separate effort setting not exposed. Read isolating-cover arguments, both F4 examples, scalar-Frobenius Lang cover, conductor reduction and explicit primitive-zero obstruction. Degree-one-allowed B.24 remains Open. Certificate reruns are recorded separately.


\begin{claim}{B.20}{Proved}{A constructive two-stage odd Lang criterion}
Let $C/L$ be proper of genus at least two over a finite field of
characteristic two, $P\in C(L)$, and let $D\to C$ be a connected
constant cyclic odd-degree \emph{\'etale} cover with a point
$P_D\in D(L)$ above $P$. Suppose $h(D)/h(C)$ is odd.
Let $E\to D$ be the Lang cover associated with the maximal odd
quotient of $J_D(L)$ and based at $P_D$. Then no point of $C(L)$
other than $P$ lifts rationally to $E$. Taking the geometric
Galois closure and its $P$-pointed twist preserves this property
and gives an odd geometric Galois cover of $C$.
\end{claim}
\begin{proof}
For $G=\operatorname{Gal}(D/C)$ of odd order, the idempotent
$e=|G|^{-1}\sum_{g\in G}g$ acts on the two-primary rational torsion
of $J_D$. The invariant part is the image of the pullback from
$J_C$: pullback and norm are inverse after division by $|G|$.
The complementary abelian factor has rational point count
$h(D)/h(C)$, since the rational Frobenius polynomial factors as
the product of the base and complementary polynomials. Therefore
the oddness assumption says that all rational two-primary points
of $J_D$ are fixed by $G$.

If $Q_D\ne P_D$ lies above $Q\ne P$, its difference class cannot be
fixed by a nonidentity odd deck transformation $a$. Otherwise
$aQ_D+P_D\sim Q_D+aP_D$. These are different effective degree-two
divisors, so the curve is hyperelliptic and uniqueness of its
degree-two pencil gives $aQ_D=\iota P_D$ and $aP_D=\iota Q_D$.
Every automorphism commutes with the hyperelliptic involution
$\iota$, whence $a^2$ fixes $Q_D$. A free deck action of odd order
makes this impossible. Consequently $[Q_D-P_D]$ is not
two-primary. In a Lang cover the rational lift criterion is
exactly vanishing of this class in the chosen quotient of
$J_D(L)$: the Lang endomorphism is $F-1$, with kernel $J_D(L)$.
Thus $Q_D$ does not lift to $E$.

The two stages have constant odd Galois groups. The geometric
Galois closure of their composite embeds in
$\operatorname{Gal}(E/D)^{|G|}\rtimes G$, hence still has odd order.
The geometric core of the pointed subgroup is stable under the
action of $s_P(G_L)$, so its pointed twist has an $L$-point over
$P$ and refines the original cover. A refinement cannot acquire
a rational lift absent from the original cover.
\end{proof}
\paragraph{Check record.} Also checked in this writing, 2026-11-18: GPT-6 Astra (OpenAI, ChatGPT Pro mode), root; separate effort setting not exposed. Read isolating-cover arguments, both F4 examples, scalar-Frobenius Lang cover, conductor reduction and explicit primitive-zero obstruction. Degree-one-allowed B.24 remains Open. Certificate reruns are recorded separately.


\begin{claim}{B.21}{Proved}{The two requested genus-two classes over $\mathbf F_4$}
Write $\alpha^2+\alpha+1=0$. The two conjugate curves
\[
 C_\alpha:\quad Y^2+(x^2+x+\alpha)Y
   =x^5+\alpha x^3+x+1
\]
are the curves denoted by the pair of class labels $203/302$ in
the supplied return. Each has seven $\mathbf F_4$-points.
For every one of these fourteen pointed cases, the file
\path{code/isolating_lang_covers.json} specifies a connected
pointed cover over $\mathbf F_4$, with odd geometric Galois group,
whose rational-point image is just the chosen point.
The representation is an exact Lang-cover and geometric-core
construction; it does not expand enormous affine equations.
\end{claim}
\begin{proof}
The program \path{code/isolating_lang_covers.py} works in
$L=\mathbf F_{64}=\mathbf F_2[t]/(t^6+t+1)$, with elements encoded
by polynomial bits; the two roots $\alpha$ have codes $58,59$.
It implements polynomial Euclidean division and the full Mumford
addition and reduction formulas for this equation, checking the
defining divisibility after every addition. Direct counts are
\[
 \#C(\mathbf F_4)=7,\qquad \#C(\mathbf F_{16})=15,\qquad
 \#C(\mathbf F_{64})=91.
\]
They give the Frobenius polynomial
$X^4+2X^3+X^2+8X+16$ and
$\#J_C(L)=6076=4\cdot7^2\cdot31$. The two independent classes
specified in the JSON generate the full $7$-primary subgroup.
Its 49-element addition table is checked against
$\mathbf F_7^2$ and all $49^3$ associativity triples.
Multiplication by $124$ projects rational point classes onto it.

For every original point $P$, all eight quotient kernels
$\chi:\mathbf F_7^2\to\mathbf F_7$ are tested. If $N_b$ is the
computed occupancy of the coset $b$ for the $P$-based Abel map,
and $c=\chi([P-\infty])$, Riemann--Roch gives exactly
\[
 L_P(\chi,T)=1+\left(\sum_{b=0}^6N_b\zeta_7^b\right)T
                    +64\zeta_7^{-2c}T^2 .
\]
The integer resultant of $L_P(\chi,1)$ and
$\Phi_7$ is the relative class number. The selected relative
values are
\[
 46713473689,\qquad120974743723,\qquad103441382829,
\]
all odd. The full occupancy vectors, character coordinates,
base-point twists, and every attempted value are retained, not
just their residues modulo two.

For $\alpha=58$ the selected characters in the recorded basis are
$(1,0)$ at $\infty,(58,1),(58,58)$ and $(1,1)$ at the other four
original points. For $\alpha=59$ they are $(1,0)$ at
$\infty,(59,1),(59,59)$ and $(1,1)$ at the other four points.
The odd degrees $m$ of the second Lang stage are respectively
\[
 70957766533591,\qquad183760635715237,\qquad157127460517251.
\]
The preceding theorem isolates $P$ over $L$.

To specify the descent, use the splitting $s_P$ of the geometric
fundamental extension over $\mathbf F_4$. Let $N$ be the geometric
core for the two-stage cover over $L$ and let $F$ be conjugation
by $s_P(\operatorname{Frob}_4)$. It satisfies $F^3N=N$. Put
$N_0=N\cap F(N)\cap F^2(N)$ and take the cover corresponding to
$N_0\rtimes s_P(G_{\mathbf F_4})$. It is an exact construction
from the computed Lang data, stable under $F$, and still of odd
geometric Galois order. After extension to $L$ it refines the
isolating cover, so it isolates $P$ already over $\mathbf F_4$.
Its degree is at most $(7m^7)^3$; we do not claim equality.
\end{proof}
\paragraph{Check record.} Also checked in this writing, 2026-11-18: GPT-6 Astra (OpenAI, ChatGPT Pro mode), root; separate effort setting not exposed. Read isolating-cover arguments, both F4 examples, scalar-Frobenius Lang cover, conductor reduction and explicit primitive-zero obstruction. Degree-one-allowed B.24 remains Open. Certificate reruns are recorded separately. | Also checked in this writing, 2026-11-18: GPT-6 Astra Pro; finite\_checks agent; effort not separately exposed. Read the full proof; checked genus-two twisted polynomial, Lang connectedness, two-stage isolation and Frobenius-core descent. Reran supplied exact divisor-class and resultant program; prior general isolation criterion is reproved by the parent reader.


\begin{claim}{B.22}{Proved}{A direct isolating cover for the curve over $\mathbf F_{2^{38}}$}
For the curve and point $P=\infty$ specified in
\statementref{B.23} below, put $m=2^{19}+1=524289$.
The pullback of $[m]:J_C\to J_C$ by
$i_P:C\to J_C$, $Q\mapsto[Q-P]$, is a connected
\emph{\'etale} Galois cover of degree $m^4$ over
$\mathbf F_{2^{38}}$. Its constant deck group is
$(\mathbf Z/m)^4$, and its rational-point image in $C$ is
exactly $\{P\}$.
\end{claim}
\begin{proof}
The explicit elliptic maps in \statementref{B.23} show that
Frobenius on $J_C$ is $[-2^{19}]$. Thus
$J_C(k)=\ker(F-1)=J_C[m](\bar k)$; all $m$-torsion is rational,
and multiplication by $m$ kills every rational point.
The Abel map induces a surjection onto the geometric abelian
fundamental group modulo $m$, so its pullback of the multiplication
cover is connected. It is \emph{\'etale}, since $m$ is odd.
A rational point $Q$ lifts precisely when $[Q-P]\in[m]J_C(k)=0$.
The Abel map is injective on geometric points, so this means $Q=P$.
The point $(P,0)$ is the specified rational lift. The exact integer
degree and the field and curve data are recorded in
\path{code/supersingular_isolating_cover.json}.
\end{proof}
\paragraph{Check record.} Also checked in this writing, 2026-11-18: GPT-6 Astra (OpenAI, ChatGPT Pro mode), root; separate effort setting not exposed. Read isolating-cover arguments, both F4 examples, scalar-Frobenius Lang cover, conductor reduction and explicit primitive-zero obstruction. Degree-one-allowed B.24 remains Open. Certificate reruns are recorded separately. | Also checked in this writing, 2026-11-18: GPT-6 Astra Pro; finite\_checks agent; effort not separately exposed. Read the full proof; checked scalar Frobenius implies J(k)=J[m], connected multiplication pullback, exact rational lifting criterion and injectivity of the Abel map.




\subsection{The nontrivial primitive-cover counterexample}

\begin{claim}{B.23}{False as stated}{Universal nontrivial primitive-zero cyclic covers}
The assertion that every genus-two curve in characteristic two has
a connected cyclic \emph{\'etale} cover of odd degree $n>1$ whose
primitive Jacobian part has two-rank zero is false. A completely
specified counterexample is the smooth proper model over
$k=\mathbf F_2[t]/(M(t))$ of
\[
\begin{split}
 M(Z)&=Z^{38}+Z^{14}+Z^{10}+Z^9+Z^8+Z^5+Z^2+Z+1,\\
 d&=t^3+1,\qquad
 C:\ y^2+y=A(u)^3,\\
 A(u)&=(u^2+t^2u+t)/d,\qquad
 B(u)=(tu^2+u+t^2)/d .
\end{split}
\]
The polynomial $M$ is irreducible, and $C$ is supersingular of
genus two with a unique rational point $P$ above infinity.
For every such cover over $\bar k$, the primitive part has positive
two-rank. This independently confirms the particular counterexample
AS-A.86. It does not refute the version allowing degree one.
\end{claim}
\begin{proof}
We supply the geometry, the finite cutoff, and the exact certificates.

\emph{The curve and its Jacobian.}
Successive squaring modulo $M$ gives
$Z^{2^{38}}\equiv Z$ and
$\gcd(M,Z^{2^2}-Z)=\gcd(M,Z^{2^{19}}-Z)=1$.
These are precisely the irreducibility tests for degree 38; every
residue is retained in \path{code/primitive_certificates/manifest.json}.
In particular $t\notin\mathbf F_4$ and $d\ne0$.
After Artin--Schreier changes removing even powers, $A(u)^3$ has
highest odd term $(t^2/d^3)u^5$, with nonzero coefficient.
The degree-two map to $\mathbf P^1_u$ is ramified only at infinity,
with different exponent six. Riemann--Hurwitz gives genus two,
one rational point $P$ at infinity, and canonical divisor $2P$.
The Artin--Schreier $p$-rank formula, or its elementary one-pole
Cartier calculation, gives two-rank zero. In this case a basis of
regular differentials is $du,u\,du$ after the reduced quintic
equation; Cartier sends $du$ to zero and $u\,du$ to $du$, so its
stable rank is zero.

On $E:y^2+y=x^3$ there are two morphisms
\[
 (x_1,y_1)=(A(u),y),\qquad
 (x_2,y_2)=(B(u),y+(u^3+1)/d).
\]
The polynomial identity
\[
 (u^2+t^2u+t)^3+(tu^2+u+t^2)^3
       =d(u^3+1)^2+d^2(u^3+1)
\]
checks the second equation. Both maps send $P$ to the origin.
The joint map is birational onto its image, since
$u=x_2+t x_1$ and $y=y_1$. Its image therefore has normalization
of genus two and cannot lie in an elliptic subvariety or its
translate. The induced map $J_C\to E^2$ is consequently an isogeny
defined over $k$. On $E$, the chord-and-tangent formula gives
$[2](x,y)=(x^4,y^4+1)$; hence its $\mathbf F_4$-Frobenius is $[-2]$.
The $\mathbf F_{2^{38}}$-Frobenius of $E^2$, and therefore that of
$J_C$, is $[-2^{19}]$. For the last implication compose
$F_{J_C}+[2^{19}]$ with the isogeny: its connected image would lie
in a finite kernel, so it is zero. This proves supersingularity
and the exact scalar Frobenius, without assuming a torsion-group
isomorphism from an unspecified isogeny.

\emph{The theta-orbit criterion.}
A cyclic cover of odd degree $n$ corresponds to a line bundle $L$
of exact order $n$ with a trivialization of $L^n$. Its coherent
cohomology decomposes as
$\bigoplus_{i\in\mathbf Z/n}H^1(C,L^i)$, and absolute Frobenius
sends the $i$-summand semilinearly to the $2i$-summand. The
nontrivial summands have dimension one by Riemann--Roch.
For a multiplication-by-two orbit $O$ of primitive indices, its
contribution to the two-rank is $|O|$ times the eventual rank of
the orbit composite. This follows from the Artin--Schreier
sequence: the Frobenius-stable subspace has dimension equal to
$\dim_{\mathbf F_2}H^1_{\et}(D,\mathbf F_2)$. The idempotent
primitive character projector, or the image of
$\prod_{\ell\mid n}(1-\sigma^{n/\ell})$, selects exactly these
primitive summands.

Let $\mathcal B$ be the cokernel of
$\mathcal O_{C^{(1)}}\to F_*\mathcal O_C$. In characteristic two it
is a line bundle with square the canonical bundle: locally write
a change of uniformizer as $v=a^2+u b^2$; its class in the
cokernel changes by $b$, whose square is $dv/du$. The cohomology
sequence tensored with $L^i$ identifies the kernel of the
Frobenius arrow with $H^0(\mathcal B\otimes L^i)$.
There is just one geometric theta characteristic because
$J_C[2](\bar k)=0$, so $\mathcal B=\mathcal O(P)$.
The corresponding theta locus is precisely the symmetric Abel
curve $\Theta=C-P$. Thus zero primitive two-rank requires at
least one theta point in every multiplication-by-two orbit of
primitive points of the cyclic subgroup $\langle L\rangle$.

\emph{A complete finite bound.}
Write $f=\operatorname{ord}_n(2)$,
$e=\gcd(f,19)$, and $M_n=\varphi(n)/e$.
Scalar Frobenius and symmetry of $\Theta$ show that one point
in a required orbit brings its entire multiplication-by-$2^{19}$
orbit into $\Theta$. Consequently there are at least $M_n$
primitive theta points, and also zero. For a nonzero Jacobian
difference, at most two ordered pairs of distinct points of
$\Theta$ realize it: two representations yield two linearly
equivalent degree-two divisors, and the unique hyperelliptic
pencil gives the sole alternative representation. Hence
\[
 M_n(M_n+1)\le2(n-1).
\]
For odd $n$, prime-power multiplication proves
$\varphi(n)\ge(2/3)n^{5/6}$. Indeed
$(\ell-1)^6\ge\ell^5$ for every prime $\ell\ge5$, and the possible
factor at three is at least $64/243>(2/3)^6$.
It follows that $n<10\cdot19^3=68590$.

The standard-library program \path{code/primitive_conductors.py}
exhaustively enumerates all odd integers in that interval and gives
\[
191,229,361,419,457,571,573,647,687,955,1083,1257,1371.
\]
It then imposes the full difference constraint on the mandatory
packets $H=\langle 2^{19},-1\rangle\subset(\mathbf Z/n)^\times$.
Already $\{0\}\cup H$ excludes all but $191,229,457,1257$.
For $457$, every coset of $U=\langle2,-1\rangle$ must contain an
$H$-packet. There are six such cosets. Normalize the first packet
to $H$ by changing a cyclic generator, and enumerate every
remaining packet. The surviving search depths are
$1,17,226,879,262,0$, so there is no admissible selection.
For $1257$, there are two $U$-cosets and each mandatory packet
has 22 elements, half in each nonzero residue class modulo three.
With zero, the residue counts are $(1,22,22)$, forcing
$924$ ordered distinct differences divisible by three. Only
$2(1257/3-1)=836$ are permitted. Thus only $191,229$ remain.
The script retains every input, packet, candidate list and count.

\emph{The two full polynomial certificates.}
For $E$ put $r_n(X)=x([n]R)$ when $X=x(R)$. The addition formula gives
\[
 r_1=X,\qquad r_{2m}=r_m^4,\qquad
 r_{2m+1}=X+(r_m+r_{m+1})^{-2}.
\]
Compute these as reduced rational functions over $\mathbf F_2$.
Independently form
\[
 \psi_1=\psi_2=\psi_4=1,\quad\psi_3=X^4+X,\quad
 \psi_{2m}=\psi_m^4,\quad
 \psi_{2m+1}=\psi_{m+2}\psi_m^3+\psi_{m-1}\psi_{m+1}^3.
\]
For each required $n$ the program verifies that the reduced
denominator of $r_n$ is exactly $\psi_n^2$, with coprime numerator.
Thus its zeros are exactly the finite coordinates of nonzero
$n$-torsion, without importing a division-polynomial formula
that excludes characteristic two.
\[
\begin{array}{c|r|r|r}
 n&\deg\psi_n&\deg\operatorname{num}r_n&
                  \deg\operatorname{den}r_n\\ \hline
191&18240&36481&36480\\
229&26220&52441&52440
\end{array}
\]
Exact arithmetic in the stated field gives
\[
 U_n(u)\psi_n(A(u))+V_n(u)\psi_n(B(u))=1
                    \quad(n=191,229).
\]
The composed polynomial degrees are respectively $36480,52440$;
both B\'ezout coefficients have degree four less in each case.
All eight polynomials are supplied as complete binary coefficient
arrays, with ascending unsigned little-endian 64-bit coefficients;
bit $i$ represents $t^i$. The manifest gives SHA-256 hashes.
The producer checks both whole identities. The separately written
\path{code/primitive_verify.py} loads the arrays, verifies their
hashes and degrees, independently composes the polynomials over
$\mathbf F_2$, and verifies the two whole identities again.

If a finite point $Q$ of $C$ gave a nonzero Abel class killed by
one of these integers, both of its finite elliptic images would
be killed by the integer. Its coordinate $u$ would therefore be a
common zero of the two composed polynomials, contradicting the
B\'ezout identity. The point at infinity is the origin and cannot
be primitive. Every possible conductor has now been excluded,
over $\bar k$ itself. This proves the counterexample.
\end{proof}
\paragraph{Check record.} Also checked in this writing, 2026-11-18: GPT-6 Astra (OpenAI, ChatGPT Pro mode), root; separate effort setting not exposed. Read isolating-cover arguments, both F4 examples, scalar-Frobenius Lang cover, conductor reduction and explicit primitive-zero obstruction. Degree-one-allowed B.24 remains Open. Certificate reruns are recorded separately. | Also checked in this writing, 2026-11-18: GPT-6 Astra Pro; finite\_checks agent; effort not separately exposed. Read full proof, including theta-orbit reduction, difference bound, conductor cutoff, elliptic-map identity and scalar Frobenius. Reran primitive\_verify on all eight full binary arrays. The entire finite conductor search was also rerun; all candidate lists and packet exclusions agree.



\noindent\emph{Statement B.24 and its produced argument are retained in Part II.}
\subsection{Rational witnesses on the explicit primitive-cover counterexample}

The geometric counterexample in \statementref{B.23} is not a counterexample
to the rational class-number criterion. We give two exact witnesses,
including the least possible prime over all constant extensions.
The computations certify the class-number parity; they do not claim
to compute the full characteristic-zero zeta numerator.

\begin{claim}{B.25}{Proved}{A normalized Cartier test in the case used here}
Let $K=\mathbf F_{2^d}$, and let
\[
 C:\quad Y^2+hY=f(x),\qquad h\in K^\times,
 \quad f\text{ monic of degree five},
\]
be smooth, with marked point $P=\infty$. Let $\ell$ be an odd prime
such that $\mu_\ell\subset K$, and let $a\in J_C(K)$ have exact
order $\ell$. Suppose every nonzero $ia$ has reduced Mumford
representative $(u_i,v_i)$ of degree two. Write $j=2i\bmod\ell$ and
\[
 \begin{split}
 w_i&=(v_i^2+hv_i+f)/u_i,\\
 V_i&=v_i+u_i\bigl((w_i/h)\bmod u_i\bigr),\\
 c_i&=(V_i)_3\,(u_j)_1.
 \end{split}
\]
Here subscripts on polynomials indicate coefficients. Then the
unique constant cyclic, $P$-split cover with character subgroup
$\langle a\rangle$ has odd relative class number if and only if
\[
 A_i:=\prod_{s=0}^{d-1}c_{2^s i}^{\,2^{d-1-s}}\ne1
       \qquad(1\le i<\ell).
\]
The index on each $c$ is reduced modulo $\ell$. The assertion includes
the specified pointing; there is no unspecified constant twist.
\end{claim}
\begin{proof}
Write $ia=[D_i-2P]$ and $L_i=\mathcal O(D_i-2P)$.
The Artin--Schreier model has a local parameter $t=x^2/Y$ at $P$;
monicity gives leading residues one for $xt^2$ and $Yt^5$.
Since $dx/h$ has divisor $2P$, Riemann--Roch and Serre duality
give dual bases
\[
 \omega_i=u_i\,dx/h\quad\text{of }H^0(C,\Omega_C\otimes L_i^{-1}),
 \qquad e_i\quad\text{of }H^1(C,L_i).
\]
For example, the divisor of $u_i\,dx/h$, after subtracting
$D_i-2P$, is $\iota D_i$, which is effective. Each space has
dimension one because $L_i$ is nontrivial of degree zero.

The displayed interpolation formula is the unique polynomial of
degree at most three congruent to $v_i$ modulo $u_i$ for which
$u_i^2\mid V_i^2+hV_i+f$. Indeed, modulo $u_i^2$ the correction
$u_i b$ changes the norm by $hu_i b$. Since the reduced double
has degree two, $V_i$ is cubic and Cantor reduction gives
\[
 V_i^2+hV_i+f=(V_i)_3^2u_i^2u_j,
 \qquad v_j\equiv V_i+h\pmod {u_j}.
\]
Consequently
\[
 g_i=\frac{Y+V_i}{(V_i)_3u_j},\qquad
 \operatorname{div}(g_i)=2D_i-D_j-2P.
\]
Its leading unit at $P$ is one. Multiplication by $g_i$ therefore
identifies $L_i^2$ with $L_j$ respecting their rigidifications at $P$.

The residue definition of Cartier gives
$\langle g_iF(z),\eta\rangle
 =\langle z,\mathcal C(g_i\eta)\rangle^2$.
Using $Y/h=(Y/h)^2+f/h^2$, the even-power terms have Cartier zero,
whereas the odd-power part of $f+hV_i$ equals that of
$(V_i)_3^2u_i^2u_j$. Thus
\[
 \mathcal C\bigl((Y+V_i)dx/h\bigr)
      =(V_i)_3\sqrt{(u_j)_1}\,u_i\,dx/h,
 \qquad
 \mathcal C(g_i\omega_j)=\sqrt{c_i}\,\omega_i.
\]
It follows that absolute Frobenius sends $e_i$ to $c_i e_{2i}$
square-semilinearly. Iterating $d$ times gives $A_i e_i$, because
$2^d\equiv1\pmod\ell$ and every representative is defined over $K$.

For completeness, the cover and the normalization can be described
without a choice of a cohomological scalar. Since $a$ has exact
order $\ell$, a rational function with divisor $\ell(D_1-2P)$
exists over $K$. Scale it to leading unit one at $P$. Adjoining
its $\ell$th root gives a connected cyclic \emph{\'etale} cover;
all its valuations are divisible by $\ell$, its character has
exact order $\ell$, and its leading unit makes the fibre over
$P$ split completely. Any two descents with the same geometric
character differ by a constant character, whose evaluation at
$P$ proves uniqueness of the split descent. Equivalently,
rigidifying all the $L_i$ at $P$ makes all their multiplication
isomorphisms canonical: any associativity discrepancy is a
constant with value one at $P$ and therefore is one. The $g_i$
used above are exactly these canonical doubling isomorphisms.

Finally, the Artin--Schreier sequence identifies the geometric
\'etale $\mathbf F_2$ cohomology with the Frobenius-stable part
of coherent cohomology; see also the characteristic-$p$
comparison in \cite[Section 9]{MilneJac}. A finite-dimensional
$K$-space with square-semilinear operator $T$ has a nonzero
$T$-fixed vector over $K$ exactly when $T^d$ has eigenvalue one:
on $\ker(T^d-1)$ the operator defines descent from $K$ to
$\mathbf F_2$, and ordinary finite-field semilinear descent
has the same dimension. Applying this on the nontrivial
character summands shows that rational new two-torsion exists
exactly when at least one $A_i$ is one. Dualizing the finite
\'etale two-torsion representation does not change the dimension
of its Frobenius-fixed space. Since the cover degree is odd,
the invariant and new parts split by the deck averaging
idempotent, and $\operatorname{Nm}\circ\pi^*=[\ell]$ is invertible
on two-primary groups. The relative class number is odd
exactly when that new rational two-primary group is zero.
\end{proof}
\paragraph{Check record.} Also checked in this writing, 2026-11-18: GPT-6 Astra (OpenAI, ChatGPT Pro mode), root; separate effort setting not exposed. Read normalized Cartier formula, rigidification at the marked point, both exact certificate interpretations, and the involution/Frobenius all-extension lower bound proving least prime 191. | Also checked in this writing, 2026-11-18: GPT-6 Astra Pro; finite\_checks agent; effort not separately exposed. Read and checked all displayed divisor, interpolation and Cartier identities, P-rigidification, and finite-field semilinear descent criterion. The general Artin–Schreier comparison is used; cited Milne section not separately reopened by this reader.


\begin{claim}{B.26}{Proved}{Exact degree-191 and degree-229 rational witnesses}
For the curve and point $(C,P)$ of \statementref{B.23}, a constant
cyclic $P$-split cover of degree $191$ with odd relative class
number exists over $\mathbf F_{2^{380}}$. A second such cover,
of degree $229$, exists over $\mathbf F_{2^{152}}$.
Their genera are respectively $192$ and $230$, and their new
Jacobian parts are ordinary of dimensions $190$ and $228$.
\end{claim}
\begin{proof}
We specify the fields, the normalization, and the exact finite
certificates. No search heuristic is used as evidence of their parity.
Put $d_0=t^3+1$ and use $A(u),B(u)$ from \statementref{B.23}.
Write $A(u)^3=\sum_{j=0}^6 a_j u^j$. Successively eliminate the
terms of degree six, four and two by adding $S(u)^2+S(u)$, with
$S$ cubic. The resulting polynomial is
\[
 f_0(u)=a_5u^5+b_3u^3+b_1u+b_0,
 \qquad a_5=t^2/d_0^3\ne0.
\]
No field extension is required for these square roots in a finite
field. Set
\[
 u=a_5x,\qquad y=a_5^3Y+S(a_5x),\qquad
 h=a_5^{-3},\qquad f(x)=a_5^{-6}f_0(a_5x).
\]
Then $Y^2+hY=f(x)$ is monic of degree five and has the same marked
point at infinity. In particular the normalization is an isomorphism
over the stated field, not a scalar extension omitted from the result.

The two certificate fields are
\[
 \begin{array}{c|c|c}
 \ell&K&[K:k]\\ \hline
 191&\mathbf F_2[z]/(z^{380}+z^{47}+1)&10\\
 229&\mathbf F_2[z]/(z^{152}+z^6+z^3+z^2+1)&4.
 \end{array}
\]
The following hexadecimal strings encode elements of $K$: binary
bit $j$ is the coefficient of $z^j$. Consecutive lines belonging
to one entry are concatenated. They give the embedding of $k$ and
the reduced representative $u=x^2+u_1x+u_0$, $v=v_1x+v_0$ of the
chosen class $a$ on the normalized model.
\begin{quote}\small
\begin{verbatim}
191:
 t  = cebd4b3dac18387ef98c70a8bc69661f941f1ab70330e1fb
      19cea93138d8ab1686e3f52a7f31d78c0d349454dd0a7a3
 u0 = 35800adb325893dcca9e4a4670173048e3e683b2ccd8383d
      42236023341647d76e22bcb0bf5df97910948458ab1fd37
 u1 = e9edfad1fe30d89fd0afb8a9cacda03aca4c455c01b4c42
      a0f679640b5f557a33700a882ea77325babe828b006e95a4
 v0 = fea3805d1c99d0cdcad5dd87afc416eb89ff67d3976e543
      0580f7152b05537b2ad904bc24804bfa1cc3fdd706890fec
 v1 = e8156100e879ff743228cc826dbe666030ae4636505cf28
      83274f15cd61a2e7e3cdd52bacc8324e95038057b70470e5
229:
 t  = 80fafae83a3d67cbb2033e0b744dcc9de8aa86
 u0 = 8f276acfd0d6f09be86235d08a7ea29b9771cb
 u1 = 47138bd65385caa096e9e047779586fa5eda15
 v0 = 8f83e3c8bccd0c1f58d6f6168617be21fab882
 v1 = 568c198f174a8035b7a6d31f2c28bce57c2db0
\end{verbatim}
\end{quote}

The files \path{code/r8_rational_witness_191.json} and
\path{code/r8_rational_witness_229.json} retain these data and
every reduced nonzero multiple of $a$, every $c_i$, and every $A_i$.
Their SHA-256 hashes are respectively
\begin{quote}\small\ttfamily
2ccc98fdd5ae31557e6f1977ceb2dd1f51face99f9a08b48f11e6d7808bee256
\par
73cee7d15efe7b6620c01660605dc4d49ab128856163b5375ea6a0fb1e93851c.
\end{quote}
The independent verifier \path{code/criterion_r8_verify.py} reads
these files and imports no producer routines. A separate binary
polynomial implementation proves both field polynomials irreducible
by Rabin's criterion, and proves the degree-38 base polynomial
irreducible. It verifies $M(t)=0$, the displayed normalization, and
the elliptic-map identity from \statementref{B.23}.

For every supplied pair $(u_i,v_i)$, the verifier checks reducedness,
$u_i\mid v_i^2+hv_i+f$, the tangent interpolation, and both the
$x$- and $Y$-coordinates of its reduced double. Independently of
the producer's general Cantor addition formula, it reconstructs
$ia+a$ by coprime Chinese-remainder interpolation for
$2\le i\le\ell-2$. All these supports are coprime in the supplied
data. It also checks $2a$ by tangent interpolation and
$(\ell-1)a=-a$. Thus $a$ is nonzero of exact prime order $\ell$,
and all the supplied indices are certified. The verifier additionally
checks the displayed class as the sum of the two saved elliptic
pullbacks, although the order and divisor-sequence verification
already suffice for the existence assertion.

It recomputes every $c_i$ and recomputes each product $A_i$ directly
from separately accumulated squares, rather than reusing the
producer's Horner recurrence. The exact results are
\[
 \begin{array}{c|r|r|r|r}
 \ell&\text{nonzero }(u_i,v_i)&\text{CRT additions}&
                 \text{checked }c_i&\#\{i:A_i=1\}\\ \hline
 191&190&188&190&0\\
 229&228&226&228&0.
 \end{array}
\]
Every $c_i$ and $A_i$ is nonzero. Extra identities, not used by
the producer, give useful checks on the indices and rigidifications:
\[
 A_{2i}=A_i^2,\qquad A_{-i}=A_i,\qquad
 A_i^{2^{19}}=A_i,\qquad
 A_{-2^{19}i}=A_i^{2^{38}}.
\]
As a compact exact record, let $m_b(X)=\sum_j b_jX^j$ for the
binary digits $b_j$ of the positive integer $b$. The characteristic
polynomial of the computed linear coherent Frobenius on the new
part is
\[
 \begin{cases}
 m_{892407}(X)^{10},&\ell=191,\\
 m_{713865}(X)^4m_{785567}(X)^4m_{913207}(X)^4,&\ell=229.
 \end{cases}
\]
Each displayed $m_b$ has degree nineteen, is the minimal polynomial
of a checked multiplier, and has $m_b(1)=1$. The verifier obtains
it by multiplying all nineteen distinct Frobenius conjugates and
checking that every coefficient is in $\mathbf F_2$.
The full results and input hashes are in
\path{code/criterion_r8_verify.json}; the concise run transcript is
\path{code/criterion_r8_verify.out}.

The scalar Frobenius of \statementref{B.23} is $s=-2^{19}$.
Direct modular multiplication gives
\[
 \operatorname{ord}_{191}(2)=95,
 \quad\operatorname{ord}_{191}(s)=10,
 \qquad
 \operatorname{ord}_{229}(2)=76,
 \quad\operatorname{ord}_{229}(s)=4.
\]
Thus the stated fields contain all the corresponding Jacobian
torsion and all the roots of unity; the constant-deck descent
condition $\sigma_{\#K}(a)=(\#K)a$ holds. Statement
\statementref{B.25} now proves both odd relative class numbers.
Nonvanishing of every $c_i$ makes the absolute Frobenius maps
invertible on all nontrivial coherent character lines, so the
new parts are ordinary. Riemann--Hurwitz gives genus $\ell+1$,
and their dimension is $\ell-1$ as asserted.
\end{proof}
\paragraph{Check record.} Also checked in this writing, 2026-11-18: GPT-6 Astra (OpenAI, ChatGPT Pro mode), root; separate effort setting not exposed. Read normalized Cartier formula, rigidification at the marked point, both exact certificate interpretations, and the involution/Frobenius all-extension lower bound proving least prime 191. | Also checked in this writing, 2026-11-18: GPT-6 Astra Pro; finite\_checks agent; effort not separately exposed. Read complete field normalization and exact-torsion proof. Reran both producer cases and independent criterion\_r8\_verify, including irreducibility, CRT additions, all Cartier multipliers and minimal polynomials.


\begin{claim}{B.27}{Proved}{The exact least prime and the two field minima}
For the fixed marked curve $(C,P)$ of \statementref{B.23}, the
least prime degree of a constant cyclic $P$-split cover satisfying
the rational class-number criterion over any finite extension
of $k=\mathbf F_{2^{38}}$ is exactly $191$.
At degree $191$ the least possible constant extension degree is ten;
at degree $229$ it is four, and each minimum is attained by
\statementref{B.26}. The second assertion is a minimum for that
specified prime, not a claimed minimum over all primes.
\end{claim}
\begin{proof}
We include the pointed-involution argument, since a geometric
positive-rank obstruction alone would not prove a rational
lower bound. Fix any geometric connected cyclic cover of odd
prime degree $\ell$. Scalar Frobenius preserves its character
line, so its geometric isomorphism class is already defined
over $k$. Choosing a point above $P$ gives its unique pointed
descent. On its deck group Frobenius acts by $s=-2^{19}$:
Cartier duality changes the scalar on $J_C[\ell]$ into
$(\#k)/s=s$. Therefore its first constant-deck field has degree
$m=\operatorname{ord}_\ell(s)$ over $k$; every later
constant-deck field has degree divisible by $m$.

The hyperelliptic involution fixes $P$ and acts as minus one
on $J_C$. Its unique lift fixing the selected point above $P$
inverts the deck group. Uniqueness makes the lift commute with
the arithmetic Frobenius of the pointed descent. Let
\[
 R=\mathbf F_2[X]/\Phi_\ell(X),\qquad
 W=H^1_{\et}(D_{\bar k},\mathbf F_2)_{\rm new}.
\]
Every nontrivial coherent character space has dimension one.
The injection of geometric \emph{\'etale} cohomology into its
Frobenius-stable coherent part consequently makes $W$ free
of rank one over a direct factor $R'$ of the finite \emph{\'etale}
algebra $R$. A zero-rank character factor is simply omitted.

Frobenius $F$ is semilinear for the automorphism $X\mapsto X^s$
of $R'$. Since $s^m=1$, its $m$th power is $R'$-linear, hence
multiplication by a unit $b\in R'$. Commuting $F^m$ with $F$
and with the lifted involution forces
\[
 b\in\bigl(R'^{\langle s,-1\rangle}\bigr)^\times.
\]
Put $f_\ell=\operatorname{ord}_\ell(2)$. If $19\nmid f_\ell$,
then $\langle2\rangle\subset\langle2^{19},-1\rangle
=\langle s,-1\rangle$. On the finite \emph{\'etale} algebra $R$,
the automorphism $X\mapsto X^2$ is the absolute Frobenius.
Its fixed ring is a product of copies of $\mathbf F_2$, whose
unit group is trivial. Therefore $b=1$. Every geometric new
two-torsion point is then rational at the first constant-deck
field and remains rational over all subsequent ones.

By \statementref{B.23}, the new part of every prime-degree
geometric cover has positive two-rank. Thus none of these
covers can have odd relative class number when
$19\nmid\operatorname{ord}_\ell(2)$. If $\ell<191$ and
$19\mid\operatorname{ord}_\ell(2)$, then $19\mid\ell-1$;
the only possibilities for an odd prime are
$\ell\in\{39,77,115,153\}$, all composite. This excludes every
odd prime below $191$ over every finite extension. Degree two
is excluded by the prime-to-characteristic requirement.

The degree-$191$ certificate proves attainment, and the necessary
constant-deck degrees computed above prove the two field minima.
For clarity, allowing a different unpointed constant twist does
not evade the lower bound: splitting above the specified point
selects the unique pointed twist, to which the commuting
involution argument applies.
\end{proof}
\paragraph{Check record.} Also checked in this writing, 2026-11-18: GPT-6 Astra (OpenAI, ChatGPT Pro mode), root; separate effort setting not exposed. Read normalized Cartier formula, rigidification at the marked point, both exact certificate interpretations, and the involution/Frobenius all-extension lower bound proving least prime 191. | Also checked in this writing, 2026-11-18: GPT-6 Astra Pro; finite\_checks agent; effort not separately exposed. Read full least-prime proof. Checked the pointed involution, Frobenius-semilinear scalar argument, fixed subring reduction and exclusion of all primes below 191; first constant-deck periods verified by rerun.


\paragraph{Consequences and scope.}
The same explicitly specified geometric curve has positive primitive
two-rank for every nontrivial odd cyclic cover, yet possesses rational
class-number witnesses. It therefore separates the geometric
primitive-zero condition from the rational criterion on a fixed
curve, with exact field and prime data. This is a positive instance
of B1, not a proof of its universal assertion, including the proposed
range $p>(2g-2)^2$.


\subsection{What varying the constant twist does and does not prove}

\begin{claim}{B.28}{Proved}{A uniform count of bad constant twists}
Let $C/K$ be smooth proper of genus $g\ge2$ over a finite field
of characteristic $p$. Let $D\to C$ be a geometrically connected
\emph{\'etale} cover with constant cyclic deck group of odd prime
order $\ell\ne p$. Put $f=\operatorname{ord}_\ell(p)$.
Among the $\ell$ constant twists of this cover, at most
\[
                \frac{(g-1)(\ell-1)}{f}
\]
have relative class number divisible by $p$. In particular, if
$f\ge g-1$, at least one twist has $p$-unit relative class number.
For genus two this last condition is automatic.
\end{claim}
\begin{proof}
Fix a generator $\delta$ of the deck group. The nontrivial factor
of its group algebra is
\[
 R=\mathbf F_p[X]/\Phi_\ell(X)
       \simeq\prod_{j=1}^{(\ell-1)/f}\mathbf F_{p^f}.
\]
Let $W$ be the new part of
$H^1_{\et}(D_{\bar K},\mathbf F_p)$. Over an algebraic closure,
the nontrivial coherent character spaces are $H^1(C,L)$ for
nontrivial degree-zero line bundles $L$ of order $\ell$; each
has dimension $g-1$. The Artin--Schreier comparison embeds
\'etale cohomology tensored with $\bar K$ into the Frobenius-stable
coherent part, equivariantly for the deck action. Consequently the rank $d_j$ of $W$ over each field
factor of $R$ satisfies $0\le d_j\le g-1$.

The deck group is constant, so arithmetic Frobenius $F$ commutes
with its action and is linear over every field factor. Twisting
by the constant cocycle with Frobenius value $\delta^t$, for
$t\in\mathbf Z/\ell$, replaces its action on $W$ by
$\delta^tF$, up to using $-t$ consistently in the contravariant
cohomology convention. On the $j$th field factor, write the image
of $\delta$ as the primitive root $\zeta_j$. A new fixed vector
exists on this factor exactly when
\[
               \det(\zeta_j^{-t}-F\mid W_j)=0.
\]
This polynomial has degree $d_j$, and its $\ell$ arguments are
distinct. At most $d_j$ twists are therefore bad on this factor.
Taking the union over the $(\ell-1)/f$ factors gives the claimed
bound. Nonsemisimplicity of $F$ causes no problem: a fixed vector
is detected exactly by this determinant. The rational two-primary
argument in \statementref{B.25} applies with $p$ in place of two:
norm and pullback split the $p$-primary groups because
$p\nmid\ell$, and rational new $p$-torsion is equivalent to
divisibility of the relative class number. Finally,
$f\ge g-1$ makes the bound at most $\ell-1$, strictly below the
number of twists.
\end{proof}
\paragraph{Check record.} Also checked in this writing, 2026-11-18: GPT-6 Astra (OpenAI, ChatGPT Pro mode), root; separate effort setting not exposed. Read unpointed twist count and the prescribed-point obstruction. This invalidates the proposed contraposition but is not a formal logical independence theorem for A.1 and universal B.1.


\begin{claim}{B.29}{Proved}{Existence with an unspecified base point}
For every smooth proper curve $C/k$ of genus $g\ge2$ over a
finite field, there are a finite extension $K/k$, a point
$P_K\in C(K)$, and a constant cyclic $P_K$-split cover of odd
prime degree, different from the characteristic, with $p$-unit
relative class number. The point $P_K$ is not prescribed in
advance. More quantitatively, for a constant degree-$\ell$ cover
as in \statementref{B.28}, put
$B=(g-1)(\ell-1)/f$ and $g_D=1+\ell(g-1)$. If $f\ge g-1$, the
number of $P\in C(K)$ for which its unique split twist is a
class-number witness is at least
\[
       \frac{\ell-B}{\ell}
             \bigl(\#K+1-2g_D\sqrt{\#K}\bigr).
\]
\end{claim}
\begin{proof}
For each twist $D_t$, let $N_t$ be the number of rational points
of $C$ that split in it. Every geometric fibre has $\ell$ points
and Frobenius is a deck translation. Thus each $P\in C(K)$
splits in precisely one twist, and
$\#D_t(K)=\ell N_t$. At least $\ell-B$ twists have $p$-unit
relative class number. Applying the Weil lower bound to each
of these geometrically connected curves gives the displayed
inequality \cite[Theorem 1.6 and formula 1.5.1]{WeilCurves}.
If its right-hand side is negative, the inequality is automatic;
otherwise summing over at least $\ell-B$ good twists proves it.

There exist odd primes $\ell\ne p$ with
$\operatorname{ord}_\ell(p)\ge g-1$: the primes with order less
than $g-1$ divide the fixed nonzero integer
$\prod_{j=1}^{g-2}(p^j-1)$ and are finite in number.
Choose such an $\ell$ and a nonzero geometric cyclic
$\ell$-cover, which exists because the prime-to-$p$ abelian
fundamental group has rank $2g$ at $\ell$.
It, its deck group, and a model of its base curve descend to a
finite extension. Enlarge that field until
$\#K+1>2g_D\sqrt{\#K}$. The quantitative assertion then produces
a good twist with a rational point above some $P_K$.
All descent and field enlargements precede the choice of a good
twist; no assertion that a previously good twist remains good
under extension is needed.
\end{proof}
\paragraph{Check record.} Also checked in this writing, 2026-11-18: GPT-6 Astra (OpenAI, ChatGPT Pro mode), root; separate effort setting not exposed. Read unpointed twist count and the prescribed-point obstruction. This invalidates the proposed contraposition but is not a formal logical independence theorem for A.1 and universal B.1.


\paragraph{The remaining pointing obstruction in B1.}
For a fixed prescribed $P$, its split twist is one specified
element of the $\ell$ twists. The counting theorem can leave
that element among the bad twists. There is no way to repair
this by simply extending a good unpointed twist until $P$
splits: this extension may make its nontrivial $p$-torsion
eigenvalues equal to one.

This failure occurs on the explicit R8 curve, not just on an
abstract representation. Consider the cyclic degree-three
character obtained by pulling back the point $(0,0)$ of exact
order three on $E:y^2+y=x^3$ along its first degree-two map
$C\to E$. Norm followed by pullback is multiplication by two,
so this character is nonzero. Its pointed cover is constant over
$k=\mathbf F_{2^{38}}$, since $s=-2^{19}\equiv1\pmod3$.
Its new geometric two-rank is positive by \statementref{B.23};
the nontrivial group algebra is the single field $\mathbf F_4$,
and the coherent character dimension is one, so its new
\'etale cohomology has rank one over that field. The commuting
pointed hyperelliptic involution used in \statementref{B.27}
forces Frobenius to be the identity. Both nontrivial constant
twists instead have primitive third-root Frobenius multiplier,
and hence odd relative class number. They do not split over $P$.
The cubic constant extension first makes that fibre split and
simultaneously cubes the multiplier to one, making the relative
class number even again. Thus the unpointed existence theorem
does not settle B1, including the range $p>(2g-2)^2$.

These three covers can be given explicitly in the original R8
coordinates by
\[
       w^3=y,\qquad w^3=ty,\qquad w^3=t^2y.
\]
On the elliptic curve, $\operatorname{div}(y)=3(0,0)-3O$;
its pullback gives the exact degree-three character above.
If $\tau=x/y$ is the elliptic local parameter at $O$, then
$y\tau^3$ has residue one. The degree-two map has
$\tau=b\eta^2+\cdots$ at $P$ for a rational local parameter
$\eta$ and $b\in k^\times$. The leading unit of $y$ at $P$
is therefore $b^{-3}$, a cube, proving that the first cover
is precisely the split twist.
The pure binary-polynomial program
\path{code/criterion_twist_obstruction.py} verifies irreducibility
of $M$ and computes
\[
        t^{(2^{38}-1)/3}=588964846
\]
in the polynomial-bit encoding of $k$. This value is a nontrivial
cube root of unity, so $t$ is a noncube and the other two covers
first split above $P$ over the cubic extension. The same program
checks the two-dimensional new cohomology matrices
\[
 I,\qquad
 \begin{pmatrix}0&1\\1&1\end{pmatrix},\qquad
 \begin{pmatrix}1&1\\1&0\end{pmatrix},
\]
up to interchanging the last two by the chosen cohomology
convention. Their determinants $\det(1-F)$ are $0,1,1$ in
$\mathbf F_2$, and all three matrices cube to $I$.
The full finite data and the geometric scope of this check are
recorded in its JSON and run transcript.


\subsection{Exact extension-field tests for plane quartics}

\begin{claim}{B.30}{Proved}{A least-prime-five quartic witness over \(\mathbf F_{3^8}\)}
Let $C/\mathbf F_3$ be the smooth plane quartic
\[
                W^4=X^3Z-XZ^3,\qquad P=(1:0:0).
\]
There is a $P$-split constant cyclic degree-five cover over
$\mathbf F_{3^8}$ whose relative class number is prime to three.
Thus five is the least odd prime different from the characteristic
which gives a witness for this pointed curve over some constant
extension. Every constant degree-five cover of this curve requires
a constant degree divisible by eight; the displayed field is
therefore also the least one for a degree-five witness.
The certificates give six such covers. For each of them the
relative class number stays prime to three exactly under further
constant extensions of odd degree.
\end{claim}
\begin{proof}
The three partial derivatives of the displayed homogeneous
polynomial have no common projective zero: in characteristic
three they are $-Z^3,W^3,X^3$, up to nonzero constants.
The unique point at infinity is $P$. On the affine chart put
$w=W/Z$ and $x=X/Z$. The degree-two morphism
\[
       C\longrightarrow E:\ y^2=x^3-x,
                    \qquad (x,w)\longmapsto(x,w^2)
\]
is branched, in particular at infinity. The elliptic curve has
four points over $\mathbf F_3$, so its Frobenius satisfies
$F^2=[-3]$. Consequently $E[5]$ is constant over
$K=\mathbf F_{3^8}$.

The program \path{code/quartic_cyclic5.py} works in
\[
 K=\mathbf F_3[z]/
       (z^8+2z^5+z^4+2z^2+2z+2).
\]
It specifies a basis of $E[5]$ and all six cyclic subgroups.
For a generator $T=(a,b)$ and $2T=(c,d)$ put
$h=(x-a)(x-c)$. The normalized isogeny has
\[
 r(x)=N(x)/h(x)^2,
 \qquad r'(x)=g(x)/h(x)^3,\qquad
       g=N'h-2Nh',
\]
where
\[
 N=xh^2+
       \sum_{(u,v)\in\{T,2T\}}
        \left(\frac{h^2}{x-u}+\frac{v^2h^2}{(x-u)^2}\right).
\]
These are polynomials. The addition formula on $E$ proves
this isogeny formula by summing $x(R+S)-x(S)$ over the four
nonzero kernel points. The program verifies its defining
identity
\[
 (x^3-x)g^2=N^3+A'Nh^4+B'h^6.
\]
It gives $s,r_0\in K$ with $s^4=-A'$ and
$r_0^3+A'r_0+B'=0$, hence the isogeny to $E$ is
\[
       x_E=(N/h^2-r_0)/s^2,
       y_E=yg/(s^3h^3).
\]
All its kernel points are $K$-rational. Pulling back
$C\to E$ gives the cover whose normalization has equation
\[
       v^4=F(x),\qquad
       F=(x^3-x)g^2h^2/s^6,\qquad
       (X_C,W_C)=((N/h^2-r_0)/s^2,v/h^2).              \tag{B.Q.1}
\]
It is connected: the branched quadratic extension defining
$C/E$ is disjoint from this degree-five unramified isogeny.
Base change makes its map to $C$ finite \emph{\'etale}, cyclic
of degree five, and constant over $K$. Infinity is a rational
point above $P$, and its deck translates split that fibre.
Its genus is $1+5(3-1)=11$.

We give the entire finite cohomological test, not a point-count
guess. Here $\deg F=19$, and $f=x^3-x$, $g$, and $h$ are
pairwise coprime and squarefree, of degrees three, six, and two.
A basis of regular differentials on the normalization is
\[
 \left\{\frac{x^jdx}{v}:0\le j\le3\right\},\quad
 \left\{\frac{gh\,dx}{v^2}\right\},\quad
 \left\{\frac{ghx^jdx}{v^3}:0\le j\le5\right\}.       \tag{B.Q.2}
\]
At a root of $f$ the ramification index is four; at a root of
$gh$ it is two. Checking the orders there proves regularity.
At infinity, $v_\infty(x)=-4$, $v_\infty(v)=-19$ and
$v_\infty(dx)=-5$, giving exactly the displayed ranges.
The eleven forms are independent by their distinct cyclic
characters and powers of $x$, so the genus calculation makes
them a basis.

For a differential $H(x)dx/v^a$, choose $b\in\{1,2,3\}$
with $3b\equiv a\pmod4$. Cartier is computed exactly by
\[
 \frac{1}{v^b}\,\mathcal C_{\mathbf P^1}
       \bigl(H(x)F(x)^{(3b-a)/4}dx\bigr),\qquad
 \mathcal C_{\mathbf P^1}\left(\sum c_ix^i dx\right)
       =\sum c_{3j+2}^{1/3}x^jdx.
\]
Division by $gh$ in the relevant character spaces is checked
with zero remainder. If $M$ is its inverse-Frobenius-semilinear
matrix, eight iterations produce its $K$-linear matrix. Every
entry of $M$, that eighth power, and the cover polynomials is
included in \path{code/quartic_cyclic5.json}. For each of the
six covers the result is
\[
     \det(T-\mathcal C^8)=T^3(T+1)^8.                 \tag{B.Q.3}
\]
The same calculation on $C$, with basis
$dx/w^2,dx/w^3,x\,dx/w^3$, gives Cartier zero.
Serre duality identifies the linear eighth powers of Cartier
and coherent Frobenius as transposes. The Artin--Schreier
sequence identifies the Frobenius-stable coherent subspace
with geometric three-torsion cohomology. Thus the cover's
new three-primary part has rank eight and Frobenius eigenvalue
$-1$ only; its relative class number is a unit. Equation
(B.Q.3) also gives the exact odd-extension condition. This
conclusion is insensitive to possible Jordan blocks.

For completeness, exact enumeration over $\mathbf F_3$,
$\mathbf F_9$, and $\mathbf F_{27}$ gives $4,4,28$ points
on $C$. Newton identities and the reciprocal functional
equation give its Frobenius polynomial
\[
                     (T^2+3)(T^2-3)^2.
\]
For an abelian variety over a finite field,
\citep[Chapter~II, Lemma~1.3 and Remark~1.4]{MilneAV2008}
prove that the algebra generated by Frobenius is a product
of fields. Its minimal polynomial is therefore squarefree,
so $F^4=[9]$ on $J_C$;
equality on a Tate module implies equality of these
endomorphisms. On every order-five character this says
$F^4=-1$ and $F^8=1$. Hence every nonzero character has
period exactly eight, proving the constant-field minimum.
Finally three is excluded in characteristic three and five
is the first admissible odd prime. The independently loaded
certificates verify the finite arithmetic in this argument.
\end{proof}
\paragraph{Check record.} Also checked in this writing, 2026-11-18: GPT-6 Astra (OpenAI, ChatGPT Pro mode), root; separate effort setting not exposed. Read quartic isogeny pullbacks, all differential bases and Cartier formulas, characteristic polynomials and exact minimality arguments. B.32 remains Open; B.34 first-field minimum is degree 3 with prime 13, distinct from least-prime cubic witness over degree 8. | Also checked in this writing, 2026-11-18: GPT-6 Astra Pro; finite\_checks agent; effort not separately exposed. Read full cyclic5 quartic construction, regular differential basis and finiteCartiertest; reran supplied verifier; literature semisimplicity citation not reopened by this reader.


\begin{claim}{B.31}{Proved}{A four-cover all-extension negative test over \(\mathbf F_5\)}
Let $C/\mathbf F_5$ have the same equation and point as in B.30.
Pull its quadratic map to $E:y^2=x^3-x$ back along the duals
of the four cyclic degree-three isogenies out of $E$, over
$K=\mathbf F_{5^8}$. These four constant $P$-split covers
have genus seven. Every one has relative class number divisible
by five over $K$ and over every further constant extension.
This assertion concerns these four covers; it does not exclude
the other order-three characters of $J_C$ or other primes.
\end{claim}
\begin{proof}
The program \path{code/quartic_cyclic3_f5.py} computes the four
kernel pairs on $E$ from the roots of its three-division
polynomial $3x^4-6x^2-1$. For a kernel pair generated by
$(a,b)$, the general odd-characteristic isogeny formula is
\[
 h=x-a,\quad N=xh^2+2(3a^2+A)h+4b^2,
 \quad g=N'h-2Nh',
\]
where $E$ has equation $y^2=x^3+Ax+B$. The identity
$fg^2=N^3+A'Nh^4+B'h^6$ is checked coefficient by
coefficient. The image of a different three-torsion point
generates the dual kernel on $E'$. Applying the formula there
gives an isogeny with target coefficients $3^4A,3^6B$;
scaling by three identifies that target with the original $E$.
Thus we have used the dual isogeny, without identifying an
ordinary isogenous elliptic curve with $E$ without proof.

With $f'=x^3+A'x+B'$ on $E'$, the pullback cover is
\[
        v^4=f'g^2h^2/3^6,
        \qquad (X_C,W_C)=(N/(9h^2),v/h^2).
\]
It is a connected constant \emph{\'etale} degree-three
$P$-split cover by the same disjointness and base-change
argument as B.30. Now $\deg F=11$, and $\deg(gh)=4$.
The regular basis consists of $x^jdx/v$ for $j=0,1$,
$gh\,dx/v^2$, and $ghx^jdx/v^3$ for $0\le j\le3$.
The same valuation calculation proves it is a full basis.
Cartier extracts the coefficients numbered $5j+4$ from
$H(x)F(x)^a$ in character $a$, and takes fifth roots.

The exact eighth powers have characteristic polynomials
$(T-1)^3$ on the old genus-three cohomology and
$(T-1)^7$ on the genus-seven cover. Polynomial division
therefore gives $(T-1)^4$ for the new factor, with value
zero at one. The old factor has been divided as a polynomial;
we have not divided the two zero values at one.
The new geometric five-torsion has a Frobenius fixed vector.
It remains fixed by every power, proving the all-extension
failure. All four input kernels, both isogenies, the cover
polynomials, and the matrices are in the accompanying JSON.
\end{proof}
\paragraph{Check record.} Also checked in this writing, 2026-11-18: GPT-6 Astra (OpenAI, ChatGPT Pro mode), root; separate effort setting not exposed. Read quartic isogeny pullbacks, all differential bases and Cartier formulas, characteristic polynomials and exact minimality arguments. B.32 remains Open; B.34 first-field minimum is degree 3 with prime 13, distinct from least-prime cubic witness over degree 8. | Also checked in this writing, 2026-11-18: GPT-6 Astra Pro; finite\_checks agent; effort not separately exposed. Read full fourcubic-isogeny construction and regular basis/Cartierproof; reran independent verifier.



\noindent\emph{Statement B.32 and its produced argument are retained in Part II.}
\subsection{A first-extension witness for the pointed Fermat quartic in characteristic five}

\begin{claim}{B.33}{Proved}{A constant cyclic thirteen-cover over \(\mathbf F_{125}\)}
Let
\[
 C/\mathbf F_5:\quad W^4=X^3Z-XZ^3,\qquad P=(1:0:0).
\]
Over $K=\mathbf F_{125}$ there is a constant cyclic, $P$-split,
\emph{\'etale} cover $D\to C$ of degree thirteen with
\[
  g_D=27,\qquad
  \frac{\#J_D(K)}{\#J_C(K)}\equiv1\pmod5.
\]
Its new Jacobian part is ordinary of dimension twenty-four.
Under an extension of $K$ of degree $n$, its relative class
number stays prime to five if and only if $4\nmid n$.
\end{claim}
\begin{proof}
All the following defining polynomials lie in $\mathbf F_5[x]$:
\[
\begin{split}
 h={}&x^6+x^2+4,\\
 N={}&x^{13}+3x^{11}+x^9+3x^7+2x^5+x,\\
 g={}&x^{18}+2x^{16}+x^{14}+3x^{12}+3x^{10}+2x^2+4.
\end{split}                                                    \tag{B.QM.1}
\]
Direct polynomial multiplication gives
\[
 g=N'h-2Nh',\qquad
       (x^3-x)g^2=N^3-Nh^4.                                    \tag{B.QM.2}
\]
The polynomials $x^3-x,g,h$ are squarefree and pairwise coprime,
and $N$ is coprime to $h$. Consequently
\[
 E:\ y^2=x^3-x\longrightarrow E,\qquad
       (x,y)\longmapsto
       \left(\frac{N(x)}{4h(x)^2},\frac{y g(x)}{2h(x)^3}\right)
                                                                  \tag{B.QM.3}
\]
is a separable degree-thirteen isogeny. Indeed (B.QM.2) verifies
the target equation, the map takes infinity to infinity, and
the rational function defining its first coordinate has degree
thirteen. A nonconstant morphism between elliptic curves taking
the origins to each other is an isogeny; its degree is thirteen
because the $x$-coordinate maps on source and target both have
degree two.

The kernel is constant over the displayed cubic field. A fully
explicit presentation and generator are
\[
 K=\mathbf F_5[z]/(z^3+4z^2+3z+1),\qquad
                    T=(z,z+3)\in E(K).                           \tag{B.QM.4}
\]
The cubic has no root in $\mathbf F_5$, so this is a field.
The elliptic addition law gives $13T=0$, with $T\ne0$.
The twelve nonzero multiples of $T$ have precisely the six
distinct roots of $h$ as their $x$-coordinates. These assertions
are also a short exact arithmetic certificate: its thirteen
points, including the identity, are listed in
\path{code/criterion_quartic13_verify.json}, in the basis
$1,z,z^2$. Thus they constitute the entire kernel of (B.QM.3).

Pull back (B.QM.3) along $C\to E$, $(X,W)\mapsto(X,W^2)$.
The normalization of the resulting connected curve is
\[
 D:\quad w^4=F(x),\qquad
 F(x)=4(x^3-x)g(x)^2h(x)^2,\qquad
 (X_C,W_C)=\left(\frac{N(x)}{4h(x)^2},\frac{w}{h(x)^2}\right).
                                                                  \tag{B.QM.5}
\]
One can verify connectedness directly: $F$ has simple zeros
at the three roots of $x^3-x$, so the fourth-root extension
has degree four even over an algebraic closure. Alternatively,
the degree-two extension $C/E$ and the odd-degree isogeny are
linearly disjoint. The map to $C$ is finite \emph{\'etale} by
base change of the isogeny; it is a constant cyclic cover over
$K$ by (B.QM.4). The fibre of $P$ is the rational kernel of
(B.QM.3), so it splits completely. Riemann--Hurwitz gives
$g_D=1+13(3-1)=27$.

We give the cohomological calculation explicitly enough to
reproduce it without an isogeny-search program. Here $\deg F=51$
and $\deg(gh)=24$. A basis of regular differentials is
\[
 \left\{\frac{x^jdx}{w}:0\le j\le11\right\},\qquad
 \left\{\frac{gh\,dx}{w^2}\right\},\qquad
 \left\{\frac{ghx^jdx}{w^3}:0\le j\le13\right\}.                \tag{B.QM.6}
\]
At a simple zero of $F$ the ramification index is four, while
at a double zero it is two. The factors $gh$ exactly remove
the poles in characters two and three at the latter zeros.
At infinity the orders of $x,w,dx$ are $-4,-51,-5$;
these give exactly the ranges in (B.QM.6). The forms are
independent by their characters and their powers of $x$, and
their number is the genus, proving that this is a full basis.

Put $h_1=1$ and $h_2=h_3=gh$. Since $5\equiv1\pmod4$,
Cartier in character $s$ is computed by
\[
 \mathcal C\!\left(\frac{h_s x^jdx}{w^s}\right)
   =\frac1{w^s}\,
      \mathcal C_{\mathbf P^1}(h_sx^jF^s dx),\qquad
 \mathcal C_{\mathbf P^1}\!\left(\sum a_i x^i dx\right)
        =\sum a_{5b+4}x^b dx.                                  \tag{B.QM.7}
\]
No coefficient roots need to be represented here because
all coefficients lie in $\mathbf F_5$. Divide the extracted
polynomial by $h_s$ and expand in the relevant part of
(B.QM.6); every division has zero remainder. This gives a
$27\times27$ matrix $M$ over $\mathbf F_5$. Its independently
computed characteristic polynomials are
\[
\begin{split}
 \det(U-M)
   &=(U-2)^5(U+2)^6
       (U^2-2U-1)^4(U^2+2U-1)^4,\\
 \det(U-M^3)&=(U-2)^{14}(U+2)^{13}.
\end{split}                                                    \tag{B.QM.8}
\]
For the old curve, the basis $dX/W^2,dX/W^3,X\,dX/W^3$ gives
\[
 M_C=\begin{pmatrix}3&0&0\\0&0&4\\0&1&0\end{pmatrix},\qquad
                \det(U-M_C^3)=(U-2)^2(U+2).                    \tag{B.QM.9}
\]
The entire matrix, its cube, and all intermediate polynomials
are in the certificate. The independent checker
\path{code/criterion_quartic13_verify.py} implements its own
prime-field polynomial arithmetic, the cubic field, the
elliptic group law, and (B.QM.7); it then computes the
characteristic polynomials by integer linear algebra reduced
modulo five. It imports no search-producer code.

Pullback and norm split old and new five-primary cohomology,
since thirteen is invertible there. Cartier and coherent
Frobenius are dual after a full finite-field cycle, and the
Artin--Schreier sequence identifies the stable part with
geometric five-torsion cohomology. Polynomial division of
(B.QM.8) by (B.QM.9) gives the new factor
\[
                  (U-2)^{12}(U+2)^{12}.                        \tag{B.QM.10}
\]
It has degree twenty-four and no zero root, proving the
claimed ordinarity. Its value at one is
$(-1)^{12}3^{12}=1$ in $\mathbf F_5$, so the relative
class number is a unit, with the asserted residue.
Both $2$ and $-2$ have multiplicative order four in
$\mathbf F_5$. The eigenvalues of a further $n$th Frobenius
power contain one exactly when $4\mid n$; possible Jordan
blocks do not affect this criterion. This proves the final
assertion.
\end{proof}
\paragraph{Check record.} Also checked in this writing, 2026-11-18: GPT-6 Astra (OpenAI, ChatGPT Pro mode), root; separate effort setting not exposed. Read quartic isogeny pullbacks, all differential bases and Cartier formulas, characteristic polynomials and exact minimality arguments. B.32 remains Open; B.34 first-field minimum is degree 3 with prime 13, distinct from least-prime cubic witness over degree 8. | Also checked in this writing, 2026-11-18: GPT-6 Astra Pro; finite\_checks agent; effort not separately exposed. Read full degree-13 isogeny construction, connected pullback, differential basis, Cartier calculation, rational parity and exact extension-degree criterion. Reran criterion\_quartic13\_verify.


\begin{claim}{B.34}{Proved}{The exact first-field minimum and the pointed Fermat identification}
For the pointed quartic of \statementref{B.33}, the least
constant extension degree over which an odd-prime constant
cyclic $P$-split class-number witness exists is exactly three.
At that first field, thirteen is the least possible prime,
and \statementref{B.33} attains it. The same assertions hold
for the pointed Fermat quartic
\[
       U^4+V^4=R^4\quad\text{over }\mathbf F_5,
                   \qquad (U:V:R)=(1:0:1).
\]
This is a field-first minimum. The least prime when arbitrary
constant extension degrees are allowed is not determined here.
\end{claim}
\begin{proof}
Direct enumeration gives
\[
 \#C(\mathbf F_5)=8,\qquad
 \#C(\mathbf F_{25})=44,\qquad
 \#C(\mathbf F_{125})=104.
\]
Newton identities and the reciprocal functional equation
give the full Jacobian Frobenius polynomial
\[
          H_C(U)=(U^2+2U+5)^2(U^2-2U+5).                       \tag{B.QM.11}
\]
The independent certificate enumerates these points, verifies
the first three Newton identities, and checks the resulting
polynomial. Its consequences for the first three fields are
\[
\begin{split}
 \#J_C(\mathbf F_5)&=8^2\cdot4=256,\\
 \#J_C(\mathbf F_{25})&=32^3=32768,\\
 \#J_C(\mathbf F_{125})&=104^2\cdot148
                         =2^8\cdot13^2\cdot37.
\end{split}                                                    \tag{B.QM.12}
\]
A constant cyclic $P$-split cover of odd prime degree
$\ell\ne5$ gives a quotient $J_C(\mathbf F_{5^m})\twoheadrightarrow
\mathbf Z/\ell$, by the pointed Abel and Lang construction.
The first two orders in (B.QM.12) exclude every such odd prime
over degrees one and two. The third order permits only
thirteen and thirty-seven. The thirteen-cover already proved
therefore establishes both claimed minima.

Finally the characteristic-five identity
\[
       (X+2Z)^4-(X-2Z)^4=X^3Z-XZ^3
\]
shows that the invertible projective linear map
\[
       (X:W:Z)\longmapsto(X-2Z:W:X+2Z)
\]
carries $C$ to the displayed Fermat curve and $P$ to
$(1:0:1)$. Thus the result concerns that exact pointed
Fermat example, not merely another curve with the same
Jacobian point counts.
\end{proof}
\paragraph{Check record.} Also checked in this writing, 2026-11-18: GPT-6 Astra (OpenAI, ChatGPT Pro mode), root; separate effort setting not exposed. Read quartic isogeny pullbacks, all differential bases and Cartier formulas, characteristic polynomials and exact minimality arguments. B.32 remains Open; B.34 first-field minimum is degree 3 with prime 13, distinct from least-prime cubic witness over degree 8. | Also checked in this writing, 2026-11-18: GPT-6 Astra Pro; finite\_checks agent; effort not separately exposed. Read full proof. Checked Newton polynomial, all three Jacobian orders, exclusion of degrees one and two, degree-three least prime 13 and exact pointed Fermat isomorphism. Supplied independent verifier rerun.



\subsection{Supersingular correspondences: a larger proved class}

The general assertion in B5 concerns a finite \'etale cover whose
geometric Galois closure has order prime to $p$. This condition is
stronger than requiring the degree of a particular non-Galois cover
to be prime to $p$. It is retained throughout the constructions below.

\begin{claim}{B.35}{Proved}{B5}
Let $E/\overline{\F}_p$ be an elliptic curve, let
$\pi:E\to E/\{\pm1\}\simeq\mathbf P^1$ be its quotient map, and
let $S\subset\mathbf P^1(\overline{\F}_p)$ be finite. There exists
$\phi\in\mathrm{PGL}_2(\overline{\F}_p)$ such that every point of
$\phi(S)$ is the image under $\pi$ of a prime-to-$p$ torsion point
of $E$. The points of $\phi(S)$ may also be required to avoid the
branch locus of $\pi$.
\end{claim}
\begin{proof}
Choose coordinates putting both $S=\{t_1,\ldots,t_N\}$ and the
branch locus $B$ of $\pi$ in $\mathbf A^1$. This is permissible even
in characteristic two; the quotient map is then separable of degree
two, although ramification may be wild. Choose $c\ne0$ outside the
finite set determined by
\[
 c(t_i-t_j)=b-b'\qquad(i\ne j,\ b,b'\in B).
\]
Thus the finite sets $B-ct_i$ are pairwise disjoint. Descend all this
data to a finite field $k=\F_q$.

If $E$ is supersingular, put $a=0$ and $v=\operatorname{id}_E$.
If $E$ is ordinary, its $p$-adic Tate module has rank one. After
enlarging $k$, write the Frobenius action as multiplication by
$\alpha\in1+p\Z_p$, and put $a=v_p(\alpha-1)<\infty$.
For every $n$ prime to $p$,
\[
 v_p(\alpha^n-1)=a,
\]
because $(\alpha^n-1)/(\alpha-1)\equiv n\pmod p$.
Thus the $p$-primary subgroup of $E(\F_{q^n})$ is killed by $p^a$.
Let $v:E'=E^{(p^a)}\to E$ be the iterated Verschiebung. It is an
\'etale isogeny of degree $p^a$ in the ordinary case, and
\[
 v(E'(\F_{q^n}))=[p^a]E(\F_{q^n})
                 =E(\F_{q^n})\{p'\},\qquad p\nmid n.
\]
Here relative Frobenius is a bijection on points over a perfect
finite field and $v\circ F^a=[p^a]$. In the supersingular case all
geometric torsion already has order prime to $p$.

The composite $h=\pi\circ v:E'\to\mathbf P^1$ is Galois with
group $G=\ker(v)\rtimes\{\pm1\}$ and branch locus $B$.
Its branch inertia groups normally generate $G$: the quotient by
their normal closure would be a connected finite \'etale cover of
$\mathbf P^1$, which must have degree one by Riemann--Hurwitz.
Consider the $N$ covers of the $s$-line given by
\[
 h(R_i)=s+ct_i.
\]
Each is connected and has branch locus $B-ct_i$. Their compositum
has geometric Galois group $H\subseteq G^N$ projecting onto each
factor. At a branch point belonging only to the $i$th factor, its
inertia subgroup is supported in that factor and projects onto the
corresponding inertia of $G$. The other covers are unramified there.
Conjugating by $H$ and using the surjection to the $i$th factor shows
that $H$ contains the normal closure of every such inertia in that
factor. These normal closures generate $G$. Consequently $H$ contains
every factor and equals $G^N$. This also proves geometric
irreducibility of the normalized fibre product; it does not assume
linear disjointness without checking it.

Let $Y/k$ be its smooth proper model. Fix a prime $l\ne p$ and take
$n=l^m$. The Weil bound makes $|Y(\F_{q^n})|$ tend to infinity.
Only finitely many geometric points lie above the branch points or
infinity, so for sufficiently large $m$ there is a point outside that
set. It supplies $s,R_1,\ldots,R_N$ over $\F_{q^n}$ with
$s+ct_i=\pi(v(R_i))$. Each $v(R_i)$ has prime-to-$p$ order by the
previous paragraph. The transformation $t\mapsto s+ct$, composed
with the initial coordinate change of $S$, is the desired $\phi$.
\end{proof}
\paragraph{Check record.} Also checked in this writing, 2026-11-18: GPT-6 Astra (OpenAI, ChatGPT Pro mode), root; separate effort setting not exposed. Read the complete torsion-placement and domination proofs: disjoint branch loci, normal generation, constant-field p-primary control, tame cancellation, and the finite-wreath-product proof of prime-to-p Galois-closure order. Opened Poonen Theorem1.7 and Lemma2.1; the stronger order condition is proved locally.


\paragraph{A finite point-count bound.}
The curve in the proof has degree $d^N$ over the $s$-line, where
$d=2p^a$. Since each individual cover has genus one and their branch
loci are disjoint, the different formula gives
$g(Y)=1+(N-1)d^N$, also when $p=2$. Thus the final existence step
can be made effective by choosing an $l$-power extension for which
the lower Weil bound exceeds the finite number of excluded geometric
points. The $p$-power isogeny in this auxiliary proof is not one of
the \'etale covers of the original curve in the next construction.

\begin{claim}{B.36}{Proved}{B5}
Let $p$ be odd and let $C/\overline{\F}_p$ be hyperelliptic of genus
at least two. For every hyperelliptic curve $H/\overline{\F}_p$ there
is a connected finite \'etale cover $D\to C$ with prime-to-$p$
geometric Galois closure and a nonconstant map $D\to H$.
In particular B5 holds for every such $C$, including characteristic
three, with a supersingular target of genus at least two.
\end{claim}
\begin{proof}
Let $C\to\mathbf P^1$ be the degree-two map, with branch set $T$
of cardinality $2g(C)+2$. Choose four points of $T$ and let
$E\to\mathbf P^1$ be the elliptic double cover branched there.
The normalization $C_0$ of $C\times_{\mathbf P^1}E$ is connected:
the two quadratic function fields are different because their
branch sets are different. By tame ramification cancellation,
$C_0\to C$ is \'etale of degree two. The map $C_0\to E$ is a
double cover ramified above the points lying over the remaining
members of $T$. Translate the group law of $E$ so that one such
branch point is its origin. Then the entire fibre of $C_0\to E$
over zero has ramification index two.

Choose the quotient map $\pi:E\to E/\{\pm1\}$ and write $S$
for the branch set of a hyperelliptic map $H\to\mathbf P^1$.
By B.35, after changing the coordinate on that target, its branch
points are $\pi(P_i)$ for finitely many prime-to-$p$ torsion points
$P_i\in E$. Choose an integer $m$ prime to $p$ killing all $P_i$.
Pull $C_0\to E$ back along $[m]:E\to E$ and normalize, obtaining
$C_1$. The map $C_1\to C_0$ is finite \'etale with group $E[m]$.
The double cover $C_1\to E$ is still connected, since its nonempty
branch divisor is the inverse image of the original branch divisor.
Above every point of $E[m]$ its ramification index is two.

Consequently every point of $C_1$ above a branch value $\pi(P_i)$
has even ramification index in the composite $C_1\to E\to\mathbf P^1$:
the two inverse images under $\pi$ are $P_i,-P_i\in E[m]$.
The normalization of $C_1\times_{\mathbf P^1}H$ is therefore
\'etale over $C_1$, by Abhyankar's lemma. Any connected component
$D$ dominates both $C_1$ and $H$, and has degree one or two over
$C_1$. The tower $D\to C_1\to C_0\to C$ consists of finite
\'etale Galois steps of prime-to-$p$ orders. Its Galois closure has
prime-to-$p$ order: inductively its group is a subgroup of an
extension of the preceding Galois group by a product of conjugate
step groups, equivalently of the corresponding finite wreath
product. This verifies the stronger Galois-closure condition.

For an explicit supersingular target choose an odd prime power
$Q=p^b\ge5$ and let
\[
 H_Q:\quad y^2=x^Q-x.
\]
Its genus is $(Q-1)/2$. For $x\in\F_{Q^2}\setminus\F_Q$, the
nonzero value $x^Q-x$ has trace zero. All such values have the same
quadratic character
$\epsilon=(-1)^{(Q+1)/2}$ in $\F_{Q^2}$: writing a trace-zero
element as $a$, one has $a^{Q-1}=-1$ and hence
$a^{(Q^2-1)/2}=(-1)^{(Q+1)/2}$; elements of $\F_Q^\times$ are
squares in $\F_{Q^2}$. Therefore
\[
 |H_Q(\F_{Q^2})|=Q^2+1+\epsilon(Q^2-Q).
\]
This attains one of the Weil bounds. All Frobenius eigenvalues
over $\F_{Q^2}$ are consequently $-\epsilon Q$, so all Newton
slopes are $1/2$. Thus $H_Q$ is supersingular. For $p=3$, take
$Q=9$, yielding genus four; for $p\ge5$, take $Q=p$.
\end{proof}
\paragraph{Check record.} Also checked in this writing, 2026-11-18: GPT-6 Astra (OpenAI, ChatGPT Pro mode), root; separate effort setting not exposed. Read the complete torsion-placement and domination proofs: disjoint branch loci, normal generation, constant-field p-primary control, tame cancellation, and the finite-wreath-product proof of prime-to-p Galois-closure order. Opened Poonen Theorem1.7 and Lemma2.1; the stronger order condition is proved locally.


The first double-cover construction is the one used in
Bogomolov--Tschinkel's hyperelliptic argument. The extra assertion
proved above is that the finite torsion positioning can be arranged
with prime-to-$p$ torsion, which keeps the final covering tower in
the required formation. We make no priority claim for this
strengthening. The ramification cancellation is the precise
one-dimensional form of SGA1 XIII Proposition~5.2 and Poonen
Lemma~2.1.\citep{LA_BogomolovTschinkel2002,LA_Poonen2005,LA_SGA1}

\begin{claim}{B.37}{Proved}{B5}
Let $p=2$. Suppose a curve $C/\overline{\F}_2$ admits a finite
separable map to $\mathbf P^1$ with three distinct fibres in which
every ramification index is divisible by three. Then a connected
finite \'etale cover of $C$ of degree one or three maps nonconstantly
to a supersingular elliptic curve. In particular this applies to
every cyclic degree-three cover of $\mathbf P^1$ of genus at least
two.
\end{claim}
\begin{proof}
Move the three values to $0,1,\infty$ and use the cyclic cover
$E:y^3=x(x-1)$. Its smooth projective model is elliptic, branched
with index three at those values. Normalized base change along
$C\to\mathbf P^1$ is \'etale by Abhyankar's lemma, and any
connected component has cyclic degree one or three over $C$.
After interchanging the coordinates, the elliptic curve is
$z^2+z=t^3$. It has nine points over $\F_4$: for each of the four
values of $t$, the right side has trace zero to $\F_2$, so there
are two values of $z$, in addition to the point at infinity.
Thus its two $\F_4$-Frobenius eigenvalues are $-2$ and it is
supersingular. A cyclic cubic cover has uniform ramification
index three over each branch point, and has at least three such
points when its genus is at least two.
\end{proof}
\paragraph{Check record.} Also checked in this writing, 2026-11-18: GPT-6 Astra (OpenAI, ChatGPT Pro mode), root; separate effort setting not exposed. Read the complete torsion-placement and domination proofs: disjoint branch loci, normal generation, constant-field p-primary control, tame cancellation, and the finite-wreath-product proof of prime-to-p Galois-closure order. Opened Poonen Theorem1.7 and Lemma2.1; the stronger order condition is proved locally.


\begin{claim}{B.38}{Proved}{B5}
Suppose $p$ is odd, $C/\overline{\F}_p$ has genus at least two,
and $G\subset\Aut(C)$ has order prime to $p$. Put $Y=C/G$ and
suppose $g(Y)\le2$. If at least one of the following holds, then
$C$ has a connected finite \'etale cover with prime-to-$p$ Galois
closure mapping to a supersingular curve of genus at least two:
\begin{enumerate}
\item $g(Y)=1$ or $2$;
\item $G$ is solvable;
\item two distinct branch values of $C\to Y$ have ramification
indices with a common prime factor.
\end{enumerate}
In particular this proves B5 for every tamely superelliptic curve
in every odd characteristic, including characteristic three.
\end{claim}
\begin{proof}
We verify the extra prime-to-$p$ condition in Poonen's
Theorem~1.7, cases~(1)--(3), rather than infer it from the word
``unramified.'' His theorem assumes precisely a prime-to-$p$
Galois action, genus at least two for the source, and quotient
genus at most two, with one of the listed extra hypotheses.
The proof below follows his \S\S3.1--3.3 and records the orders
of the covering groups.\citep[Theorem~1.7 and \S3]{LA_Poonen2005}

First suppose $Y$ has genus two. It is already hyperelliptic,
so no \'etale covering is needed to dominate a hyperelliptic curve.
Suppose instead that $Y=E$ is elliptic. The map $C\to E$ is
ramified, since $g(C)>1$. Translate a branch value to the origin,
and choose a prime $e\ne p$ dividing its ramification index.
The Galois hypothesis ensures that this divisibility holds at
every point above the origin. Pull back along $[l]:E\to E$ for
an odd prime $l\ge5$ distinct from $p$, and choose a connected
component $C'$. Then $C'\to C$ is finite \'etale Galois with
group a subgroup of $E[l]$. Every point of $C'$ above $E[l]$
has ramification index divisible by $e$ in $C'\to E$.

Under $E\to E/\{\pm1\}=\mathbf P^1$, the nonzero $l$-torsion
points give $(l^2-1)/2$ distinct values $a_i$. If $e\ge3$, take
the cyclic cover
\[
 H:\quad z^e=\frac{(x-a_1)(x-a_2)}{(x-a_3)(x-a_4)}.
\]
It has genus $e-1\ge2$, and its $z$-map has degree two, so it
is hyperelliptic. If $e=2$, use three factors in the numerator
and three in the denominator, giving a genus-two hyperelliptic
curve instead. At every branch value of $H\to\mathbf P^1$,
all ramification indices of $C'\to\mathbf P^1$ are divisible
by $e$. A connected component of the normalized pullback of
$H$ is therefore \'etale over $C'$, with cyclic group of order
one or $e$, and dominates $H$.

Now let $Y=\mathbf P^1$ and $G$ be solvable. Induct on $|G|$.
For a nontrivial proper normal subgroup $N$, if $g(C/N)\ge2$
apply induction to $C/N\to\mathbf P^1$ and pull its \'etale
cover back to $C$. If $g(C/N)=1$, apply the elliptic case to
$C\to C/N$. If $g(C/N)=0$, apply induction to $C\to C/N$.
It remains to treat $G$ cyclic of prime order $e\ne p$.
For $e=2$, $C$ itself is hyperelliptic. For $e\ge3$, pull
back along a cyclic degree-$e$ map $\mathbf P^1\to\mathbf P^1$
branched at two branch values of $C$. The two degree-$e$
extensions are different (the original curve has positive genus),
so the pullback is connected and gives an \'etale cyclic
degree-$e$ cover $C'\to C$. Its new cyclic map to $\mathbf P^1$
has at least four branch values, since its genus has increased.
Using four of these values in the preceding displayed equation
again gives an \'etale cyclic pullback dominating a hyperelliptic
curve.

Finally suppose two branch indices have a common prime divisor
$e$. Pull back along a cyclic degree-$e$ map of projective lines
branched at those values. A third branch value has at least two
inverse images; a prime factor $e'$ of its ramification index
therefore occurs at two new branch values. Repeat this construction.
Every step on the source is cyclic \'etale of order dividing the
chosen prime, hence prime to $p$. Every new map to the projective
line is Galois of degree at most $|G|$, with ramification indices
dividing $|G|$. The degrees of the composite maps of projective
lines tend to infinity, forcing the degrees of the covers of $C$
and hence their genera to tend to infinity. Riemann--Hurwitz then
forces the number of branch values to tend to infinity. There are
only finitely many possible indices, so at some stage at least six
branch values have a common index $a>1$. Choose a prime $e\mid a$
and a cyclic degree-$e$ cover of $\mathbf P^1$ branched at these
six points, using three exponents $1$ and three exponents $-1$.
It has genus $2(e-1)\ge2$. Its normalized pullback is \'etale
over the stage curve, and its solvable cyclic group puts the target
in the preceding case.

Thus in every case a hyperelliptic curve is dominated after a
finite tower of \'etale Galois prime-to-$p$ covers and their base
changes. A connected base change of such a Galois cover has a
subgroup of the original group; finite composition has prime-to-$p$
Galois closure by the wreath-product argument in B.36. Composing
with B.36 (and pulling its cover back) gives the claimed
supersingular target.
\end{proof}
\paragraph{Check record.} Also checked in this writing, 2026-11-18: GPT-6 Astra (OpenAI, ChatGPT Pro mode), root; separate effort setting not exposed. Read the complete torsion-placement and domination proofs: disjoint branch loci, normal generation, constant-field p-primary control, tame cancellation, and the finite-wreath-product proof of prime-to-p Galois-closure order. Opened Poonen Theorem1.7 and Lemma2.1; the stronger order condition is proved locally.


\paragraph{The modular fourth case, with its group condition.}
Poonen's fourth condition asks for three branch indices divisible
by $2,3,l$, respectively, where $l\le89$ is prime or
$l\in\{101,103,107,131,167,191\}$. For $l=2,3$ the preceding
common-factor case applies. For $l\ge5$ his proof uses
$X(l)\to\mathbf P^1$ with geometric group
$\mathrm{PSL}_2(\F_l)$, then solvable covers of $X_0(l)$ and the
Atkin--Lehner quotient $X_0^*(l)$ of genus at most two.
Consequently the same prime-to-$p$ conclusion follows if additionally
$p\nmid l(l^2-1)$: the modular pullback, its Borel-subgroup steps,
and the order-two Atkin--Lehner step all have prime-to-$p$ Galois
closure. In the $l=5$ case, the new Galois map to $X(5)\simeq
\mathbf P^1$ has repeated branch indices and the third case applies.
Without the additional group-order check, tameness of the local
indices $2,3,l$ alone is insufficient to certify the required
Galois closure.\citep[Lemma~2.3 and \S3.4]{LA_Poonen2005}

\paragraph{The general reduction still required.}
B.36 also applies whenever a prime-to-$p$ \'etale cover of the
given odd-characteristic curve maps to a hyperelliptic curve.
B.38 establishes this under the listed low-genus Galois
hypotheses, but a general curve is not assumed to have such a map.

For a general curve, a tame map to $\mathbf P^1$ does not suffice.
Its ramification over a branch value may include unramified points,
where a branched target cover would remain ramified after base
change. Nor need the Galois closure of a tame map have
prime-to-$p$ order. Thus the tame Belyi theorem does not provide
the missing divisible fibres or the missing prime-to-$p$ covering
tower. The general universal-correspondence assertion is explicitly
Conjecture~2.6 in Bogomolov--Tschinkel; the theorem there is for
hyperelliptic sources and $p\ge5$. Neither that conjecture nor the
present special cases settle B5 for every nonhyperelliptic curve
or for all curves in characteristic two.
\citep[Theorem~1.7 and \S3]{LA_Poonen2005}
\citep[Theorem~1.7 and Conjecture~2.6]{LA_BogomolovTschinkel2002}
\citep[Theorem~1.5]{LA_KedlayaLittWitaszek2023}

The ordinarity literature does not supply the missing conclusion.
Raynaud's Theorem~2 constructs, for a proper smooth curve of genus
at least two in characteristic $p$, an \'etale cover with a solvable
prime-to-$p$ group which is nonordinary. Nonordinarity is weaker than
having a supersingular isogeny factor, and still weaker than the
required map to a supersingular curve. Thus this result cannot be
substituted for B5. The Prym results of Ozman--Pries likewise give
$p$-rank information, with no universal supersingular-target assertion.
\citep[Theorem~2]{LA_Raynaud2000}
\citep[Theorem~1.1 and \S8.2]{LA_OzmanPries2019}


\subsection{A positive-rank primitive-zero cover without a supersingular factor}

\begin{claim}{B.39}{Proved}{An exact cyclic-thirteen distinction between B4 and B5}
Let $C/\F_2$ be the smooth proper model
\[
 C:\quad Y^2+X^2Y=X^5+X,
 \qquad P=\infty.
\]
It has genus two, $2$-rank one, and no nonzero supersingular
isogeny factor in its geometric Jacobian. Over $k=\F_{64}$
there is a connected $P$-pointed cyclic \'etale cover
$D\to C_k$ of degree $13$, with zeta numerator
\[
 P_D(T)=
 (1-T-80T^2-64T^3+4096T^4)
 (1+12T+64T^2)^{12}.
 \tag{B.LIT.7}
\]
Its primitive degree-$13$ factor has Newton slopes
$1/3$ and $2/3$, each with multiplicity $12$, and hence
has $2$-rank zero. Its relative class number is
\[
 \frac{|J_D(k)|}{|J_C(k)|}=77^{12}
       =43439888521963583647921,
\]
which is odd. Nevertheless $D$ has no nonconstant map,
even over $\overline{\F}_2$, to any supersingular curve
of positive genus. Thus this is a B1 witness and a
nontrivial positive-$2$-rank case of B4, but this particular
cover does not supply the correspondence asked for in B5.
\end{claim}

\begin{proof}
The substitution $x=X^{-1}$, $y=Y/X^2$ identifies the
function field with
\[
 y^2+y=x^3+x^{-1}.
\]
The Artin--Schreier poles have odd orders three and one.
Their different exponents are four and two, so
Riemann--Hurwitz gives genus two. In the polynomial model
the regular differential $dX/X^2$ has divisor $2P$.
Direct counts give $|C(\F_2)|=4$ and $|C(\F_4)|=4$;
the genus-two functional equation therefore gives the
Frobenius characteristic polynomial
\[
 H_2(U)=U^4+U^3+2U+4.
\]
It is irreducible modulo three and hence over $\Q$.
Its Newton slopes are $0,1/2,1/2,1$, which also verifies
the asserted $2$-rank. Taking sixth powers of its roots
gives
\[
 H_{64}(U)=U^4-U^3-80U^2-64U+4096,
 \qquad |J_C(k)|=H_{64}(1)=3952=16\cdot13\cdot19.
\]

Here is a complete definition of the cover used in the
computation. Write $k=\F_2[z]/(z^6+z+1)$ and encode
elements by the binary coefficients in the basis
$1,z,\ldots,z^5$. In reduced Mumford coordinates for the
displayed polynomial model, put
\[
 G=\bigl(X^2+25X+14,\;17X+26\bigr)\in J_C(k).
\]
The numbers in this formula are binary field encodings,
not integers embedded in a characteristic-two field.
The certificate verifies $G\ne0$ and $13G=0$.
Since the $13$-primary subgroup has order $13$, there is
a unique surjection
\[
 \chi:J_C(k)\longrightarrow\Z/13,
 \qquad [304]A=\chi(A)G.
\]
Pull back the Lang isogeny $F_{64}-1:J_C\to J_C$
along $R\mapsto[R-P]$ and quotient its deck group
$J_C(k)$ by $\ker\chi$. This defines $D$. The point
above $P$ represented by zero makes it pointed. The
Abel map induces the isomorphism from the abelianized
geometric fundamental group of $C$ to the fundamental
group of $J_C$; thus the pullback of this connected
cyclic isogeny cover is geometrically connected.
This is the precise Jacobian covering statement of
Milne, \S9, Proposition~9.1, applied after base change
to $\overline{k}$: that proposition assumes a separably
closed field, a smooth proper curve, and a chosen point.
\citep[\S9, Proposition~9.1]{MilneJac}

For every $R\in C(k)$ the program computes
$\chi([R-P])$ by exact Cantor arithmetic. Its complete
table, including every point and every element of
$\langle G\rangle$, gives
\[
 \#\{R:\chi([R-P])=0\}=16,
 \qquad
 \#\{R:\chi([R-P])=j\}=4
 \quad(1\le j\le12).
\]
Every polynomial division in the group law is checked
to be exact, and every output Mumford pair is checked
against the curve equation. The producer
\path{code/b5_char2_13_certificate.py} writes these finite
tables and the entire expanded degree-$28$ numerator to
\path{code/b5_char2_13_certificate.json}; its textual
output is also retained.

The independent verifier
\path{code/c0_cyclic13_verify.py} imports no Cantor
implementation. It uses fixed-point addition by the
Chinese remainder theorem and tangent Taylor expansion,
and independently reproduces all $64$ point projections,
all $13$ subgroup elements, and $H_{64}$ by a Newton-sum
recurrence. Its exact output and certificate are retained
beside the producer.

Let $\zeta$ be a primitive thirteenth root of unity.
The complete twisted polynomial has degree two. Its
linear coefficient, for every $i\ne0$, is
\[
 \sum_{R\in C(k)}\zeta^{i\chi([R-P])}
       =16+4\sum_{j=1}^{12}\zeta^{ij}=12.
\]
Because the canonical divisor is $2P$, its quadratic
coefficient is $64$. These are exactly the effective-divisor
and functional-equation formulas proved in B.2, so no
single evaluation at $T=1$ is substituted for a Frobenius
polynomial. Multiplying the twelve complete twisted
polynomials and the numerator of $C_k$ proves (B.LIT.7).
It also gives the relative class number and
$g(D)=13(g(C)-1)+1=14$.

The two roots of $U^2+12U+64$ have $2$-adic valuations
two and four: the point $(1,2)$ lies strictly below the
segment joining $(0,6)$ and $(2,0)$. Dividing valuations
by six gives slopes $1/3$ and $2/3$. Thus the primitive
factor has $2$-rank zero and no supersingular factor.
For completeness, $J_C$ has no geometric supersingular
factor either. If a Frobenius root $\alpha$ belonged to
one after a finite constant extension, then
$\alpha/\sqrt2$ would be a root of unity. One elementary
way to see this criterion is that a supersingular Weil
number divided by the appropriate square root of the
field size is an algebraic unit at every finite place
and has absolute value one at every complex place;
Kronecker's argument then makes it a root of unity.
The property descends from $\alpha^m$ to $\alpha$.
Irreducibility of $H_2$ would make all four of its roots
have this property, contradicting its slope-zero root.
This also proves geometric simplicity here: a geometric
elliptic decomposition of a $2$-rank-one surface would
include one supersingular elliptic factor.

Finally $J_D$ is isogenous to $J_C\times\Pr(D/C)$
by pullback and norm. A nonconstant map $D\to S$ of
smooth proper curves induces pullback and norm with
composition multiplication by its degree on $J_S$.
Consequently a positive-dimensional supersingular
$J_S$ would give a supersingular factor of $J_D$,
which the preceding two paragraphs exclude. This holds
for inseparable as well as separable maps, since a
nonzero multiplication map is invertible in the
isogeny category.
\end{proof}
\paragraph{Check record.} Also checked in this writing, 2026-11-18: GPT-6 Astra (OpenAI, ChatGPT Pro mode), root; separate effort setting not exposed. Read complete proofs: cyclic13 Frobenius factor and geometric non-supersingularity; theta-divisor exclusion and characteristic-two reflection Prym; ramification gcd obstruction; Hermitian/Fermat supersingular pullback. Finite certificates are separate events; the g=2, 3 restriction in B40 is essential. | Also checked in this writing, 2026-11-18: GPT-6 Astra Pro; finite\_checks agent; effort not separately exposed. Read full cyclic-13 construction and proof. Checked all Artin–Schreier and Newton-slope arguments, geometric supersingular-factor exclusion and map obstruction. Reran producer and independent c0\_cyclic13\_verify.


\paragraph{The scope of the search.}
This does not refute the existence of a different odd
\'etale cover of $C$ with a supersingular factor. The
independent theta-Taylor elimination in
\path{code/b5_theta_division.py} excludes torsion orders
$3,5,7,9,11$ on the Raynaud theta curve of this example;
the order-$13$ candidate locus is $r^2+r+1=0$.
The companion probe records complete character sums
for exponent-one prime quotients over all
$\F_{2^m}$ with $1\le m\le10$. These bounded searches
are retained as evidence about the attempted route,
not as a classification of all odd covers or all
constant extensions.

\subsection{Two exact limits of the proposed correspondence constructions}

\begin{claim}{B.40}{Proved}{Ordinary cubic Pryms in two small genera}
Let $C$ be a smooth hyperelliptic curve over an algebraically
closed field of characteristic two, of genus $g\in\{2,3\}$
and $2$-rank zero. For every connected cyclic \'etale
cover $D\to C$ of degree three, the Prym variety is
ordinary of dimension $2(g-1)$. In particular,
$f_D=2(g-1)$. In genus three its genus-two reflection
quotient is ordinary, so this cubic-reflection construction
does not yield a supersingular target.
\end{claim}

\begin{proof}
The hyperelliptic Artin--Schreier map has exactly one
branch point, since its number of branch points is
$f_C+1=1$. Put this point at infinity, denote its
unique preimage by $P$, and use an equation
$y^2+y=f(x)$ with $\deg f=2g+1$. Then
\[
 \operatorname{div}(dx)=(2g-2)P,
 \qquad B^1\simeq\mathcal O_C((g-1)P),
\]
where $B^1$ is Raynaud's line bundle of locally exact
differentials on the Frobenius twist. We identify that
twist with its corresponding equation over the perfect
ground field; all arguments below are unchanged on it.

Suppose a nonzero three-torsion class $L$ satisfies
$H^0(B^1\otimes L)\ne0$. Removing copies of $P$ and
hyperelliptic conjugate pairs from an effective divisor
of degree $g-1$ gives
\[
 L\simeq\mathcal O_C(E-dP),\qquad 1\le d\le g-1,
\]
with $E$ reduced in the hyperelliptic divisor sense:
it contains no $P$ and no conjugate pair. There is a
rational function $h$ with
$\operatorname{div}(h)=3E-3dP$. For $g=2$ or $3$,
\[
 3d\le3(g-1)\le2g<2g+1.
\]
The pole orders of $x,y$ at $P$ are $2,2g+1$.
Consequently every function whose only pole is at
$P$ of order at most $2g$ belongs to $k[x]$.
Every finite zero of such a function occurs together
with its hyperelliptic conjugate and with the same
multiplicity; there are no finite branch points.
This contradicts the form of $E$. Thus neither $L$
nor $L^2$ lies on the theta divisor of $B^1$.

For clarity, the ordinarity implication uses the exact
sequence
\[
 0\longrightarrow L\longrightarrow F_*F^*L
   \longrightarrow B^1\otimes L\longrightarrow0.
\]
Since $H^0(L)=H^0(L^2)=0$, the just-proved vanishing
makes Frobenius an isomorphism between the two
$g-1$ dimensional twisted $H^1(\mathcal O)$ spaces.
These are precisely the nontrivial character summands
of $H^1(D,\mathcal O_D)$, as
$\pi_*\mathcal O_D=\mathcal O_C\oplus L\oplus L^2$.
Their stable Frobenius rank is $2(g-1)$, which proves
the assertion. This is also the rank-one instance of
Raynaud's ordinarity criterion in \S1.
\citep[\S1, before Proposition~1]{LA_Raynaud2000}

Here are the details of the reflection assertion in
characteristic two. Let $\delta$ generate the cyclic
deck group, and let $\iota$ be the hyperelliptic
involution on $C$. Since $\iota^*L=L^{-1}$,
$\iota$ lifts to an automorphism $\tau$ of $D$ with
$\tau\delta\tau^{-1}=\delta^{-1}$.
Write $\tau^2=\delta^a$. The element $\tau$ commutes
with its square and also sends $\delta^a$ to
$\delta^{-a}$, so $a=-a$ in $\Z/3$ and $a=0$.
Thus $\tau$ is an involution. It fixes exactly one
of the three points over $P$, since it acts there
as a reflection of the three-element deck torsor.
It has no other fixed points, as $\iota$ has none.
The \'etale map $D\to C$ identifies the completed
local rings at that fixed point and at $P$.
The local different for $D\to H:=D/\langle\tau\rangle$
therefore has exponent $2g+2$, the same as that of
$C\to\mathbb P^1$. Consequently
\[
 6g-6=2g_D-2=2(2g_H-2)+(2g+2),
 \qquad g_H=g-1.
\]

Let $A,A'$ be the pullback images in $J_D$ of the
Jacobians of the quotients by $\tau$ and by a distinct
conjugate reflection. Both lie in the Prym: their
norm to $J_C$ is fixed by the induced hyperelliptic
action $[-1]$, hence has finite image and is zero
on these connected abelian varieties. Their
intersection is finite. Indeed an element fixed by
the two reflections is fixed by $\delta$ as well.
Its norm $z$ to $J_C$ satisfies $2z=0$, while the
pullback of $z$ equals $3$ times that element; hence
the intersection is killed by six. Pullback from
each quotient has finite kernel by its pullback--norm
identity. Both images thus have dimension $g-1$,
and their addition gives an isogeny
\[
 J_H\times J_H\longrightarrow\Pr(D/C).
\]
For $g=3$ the quotient therefore has genus two and
is ordinary.
The inequality used above fails for $g\ge4$; no
assertion in those genera follows from this proof.
\end{proof}
\paragraph{Check record.} Also checked in this writing, 2026-11-18: GPT-6 Astra (OpenAI, ChatGPT Pro mode), root; separate effort setting not exposed. Read complete proofs: cyclic13 Frobenius factor and geometric non-supersingularity; theta-divisor exclusion and characteristic-two reflection Prym; ramification gcd obstruction; Hermitian/Fermat supersingular pullback. Finite certificates are separate events; the g=2, 3 restriction in B40 is essential.


\begin{claim}{B.41}{Proved}{A fixed simply branched function cannot supply the missing fibres}
Let $k$ be algebraically closed, let $f:C\to\mathbb P^1_k$
be finite separable, and assume that for every
$t\in\mathbb P^1(k)$,
\[
 \gcd\{e_P(f):P\in f^{-1}(t)\}=1.
\]
Let $C'\to C$ be any finite \'etale cover and
$h:\mathbb P^1\to\mathbb P^1$ any finite separable map.
The composite $h\circ f\circ\pi$ cannot have three
distinct fibres in each of which all ramification
indices are divisible by three. It cannot have four
distinct fibres in each of which they are all even.
More generally, uniform fibre divisibilities $m_i$
at distinct values require
\[
 \sum_i(1-1/m_i)<2.
\]
\end{claim}

\begin{proof}
The \'etale pullback preserves every ramification index
over the algebraically closed ground field. For a
fixed $t$ over $u$ the gcd of the indices of the
composite at points over $t$ is therefore $e_t(h)$.
Taking the gcd over all $t\in h^{-1}(u)$ shows that
any uniform divisor $m_i$ of the composite fibre
already divides every $e_t(h)$ in that fibre of $h$.
Write $n=\deg h$. There are at most $n/m_i$ geometric
points in such a fibre. Their contribution to the
different is at least
$\sum_{t\mid u_i}(e_t(h)-1)\ge n(1-1/m_i)$,
also in the presence of wild ramification.
Riemann--Hurwitz for $h:\mathbb P^1\to\mathbb P^1$
now gives
\[
 n\sum_i(1-1/m_i)\le2n-2<2n.
\]
The claimed patterns each give a sum of two and are
impossible. In particular, the hypothesis holds for
a simply branched map of degree at least three:
every branch fibre has an unramified point. Thus
rational postcomposition and \'etale pullback of that
fixed function cannot complete the general B5 route;
a new function on the covering curve would be needed.
\end{proof}
\paragraph{Check record.} Also checked in this writing, 2026-11-18: GPT-6 Astra (OpenAI, ChatGPT Pro mode), root; separate effort setting not exposed. Read complete proofs: cyclic13 Frobenius factor and geometric non-supersingularity; theta-divisor exclusion and characteristic-two reflection Prym; ramification gcd obstruction; Hermitian/Fermat supersingular pullback. Finite certificates are separate events; the g=2, 3 restriction in B40 is essential.


\begin{claim}{B.42}{Proved}{Three uniformly divisible fibres with a supersingular Fermat quotient}
Let $C/\overline{\F}_p$ be a smooth proper curve admitting
a finite separable map to $\mathbb P^1$ with three
distinct fibres in which every ramification index is
divisible by the same odd integer $r\ge3$.
Suppose $r\mid p^a+1$ for some $a\ge1$.
Then $C$ has a connected cyclic \'etale cover of degree
dividing $r$ which dominates a supersingular curve of
genus $(r-1)/2$. Its geometric Galois closure therefore
has order prime to $p$.
In characteristic two this applies to every cyclic
prime-degree cover of $\mathbb P^1$ of genus at least
two whose degree $\ell$ satisfies that
$\operatorname{ord}_\ell(2)$ is even.
\end{claim}

\begin{proof}
Send the three specified values to $0,1,\infty$.
The smooth proper model
\[
 H_r:\quad y^r=x(x-1)
\]
is a cyclic cover of degree $r$, totally ramified
at those three values and nowhere else. Since $r$
is odd, its genus is $(r-1)/2$. Normalizing
$C\times_{\mathbb P^1}H_r$ and taking any connected
component gives a cyclic cover $D\to C$ of degree
dividing $r$. At each of the three relevant fibres,
the tame local ramification index $r$ divides that
of $C\to\mathbb P^1$, so the normalized pullback is
unramified. It is unramified elsewhere as well.
Thus $D\to C$ is \'etale, and $D\to H_r$ is dominant.

Put $q=p^a$. The Fermat curve of degree $r$ maps
to $H_r$ by
\[
 U^r+V^r=1,\qquad x=U^r,\quad y=-UV.
\]
It is in turn dominated by the Hermitian Fermat
curve $U^{q+1}+V^{q+1}=1$, by taking coordinate
powers $(q+1)/r$. This latter curve has genus
$q(q-1)/2$ and exactly $q^3+1$ points over
$\F_{q^2}$. Indeed the norm map to $\F_q$ has
one preimage of zero and $q+1$ preimages of each
nonzero value. Counting affine solutions gives
$2(q+1)+(q-2)(q+1)^2=q^3-q$, and there are
$q+1$ points at infinity. This attains the upper
Weil bound, so every Frobenius root equals $-q$.
The curve is therefore supersingular. Pullback
and norm inject the Jacobian of every dominated
curve in the isogeny category; hence $H_r$ is
supersingular too.

For the last assertion, a cyclic cover of odd prime
degree $\ell$ has all ramification indices equal
to $\ell$ at its branch values. Genus at least two
forces at least three such values. An even order
of $2$ modulo $\ell$ gives
$2^{\operatorname{ord}_\ell(2)/2}=-1$ modulo
$\ell$, so the proved criterion applies. Degrees
$5,11,13,17,19$ give examples beyond the cubic
case. Degree $7$ has order three and is outside
this criterion; no inference for arbitrary
characteristic-two curves follows.
\end{proof}
\paragraph{Check record.} Also checked in this writing, 2026-11-18: GPT-6 Astra (OpenAI, ChatGPT Pro mode), root; separate effort setting not exposed. Read complete proofs: cyclic13 Frobenius factor and geometric non-supersingularity; theta-divisor exclusion and characteristic-two reflection Prym; ramification gcd obstruction; Hermitian/Fermat supersingular pullback. Finite certificates are separate events; the g=2, 3 restriction in B40 is essential.



% Configuration audit contribution to Part B. AI-produced by GPT-6 Astra.
% Provisional statement identifiers B.CF*; parent editor renumbers.
\subsection{Ordinary pro-reconstruction in finite transcendence degree}

We use an ordinary category here, not a category in which coherence is an
additional hypothesis. Let \(\mathsf T\) be the category whose objects are
Grothendieck topoi and whose morphisms are isomorphism classes of geometric
morphisms. Its ordinary pro-category has
\[
 \operatorname{Hom}_{\operatorname{Pro}(\mathsf T)}
       (\{\mathcal X_i\},\{\mathcal Y_j\})
   =\varprojlim_j\varinjlim_i
       \operatorname{Hom}_{\mathsf T}(\mathcal X_i,\mathcal Y_j).
                                                        \tag{B.CF.1}
\]
Fix a field \(k\) of characteristic zero and models of finite-type
\(k\)-schemes \(X,Y\) over a finitely generated subfield \(F_0\).
For finitely generated \(F\subset k\) containing \(F_0\), put
\[
 \mathcal X=\{X_{F,\mathrm{et}}\}_F,
 \qquad \mathcal Y=\{Y_{F,\mathrm{et}}\}_F,
 \qquad \mathcal B=\{\operatorname{Spec}(F)_{\mathrm{et}}\}_F.
\]
These are pro-objects of \(\mathsf T\), with structural arrows to
\(\mathcal B\). A morphism in the slice
\(\operatorname{Pro}(\mathsf T)_{/\mathcal B}\) induces a map
\(|X|\to |Y|\): apply the underlying-space functor and take inverse
limits. The identification with the scheme spaces follows from
\citep[Lemma~32.4.2, Tag~0CUF]{StacksProject}; its hypotheses are a directed
inverse system of schemes with affine transition maps, as here.
We call this pro-morphism \emph{pinned} when that induced map sends
closed points to closed points. This definition does not impose any
same-field representatives or coherence data.

\begin{claim}{B.43}{Proved}{Ordinary pro-reconstruction for every field of finite transcendence degree}
Suppose \(\operatorname{trdeg}_{\Q}k<\infty\). For finite-type
\(k\)-schemes, the functor to the preceding pro-systems is fully faithful
after universal homeomorphisms are inverted, with pinned morphisms in
\(\operatorname{Pro}(\mathsf T)_{/\mathcal B}\) as target morphisms.
Equivalently, for absolutely weakly normal finite-type \(X,Y\), every
such pinned ordinary pro-morphism is induced by a unique \(k\)-morphism
\(X\to Y\). This includes all finite-transcendence-degree fields of B8,
regardless of whether reconstruction from the single topos over \(k\)
holds.
\end{claim}
\begin{proof}
Enlarge \(F_0\) to contain a transcendence basis of \(k/\Q\).
Then \(k/F_0\) is algebraic. Every finitely generated subfield of \(k\)
is contained in one containing this \(F_0\), and every such field is
finite over \(F_0\). Thus we have made a cofinal change of the indexing
system, and all its field transitions are finite separable. In
particular \(X_E\) and \(Y_F\), for \(E\supset F\supset F_0\),
are finite-type schemes over the same finitely generated field \(F\).

We first verify that forgetting structural 2-cells causes no ambiguity
at these levels. For every finitely generated characteristic-zero
field \(F\), its absolute Galois group \(G_F\) is slim, meaning that
all its open subgroups have trivial centre. Indeed, if
\(\operatorname{trdeg}_{\Q}F>0\), apply
\citep[Corollary~4.5(1),(3)]{HoshiTsujimura2023}; for a number field apply
Corollary~4.5(3) using any of its discrete mixed-characteristic valuation
rings. The same source, Definition~4.3, states the equivalent condition
that the centralizer in \(G_F\) of every normal open subgroup is trivial.
It follows also for every open subgroup, by taking its normal core.
For each connected component \(Z\) of a nonempty finite-type
\(F\)-scheme, the image of \(\pi_1(Z)\to G_F\) is open: it is the
stabilizer of a geometric connected component, of which there are
finitely many. An automorphism of the structural geometric morphism
\(Z_{\mathrm{et}}\to BG_F\) is an element of the centralizer of this
image. This follows by evaluating its action on every finite
\(G_F\)-set at a geometric base point; naturality and transport require
and suffice for commutation with the fundamental-group image.
Consequently that automorphism group is trivial. The same conclusion
holds component by component. Therefore, whenever the structural
triangle commutes as isomorphism classes, its commutation 2-cell is
unique. Moreover, any absolute isomorphism between two geometric
morphisms over \(BG_F\) automatically respects their structural 2-cells.

Let \(\alpha:\mathcal X\to\mathcal Y\) be a pinned pro-morphism over
\(\mathcal B\). At a target level \(F\), formula~(B.CF.1) represents
it by a geometric morphism
\[
             a:X_{E,\mathrm{et}}\longrightarrow Y_{F,\mathrm{et}}
                                                        \tag{B.CF.2}
\]
for some \(E\supset F\). The equality of structural pro-morphisms,
applied at the target level \(BG_F\), means that after enlarging \(E\)
once more, (B.CF.2) is a morphism over \(BG_F\). This is precisely the
equality criterion in the filtered colimit in (B.CF.1); it uses no
uniform choice over all target levels.

The representative (B.CF.2) is pinned. To prove this, let \(x\) be a
closed point of \(X_E\). The nonzero finite \(k\)-algebra
\(\kappa(x)\otimes_E k\) has a closed point, so \(x\) lifts to a
closed point \(\widetilde x\) of \(X\). The point
\(\widetilde y=|\alpha|(\widetilde x)\) is closed by assumption.
Its image in \(Y_F\) is closed: its residue field is generated by
finitely many coordinates over \(F\), is contained in a finite
extension of \(k\), and is algebraic over \(F\), since \(k/F\) is
algebraic. It is therefore finite over \(F\). Compatibility of the
underlying-space pro-morphism identifies that image with \(|a|(x)\).
This proves pinnedness at the actual finite stage.

We now apply the opened theorem
\citep[Theorem~0.1 and the equivalent formulation~(0.2)]{CarlsonHaineWolf2024}:
for a field finitely generated over its prime field, with the additional
hypothesis of positive transcendence degree only in positive
characteristic, the functor from finite-type schemes localized at
universal homeomorphisms to pinned relative geometric morphisms is
fully faithful. The field here is \(F\), of characteristic zero;
both schemes in (B.CF.2) are finite type over it. It follows that
(B.CF.2) comes from a unique localized \(F\)-scheme morphism
\(g_F:X_E\to Y_F\). Pairing this morphism with the structural map
to \(\operatorname{Spec}E\) gives a localized \(E\)-scheme morphism
\(X_E\to Y_E\), and then a localized \(k\)-scheme morphism
\(f_F:X\to Y\).

Choose this construction first at \(F_0\), and write \(f=f_{F_0}\).
For another target field \(F\), compatibility of the pro-morphism
means that \(g_F\) followed by \(Y_F\to Y_{F_0}\), and \(g_{F_0}\),
induce isomorphic geometric morphisms after a common finite source
extension \(H\subset k\). They are morphisms of finite-type
\(F_0\)-schemes. The uniqueness of structural 2-cells just proved
makes this a relative isomorphism; the faithfulness in Theorem~0.1
therefore makes the two localized scheme morphisms equal.
Their lifts to \(Y_F=Y_{F_0}\times_{F_0}F\) agree as well, since they
also have the same structural morphism to \(\operatorname{Spec}F\).
It follows that \(f_F=f\), and the pro-morphism induced by \(f\)
is exactly \(\alpha\), at every target level.

Conversely a localized \(k\)-scheme morphism descends to some finite
extension of \(F_0\): a representative finite diagram, its equations,
and its inverse universal homeomorphisms involve only finitely many
coefficients. Being a universal homeomorphism descends under the
faithfully flat field extension, and any necessary enlargement is
still finite over \(F_0\). Thus it defines the required ordinary
pro-morphism, and its map on limits is pinned. If two such morphisms
induce the same pro-morphism, compare them at target \(F_0\), refine
once, and apply the same finite-stage faithfulness and base change.
They are equal. This proves full faithfulness in the localized
category. The final formulation uses the equivalence of that
localization with absolutely weakly normal schemes, as in
\citep[Theorem~1.13, Lemma~1.32, Proposition~2.28]{CarlsonHaineWolf2024}.
\end{proof}
\paragraph{Check record.} Also checked in this writing, 2026-11-18: GPT-6 Astra (OpenAI, ChatGPT Pro mode), root; separate effort setting not exposed. Full targeted reading of ordinary finite-transcendence pro-reconstruction, all formal-arc avoidance cases, the henselian centralizer example, stable marked logarithmic diagrams and uniqueness, full finite-cover-to-small-site gluing, and arbitrary rank-one logarithmic approximation. Configuration reader separately opened all exact primary inputs. Inverse-image orientation is made explicit; the fixed log point retains the entire arithmetic Galois group.


\begin{claim}{B.44}{Proved}{The remaining ordinary pro-category issue is an infinite-transcendence issue}
The mechanism in B.43 does not extend by the same argument to an
arbitrary characteristic-zero field: transitions corresponding to
transcendental field extensions need not preserve closed points.
Specifically, for \(F\) of characteristic zero and \(u\) transcendental,
the point \(T=u\) of \(\mathbb A^1_{F(u)}\) is closed, whereas its
image in \(\mathbb A^1_F\) is generic. Consequently the full model
system is not in general an object of the ordinary pro-category of
pinned geometric morphisms. This is an obstruction to that categorical
formulation, and is not a counterexample to reconstruction for pinned
morphisms in \(\operatorname{Pro}(\mathsf T)_{/\mathcal B}\).
\end{claim}
\begin{proof}
The contraction of the ideal \((T-u)\) from \(F(u)[T]\) to \(F[T]\)
is zero, because \(u\) is transcendental. This proves the assertion
on points. The scheme projection induces the transition geometric
morphism of the model system and has that same underlying-space map.
An inverse system in a category must have all its transition arrows
in that category. Finally, no assertion in this example constructs
a pinned ordinary pro-morphism failing to come from a scheme map;
thus it supplies no such counterexample.
\end{proof}
\paragraph{Check record.} Also checked in this writing, 2026-11-18: GPT-6 Astra (OpenAI, ChatGPT Pro mode), root; separate effort setting not exposed. Full targeted reading of ordinary finite-transcendence pro-reconstruction, all formal-arc avoidance cases, the henselian centralizer example, stable marked logarithmic diagrams and uniqueness, full finite-cover-to-small-site gluing, and arbitrary rank-one logarithmic approximation. Configuration reader separately opened all exact primary inputs. Inverse-image orientation is made explicit; the fixed log point retains the entire arithmetic Galois group.


\begin{claim}{B.45}{Proved}{The infinite-transcendence obstacle and its resolution}
For \(\operatorname{trdeg}_{\Q}k=\infty\), ordinary relative
pro-reconstruction is proved by B.62--B.66 below, culminating in
B.66. The shorter argument of B.43 alone does not prove it:
a representative may have the form
\(X_{E,\mathrm{et}}\to Y_{F,\mathrm{et}}\) with \(E/F\)
transcendental. Pinnedness of the map on limits does not, by the
argument of B.43, give pinnedness of a finite-stage representative.
The coherent same-field result AS-FF.13 supplies neither that
finite-stage assertion nor an equivalence with this ordinary
pro-category. The new proof resolves the obstacle by deriving
coherence from compactness, comparing whole Kummer systems after
one finitely generated refinement, and reconstructing every étale
slice over a fixed source field. It does not assert that those
finite-stage transition morphisms preserve closed points.
\end{claim}


\begin{proof}
B.62 supplies coherent relative realization and its slice compatibility
from the ordinary pro-category data. B.63 compares the whole Kummer
system after one finitely generated refinement. B.64 then recovers an
actual value over the fixed source field. B.65 reconstructs that
representative and its entire inverse-image functor, and B.66 descends
the resulting morphism independently of the target level. Thus the
ordinary pro-category formulation is fully faithful in the asserted
arbitrary transcendence degree. The example in B.44 concerns only the
different category with pinned transition arrows.
\end{proof}
\paragraph{Check record.} Also checked in this writing, 2026-11-18: GPT-6 Astra (OpenAI, ChatGPT Pro mode), root; separate effort setting not exposed. Read complete proofs: ordinary pro-form consequence, solvable strict-local closure via exact tame/wild sequence, metabelian stable Kummer kernel and complementary-unit intersection, nilpotent Kummer theorem including p2/p3 exceptions, rational residues, affine similarity and polynomial point maps. B56 nontriviality is now proved directly by odd-prime affine groups and the p2 dihedral splitting field, without wild freeness. B57 reconstruction scope updated by B106–110.

\subsection{A henselian field with no centralizing coefficient automorphisms}

\begin{claim}{B.46}{Proved}{A formal arc with no nontrivial algebraic self-correspondence}
There exists \(u\in t\Q[[t]]\), transcendental over \(\Q(t)\), with
the following property. Let \(\Omega\) be any complete rank-one valued extension
of \(\Q((t))\), with the valuation trivial on \(\Q\), and let
\(h\in\Omega\) satisfy \(v(h)=v(t)=1\). If both \(h\) and
\(u(h)\) are algebraic over \(\Q(t,u(t))\), and \(h,u(h)\) are
algebraically independent over \(\Q\), then
\[
                         h=t.                            \tag{B.CF.3}
\]
The series \(u\) can be specified by a recursive sequence of rational
finite jets; the proof gives a finite elimination certificate at every
stage. The property holds simultaneously in all finite ramified
extensions and their completions.
\end{claim}
\begin{proof}
Let \(Z\subset\mathbb A^2_{\Q}\times\mathbb A^2_{\Q}\) be an
irreducible closed subvariety of dimension two, dominant under both
projections, and different from the diagonal. We first prove that
in every cylinder of prescribed rational coefficients in \(t\Q[[t]]\)
there is a smaller cylinder none of whose series has a pair
\[
                 (t,u(t),h,u(h))\in Z(\Omega),\qquad v(h)=1.
                                                        \tag{B.CF.4}
\]
Write the prescribed jet as \(J(t)\) modulo \(t^N\), \(N\geq1\).
Consider the two-parameter family
\[
                       P(t)=J(t)+a t^N+b t^{N+1}.
\]
A generic member has finite intersection
\[
                         Z\cap(\Gamma_P\times\Gamma_P), 
                                                        \tag{B.CF.5}
\]
where \(\Gamma_P=\{(x,P(x))\}\). Here is a dimension proof, including
the loci where evaluation parameters coincide. On
\(xx'(x-x')\ne0\), evaluation at \(x,x'\) is a surjective linear
map from the two parameters \((a,b)\) to two values; the determinant
is \(x^N(x')^N(x'-x)\). The incidence variety of (B.CF.5) consequently
has dimension at most two over this locus. When \(x=x'\ne0\), the
point must lie in \(Z\cap\mathrm{Diag}\), of dimension at most one;
one independent evaluation condition again leaves total dimension at
most two. When \(x=0\), \(x'\ne0\), the first point is the fixed
point \((0,J(0))\). A fibre of the first dominant projection of the
irreducible surface \(Z\) has dimension at most one, and evaluation
at \(x'\) imposes one independent parameter condition. The same
argument applies with the factors exchanged. When both \(x=x'=0\),
only the single pair of fixed points is possible. Thus the full
incidence variety has dimension at most the two-dimensional parameter
space. Generic fibres are zero-dimensional or empty. We may choose
\(a,b\in\Q\), because rational points are Zariski dense in that
parameter space, including after passage to an algebraic closure.
This proves the assertion for the entire affine intersection; no
localization or denominator is needed.

Choose polynomial generators \(Q_j(x,y,x',y')\) of the ideal of
\(Z\). Since (B.CF.5) is zero-dimensional, its coordinate algebra is
finite-dimensional over \(\Q\). Elimination therefore supplies a
nonzero polynomial \(R(x)\in\Q[x]\) and polynomials
\(A_j(x,x')\in\Q[x,x']\) with
\[
 R(x)=\sum_j A_j(x,x')Q_j\bigl(x,P(x),x',P(x')\bigr).
                                                        \tag{B.CF.6}
\]
If \(u\equiv P\pmod{t^M}\), \(v(h)=1\), and (B.CF.4) holds,
then each summand on the right after \((x,x')=(t,h)\) has valuation
at least \(M\). Indeed \(u(t)-P(t)\) and \(u(h)-P(h)\) have
valuation at least \(M\), all rational coefficients have valuation
zero, and all remaining arguments have nonnegative valuation.
Taking \(M>\operatorname{ord}_{t=0}R(t)\), as well as
\(M>\deg P\), contradicts (B.CF.6). The same finite coefficient
cylinder is excluded for every valued extension \(\Omega\), not
just for one choice of Puiseux series field.

There are countably many such rational varieties \(Z\). Each
forbidden set is nowhere dense by the cylinder argument. The series
algebraic over \(\Q(t)\) also form a countable union of nowhere dense
sets: for every nonzero \(A(t,U)\in\Q[t,U]\), a polynomial \(P\)
with any prescribed jet can be chosen so that \(A(t,P(t))\ne0\),
and a sufficiently long further jet then preserves its nonzero
leading coefficient. The complete metric space \(t\Q[[t]]\), with
the coefficient topology, is a Baire space. Choose \(u\) outside
all these forbidden sets.

For the claimed pair \(h,u(h)\), take the rational Zariski closure of
\((t,u(t),h,u(h))\). Algebraicity and the two algebraic-independence
conditions make this an irreducible surface dominant on both factors.
If it were different from the diagonal it would be among the excluded
varieties. It is therefore the diagonal, which gives (B.CF.3).

For a recursive specification, enumerate rational polynomial ideals
which define the indicated irreducible dominant surfaces, and
interleave all nonzero polynomials \(A(t,U)\). At each step choose
the first rational pair \((a,b)\) in a fixed enumeration for which
(B.CF.5) is zero-dimensional, compute an elimination polynomial in
(B.CF.6), and extend the fixed jet beyond its order. Increase the
length strictly at every step. Polynomial ideal dimension,
irreducibility and elimination over \(\Q\) are finite algorithms;
the preceding proof guarantees that every search terminates. The
nested jets define a single series and prove its property. No finite
initial segment alone is asserted to establish all the infinitely
many exclusions.
\end{proof}
\paragraph{Check record.} Also checked in this writing, 2026-11-18: GPT-6 Astra (OpenAI, ChatGPT Pro mode), root; separate effort setting not exposed. Full targeted reading of ordinary finite-transcendence pro-reconstruction, all formal-arc avoidance cases, the henselian centralizer example, stable marked logarithmic diagrams and uniqueness, full finite-cover-to-small-site gluing, and arbitrary rank-one logarithmic approximation. Configuration reader separately opened all exact primary inputs. Inverse-image orientation is made explicit; the fixed log point retains the entire arithmetic Galois group.


\begin{claim}{B.47}{Proved}{Trivial centralizers at every finite extension of a transcendence-degree-two henselian field}
Choose a series \(u\) as in B.46, and let
\[
             K=\Q(t,u(t))^h\subset\Q((t))               \tag{B.CF.7}
\]
be the henselization for the induced \(t\)-adic valuation. Then
\(K\) is countable, has transcendence degree two over \(\Q\), value
group \(\Z\), residue field \(\Q\), and completion \(\Q((t))\).
For every finite extension \(E/K\), in a fixed algebraic closure
\(\overline K\), one has
\[
              C_{\operatorname{Aut}(\overline K)}(G_E)=1.
                                                        \tag{B.CF.8}
\]
Thus the centralizer hypothesis of B7 is attained among the
residue-characteristic-zero fields with positive transcendence defect
specified by B8.
\end{claim}
\begin{proof}
The series is transcendental, so the initial field has transcendence
degree two. It lies in \(\Q((t))\), contains \(\Q(t)\), and has
residue field \(\Q\) and value group \(\Z\). Henselization is
algebraic and immediate. These properties, including its universal
property as a valued-field extension, are given in the opened
\citep[Proposition~3.3.25 and Corollary~3.3.26]{ADH2017}.
The polynomial truncations show that \(\Q(t)\) is dense in
\(\Q((t))\); hence the completion of \(K\) is that same field.
Countability follows from algebraicity over a countable field.

We record explicitly the group calculation used below. Inside a fixed
valued algebraic closure, put
\[
                  K^\dagger=\bigcup_{n\geq1}
                         K\overline{\Q}(t^{1/n}),
                                                        \tag{B.CF.9}
\]
with compatible roots. It is henselian, has algebraically closed
residue field \(\overline\Q\), and has divisible value group
\(\Q\). By the opened
\citep[Corollary~3.3.22]{ADH2017}, a valued field of equicharacteristic
zero is algebraically closed exactly when it is henselian with
algebraically closed residue field and divisible value group.
Thus (B.CF.9) is \(\overline K\). The constants \(\Q\) are
algebraically closed in \(K\), since \(K\subset\Q((t))\).
Adjoining a number field changes only the residue field, and adjoining
\(t^{1/n}\) then has degree \(n\), by its distinct valuation classes.
These facts show that restriction to constants and action on the
compatible roots give
\[
        G_K\simeq\widehat\Z(1)\rtimes G_\Q,             \tag{B.CF.10}
\]
where a splitting acts on the constants and fixes the chosen roots
of \(t\). The action on inertia is cyclotomic.

The centralizer in \(G_K\) of every open subgroup \(H\) is trivial.
An element of that centralizer projects to the centralizer of an
open subgroup of \(G_\Q\), which is trivial by
\citep[Corollary~4.5(3) and Definition~4.3]{HoshiTsujimura2023}.
It consequently lies in \(\widehat\Z(1)\). In each \(\ell\)-adic
coordinate, an element fixed by an open subgroup of \(G_\Q\) is
zero: the \(\ell\)-adic cyclotomic character is nontrivial on the
absolute Galois group of every number field, because a number field
contains only finitely many roots of unity, and multiplication by a
nonzero element of \(\Z_\ell\) is injective on \(\Z_\ell\).
This proves the asserted centralizer calculation without using the
conclusion (B.CF.8).

Now let \(E/K\) be finite and let
\(\sigma\in\operatorname{Aut}(\overline K)\) centralize \(G_E\).
Its fixed-field action preserves \(E\). The unique prolongation of
our valuation to \(E\) is henselian and discrete, with residue field
a number field, so \(E\) is not separably closed. The opened
\citep[Lemma~2.1]{KhassaKhanduja2010} states that every automorphism
of a non-separably-closed henselian valued field is continuous for
its valuation topology; its proof applies the theorem of F.\ K.
Schmidt to the original and transported valuations. For a discrete
rank-one valuation, equivalent valuation topologies give the same
valuation ring. An automorphism of its ordered value group is the
identity. Hence \(v(\sigma a)=v(a)\) on \(E\), and uniqueness of
prolongation over a henselian field gives the same equality on
\(\overline K\).

The field \(\overline\Q\subset\overline K\) is invariant under
\(\sigma\). The restriction of \(G_E\) to it is the open group
\(G_{E\cap\overline\Q}\); this follows from the Galois
correspondence, and the intersection is a number field since
\(E/K\) is finite and \(K\cap\overline\Q=\Q\). Therefore
\(\sigma|_{\overline\Q}\in G_\Q\) centralizes an open subgroup
of \(G_\Q\), and is the identity. In particular all coefficients
of \(u\) are fixed.

Embed \(\overline K\) as a valued field in an algebraic closure of
\(\Q((t))\), and work in its completion. Continuity and convergence
of the rational polynomial truncations give
\[
          \sigma u(t)=u(\sigma t),\qquad v(\sigma t)=1.
                                                        \tag{B.CF.11}
\]
Both \(\sigma t\) and \(\sigma u\) are algebraic over
\(\Q(t,u)\), because \(\overline K\) is algebraic over that
field; they are algebraically independent over \(\Q\), because
\(\sigma\) is a field automorphism. B.46 applied to
\(h=\sigma t\) yields \(\sigma t=t\), and (B.CF.11) gives
\(\sigma u=u\). The uniqueness of a valued embedding of a
henselization now makes \(\sigma\) the identity on \(K\).
It lies in \(G_K\) and centralizes its open subgroup \(G_E\).
The calculation following (B.CF.10) forces \(\sigma=1\), proving
(B.CF.8) for every finite \(E\).
\end{proof}
\paragraph{Check record.} Also checked in this writing, 2026-11-18: GPT-6 Astra (OpenAI, ChatGPT Pro mode), root; separate effort setting not exposed. Full targeted reading of ordinary finite-transcendence pro-reconstruction, all formal-arc avoidance cases, the henselian centralizer example, stable marked logarithmic diagrams and uniqueness, full finite-cover-to-small-site gluing, and arbitrary rank-one logarithmic approximation. Configuration reader separately opened all exact primary inputs. Inverse-image orientation is made explicit; the fixed log point retains the entire arithmetic Galois group.


\subsection{A residue-characteristic-zero obstruction beyond coefficient automorphisms}

\begin{claim}{B.48}{Proved}{Compatible logarithmic reduction under a supported infinitesimal scaling}
Let \(R\) be a henselian discrete valuation ring of residue
characteristic zero, with fraction field \(K\), uniformizer \(t\),
and the unique extension \(v\) of its valuation to \(\overline K\).
Put \(\lambda=1+t\), and define
\[
 F(a)=
 \begin{cases}
   \lambda a,&v(a)>0,\\
   a,&v(a)\leq0,
 \end{cases}
 \qquad F(0)=0,
 \qquad F(\infty)=\infty.                                \tag{B.CF.12}
\]
For every finite \(G_K\)-stable set \(S\subset\mathbb P^1(\overline K)\)
containing \(0,1,\infty\), the stable marked genus-zero models for
\(S\) and \(F(S)\), over a finite Galois extension splitting their
marks, have isomorphic logarithmic special fibres over the same
standard log point. These isomorphisms preserve the labels and are
compatible with forgetting marks, finite extensions of the field of
definition, and the Galois descent data. At a mark \(a\), the induced
comparison of the arithmetic cusp group preserves its inertia and
induces the identity on its quotient \(G_{K(a)}\), with
\(K(F(a))=K(a)\) as actual subfields of \(\overline K\).
\end{claim}
\begin{proof}
The map (B.CF.12) is a \(G_K\)-equivariant bijection, whose inverse
scales the same open valuation ball by \(\lambda^{-1}\).
For two distinct finite points,
\[
          \frac{F(a)-F(b)}{a-b}\in1+\mathfrak m_{\overline K}.
                                                        \tag{B.CF.13}
\]
When both points belong to the ball this ratio is \(\lambda\), and
when neither does it is one. In the remaining case, say
\(v(a)>0\geq v(b)\), the perturbation is \(ta\), which has
valuation strictly larger than \(v(a-b)=v(b)\). This proves
(B.CF.13). In particular every difference has the same valuation and
the same initial form in the associated graded valued field before
and after applying \(F\). The assertion also holds for differences
in the reciprocal coordinate near infinity: \(F\) fixes all points
in that branch.

Choose a finite Galois extension \(L/K\) containing all the marks,
and write \(\pi\) for a uniformizer of its valuation ring. Distinct
rational sections of \(\mathbb P^1_{R_L}\) can be separated by
successive blowups at their common reductions. Each blowup is at a
smooth point of the special fibre. Its exceptional component has an
affine coordinate
\[
                          z=(x-a)/\pi^m,                 \tag{B.CF.14}
\]
for a centre \(a\) in the corresponding cluster of marks, with
\(m\) the depth of that cluster. Equivalently one may use the
reciprocal coordinate on the infinity branch. The finite tree of
components, their depths, and the locations of the marked and nodal
points on each component are obtained from the valuations and
normalized initial forms of the differences between the marks.
Contracting components with fewer than three special points gives
the stable marked model. This description is also a direct
construction of it over \(R_L\).

Equation (B.CF.13) identifies the two trees, their depths, and all
component configurations. In the coordinates (B.CF.14), the
component isomorphisms have the same reduced coordinates. They are
unique on stable components because three labelled special points
fix an automorphism of \(\mathbb P^1\). Equip the models with the
fs divisorial log structures defined by the special fibre and the
marked sections. A node on the stable model has a toric chart
\(uv=\epsilon\pi^d\), with \(d\geq1\) and \(\epsilon\) a
unit. The logarithmic special fibre retains \(d\) and the residue
of the smoothing unit, with its two tangent directions. In cluster
coordinates before contraction one may write
\(uv=\pi^{n-m}\); after changing normalized coordinates the
smoothing unit is a product and quotient of initial forms of the
same differences. Thus (B.CF.13) preserves the smoothing data and
the component maps lift to log isomorphisms over the fixed log point.
These toric charts show at the same time that the total marked model
is log regular; contracting a chain merely replaces thickness one
by a positive integral thickness.

We give the compatibility calculation, since uniqueness of the
isomorphisms by itself would not suffice for contractions. If
\(z=(x-a)/\pi^m\) is a coordinate for a target component and
\(w=(x-b)/\pi^n\) is one for a source component, the transition is
\[
       z=\frac{b-a}{\pi^m}+\pi^{n-m}w.                   \tag{B.CF.15}
\]
Extract the smaller valuation of the two coefficients. Its first
nonzero normalized reduction is a nonzero constant, \(w\), or
\(c+w\); it depends only on the initial form of \(b-a\) and the
two depths. The transformed expression has exactly the same
exponent and the same reduced factor, by (B.CF.13). This checks
both the characteristic monoid and the unit part of the pullback
of each logarithmic generator, including a generator which
vanishes identically on a contracted tail. All node, mark, and
infinity charts are obtained by products, quotients, and
reciprocals of these linear expressions. Hence the logarithmic
isomorphisms commute with every mark-forgetting contraction.
Equality of the unit factors may be checked on the irreducible
components, because the special fibre is reduced.

For completeness, the lift of a fixed component isomorphism is
unique once the characteristic monoids and log base are fixed.
At a node its possible change multiplies the two generators by
units \(u_1,u_2\) with \(u_1u_2=1\), \(u_1x=x\), and
\(u_2y=y\) in the reduced nodal ring. Thus
\(u_1-1\in(x)\cap(y)=0\). On the smooth locus the vertical
generator comes from the fixed log base, and a horizontal generator
maps to a non-zero-divisor. There is no additional freedom there.
The same coordinate check as in (B.CF.15) works under finite base
extension: if \(\pi=\epsilon'(\pi')^e\), the same unit
\(\epsilon'\) occurs on both sides, and initial forms commute
with the map between the associated graded valued fields.
This proves base-extension compatibility. Since (B.CF.12) is
\(G_K\)-equivariant, the uniquely labelled component maps and their
log lifts commute with \(\operatorname{Gal}(L/K)\).

At a marked section (a), the local horizontal parameter is
\(x-a\). Its image is \(x-F(a)\); inside the ball their
comparison is the scaling \(\lambda\), and outside it the
identity. The vertical log structure and the field \(K(a)\)
are unchanged, since multiplication by the nonzero element
\(\lambda\in K\) preserves \(K(a)\). The local Kummer chart
therefore identifies the horizontal inertia \(\widehat\Z(1)\)
and its arithmetic quotient \(G_{K(a)}\), with the identity on
the latter. This is the usual cusp description obtained from the
punctured strict henselian normal trait, or directly from the
Kummer chart in SGA~1, XIII, Proposition~5.2. Choices of geometric
fibre functors change the resulting group maps by the corresponding
cuspidal conjugacy, not by an unspecified automorphism of the
residue field.
\end{proof}
\paragraph{Check record.} Also checked in this writing, 2026-11-18: GPT-6 Astra (OpenAI, ChatGPT Pro mode), root; separate effort setting not exposed. Full targeted reading of ordinary finite-transcendence pro-reconstruction, all formal-arc avoidance cases, the henselian centralizer example, stable marked logarithmic diagrams and uniqueness, full finite-cover-to-small-site gluing, and arbitrary rank-one logarithmic approximation. Configuration reader separately opened all exact primary inputs. Inverse-image orientation is made explicit; the fixed log point retains the entire arithmetic Galois group.


\begin{claim}{B.49}{False as stated}{Reconstruction over an equicharacteristic-zero henselian discretely valued field}
For every field \(K\) admitting a nontrivial henselian discrete
valuation with residue characteristic zero, closed-point-preserving
reconstruction of finite-type \(K\)-schemes from their relative
small étale topoi is false, even with universal homeomorphisms
inverted. The supported scaling (B.CF.12) is the geometric-point
function of a pinned autoequivalence
\[
          (\mathbb A^1_K)_{\mathrm{et}}
             \longrightarrow(\mathbb A^1_K)_{\mathrm{et}}
                  \quad\hbox{over }BG_K                 \tag{B.CF.16}
\]
which is not induced by a morphism of \(K\)-schemes.
\end{claim}
\begin{proof}
We use the inverse-image orientation: a covering with target puncture
set $F(S)$ is sent to the corresponding covering with source puncture
set $S$. With this orientation the induced map on underlying points
is $F$. Transport in the reverse direction constructs its inverse.

We first state precisely the opened comparison theorem.
\citep[Theorem~1.6]{Lepage2013} assumes that \(\mathscr X\) is a
log-regular fs log scheme and every prime in \(L\) is invertible
on it. Restriction identifies its finite Kummer-étale covers of
pro-\(L\) type with the corresponding finite étale covers of its
log-trivial locus. Theorem~1.8 in the same paper assumes a
noetherian henselian base \(B\) with closed point \(s\), a
connected fs log scheme \(\mathscr X\), and properness of its
underlying scheme over \(B\); restriction identifies finite
Kummer-étale covers of \(\mathscr X\) with those of its log
special fibre. These are equivalences of covering categories,
functorial in the log schemes. For our models all primes are
invertible: their valuation rings contain \(\Q\). Taking
\(L\) to be the set of all primes gives, for each marked model,
\[
 \operatorname{FEt}(\mathbb P^1_L\setminus S)
       \simeq\operatorname{KCov}(\mathscr X_S)
       \simeq\operatorname{KCov}((\mathscr X_S)_s).
                                                        \tag{B.CF.17}
\]
Here \(L/K\) now denotes a finite splitting field for the marks;
the prime-set letter in the quoted theorem has served its purpose.
The model is proper over a henselian DVR and log regular by the
charts in B.48. Thus every hypothesis has been checked. This
application retains all finite covers, with no prime omitted.

Apply (B.CF.17) to the logarithmic special-fibre isomorphisms of
B.48. We obtain compatible inverse-image equivalences
\[
 \operatorname{FEt}(\mathbb P^1_L\setminus F(S))
           \longrightarrow
 \operatorname{FEt}(\mathbb P^1_L\setminus S).            \tag{B.CF.18}
\]
Their compatibility with mark-forgetting, finite constant
extensions, and Galois descent was checked in (B.CF.15), so
finite étale descent gives the same system over \(K\).
These equivalences are over \(\operatorname{FEt}(K)\), rather
than merely geometric: apply the two quoted theorems also to
\(\operatorname{Spec}R\) with its vertical log structure.
They identify covers of the fixed log point with finite étale
\(K\)-algebras, and every special-fibre isomorphism used above
lies over that same log point. At each puncture (B.CF.18)
identifies the full decomposition and inertia pair, and acts as
the identity on the actual quotient \(G_{K(a)}\), as proved
at the end of B.48.

We now pass from this compatible system of finite covering
categories to the entire small étale site. This step is necessary;
an isomorphism at one punctured fundamental group alone would
not suffice. A quasi-compact separated étale object
\(V\to\mathbb P^1_K\) is a disjoint union of normal curves.
Zariski's Main Theorem identifies it with an open subscheme of
the finite normal compactification \(\overline V\to\mathbb P^1_K\)
in its generic finite separable algebras. Its complement is a
finite set of points. Choose a finite set \(S\) containing the
branch locus, the images of every point of
\(\overline V\setminus V\), and \(0,1,\infty\).
Over its complement the compactification is finite étale.
Use (B.CF.18) to transport this
finite étale cover, and normalize \(\mathbb P^1_K\) in the
resulting generic algebra. The full cusp pairs give the
bijections between the points over the omitted base points:
ramification is measured by inertia, and the unramified residue
extension by the quotient decomposition group over the fixed
\(G_{K(a)}\). Therefore they transport the removed points of
\(\overline V\setminus V\), preserve the unramified locus, and
produce a quasi-compact separated étale object again.
The construction is independent, up to its compatible natural
isomorphisms, of enlarging \(S\), by the restriction
compatibility in (B.CF.18).

A morphism between such étale objects is determined by its map
of generic finite algebras, extended uniquely to their normal
compactifications, together with the condition that its image
lies in the chosen open subsets. The covering-category
comparisons preserve these generic maps and the conditions on
points. Thus they give an equivalence of these categories of
étale objects. They preserve covering families: joint
surjectivity at a geometric generic point is detected by the
finite covering category; at a geometric closed point it is
detected by the same local decomposition/inertia comparison.
These objects form a generating small étale site. We have
therefore constructed an autoequivalence of
\((\mathbb P^1_K)_{\mathrm{et}}\), over \(BG_K\), whose
function on \(\overline K\)-valued points is exactly (B.CF.12). It is pinned,
since that function is a \(G_K\)-equivariant bijection and
preserves the finite field \(K(a)\) at every algebraic point.
It fixes infinity and hence restricts to (B.CF.16).

Every nonzero rational number is a valuation unit and is fixed
by (B.CF.12), whereas \(t\) is sent to \((1+t)t\ne t\).
If a \(K\)-scheme morphism \(\mathbb A^1_K\to\mathbb A^1_K\)
induced (B.CF.16), its polynomial would agree with \(T\) on the
infinite subset \(\Q^\times\), and would consequently be \(T\)
identically. Its value at \(t\) gives the contradiction.
The scheme \(\mathbb A^1_K\) is already absolutely weakly
normal, so inverting universal homeomorphisms introduces no
additional morphism which could realize this point function.
\end{proof}
\paragraph{Check record.} Also checked in this writing, 2026-11-18: GPT-6 Astra (OpenAI, ChatGPT Pro mode), root; separate effort setting not exposed. Full targeted reading of ordinary finite-transcendence pro-reconstruction, all formal-arc avoidance cases, the henselian centralizer example, stable marked logarithmic diagrams and uniqueness, full finite-cover-to-small-site gluing, and arbitrary rank-one logarithmic approximation. Configuration reader separately opened all exact primary inputs. Inverse-image orientation is made explicit; the fixed log point retains the entire arithmetic Galois group.


\begin{claim}{B.50}{False as stated}{The centralizer condition in B7 is insufficient}
There exists a characteristic-zero field \(K\), countable and of
transcendence degree two over \(\Q\), such that
\[
 C_{\operatorname{Aut}(\overline K)}(G_E)=1
           \quad\hbox{for every finite }E/K,
\]
but closed-point-preserving reconstruction of finite-type
\(K\)-schemes from their relative étale topoi fails. One may take
exactly the recursively specified field (B.CF.7). The obstruction
is a compatible nonalgebraic permutation of points and their
whole covering categories, not a centralizing field automorphism.
\end{claim}
\begin{proof}
B.47 proves the displayed condition for this field, with all
finite extensions included. Its valuation is henselian and
discrete, with residue field \(\Q\), so B.49 supplies the
pinned nonalgebraic autoequivalence of its affine-line topos.
These are the same field and the same entire arithmetic base
\(BG_K\). This is a counterexample to the asserted implication.
We also verify the surviving necessary direction, rather than
accepting AS-FF.6 on authority. If a field satisfies reconstruction,
then every finite extension does: finite étale slicing identifies
objects and relative morphisms over that extension, and faithfulness
forces the reconstructed map to preserve the extension's structural
map. Let \(E/K\) be finite and let \(\sigma\) centralize
\(G_E\). It preserves \(E\). The coefficient automorphism on
\(\mathbb G_{m,E}\), together with the identification of its
base Galois action supplied by \(\sigma G_E\sigma^{-1}=G_E\)
pointwise, gives a pinned relative topos automorphism whose geometric
point function is \(a\mapsto\sigma(a)\). This includes the case
where \(\sigma\) fixes \(E\): then it is the identity absolute
morphism with its relative structural 2-cell twisted by the central
element. The point-function description follows directly by pulling
back to finite separable extension points, as in the opened
\citep[Proposition~3.6]{CarlsonHaineWolf2024}.
Reconstruction over \(E\) would make this an \(E\)-automorphism
of \(\mathbb G_m\), hence \(T\mapsto cT^{\pm1}\).
Evaluation at one gives \(c=1\); evaluation at two excludes
inversion. Consequently \(\sigma(a)=a\) for every
\(a\in\overline K^\times\), so \(\sigma=1\).
This proves necessity for all finite extensions. B.50 shows that
excluding this particular obstruction is not sufficient.
A precise valid sufficient-and-necessary replacement is the
point-function condition of the opened
\citep[Theorem~4.18(5)]{CarlsonHaineWolf2024}: for every affine
regular connected finite-type source with a rational point,
every pinned relative map to \(\mathbb G_m\) must have the
geometric-point function of a regular unit. The supported
scaling demonstrates why a condition solely on coefficient
field automorphisms cannot replace that requirement.
\end{proof}
\paragraph{Check record.} Also checked in this writing, 2026-11-18: GPT-6 Astra (OpenAI, ChatGPT Pro mode), root; separate effort setting not exposed. Full targeted reading of ordinary finite-transcendence pro-reconstruction, all formal-arc avoidance cases, the henselian centralizer example, stable marked logarithmic diagrams and uniqueness, full finite-cover-to-small-site gluing, and arbitrary rank-one logarithmic approximation. Configuration reader separately opened all exact primary inputs. Inverse-image orientation is made explicit; the fixed log point retains the entire arithmetic Galois group.


\begin{claim}{B.51}{False as stated}{The specified transcendence-defect field in AS-FF.18 does not reconstruct}
Let
\[
     u_*(t)=\sum_{n\geq1}t^{n!},\qquad
     K_* =\Q(t,u_*(t))^h\subset\Q((t)).                  \tag{B.CF.19}
\]
This field has transcendence degree two, residue field \(\Q\),
value group \(\Z\), and positive transcendence defect one.
It does not satisfy closed-point-preserving étale reconstruction.
This settles the particular residue-characteristic-zero case
left undecided in AS-FF.18, without claiming that its centralizers
have the property proved for the different series in B.46.
\end{claim}
\begin{proof}
For a nonzero \(P(t,U)\in\Q[t,U]\), write \(e=\deg_U P\)
and \(m=\deg_t P\), and put
\(u_N=\sum_{n=1}^N t^{n!}\). If \(P(t,u_*)=0\), then a
nonzero \(P(t,u_N)\) has order at least \((N+1)!\) but degree
at most \(m+eN!\), impossible for sufficiently large \(N\).
Thus infinitely many distinct \(u_N\) would be roots of the
nonzero polynomial \(P(t,U)\) over \(\Q(t)\), again impossible.
The series is transcendental. The field and its henselization
therefore have the asserted transcendence degree, and
henselization preserves residue field and value group.
The defect is \(2-0-1=1\). B.49 applies directly to this
henselian discrete valuation and refutes reconstruction.
\end{proof}
\paragraph{Check record.} Also checked in this writing, 2026-11-18: GPT-6 Astra (OpenAI, ChatGPT Pro mode), root; separate effort setting not exposed. Full targeted reading of ordinary finite-transcendence pro-reconstruction, all formal-arc avoidance cases, the henselian centralizer example, stable marked logarithmic diagrams and uniqueness, full finite-cover-to-small-site gluing, and arbitrary rank-one logarithmic approximation. Configuration reader separately opened all exact primary inputs. Inverse-image orientation is made explicit; the fixed log point retains the entire arithmetic Galois group.



\noindent\emph{Statement B.52 and its produced argument are retained in Part II.}


\begin{claim}{B.53}{False as stated}{The unbounded solvable closure of the strict-local field}
Fix an embedding of the strict henselization
\(A=\Z_{(p)}^{sh}\) in \(\overline\Q\), put
\(K_p=\operatorname{Frac}(A)\), and write \(I=G_{K_p}\).
Let \(I^{(0)}=I\) and
\(I^{(d+1)}=\overline{[I^{(d)},I^{(d)}]}\) be the closed
derived series. Then
\[
 K_p^{\mathrm{sol}}
   :=\bigcup_{d\geq 0}\overline\Q^{\,I^{(d)}}
   =\overline\Q.                                      \tag{B.CF.20}
\]
Consequently this maximal solvable closure does not satisfy
closed-point-preserving \(\acute etale\) reconstruction. The statement
concerns the actual field in AS-FF.39, not a substitution of a
completed local field. The bounded derived-length field $B_p$ and the lower-central union are different fields; B.136 and B.138 respectively prove their reconstruction below.
\end{claim}
\begin{proof}
The ring \(A\) is a strictly henselian discrete valuation ring
with residue field \(\overline{\F}_p\). For a finite Galois
extension \(M/K_p\), henselianity gives a unique extension of
the valuation. Thus its decomposition group is the whole
\(\operatorname{Gal}(M/K_p)\). Its residue extension is trivial,
since it is finite over \(\overline{\F}_p\), so the whole
Galois group is also its inertia group. The opened
\citep[Lemma~15.114.5, Tag~09EE]{StacksProject} gives an exact
sequence with a finite \(p\)-group as kernel and a cyclic
prime-to-\(p\) group as quotient. Both are solvable, hence every
finite quotient of \(I\) is solvable. Equivalently one may apply
the profinite version, Lemma~58.13.5, Tag~0BUA, to this same DVR.

Every nonidentity element of a profinite group survives in some
finite quotient. If that quotient is solvable of derived length
\(d\), the image of \(I^{(d)}\) in it is trivial. It follows
that \(\bigcap_d I^{(d)}=1\). Alternatively, every element of
\(\overline\Q\) lies in a finite Galois extension of \(K_p\),
whose finite solvable Galois group has some finite derived length;
that extension lies in \(\overline\Q^{I^{(d)}}\). This proves
(B.CF.20), without asserting any uniform bound on derived length.

Choose a nonidentity automorphism \(\sigma\) of
\(\overline\Q/\Q\), for example complex conjugation after
choosing a complex embedding. Acting on coefficients induces an
autoequivalence of \(\mathbb A^1_{\overline\Q,\mathrm{et}}\)
over \(B G_{\overline\Q}=\mathsf{Sets}\); it preserves closed
points and acts on their values by \(a\mapsto\sigma(a)\).
If induced by a morphism of \(\overline\Q\)-schemes, its
polynomial would fix every rational number, and hence would be
\(T\). It moves an algebraic number, a contradiction. Inverting
universal homeomorphisms does not change this conclusion for the
normal characteristic-zero affine line. This is a genuine
negative reconstruction result for the exact unbounded derived
tower, rather than a failure of a proposed proof of reconstruction.
\end{proof}
\paragraph{Check record.} Also checked in this writing, 2026-11-18: GPT-6 Astra (OpenAI, ChatGPT Pro mode), root; separate effort setting not exposed. Read complete proofs: ordinary pro-form consequence, solvable strict-local closure via exact tame/wild sequence, metabelian stable Kummer kernel and complementary-unit intersection, nilpotent Kummer theorem including p2/p3 exceptions, rational residues, affine similarity and polynomial point maps. B56 nontriviality is now proved directly by odd-prime affine groups and the p2 dihedral splitting field, without wild freeness. B57 reconstruction scope updated by B106–110.



\subsection{Nondiscrete valuations of residue characteristic zero}

The logarithmic argument can be extended beyond DVRs. The full
valuative log structure is essential here: replacing it by a fixed
finitely generated submonoid would, in general, change its Kummer
coverings and its arithmetic fundamental group.

\begin{claim}{B.54}{Proved}{All-prime logarithmic comparison over a henselian rank-one valuation ring}
Let $R$ be a henselian rank-one valuation ring with fraction field $K$
and residue field $\kappa$ of characteristic zero. Let
$S\subset\mathbb P^1(K)$ be a finite set containing $0,1,\infty$,
and let $\mathscr X/R$ be its stable marked genus-zero model.
Endow $\operatorname{Spec}R$ with the canonical valuative log
structure $R\setminus\{0\}\to R$, and $\mathscr X$ with the
log structure of its vertical boundary and its marked sections.
Write $\mathscr X_s$ for its log closed fibre and
$U=\mathbb P^1_K\setminus S$. Then restriction gives equivalences
\[
 \operatorname{FEt}(U)
 \ \longleftarrow\
 \operatorname{KCov}(\mathscr X)
 \ \longrightarrow\
 \operatorname{KCov}(\mathscr X_s).                 \tag{B.LIT.1}
\]
Here $\operatorname{KCov}$ denotes all finite Kummer \'etale
coverings in the saturated, not necessarily fs, sense of
Achinger--H\"ubner--Lara--Stix. In that definition the underlying
scheme map is integral and need not be finite. The equivalences
are over the entire $BG_K$, and are compatible with finite valued
extensions and with morphisms between stable models obtained by
forgetting marks. Neither discreteness nor completeness is assumed.
\end{claim}
\begin{proof}
We state the exact source inputs. AHLS Lemma~5.3.2 says that a
microbial valuation ring of residue characteristic zero is a
filtered colimit of local excellent regular subalgebras
$\Lambda_i$ with local maps to $R$, containing a fixed
pseudouniformizer $\pi$, for which $V(\pi)_{\mathrm{red}}$ is
a strict normal crossings divisor. If $R$ is henselian the
$\Lambda_i$ may be chosen henselian. A nontrivial rank-one
valuation is microbial, because $R[1/\pi]=K$ for every
$0\ne\pi\in\mathfrak m_R$. Put
\[
 M=R\setminus\{0\},\qquad
 M_i=\Lambda_i\cap\Lambda_i[1/\pi]^\times.
\]
The paragraph following that lemma gives
$(M\to R)=\varinjlim_i(M_i\to\Lambda_i)$ as saturated prelog
rings; the associated log schemes $B_i$ are fs and log regular.
AHLS Proposition~4.2.6 identifies finite Kummer \'etale covering
categories with their filtered colimit for an affine charted
inverse system of qcqs saturated log schemes, in the sense of its
Definition~3.4.4. These are the precise continuity and approximation
theorems used below.\citep[Definition~3.4.4,
Proposition~4.2.6, Lemma~5.3.2]{LA_AHLS2024}

For completeness, the required stable family exists over $R$,
without a noetherian assumption. Hassett Theorem~2.1, with genus
zero and all weights one, gives a proper moduli stack and its
projective coarse scheme over $\mathbb Z$. Every stable labelled
genus-zero curve has trivial automorphism group, so in this case
the stack is the fine projective moduli scheme
$\overline M_{0,|S|}$. Its universal family and the valuative
criterion give $\mathscr X$. The finite-presentation moduli
problem over arbitrary bases is obtained from the initially
noetherian formulation by descent of its finite presentation.
The model and its generic identification are unique, and its base
change to an extension of valuation rings is the corresponding
stable model.\citep[Theorem~2.1]{LA_Hassett2003}

One can describe its finite chart data directly. For a finite set
of split marks, the finite tree of branching clusters is determined
by the valuations of pairwise differences. At a branching radius,
choose a difference $d\ne0$ of that valuation and blow up the ideal
$(x-a,d)$ in the appropriate smooth coordinate. Its charts are
$x-a=dz$ and $(x-a)w=d$. Repeat at the finitely many branching
radii, omitting or contracting components with fewer than three
special points. This is the stable model just described. Its node
charts have the form $xy=a$, for $0\ne a\in\mathfrak m_R$;
its marked charts have a smooth horizontal parameter. All nodes
and their branches are split, since the marks are split and their
dual graph is a labelled tree.

Over the base monoid $M$, the log charts are $M$, $M\oplus\mathbb N$
at a mark, and
\[
       M[e_1,e_2]/(e_1+e_2=a)
                                                        \tag{B.LIT.2}
\]
at a node. AHLS Example~2.4.5 proves that these node monoids are
saturated and that the maps from $M$ are smooth, vertical and
saturated for every saturated $M$ and every $a\in M$.
The standard vertical log structure agrees with these charts by
AHLS Proposition~5.3.1; the disjoint smooth horizontal sections
add the displayed $\mathbb N$ factors. The log-trivial open is $U$.

Descend the marked stable family, its finite chart cover, and the
chart gluings to a sufficiently large $\Lambda_i$. This is ordinary
finite-presentation descent; equivalently, descend the $R$-point
of the fine moduli scheme and pull back its universal family.
After increasing $i$, all the finitely many node parameters and
chart identities descend. For each such nonzero parameter $a$,
rank one gives $N$ with $\pi^N/a\in R$. Require $a$ and
$\pi^N/a$ to belong to $\Lambda_i$. Then $a\in M_i$; on the
regular snc base it is a unit times a monomial in the local
equations of $V(\pi)_{\mathrm{red}}$. Include also the inverses
of the finitely many unit coefficient differences which separate
sections and chart overlaps. All these requirements persist on a
cofinal subsystem.

Let $\mathscr X_i/B_i$ denote the resulting marked log curve.
The proper stable underlying scheme is an ordinary base change
of one fixed descended family. Its local log charts are the
versions of (B.LIT.2) and $M_i\oplus\mathbb N$ over $M_i$.
It is thus fs and log smooth over the log-regular $B_i$, hence
log regular by the log-regularity theorem recalled immediately
after AHLS Lemma~5.3.2. All integers are invertible: the
$\Lambda_i$ can be taken as $\mathbb Q$-algebras in its
characteristic-zero construction. Let $U_i$ be the log-trivial
open; its limit is $U$.

The node and marked monoid maps are saturated. Their saturated
log base changes consequently have the same underlying schemes
as ordinary base changes. The compatible finite strict \'etale
chart covers make $\{\mathscr X_i\}$ an affine charted system
in exactly AHLS Definition~3.4.4, with limit $\mathscr X$.
The same applies to the log closed fibres
$\mathscr X_{i,s_i}$. Indeed, all the maps $\Lambda_i\to R$
and the transition maps are local inclusions. Their induced
residue-field maps are injective: the inverse image of
$\mathfrak m_R$ in each ring is its maximal ideal. Those
residue fields therefore have colimit $\kappa$. Taking the same charts modulo their maximal
ideals therefore gives the log fibre $\mathscr X_s$ as their
saturated inverse limit.

Lepage Theorem~1.6 identifies finite Kummer \'etale covers of an
fs log-regular scheme with finite \'etale covers of its log-trivial
open, for primes invertible on that scheme. His Theorem~1.8
identifies finite Kummer \'etale covers of a connected fs log
scheme proper over a noetherian henselian base with those of its
log closed fibre. Apply both theorems to every $\mathscr X_i$,
with all primes. They give
\[
 \operatorname{FEt}(U_i)\ \longleftarrow\
 \operatorname{KCov}(\mathscr X_i)\ \longrightarrow\
 \operatorname{KCov}(\mathscr X_{i,s_i}).
\]
Ordinary finite \'etale continuity on the left and AHLS
Proposition~4.2.6 on the other two terms prove (B.LIT.1).
This is a limit argument using the full valuative monoid, not an
application of the fs theorems directly to $R$.
\citep[Theorems~1.6 and 1.8]{Lepage2013}

Apply the same argument to the log base, proper over itself.
It identifies the covering categories of that base and its log
closed point with $\operatorname{FEt}(K)$. Independently, AHLS
Corollary~4.4.9 says that the Kummer \'etale fundamental group of
the canonical log valuation ring is the maximal tame quotient of
$G_K$. Every finite extension is tame in residue characteristic
zero, so this is all of $G_K$. Since all the restriction functors
commute with pullback of base covers, (B.LIT.1) is over $BG_K$.

For a finite valued extension $L/K$, AHLS Proposition~4.4.3
identifies the canonical valuative log map as a finite Kummer
\'etale map in the stated integral sense. Stable models base
change as schemes, and their saturated node charts base change
as just checked. This yields compatibility with the actual field
extension. For a finite diagram of mark-forgetting morphisms,
descend the diagram and all its charts to a common stage; the
restriction squares commute at that stage and hence at the limit.
Equivalently, a transported cover with its specified restriction
has a contractible groupoid of extensions by full faithfulness.
The resulting comparison isomorphisms therefore satisfy their
cocycle identities. This proves the asserted functoriality.
\citep[Proposition~4.4.3, Corollary~4.4.9]{LA_AHLS2024}
\end{proof}
\paragraph{Check record.} Also checked in this writing, 2026-11-18: GPT-6 Astra (OpenAI, ChatGPT Pro mode), root; separate effort setting not exposed. Full targeted reading of ordinary finite-transcendence pro-reconstruction, all formal-arc avoidance cases, the henselian centralizer example, stable marked logarithmic diagrams and uniqueness, full finite-cover-to-small-site gluing, and arbitrary rank-one logarithmic approximation. Configuration reader separately opened all exact primary inputs. Inverse-image orientation is made explicit; the fixed log point retains the entire arithmetic Galois group.


\begin{claim}{B.55}{False as stated}{Reconstruction over all finite-transcendence-degree henselian fields of residue characteristic zero}
Every field $K$ admitting a nontrivial henselian rank-one valuation
of residue characteristic zero fails closed-point-preserving
reconstruction of finite-type $K$-schemes from their \'etale topoi
over $BG_K$. In particular this holds for every field of finite
transcendence degree over $\mathbb Q$ with a nontrivial henselian
valuation of residue characteristic zero, regardless of its rank,
discreteness, or transcendence defect.
\end{claim}
\begin{proof}
We first work in rank one. Choose $0\ne t\in\mathfrak m_K$ and
put $\lambda=1+t$. On $\mathbb P^1(\overline K)$ set
\[
 F(a)=\begin{cases}\lambda a,&v(a)>0,\\a,&v(a)\leq0,\end{cases}
 \qquad F(0)=0,\qquad F(\infty)=\infty.
                                                        \tag{B.LIT.3}
\]
Henselianity makes the extended valuation $G_K$-invariant.
Thus $F$ is an equivariant bijection, with inverse scaling the
same ball by $\lambda^{-1}$. As in \statementref{B.48},
\[
       \frac{F(a)-F(b)}{a-b}\in1+\mathfrak m_{\overline K}
       \quad(a\ne b).                                  \tag{B.LIT.4}
\]
When both points lie in the ball the ratio is $\lambda$; when
neither does it is one. In the remaining case
$v(a)>0\geq v(b)$, the change $ta$ has strictly greater
valuation than $a-b$. This proves (B.LIT.4) for an arbitrary
ordered rank-one value group. Hence every difference has the same
valuation and initial form before and after $F$.

Take a finite $G_K$-stable marked set containing $0,1,\infty$
and a finite Galois extension $L/K$ splitting it. The stable
models for $S$ and $F(S)$ have the same labelled cluster tree.
For a cluster with centre $a$ choose a scalar $d\in L^\times$
realizing its radius; use that \emph{same} scalar for the
transformed cluster with centre $F(a)$. Instead of the integral
power coordinates of the DVR proof, write $z=(x-a)/d$. If
$w=(x-b)/e$ is another component coordinate, the transition is
\[
                  z=\frac{b-a}{d}+\frac ed w.            \tag{B.LIT.5}
\]
Extract the coefficient of smaller valuation. Its first nonzero
normalized reduction is a nonzero constant, $w$, or $c+w$.
Equation (B.LIT.4) preserves this reduced factor and its value.
It therefore preserves both the ghost-monoid element and the
unit part of every logarithmic pullback. The smoothing parameter
$e/d$ can be chosen literally the same on the two sides.
Products, quotients and reciprocal coordinates give all node,
mark and infinity charts, including generators whose pullback
vanishes identically on a contracted tail. This proves
compatibility with every mark-forgetting contraction.

The labelled component maps lift uniquely once these ghost maps
and the log base are fixed. At a node a possible difference of
two lifts multiplies the branch generators by units $u_1,u_2$
with $u_1u_2=1$, $u_1x=x$, and $u_2y=y$. Thus
$u_1-1\in(x)\cap(y)=0$ in the reduced nodal ring. On the
smooth locus the vertical log structure comes from the fixed
base, while the horizontal generator is a non-zero-divisor.
Initial forms commute with finite extensions of valued fields,
so the comparisons commute with those extensions; their unique
labelled description and equivariance of $F$ give Galois descent.
At a mark $a$, the horizontal parameter changes by $\lambda$
or one. In particular $K(F(a))=K(a)$ as actual subfields,
the cusp inertia is preserved, and the quotient $G_{K(a)}$ is
identified by the identity. This is the same local Kummer-chart
calculation as in \statementref{B.48}; principal units have all
their roots by Hensel's lemma, since every integer is invertible.

Apply \statementref{B.54} to these split models and descend.
It gives compatible equivalences between the finite \'etale
cover categories of $\mathbb P^1_K\setminus S$ and
$\mathbb P^1_K\setminus F(S)$ over the entire $BG_K$.
They respect all inclusions of marked sets and the actual
arithmetic cusp residue extensions. We recall why this transports
the whole small \'etale site, as required in
\statementref{B.49}. A quasi-compact separated \'etale object
over $\mathbb A^1_K$ is an open subscheme of the normalization
in its generic finite algebra. Its data consist of that finite
algebra, its branching over finitely many closed points, and the
retained unramified points with their residue extensions.
The cover equivalences and cusp comparisons carry all these
data. Maps extend uniquely from their generic algebra maps,
with the condition on retained points preserved. Geometric
point surjectivity, which detects covering families, is preserved
as well. Such objects form a basis of the small \'etale site.
The compatibility already proved therefore yields a pinned
topos self-equivalence over $BG_K$, whose function on
$\overline K$-valued points is (B.LIT.3).

Residue characteristic zero implies $v(q)=0$ for every
$q\in\mathbb Q^\times$. Thus $F$ fixes all rational numbers
and moves $t$. No $K$-polynomial has that point function: a
polynomial fixing infinitely many rational numbers is the
identity. The constructed topos morphism is consequently not
induced by a $K$-scheme morphism. This proves failure of
reconstruction in rank one.

Finally let $\operatorname{trdeg}_{\mathbb Q}K=d<\infty$
and let the original valuation have residue characteristic zero.
It is trivial on $\mathbb Q$. If
$v(a_1),\ldots,v(a_m)$ are rationally independent, distinct
monomials in the $a_i$ have distinct valuations, so no nonzero
$\mathbb Q$-polynomial in them can vanish. Hence $m\leq d$:
the value group has finite rational rank. Any strict chain of
convex subgroups gives rationally independent elements chosen
in successive differences; the highest nonzero coefficient in
a rational relation cannot vanish modulo the preceding convex
subgroup. The ordered group therefore has finite rank and a
maximal proper convex subgroup. Its quotient is a nonzero
rank-one ordered group. Coarsen by this subgroup.
ADH Lemma~3.4.2 states that a valued field is henselian if and
only if both a coarsening and its induced residue valuation are
henselian. The coarsening is therefore henselian; its residue
characteristic is still zero because $\mathbb Q^\times$
remains in its unit group. The rank-one result applies to the
same field $K$.\citep[Lemma~3.4.2]{ADH2017}
\end{proof}
\paragraph{Check record.} Also checked in this writing, 2026-11-18: GPT-6 Astra (OpenAI, ChatGPT Pro mode), root; separate effort setting not exposed. Full targeted reading of ordinary finite-transcendence pro-reconstruction, all formal-arc avoidance cases, the henselian centralizer example, stable marked logarithmic diagrams and uniqueness, full finite-cover-to-small-site gluing, and arbitrary rank-one logarithmic approximation. Configuration reader separately opened all exact primary inputs. Inverse-image orientation is made explicit; the fixed log point retains the entire arithmetic Galois group.


This extension removes the nondiscrete residue-characteristic-zero
cases from the unresolved portion of B8. It does not replace the
all-prime comparison by a tame quotient in mixed characteristic.
Fields admitting no nontrivial henselian residue-characteristic-zero
coarsening require a separate argument.


% Configuration audit continuation: exact mixed-characteristic frontiers.
% AI-produced by GPT-6 Astra; provisional statement identifiers B.56--B.57.
\subsection{Wild Kummer obstructions and the unbounded nilpotent tower}

Throughout this subsection,
\[
 K_p=\operatorname{Frac}(\Z_{(p)}^{sh})\subset\overline\Q,
 \qquad I=G_{K_p}=I_{\Q_p},\qquad
 F_p=\overline\Q^{I'},\qquad B_p=\overline\Q^{I''}.
                                                        \tag{B.MIX.1}
\]
The identification with local inertia is the henselization--completion
comparison, with the chosen prolongation of the valuation. It concerns
absolute Galois groups; the global field \(K_p\) has not been replaced
by \(\Q_p^{\mathrm{ur}}\). For clarity, a finite collection of global
algebraic numbers and its local Galois conjugates lie in a number field
\(D\) whose completion \(M\) is finite Galois over \(\Q_p\).
Then \(I\cap G_M=I_M\) fixes the compositum \(K_pD\).
These are the identifications used below.

\begin{claim}{B.56}{Proved}{Wild covers obstruct supported scaling at the metabelian field}
For a finite extension \(E/B_p\), let \(\kappa_E\) be the rational
\(p\)-primary Kummer map. Put
\[
 W_p=\{a\in B_p^\times:a^m\in F_p^\times
                         \text{ for some }m\geq1\}.
\]
Then
\[
 \ker\kappa_E=\bigcap_{n\geq1}(E^\times)^{p^n}=W_p,
 \qquad W_p\cap(1-W_p)=F_p\setminus\{0,1\}.             \tag{B.MIX.2}
\]
In particular the supported scaling used in B.48--B.49 cannot be the
point function of a pinned relative autoequivalence over \(B_p\).
Precisely, for \(\lambda=1+p\), the bijection
\[
 f(a)=\begin{cases}\lambda a,&v(a)>0,\\a,&v(a)\leq0,
                         \end{cases}\qquad f(0)=0
                                                        \tag{B.MIX.3}
\]
of \(\overline\Q\) is excluded by two one-puncture
\(p\)-Kummer comparisons. This is an obstruction to this proposed
counterexample, and not a proof of reconstruction over \(B_p\).
\end{claim}
\begin{proof}
We independently reprove the kernel input of AS-FF.42. Suppose
\(\kappa_E(a)=0\). Enlarge to \(E'=B_pD\), with \(D\) a number
field stable under the local Galois group, containing \(a\) and the
finitely many generators for \(E/B_p\). Its local completion
\(M/\Q_p\) is finite Galois. Write
\[
 J=I\cap G_M=I_M,\quad H=I''\cap J,
 \quad\Gamma=J/H,\quad C=(I'\cap J)/H.
\]
The Kummer class of \(a\) on \(J\) inflates from
\(H^1(\Gamma,\Q_p(1))\). The group \(C\) and the quotient
\(\Gamma/C\) are abelian; the cyclotomic character is trivial
on \(C\) and has infinite image on \(\Gamma/C\).
For an abelian group \(A\) with nontrivial scalar character \(\chi\),
the cocycle equation gives
\((\chi(g)-1)c(h)=(\chi(h)-1)c(g)\). Choosing one
\(g\) with \(\chi(g)\ne1\) proves \(H^1(A,\Q_p(\chi))=0\).
Inflation--restriction therefore injects
\(H^1(\Gamma,\Q_p(1))\) into \(H^1(C,\Q_p(1))\).
Every \(g\in I'\) centralizes \(C\), because the commutator
lies in \(I''\cap J\), and fixes its coefficient module.
Thus it fixes the class, so
\(\kappa_J(g(a)/a)=0\).

Both \(a\) and \(g(a)\) belong to \(M\). Restriction from
\(H^1(G_M,\Q_p(1))\) to \(H^1(I_M,\Q_p(1))\) is injective,
since \(\Q_p(1)^{I_M}=0\). Over the finite local field \(M\),
the rational Kummer kernel is exactly \(\mu(M)\): valuation
detects the uniformizer component, and the logarithm identifies a
sufficiently small principal-unit group with an additive lattice
in \(M\). Hence \(g(a)/a\in\mu(M)\) for every \(g\in I'\).
An integer \(m\) killing the finite group \(\mu(M)\) makes
\(a^m\) fixed by \(I'\), so \(a^m\in F_p\).
Since \(F_p\) contains all roots of unity,
\(F_p(a)/F_p\) is cyclic and belongs to its maximal abelian
extension \(B_p\). Thus \(a\in W_p\).

Conversely, if \(a^m\in F_p\), every \(p^n\)-th root of
\(a\) is a radical of an element of \(F_p\). Its extension
over \(F_p\) is cyclic, so it belongs to \(B_p\).
Compatible roots exist. Since \(G_E\) fixes \(\mu_{p^\infty}\),
integral Kummer classes are homomorphisms to the torsion-free
module \(\Z_p(1)\); rational vanishing is the same as integral
vanishing. This proves the first equality in (B.MIX.2).
If \(a,1-a\in W_p\), then for each \(g\in I'\) one has
\(g(a)=\zeta a\), \(g(1-a)=\eta(1-a)\), with
\(\zeta,\eta\in\mu_\infty\subset F_p\).
If \(g(a)\ne a\), their sum gives
\(a=(1-\eta)/(\zeta-\eta)\in F_p\), a contradiction;
the case \(\zeta=\eta\) forces both to be one.
This proves the second equality.

Choose \(x\in B_p\setminus F_p\) with \(v(x)>0\), and
\(c\in K_p\) with both \(c\) and \(1-c\) units. Such an \(x\) exists by the following elementary argument.
The commutator group \(I'\) is pro-\(p\), since tame inertia
is abelian. It is nontrivial. For odd \(p\), the splitting
field of \(X^p-p\) over \(K_p\) has nonabelian group
\(C_p\rtimes(\mathbf Z/p)^\times\): Eisenstein gives the
radical degree \(p\), the cyclotomic degree is \(p-1\),
and the scalar action is faithful. For \(p=2\), put
\(\beta=2^{1/4}\). The polynomial \(X^4-2\) is Eisenstein,
so \([K_2(\beta):K_2]=4\). The distinct quadratic fields
\(K_2(\sqrt2)\) and \(K_2(i)\) have compositum of degree
four: their equality would make \(-2\) a square in \(K_2\),
contrary to its odd valuation. This compositum does not contain
\(\beta\). Indeed, writing \(\beta=a+bi\) with
\(a,b\in K_2(\sqrt2)\) and squaring would force \(ab=0\);
then \(\sqrt2\) or \(-\sqrt2\) would be a square in that
field, contrary to its value group \(\tfrac12\mathbf Z\).
Hence \(i\notin K_2(\beta)\), and the splitting field has
the nonabelian dihedral group of order eight. Every nontrivial
pro-\(p\) group has a nontrivial finite \(p\)-group quotient,
and then a quotient of order \(p\). Thus \(I'/I''\ne1\),
so \(B_p\ne F_p\). Multiplication by a large power of \(p\)
makes an element of \(B_p\setminus F_p\) have positive value. The residue
field \(\overline\F_p\) supplies \(c\), including when \(p=2\).
Put \(y=1+c(x-1)\). Then \(y\) is a unit,
\(f(1)=1\), \(f(y)=y\), and \(f(x)=\lambda x\).

Suppose the stated relative autoequivalence existed. On the line
punctured at a point \(a\in B_p\), it sends the rank-one geometric
\(p\)-Kummer line to itself by a scalar
\(\alpha_a\in\Z_p^\times\). The remaining summand is an
arithmetic Kummer class, which cancels on subtracting two point
evaluations. The puncture at \(1\), evaluated at \(y,x\), gives
\[
 \kappa_{B_p}\!\left(\frac{y-1}{\lambda x-1}\right)
     =\alpha_1\kappa_{B_p}\!\left(\frac{y-1}{x-1}\right)
     =\alpha_1\kappa_{B_p}(c)=0.
\]
The puncture at \(x\), evaluated at \(y,1\), similarly gives
\[
 \kappa_{B_p}\!\left(\frac{y-\lambda x}{1-\lambda x}\right)
     =\alpha_x\kappa_{B_p}(1-c)=0.
\]
The two left arguments are \(A\) and \(1-A\), for
\(A=c(x-1)/(\lambda x-1)\). Equation (B.MIX.2) forces
\(A\in F_p\). But this fractional linear transformation has
coefficients in \(F_p\) and nonzero determinant
\(c(\lambda-1)\), so it cannot carry
\(x\notin F_p\) into \(F_p\). This contradiction pinpoints the
wild arithmetic obstruction absent from the residue-zero proof.
\end{proof}
\paragraph{Check record.} Also checked in this writing, 2026-11-18: GPT-6 Astra (OpenAI, ChatGPT Pro mode), root; separate effort setting not exposed. Read complete proofs: ordinary pro-form consequence, solvable strict-local closure via exact tame/wild sequence, metabelian stable Kummer kernel and complementary-unit intersection, nilpotent Kummer theorem including p2/p3 exceptions, rational residues, affine similarity and polynomial point maps. B56 nontriviality is now proved directly by odd-prime affine groups and the p2 dihedral splitting field, without wild freeness. B57 reconstruction scope updated by B106–110.



\begin{claim}{B.57}{Proved}{A prime-dependent Kummer theorem for the unbounded nilpotent tower}
Let
\[
 N=\bigcap_{c\geq1}\gamma_{c+1}(I),\qquad
 K_{p,\mathrm{nil}}=\overline\Q^N
      =\bigcup_{c\geq1}\overline\Q^{\gamma_{c+1}(I)}.
                                                        \tag{B.MIX.4}
\]
For every \(p\geq5\) and every finite extension
\(E/K_{p,\mathrm{nil}}\),
\[
 \ker\!\left(E^\times\longrightarrow H^1(G_E,\Q_p(1))\right)
       =\bigcap_{n\geq1}(E^\times)^{p^n}
       =\mu_\infty.                                    \tag{B.MIX.5}
\]
For \(p=2\) and \(p=3\), the corresponding kernel is already
strictly larger than \(\mu_\infty\) over
\(K_{p,\mathrm{nil}}\) itself. Explicit nontorsion elements in it are
\[
       2\quad(p=2),\qquad
       \frac{2+\sqrt{-3}}{2-\sqrt{-3}}\quad(p=3).         \tag{B.MIX.6}
\]
These conclusions decide this exact Kummer input to AS-FF.39 at its
unbounded union. B.138 below proves reconstruction at that union and every finite extension for every prime, including $2,3$; the kernel assertions here are retained unchanged.
\end{claim}
\begin{proof}
Write \(K=K_{p,\mathrm{nil}}\) and \(Q=I/N\).
The quotient \(Q\) is the maximal pronilpotent quotient of
\(I\): a finite quotient is nilpotent exactly when it kills
some lower-central term, and any finite compositum of such
quotients is again nilpotent. Its pro-\(p\) Sylow subgroup
\(Q_{(p)}\) is normal, and
\[
                    Q=Q_{(p)}\times T,
                                                        \tag{B.MIX.7}
\]
where \(T\) is the tame prime-to-\(p\) procyclic group.
Indeed every finite quotient of \(I\) has a normal wild
\(p\)-subgroup and cyclic prime-to-\(p\) quotient, by the
inertia sequence verified in B.53. In a nilpotent quotient
the different Sylow subgroups commute. Taking inverse limits
gives (B.MIX.7); the full tame quotient survives, since it
is already abelian. All roots of unity belong to \(K\).
For odd \(p\), the cyclotomic character on \(T\) surjects
onto \(\mu_{p-1}\subset\Z_p^\times\), and its pro-\(p\)
part comes from \(Q_{(p)}\).

Let \(a\in E^\times\) have zero rational Kummer class.
As above, \(G_E\) fixes the coefficient module, so this is
equivalent to having all compatible \(p\)-power roots of
\(a\) in \(E\). Enlarge to a finite extension \(E'/K\)
normal under the local absolute Galois group, containing the
finite local orbit \(a_1=a,\ldots,a_m\). This is possible
by adjoining the finitely many local conjugates of generators
for \(E/K\); the field \(K\) is stable under \(G_{\Q_p}\),
since \(N\) is characteristic in its normal inertia subgroup.
Choose a number field \(D\subset E'\), stable under
\(G_{\Q_p}\), containing these generators and conjugates,
and let \(M=D_v\), a finite Galois extension of \(\Q_p\).
Put
\[
 F=K_pD(\mu_{p^\infty}),\qquad
 R=F(a_j^{1/p^n}:1\leq j\leq m,\ n\geq1).
\]
Then \(R/K_p\) is Galois and \(R\subset E'\).
The group \(D_R=\operatorname{Gal}(R/F)\) embeds by
Kummer characters into \(\Z_p(1)^m\); in particular it is
torsion-free. The image \(C_R\) of \(N\) in
\(\operatorname{Gal}(R/K_p)\) is finite, because
\(E'/K\) is finite.

Choose \(t\in T\) whose cyclotomic character is a primitive
\((p-1)\)-st root \(\zeta\), and choose a lift
\(\tau\in I\) of \((1,t)\in Q_{(p)}\times T\).
The cyclotomic character factors through \(Q\), so
\(\chi(\tau)=\zeta\) exactly. Its image is central modulo
\(C_R\). Consequently
\[
             [\tau,D_R]\subset C_R\cap D_R=1.           \tag{B.MIX.8}
\]
Here the last equality uses finite versus torsion-free, not
an assertion that the entire radical extension is abelian
over \(K_p\).

Set \(J=I_M=I\cap G_M\) and \(J_0=J\cap\ker\chi\).
The group \(J_0\) fixes \(K_p,D,\mu_{p^\infty}\) and
maps onto \(D_R\). Kummer conjugation and (B.MIX.8) give
\[
       \kappa_{J_0}(\tau a)=\zeta\,\kappa_{J_0}(a).
                                                        \tag{B.MIX.9}
\]
There is no lost subgroup here: \(J\) already fixes the
strict-local field \(K_p\). Restriction
\(H^1(J,\Q_p(1))\to H^1(J_0,\Q_p(1))\) is injective,
because \(J/J_0\) is an abelian group with nontrivial
cyclotomic character and its first cohomology vanishes by
the displayed cocycle calculation in B.56. Restriction
from \(G_M\) to \(I_M\) is injective as well, since the
inertia invariants of \(\Q_p(1)\) vanish. Both \(a\) and
\(\tau a\) belong to \(M\), so (B.MIX.9) implies
\[
                \kappa_M(\tau a)=\zeta\kappa_M(a).
                                                        \tag{B.MIX.10}
\]

Valuation in (B.MIX.10) gives \(v(a)=\zeta v(a)\);
as \(\zeta\ne1\), the element \(a\) is a unit.
The principal-unit logarithm then gives
\(\log(\tau a)=\zeta\log(a)\).
The opened Baker--Brumer theorem, with proof in
\citep[Théorème, pp.~1--2]{Rousseau1969}, states that
\(p\)-adic logarithms of algebraic units which are linearly
independent over \(\Q\) remain linearly independent over
\(\overline\Q\). If \(\log(a)\ne0\), the preceding
algebraic relation therefore makes the two logarithms
\(\Q\)-dependent, and forces \(\zeta\in\Q\).
For \(p\geq5\), a primitive \((p-1)\)-st root is not
rational. Thus \(\log(a)=0\); the unit \(a\) in the
finite local field \(M\) is a root of unity. Conversely
all roots of unity have all their divisions in \(K\).
This proves (B.MIX.5), with every finite extension included.

For \(p=2\), each splitting field over \(K_2\) of
\(X^{2^n}-2\) has a finite \(2\)-group as Galois group:
adjoin the \(2^n\)-th roots of unity and then a Kummer
root. These are nilpotent groups, so all those roots lie
in \(K_{2,\mathrm{nil}}\), proving the first exception.

For \(p=3\), put \(z=\sqrt{-3}\),
\(F_0=K_3(z)=K_3(\mu_3)\), and
\(a=(2+z)/(2-z)\). The nonidentity element of
\(\operatorname{Gal}(F_0/K_3)\) sends \(a\) to \(a^{-1}\).
Let \(R_n=F_0(\mu_{3^n},b_n)\), where \(b_n^{3^n}=a\).
This is Galois over \(K_3\), and its kernel over \(F_0\)
is a \(3\)-group \(P_n\). Cauchy's theorem supplies an
involution \(s\) in the Galois group. It inverts
\(\mu_{3^n}\). Writing \(s(b_n)=\xi b_n^{-1}\), its
order two forces \(\xi^2=1\), hence \(\xi=1\).
For \(g\in P_n\), write
\(g(b_n)=\eta b_n\) and \(g(\zeta_{3^n})=\zeta_{3^n}^{u}\).
Direct substitution on these two generators gives
\(sgs^{-1}=g\). Thus
\(\operatorname{Gal}(R_n/K_3)=P_n\times\langle s\rangle\)
is nilpotent. All \(3\)-power roots of \(a\) lie in
\(K_{3,\mathrm{nil}}\). Finally
\(a+a^{-1}=2/7\), which is not an algebraic integer,
so \(a\) is not a root of unity. This proves the second
exception and completes the stated Kummer classification.
\end{proof}
\paragraph{Check record.} Also checked in this writing, 2026-11-18: GPT-6 Astra (OpenAI, ChatGPT Pro mode), root; separate effort setting not exposed. Read complete proofs: ordinary pro-form consequence, solvable strict-local closure via exact tame/wild sequence, metabelian stable Kummer kernel and complementary-unit intersection, nilpotent Kummer theorem including p2/p3 exceptions, rational residues, affine similarity and polynomial point maps. B56 nontriviality is now proved directly by odd-prime affine groups and the p2 dihedral splitting field, without wild freeness. B57 reconstruction scope updated by B106–110.


\paragraph{Why the new Kummer theorem does not invoke an existing reconstruction theorem.}
The amenability hypothesis in
\citep[Definition~5.15(5.15.1)]{CarlsonHaineWolf2024}
requires each torus to have no nonzero divisible rational
points, with all primes included in characteristic zero.
It fails already for \(\mathbb G_m\), since
\(\mu_\infty\subset K_{p,\mathrm{nil}}\).
The independent condition~(5.15.2), requiring the Tate module
of \(\operatorname{Pic}(X)\) to vanish for every regular
finite-type \(X\), fails as well. For an explicit example take
\(A:y^2+xy=x^3-p\), with \(c_4=1\) and
\(\Delta=p(1-432p)\). Its node has the two rational tangent
lines \(y=0\), \(y+x=0\), hence split multiplicative reduction.
The opened Tate uniformization and torsion sequence
\citep[Theorem~1, author p.~6; pp.~18--19]{Tate1995}
give a compatible torus Tate line \(\Z_p(1)\subset T_pA\).
The subgroup \(I'\) fixes this line pointwise. Torsion points
of the globally defined elliptic curve are algebraic over
\(\Q\), so their coordinates lie in \(F_p\), and consequently
in \(K_{p,\mathrm{nil}}\). Their divisor classes give
\(\Z_p\subset T_p\operatorname{Pic}(A_{K_{p,\mathrm{nil}}})\).
Thus even the Picard condition has an actual counterexample;
it is not merely absent from the proof of (B.MIX.5).

The explicit positive argument for a fixed bounded nilpotency
class in AS-FF.39 also uses an inertia-invariant lift of every
visible generalized-Albanese functional. Passing to a
pronilpotent quotient does not retain its bounded-nilpotency
Lie-algebra argument: arbitrary pro-\(p\) groups are
pronilpotent. Formula (B.MIX.5) alone supplies neither that
lift nor compatible morphisms on all finite étale covers.
The separate reconstruction argument is supplied for every prime in B.137--B.138 below. The larger kernels at \(p=2,3\)
likewise refute a particular input, not reconstruction.
No single-topos reconstruction conclusion is deduced from
ordinary pro-reconstruction.

\subsection{Affine point functions over the unbounded tower when \texorpdfstring{\(p\geq5\)}{p >= 5}}

\begin{claim}{B.58}{Proved}{An actual Kummer value detects a rational scalar}
Let \(p\geq5\), \(K=K_{p,\mathrm{nil}}\), and let \(E/K\)
be finite. If \(b\in E^\times\) and \(\alpha\in\Q_p\)
satisfy
\[
                 \kappa_E(b)=\alpha\kappa_E(p),           \tag{B.MIX.11}
\]
then \(\alpha\in\Q\).
\end{claim}
\begin{proof}
Pass to a finite extension \(E'/K\) normal under
\(G_{\Q_p}\), as in B.57, and put \(N'=G_{E'}\).
It is normal in \(I\), contained in \(N=G_K\), and of
finite index in \(N\). For \(g\in I\), conjugate (B.MIX.11)
on \(N'\). The restricted Kummer class of the rational
number \(p\) is \(I\)-invariant, since it comes from
\(H^1(I,\Q_p(1))\). Hence
\(\kappa_{E'}(g(b)/b)=0\). B.57 makes every
\(g(b)/b\) a root of unity. The finite local orbit of
the globally algebraic number \(b\) lies in one number
field, whose roots of unity form a finite group. A single
integer \(m>0\) therefore gives
\(b^m\in\overline\Q^I=K_p\).

The class
\(\kappa_I(b^m)-m\alpha\kappa_I(p)\) restricts to zero
on \(N'\), so it inflates from \(H^1(I/N',\Q_p(1))\).
This group is zero. Indeed the kernel \(N/N'\) of
\(I/N'\to I/N=Q\) is finite and fixes the coefficient
module, so its rational first cohomology is zero.
Inflation--restriction identifies the group with
\(H^1(Q,\Q_p(1))\). The group \(Q\) has a central
tame element \(t\) with \(\chi(t)\ne1\).
For every cocycle \(c\), centrality gives
\((\chi(t)-1)c(g)=(\chi(g)-1)c(t)\), so \(c\) is a
coboundary. This proves the claimed vanishing.

The element \(b^m\in K_p\) lies in a finite unramified
local extension \(M/\Q_p\), with \(I_M=I\).
The local Kummer class comparison is injective upon
restriction to inertia, because \(\Q_p(1)^I=0\).
Thus the preceding zero class already vanishes on
\(G_M\). Normalizing its valuation by \(v_M(p)=1\)
now yields
\(m\alpha=v_M(b^m)\in\Z\). Hence \(\alpha\) is rational.
\end{proof}
\paragraph{Check record.} Also checked in this writing, 2026-11-18: GPT-6 Astra (OpenAI, ChatGPT Pro mode), root; separate effort setting not exposed. Read complete proofs: ordinary pro-form consequence, solvable strict-local closure via exact tame/wild sequence, metabelian stable Kummer kernel and complementary-unit intersection, nilpotent Kummer theorem including p2/p3 exceptions, rational residues, affine similarity and polynomial point maps. B56 nontriviality is now proved directly by odd-prime affine groups and the p2 dihedral splitting field, without wild freeness. B57 reconstruction scope updated by B106–110.


\begin{claim}{B.59}{Proved}{Affine actual point functions for all pinned affine-line automorphisms}
Let \(p\geq5\), and let \(E\) be any finite extension
of \(K_{p,\mathrm{nil}}\). Every pinned autoequivalence
of \(\mathbb A^1_{E,\mathrm{et}}\) over \(BG_E\), with
its relative structural isomorphism retained, has actual
geometric point function
\[
                         a\longmapsto ua+v,
                 \qquad u\in E^\times,\ v\in E.          \tag{B.MIX.12}
\]
The coefficients are unique. This conclusion concerns the
actual point function; it does not identify the remaining
pointwise-trivial automorphism kernel on the full étale site,
or assert existence for arbitrary noninvertible pinned maps.
\end{claim}
\begin{proof}
Let \(f:\overline\Q\to\overline\Q\) be the actual point
function. Relative point evaluation
\citep[Proposition~3.6]{CarlsonHaineWolf2024}, together with
the inverse equivalence, makes it a \(G_E\)-equivariant
bijection. On every finite extension \(L/E\), it gives a
bijection of \(L\). Compose with the unique affine scheme
automorphism taking \(f(0),f(1)\) to \(0,1\); henceforth
\(f(0)=0\), \(f(1)=1\).

For \(a\in\overline\Q\), the equivalence after a finite
scalar extension restricts to the lines punctured at
\(a\) and \(f(a)\). On their rank-one geometric
\(p\)-Kummer lines it acts by
\(\alpha_a\in\Z_p^\times\). The arithmetic summand
cancels when two point evaluations are subtracted. Thus,
over any common finite extension containing \(a,b,x\),
\[
 \kappa_L\!\left(\frac{f(x)-f(a)}{f(b)-f(a)}\right)
       =\alpha_a\kappa_L\!\left(\frac{x-a}{b-a}\right)
       \qquad (b,x\ne a).                              \tag{B.MIX.13}
\]
This only uses the one-puncture cohomology calculation;
no preservation theorem for the place at infinity and
no equality of different inertia exponents is assumed.
Set \(b=a+1\), \(x=a+p\). B.58 applies to the left
argument in (B.MIX.13), and proves
\(\alpha_a\in\Q^\times\) for every \(a\).

Fix a complex embedding of \(\overline\Q\), and write
\(|\cdot|\) for its ordinary complex absolute value.
Clear the rational denominator of \(\alpha_a\) in
(B.MIX.13). B.57 says that the resulting multiplicative
error is a root of unity. Its complex absolute value is one.
Consequently there is a positive real constant \(C_a\)
such that
\[
                 |f(x)-f(a)|=C_a|x-a|^{\alpha_a}
                       \qquad(x\ne a).                  \tag{B.MIX.14}
\]
One can take \(C_a=|f(a+1)-f(a)|\).

We spell out how (B.MIX.14) forces the exponents and constants
to coincide. If \(\alpha_a<0\), the sequence \(f(n)\), for
positive rational integers \(n\to\infty\), tends to \(f(a)\).
For another centre \(b\ne a\), its distance from \(f(b)\)
would therefore tend to the finite positive number
\(|f(a)-f(b)|\). The formula at \(b\) makes that distance
tend to either zero or infinity, since \(\alpha_b\ne0\).
This contradiction proves \(\alpha_a>0\) for every \(a\).
The reverse triangle inequality gives
\[
 \left|C_a|n-a|^{\alpha_a}-C_b|n-b|^{\alpha_b}\right|
                      \leq |f(a)-f(b)|.
\]
Divide by the larger power of \(n\) and let \(n\to\infty\).
Unequal positive exponents give a contradiction, so all
\(\alpha_a\) equal one positive number \(\alpha\).
Dividing again by \(n^\alpha\) gives \(C_a=C_b\).
Normalization at zero and one makes the common constant one.
Thus
\[
                       |f(x)-f(y)|=|x-y|^\alpha
                       \qquad(x,y\in\overline\Q).       \tag{B.MIX.15}
\]
The triangle with vertices \(0,1,2\) gives
\(2^\alpha\leq2\), hence \(\alpha\leq1\).
Since \(f\) is bijective, its inverse satisfies the same
identity with exponent \(1/\alpha\), and gives
\(1/\alpha\leq1\). Hence \(\alpha=1\).

An isometry of \(\overline\Q\subset\mathbb C\) fixing
zero and one is the identity or complex conjugation.
Indeed distances to zero and one give \(f(i)=i\) or
\(f(i)=-i\); distances to the three noncollinear points
\(0,1,i\) then determine the image of every point, with
the second case giving the reflected point.

The conjugation case cannot be \(G_E\)-equivariant.
First \(E\) is not separably closed. For \(p\geq5\),
the splitting field of \(X^p-p\) over \(K_p\) has
non-nilpotent Galois group
\(\mathbb F_p\rtimes\mathbb F_p^\times\): the radical
degree is \(p\), the cyclotomic tame degree is \(p-1\),
and the natural scalar action is faithful. Thus \(N\ne1\).
It is a closed subgroup of the free wild pro-\(p\) group,
so is torsion-free and infinite; every finite-index subgroup
\(G_E\) is nontrivial. The field \(E\) carries its
nontrivial rank-one henselian mixed-characteristic valuation.
If complex conjugation centralized \(G_E\), it would
preserve \(E\). The opened
\citep[Lemma~2.1]{KhassaKhanduja2010} makes this automorphism
continuous for the valuation; in rank one it preserves the
valuation ring. Uniqueness of prolongation then makes it
preserve the chosen valuation on \(\overline\Q\), so it
lies in the local decomposition group \(G_{\Q_p}\).
Its element of order two maps trivially to the unramified
quotient \(\widehat\Z\), then trivially to the torsion-free
tame inertia quotient. It would lie in wild pro-\(p\)
inertia, which has no element of order two for \(p\geq5\).
This is impossible. The normalized \(f\) is therefore the
identity. Undoing normalization gives (B.MIX.12).
\end{proof}
\paragraph{Check record.} Also checked in this writing, 2026-11-18: GPT-6 Astra (OpenAI, ChatGPT Pro mode), root; separate effort setting not exposed. Read complete proofs: ordinary pro-form consequence, solvable strict-local closure via exact tame/wild sequence, metabelian stable Kummer kernel and complementary-unit intersection, nilpotent Kummer theorem including p2/p3 exceptions, rational residues, affine similarity and polynomial point maps. B56 nontriviality is now proved directly by odd-prime affine groups and the p2 dihedral splitting field, without wild freeness. B57 reconstruction scope updated by B106–110.


\begin{claim}{B.60}{Proved}{Actual Kummer values force rational boundary residues}
Let \(p\geq5\), let \(E/K_{p,\mathrm{nil}}\) be finite,
and let \(U\) be a smooth geometrically connected curve over
\(E\), with smooth proper compactification \(\overline U\).
Suppose \(\xi\in H^1(U,\Z_p(1))\) evaluates at every point
\(x\in U(L)\), for every finite extension \(L/E\), to the
Kummer class of an actual element of \(L^\times\).
Every geometric boundary residue of \(\xi\) belongs to
\(\Q\cap\Z_p\).
\end{claim}
\begin{proof}
Split the finitely many boundary points over a finite extension
and fix one of them, \(q\). Riemann--Roch supplies a nonconstant
rational function \(g\) with its only pole at \(q\) and with a
zero at every other boundary point: prescribe those finitely
many zeros and allow a sufficiently large pole at \(q\).
The map \(g:\overline U\to\mathbb P^1\) is finite.
Remove zero, infinity, and all branch values from the target,
obtaining \(V\); its inverse image \(U'\) lies in \(U\)
and is finite étale over \(V\). Corestrict \(\xi|_{U'}\)
to a class \(\eta\in H^1(V,\Z_p(1))\).

Its geometric residues vanish away from zero and infinity.
At a removed branch value the original class has zero
residues at all its preimages. The corestriction residue
is the sum of source residues multiplied by residue-field
degrees, not by ramification indices. This formula follows
at finite Kummer levels from valuation of the field norm,
and then by passage to the compatible inverse limit.
The residue at infinity is precisely the residue at the
sole point \(q\), since \(q\) is rational over the chosen
base. Thus the class has the form
\(\eta=\eta_0+\alpha\kappa(T)\), with \(\eta_0\) arithmetic
and \(\alpha=-\operatorname{res}_q\xi\). This also follows
directly from the Kummer sequence on the punctured rational
line, whose Picard group is zero; zero residues at the
additional punctures remove their basis terms.

At every \(t\in V(E)\), the evaluation of \(\eta\) is
an actual Kummer value. Indeed the fibre is a product of
finite separable extensions, and corestriction evaluates
as the product of the field norms of the actual values
of \(\xi\) on these points. Choose \(t_0\in V(E)\) and
a positive integer \(n\) for which \(p^nt_0\in V(E)\);
only finitely many choices are excluded. Subtracting the
two evaluations cancels \(\eta_0\) and gives an actual
value \(b\in E^\times\) with
\(\kappa_E(b)=n\alpha\kappa_E(p)\).
B.58 gives \(n\alpha\in\Q\), hence the desired
rationality of the residue. The argument works for each
boundary point and is unchanged by the finite scalar
extension used to split them.
\end{proof}
\paragraph{Check record.} Also checked in this writing, 2026-11-18: GPT-6 Astra (OpenAI, ChatGPT Pro mode), root; separate effort setting not exposed. Read complete proofs: ordinary pro-form consequence, solvable strict-local closure via exact tame/wild sequence, metabelian stable Kummer kernel and complementary-unit intersection, nilpotent Kummer theorem including p2/p3 exceptions, rational residues, affine similarity and polynomial point maps. B56 nontriviality is now proved directly by odd-prime affine groups and the p2 dihedral splitting field, without wild freeness. B57 reconstruction scope updated by B106–110.


\begin{claim}{B.61}{Proved}{Polynomial actual point functions for arbitrary affine-line morphisms}
Let \(p\geq5\), and let \(E\) be any finite extension
of \(K_{p,\mathrm{nil}}\). Every pinned geometric morphism
\[
       \mathbb A^1_{E,\mathrm{et}}
              \longrightarrow\mathbb A^1_{E,\mathrm{et}}
\]
over \(BG_E\) has actual geometric point function
\(x\mapsto P(x)\) for a unique polynomial \(P\in E[T]\).
Invertibility of the geometric morphism is not assumed.
This does not identify the whole inverse-image functor
with the functor induced by \(P\), and does not prove
existence for an arbitrary source scheme.
\end{claim}
\begin{proof}
Write \(f:\overline\Q\to\overline\Q\) for the actual
point function supplied by relative pinning. If the
underlying generic point maps to a closed point, continuity
makes the entire image that closed point. Evaluation at
zero makes it an \(E\)-rational point, and the assertion
holds with a constant polynomial. Otherwise the generic
point maps to the generic point. Each geometric fibre
\(S_a=f^{-1}(a)\) is then finite: after a finite constant
extension containing \(a\), it is a proper closed subset
of the source affine line. In particular the image of
\(f\) is infinite.

For \(a\in\overline\Q\), restrict to
\(U_a=\mathbb A^1\setminus S_a\) and pull back the
Kummer class of \(T-a\). Split \(S_a\) after a finite
extension. The resulting class is an arithmetic class
plus
\(\sum_{r\in S_a}q_{a,r}\kappa(T-r)\), because a
localization of the polynomial ring has trivial Picard
group. B.60 gives \(q_{a,r}\in\Q\cap\Z_p\).
Subtract the evaluation at any fixed \(x_0\notin S_a\),
and clear these finitely many rational denominators.
B.57 identifies the resulting multiplicative error
with a root of unity, over every finite field of
point definition. In a fixed complex embedding, this gives
\[
 |f(x)-a|^2=P_a(x),\qquad
 P_a(z)=C_a^2\prod_{r\in S_a}|z-r|^{2q_{a,r}},
 \quad C_a>0,
 \qquad x\in\overline\Q\setminus S_a.                   \tag{B.MIX.16}
\]
No sign or integrality of the exponents has been assumed.
The function \(P_a\) is positive and real analytic off
\(S_a\), and \(\log P_a\) is harmonic there.

Put \(S=S_0\cup S_1\cup S_i\) and
\(\Omega=\mathbb C\setminus S\). Define a real-analytic
function on this connected open set by
\[
 F(z)=\frac{P_0(z)-P_1(z)+1}{2}
          +i\frac{P_0(z)-P_i(z)+1}{2}.                  \tag{B.MIX.17}
\]
The elementary squared-distance identities show that
\(F(x)=f(x)\) at every algebraic \(x\in\Omega\).
Density and continuity then give
\(|F(z)-a|^2=P_a(z)\) on \(\Omega\setminus S_a\)
for every algebraic \(a\). The function \(F\) is
nonconstant, since otherwise \(f\) would have finite image.

Here is the analytic implication in full. At a point
\(z\in\Omega\), write
\(w=F(z)\), \(A=F_z(z)\), \(B=F_{\bar z}(z)\), and
\(C=F_{z\bar z}(z)\). For all algebraic \(a\ne w\)
for which \(z\notin S_a\), harmonicity yields
\[
 \operatorname{Re}\left(\frac{C}{w-a}
                  -\frac{AB}{(w-a)^2}\right)=0.         \tag{B.MIX.18}
\]
There is at most one excluded algebraic centre: all
\(S_a\) consist of algebraic points and are disjoint.
Thus density extends (B.MIX.18) to every complex
\(a\ne w\). Substituting \(b=1/(w-a)\) gives
\(\operatorname{Re}(Cb-ABb^2)=0\) for every nonzero
\(b\in\mathbb C\). The two polynomial coefficients
must vanish, so \(C=0\) and \(AB=0\).
Consequently \(F\) is harmonic and conformal.
More explicitly, if \(F_z\) is nonzero anywhere,
\(F_{\bar z}\) vanishes on an open neighbourhood, and
real-analytic continuation makes it vanish on all of
\(\Omega\). Otherwise \(F_z\) vanishes identically.
Thus \(F\) is holomorphic or antiholomorphic on \(\Omega\).

For \(r\in S\), choose an algebraic centre
\(a\ne f(r)\). Its fibre omits \(r\), so \(P_a\)
is finite and positive near \(r\). Equation (B.MIX.16)
makes \(F\) bounded on a punctured neighbourhood.
The singularity is removable. Choose three noncollinear
such centres; their distance identities at \(r\),
including the original identities for \(f(r)\), show
that the extended value \(F(r)\) equals \(f(r)\).
This gives an entire holomorphic or antiholomorphic
function agreeing with \(f\) on every algebraic point.
For any fixed centre \(a\), (B.MIX.16) bounds its growth
at infinity by a finite power of \(|z|\). Cauchy's
estimates make it a polynomial or an antipolynomial:
\(F(z)=P(z)\) or \(F(z)=P(\bar z)\).
Interpolation at distinct rational integers, whose
images lie in \(E\), gives \(P\in E[T]\) in either case.

The nonconstant antipolynomial case is impossible.
Let \(c\) denote complex conjugation on \(\overline\Q\).
For \(g\in G_E\), equivariance would give
\(P(cg x)=P(gc x)\) for every \(x\). Thus the field
commutator \(h=cgcg^{-1}\) satisfies \(P(hz)=P(z)\)
for every \(z\). Apply this to \(z+n\), with
\(n\in\Z\). If \(d>0\) is the degree of \(P\),
comparison of the coefficients of \(n^{d-1}\) in
\(P(hz+n)=P(z+n)\) gives \(hz=z\).
So \(c\) centralizes \(G_E\), which is excluded by
the valuation and ramification argument in B.59.
Hence \(F=P\), proving the asserted actual point
function. Its polynomial is unique because the infinite
set of rational inputs already determines it.
\end{proof}
\paragraph{Check record.} Also checked in this writing, 2026-11-18: GPT-6 Astra (OpenAI, ChatGPT Pro mode), root; separate effort setting not exposed. Read complete proofs: ordinary pro-form consequence, solvable strict-local closure via exact tame/wild sequence, metabelian stable Kummer kernel and complementary-unit intersection, nilpotent Kummer theorem including p2/p3 exceptions, rational residues, affine similarity and polynomial point maps. B56 nontriviality is now proved directly by odd-prime affine groups and the p2 dihedral splitting field, without wild freeness. B57 reconstruction scope updated by B106–110.



% Pro-category completion, developed during the independent configuration audit.
% Provisional B.PRO identifiers; parent editor may renumber.
\subsection{Completion of ordinary pro-reconstruction in arbitrary transcendence degree}

We retain the ordinary category \(\mathsf T\), the model pro-objects,
the slice over \(\mathcal B\), and the limit-point definition of
pinnedness in B.43. In particular no compatible structural 2-cells
are assumed in the definition of a pro-morphism. The following
argument supplies the coherence that is used, and then reconstructs
at a fixed source field. This avoids trying to make the transition
maps pinned.

\begin{claim}{B.62}{Proved}{Compact coherence for ordinary pro-morphisms}
Let \(k\) have characteristic zero, let \(X,Y\) be finite-type
\(k\)-schemes, and let
\(\alpha:\mathcal X\to\mathcal Y\) be an ordinary pro-morphism
over \(\mathcal B\). There is a geometric morphism
\[
                 \alpha_k:X_{\mathrm{et}}\longrightarrow
                           Y_{\mathrm{et}}                 \tag{B.PRO.1}
\]
over \(BG_k\), compatible with the finite-level representatives of
\(\alpha\). Its underlying-space map is the map on inverse limits
defined by \(\alpha\), so it is pinned if \(\alpha\) is pinned.
Moreover finite constant base change, precomposition by a scheme
morphism, restriction to open subobjects, and slicing by a
quasi-compact separated étale object preserve the class of maps
obtained from ordinary pro-morphisms. For a slice defined at a
particular target level, the source model can be taken to be its
pullback by that level's chosen representative.
\end{claim}
\begin{proof}
Nilpotents do not affect the étale sites, so reduce the source.
At each target field \(F\), choose a representative
\(a_F:X_{E_F,\mathrm{et}}\to Y_{F,\mathrm{et}}\), where
\(E_F\supset F\) is finitely generated. Equality over the pro-base
permits one further source refinement making this a relative
morphism over \(BG_F\). Its structural 2-cell is unique:
\(G_F\) is slim by the opened Hoshi--Tsujimura Corollary~4.5,
and the image of each connected component's fundamental group in
\(G_F\) is open. For the latter assertion first map to \(G_{E_F}\);
the relative algebraic closure of \(F\) in the finitely generated
field \(E_F\) is finite. The same argument applies after every
finitely generated source refinement. Thus every absolute
comparison isomorphism at such a refinement is automatically
relative to \(BG_F\).

The opened \citep[Theorem~A.8 and Corollary~A.9]{CarlsonHaineWolf2024}
apply: the source is topologically noetherian and the target qcqs.
Every representative is coherent, and its inverse image takes
quasi-compact separated étale schemes to quasi-compact separated
étale schemes. Compose with \(X\to X_{E_F}\), obtaining functors
from the qcs étale sites of \(Y_F\) to that of \(X\).

Here is the finite-choice point that permits coherence. For two
qcs separated étale schemes \(U,V\) over a reduced noetherian
scheme \(X\) with finitely many irreducible components,
\(\operatorname{Isom}_X(U,V)\) is finite. An isomorphism is
determined on the finite fibres over the generic points: equality
there implies equality on \(U\), since \(U\) is reduced, those
fibres are dense, and \(V\to X\) is separated. The possible maps
of the finite generic étale algebras form a finite set.
Consequently the set of natural isomorphisms between two of the
preceding inverse-image functors is a closed subset of a product
of finite sets. It is compact, and the relative condition over
\(BG_F\) is closed as well.

For \(F\subset H\), ordinary pro-compatibility supplies an
isomorphism between the functor at \(F\) and the restriction of
the functor at \(H\), after a single finitely generated source
refinement and hence after base change to \(k\). The compact
space of relative comparison isomorphisms is therefore nonempty.
In each such compact space use the \emph{closure of the subset
of isomorphisms which descend after a finitely generated source
refinement}. This refinement is important: an arbitrary isomorphism
over \(k\) need not descend as an isomorphism of the entire
inverse-image functors. Membership in this closure says that on
every finite list of étale objects its components agree with those
of one whole comparison isomorphism descending at a finite stage.
Take the product of these closed spaces over all pairs of fields, and
impose naturality, identity, and the triple cocycle equations.
These are closed conditions. Every finite collection of them is
satisfiable: choose one common upper target field, compare its
representative with all the finitely many representatives after
one common finitely generated source refinement, and obtain all
comparisons by passing through this one representative. They
then satisfy every required finite cocycle equation. Compactness
gives a simultaneous coherent choice. This proves existence of
coherence; it has not been inserted into the hypothesis.

Every qcs étale object, every morphism between such objects, and
every finite diagram over \(Y\) descends to a finitely generated
field, by finite presentation. Finite covering families descend
too; the exact site-continuity statement is the opened
\citep[Lemma~59.51.2, Tag~0EYL, together with
Lemma~32.10.1, Tag~01ZM]{StacksProject}. We use the canonical
associative base-change identifications, or equivalently first
strictify this fixed diagram of sites. The coherently identified
inverse-image functors consequently
define a left exact covering-preserving functor on this site of
\(Y\), hence (B.PRO.1). The identical construction on finite étale
constant objects gives its structural isomorphism over \(BG_k\).
The comparison on open subobjects proves the assertion on spaces.

For clarity, slicing entails no additional coherence assumption.
Suppose \(V_F\to Y_F\) is qcs separated étale and put
\(U_{E_F}=a_F^*V_F\). For any higher target level, the preceding
comparisons identify the pullback of \(V\) with \(U\) over \(X\).
This is an isomorphism between finitely presented étale schemes,
so it descends after one finitely generated source refinement.
The resulting slice map represents a map of the model systems
\(\mathcal U\to\mathcal V\). For ordinary pro-compatibility,
approximate a selected comparison on the finite list consisting
of \(V\) and any specified arrows between slice objects by a
\emph{whole} comparison isomorphism descending at a finite stage.
This is exactly the closure property imposed above. Its component
on \(V\) agrees with the selected identification, so it induces
the required comparison of the entire slice geometric morphisms
after one refinement. Thus descending an individual isomorphism
of étale objects is not being substituted for descending a
comparison of functors. The same construction
over a finite étale arithmetic base gives finite constant base
change; cofinal fields have the form \(F'L_0\) for a fixed finite
field extension \(L_0/F_0\) defining that base change. Open
restriction and scheme precomposition are immediate. If the
original map is pinned, so are the slice maps: an étale
finite-type scheme has only closed points above a closed point.
\end{proof}
\paragraph{Check record.} Also checked in this writing, 2026-11-18: GPT-6 Astra (OpenAI, ChatGPT Pro mode), root; separate effort setting not exposed. Read complete ordinary pro-category proofs, including compact choices of comparisons descending after one refinement, slice coherence, simultaneous whole-Kummer comparison, completion saturation, fixed-field reconstruction, unique functor isomorphisms and h-descent. Opened CHW Proposition5.22, Proposition6.22 and Claim6.23, TheoremA.8, CorollaryA.9, Proposition4.8 and Stacks Tags0EYL, 01ZM, 0CUF with exact hypotheses.


\begin{claim}{B.63}{Proved}{Pro-points and a single-stage Kummer comparison}
For every characteristic-zero \(k\), pinned ordinary pro-morphisms
\(\mathcal B\to\mathcal Y\) over \(\mathcal B\) are precisely
the \(k\)-rational points of \(Y\). This statement is compatible
with finite extensions of \(k\).

In particular, suppose a fixed representative of a pinned
pro-morphism pulls the universal Kummer class back to
\(\xi\in H^1(X_E,\widehat\Z(1))\). If \(x\in X_E(L)\), with
\(L/E\) finite, and \(b\) is the value of its induced pro-point
in \(\mathbb G_m\), then there is \emph{one} finitely generated
field \(M\) containing \(L\) and \(b\) for which
\[
        \operatorname{res}_{M/L}(x^*\xi)=\kappa_M(b)
                \quad\hbox{in }\widehat{M^\times}.       \tag{B.PRO.2}
\]
The field \(M\) works simultaneously for all Kummer levels.
\end{claim}
\begin{proof}
Let the limit point of a pinned pro-point be \(y\in Y\).
It is closed, and a finite closed subscheme containing it descends
to some \(Y_{F_0}\). At every sufficiently large target field
\(F\), its image \(y_F\) is therefore closed, with
\(\kappa(y_F)/F\) finite. A representative
\(BG_E\to Y_{F,\mathrm{et}}\) has underlying image \(y_F\).
It factors through the closed subtopos
\(B G_{\kappa(y_F)}\): the inverse image of the complementary
open subobject is empty. After one source refinement it is
relative over \(BG_F\), with the unique structural 2-cell
established in B.62. Since \(\kappa(y_F)/F\) is finite
separable, the finite étale slice description says that this
relative morphism is exactly an \(F\)-embedding
\(\kappa(y_F)\hookrightarrow E\). Thus it is an actual
\(E\)-valued point of \(Y_F\). Compatibility at a common
source refinement makes these points the same \(k\)-point;
the finite residue extension of the limit point consequently
has degree one. This also proves uniqueness of the pro-morphism.

For the second assertion base change the pro-morphism to the
finite extension \(kL/k\), using B.62, and compose with the
point \(x\). The first assertion gives an actual value
\(b\in(kL)^\times\). Choose a finitely generated target field
\(H\subset kL\) containing \(b\) and all fields needed for the
displayed point. At target \(\mathbb G_{m,H}\), the pro-point's
representative has the closed underlying image \(T=b\).
After one refinement making it relative to \(BG_H\), it is
the constant-point morphism, as just proved. Its comparison
with the original representative at the original target is an
isomorphism of \emph{whole geometric morphisms} after one common
finitely generated source refinement \(M\). Pulling back all
the compatible universal Kummer torsors gives (B.PRO.2).
This proves simultaneous equality in the completed multiplicative
group, rather than a sequence of equalities over fields depending
on \(n\), or an equality only after restriction to \(G_{kL}\).
\end{proof}
\paragraph{Check record.} Also checked in this writing, 2026-11-18: GPT-6 Astra (OpenAI, ChatGPT Pro mode), root; separate effort setting not exposed. Read complete ordinary pro-category proofs, including compact choices of comparisons descending after one refinement, slice coherence, simultaneous whole-Kummer comparison, completion saturation, fixed-field reconstruction, unique functor isomorphisms and h-descent. Opened CHW Proposition5.22, Proposition6.22 and Claim6.23, TheoremA.8, CorollaryA.9, Proposition4.8 and Stacks Tags0EYL, 01ZM, 0CUF with exact hypotheses.


\begin{claim}{B.64}{Proved}{Intersection with a finitely generated field's Kummer completion}
Let \(A\subset B\) be finitely generated fields of characteristic
zero. If \(\xi\in\widehat{A^\times}\), \(b\in B^\times\), and
\[
                    \operatorname{res}_{B/A}\xi=\kappa_B(b),
                                                               \tag{B.PRO.3}
\]
then \(b\in A^\times\) and \(\xi=\kappa_A(b)\).
\end{claim}
\begin{proof}
Let \(A'\) be the relative algebraic closure of \(A\) in
\(B\); it is finite over \(A\). Take a normal projective
\(A'\)-model of the finitely generated regular extension
\(B/A'\). Every codimension-one valuation of this model is
trivial on \(A\). Applying its completed valuation to
(B.PRO.3) shows that the ordinary integer valuation of \(b\)
is zero. Normality then makes \(b\) a global invertible function
on this proper model. Such functions belong to its constant
field \(A'\), so \(b\in A'^\times\).

The map \(\widehat{A'^\times}\to\widehat{B^\times}\) is
injective: for every \(n\), an \(n\)-th root in \(B\) of an
element of \(A'\) is algebraic over \(A'\), hence already in
\(A'\). Thus (B.PRO.3) already holds over \(A'\). Taking the
norm to \(A\), with \(d=[A':A]\), gives
\[
                       d\xi=\kappa_A(N_{A'/A}b).          \tag{B.PRO.4}
\]

We record the elementary saturation argument, so no unsupported
finite-generation claim about \(A^\times\) is needed. In the
completion \(\widehat{A^\times}=\varprojlim_n
A^\times/(A^\times)^n\), if \(d\eta=0\), then its component
modulo \(n\) is represented by a \(d\)-th root of unity.
Indeed lift to level \(nd\): if \(a^d=c^{nd}\), then
\(a=c^n\zeta\) with \(\zeta^d=1\). The finite set
\(\mu_d(A)\) and compatibility over common multiples imply
\(\eta=\kappa_A(\zeta)\) for one \(\zeta\in\mu_d(A)\).
If \(d\xi=\kappa_A(a)\), reduction modulo \(d\) first gives
\(a=c^d\) with \(c\in A^\times\). Apply the preceding
torsion assertion to \(\xi-\kappa_A(c)\); hence \(\xi\)
is the Kummer class of an actual element \(c'\in A^\times\).
Apply this to (B.PRO.4). Finally Kummer maps of finitely generated
characteristic-zero fields are injective, by
\citep[Proposition~5.24 and Observation~6.9]{CarlsonHaineWolf2024}.
Equation (B.PRO.3) therefore makes \(c'=b\) in \(B\), proving
both assertions.
\end{proof}
\paragraph{Check record.} Also checked in this writing, 2026-11-18: GPT-6 Astra (OpenAI, ChatGPT Pro mode), root; separate effort setting not exposed. Read complete ordinary pro-category proofs, including compact choices of comparisons descending after one refinement, slice coherence, simultaneous whole-Kummer comparison, completion saturation, fixed-field reconstruction, unique functor isomorphisms and h-descent. Opened CHW Proposition5.22, Proposition6.22 and Claim6.23, TheoremA.8, CorollaryA.9, Proposition4.8 and Stacks Tags0EYL, 01ZM, 0CUF with exact hypotheses.


\begin{claim}{B.65}{Proved}{Reconstruction at one fixed source field}
Let \(\alpha\) be a pinned ordinary pro-morphism and choose a
relative representative
\[
                   a:X_{E,\mathrm{et}}\longrightarrow
                     Y_{F,\mathrm{et}}                 \tag{B.PRO.5}
\]
over \(BG_F\), with \(F\subset E\subset k\) finitely
generated. If \(X_E\) is normal, then (B.PRO.5) is induced by a
unique \(F\)-scheme morphism \(f_E:X_E\to Y_F\).
Its lift to \(Y_E\), after extension to \(k\), realizes the
point function of \(\alpha_k\). The same conclusion holds with
the absolutely weak normalization of \(X_E\) as source.
\end{claim}
\begin{proof}
\emph{1. Units over the fixed field.}
First take \(Y_F=\mathbb G_{m,F}\), and suppose \(X_E\) is
affine, geometrically integral, normal, and has an \(E\)-point.
The pullback \(\xi\) of the compatible Kummer classes lies in
\(\widehat{\mathcal O(X_E)^\times}\). For normal schemes this
follows from the Kummer sequence and finite generation of
\(\operatorname{Pic}(X_E)\), as stated in the opened
\citep[Proposition~5.22]{CarlsonHaineWolf2024}.
For every finite \(L/E\) and \(x\in X_E(L)\), B.63 gives
one finitely generated \(M\supset L\) and a point value \(b\)
satisfying (B.PRO.2). B.64 gives
\[
                         b\in L^\times,\qquad
                       x^*\xi=\kappa_L(b).              \tag{B.PRO.6}
\]
In particular the value at the chosen \(E\)-point is an actual
element of \(E^\times\); divide it out.

The normalized unit group is finitely generated free, by
\citep[Corollary~6.18]{CarlsonHaineWolf2024}. The evaluation of
the normalized completed unit at every algebraic point is an
actual algebraic unit by (B.PRO.6). The arithmetic argument of
Proposition~6.22 of that source now makes the completed unit
an actual unit of \(X_E\). Its hypotheses are precisely a
finitely generated characteristic-zero field (hence amenable
and admitting a nontrivial discrete valuation), the stated
finitely generated free unit group, and these actual point
evaluations. The proof works for normal sources as well:
Picard finite generation is Proposition~5.22, the unit assertion
is Corollary~6.18, and the nonconstant unit maps used in Claim~6.23
have image containing a nonempty open of \(\mathbb G_m\).
The source explicitly records this normal-source extension in
Remark~4.20. Equivalently, its full-rank integer valuation matrix
makes an integer multiple of each profinite exponent integral;
\(\widehat\Z\cap\Q=\Z\) then makes each exponent integral.
Thus \(\xi=\kappa(f)\) for \(f\in\mathcal O(X_E)^\times\).

This \(f\) realizes the value at \emph{every} point of
\(X(\overline k)\), not only the points algebraic over \(E\).
For such a point choose a finitely generated field of definition
\(L\supset E\) inside \(\overline k\). This extension need
not be finite. First base change the pro-system to the finite
extension \(kL/k\), then refine its source model from \(E\)
to \(L\). The point is now \(L\)-rational, so B.63 applies
with the identity finite extension of the new source field
\(L\); B.64 applies to the resulting finitely generated
comparison field. Its value \(b\) satisfies
\(\kappa_L(b)=\kappa_L(f(x))\); injectivity gives
\(b=f(x)\). The finite extension of \(k\) needed to contain
the point is permitted by B.62.

Remove the rational-point and geometric-connectedness assumptions
by a finite Galois extension of \(E\) splitting the finitely
many geometric components and supplying a point on each. The
preceding proof gives units on these components. Their Kummer
classes are the pullback of the one fixed \(\xi\), so Galois
conjugates agree by Kummer injectivity. They descend to a unit
over \(E\). Passing to this finite extension is justified by
finite constant base change to its compositum with \(k\).
Finally a finite affine open cover handles a nonaffine normal
source; equality of the point functions glues the units uniquely.

\emph{2. Coordinates and the point function.}
For general \(Y_F\), inverse images of Zariski opens under
\(a\) are Zariski opens of \(X_E\), and the corresponding pro
restrictions are pinned. For a regular function \(h\) on an
affine target open, cover that open by \(D(h)\) and \(D(h-1)\).
Step~1 reconstructs the pulled-back units \(h\) and \(h-1\)
on their respective source opens over the same field \(E\).
On the overlap their difference is one, since this holds on
all \(\overline k\)-points and the source is reduced. They
therefore give a regular function on the entire inverse-image
open. The ring operations and all defining relations hold on
\(\overline k\)-points, hence as functions. The resulting maps
on affine target opens glue uniquely to
\(f_E:X_E\to Y_F\). Their point function is that of
\(\alpha_k\). The open inverse-image conditions hold as well,
because finite-type schemes over a field are Jacobson and
membership in an open is determined by the closed points of
its complement. This argument does not require the target to
be separated.

\emph{3. The entire functor at the same source level.}
For each qcs separated étale \(V_F\to Y_F\), the sheaf
\(a^*V_F\) is represented by a qcs separated étale scheme
\(U_E\to X_E\), by CHW A.8--A.9. It is normal. Slice the
pro-morphism, as in B.62, using the fixed representative
\(a_V:U_{E,\mathrm{et}}\to V_{F,\mathrm{et}}\).
Step~2 gives a scheme map \(g_V:U_E\to V_F\) over this
same field \(E\), with the slice point function. This step
does not require the whole topos map of the slice to have been
reconstructed already. Composing to \(Y_F\) and comparing
point functions gives
\[
                 g_V\text{ over }Y_F
                    = f_E\circ(U_E\to X_E).
\]
It induces \(h_V:U_E\to X_E\times_{Y_F}V_F\).

The coherent realization (B.PRO.1) supplies the usual étale
slice fibre bijection: above every \(x\in X(\overline k)\),
the points of \(U\) are precisely the points of \(V\) above
\(\alpha_k(x)\). The map \(h_V\) realizes that bijection.
An étale morphism between étale schemes over \(X\) which is
bijective on every geometric closed fibre is an isomorphism:
the failure loci are constructible, so any nonempty such locus
over a field would contain a closed point. Thus \((h_V)_k\)
is an isomorphism, and faithful flatness makes \(h_V\) an
isomorphism over \(E\). This argument requires no additional
extension depending on \(V\).

For a morphism \(V_F\to W_F\), the two maps between the
reconstructed pullbacks agree on \(\overline k\)-points by
the natural slice functor of \(\alpha_k\); reducedness makes
them equal as scheme maps. The \(h_V\) consequently form a
natural isomorphism \(a^*\simeq f_E^*\) on the entire qcs
étale site. That site is a basis, so (B.PRO.5) is the induced
geometric morphism. This is equality at the fixed source field,
not just an equality of limit point functions or an objectwise
collection of equalities after unrelated extensions.

\emph{4. Faithfulness, automorphisms, and a seminormal source.}
We spell out the descent needed for the last assertion. For two
\(F\)-scheme maps from a reduced finite-type \(E\)-scheme to
\(Y_F\), an isomorphism of their relative geometric morphisms
implies equality of their pullbacks of all open subobjects and
all Kummer classes of units on target opens. Unit Kummer maps
on the source opens are injective: evaluate at closed points,
whose residue fields are finitely generated of characteristic
zero, and use their injectivity, followed by reducedness.
The coordinate argument of Step~2 therefore makes the two
scheme maps equal. This proves the required faithfulness even
when \(E/F\) is transcendental.

An actual such scheme map has no automorphisms as a geometric
morphism. At any source point \(x\), its residue field and
that of its image are finitely generated over \(\Q\); the
image of the source residue Galois group in the target residue
Galois group is open. The latter group is slim. Thus the
centralizer describing an automorphism on that stalk is trivial.
Geometric stalks are conservative, proving the assertion for
the whole inverse-image functor.

Replace \(X_E\) by its absolute weak normalization; in
characteristic zero this is finite-type seminormal. Its finite
normalization \(Z_E\to X_E\) is an h-cover. Apply the normal
case to \(Z_E\) and to a normalization of the reduced fibre
product \(Z_E\times_{X_E}Z_E\). The preceding faithfulness
makes the two maps agree after this normalization. Since that
normalization is an h-cover, h-separatedness in CHW Theorem~4.1
makes them agree on \((Z_E\times_{X_E}Z_E)^{\mathrm{awn}}\)
before applying the h-sheaf equalizer. The opened
\citep[Theorem~4.1]{CarlsonHaineWolf2024}, which states that
\(T\mapsto\operatorname{Hom}(T^{\mathrm{awn}},Y_E)\)
is an h-sheaf, descends them to a scheme map on \(X_E\).
Their topos isomorphisms are unique by the just-proved absence
of automorphisms and hence satisfy the cocycle condition.
Étale-topos h-descent, the opened Proposition~4.8 of the same
source, descends the isomorphism to (B.PRO.5) itself. This proves
the last assertion and uniqueness.
\end{proof}
\paragraph{Check record.} Also checked in this writing, 2026-11-18: GPT-6 Astra (OpenAI, ChatGPT Pro mode), root; separate effort setting not exposed. Read complete ordinary pro-category proofs, including compact choices of comparisons descending after one refinement, slice coherence, simultaneous whole-Kummer comparison, completion saturation, fixed-field reconstruction, unique functor isomorphisms and h-descent. Opened CHW Proposition5.22, Proposition6.22 and Claim6.23, TheoremA.8, CorollaryA.9, Proposition4.8 and Stacks Tags0EYL, 01ZM, 0CUF with exact hypotheses.


\begin{claim}{B.66}{Proved}{The ordinary pro-form for every characteristic-zero field}
For every field \(k\) of characteristic zero, the functor from
finite-type \(k\)-schemes localized at universal homeomorphisms
to their model pro-systems is fully faithful, with pinned
morphisms in the ordinary slice
\(\operatorname{Pro}(\mathsf T)_{/\mathcal B}\) as defined
above. In particular B9 is proved in arbitrary transcendence
degree; no coherent-relative morphism condition has been added
to its input.
\end{claim}
\begin{proof}
For a pinned ordinary pro-morphism \(\alpha\), choose a target
field \(F_0\) and a relative representative at a finitely
generated source field \(E\). B.65 reconstructs it as a
localized scheme map \(X_E\to Y_{F_0}\). Pairing with the
structural map to \(E\) and extending to \(k\) gives
\(f:X^{\mathrm{awn}}\to Y\). At every other target field
\(F\), B.65 reconstructs that level as well. Pro-compatibility
compares the two maps after one common finitely generated
source extension. The faithfulness proved in Step~4 of B.65
makes their projections to \(Y_{F_0}\) equal; their structural
maps agree, so their lifts to \(Y_F\) and to \(Y\) agree.
At each target field the ordinary pro-class is therefore that
of \(f\). This is surjectivity on the Hom set in the ordinary
pro-category, rather than only on the corresponding limit
point functions.

Conversely, every localized finite-type \(k\)-scheme morphism
and its finite diagram of equations descend to a finitely
generated subfield, and hence define the model pro-morphism.
Its limit map preserves closed points. If two scheme morphisms
give the same pro-morphism, compare at a common target field,
refine once, and use the same fixed-field faithfulness; their
base changes to \(k\) are equal. Absolute weak normalizations
identify the localization, as in CHW Theorem~1.13 and
Proposition~2.28. This proves full faithfulness.
\end{proof}
\paragraph{Check record.} Also checked in this writing, 2026-11-18: GPT-6 Astra (OpenAI, ChatGPT Pro mode), root; separate effort setting not exposed. Read complete ordinary pro-category proofs, including compact choices of comparisons descending after one refinement, slice coherence, simultaneous whole-Kummer comparison, completion saturation, fixed-field reconstruction, unique functor isomorphisms and h-descent. Opened CHW Proposition5.22, Proposition6.22 and Claim6.23, TheoremA.8, CorollaryA.9, Proposition4.8 and Stacks Tags0EYL, 01ZM, 0CUF with exact hypotheses.



\subsection{The remaining universal supersingular-target assertion}


\noindent\emph{Statement B.67 and its produced argument are retained in Part II.}
\subsection{Point sections and generalized-Jacobian Kummer classes}

We use left Galois cocycles: $z(\sigma\tau)=z(\sigma)+\sigma z(\tau)$.
For a geometric path $\gamma$ from a rational base point $b$ to a
rational point $x$, transport back to $b$ defines the section $s_x^\gamma$.
Our section-difference convention is $s_x^\gamma(\sigma)s_b(\sigma)^{-1}$.
The geometric abelian quotient is identified with torsion of the Jacobian
by the monodromy of multiplication covers. These choices fix the sign.

\begin{claim}{B.68}{Proved}{The point-to-section Kummer comparison, including a modulus}
Let $K$ be any field, let $X/K$ be a smooth, proper, geometrically
connected curve, and let $\mathfrak m$ be an effective divisor on $X$.
Put $U=X-\operatorname{supp}(\mathfrak m)$ and assume $b\in U(K)$.
Let $J_{\mathfrak m}$ be the generalized Jacobian, with $J_0=J_X$, and
let $j_b:U\to J_{\mathfrak m}$ send $x$ to the modulus class $[x-b]$.
For every positive integer $n$ invertible in $K$, the sequence
\[
 0\longrightarrow J_{\mathfrak m}[n]\longrightarrow
 J_{\mathfrak m}\xrightarrow{[n]}J_{\mathfrak m}\longrightarrow0
\]
is exact as a sequence of fppf sheaves, and $[n]$ is finite \emph{\'etale}.
The Abel map induces a $G_K$-equivariant isomorphism of finite
abelian groups
\[
 \pi_1(U_{\overline K},\bar b)^{\mathrm{ab}}/n
 \ \xrightarrow{\sim}\ J_{\mathfrak m}[n](\overline K).
\]
For every $x\in U(K)$ the abelianized section-difference class is
\[
 [s_x-s_b]_n=\delta_n(j_b(x))\quad\hbox{in}\quad
 H^1(G_K,J_{\mathfrak m}[n](\overline K)).                 \tag{N1.1}
\]
Thus the sign is positive for $x-b$ and the stated difference convention.
Changing the geometric path changes the cocycle by a coboundary;
conjugacy throughout is conjugacy by the geometric kernel. Reversing
the section difference gives $-\delta_n(j_b(x))$.
\end{claim}
\begin{proof}
We give the geometric comparison and the arithmetic sign separately.
Over an algebraically closed field a generalized Jacobian is an
extension of the ordinary Jacobian by the connected affine group
of units of the modulus, modulo constant units. For reduced modulus
with $d>0$ geometric points its affine part is $\mathbf G_m^{d-1}$.
Increasing the multiplicities adds a connected unipotent kernel.
Multiplication by $n$ is an automorphism on that unipotent kernel:
in characteristic zero one uses a filtration by additive groups;
in characteristic $p$ one uses the finite $p$-power exponent
filtration of the units $1+\mathfrak a$ of a finite local algebra.
For the torus and the abelian quotient multiplication by $n$ is
finite, surjective and \emph{\'etale}. The same follows for
$J_{\mathfrak m}$ by descent and the extension sequence. In
particular all its geometric $n$-division points are separable.

Here is a direct verification that the indicated abelian quotient
is the one represented by these covers. For a nonempty reduced
geometric boundary $D$ the localization sequence gives
\[
 0\to H^1(X_{\overline K},\mathbf Z/n)\to
 H^1(U_{\overline K},\mathbf Z/n)\to
 (\mathbf Z/n)(-1)^D\xrightarrow{\sum}
 (\mathbf Z/n)(-1)\to0.                                \tag{N1.2}
\]
The residue map is the tame-inertia character at each boundary point;
the last map is degree. Its exactness follows from the Kummer
localization sequence, purity for a smooth curve, and the degree
isomorphism in degree two. The proper Abel map identifies the first
term with the character group of $J_X[n]$: this is
\citet[Proposition 9.1]{MilneJac}, stated there for a complete smooth
curve over a separably closed field, with all connected abelian
\emph{\'etale} covers included. On the generalized-Jacobian side the
exact sequence of $n$-torsion is
\[
 0\to (\boldsymbol\mu_n^D/\boldsymbol\mu_n)
 \to J_D[n]\to J_X[n]\to0.
\]
Its character sequence is (N1.2). To check the comparison at the
boundary, trivialize a modulus line bundle near the boundary and
move its zero once around a local parameter. The trivialization
changes in precisely that boundary factor of $\boldsymbol\mu_n^D$;
simultaneous change of all trivializations is a constant unit and
is quotiented out. Equivalently the local multiplication cover is
$w^n=z$, with residue $+1$ in the corresponding factor. Thus the
residue maps agree, and the two end isomorphisms give the middle
isomorphism. If $D$ is empty this is just the proper assertion.
The unipotent kernel contributes no $n$-torsion and no finite
abelian cover of exponent dividing $n$, so the answer is unchanged
by multiplicities. All constructions are canonical under field
automorphisms and hence $G_K$-equivariant. For an imperfect ground
field this verification can be made after its perfect closure;
purely inseparable extension induces equivalences on finite
\emph{\'etale} schemes and identifies absolute Galois groups, so
the same assertion descends to $K$.

Put $a=j_b(x)$ and choose $t\in J_{\mathfrak m}(K^{\mathrm{sep}})$
with $nt=a$. By definition of the connecting map,
\[
               \delta_n(a)(\sigma)=\sigma(t)-t.           \tag{N1.3}
\]
Pull back $[n]$ by $j_b$ to a finite \emph{\'etale}
$J_{\mathfrak m}[n]$-torsor on $U$. Its fibre at $b$ is pointed
by $0$ and its fibre at $x$ contains $t$. Choose a geometric path
$\gamma$ carrying $0$ to $t$. Such a path exists because a
geometric path identifies the fibre functors, and its possible
values on this torsor are all its geometric fibre points by the
just-proved geometric quotient isomorphism. The conjugate path
$\sigma\gamma$ carries $0$ to $\sigma(t)$. Consequently
$\gamma^{-1}(\sigma\gamma)$, acting on the fibre at $b$, has
translation coordinate $\sigma(t)-t$. With transport of sections
back to $b$, this path difference is exactly the image of
$s_x^\gamma(\sigma)s_b(\sigma)^{-1}$. This proves (N1.1).
Replacing $t$ by $t+u$ adds $\sigma(u)-u$, which also verifies
explicitly the ambiguity under a different path. Restriction to
any field extension, maps of pointed curves, and transition maps
$n'\to n$ commute with this calculation. Hence it proves the
compatible Tate-module statement as well as each finite level.
\end{proof}
\paragraph{Check record.} Checked only in this writing, GPT-6 Astra Pro, 18 November 2026: complete proof written, root targeted reading and configuration-reader second reading; separate effort setting not exposed.


\paragraph{Check record.}
Checked only in this writing: GPT-6 Astra (OpenAI, ChatGPT Pro mode),
18 November 2026, independent cocycle derivation and direct checking
of the generalized-Jacobian comparison. The source's proper-curve
Proposition 9.1 and its generalized-Jacobian discussion, including
Proposition 9.7, were opened. A separate reasoning-effort setting
was not exposed. No checking by another model family is claimed.

\paragraph{Manuscript scope and the Kummer kernels.}
Cote 43, scan 113 (manuscript page 112), contains the local
decomposition and inertia homomorphisms in a deleted passage; it
does not specify the path or cocycle sign in (N1.1). Scans 132--133
(pages 131--132) contain the Kummer sequence and the $T(J)$ bridge.
These are compatible with (N1.1), but we do not attribute our
modern sign convention to an undeciphered symbol. Cote 149,
scan 33 (manuscript page 16), explicitly identifies the
point homomorphism with the inverse-limit Kummer homomorphism.
Scans 34--36 then ask what remains when the characteristic prime
is omitted. Our proposition supplies exactly that comparison.

For a finite field $k$ the group $J_X(k)$ is finite. The kernel
of its map to $\prod_{\ell\ne p}H^1(G_k,T_\ell J_X)$ is therefore
its $p$-primary subgroup. Equality of proper point sections
implies $[P-Q]\in J_X(k)[p^\infty]$. This is the principal-divisor
step used in A.19, with $P-Q$ obtained by subtracting the two
classes relative to the same base point; replacing it by $Q-P$
changes the sign but not its order.

More generally the proof of A.30 applies to \emph{semiabelian}
groups over fields finitely generated over their prime field.
We have rechecked it: after splitting the torus over a finite
separable extension, the divisor sequence makes $K^\times$
an extension of its constant units by a free abelian group;
the constants have finite torsion and a free abelian quotient.
The abelian-variety rational points are finitely generated by
Lang--N\'eron (\citealt[Theorem 2.1]{Conrad2006}) and the usual
Mordell--Weil theorem for number-field constants. Lifts of finitely
many generators and relations use only finitely many basis
elements of the free torus subgroup. Thus the semiabelian
rational-point group is a direct sum of a free abelian group
and a finitely generated group. Its intersection of all
prime-to-$p$ multiples is precisely its $p$-primary torsion,
and in characteristic zero the intersection of all multiples
is zero. Descent of the resulting torsion assertion through
the finite splitting extension is valid because rational-point
inclusion is injective. This proves the kernels asserted in A.30.

The claim on cote 149, scan 34, that all such $G(K)$ are finitely
generated is false for a torus: $\mathbf G_m(\mathbf Q)=\mathbf Q^\times$
has the independent valuations at all rational primes. The
divisibility conclusion above repairs the argument without that
claim. For a nonreduced modulus $J_{\mathfrak m}$ need not be
semiabelian; its unipotent part is invisible to prime-to-characteristic
Kummer theory. We therefore do not extend A.30 to that larger
class of groups. For example $J_{2\infty}$ of $\mathbf P^1$
is $\mathbf G_a$, on which every invertible $n$ acts bijectively.

\subsection{The complete cubic test for one pointed quartic in characteristic five}

\begin{claim}{B.80}{Proved}{All cubic covers of the pointed quartic $W^4=X^3Z-XZ^3$}
Let
\[
 C/\mathbf F _5:\quad W^4=X^3Z-XZ^3,
 \qquad P=(1:0:0),\qquad K=\mathbf F _{5^8}.
\]
Every nonzero geometric order-three character first becomes constant
over $K$. There are $364$ connected constant cyclic degree-three
$P$-split covers of $C_K$, up to inversion of the deck character.
Exactly $96$ are class-number witnesses; the other $268$ fail the
test over $K$ and every further constant extension. The successful
ones remain witnesses over $\mathbf F _{5^{8n}}$ exactly when
$3\nmid n$. Consequently the least admissible odd prime giving a
witness for this pointed quartic over some extension is three, and
the least field for an order-three witness is $\mathbf F _{5^8}$.

More explicitly, the three elliptic quotients identify
$J_C[3]$ over $\overline{\mathbf F}_5$ with
$E_0[3]\oplus E_1[3]\oplus E_2[3]$. A cover represented by
$(T_0,T_1,T_2)$ is successful precisely when exactly two entries
are nonzero. On its new coherent cohomology, the characteristic
polynomial of the $K$-linear Cartier iterate is
\[
 \begin{cases}
 (T^2+T+1)^2,&\text{exactly two nonzero entries},\\
 (T-1)^4,&\text{one or three nonzero entries}.
 \end{cases}                                                   \tag{B.80.1}
\]
The different optimization in B.34, minimizing the extension degree
first among all primes, has degree three and least prime thirteen
there. It is not the prime-first minimum asserted here.
\end{claim}

\begin{proof}
We first justify that the calculation covers every geometric character.
Choose $\eta$ with $\eta^4=-1$, so $\eta\in\mathbf F _{25}$.
The change of variables
\[
 U=X-2Z,\qquad V=W,\qquad T=\eta(X+2Z)
\]
takes $C$ to the Fermat quartic $U^4+V^4+T^4=0$.
Its three sign involutions generate a Klein four group, whose common
quotient is the genus-zero conic in $U^2,V^2,T^2$.
The three individual quotients are elliptic curves. Their quotient
maps $q_i:C\to E_i$ give homomorphisms
\[
 \phi:\prod_iE_i\to J_C,\quad (a_i)\mapsto\sum_iq_i^*a_i,
 \qquad
 \psi:J_C\to\prod_iE_i,\quad a\mapsto(q_{i*}a)_i.
\]
On $J_C$, $q_i^*q_{i*}=1+\iota_i$. The norm of the Klein four
group factors through the genus-zero quotient, and is zero. Hence
$\phi\psi=\sum_i(1+\iota_i)=2$. The mixed maps
$q_{i*}q_j^*$ vanish for $i\ne j$: the involution $\iota_i$ acts
by $-1$ on the elliptic quotient attached to $\iota_j$, since its
quotient there is the genus-zero common quotient. The mixed map is
therefore killed by two and must vanish as a homomorphism of abelian
varieties. Thus $\psi\phi=2$ as well. In particular the maps induce
inverse isomorphisms, up to multiplication by two, on three-torsion.

All three elliptic quotients become isomorphic over $\mathbf F _{25}$
to $E:y^2=x^3-x$. This follows explicitly by permuting the Fermat
coordinates and reversing the displayed linear change. The three
matrices implementing it are retained, and substitution checks each
quartic equation identically. The curve $E$ has eight points over
$\mathbf F _5$, so its Frobenius satisfies
$F_E^2+2F_E+5=0$. Modulo three this gives $F_E^4=-1$ and
$F_E^8=1$. Since the isogeny above is defined over $\mathbf F _{25}$,
the same identities hold on $J_C[3]$. No nonzero vector can have
period dividing four. Every period divides eight, so every nonzero
vector has period exactly eight. On the character space the dual
contragredient action has the same property. Thus degree eight is
the first field on which any nonzero order-three character is constant.

We describe all of the covers and their pointings. For $T_i=(a_i,b_i)$
in the nonzero three-torsion of $E_i$, let
\[
 m_i=\frac{3a_i^2-1}{2b_i},\qquad c_i=b_i-m_i a_i,
 \qquad g_i=y_i-m_ix_i-c_i.
\]
The tangent is a flex, so
$\operatorname{div}(g_i)=3T_i-3O_i$. Directly,
$(m_i(x-a_i)+b_i)^2-(x^3-x)=-(x-a_i)^3$.
For an active set of entries put $G=\prod_iq_i^*g_i$ and divide
it by its leading unit at $P$. The normalization is explicit in the
certificate. If the identity quotient is active, it contributes a
pole of order six at $P$, with leading coefficient one; this pole
is a cube and gives no ramification. All other active factors are
units at $P$. The normalization of
\[
                             z^3=G
\]
is therefore an \emph{\'etale}, constant cyclic, $P$-split cover.
It is connected because the corresponding nonzero three-torsion
line bundle is not trivial. Replacing every $T_i$ by $-T_i$ reverses
the deck generator and gives the same cover. These are the only
identifications. Consequently the $9^3-1$ nonzero triples give
$(9^3-1)/2=364$ covers, with respectively
\[
               12,\quad 96,\quad 256
\]
having one, two, or three nonzero entries. The equality with all
constant characters uses $\mu_3\subset K$ and the Kummer
identification with $J_C[3]$; the chosen unit normalization removes
the constant twist ambiguity.

We now justify the differential calculation. In the affine coordinate
ring $K[x,w]/(w^4-x^3+x)$, write
$q_i=(X_i/Z_i,W_i^2/Z_i^2)$, with the displayed linear forms coming
from the three coordinate matrices. Set
$Q_i=W_i^2-m_iX_iZ_i-c_iZ_i^2$. If $s$ entries are active, take
homogeneous forms $H$ of degree $s+1$. In the two nontrivial deck
characters, regular differentials have the form
\[
                \frac{H\,dx}{w^3\prod_iZ_i\,z^a},
                         \qquad a=1,2.                       \tag{B.80.2}
\]
To see the precise regularity conditions, let
$D_i=q_i^{-1}(T_i)$ and $H_i=q_i^{-1}(O_i)$. The line $Z_i=0$
cuts the divisor $2H_i$ on $C$, and $\operatorname{div}(G)$ is
$3\sum_i(D_i-H_i)$. The differential $dx/w^3$ represents the
canonical plane section. Therefore (B.80.2) is regular exactly when
the section $H$ vanishes along
\[
                         \sum_i(aD_i+(2-a)H_i).
\]
For $a=1$ these are simple conditions at $D_i$ and double conditions
at the hyperflex supporting $H_i$; for $a=2$ they are double conditions
at $D_i$. All these conditions are explicit linear equations in the
coefficients of $H$. Each solution space has dimension two:
Riemann--Roch gives $g(C)-1=2$ for the canonical bundle tensored
with a nontrivial degree-zero three-torsion line bundle.
The programs construct two independent forms, verify every vanishing
and derivative condition, and hence obtain a full basis for each
character. Together they give the new four-dimensional differential
space; the genus of each cover is seven.

The first program computes Cartier by extracting the coefficients
numbered $5j+4$ after clearing denominators and taking fifth roots.
Every division and basis-coordinate identity is checked with zero
remainder. A separate program imports no producer code and checks
the different rational-function identity
\[
 \left(\frac{\mathcal C(v\,dx)}{dx}\right)^5
                          =-\left(\frac d{dx}\right)^4v.
                                                               \tag{B.80.3}
\]
Identity (B.80.3) follows by writing
$v=\sum_{j=0}^4v_j^5x^j$: the fourth derivative kills every term
except $4!v_4^5=-v_4^5$. Thus it determines Cartier exactly.
In our function field $dw/dx=(3x^2-1)/(4w^3)$; the verifier applies
this derivation to every component in the basis $1,w,w^2,w^3$.
All $1456$ Cartier columns and all $5824$ component identities pass.
It also checks the eightfold semilinear product in the opposite
implementation order and calculates the characteristic polynomial
by a separate recursive adjugate formula. The result is (B.80.1) for
every one of the $364$ covers. The exact field modulus, torsion
points, coordinate matrices, cover normalizations, differential
bases, Cartier matrices and their eighth iterates are all retained.

By duality, the linear Cartier iterate and coherent Frobenius have
the same characteristic polynomial. The Artin--Schreier sequence
identifies the Frobenius-stable coherent part with the geometric
five-torsion cohomology, while its nilpotent part contributes no
factor to the determinant at one. Since the cover degree is prime
to five, pullback and norm split the old and new five-primary data.
Thus the relative class number is a unit exactly when the polynomial
in (B.80.1) is nonzero at one. This is true precisely in the two-entry
case. Under a further constant extension of degree $n$, the two
primitive cube-root eigenvalues acquire an eigenvalue one exactly
when $3\mid n$. A pre-existing eigenvalue one in either of the
other cases stays one under every power. This proves all extension
claims, and three is the smallest admissible odd prime in
characteristic five.
\end{proof}
\paragraph{Check record.} Checked only in this writing as a promoted statement; GPT-6 Astra Pro; 2026-11-18. Both supplied sources inspected; producer and separate derivative verifier rerun; mathematical interpretation checked. Effort not separately exposed.


\paragraph{Check record.}
Checked only in this writing as a promoted statement: GPT-6 Astra Pro,
finite-checks reader, 18 November 2026. We read both complete program
sources and verified the mathematical interpretation above, reran the
producer and its independent fourth-derivative verifier, and compared
all $364$ results. These are reruns of the supplied programs, not a
claim that their code was newly independently reimplemented here.
The effort setting is not separately exposed by the runtime.

\paragraph{Remaining scope of B6.}
This theorem concerns exactly the displayed pointed quartic. It does
not classify every quartic over $\mathbf F _3$ or $\mathbf F _5$,
and does not compare the unrestricted existence of abelian and
nonabelian witnesses for arbitrary curves. Those quantifiers are
retained in B.32.

\subsection{A complete pointed input census and first character fields}

\begin{claim}{B.81}{Proved}{Pointed quartic representatives and the first possible character fields}
There are exactly $3001$ pointed smooth plane-quartic isomorphism
classes over $\mathbf F _3$, and exactly $94283$ over $\mathbf F _5$.
Here a point is part of the object and is rational over the stated
base field; isomorphisms preserve it. Complete equations, point
counts over the first three extensions, and Weil polynomials are
retained for every class.

For each such pointed curve let $m_\ell$ be the least positive
integer $m$ for which a nonzero character
$J_C(\mathbf F _{q^m})\twoheadrightarrow\mathbf Z/\ell$ exists.
For the smallest admissible odd prime, namely $\ell=5$ when $q=3$
and $\ell=3$ when $q=5$, the exact distributions are
\[
\begin{array}{c|rrrrrr}
q=3:\ m_5&1&2&3&6&8&24\\ \hline
\text{classes}&650&515&49&86&327&350\\[2pt]
m_5&31&56&62&72&168&504\\ \hline
\text{classes}&255&47&261&79&73&309
\end{array}
\]
and
\[
\begin{array}{c|rrrrr}
q=5:\ m_3&1&4&8&13&56\\ \hline
\text{classes}&41623&5163&19699&14400&13398.
\end{array}
\]
These are first fields of existence of a nonzero constant character,
not first fields of a class-number witness.
\end{claim}

\begin{proof}
A smooth plane quartic is canonically embedded, so its abstract
isomorphisms are exactly the projective linear transformations of
its equation. Send the marked point to $P=(0:0:1)$ and its tangent
to $Y=0$, then normalize the coefficient of $YZ^3$ to one. Its
unique general form is
\[
\begin{split}
 F={}&YZ^3+(aX^2+bXY+cY^2)Z^2\\
     &+(dX^3+eX^2Y+fXY^2+gY^3)Z+Q_4(X,Y).
\end{split}                                                   \tag{B.81.1}
\]
The stabilizer is
\[
 (X,Y,Z)\longmapsto(AX+BY,CY,DX+EY+Z),\qquad A,C\ne0,
                                                               \tag{B.81.2}
\]
followed by division of the polynomial by $C$. Over $\mathbf F _3$
this group has order $108$. Enumerate all $3^{12}=531441$ coefficient
vectors and mark their full orbits under (B.81.2).

For $\mathbf F _5$ the following three slices suffice:
\[
\begin{array}{c|c|c|c}
\text{tangent contact}&\text{fixed coefficients}&\text{free}&w\\ \hline
2&a=1,\ b=c=d=0&8&2\\
3&a=0,\ d=1,\ b=c=e=0&7&3\\
4&a=d=0,\ u=1,\ b=c=v=0&6&4,
\end{array}
\]
where $u,v$ are the coefficients of $X^4,X^3Y$ in $Q_4$.
Here is a proof of completeness and uniqueness up to the stated
residual group. The transformed first coefficients are
\[
\begin{aligned}
a'&=aA^2/C,&b'&=3D+2aAB/C+bA,\\
c'&=3E+aB^2/C+bB+cC,&d'&=2aA^2D/C+dA^3/C.
\end{aligned}
\]
If $a\ne0$, choose $C=aA^2$, $D=-dA/(2a)$,
$B=-A(3D+bA)/2$, then the unique $E$ for which $c'=0$.
If $a=0,d\ne0$, choose $C=dA^3$, $D=-bA/3$, then the unique
$B$ making $e'=3D^2+2bAD+3B/A+eA^2$ zero, and again the
unique $E$ making $c'=0$.
If $a=d=0$, smoothness forces $u\ne0$, since otherwise the tangent
line is a component. Choose $C=uA^4$, $D=-bA/3$, then the unique
$B$ making $v'=D^3+bAD^2+eA^2D+4B/A+vA^3$ zero and the unique
$E$ making $c'=0$. All denominators are nonzero in characteristic
five. Conversely these coefficient equations show that a map between
two forms of the same slice has $C=A^w$ and $D=B=E=0$.
Its action on the coefficient of $X^iY^jZ^{4-i-j}$ is therefore
multiplication by $A^{i+wj-w}$. The residual group has order four.
This reduces the enumeration to $5^8+5^7+5^6=484375$ forms under
four diagonal transformations. The contact orders distinguish the
three slices intrinsically.

Smoothness is tested over the algebraic closure, not merely at
rational points. In three variables, multiply each of the three
cubic partial derivatives by every degree-four monomial and form
the resulting $45\times36$ coefficient matrix in degree seven.
The curve is smooth exactly when this matrix has rank $36$.
Indeed Euler's identity has invertible coefficient four in both
characteristics. A common nonzero zero of the partials makes degree
seven nonzero in their homogeneous quotient. In the absence of such
a zero the three cubic partials form a regular sequence in the
three-variable polynomial ring. Its quotient has Hilbert series
$(1+T+T^2)^3$ and hence no degree-seven part. These arguments prove
both directions of the rank test.

The complete enumeration gives $3001$ smooth orbits, with total
orbit size $315360$, over $\mathbf F _3$. Over $\mathbf F _5$ the
three slices give $75418,15626,3239$ smooth orbits. Their sum is
$94283$. An orbit of size $s$ in a slice corresponds to $500s$
forms in (B.81.1), because the full stabilizer has order $2000$.
The weighted total is $187512000$. These identities are checked
by the program itself, and every smooth representative, its orbit
size and contact order is retained in the output tables.

For each representative, direct point enumeration gives
$N_i=\#C(\mathbf F _{q^i})$ for $i=1,2,3$. The implementation
uses the chart $Y=1$, in which (B.81.1) is monic cubic in $Z$;
precomputed exact cubic-root counts over the six small fields
avoid repeated polynomial solving. The chart $Y=0$ is counted
separately. The complete genus-three Frobenius polynomial is
\[
 H(U)=U^6+a_1U^5+a_2U^4+a_3U^3+qa_2U^2+q^2a_1U+q^3,
\]
where, putting $s_i=q^i+1-N_i$, Newton's identities give
\[
 a_1=-s_1,\qquad
 2a_2=s_1^2-s_2,\qquad
 3a_3=-(s_3+a_1s_2+a_2s_1).
\]
These are integer identities, so the divisions are taken before
reduction modulo $\ell$.

A nonzero quotient to $\mathbf Z/\ell$ exists exactly when
$\ell$ divides the finite abelian group order
$\#J_C(\mathbf F _{q^m})$. Modulo $\ell$ this order is
$\prod_j(1-\lambda_j^m)$, where the $\lambda_j$ are the roots
of $H$ in an algebraic closure of $\mathbf F _\ell$. Therefore
\[
 m_\ell=\min_f\operatorname{ord}(U\bmod f),                  \tag{B.81.3}
\]
with $f$ ranging over the distinct monic irreducible factors of
$H$ modulo $\ell$. This does not assume semisimplicity: the
singularity of $F^m-1$, or simply its determinant, is enough.
For a factor of degree $d$, the root order divides $\ell^d-1$;
exact modular powering, removing prime divisors from that integer,
computes the order. Every factor and order is retained. Applying
(B.81.3) to the full tables gives the stated distributions.
The pointed Lang pullback realizes each resulting quotient with
its unique split fibre at $P$, so the character existence has its
usual geometric interpretation.
\end{proof}
\paragraph{Check record.} Census totals also checked in this writing; first-character-field distributions checked only in this writing. GPT-6 Astra Pro, 2026-11-18; independently reconstructed orbit and smoothness enumeration, exact extension point counts and polynomial-order calculations; effort not separately exposed.


\paragraph{Check record and boundary.}
Also checked in this writing for the census totals carried from
AS-A.72--AS-A.73, and checked only in this writing for the complete
first-character-field distributions: GPT-6 Astra Pro, 18 November
2026, finite-checks reader. The programs
\path{code/pointed_quartic_census.cpp},
\path{code/quartic_weil_census.cpp}, and
\path{code/quartic_character_fields.py} were written from the
mathematical descriptions, without access to the absent prior
census programs or their large data files. Their equations, all
point counts, Weil polynomials, factors and root orders are retained.
The effort setting is not separately exposed.

The large first degrees $56$ and $504$ show why a search through a
few constant extensions cannot decide B6. This statement does not
compute the required Cartier action on the new Jacobian parts at
those fields. It also concerns quartics with a point over the base
field; curves acquiring their first point only after extension
remain part of the uncompleted universal question.

\subsection{The completed characteristic-two cyclic-seven search}

\begin{claim}{B.82}{Proved}{No supersingular factor in the bounded rational-branch family}
Let $H/\overline{\mathbf F}_2$ be a smooth proper connected curve
with a cyclic degree-seven map to the projective line. Suppose its
branch locus is contained in $\mathbf P^1(\mathbf F _8)$ and has at
most seven points. Then $J_H$ has no positive-dimensional
supersingular isogeny factor, even geometrically. Consequently $H$
has no nonconstant map to a supersingular curve of positive genus.
For the normalized branch patterns with four through seven branch
points, the exhaustive counts are respectively
$620,2775,6666,6665$, and every one fails the supersingular-factor
test. This settles precisely the finite search family, not B5.
\end{claim}

\begin{proof}
Put $\mathbf F _8=\mathbf F _2[\alpha]/(\alpha^3+\alpha+1)$.
There are at least two unbranched rational points on the projective
line. A projective linear transformation over $\mathbf F _8$ can
send two of them to $0,\infty$ and one branch point to $1$.
Kummer theory, since seven is prime to two, and a change of cyclic
generator give the geometric model
\[
 y^7=\prod_{i=0}^6(x-\alpha^i)^{e_i},\qquad
 e_0=1,\quad e_i\in\{0,1,\ldots,6\},\quad
                       \sum_i e_i=0\pmod7.                 \tag{B.82.1}
\]
Indeed valuations give the exponents modulo seven, and a rational
function with all valuations divisible by seven is a constant times
a seventh power on the projective line. A constant becomes a seventh
power over the algebraic closure. Thus this normalization loses no
geometric isomorphism class in the stated family. Let $r$ be the
number of nonzero exponents. The cases $r\le2$ have genus zero or
are impossible; the case $r=3$ is treated below without a search.

Let $\zeta$ be a primitive seventh root of unity and let
$\chi_m:\mathbf F _{8^m}^{\times}\to\langle\zeta\rangle$
be the norm-compatible character with $\chi_1(\alpha)=\zeta$.
At an unramified finite point use $\chi_m(f(x))$, at a ramified
point use zero, and at infinity use one, because the degree in
(B.82.1) is divisible by seven and its leading coefficient is one.
Write $S_m$ for the sum of these values. The nontrivial character
factor is
\[
 L(T)=\exp\left(\sum_{m\ge1}S_mT^m/m\right)
     =\sum_{j=0}^{r-2}a_jT^j,
 \qquad a_0=1,\qquad
             j a_j=\sum_{m=1}^jS_ma_{j-m}.                 \tag{B.82.2}
\]
Here is a direct degree bound. The finite Euler product is the
Dirichlet $L$-series of the character obtained by evaluating monic
polynomials at the $r$ branch points and taking the prescribed
powers of the residue characters. For degree at least $r$, the
values at these points range uniformly and independently, so the
coefficient sum vanishes by the nontrivial character sum in any one
coordinate. The character is trivial on constants because
$\sum e_i=0$ modulo seven. Its finite $L$-polynomial therefore has
the factor $1-T$: sum over all polynomials of degree below $r$,
first grouping by their nonzero scalar multiple, to obtain zero
again by uniform evaluation. Restoring the unramified Euler factor
at infinity divides by $1-T$. This proves $\deg L\le r-2$.
The six nontrivial factors are its conjugates under
$\zeta\mapsto\zeta^j$, $1\le j\le6$; their product is the zeta
numerator of $H$, since the invariant factor is the zeta function
of the projective line. These statements also follow directly by
multiplying the seven local Kummer Euler factors. Thus the finite
sums through $m=r-2$ determine every root that needs testing.

We next prove that the test list is complete. A supersingular
Frobenius root $\lambda$ over $\mathbf F _8$ has
$\lambda/\sqrt8$ a root of unity. To justify this implication,
all its finite valuations, after this division, are zero: the
Newton slopes at two are $1/2$, and the other valuations vanish
because a Weil number has norm a power of two. All complex
absolute values are one. An algebraic unit with this property is
a root of unity: the monic polynomials of its powers have bounded
integer coefficients and bounded degree, so only finitely many
powers can occur among their roots. This is the needed elementary
form of Kronecker's argument.

Put $K=\mathbf Q(\zeta)$. If $\lambda$ is a root of one of our
character factors, then $[K(\lambda):K]\le5$. The root of unity
$\beta=\lambda^2/8$ has order $N$, whence
\[
 [K(\beta):K]=\frac{\varphi(\operatorname{lcm}(7,N))}{6}\le5.
\]
Write $N=7^a m$ with $7\nmid m$. Necessarily $a\le1$ and
$\varphi(m)\le5$, so
\[
                 m\in\{1,2,3,4,5,6,8,10,12\}.
\]
After multiplying $\lambda$ by a power of $\zeta$, its square
divided by eight has order $m$. The resulting number belongs to
a cyclotomic field ramified only at primes dividing $2m$. That
field is linearly disjoint from $K$: every nontrivial subfield of
$K$ is ramified at seven. Its degree over $K$ is therefore its
degree over $\mathbf Q$.

Factoring
$\operatorname{Res}_Z(\Phi_m(Z),U^2-8Z)$ for the nine displayed
values of $m$ now gives exactly the following possible minimal
polynomials of degree at most five:
\[
\begin{gathered}
 U^2-8,\quad U^2+8,\quad U^2\pm4U+8,\\
 U^4\pm8U^2+64,\qquad
 U^4\pm4U^3+8U^2\pm32U+64,
\end{gathered}                                                \tag{B.82.3}
\]
where the two signs in the last expression agree. The remaining
factors, for $m=5,8,10$, have degree eight and are irreducible.
These small rational factorizations, with every coefficient, are
also checked symbolically in the certificate. The listed factors
are irreducible over $K$ by the same disjointness argument.
They all have degree at least two, which proves the assertion for
$r=3$ immediately.

For each of the $16726$ normalized exponent vectors with $r\ge4$,
the program computes (B.82.2) in the integral basis
$1,\zeta,\ldots,\zeta^5$ and tests divisibility by the reciprocal
polynomial of each of (B.82.3), after all seven changes
$T\mapsto\zeta^jT$. No division succeeds. The complete character
polynomials and every power-sum residue count are retained, not
merely the absence of a search hit. The exact arithmetic uses the
fields of degrees $3,6,9,12,15$ over $\mathbf F _2$ with respective
binary moduli $11,67,529,4179,32771$; the full nonzero power table
checks that each proposed modulus gives a field and each chosen
generator has full order. The embeddings of $\alpha$ and norm
characters are recorded as well. This exhausts (B.82.1) and (B.82.3),
so no Frobenius root is supersingular.

The conclusion persists under every finite constant extension:
if $\lambda^n/8^{n/2}$ were a root of unity, then
$\lambda/\sqrt8$ would be one. Finally a nonconstant map to a
positive-genus supersingular curve induces pullback and norm on
Jacobians whose composition is multiplication by the degree.
Its Jacobian would therefore be a positive-dimensional
supersingular isogeny factor, a contradiction. This argument also
covers inseparable morphisms.
\end{proof}
\paragraph{Check record.} Checked only in this writing; GPT-6 Astra Pro, 2026-11-18. Supersingular factor-list completeness proved; original search inspected, completed and rerun; full character polynomials retained. This is an adapted supplied implementation, not a blind independent program; effort not separately exposed. Independent verifier checked all Newton identities and 841596 nondivisibility cases modulo 29 or 43.


\paragraph{Check record and exact boundary.}
Checked only in this writing: GPT-6 Astra Pro, finite-checks reader,
18 November 2026. We completed the supplied search, proved completeness
of the previously heuristic eight-factor list, and retained all
$16726$ character polynomials. The program
\path{code/b5_cyclic7_complete.py} is an inspected adaptation of the
supplied search, not a blind independent implementation. Its full
output and compressed polynomial certificate are retained. A separate verifier, \path{code/b5_cyclic7_verify.py}, imports no
producer code. It checks all Newton identities by cyclic convolution,
then proves nondivisibility by descending polynomial division after
specializing $\zeta$ at split primes. Of its $841596$ nondivisibility
checks, $840805$ have a nonzero remainder modulo $29$ and the remaining
$791$ have one modulo $43$. The leading coefficients are units in
these fields, so each nonzero remainder is a valid characteristic-zero
nondivisibility certificate. The effort
setting is not separately exposed. Branch sets with eight or nine
$\mathbf F _8$-points, branch points over larger fields, and arbitrary
prime-to-two correspondences are outside this statement. In particular
this negative finite family neither proves nor refutes B5.

\subsection{Every rational branch pattern over the field of eight elements}

\begin{claim}{B.83}{Proved}{The complete cyclic-seven rational-branch exclusion}
Let $H/\overline{\mathbf F}_2$ be smooth, proper and connected, with
a cyclic degree-seven map to $\mathbf P^1$. If, after a projective
change of coordinate, its branch locus is contained in
$\mathbf P^1(\mathbf F _8)$, then $J_H$ has no positive-dimensional
geometric supersingular isogeny factor. In particular $H$ has no
nonconstant map to a positive-genus supersingular curve.

For eight and nine branch points there are respectively $745$ and
$511$ branch-exponent orbits under $\operatorname{PGL}_2(\mathbf F _8)$
and changes of cyclic generator. Their character polynomials all
fail the complete supersingular-factor test. Together with B.82,
these exhaust every possible rational branch pattern over
$\mathbf F _8$.
\end{claim}

\begin{proof}
Cases with at most seven branch points are B.82. Label the nine
points of $\mathbf P^1(\mathbf F _8)$ by $0,1,\alpha,1+\alpha,
\alpha^2,1+\alpha^2,\alpha+\alpha^2,1+\alpha+\alpha^2,\infty$,
where $\alpha^3+\alpha+1=0$. A Kummer cover is specified geometrically
by a nonzero vector
\[
 (e_0,\ldots,e_8)\in\mathbf F _7^9,
 \qquad \sum_i e_i=0,
\]
up to multiplication by $\mathbf F _7^\times$; zero entries are the
unbranched points. Constants in the Kummer equation do not change
the geometric isomorphism class. The projective group permutes the
nine coordinates. For eight branch points move the unique unbranched
point to infinity, and then normalize $e_0=1$. For nine branch points
all entries are nonzero, so the same scalar normalization is available.
Thus representatives have equations
\[
                 y^7=\prod_{b\in\mathbf F _8}(x-b)^{e_b},
                 \qquad e_\infty=-\sum_b e_b\pmod7.        \tag{B.83.1}
\]

The orbit enumeration is a finite certificate with an elementary
coverage proof. Enumerate the $7^8$ assignments to the finite
coordinates and recover the infinity coordinate from their sum.
Keep vectors of the required support size and normalize their first
nonzero coordinate to one. Enumerate all invertible two-by-two
matrices over $\mathbf F _8$, identify their induced permutations,
and obtain exactly $8(8^2-1)=504$ permutations. Mark all normalized
images of each unmarked vector. This partitions the normalized set
into the listed orbits. The independent size check is
\[
 \#\{\text{scalar-normalized vectors with support }r\}
  =\binom9r\frac{6^r+6(-1)^r}{42}.                          \tag{B.83.2}
\]
Indeed the number of nonzero $r$-tuples of sum zero is
$(6^r+6(-1)^r)/7$, either by additive-character orthogonality or the
recurrence separating the last entry; scalar multiplication is free.
The orbit-size sums are $359919$ for $r=8$ and $239945$ for $r=9$,
in exact agreement with (B.83.2). The complete $1256$ representatives
and their orbit sizes are retained.

Put $K=\mathbf Q(\zeta_7)$. Use the norm-compatible character of
B.82 to calculate the power sums $S_m$ and the character polynomial
$L(T)=\sum_{j=0}^{r-2}a_jT^j$. The degree bound in B.82 applies
unchanged when infinity is unramified. If infinity is ramified,
there are eight finite branch points and the residue character on
constants is nontrivial: its exponent sum is
$-e_\infty\ne0$. Uniform evaluation makes every finite
Dirichlet coefficient of degree at least eight vanish. There is no
Euler factor at infinity to restore, and the degree bound is seven,
again $r-2$. Thus in both cases the sums through $m=r-2$ determine
the entire polynomial by
\[
                j a_j=\sum_{m=1}^j S_m a_{j-m}.             \tag{B.83.3}
\]
At a ramified point the character value is zero. In the eight-point
case the infinity contribution is one because (B.83.1) is monic
of degree divisible by seven; in the nine-point case it is zero.

We must extend, rather than assume, the completeness of the factor
list in B.82. If a supersingular Frobenius root $\lambda$ occurs,
then $[K(\lambda):K]\le7$, and $\beta=\lambda^2/8$ is a root of
unity by the algebraic-unit argument in B.82. If its order is
$N=7^a m$ with $7\nmid m$, then
\[
             \frac{\varphi(\operatorname{lcm}(7,N))}{6}\le7.
\]
For $a\le1$, one has $\varphi(m)\le7$, which gives exactly
\[
                 m\in\{1,2,3,4,5,6,8,9,10,12,18\}.
\]
Remove the seventh-root phase from $\lambda$. As in B.82, its
remaining cyclotomic field is ramified away from seven and is
linearly disjoint from $K$. The only new possibilities are $m=9,18$.
For $m=9$, the field generated by the remaining root is
$\mathbf Q(\zeta_9,\sqrt2)$; for $m=18$ it is
$\mathbf Q(\zeta_9,\sqrt{-2})$. To see the equality, its square
generates $\mathbf Q(\zeta_9)$, and a square root of the root-of-unity
factor already belongs to that field, up to the indicated sign.
Both fields have degree twelve: the quadratic radical is ramified
at two whereas $\mathbf Q(\zeta_9)$ is not. They are excluded.

If $a=2$, the degree inequality forces $m=1$ or $2$.
The same argument now gives
$K(\lambda)=\mathbf Q(\zeta_{49},\sqrt{\pm2})$,
of degree fourteen over $K$, again impossible. Values $a\ge3$
already give degree at least forty-nine for $\beta$ over $K$.
Consequently the only possible minimal polynomials are still the
eight polynomials (B.82.3), with all seven possible phases. This is
a complete list for character degree at most seven.

For all $1256$ representatives the program evaluates every required
power sum directly, using the fields of degrees $3m$ over
$\mathbf F _2$, for $1\le m\le7$. The binary moduli are
\[
              11,\ 67,\ 529,\ 4179,\ 32771,\ 262273,
              \ 2097157.
\]
The last two are $X^{18}+X^7+1$ and $X^{21}+X^2+1$.
Full nonzero power tables certify field construction and generator
orders; the embeddings and character normalizations are saved.
All coefficients, all residue counts and all orbit representatives
are retained. No polynomial has any factor in the complete list.

An independent verifier checks every Newton identity and proves
all $70336$ required nondivisibilities by descending division at
split primes of $K$: $70279$ nonzero remainders occur modulo $29$,
and the remaining $57$ occur modulo $43$. The test-polynomial
leading coefficients are units at both primes. Thus any actual
characteristic-zero divisor would remain a divisor after either
reduction, so each nonzero remainder is conclusive.

This proves absence of supersingular Frobenius roots. The same
unit argument as in B.82 shows the conclusion persists under every
constant extension, and its pullback--norm argument excludes a
map to a positive-genus supersingular curve. All geometric twists
were already accounted for by the Kummer normalization.
\end{proof}
\paragraph{Check record.} Checked only in this writing; GPT-6 Astra Pro, finite-checks reader, 2026-11-18; effort not separately exposed. New PGL2(F8) orbit enumeration; independently verified all 599864 normalized vectors form disjoint orbits; proved completeness of the supersingular factor list through character degree seven. Every power sum was computed directly; independent verifier checked all Newton identities and 70336 modular nondivisibility certificates.


\paragraph{Check record and boundary.}
Checked only in this writing: GPT-6 Astra Pro, finite-checks reader,
18 November 2026; effort setting not separately exposed. New exact
orbit enumeration, complete character-factor computation, and an
independent modular verifier are in
\path{code/b5_all_f8_orbits.cpp},
\path{code/b5_all_f8_direct.py}, and
\path{code/b5_all_f8_direct_verify.py}, with full certificates.
The earlier faster calculation recovers high coefficients from
the functional equation; the direct calculation recomputes every
power sum and does not rely on that recovery. Branch loci over
larger fields and arbitrary odd-degree correspondences are outside
this finite family. B5 remains open in its full quantifiers.

\subsection{A branch-field-independent obstruction for cyclic-seven curves}

\begin{claim}{B.84}{Proved}{No supersingular factor for any four-branch cyclic-seven curve in characteristic two}
Let $k$ be an algebraically closed field of characteristic two, and
let $H/k$ be a smooth proper connected curve with a cyclic degree-seven
map $H\to\mathbf P^1_k$ having at most four branch points. Then the
Newton polygon of $J_H$ has no slope $1/2$. In particular $J_H$ has
no positive-dimensional supersingular isogeny factor, and $H$ has no
nonconstant map to a positive-genus supersingular curve.

In the four-branch case, the only inertia multisets whose signatures
could permit a supersingular Jacobian are opposite pairs
$\{u,-u,v,-v\}$ in $(\mathbf Z/7)^\times$. Every such curve is
ordinary. More explicitly, after normalization it has equation
\[
 y^7=x(x-1)^a(x-\lambda)^{7-a},\qquad
 a\in\{1,\ldots,6\},\quad \lambda\in k\setminus\{0,1\},
                                                               \tag{B.84.1}
\]
and the Cartier operator on its six regular differentials is
bijective for every allowed $\lambda$.
\end{claim}

\begin{proof}
We first make the Newton restriction precise, including characteristic
two. For a cyclic degree-seven cover of genus $3(r-2)$ with $r$
branch points, the nontrivial character spaces in first cohomology
all have dimension $r-2$, and the invariant space is zero. The latter
space is the cohomology of the projective line; the equal dimensions
also follow from the tame Euler characteristic, or from the explicit
differential dimensions below and their duals.

Modulo two,
\[
 \Phi_7(U)=(U^3+U+1)(U^3+U^2+1).
\]
Thus the nontrivial part of the two-divisible group and its rational
Dieudonn\'e module split into the two orbit pieces indexed by
\[
                A=\{1,2,4\},\qquad B=\{3,5,6\}.
\]
Each has height $3(r-2)$. Each Newton-slope multiplicity in one
piece is divisible by three, and the two pieces are dual, with
slopes changed by $s\mapsto1-s$. These are exactly the necessary
conditions (M) and (D) in Li--Mantovan--Pries--Tang, Theorem~4.6;
the standing hypothesis there is $p\nmid m$, so $p=2,m=7$ is
included. The slope-denominator divisibility is stated
in their \S2.3. We use only these necessary conditions and their
orbit decomposition, not a realization theorem.
\citep[\S2.3 and Theorem~4.6, conditions (D), (M)]{LMPT2019}

For clarity, the particular restrictions also have a direct explanation.
The two integral idempotents in $\mathbf Z_2[\zeta_7]$ split the
two-divisible group. On each piece, Frobenius permutes the three
embeddings of its unramified cubic coefficient ring. It therefore
permutes three equal-rank spaces of each fixed Newton slope, giving
the factor three. The polarization sends the action of a deck
generator to its inverse under Rosati, so it interchanges the two
pieces and gives the duality. Finally a slope with denominator two
has even multiplicity in any two-divisible group, by its integral
Newton break points. Consequently a slope $1/2$ in either piece
must have multiplicity at least six.

For $r=3$ each piece has height three, so there is no such slope.
For $r\le2$ the curve has genus zero or the branch pattern is
impossible. For $r=4$ both pieces have height six. If slope $1/2$
occurred, it would fill one entire piece and then its dual, making
$J_H$ supersingular. It remains to exclude this last possibility.

Write the four nonzero inertia exponents as $e_1,\ldots,e_4$,
with $\sum_i e_i=0$ modulo seven. The regular differential character
dimensions, with the character index sign chosen as below, are
\[
             n_j=-1+\sum_{i=1}^4
                        \left\langle\frac{j e_i}{7}\right\rangle,
                   \qquad 1\le j\le6.                       \tag{B.84.2}
\]
This formula can be checked without a characteristic-zero comparison:
in a Kummer equation with finite branch exponents $e_i$, multiply
$dx/y^j$ by $\prod_i(x-t_i)^{\lfloor je_i/7\rfloor}$, and count
the permitted powers of $x$ at infinity. At a finite branch the
orders of $x-t_i,y,dx$ are $7,e_i,6$; at infinity the same count
includes the fourth exponent. It gives exactly (B.84.2).

Put $n_A=\#\{i:e_i\in A\}$. Since the sum of the least positive
residues in $A$ is seven and that in $B$ is fourteen, (B.84.2) gives
\[
       \sum_{j\in A}n_j=5-n_A,
       \qquad \sum_{j\in B}n_j=1+n_A.                         \tag{B.84.3}
\]
These are, up to interchange arising from the covariant convention,
the dimensions of the two height-six two-divisible-group pieces.
One can see the endpoint statement directly on their Dieudonn\'e
lattices: the valuation of the determinant of Frobenius is the
length of its lattice cokernel, which is the dimension of the
corresponding piece. Equivalently the Hodge and Newton polygons
of each orbit piece have the same endpoints, as in condition (WA)
of the same Theorem~4.6. A height-six supersingular piece has
dimension three. Therefore supersingularity would force $n_A=2$.

There are exactly two exponents from $A$ and two from $B$. Their
integer sum lies strictly between seven and twenty-one, so it is
fourteen. The six sums of unordered pairs from $A$ are
$2,3,4,5,6,8$, all distinct. The pair from $B=-A$ must therefore
be the negatives of the pair from $A$. This proves that the multiset
is $\{u,-u,v,-v\}$. Changing the cyclic generator and moving three
branch points to $0,1,\infty$ gives (B.84.1).

We now compute Cartier over the rational function field in the
arbitrary parameter $\lambda$. For $1\le j\le6$ put
\[
 b_j=\left\lfloor\frac{aj}{7}\right\rfloor,
 \qquad h_j=(x-1)^{b_j}(x-\lambda)^{j-1-b_j},
 \qquad \omega_j=\frac{h_j\,dx}{y^j}.
\]
The finite-branch valuation check just given proves regularity there.
At infinity, $x,y,dx$ have orders $-7,-8,-8$, respectively, and
$\deg h_j=j-1$. Thus the order of $\omega_j$ at infinity is
$j-1\ge0$. The six forms have distinct deck characters, and
Riemann--Hurwitz gives genus six, so they are a full basis.

Let $j'=j/2$ if $j$ is even and $j'=(j+7)/2$ if $j$ is odd;
put $\epsilon=0$ or $1$, respectively. With
$f=x(x-1)^a(x-\lambda)^{7-a}$ we have
\[
 \mathcal C(\omega_j)
   =\frac1{y^{j'}}\mathcal C(h_jf^\epsilon dx)
   =c_{a,j}\,\omega_{j'},
 \qquad
        c_{a,j}^{\,2}h_{j'}^{\,2}
                   =\frac{d}{dx}(h_jf^\epsilon).              \tag{B.84.4}
\]
The last identity is the characteristic-two polynomial Cartier
formula: differentiation retains the odd powers, and their square
root is the Cartier coefficient polynomial. Every entry in the
following table is a polynomial identity in $\mathbf F _2[\lambda,x]$,
not a test at finitely many values of $\lambda$:
\[
\begin{array}{c|cccccc}
a\backslash j&1&2&3&4&5&6\\ \hline
1&1&1&1&1&1&1\\
2&\lambda&1&\lambda&1&1&1\\
3&1&1&\lambda&1&1&1\\
4&\lambda&1&1&1&\lambda&1\\
5&1&1&1&1&\lambda&1\\
6&\lambda&1&\lambda&1&\lambda&1
\end{array}                                                   \tag{B.84.5}
\]
The entries are $c_{a,j}^2$. Expanding the two factors in $h_j$,
differentiating, and using $j+7\epsilon=2j'$ proves all thirty-six
identities; their complete monomial lists and exact checks are
retained in the symbolic certificate. Since $\lambda\ne0$, every
scalar is nonzero. Cartier permutes the six lines bijectively, so
its stable rank is six. Here is the comparison being used. Serre duality
identifies Cartier with the dual of Frobenius on $H^1(H,\mathcal O_H)$.
The Artin--Schreier sequence identifies
$H^1_{\text{ét}}(H,\mathbf F _2)$ with the fixed vectors of
that Frobenius, because $u\mapsto u^2-u$ is surjective on $k$.
For an invertible Frobenius-semilinear map on a $d$-dimensional
$k$-space, its fixed-point map has exactly $2^d$ solutions:
in coordinates $A z^{[2]}-z$ is finite of degree $2^d$ (its leading
homogeneous equations have no nonzero common zero) and is \'etale
with derivative $-I$. Thus the two-rank is six. The slope-zero
multiplicity is the two-rank; polarization supplies six slope-one
dimensions as well. Hence $H$ is ordinary and cannot be supersingular.

We have excluded slope $1/2$ in every case. A supersingular isogeny
factor would contribute that slope. A nonconstant map to a
positive-genus supersingular curve would provide such a factor by
pullback and norm, whose composition is multiplication by the
nonzero degree, including in the inseparable case. This proves
all assertions.
\end{proof}
\paragraph{Check record.} Checked only in this writing; GPT-6 Astra Pro, finite-checks reader, 2026-11-18; effort not separately exposed. Reproved orbit-height and inertia-signature obstruction; all 186 ordered signatures and 36 symbolic Cartier identities checked over F2[lambda, x]. Opened LMPT2019 primary full text and checked p not dividing m, hence characteristic two is included. Additional finite checks over F64 and F512 are not used for the universal deduction.


\paragraph{Check record and scope.}
Checked only in this writing: GPT-6 Astra Pro, finite-checks reader,
18 November 2026. The primary theorem was opened with its exact
hypothesis $p\nmid m$. We wrote the orbit-height and signature
argument and verified all $186$ ordered inertia signatures and all
$36$ symbolic Cartier identities in
\path{code/b5_fourbranch_cartier.py} and its JSON certificate.
The effort setting is not separately exposed. Independent finite
character sums over $\mathbf F _{64}$ and $\mathbf F _{512}$,
covering $1922$ and $15810$ normalized cases, also give no
supersingular factors; these are checks of the theorem, not its
universal proof. The result excludes this entire cyclic-seven
four-branch construction. It does not exclude other odd \'etale
covers or settle the universal correspondence assertion B5.

\subsection{Two exact genus-two loci with supersingular correspondences}

\begin{claim}{B.85}{Proved}{The ordinary cubic locus and an explicit cubic witness for the curve of S2}
Let $k$ be algebraically closed of characteristic two, and let $C$ be
the smooth proper model of
\[
 z^2+z=a x+\frac b x+\frac c{x+1},\qquad abc\ne0.                 \tag{B.85.1}
\]
This is an ordinary genus-two curve. It has a connected cyclic
\'etale cover of degree three with primitive two-rank zero if and
only if
\[
                         (a+b+c)^3=abc.                        \tag{B.85.2}
\]
When (B.85.2) holds and $a,b,c$ are not all equal, exactly one of
the forty geometric cyclic-three covers has primitive two-rank
zero. When $a=b=c$, exactly two do. Every other cyclic-three
cover has ordinary primitive part. Thus every curve in (B.85.2)
satisfies B4 and B5, with an actual map from its cubic cover to a
supersingular elliptic curve.

Here are explicit defining functions. Choose $r\notin\{0,1\}$
and $\sigma$ with
\[
 b=a r^3,\qquad c=a(r+1)^3,\qquad
 \sigma^2+\sigma=a(r+1),\qquad h=x(x+1).
\]
Then the normalization of
\[
                    w^3=\frac{b(z+\sigma)}h                   \tag{B.85.3}
\]
is one of the required covers. If the displayed coefficients and
$\mu_3$ belong to a finite field $L$, each of the three points
$W_0,W_1,W_\infty$ above $0,1,\infty$ splits over $L$.
Its relative class number is odd over $L$ and every finite
extension of $L$.

In particular, for the S2 curve
$y^2+x(x+1)y=x^5+1$, set
\[
 z=\frac{y+x+1}{x(x+1)},\qquad
 L=\mathbf F _2(\sigma),\quad \sigma^4+\sigma+1=0,
 \qquad r=\sigma^2+\sigma+1.
\]
The cubic cover (B.85.3), with $a=b=c=1$, is defined over
$L=\mathbf F _{16}$, is split at both $P=(0,1)$ and $Q=(1,0)$,
and has genus four. Its primitive zeta numerator is
$(1-4T)^4$, so its exact relative class number is $81$.
An explicit nonconstant map to the supersingular elliptic curve
$E:Y^2+Y=T^3$ is
\[
 U=w+\frac{x+r}{x(x+1)w},\qquad
 T=\frac1{U+1},\qquad Y=\frac{xU^2+r}{(U+1)^2}.               \tag{B.85.4}
\]
All these statements concern the smooth normalizations; the
affine Kummer equation alone has singularities at its zero and
poles.
\end{claim}

\begin{proof}
We first give the local and cohomological facts used in both this
claim and B.86. For an Artin--Schreier double cover with a reduced
pole of odd order $m$, write the base parameter as $t$. Its
valuations upstairs are $v(t)=2$ and $v(z)=-m$, and
$\pi=t^{(m+1)/2}z$ is a uniformizer. The involution changes
$\pi$ by $t^{(m+1)/2}$; the different exponent is consequently
$m+1$, by the derivative formula for a monogenic extension of
complete discrete valuation rings. In (B.85.1) the three exponents
are two. Riemann--Hurwitz gives genus two and
\[
             \operatorname{div}(dx)=2W_0+2W_1-2W_\infty.
\]
The regular forms $dx/x$ and $dx/(x+1)$ are linearly independent
and fixed by Cartier. They form a basis, so $C$ is ordinary.

Let $F$ denote absolute Frobenius, and let $B$ be the line bundle
of locally exact differentials, with its $\mathcal O_C$-action
through squares. Locally $dt$ generates this line. A rational
exact differential $f^2dt$ has order $\operatorname{ord}(f)$ as
a section of $B$, half its usual differential order. Hence
\[
 B\simeq\mathcal O_C(W_0+W_1-W_\infty),\qquad B^2\simeq\omega_C.
                                                               \tag{B.85.5}
\]
We use the same letter for the divisor in (B.85.5). For every
nontrivial odd-order line bundle $M$ of degree zero, the exact
sequence of locally exact differentials, tensorized by $M$, is
\[
 0\longrightarrow M\longrightarrow F_*M^2
     \longrightarrow B\otimes M\longrightarrow0.
\]
Both left-hand degree-zero bundles have no nonzero global section.
Thus
\[
 H^0(C,B\otimes M)
   =\ker\bigl(F:H^1(C,M)\longrightarrow H^1(C,M^2)\bigr).     \tag{B.85.6}
\]
This is the rank-one form of the exact sequence and criterion
immediately preceding Proposition~1 in Raynaud's paper; the
absolute-Frobenius notation incorporates the relative twist.
The local calculation above also identifies the particular
theta characteristic used here. \citep[\S1, p.~74]{LA_Raynaud2000}

For a cyclic cover of prime odd degree $\ell$, Kummer theory gives
\[
 H^1(D,\mathcal O_D)=\bigoplus_{j=0}^{\ell-1}H^1(C,M^j),
\]
where $M$ has exact order $\ell$. Each nontrivial summand has
dimension one, by Riemann--Roch. Frobenius sends index $j$ to
$2j$. Within an index orbit its stable rank is the orbit length
if every arrow is nonzero, and zero if any arrow is zero. The
stable Frobenius rank is the two-rank: this follows from the
Artin--Schreier sequence, by restricting Frobenius to its eventual
image, where an invertible semilinear map in dimension $d$ has
$2^d$ fixed vectors. Indeed its fixed-point equations have
invertible linear derivative and degree $2^d$, with no nonzero
common zero at infinity. Since multiplication by two is transitive
on $(\mathbf Z/\ell)^\times$ for $\ell=3$ and $5$, the primitive
two-rank for either of those degrees is zero exactly when
\[
             H^0(C,B\otimes M^j)\ne0
                 \quad\hbox{for some }1\le j<\ell.             \tag{B.85.7}
\]
Otherwise the primitive part is ordinary. This is a full-orbit
argument; a single vanishing arrow need not kill every primitive
orbit for a general odd degree.

We now determine the three-torsion in (B.85.7). Put
$D_0=W_0+W_1+W_\infty$. The divisor $B$ is not effective.
Otherwise $B\sim R$ would give
$R+W_\infty\sim W_0+W_1$. The latter degree-two divisor is not
canonical, since the canonical pencil is the hyperelliptic pencil
and its ramified fibres are $2W_i$. Its space of sections therefore
has dimension one. The two effective divisors would have to be
equal, which their supports exclude. In particular every line
bundle $\mathcal O(R-B)$ is nontrivial.

We have
\[
 \operatorname{div}(h)=3B-D_0,
 \qquad L(D_0)=\langle1,z\rangle,
 \qquad L(3B)=\langle1/h,z/h\rangle.                          \tag{B.85.8}
\]
The dimensions are two by Riemann--Roch, and the asserted functions
have exactly the required poles. A point $R$ makes
$\mathcal O(R-B)$ three-torsion exactly when there is a function
with divisor $3R-3B$. The constant numerator $1/h$ in (B.85.8)
has effective residual divisor $D_0$, not a triple point. Thus
such a function must be a nonzero scalar times $(z+\sigma)/h$,
and the fibre $z=\sigma$ must be a triple point. With
$s=\sigma^2+\sigma$, its equation in $x$ is
\[
             a x^3+(a+s)x^2+(b+c+s)x+b=0.                    \tag{B.85.9}
\]
Comparison with $a(x+r)^3$ is exactly
$b=a r^3$, $c=a(r+1)^3$, and $s=a(r+1)$.
Since $abc\ne0$, the resulting $r$ is neither zero nor one.
The point $x=r$ is away from the branch locus, so a triple root
in (B.85.9) is a zero of order three on the smooth curve.

Eliminating $r$ gives (B.85.2). We verify sufficiency, including
the exceptional denominator. Put $d=b+c$, $n=a+c$, and
$H=(a+b+c)^3+abc$. Direct polynomial identities in characteristic
two are
\[
 a n^3+b d^3=(a+b)H,\qquad
 a(n+d)^3+c d^3=(a+c)H.                                      \tag{B.85.10}
\]
If $d\ne0$, set $r=n/d$ and obtain both required cubes.
This $r$ is unique, since
$a+c=r(b+c)$ follows from the two cube equations. If $d=0$,
(B.85.2) says $a^3=ab^2$, hence $a=b=c$. Then the two possible
$r$ are exactly the roots of $r^2+r+1=0$. For each $r$ there
are two choices of $\sigma$, differing by one; their points are
hyperelliptic conjugates. Since $B^2\simeq\omega_C$, their
classes $R-B$ are negatives. The Abel map on points is injective:
linear equivalence of two distinct degree-one points would give
a degree-one map to the projective line. We have therefore found
exactly one or two distinct cyclic-three subgroups, respectively.
There are $(3^4-1)/(3-1)=40$ such subgroups in the genus-two
Jacobian. Equations (B.85.6)--(B.85.7) prove all the asserted
primitive two-ranks.

For the explicit cover, (B.85.8)--(B.85.9) give
\[
 \operatorname{div}\left(\frac{b(z+\sigma)}h\right)=3(R-B).
\]
Its Kummer normalization is unramified everywhere: in a completed
local ring the exponent is divisible by three, and after removing
that power the remaining equation has unit right side and
invertible derivative. It is connected because a cube in $k(C)$
would make $R-B$ principal, already excluded. At each $W_i$ use
the uniformizer $t=1/z$. The leading coefficients of the right
side, after removing its valuation, are
\[
 1\quad(W_0),\qquad b/c=(r/(r+1))^3\quad(W_1),\qquad
 a^2b=(ar)^3\quad(W_\infty).
\]
They are cubes in $L$, so all three fibres split when $\mu_3
\subset L$. The connected kernel of the norm on Jacobians has
geometric two-rank zero. For a cyclic cover of odd prime degree,
its isogeny decomposition with $J_C$ has odd degree: pullback
and norm compose to multiplication by that prime. Hence their
two-divisible groups split. The relative zeta factor is the
Frobenius polynomial of the norm kernel, and its value at one
is odd over every finite extension, since a two-rank-zero
abelian variety has no nonidentity geometric point of order two.

For use in B5, we give the reflection argument in full. Let
$D\to C$ be any connected cyclic \'etale cover of odd prime degree
$\ell$ of a genus-two curve in characteristic two, with deck
generator $\delta$. The hyperelliptic involution $\iota$ acts as
inversion on $J_C$, since $R+\iota R$ is a canonical fibre.
It therefore preserves the cyclic Kummer subgroup and lifts to
$\tau$ on $D$ with $\tau\delta\tau^{-1}=\delta^{-1}$.
Write $\tau^2=\delta^e$. Conjugation by $\tau$ fixes its own
square but inverts $\delta^e$, so $2e=0$ modulo $\ell$ and
$\tau^2=1$.

Above each ramification point of $C\to\mathbf P^1$, this
reflection fixes exactly one of the $\ell$ points: its permutation
of the fibre has the form $i\mapsto c-i$. There are no other
fixed points. Because $D\to C$ is \'etale, the completed local
ring at a fixed point is equivariantly the same as the corresponding
completed local ring of $C$. Thus the total different of
$D\to H:=D/\langle\tau\rangle$ is the same six as for
$C\to\mathbf P^1$. Since $g_D=\ell+1$, Riemann--Hurwitz gives
\[
                         g_H=(\ell-1)/2.                     \tag{B.85.11}
\]
The pullback of $J_H$ to $J_D$ lies in the connected norm kernel:
its norm image is fixed by $\iota=-1$, hence killed by two, and
a connected abelian-variety image contained in a finite group
scheme is zero. Take also the reflection $\delta\tau$. The
intersection of the two pullback images is finite, since a point
fixed by both reflections is fixed by $\delta$ and, in the norm
kernel, is killed by $\ell$. Their dimensions sum to $\ell-1$,
the dimension of the norm kernel. The two reflection quotients
are isomorphic by a deck transformation, as two is invertible
modulo $\ell$. Consequently
\[
                         P(D/C)\sim J_H^2.                    \tag{B.85.12}
\]
For $\ell=3$ the target $H$ is elliptic; primitive two-rank zero
therefore makes it supersingular. For $\ell=5$ it has genus two,
and the same conclusion holds: in height four a nonzero slope
other than $1/2$ and $1$ would, together with its dual and the
denominator divisibility of slope multiplicities, require height
at least six. A genus-two two-rank-zero Jacobian has no slopes
zero or one, so all slopes are $1/2$. This proves the B5 assertion
here and the implication needed below, with odd-order Galois
closure already furnished by $D\to C$.

For S2, direct substitution gives (B.85.1) with $a=b=c=1$.
The points $P,Q$ are $W_0,W_1$. The polynomial
$\sigma^4+\sigma+1$ is irreducible over $\mathbf F _2$ and
$r^2+r+1=0$, so (B.85.3) is defined over the stated field and its
three ramified base points split. Its defining equations imply
\[
 x(x+1)U^3+(x+r)U+1=0.
\]
Indeed the two summands defining $U$ have cubes summing to
$1/h$ and product $(x+r)/h$. Substitution in (B.85.4), using
$r^2+r+1=0$, gives $Y^2+Y=T^3$. This rational map of smooth
proper curves extends everywhere. It is nonconstant: a constant
$U$ in the displayed cubic would force $x$ to satisfy a nonzero
constant-coefficient polynomial. It has degree two, as adjoining
$w$ to $k(x,U)$ is the reflection quadratic. More explicitly,
$U=1/T+1$ and $x=(Y(U+1)^2+r)/U^2$, so
$k(x,U)=k(T,Y)$. The target is a
smooth genus-one Artin--Schreier curve with a single pole of order
three; its regular differential is killed by Cartier, so it is
supersingular.

For the exact finite certificate, direct counts on the smooth
normalizations are
\[
\begin{array}{c|rrrr}
 m&1&2&3&4\\\hline
 \#C(\mathbf F _{16^m})&31&223&4111&65983\\
 \#D(\mathbf F _{16^m})&15&159&3855&64959.
\end{array}
\]
Each of $W_0,W_1,W_\infty$ contributes three points. At the
triple zero $(r,\sigma)$, replace $w$ locally by $w/(x+r)$;
the residue of its cube is $1/[r(r+1)]^2=1$, so there are
three points there as well. All other fibres are counted by
testing cubes. Newton's identities and the functional equation give
\[
 L_C(T)=1+14T+81T^2+224T^3+256T^4,
 \qquad L_D(T)=L_C(T)(1-4T)^4.
\]
Thus $\#J_C(L)=576$, $\#J_D(L)=46656$, and the relative
class number is $81$. The explicit reflection maps give an
independent interpretation: $P(D/C)$ is $L$-isogenous to $E^2$,
and $E/\mathbf F _{16}$ has Frobenius polynomial $(X-4)^2$.
No implication from section injectivity, and no converse to a
class-number criterion, is used in this construction.
\end{proof}
\paragraph{Check record.} Also checked in this writing; GPT-6 Astra Pro, 2026-11-18. Full local theta/Frobenius/reflection argument; exact symbolic identities; normalized Kummer counts; independent FLINT trace counts of base and target and integer zeta convolution. Independent verifier does not re-enumerate Kummer fibres. Effort not separately exposed. Root second reader completed a full proof reading, including theta exact sequence, elimination, normalized residues, reflection isogeny and finite zeta identities, on 2026-11-18.


\begin{claim}{B.86}{Proved}{The two-rank-one quintic locus and a supersingular genus-two target}
Let $k$ be algebraically closed of characteristic two and let $C$
be the smooth proper model of
\[
                  z^2+z=x^3+a x+b/x,\qquad b\ne0.              \tag{B.86.1}
\]
It has genus two and two-rank one. Every connected cyclic
\'etale cubic cover of $C$ has ordinary primitive part. A connected
cyclic \'etale quintic cover with primitive two-rank zero exists
if and only if
\[
                              b a^{16}=1.                    \tag{B.86.2}
\]
If (B.86.2) holds, there is exactly one such geometric cyclic-five
cover; every other cyclic-five cover has ordinary primitive part.
In particular the entire locus (B.86.2) satisfies B4 and B5.

For an explicit cover, put $u=1/x$, $v=z u^2$, $r=a^4$, choose
$d$ with $d^2+d=b a^4$, and use
\[
 v^2+u^2v=b u^5+a u^3+u,\qquad
              w^5=v+d u^2+a u+a^2.                           \tag{B.86.3}
\]
Its smooth normalization has genus six and a reflection quotient
of genus two with supersingular Jacobian.

For $a=b=1$, take $d=\rho$ with $\rho^2+\rho+1=0$.
There is the following explicit map to a supersingular genus-two
curve:
\[
 H:Y^2+Y=T^5+T^3+T+1,\qquad
 U=w+\frac{u+1}{w},\quad
 T=\frac1U,\quad Y=\frac{(U+1)u}{U^3}+\frac1{U^2}.           \tag{B.86.4}
\]
Over $\mathbf F _{16}$ the quintic deck group is constant, both
ramified points of $C\to\mathbf P^1$ split, and the exact relative
class number is $194481=441^2$, an odd integer.
\end{claim}

\begin{proof}
The reduced Artin--Schreier poles in (B.86.1) have orders one at
$x=0$ and three at infinity. The local different calculation in
B.85 gives genus two. Write $P$ for the point above $x=0$.
Then $\operatorname{div}(dx)=2P$, so the locally exact-differential
line bundle is $B\simeq\mathcal O_C(P)$. The regular basis
$dx,dx/x$ is sent by Cartier to $0,dx/x$, respectively. Its
stable rank, and hence the two-rank, is one.

In the $u,v$ model in (B.86.3), $P$ is the unique point at infinity,
and
\[
       v_P(u)=-2,\qquad v_P(v)=-5,\qquad
 L(3P)=\langle1,u\rangle,\quad
 L(5P)=\langle1,u,u^2,v\rangle.                              \tag{B.86.5}
\]
The displayed functions are regular away from $P$, have distinct
pole orders, and have the full Riemann--Roch dimensions two and
four. If a nontrivial line bundle $\mathcal O(R-P)$ had order
three, a function with divisor $3R-3P$ would have pole order
three, contrary to (B.86.5). Thus the theta divisor has no
nontrivial three-torsion. The orbit argument (B.85.6)--(B.85.7)
proves the cubic assertion.

For order five, a function with divisor $5R-5P$ must, after
scaling, be
$f=v+c_2u^2+c_1u+c_0$; a polynomial in $u$ alone has even
pole order at most four. Its norm to $k(u)$ is
\[
 b u^5+(c_2^2+c_2)u^4+(a+c_1)u^3
                 +(c_1^2+c_0)u^2+u+c_0^2.                  \tag{B.86.6}
\]
If its zeros are concentrated at a point $R$ with $u$-coordinate
$r$, this norm must equal $b(u+r)^5$. Coefficient comparison
gives, successively,
\[
 c_1=a,\quad c_0=a^2,\quad br^4=1,\quad a^4=br^5=r,\quad
                         c_2^2+c_2=br.                      \tag{B.86.7}
\]
They imply (B.86.2). Conversely (B.86.2), $r=a^4$ and
$c_2=d$ give the identity
\[
 \bigl(v+d u^2+a u+a^2\bigr)
 \bigl(v+u^2+d u^2+a u+a^2\bigr)=b(u+r)^5.
\]
Here $r\ne0$, so the hyperelliptic map is unramified at $u=r$,
and the two factors have different values there. One has a zero
of order five at the point
$R=(r,d r^2+a r+a^2)$ and no other zero, while its only pole
has order five at $P$. Thus $\operatorname{div}(f)=5(R-P)$.
The class $R-P$ is nonzero, since $R\ne P$ and a principal
degree-one difference would make the curve rational. The Kummer
cover (B.86.3) is therefore connected and \'etale of degree five.

There are precisely two choices of $d$, differing by one. Their
points $R$ are exchanged by the hyperelliptic involution, and
their classes are negatives because $\omega_C\simeq\mathcal
O_C(2P)$. The preceding necessity proof lists all nontrivial
five-torsion points on the theta divisor. Consequently they
belong to exactly one cyclic-five subgroup. Multiplication by
two is a four-cycle on the nonzero residues modulo five, so
(B.85.7) gives primitive two-rank zero for this cover and ordinary
primitive part for all others. Equations (B.85.11)--(B.85.12) now give
an actual genus-two reflection quotient $H$, with
$P(D/C)\sim J_H^2$ and $J_H$ supersingular. This is a B5
correspondence with cyclic Galois closure of degree five over $C$.

For $a=b=1$, the involution of (B.86.3) is explicitly
$v\mapsto v+u^2$, $w\mapsto(u+1)/w$.
The identity
$A^5+B^5=(A+B)^5+AB(A+B)^3+(AB)^2(A+B)$ in characteristic
two gives
\[
                   (U+1)u^2+U^3u+U^5+U^3+U=0.
\]
Substitution of (B.86.4) yields the claimed equation of $H$.
The map is nonconstant for the same function-field reason as
in B.85, and it extends to the proper smooth models. Here
$U=1/T$ and $u=(YU^3+U)/(U+1)$, so $k(u,U)=k(T,Y)$;
adjoining $w$, a root of $w^2+Uw+(u+1)$, is exactly the
nontrivial reflection quadratic. The right
side of its Artin--Schreier equation has a single pole of order
five, giving genus two. Its regular basis $dT,TdT$ has Cartier
images $0,dT$, so stable Cartier rank is zero. The height-four
Newton argument given after (B.85.12) proves supersingularity.

Over $\mathbf F _{16}$ all fifth roots of unity and $\rho$ are
present. At the point $u=0$ the Kummer right side equals one.
At $P$, use $t=v/u^3$, for which $u=t^{-2}+\text{higher terms}$
and $v=t^{-5}+\text{higher terms}$; the right side times $t^5$
has residue one. Thus both fibres split. At its order-five zero
$(u,v)=(1,\rho)$ the normalized unit $f/(u+1)^5$ also has
residue one. Direct normalized counts are
\[
\begin{array}{c|rr}
 &\mathbf F _{16}&\mathbf F _{256}\\\hline
 C&24&256\\ D&40&320\\ H&25&289.
\end{array}
\]
They give
\[
 L_C=1+7T+24T^2+112T^3+256T^4,\qquad
 L_H=1+8T+48T^2+128T^3+256T^4.
\]
Both reflections and their conjugating deck transformations are
defined over $\mathbf F _{16}$, so (B.85.12) is an isogeny over
that field. Hence $L_D/L_C=L_H^2$ and the relative class number
is $L_H(1)^2=441^2=194481$. Oddness also follows over every
finite extension from primitive two-rank zero.
\end{proof}
\paragraph{Check record.} Also checked in this writing; GPT-6 Astra Pro, 2026-11-18. Full local theta/Frobenius/reflection argument; exact symbolic identities; normalized Kummer counts; independent FLINT trace counts of base and target and integer zeta convolution. Independent verifier does not re-enumerate Kummer fibres. Effort not separately exposed. Root second reader completed a full proof reading, including theta exact sequence, elimination, normalized residues, reflection isogeny and finite zeta identities, on 2026-11-18.


\paragraph{Check records, attribution, and boundary.}
Also checked in this writing: GPT-6 Astra Pro, finite-checks reader
and root second reader, 18 November 2026. The root second reading
checked the entire theta, elimination, local-residue and reflection
arguments and the finite zeta identities. The exact sequence used from Raynaud (2000),
\S1, p.~74, was opened and checked, as was his Theorem~2 on
p.~76. The latter asserts the existence of a nonordinary solvable
prime-to-characteristic \'etale cover; it does not assert
primitive rank zero or a supersingular factor, and is not used
to obtain either here. The polynomial elimination, norm and
explicit-map identities above are exact symbolic identities over
characteristic two. They and the finite-field normalized point
counts are retained in
\texttt{code/b5\_genus2\_positive.py}, its JSON certificate and
its output. An independent finite-field check and zeta reconstruction
are recorded separately. Nakajima's 1983 \S6 was identified as
prior work on the ordinary cubic example; the available primary
host returned a security challenge, so that section is not
represented as read and no novelty claim is made for (B.85.2).
The proofs here are self-contained apart from the explicitly
stated standard curve and abelian-variety facts. These exact
degree-three and degree-five classifications leave the other
odd degrees, the other genus-two parameters and the universal
B4 and B5 quantifiers unresolved.

\subsection{A finite-degree avoidance boundary}

\begin{claim}{B.87}{Proved}{No uniform finite list of odd cyclic degrees suffices for ordinary genus two}
For every finite set $\mathcal N$ of odd integers greater than one,
there is an ordinary smooth proper genus-two curve
$C/\overline{\mathbf F}_2$ such that every connected cyclic
\'etale cover of $C$ of degree in $\mathcal N$ is ordinary.
More precisely, if $M$ is any odd common multiple of $\mathcal N$,
one can choose $C$ so that its maximal abelian \'etale cover of
exponent dividing $M$, of degree $M^4$, is ordinary. All its
intermediate covers are then ordinary.

Consequently no bound depending only on $g=2,p=2$ for the degree
of a B4 witness can hold for all ordinary curves, if every such
curve has a witness at all. Degree one does not evade this
conclusion, since $J_C$ has two-rank two. This is a finite-degree
avoidance result, not a curve avoiding every odd degree, a
counterexample to B4, or a resolution of any universal B4 or B5
gate.
\end{claim}

\begin{proof}
Let $k=\overline{\mathbf F}_2$. Choose two ordinary elliptic
curves $E_1,E_2/k$ with origins $O_1,O_2$, and form the nodal
compact-type curve
\[
                         C_0=E_1\cup_{O_1=O_2}E_2.
\]
For example both elliptic curves may be
$y^2+xy=x^3+1$. Its projective equation is smooth, as the
partial derivatives show; the point $(0,1)$ is a nonidentity
point fixed by inversion, hence a point of order two. Thus the
elliptic curve is ordinary. Each genus-one component of $C_0$
has a node, so $C_0$ is stable of arithmetic genus two.

There is a regular one-parameter smoothing
$\mathcal C\to\operatorname{Spec}R$, $R=k[[t]]$, with special
fibre $C_0$ and smooth generic fibre. We specify the deformation
fact being used. For a nodal curve, the obstruction group
$\operatorname{Ext}^2(\Omega_{C_0},\mathcal O_{C_0})$ vanishes:
in the local-to-global sequence its possible terms are a second
cohomology group on a curve, a first cohomology group of a sheaf
supported at nodes, and local second Ext of a hypersurface, all
zero. The local first Ext at the one node is one-dimensional,
and the map from global first Ext to it is surjective for the
same reason. We may therefore deform that node by $uv=t$.
The relatively ample dualizing bundle algebraizes the formal
deformation. These are precisely Deligne--Mumford's Lemma~1.3,
Proposition~1.5 and the algebraization paragraph immediately
following it, including characteristic two.
\citep[pp.~79--82]{DeligneMumford1969Finite}

Here is one finite cover which simultaneously controls all the
required cyclic covers. The morphism
\[
 i:C_0\longrightarrow E_1\times E_2
\]
sends $E_1$ to $E_1\times\{O_2\}$ and $E_2$ to
$\{O_1\}\times E_2$; the two definitions agree at the node.
Pull back multiplication by $M$ on $E_1\times E_2$. Its pullback
$\mathcal Y_0\to C_0$ is a finite \'etale torsor for
\[
 G=E_1[M]\times E_2[M]\simeq(\mathbf Z/M)^4.
\]
There are $M^2$ copies of $E_1$ in its normalization, indexed
by $E_2[M]$, and $M^2$ copies of $E_2$, indexed by $E_1[M]$.
For every pair $(a,b)\in E_1[M]\times E_2[M]$, the point $a$
on the copy of $E_1$ labelled $b$ is joined to the point $b$
on the copy of $E_2$ labelled $a$. Its dual graph is therefore
the complete bipartite graph on $M^2$ vertices in each part.
In particular $\mathcal Y_0$ is connected, with $M^4$ nodes.
Its arithmetic genus is
\[
       2M^2+(M^4-2M^2+1)=M^4+1.                            \tag{B.87.1}
\]

The torsor extends uniquely to a finite \'etale $G$-torsor
$\mathcal Y\to\mathcal C$. The exact extension theorem is
the equivalence between finite \'etale covers of a proper
scheme over a henselian local ring and of its special fibre,
in Stacks Project Lemma~58.9.1, Tag~0A48. Its hypotheses hold
because $R$ is complete local and $\mathcal C$ is proper.
In the present case its proof is especially direct: finite
\'etale algebras and their homomorphisms lift uniquely through
the nilpotent thickenings $\mathcal C\bmod t^j$; proper
algebraization gives the finite algebra on $\mathcal C$, and
its \'etaleness follows from the special fibre and properness.
The deck action and torsor identities lift as morphisms.
\citep[Lemma~58.9.1, Tag~0A48]{Stacks0A48Finite}
The generic fibre of $\mathcal Y$ is geometrically connected,
since the special fibre is geometrically connected and proper
flat fibres have upper semicontinuous $h^0(\mathcal O)$.
Equivalently its Stein algebra is $R$, by cohomology and base
change; after any finite extension of $R$, the same argument
applies.

This construction also checks the line-bundle extension
explicitly. Since $M$ is odd and $\mu_M\subset k$, the character
idempotents in $R[G]$ decompose the finite locally free algebra
of the torsor into invertible character eigensheaves
$\mathcal L_\chi$ on $\mathcal C$. Their multiplication gives
$\mathcal L_\chi^{\,\operatorname{ord}(\chi)}\simeq\mathcal O$.
Over $C_0$, they are exactly all $M^4$ bundles in
\[
 \operatorname{Pic}^0(C_0)[M]
       =\bigl(E_1\times E_2\bigr)[M].
\]
Over the geometric generic fibre they remain distinct, because
the connected $G$-torsor has full monodromy $G$. A smooth
genus-two Jacobian has exactly $M^4$ geometric $M$-torsion
points, so these are again all of them. Thus no cyclic cover
of degree dividing $M$ is omitted: every such geometric cover
is a quotient of this $G$-cover. This verifies the full
line-bundle and character quantifier, not merely the extension
of one chosen cover.

We next verify ordinarity without treating graph cohomology
as the geometric $2$-torsion of a torus. For a connected nodal
curve $Y$ with normalization components $Y_v$ and dual graph
$\Gamma$, the normalization exact sequence for $\mathcal O_Y$
gives a Frobenius-compatible exact sequence
\[
 0\longrightarrow H^1(\Gamma,k)
 \longrightarrow H^1(Y,\mathcal O_Y)
 \longrightarrow\bigoplus_v H^1(Y_v,\mathcal O_{Y_v})
 \longrightarrow0.                                        \tag{B.87.2}
\]
On the first term, Frobenius is coefficientwise squaring in
a basis defined over $\mathbf F_2$, so is bijective. Every
normalization component of $\mathcal Y_0$ is an ordinary
elliptic curve, so Frobenius is bijective on the last term too.
It is consequently bijective on the middle term. Its dimension
is $M^4+1$ by (B.87.1).

For the proper flat family $\mathcal Y/R$,
$R^1f_*\mathcal O_{\mathcal Y}$ is locally free of rank $M^4+1$
and commutes with base change. The linearization of Frobenius
is a morphism from its pullback by Frobenius on the base to
itself. Its determinant is a unit at the special fibre by
(B.87.2), and hence is a unit over $R$. Frobenius on the
geometric generic fibre is bijective. That smooth curve is
therefore ordinary, by the Artin--Schreier stable-rank
criterion. The base curve and all intermediate covers are
ordinary as well: pullback on $H^1(\mathcal O)$ is a
Frobenius-compatible split injection, with left inverse the
trace divided by the odd covering degree. An invariant
subspace of a finite-dimensional bijective semilinear map
over a perfect field is again bijective.

Finally the smooth generic curve, its finite torsor, all deck
transformations and the preceding identities involve finitely
many coefficients in an algebraic closure of $k((t))$. They
spread out over an integral finite-type $\mathbf F_2$-algebra
$A$. After restricting to a nonempty open subset of
$\operatorname{Spec}A$, the curves are smooth proper and
geometrically connected, the cover is a $G$-torsor, and its
Frobenius determinant is nowhere zero. These are open
conditions or finitely many equalities already satisfied at
the generic point. Every nonempty finite-type scheme over a
finite field has a closed point with finite residue field.
Such a point gives the required curve over
$\overline{\mathbf F}_2$ and an ordinary maximal exponent-$M$
abelian cover. Its $M^4$ character bundles still exhaust
geometric $M$-torsion, so the conclusion holds for every
cyclic cover of each degree in $\mathcal N$.

For any one of these odd cyclic covers, the primitive factor
is an abelian subquotient up to isogeny of an ordinary
Jacobian and is ordinary of positive dimension. Its two-rank
cannot be zero. Taking $\mathcal N$ to contain all odd
integers between three and any proposed bound proves the
stated limitation. The curve obtained depends on this finite
set. No intersection of all the open conditions as the set
grows has been asserted to contain a closed point.
\end{proof}
\paragraph{Check record.} Also checked in this writing; GPT-6 Astra Pro; 2026-11-18. Full graph, character, deformation, lift, Frobenius openness and specialization proof. No computation needed; effort not separately exposed. Root completed a full second proof reading on 2026-11-18, including smoothing, full torsor, all character bundles, graph cohomology, Frobenius openness and specialization.


\paragraph{Check record.}
Also checked in this writing: GPT-6 Astra Pro, finite-checks
reader and root second reader, 18 November 2026. The complete
smoothing, torsor, character, graph cohomology, Frobenius and
specialization arguments passed the second reading. The full local deformation and
finite-torsor lift hypotheses were compared with the opened
primary sources cited above. The bipartite special cover,
all character line bundles, the normalization cohomology
sequence, Frobenius openness, and finite-field specialization
are included in the proof. There is no numerical search and
no assertion of an all-prime counterexample.

\subsection{A finite pro-three separator for the specified S2 pair}

\begin{claim}{B.88}{Proved}{The S2 pair is separated by a degree-\(3^7\) pointed cover}
Let $C/\mathbf F_2$ and $P,Q$ be the curve and points of S2.
Over $L=\mathbf F_{16}$, there is a specified connected finite
\'etale cover $E\to C_L$ of degree $3^7=2187$ with an
$L$-rational point over $P$ and no $L$-rational point over $Q$.
Its geometric Galois closure is a $3$-group; indeed $E_{\bar L}
\to C_{\bar L}$ is already Galois of degree $3^7$.
Consequently the sections of $P$ and $Q$ are distinct in the
geometrically pro-$3$ arithmetic quotient, over both $L$ and
$\mathbf F_2$, and in every geometrically pro-$\Sigma$ quotient
with $3\in\Sigma$. This does not assert separation of all points
for the singleton prime set $\{3\}$.

For the cover's complete definition, take the cubic cover
$\pi:D\to C_L$ of B.85, choose any $P'\in D(L)$ above $P$,
and put
\[
 J=J_D,\qquad A=J(L),\qquad H=A[2^\infty].
\]
Factor the Lang isogeny as
\[
 F_{16}-1:J\longrightarrow J/H
       \mathop{\longrightarrow}^{\lambda}J.
\]
Then
\[
 E=D\times_{J,\,[R\mapsto R-P']} (J/H)                       \tag{B.88.1}
\]
with point $(P',0+H)$ is the required cover. Thus its definition
uses an explicit smooth curve, a specified finite group scheme
of rational translations and a specified isogeny; equations for
a genus-$2188$ plane model are not required.
\end{claim}

\begin{proof}
B.85 gives
\[
 \#J_C(L)=576=2^6 3^2,\qquad
 \#J_D(L)=46656=2^6 3^6,
\]
and all three points over each of $P,Q$ are rational. We first
prove that
\[
                [Q'-P']\notin J_D(L)[2^\infty]
                  \quad\text{for every }Q'\in\pi^{-1}(Q).     \tag{B.88.2}
\]
Pullback and norm satisfy $\pi_*\pi^*=[3]$. Pullback is
therefore injective on the rational two-primary subgroup of
$J_C$, and the equal two-primary orders show that it maps
onto the rational two-primary subgroup of $J_D$. Every class
in that subgroup is fixed by the cubic deck group.

If $T=[Q'-P']$ were two-primary, a generator $\delta$ of the
deck group would fix it. The effective divisors
$\delta Q'+P'$ and $Q'+\delta P'$ would then be linearly
equivalent, distinct and disjoint. The points over $P$ cannot
equal points over $Q$, and the nonidentity deck transformation
fixes no point. The resulting nonconstant degree-two function
makes $D$ hyperelliptic. This degree-two map is separable:
a purely inseparable degree-two extension of $\bar L(x)$, over
the perfect field $\bar L$, is $\bar L(x^{1/2})$ and is rational,
whereas $g_D=4$. Its unique hyperelliptic involution is
central, so $\delta$ descends to an automorphism of order three
of the projective line. Such an automorphism has a geometric
fixed point. Its permutation of the fibre above that point,
which has at most two points, is trivial, giving a fixed point
of $\delta$ on $D$. This contradicts the \'etaleness of
$D\to C$. This proves (B.88.2) without invoking A1.

The Lang morphism $F_{16}-1$ is a finite \'etale isogeny whose
kernel is the constant group $A=J(L)$: its differential is
$-1$, and its geometric kernel is exactly the rational points.
Quotienting its source by $H$ gives a connected \'etale isogeny
$\lambda$ with constant deck group
\[
                         A/H,\qquad |A/H|=3^6.
\]
After base change to $\bar L$, the Abel map on the smooth
proper curve $D$ induces the isomorphism on abelianized
fundamental groups with its Jacobian. Therefore (B.88.1) is
geometrically connected. This is exactly Milne, Proposition~9.1,
whose separably closed field hypothesis is met after this base
change. In the prime-to-characteristic part used here, it also
follows from Kummer theory and the Jacobian principal
polarization: pullback identifies the prime-power torsion in
$\operatorname{Pic}^0(J)$ with that in $\operatorname{Pic}^0(D)$,
so no nontrivial character of a connected isogeny cover becomes
trivial on $D$. \citep[Proposition~9.1 and Lemmas~9.2--9.3]{MilneJac}

For $T\in J(L)$, the Frobenius translation on the fibre of
$\lambda$ above $T$ is $T+H\in A/H$. Indeed, choose $y$ with
$F_{16}y-y=T$; then
$F_{16}(y+H)-(y+H)=T+H$. Thus this fibre has an $L$-rational
point if and only if $T\in H$. In (B.88.1) the point above $P'$
is rational by construction, while (B.88.2) excludes rational
points above every $Q'$. Any rational point of $E$ above $Q$
would map to one of these rational $Q'$, so none exists.
The degree is $3\cdot3^6=3^7$, and Riemann--Hurwitz gives
$g_E=1+3^7=2188$.

For completeness we verify the geometric group assertion,
which is essential to the pro-$3$ conclusion. Already the
normal-core argument suffices: the geometric subgroup for
$D$ has index three, the subgroup for $E$ is normal in it
with $3$-group quotient, and the intersection of its three
conjugates has $3$-group quotient in that subgroup. Its full
geometric normal closure is consequently a $3$-group.

In fact the cover is geometrically Galois. The automorphism
$\delta$ of $D$ induces an automorphism $\delta_J$ of $J$
commuting with $F_{16}$ and preserving $H$, hence an
automorphism $\delta_{J/H}$ with
$\lambda\delta_{J/H}=\delta_J\lambda$. Choose
$z_0\in(J/H)(\bar L)$ with
$\lambda(z_0)=[\delta P'-P']$. Then
\[
 (R,z)\longmapsto
           \bigl(\delta R,\delta_{J/H}z+z_0\bigr)
\]
is a lift of $\delta$ to $E_{\bar L}$. Together with the
$3^6$ translations by $\ker\lambda$, it generates a group
surjecting onto $\langle\delta\rangle$, with that translation
group as kernel. Its order is $3^7$, equal to the degree.
The lift need not be defined over $L$, which is why we assert
geometric Galoisness and do not assert a constant deck group
for $E\to C_L$.

Finally an $L$-rational point of a finite \'etale cover above
a rational point $R$ is a fixed point of $R$'s section on
the finite geometric fibre. Conjugate sections have fixed
points in precisely the same finite covering sets. The cover
(B.88.1) factors through the geometrically pro-$3$ quotient,
has such a fixed point for $P$, and none for $Q$. Their
sections cannot be conjugate there. Conjugacy over
$\mathbf F_2$ would remain conjugacy on restriction to
$G_L$, so it is excluded as well. Projecting any quotient
with $3\in\Sigma$ to the pro-$3$ quotient gives the last
assertion.
\end{proof}
\paragraph{Check record.} Also checked in this writing; GPT-6 Astra Pro, 2026-11-18. B85 finite inputs have producer, independent verifier and root second reading. Complete norm, hyperelliptic obstruction, Lang, connectedness, geometric Galois lift and finite covering-set proof here. Effort not separately exposed. Root completed the full second proof reading, including geometric Galoisness, on 2026-11-18.


\paragraph{Check record.}
Also checked in this writing: GPT-6 Astra Pro, finite-checks
reader and root second reader, 18 November 2026. The root
second reading passed the full proof, including the stronger
geometric Galoisness assertion. The exact finite input is B.85's
producer and independent verifier, already checked twice.
The prime decompositions, Lang fibre criterion, norm
argument, connectedness, geometric Galois lift and
point-section fixed-point implication are proved above.
Milne's Proposition~9.1 was opened at its statement and
prime-to-characteristic proof, with its field hypothesis
recorded. No computation of all points of $E$ or of its
Jacobian is claimed or needed.

\subsection{The large-characteristic condition cannot bound the witness prime}

\begin{claim}{B.89}{Proved}{Arbitrarily large finite sets of primes can fail permanently in B1}
Fix a prime $p$, an integer $g\ge2$, and a finite set
$\mathcal S$ of odd primes different from $p$. There exist a
finite field $k$ of characteristic $p$, an ordinary smooth
proper geometrically connected genus-$g$ curve $C/k$, and
$P\in C(k)$ such that, for every finite extension $L/k$, every
$\ell\in\mathcal S$, and every nontrivial character
$\chi:J_C(L)\to\mathbf Z/\ell$, one has
\[
 p^{(\ell-1)(g-1)}
 \ \bigm|\ 
       \prod_{i=1}^{\ell-1} L_P(\chi^i,1).                    \tag{B.89.1}
\]
Thus none of those primes is a B1 witness over any constant
extension. In particular no successful witness prime can be
bounded in terms of $(p,g)$ alone, even within
$p>(2g-2)^2$, if unrestricted B1 is true. This statement
does not refute unrestricted B1, and does not decide its
large-characteristic cases: the curve and field in the
construction depend on the prescribed finite set of primes.
\end{claim}

\begin{proof}
Put $M=\prod_{\ell\in\mathcal S}\ell$; if the set is empty the
assertion is vacuous. We first construct an ordinary
geometrically connected maximal exponent-$M$ abelian cover
$Y\to C$ in characteristic $p$, of group
$G=(\mathbf Z/M)^{2g}$, with $C$ of genus $g$. The construction
of B.87 extends to every $p,g$; we check the changes and the
full character quantifier explicitly.

Ordinary elliptic curves exist over $\bar{\mathbf F}_p$.
For $p=2$ the explicit curve in B.87 supplies one. For odd
$p$, put $m=(p-1)/2$ and consider
$y^2=x(x-1)(x-\lambda)$ with $\lambda\ne0,1$. In the Cartier
formula for $dx/y$, its scalar is the $p$th root of the
coefficient of $x^{p-1}$ in
$[x(x-1)(x-\lambda)]^m$. Up to a nonzero sign this is
\[
                         \sum_{i=0}^m\binom mi^2\lambda^i.
\]
Its constant term is one, so it is not the zero polynomial.
Choose $\lambda$ outside its roots and $\{0,1\}$ in the
infinite algebraic closure. Cartier is then bijective on
the one-dimensional regular differential space, proving
ordinarity.

Choose $g$ ordinary elliptic curves $E_i$. When $g=2$, use
their one-node union as in B.87. When $g\ge3$, use a central
projective line with $g$ distinct attachment points, and join
the origin of $E_i$ to the $i$th point. This is a stable
compact-type nodal curve $C_0$ of arithmetic genus $g$.
There is a morphism
$C_0\to A:=\prod_i E_i$ which is constant zero on the central
component and is the $i$th coordinate-axis inclusion on $E_i$.
Pulling back $[M]:A\to A$ gives a finite \'etale $G$-torsor
$Y_0\to C_0$.

For $g\ge3$, the central components of the normalization of
$Y_0$ are projective lines indexed by $G$. The elliptic
components of type $i$ are indexed by the cosets of
$A_i:=E_i[M]$ in $G$, and are copies of $E_i$. A central
vertex $a$ is joined to the elliptic vertex $a+A_i$ for each
$i$. This graph is connected: moving through a component of
type $i$ changes $a$ by any element of $A_i$, and the $A_i$
generate $G$. It has
\[
 |V|=M^{2g}+gM^{2g-2},\qquad |E|=gM^{2g}.
\]
Adding its first Betti number to the sum of component genera
gives
\[
                    p_a(Y_0)=1+(g-1)M^{2g}.                 \tag{B.89.2}
\]
For $g=2$, the bipartite graph of B.87 gives the same formula.
Every normalization component is ordinary or rational.

The nodal deformation calculation in B.87 works in every
characteristic: its obstruction group vanishes and its local
smoothing parameters can all be set to the same parameter
$t$. The ample dualizing bundle algebraizes the resulting
regular smoothing over $\bar{\mathbf F}_p[[t]]$. The
proper henselian equivalence for finite \'etale covers lifts
$Y_0$ and the whole $G$-action uniquely. Both inputs have
exactly the characteristic-independent hypotheses checked
in B.87. \citep[Lemma~1.3, Proposition~1.5 and following
paragraph]{DeligneMumford1969Finite}
\citep[Tag~0A48]{Stacks0A48Finite}

The normalization cohomology sequence (B.87.2) now has
coefficientwise $p$th-power Frobenius on its graph term,
and bijective Frobenius on every elliptic component term.
It follows that Frobenius on $H^1(Y_0,\mathcal O_{Y_0})$
is bijective. Its determinant on relative $R^1f_*\mathcal O$
is a unit, so the smooth geometric generic fibre $Y$ is
ordinary. The degree of $Y\to C$ is prime to $p$, hence
trace divided by that degree splits the pullback of
$H^1(C,\mathcal O_C)$; therefore $C$ is ordinary too.
Spreading out the curves, the $G$-torsor and the nonzero
Frobenius determinant, and choosing a finite-field closed
point as in B.87, gives such an ordinary cover over a
finite field.

No geometric cyclic $\ell$-cover for $\ell\in\mathcal S$
is missing. Since $Y_{\bar k}\to C_{\bar k}$ is a connected
$G$-torsor, its character eigensheaves give $M^{2g}$ distinct
line bundles of $M$-torsion on $C_{\bar k}$. A genus-$g$
Jacobian has exactly $M^{2g}$ such geometric points, because
$M$ is prime to $p$. These are all of them. Equivalently
the corresponding quotient of the geometric abelian
fundamental group is the full exponent-$M$ quotient.
Every connected cyclic cover of degree $\ell\in\mathcal S$
is thus $D_\chi=Y/\ker\chi$ for some surjection
$\chi:G\to\mathbf Z/\ell$.

We now choose the finite-field descent and the pointing in
the order that avoids the twist error. Enlarge the finite
field until the deck group of $Y\to C$ is constant, a
chosen geometric point $y\in Y$ is rational, and every
geometric point of $J_Y[p]$ is rational. This is a finite
extension: the group of geometric $p$-torsion points is
finite. Denote the resulting field by $k$ and put $P=\pi(y)$.
All the quotients $D_\chi\to C$ are now constant and
$P$-split. We do not first choose a good twist and later
extend its field; the purpose of this enlargement is
instead to make every listed quotient permanently bad.

For each such quotient, the map $Y\to D_\chi$ has degree
prime to $p$. Its pullback on geometric $p$-torsion is
injective, since norm composed with pullback is multiplication
by that degree. As all geometric $p$-torsion on $J_Y$
is rational, all of it on $J_{D_\chi}$ is rational too.
The same assertion holds after every extension $L/k$.
Moreover $D_\chi$ is ordinary, because its coherent
cohomology is again a Frobenius-compatible direct summand
of that of $Y$.

Write $\pi_\chi:D_\chi\to C$ and
$B_\chi=(\ker(\pi_{\chi*}:J_{D_\chi}\to J_C))^0$.
Pullback and norm satisfy
$\pi_{\chi*}\pi_\chi^*=[\ell]$. On the $p$-divisible
groups the operator
$\ell^{-1}\pi_\chi^*\pi_{\chi*}$ is an idempotent, whose
complement is the $p$-divisible group of $B_\chi$.
Thus $B_\chi$ is ordinary of dimension
$(\ell-1)(g-1)$ and all its geometric points of order
$p$ are $L$-rational. In particular
\[
             p^{(\ell-1)(g-1)}\mid\#B_\chi(L).
\]
The pullback--norm decomposition gives an isogeny over $L$
from $J_C\times B_\chi$ to $J_{D_\chi}$. Finite-field
point counts are invariant under isogeny, by their
Frobenius characteristic polynomials, so
$\#B_\chi(L)=\#J_{D_\chi}(L)/\#J_C(L)$.
Factoring each Euler factor of the cyclic cover into its
character factors gives
\[
 \frac{\#J_{D_\chi}(L)}{\#J_C(L)}
              =\prod_{i=1}^{\ell-1}L_P(\chi^i,1).
\]
This proves the required divisibility for these pointed
quotients.

Finally, take any other constant cyclic $\ell$-cover over
$L$ whose fibre at $P$ splits. Its geometric character is
one of those just exhausted. Two extensions of that same
geometric character to the arithmetic fundamental group
differ by a constant character of $G_L$. The splitting
at the prescribed point forces this difference to vanish:
both extensions are zero on the section of $P$. The cover
is therefore exactly the base change of the corresponding
$P$-split $D_\chi$. Hence the divisibility applies to
every nontrivial character occurring in (B.89.1), over
every finite extension.

Taking $\mathcal S$ to contain all admissible odd primes
up to any proposed bound proves the final assertion,
without any restriction on $p$ relative to $g$. The
construction depends on that finite set and supplies
no curve for which all primes fail at once.
\end{proof}
\paragraph{Check record.} Also checked in this writing; GPT-6 Astra Pro, 2026-11-18. Full geometric construction and all finite-field/pointed-cover quantifiers checked; no computation or unproved generic-to-closed-point inference. Effort not separately exposed. Root completed a full second proof reading on 2026-11-18, including the compact-type graph, ordinary deformation, every-character exhaustion, rational p-torsion field choice before pointing, unique split twist over every extension and exact divisibility.


\paragraph{Check record and threshold boundary.}
Also checked in this writing: GPT-6 Astra Pro, finite-checks
reader and root second reader, 18 November 2026. The root
second reader checked the complete proof independently,
including all constant-extension and pointing quantifiers. The genus-$g$ special graph,
complete character count, ordinary smoothing, finite-field
rationalization of the full geometric $p$-torsion, norm
splitting and uniqueness of the pointed twist have each
been checked in the proof. No bounded search or generic
point is substituted for a curve over a finite field.
The inequality $p>(2g-2)^2$ remains an unresolved
unrestricted B1 range; this claim rules out only a
bounded-prime conclusion depending on those two integers.

% AI-produced by GPT-6 Astra (OpenAI), ChatGPT Pro mode, 18 November 2026.
% Working identifiers; the editor assigns the continuing Part B numbers.
\subsection{Faithfulness, good reduction, and the marked groupoid}

We fix the following convention.  A marked groupoid has a nonempty finite
set $I\subset X(K)$ of rational points as objects; its arrows are the
isomorphisms between the corresponding geometric fibre functors on finite
\'etale covers of $X_{\overline K}$.  The natural $G_K$-action fixes its
objects.  It acts on the actual arrows, not on their classes modulo inner
automorphisms.  Changing the geometric fibre functors conjugates this
action by a groupoid isomorphism and does not change its kernel.  The
singleton convention is already sufficient to test the expectation in
cote~142, scan~096, manuscript page~95.

\begin{claim}{B.90}{Known}{The constant-field restriction on an outer kernel}
Let $k$ be a perfect field of characteristic zero and $X/k$ a hyperbolic
curve.  For the full geometric profinite group, the kernel of
$G_k\to\operatorname{Out}(\pi_1(X_{\overline k}))$ is contained in the
kernel of restriction $G_k\to G_{\mathbb Q}$, up to the usual choice of
an embedding of $\overline{\mathbb Q}$ in $\overline k$.  If $k$ is a
number field, the outer action, and hence every nonempty marked-groupoid
action, is faithful.
\end{claim}
\begin{proof}
This is Hoshi--Mochizuki, Theorem~C(ii), equivalently Corollary~6.2(ii),
specialized to configuration level $n=1$ and to the set of all primes.
Corollary~6.3(i) states the resulting injectivity for number fields and
$p$-adic local fields.  The source assumes a hyperbolic curve over a
perfect field and that every selected prime is invertible; these
hypotheses hold here.  Restriction $G_k\to G_{\mathbb Q}$ is injective
when $k/\mathbb Q$ is algebraic, which proves the number-field conclusion.
The kernel of an actual groupoid action lies in its outer-action kernel.
The proof of Corollary~6.2(ii) uses Belyi faithfulness on the
thrice-punctured line after Corollary~6.2(i).  It does not assert that
restriction $G_k\to G_{\mathbb Q}$ is injective when $k$ has positive
transcendence degree.  The attached author's PDF, printed pages~7 and
63--65, and the author's online copy were both opened.
\end{proof}
\paragraph{Check record.} checked only in this writing, 2026-11-18: GPT-6 Astra (OpenAI), ChatGPT Pro mode; internal effort value not exposed. Own detailed proof written and reread; primary source hypotheses matched. All seven N4 claims read; targeted checks passed, as root message records


\begin{claim}{B.91}{Proved}{Good reduction kills the actual marked action}
Let $R$ be a henselian DVR of residue characteristic zero, with fraction
field $F$.  Let $\overline{\mathcal X}/R$ be a smooth proper curve and
$\mathcal D\subset\overline{\mathcal X}$ a relative finite \'etale
divisor.  Put $\mathcal X=\overline{\mathcal X}\setminus\mathcal D$.
Suppose every point of a finite set $I\subset\mathcal X_F(F)$ extends
to an $R$-section of $\mathcal X$.  Then the inertia subgroup $I_R$ of
$G_F$ acts trivially on the full profinite groupoid
$\widehat\Pi_1(\mathcal X_{\overline F},I)$.
\end{claim}
\begin{proof}
Pass to the strict henselization $R^{\rm sh}$.  Its fraction field is
the maximal unramified extension of $F$ and its absolute Galois group
is $I_R$.  SGA~1, Expos\'e~XIII, Corollary~2.9 and the specialization
construction~2.10 apply to this smooth proper family with relative
normal-crossings boundary.  All primes are invertible in $R$, so the
prime-to-residue-characteristic restriction disappears.  They give the
equivalence of finite-cover categories between the geometric generic
fibre and the strictly local family (and the special fibre).  Equivalently,
every finite geometric cover is the restriction of a finite \'etale
cover of $\mathcal X_{R^{\rm sh}}$.

For an extending section $i$, the pullback of any such cover to
$\operatorname{Spec}R^{\rm sh}$ is a constant finite \'etale scheme.
Consequently its geometric generic fibre has trivial $I_R$-action.
The category of these extending covers is cofinal in all finite
geometric covers, and the same assertion holds simultaneously for each
of the finitely many fibre functors.  Naturality shows that inertia acts
identically on every isomorphism between those fibre functors.  Taking
the inverse limit proves triviality on every arrow of the profinite
groupoid.  In the singleton case this can also be seen from the based
specialization isomorphism: the rational section supplies its basepoint
transport, and $I_R$ acts trivially on the special-fibre group.
Thus an outer triviality statement has not been substituted for the
required actual-action statement.
\end{proof}
\paragraph{Check record.} checked only in this writing, 2026-11-18: GPT-6 Astra (OpenAI), ChatGPT Pro mode; internal effort value not exposed. Own detailed proof written and reread; primary source hypotheses matched. All seven N4 claims read; targeted checks passed, as root message records


\begin{claim}{B.92}{Proved}{Every finite marking over positive transcendence degree has a kernel}
Let $K$ be finitely generated over $\mathbb Q$, with
$\operatorname{trdeg}_{\mathbb Q}K>0$, and let $X/K$ be a smooth
hyperbolic curve.  For every nonempty finite set $I\subset X(K)$, the
actual action of $G_K$ on $\widehat\Pi_1(X_{\overline K},I)$ is not
faithful.  In fact its kernel contains a subgroup isomorphic to
$\widehat{\mathbb Z}(1)$.
\end{claim}
\begin{proof}
Choose a smooth integral model $S$ of $K$ over a number field.  Spreading
out the proper smooth completion, its finite reduced boundary, and the
finitely many marked points, and then shrinking $S$, gives a smooth
proper family with relative finite \'etale boundary and all marked
sections inside its complement.  There is a codimension-one point
$\eta$ of this nonempty positive-dimensional smooth scheme; its local
ring $R=\mathcal O_{S,\eta}$ is a DVR with residue characteristic zero
and fraction field $K$.  Its inertia subgroup embeds in $G_K$ after a
prolongation is fixed.  Tame Kummer theory gives
$I_R\simeq\widehat{\mathbb Z}(1)$, since the strict henselization has
value group $\mathbb Z$, separably closed residue field, and residue
characteristic zero.  In particular the subgroup is nontrivial: if
$\varpi$ is a uniformizer, adjoining $\sqrt[n]{\varpi}$ has ramification
index $n$ for every $n>1$.  Claim~B.91 puts this whole subgroup in the
marked-action kernel.  Every comparison uses restriction from the same
chosen algebraic closure of $K$; changing the prolongation merely
conjugates the displayed subgroup.
\end{proof}
\paragraph{Check record.} checked only in this writing, 2026-11-18: GPT-6 Astra (OpenAI), ChatGPT Pro mode; internal effort value not exposed. Own detailed proof written and reread; primary source hypotheses matched. All seven N4 claims read; targeted checks passed, as root message records


\begin{claim}{B.93}{False as stated}{The moduli condition does not imply uniqueness}
The uniqueness expectation of cote~142, scan~096 is false for finite
markings, even when the field of definition is finite over the field of
moduli.  A counterexample is
\[
 K=\mathbb Q(t),\qquad
 X=\mathbb P^1_K\setminus\{0,1,t,\infty\},\qquad I=\{2\}.
\]
The inertia at $t=3$ is a nontrivial subgroup of the kernel of the actual
marked action.  The assertion becomes true for number fields, and for
finite markings positive transcendence degree always prevents it, as
specified in Claims~B.90 and B.92.
\end{claim}
\begin{proof}
The curve has type $(0,4)$ and is hyperbolic.  Set
\[
 j(t)=256\frac{(1-t+t^2)^3}{t^2(1-t)^2}.
\]
The six cross-ratio transforms of the unordered four-point set are
$t,1-t,1/t,1/(1-t),t/(t-1),(t-1)/t$.  They form its generic orbit,
and $j$ is their common invariant.  The rational map $t\mapsto j(t)$
has degree six (its coprime numerator and denominator have degrees six
and four).  Thus the fixed field of this order-six group is exactly
$\mathbb Q(j)$, the field of moduli of the unlabelled punctured curve,
and $[K:\mathbb Q(j)]=6$.  If the punctures are ordered, the moduli
field is already $K$; adding the rational marking does not invalidate
the finiteness assertion.

Over $R=\mathbb Q[t]_{(t-3)}$, the sections $0,1,t,\infty$ are disjoint
and the section $2$ avoids them.  Claim~B.91 therefore kills the
inertia at $t-3$ in the actual based action.  This inertia is nontrivial
because $K(\sqrt{t-3})/K$ is ramified of degree two at that valuation;
indeed all Kummer degrees occur after the unramified constant extension.
This proves a counterexample in the full profinite group and does not
rely on a finite nilpotent approximation.  The program
\texttt{n4\_moduli\_certificate.py} checks the six rational invariants,
the polynomial degree and coprimality, and the specialization of the
marked divisor.  The ramification and trivial-action assertions are the
preceding proofs, not claims made by that finite calculation.
\end{proof}
\paragraph{Check record.} checked only in this writing, 2026-11-18: GPT-6 Astra (OpenAI), ChatGPT Pro mode; internal effort value not exposed. Own detailed proof written and reread; primary source hypotheses matched. All seven N4 claims read; targeted checks passed, as root message records


\begin{claim}{B.94}{Proved}{The quotient that good-reduction inertia actually supplies}
Let a marked curve as above spread to a smooth relative pair over a
smooth integral characteristic-zero scheme $S$, with function field
$K$.  Let $N_S$ be the closed normal subgroup of $G_K$ generated by
the inertia groups at all codimension-one points of $S$.  Then
\[
 G_K/N_S\simeq\pi_1(S),\qquad
 \ker(\rho_I)/N_S=\ker(\rho_{S,I}),
\]
where $\rho_{S,I}:\pi_1(S)\to\operatorname{Aut}
(\widehat\Pi_1(X_{\overline K},I))$ is the actual family monodromy.
Thus uniqueness modulo $N_S$ is precisely injectivity of this family
monodromy; it is not a consequence of the moduli-field condition alone.
\end{claim}
\begin{proof}
The restriction functor from finite \'etale covers of normal $S$ to
finite separable extensions of $K$ is fully faithful, and its image
consists of the extensions whose normalizations are unramified over
$S$.  Since $S$ is regular, purity of the branch locus identifies this
image with extensions unramified at every codimension-one point.
Equivalently the corresponding finite quotients of $G_K$ are exactly
those killing all of the indicated inertia subgroups.  This proves the
first displayed isomorphism by the universal property of the profinite
quotient.  The smooth pair supplies finite locally constant systems of
geometric covers and of their fibres at the marked sections; hence its
groupoid monodromy is defined on $\pi_1(S)$.  Restriction to the generic
point is $\rho_I$.  Kernels in this factorization give the second
identity.  Purity is used here in its usual regular excellent form:
Stacks Project, Lemma~58.21.4 (Tag~0BMB), applied to the finite
normalization.  This result gives a factorization and the exact
remaining monodromy kernel, not an unproved assertion that the latter
is always zero.
\end{proof}
\paragraph{Check record.} checked only in this writing, 2026-11-18: GPT-6 Astra (OpenAI), ChatGPT Pro mode; internal effort value not exposed. Own detailed proof written and reread; primary source hypotheses matched. All seven N4 claims read; targeted checks passed, as root message records


\begin{claim}{B.95}{Proved}{A faithful corrected monodromy in the counterexample}
For the curve and marking of Claim~B.93, put
$S=\mathbb P^1_{\mathbb Q}\setminus\{0,1,2,\infty\}$.
Its actual family monodromy
\[
 \pi_1(S)\longrightarrow
 \operatorname{Aut}\bigl(\pi_1(X_{\overline{\mathbb Q(t)}},2)\bigr)
\]
is injective.  Consequently the kernel of the $G_{\mathbb Q(t)}$-action
is exactly $N_S$ in Claim~B.94.
\end{claim}
\begin{proof}
Write $C=\mathbb P^1\setminus\{0,1,\infty\}$ and
$Y=(C\times C)\setminus\operatorname{diag}$.  Work first over
$\overline{\mathbb Q}$ with the point $(x,t)=(2,3)$.  In
$E=\pi_1(Y)$ let $F$ be the fibre subgroup of the second projection
and $H$ the fibre subgroup of the first projection.  Both are free
profinite groups of rank three and both inject in $E$.  Here full
profinite fibre exactness can be verified directly.  Over the complex
numbers the projection is a locally trivial punctured-surface fibration
with base having free discrete group of rank two.  Its exact sequence
splits because that quotient is free.  For every finite quotient of
the finitely generated free fibre group, choose a characteristic
finite-index subgroup inside its kernel.  The resulting finite fibre
group $Q$ has finite automorphism group; the semidirect product maps
to $Q\rtimes\operatorname{Aut}(Q)$, proving that every finite fibre
quotient is detected in a finite quotient of the total group.  Right
exactness of completion proves the assertion.  Riemann existence and
invariance under algebraically closed characteristic-zero extension
identify these completions with the groups in use.

The geometric monodromy of the marked family is the conjugation of
$H$ on $F$: the map $S\to Y$, $t\mapsto(2,t)$, is exactly the fibre
inclusion for the first projection and supplies the section for the
second projection after restricting its base to $S$.  The intersection
$F\cap H$ is the kernel of
$H=\pi_1(C\setminus\{2\})\to\pi_1(C)$.  It contains the cusp
inertia group $J$ at $2$ and every $H$-conjugate of $J$.

The cusp facts needed next are Stix (2012), Proposition~16(1)--(3):
distinct geometric cusp inertia groups have trivial intersection,
their geometric normalizers equal themselves, and an abelian subgroup
meeting such inertia nontrivially lies in it.  They hold for this
hyperbolic curve.  A topological generator $c$ of $J$ has
$C_H(c)=J$: its centralizer normalizes the procyclic closure of
$\langle c\rangle$, so lies in $J$, and the reverse inclusion is clear.
There exists $h\in H\setminus J$, because $H$ has rank three and is
not procyclic.  Then $J\cap hJh^{-1}=1$.  Any element of $H$ acting
trivially on $F$ centralizes both $c$ and $hch^{-1}$, and hence lies in
this trivial intersection.  Geometric monodromy is injective.

Finally filling the moving puncture gives an equivariant surjection
$F\twoheadrightarrow\pi_1(C_{\overline{\mathbb Q}},2)$.
The induced arithmetic monodromy on this quotient is the constant
curve's $G_{\mathbb Q}$-action.  Claim~B.90 makes its outer action
faithful.  An element of the full $\pi_1(S)$-kernel therefore has
trivial image in $G_{\mathbb Q}$, so belongs to the geometric subgroup,
where injectivity was just proved.  This establishes the stated
arithmetic injection and then the kernel equality.
\end{proof}
\paragraph{Check record.} checked only in this writing, 2026-11-18: GPT-6 Astra (OpenAI), ChatGPT Pro mode; internal effort value not exposed. Own detailed proof written and reread; primary source hypotheses matched. All seven N4 claims read; targeted checks passed, as root message records


\begin{claim}{B.96}{Proved}{The constant example and the scope of image conditions}
For a constant hyperbolic curve $X_0/\mathbb Q$ with constant nonempty
finite marking $I_0$, its base change to $K=\mathbb Q(t)$ has marked
action factoring through $G_K\to G_{\mathbb Q}$; its kernel is exactly
the kernel of that restriction.  This explains the isotrivial
obstruction, but Claims~B.92--B.93 show that eliminating
isotriviality does not eliminate good-reduction inertia.

The additional relations of Nakamura--Schneps (2000), Theorem~1.1 and
Definition~1.1, define $\mathrm I\Gamma\subset\widehat{GT}$;
Theorem~1.2 makes it a subgroup containing the Galois image.  None of
these statements identifies it with all of $\widehat{GT}$.
Likewise Hoshi--Mochizuki--Tsujimura (2025), Definition~3.3(ii),(iii),(v),
imposes the COF and RGC properties on closed subgroups before denoting
them by $\mathrm{BGT}$.  These are hypotheses on the selected subgroups,
not assertions about every element of the ambient group.
\end{claim}
\begin{proof}
In the constant example, covers and the constant fibre functors descend
from the algebraic closure of the constant field.  Thus the action
factors as asserted.  Faithfulness over $\mathbb Q$ from Claim~B.90
identifies its kernel.  The source distinctions in the second paragraph
follow by reading the stated definitions and theorems in the opened
primary texts.  They cannot be used to prove that the necessary
conditions a), b) of cote~142, scan~095 characterize the whole ambient
automorphism group.  Those conditions concern existence of a Galois
element; the counterexample in Claim~B.93 concerns uniqueness of an
element already known to act.  The two assertions are logically
separate.  The Belyi argument in Claim~B.90 proves faithfulness for
the constant-field Galois quotient and does not remove geometric
inertia over a field of positive transcendence degree.
\end{proof}
\paragraph{Check record.} checked only in this writing, 2026-11-18: GPT-6 Astra (OpenAI), ChatGPT Pro mode; internal effort value not exposed. Own detailed proof written and reread; primary source hypotheses matched. All seven N4 claims read; targeted checks passed, as root message records


\paragraph{Check record.}
Claims B.91--B.96 were checked only in this writing by GPT-6 Astra
(OpenAI, ChatGPT Pro mode), 18 November 2026: direct re-proof of the
specialization and fibre-functor comparisons, independent construction
of the counterexample, and the explicit symbolic certificate specified
above.  The separately selected internal reasoning-effort value is not
exposed.  B.90 is an opened published theorem with its hypotheses
matched here.  No human or different-model-family verification is
claimed.  The field-of-moduli condition and the finite-marking
convention are stated explicitly; no assertion about arbitrary infinite
marking sets is silently inferred.

% AI-produced by GPT-6 Astra (OpenAI), ChatGPT Pro mode, 18 November 2026.
% Working identifiers; editor assigns the continuing Part B numbers.
\subsection{Valuation inertia, centres, and boundary characters}

Here $K$ has characteristic zero, $L/K$ is finitely generated, and $v$
is a valuation of $L$ trivial on $K$.  Fix a prolongation to
$\overline L$; all decomposition and inertia groups are actual closed
subgroups of $G_L$.  Put $T=\widehat{\mathbb Z}(1)$ and
$\Gamma=\Gamma_v$.  We distinguish
\[
 \Gamma_{\rm div}=\bigcap_{n\geq1}n\Gamma
 \quad\hbox{from the largest divisible \emph{convex} subgroup}
 \quad\Delta\subset\Gamma.
\]
The first subgroup is divisible: if $a\in\Gamma_{\rm div}$, the
unique element $a/m$ in $\Gamma\otimes\mathbb Q$ belongs to
$n\Gamma$ for every $n$, since $a\in mn\Gamma$.  Conversely every
divisible subgroup is contained in it.  It need not be convex.

The page marks in the supplied second-pass transcription put scan~119
at manuscript page~59 and scan~121 at page~60; scan~120 continues the
former page and ends with the margin number~119.  Thus the question
in scan~120 belongs to manuscript page~59.  This is a check of the
supplied page-mark transcription, not a claim of a new inspection of an
unprovided image of that scan.

\begin{claim}{B.97}{Proved}{The inertia formula in every rank}
There are canonical Kummer identifications
\[
 I_v=\operatorname{Hom}((\Gamma\otimes\mathbb Q)/\Gamma,\mu_\infty)
     =\operatorname{Hom}(\Gamma,T).
\]
The topology is the inverse-limit topology of finite Kummer quotients.
Furthermore
\[
 1\longrightarrow I_v\longrightarrow D_v\longrightarrow
 G_{\kappa(v)}\longrightarrow1
\]
splits, noncanonically.  Conjugation on inertia is the cyclotomic action.
If $r_\ell=\dim_{\mathbb F_\ell}\Gamma/\ell\Gamma$, then
\[
 I_v\simeq\prod_\ell\mathbb Z_\ell(1)^{r_\ell},\qquad
 r_\ell\leq\dim_{\mathbb Q}(\Gamma\otimes\mathbb Q)
 \leq\operatorname{trdeg}_K L-\operatorname{trdeg}_K\kappa(v).
\]
For $\gamma\in\Gamma$, evaluation
$\chi\mapsto\chi(\gamma)$ is identically zero exactly when
$\gamma\in\Gamma_{\rm div}$.  In particular $I_v=1$ exactly when
$\Gamma$ is divisible.  These statements re-prove FLAGGED\_149 D4
in the characteristic-zero setting at issue here.
\end{claim}
\begin{proof}
Let $H=L_v^h$ be the chosen henselization, and let $F=H^{\rm nr}$
be its maximal unramified extension.  Its residue field is
$\kappa(v)^{\rm sep}$ and its value group is $\Gamma$.  Units of
$F$ have $n$th roots in $F$: choose an $n$th root of the residue and
apply Hensel's lemma, whose derivative is invertible in residue
characteristic zero.  Tame extension theory, in the form of Temkin,
Theorem~3.4.4, and the absence of wild inertia in Corollary~3.5.3,
therefore identifies the absolute Galois group of $F$ with the stated
character group.  Explicitly its pairing is
\[
 (\sigma,v(a)/n\bmod\Gamma)\longmapsto
 \sigma(a^{1/n})/a^{1/n}.
\]
Independence of $a$ with a fixed value follows from the assertion about
units.  The finite Kummer pairings are perfect, since a value not in
$n\Gamma$ gives a nontrivial root extension.  Their union and dual
give the first formula.  Restriction to
$(n^{-1}\Gamma)/\Gamma\simeq\Gamma/n\Gamma$ gives the second.
Thus evaluation at $\gamma$ vanishes exactly when
$\gamma/n\bmod\Gamma=0$ for every $n$, which is the stated
divisibility criterion.  The exact decomposition sequence follows
from lifting finite separable residue extensions to unramified
extensions of the henselian valuation ring.

Here is a complement argument that does not assume a coefficient
field.  Among algebraic extensions $M/H$ in $\overline L$ with
unchanged residue field choose a maximal one.  Unions of chains
preserve this property.  If a value in $\Gamma_M$ is not divisible by
a prime $\ell$, adjoining an $\ell$th root of an element with that
value has ramification index $\ell$, hence unchanged residue field,
contrary to maximality.  Thus $\Gamma_M$ is divisible.  All these
extensions are henselian.  Residue-characteristic-zero henselian
fields are defectless (Temkin, Corollary~3.5.8), so $M^{\rm nr}$,
which has divisible value group and algebraically closed residue
field, is algebraically closed.  Hence restriction identifies
$G_M$ with $G_{\kappa(v)}$, and $G_M\hookrightarrow G_H=D_v$
is a complement.  Applying a decomposition element to the displayed
Kummer pairing fixes the value and changes only the root of unity;
this proves the cyclotomic assertion.

The torsion group $(\Gamma\otimes\mathbb Q)/\Gamma$ is divisible;
its $\ell$-torsion is $\Gamma/\ell\Gamma$.  The classification of
divisible torsion groups gives the displayed product after duality.
An $\mathbb F_\ell$-independent set of values is rationally
independent: in an integer relation divide by the greatest common
power of $\ell$ and reduce modulo $\ell$.  For the last bound,
choose elements with rationally independent values and units with
algebraically independent residues.  In a polynomial relation, terms
with distinct value exponent vectors have distinct values, and the
coefficient of the least-value vector is a nonzero polynomial in the
independent residues.  Cancellation is impossible.  These elements
are therefore algebraically independent, proving the bound.
\end{proof}
\paragraph{Check record.} checked only in this writing, 2026-11-18: GPT-6 Astra (OpenAI), ChatGPT Pro mode; internal effort value not exposed. Own detailed proof written and reread; primary source hypotheses matched. Through N3.7 read; targeted checks passed, as root message records


\begin{claim}{B.98}{Proved}{The inertia normalizer and its necessary coarsening}
Suppose $I_v\ne1$.  Let $w=v/\Delta$ be the coarsening defined by
the largest divisible convex subgroup of $\Gamma$.  Then
\[
 I_w=I_v,\qquad N_{G_L}(I_v)=D_w.
\]
Writing $P_{L/K}=\operatorname{Gal}(\overline L/L\overline K)$ gives
\[
 N_{P_{L/K}}(I_v)=D_w\cap P_{L/K},\qquad
 1\to I_v\to D_w\cap P_{L/K}\to
 G_{\kappa(w)\overline K}\to1.
\]
Moreover $D_w=D_v$ exactly when the residual valuation $v/w$ on
$\kappa(w)$ is henselian.  In particular $N_{G_L}(I_v)=D_v$ for
every rank-one valuation with nontrivial inertia.  If $I_v=1$, its
normalizer is all of $G_L$.
\end{claim}
\begin{proof}
Convex subgroups form a chain, so their divisible members have a
largest union $\Delta$.  For any convex subgroup $A$, both $A$ and
$\Gamma/A$ are torsion free.  The exact sequence of divisible torsion
groups associated with $0\to A\to\Gamma\to\Gamma/A\to0$ and
Kummer duality gives
\[
 1\to I_{v/A}\to I_v\to\operatorname{Hom}(A,T)\to1.
\]
The first map is the actual inclusion for the induced prolongations:
its characters are pullbacks along $\Gamma\to\Gamma/A$.  It is an
equality exactly when $A$ is divisible.  Thus $I_w=I_v$.

On $F=\overline L^{I_v}$ the valuation $v$ is henselian, with
separably closed residue and value group $\Gamma$.  Its coarsening
by $A$ has separably closed residue exactly when $A$ is divisible.
Indeed its residual refinement is henselian, has value group $A$
and separably closed final residue; Claim~B.97 computes the absolute
Galois group of the intermediate residue field as
$\operatorname{Hom}(A,T)$.  Consequently $w|_F$ is the coarsest
such coarsening.

It is also the coarsest among \emph{all} henselian valuations of $F$
with separably closed residue.  The precise comparison theorem used
here is Kuhlmann, Theorem~9.44 and Corollary~9.45: two independent
nontrivial henselian valuations force the field to be separably
closed, and two incomparable henselian valuations have a finest common
coarsening with separably closed residue.  If another such valuation
were incomparable with $w|_F$, their common coarsening would contradict
the preceding maximality of $\Delta$.  If it were a proper coarsening,
the same maximality would give a contradiction.  It must therefore
refine $w|_F$.  Since $I_v\ne1$, the field $F$ is not separably
closed, and this distinguished valuation is nontrivial.

An element normalizing $I_v$ acts on $F$ as an $L$-automorphism, so
preserves this intrinsically characterized valuation.  Its extension
to $\overline L$ is unique by henselianity, hence the element belongs
to $D_w$.  Conversely $D_w$ normalizes its inertia.  Intersection with
$P_{L/K}$ gives the geometric identity.  Constants have value zero;
passing from $L$ to $L\overline K$ preserves the value group and
adjoins $\overline K$ to the residue field.  The residue exact
sequence gives the last displayed sequence.  Finally $D_v$ is the
inverse image in $D_w$ of the decomposition group of $v/w$; that
group is the whole residue Galois group exactly for a henselian
valuation.  For rank one the only convex subgroups are $0,\Gamma$,
and $I_v\ne1$ excludes a divisible $\Gamma$.

This verifies FLAGGED\_149 D5, rather than replacing $\Delta$ by
the generally nonconvex subgroup $\Gamma_{\rm div}$.  It also checks
the two counterexamples in D6.  For the $x$-adic valuation of
$\mathbb Q(x,y)$, the geometric quotient of the normalizer is
$G_{\overline{\mathbb Q}(y)}\ne1$.  For an $x$-adic valuation of
$\mathbb Q(x,y,z)$ refined by the divisible-value valuation of
Claim~B.99 on $\mathbb Q(y,z)$, the value group is
$\mathbb Z\oplus_{\rm lex}\mathbb Q$, so the two inertias agree.
The refinement is not henselian: the two prolongations to
$\mathbb Q(y,z)(\sqrt{1+y})$ distinguish residues $1,-1$ of the
root, while the quadratic polynomial is irreducible.  Thus its
decomposition subgroup, and hence $D_v$, is proper in $D_w$.
\end{proof}
\paragraph{Check record.} checked only in this writing, 2026-11-18: GPT-6 Astra (OpenAI), ChatGPT Pro mode; internal effort value not exposed. Own detailed proof written and reread; primary source hypotheses matched. Through N3.7 read; targeted checks passed, as root message records


\begin{claim}{B.99}{Proved}{Divisible-value counterexamples, including nonzero invisible inertia}
FLAGGED\_149 D3 is valid: $K(x,y)$ has valuations with residue $K$
and value group $\mathbb Q$, or $\mathbb Z[1/2]$, centred at $(0,2)$
on $(\mathbb P^1)^2$.  Consequently the centre converse for arbitrary
valuations is false on $(\mathbb P^1-\{0,1,\infty\})^2$.

More strongly, on
$U=(\mathbb P^1-\{0,1,\infty\})^3$ there is a rank-one valuation
of $K(U)$ with $I_v\simeq T\ne1$, whose image in $\pi_1(U)$ is
trivial, but which has no centre on $U$.  For this example
$\Delta=0$; the inertia-preserving coarsening of Claim~B.98 does
not repair the centre criterion.
\end{claim}
\begin{proof}
In the Hahn field $K((t^{\mathbb Q}))$ put
\[
 N_j=j!\quad\hbox{or}\quad N_j=2^j,\qquad
 a_j=j+1/N_j,\qquad x=t,\quad y=2+\sum_{j\ge1}t^{a_j}.
\]
The increasing supports are well ordered.  Their rational expressions
are series with support finite below any real bound, so sums,
products, and inverses obtained by the geometric-series expansion
are defined.  The value group is exactly
$\bigcup_m N_m^{-1}\mathbb Z$.  To prove the nontrivial inclusion,
let
\[
 P_m(X,Y)=\operatorname{Norm}_{K(X^{1/N_m})/K(X)}
 \left(Y-2-\sum_{j\le m}X^{a_j}\right).
\]
The $N_m$ conjugate partial sums are distinct: the numerator of the
last exponent is coprime to $N_m$.  At $(t,y)$ the identity factor
has value $a_{m+1}$; all other factors have values in
$N_m^{-1}\mathbb Z$.  Hence
$v(P_m(t,y))=a_{m+1}+c_m$ with $c_m$ in that group.  Starting from
$v(x)=1$ and $v(y-2)=a_1$, induction yields every $1/N_m$.
The reverse inclusion follows by examining supports.  Algebraicity
of $y$ over $K(t)$ would bound the value-group index over
$\mathbb Z$ by the finite field degree, contradicting these
unbounded denominators.  Thus $K(x,y)$ is a rational function field.
Leading coefficients lie in $K$, so the residue field is $K$.
The centre is $(0,2)$, on the single boundary branch $x=0$.
For the factorial choice $\Gamma=\mathbb Q$ and
Claim~B.97 gives $I_v=1$.  The uniqueness of a centre on the
proper separated compactification rules out a centre on $U$.

For the stronger example choose a positive irrational real number
$\alpha$, and work in $K((t^{\mathbb Q\alpha+\mathbb Z}))$ with
the order inherited from $\mathbb R$.  Use the factorial sequence and
put
\[
 x=t^\alpha,\qquad
 y=2+\sum_{j\ge1}t^{\alpha a_j},\qquad z=2+t.
\]
The preceding norm argument gives value group $\mathbb Q\alpha$
on $K(x,y)$.  The element $t$ is transcendental over this field:
in a polynomial $\sum b_jt^j$, the nonzero terms have pairwise
distinct values $q_j\alpha+j$.  No cancellation is possible.
Therefore $x,y,z$ are algebraically independent, the total value
group is $\mathbb Q\alpha\oplus\mathbb Z$, and the residue is $K$.
Its largest divisible subgroup is $\mathbb Q\alpha$ (project to the
second summand), whereas its only convex subgroups are $0,\Gamma$
because it is archimedean.  Thus $\Delta=0$ and
$I_v=\operatorname{Hom}(\Gamma,T)=T$.
Its centre is $(0,2,2)$, with only $x=0$ on the boundary.  The local
inertia map is evaluation at $v(x)=\alpha$, which vanishes because
$\alpha\in\Gamma_{\rm div}$.  The other two curve coordinates have
interior centres.  The product formula for geometric fundamental
groups, or Claim~B.103 below, now proves that its entire global
inertia image is trivial.  This is a nonzero inertia counterexample
and does not use a zero source group.
\end{proof}
\paragraph{Check record.} checked only in this writing, 2026-11-18: GPT-6 Astra (OpenAI), ChatGPT Pro mode; internal effort value not exposed. Own detailed proof written and reread; primary source hypotheses matched. Through N3.7 read; targeted checks passed, as root message records


\begin{claim}{B.100}{Proved}{The local calculation and its exact image}
Let $X/K$ be smooth proper and $D=X\setminus U$ a simple normal
crossings divisor.  At a geometric centre $\bar x$ of $v$, write
$z_1,\ldots,z_r$ for the branch parameters and $\gamma_i=v(z_i)$.
The local geometric group is $T^r$ and the inertia homomorphism is
\[
 \rho_v:\operatorname{Hom}(\Gamma,T)\longrightarrow T^r,
 \qquad\chi\longmapsto(\chi(\gamma_1),\ldots,\chi(\gamma_r)).
\]
Its dual is
\[
 \psi_v:(\mathbb Q/\mathbb Z)^r\longrightarrow
 (\Gamma\otimes\mathbb Q)/\Gamma,
 \qquad(q_i)\longmapsto\sum q_i\gamma_i\bmod\Gamma.
\]
Its image is the annihilator of $\ker\psi_v$.  It is open exactly
when that kernel is finite, and surjective exactly when that kernel
vanishes.  Openness requires rational independence of the
$\gamma_i$; under independence its exact criterion is
\[
 [\Gamma\cap\operatorname{span}_{\mathbb Q}(\gamma_i):
       \textstyle\sum_i\mathbb Z\gamma_i]<\infty.
\]
Surjectivity means this index is one.  This re-proves
FLAGGED\_149 E3, including valuations of every rank.
\end{claim}
\begin{proof}
In the strict henselization of the local ring at $\bar x$, adjoining
all $n$th roots of the branch equations gives covers with group
$(\mu_n)^r$.  They are cofinal.  Indeed normalize in an arbitrary
finite cover of the punctured strictly local scheme and choose $n$
divisible by its finite tame inertia orders at the branches.
Abhyankar's lemma kills that codimension-one ramification after the
Kummer base change, whose target is regular.  Purity then makes the
normalization finite \emph{\'etale}, and strict henselianity makes it
trivial.  Normalization is finite here by excellence.  This proves
the local $T^r$ description without an assertion about its global
injectivity.

Pulling back $z_i^{1/n}$ to the valuation field gives the character
$\gamma_i/n\bmod\Gamma$ of Claim~B.97.  Passing through every
$n$ proves both formulas.  Replacing $z_i$ by a unit multiple does
not change them, because units have all roots in the strictly
henselian valuation field.  Character duality is exact, since
$\mu_\infty$ is divisible and hence injective as an abelian group.
It identifies the quotient of $T^r$ by the compact image with the
dual of $\ker\psi_v$, giving the two image tests.

An integer relation among the values supplies an infinite copy of a
quotient of $\mathbb Q/\mathbb Z$ in the kernel.  Conversely, under
independence, the map $\mathbb Q^r\to\Gamma\otimes\mathbb Q$
is injective, so its preimage of $\Gamma$, modulo $\mathbb Z^r$,
is exactly the saturation quotient displayed above.  For a DVR
the formula reads $a\mapsto(m_1a,\ldots,m_ra)$ with $m_i>0$.
It is not open for $r>1$.  For the $\mathbb Z[1/2]$ example of
Claim~B.99 on the one-branch stratum, the image is
$\prod_{\ell\ne2}\mathbb Z_\ell(1)$, also not open.  Thus
rational rank equal to the number of branches is not enough.
\end{proof}
\paragraph{Check record.} checked only in this writing, 2026-11-18: GPT-6 Astra (OpenAI), ChatGPT Pro mode; internal effort value not exposed. Own detailed proof written and reread; primary source hypotheses matched. Through N3.7 read; targeted checks passed, as root message records


\begin{claim}{B.101}{False as stated}{Local boundary groups need not inject globally}
Even for a product of two hyperbolic curves, a smooth proper simple
normal crossings compactification can have a noninjective boundary
homomorphism $T^r\to\pi_1(U)$.  For
$U=(\mathbb P^1-\{0,1,\infty\})^2$, blow up $(0,2)$ in
$(\mathbb P^1)^2$.  At the intersection of the exceptional divisor
and the strict transform of $x=0$, the homomorphism is
\[
 T^2\longrightarrow T\hookrightarrow\pi_1(U),\qquad
 (a,b)\longmapsto a+b.
\]
Its kernel is the antidiagonal.  This counterexample concerns good
smooth SNC compactifications; it does not assert that the blow-up is
a compactification adapted as a smooth relative curve at every level
of the original tower.
\end{claim}
\begin{proof}
The chart $x=uv$, $y=2+v$ has boundary $uv=0$.  Each positive local
meridian goes once around $x=0$.  The $y$ coordinate lies in a
neighbourhood of the interior point $2$, so gives no meridian in the
second factor.  These identities hold in every finite Kummer
quotient, hence in $T$.  Cusp inertia embeds in the first punctured
line: its standard generators give a free profinite group of rank
two, and a meridian is a primitive generator.  The kernel is
therefore exactly as stated.  Equivalently the calculation can be
made with complex positive loops and Riemann existence.

The nonzero kernel does \emph{not} contradict the discrete-valuation
centre criterion.  A DVR centred at this crossing has positive
weights $m_1,m_2$, and its inertia maps by multiplication by
$m_1+m_2>0$, an injective endomorphism of $T$.  The supplied program
\texttt{n3\_boundary\_certificate.py} checks the coordinate matrix,
its integer kernel, and the stated finite Smith calculations; the
positive-weight assertion for all integers is the displayed formula.
\end{proof}
\paragraph{Check record.} checked only in this writing, 2026-11-18: GPT-6 Astra (OpenAI), ChatGPT Pro mode; internal effort value not exposed. Own detailed proof written and reread; primary source hypotheses matched. Through N3.7 read; targeted checks passed, as root message records


\begin{claim}{B.102}{Proved}{A larger valid form of the centre margin}
Let $U=U_d\to\cdots\to U_0=\operatorname{Spec}K$ be a tower of
smooth geometrically connected curves, each with a smooth proper
relative compactification and finite \emph{\'etale} relative boundary.
Suppose no factor has type $(0,0)$ or $(0,1)$.  If
$\Gamma_{v,\rm div}=0$, then
\[
 I_v\longrightarrow\pi_1(U)\text{ is trivial}
 \quad\Longleftrightarrow\quad
 v\text{ has a centre on }U.
\]
This holds in every dimension and for every ordered rank.  In
particular it proves the DVR assertion of FLAGGED\_149 D7 and
extends it to all finitely generated value groups and to groups such
as $\mathbb Z[1/2]$.  The forward implication is false without
the value-group restriction, even when $I_v\ne1$, by
Claim~B.99.  The reverse implication needs no type or value-group
restriction.
\end{claim}
\begin{proof}
We first supply the fibre-injectivity fact needed when elliptic and
$\mathbb G_m$ factors occur; using a hyperbolic-only theorem there
would omit these cases.  Over an algebraically closed field of
characteristic zero, every curve in the tower is an aspherical
complex surface after spreading out and embedding a field of
definition into $\mathbb C$.  Here ``surface'' means a real
two-manifold.  Compact curves of genus at least one and punctured
curves other than the projective line are aspherical.  A smooth proper
pair with disjoint boundary sections is a locally trivial topological
bundle, by local triviality of the proper smooth family together with
isotopy of its finite sections.  The homotopy sequence therefore gives
a discrete exact sequence of fundamental groups with injective fibre
and an aspherical total space, inductively along the tower.

We verify that the fibre remains injective after \emph{full}
profinite completion.  Recall that a group is good if its cohomology
with every finite module agrees with continuous cohomology of its
profinite completion.  Free groups of finite rank and oriented compact
surface groups of positive genus are good.  Here is the usual
effacement proof, included to specify this input.  After a finite
cover the module action is trivial.  A class in degree one is killed
by the kernel of its homomorphism to the finite module.  A class in
degree two for a compact surface is killed by a cover whose degree is
divisible by the exponent of the module; such covers exist in every
positive degree, and pullback multiplies the orientation class by the
degree.  Higher cohomology vanishes; for free groups degrees at least
two vanish.  The same argument holds in every finite cover.  This
effacement criterion implies goodness: the injections into finite
coinduced modules identify the two universal cohomological
$\delta$-functors, since continuous finite-module cocycles are
effaced by passage to a sufficiently small open subgroup.  Thus the
comparison is an isomorphism in every degree.

For clarity, the extension step is also proved.  Let
$1\to A\to B\to C\to1$ be one of the discrete sequences, with
$A$ a finitely generated surface or free group and with $C$ good,
as are its finite-index subgroups.  Choose any finite-index
characteristic subgroup $A_0\subset A$ and set $Q=A/A_0$.
The quotient $B/A_0$ is an extension of $C$ by finite $Q$.
On a finite-index subgroup $C_0$ the outer action on $Q$ is trivial.
The centralizer of $Q$ in its inverse image maps onto $C_0$, with
kernel $Z(Q)$, so is a central extension classified by
$H^2(C_0,Z(Q))$.  Goodness makes this class continuous.  Its cocycle
factors through a finite quotient, hence becomes zero on a further
finite-index subgroup $C_1$.  A splitting there gives a finite-index
subgroup of $B/A_0$ meeting $Q$ trivially.  Taking its normal core
produces a finite quotient of $B$ which is injective on $Q$.
Because $A$ is finitely generated, characteristic finite-index
subgroups are cofinal.  Hence the induced topology on $A$ is its
full profinite topology, and
\[
 1\to\widehat A\to\widehat B\to\widehat C\to1
\]
is exact.  The Lyndon--Hochschild--Serre spectral sequences for the
discrete and profinite extensions now prove goodness of $B$:
$H^q(A,M)$ is finite for finite $M$, and goodness of $A$ and $C$
identifies their entire $E_2$ pages.  Finite-index subgroups have
the same extension structure and satisfy the same argument.  Induction
therefore proves the required injectivity at every stage.  Riemann
existence and invariance under algebraically closed characteristic-zero
extension transfer it to the geometric \emph{\'etale} groups.
For hyperbolic-only towers the same result is Hoshi (2014),
Proposition~2.4(i)--(ii); its published hypotheses are connected
noetherian separated normal bases over a characteristic-zero field.

If $v$ has a centre, pullback of every finite \emph{\'etale} cover
to the strict henselization of its valuation ring is constant.
Consequently inertia acts trivially.  Conversely suppose the image
of $I_v$ is trivial and extend the generic map successively along
the tower.  It extends at level zero.  Suppose it has extended to
$\operatorname{Spec}V\to U_{i-1}$.  The proper relative
compactification of $U_i$ extends the generic lift uniquely to
$\operatorname{Spec}V$.  If its closed point meets the relative
boundary, then after strict henselization a single boundary section
has local equation $z$ and $\gamma=v(z)>0$.  The induced inertia
map to its cusp group is evaluation at $\gamma$, by the Kummer
calculation in Claim~B.100.  It is nonzero because
$\Gamma_{\rm div}=0$.  This cusp group embeds in the geometric
fibre group, which embeds in $\pi_1(U_i)$ by the preceding argument,
contradicting the assumed trivial image under
$\pi_1(U)\to\pi_1(U_i)$.

Here is a local-ring verification of that comparison which avoids
any noetherian hypothesis on $V$, and does not apply a DVR theorem
to it.  Let $b$ be the centre in $U_{i-1}$ and $x$ its lift on the
relative compactification.  The total space is smooth over $K$ and
therefore noetherian and regular.  Its strictly henselian local ring
$R=\mathcal O^{\rm sh}_{\overline U_i,\bar x}$ maps to
$V^{\rm sh}$ by the universal property of strict henselization,
after fixing the compatible separably closed residue fields.  Near
$x$ the relative boundary has one branch, so
$\pi_1(\operatorname{Spec}R[1/z])=T$ by the noetherian local
Kummer--purity argument of Claim~B.100.  The punctured strictly
local ring in the geometric fibre at $\bar b$ also has group $T$;
its map to the total punctured local scheme induces the identity on
these $T$'s, since in both cases their finite characters are the
roots of the same $z$.  Thus the total local $T\to\pi_1(U_i)$
is exactly the cusp map in that geometric fibre followed by its
already proved injection.  Pullback along $R\to V^{\rm sh}$
computes it by evaluation at $v(z)$ using Claim~B.97.
Strict henselization preserves the value group.  This proves the
comparison for every valuation ring directly through noetherian
local rings on the given tower.

For a proper elliptic factor there is no boundary, so properness
already extends the map.  For a $\mathbb G_m$ factor its cusp group
is $T$ and the same Kummer calculation applies.  For every hyperbolic
affine factor cusp inertia is injective.  One direct verification of
the latter is to pass to a finite cover with at least two punctures;
then the cusp has a nonzero primitive class in the free abelianization.
Such a cover exists for a once-punctured positive-genus curve by any
nontrivial unramified cover of its proper completion; genus-zero
hyperbolic curves already have at least three punctures.  A
nontrivial multiple of a cusp generator is still detected.  This
checks every allowed type.  Avoiding the boundary at the closed
point means avoiding it on all of the local scheme $\operatorname{Spec}V$.
The induction gives the required centre without a descent gap.
\end{proof}
\paragraph{Check record.} checked only in this writing, 2026-11-18: GPT-6 Astra (OpenAI), ChatGPT Pro mode; internal effort value not exposed. Own detailed proof written and reread; primary source hypotheses matched. Through N3.7 read; targeted checks passed, as root message records


\begin{claim}{B.103}{Proved}{An exact answer for products, with every factor type}
Let $U=\prod_{j=1}^d C_j$ be a product of smooth geometrically
connected curves in characteristic zero, with proper smooth
completions.  Let $x$ be the centre of $v$ on their product, and for
each coordinate that is a missing point choose its local parameter
$z_j$ and write $\gamma_j=v(z_j)>0$.  Then
\[
 I_v\to\pi_1(U)\text{ is trivial}
 \quad\Longleftrightarrow\quad
 \gamma_j\in\Gamma_{\rm div}\quad(j\in B_{\ne\mathbb A^1}).
\]
Here $B_{\ne\mathbb A^1}$ is the set of boundary coordinates whose
curve factor is not geometrically isomorphic to $\mathbb A^1$.
Proper factors contribute no condition; affine-line factors
contribute no condition.  In particular projective-line and elliptic
proper factors, and punctured-line factors of type $(0,2)$, are all
included.
\end{claim}
\begin{proof}
Inertia is geometric because the valuation is trivial on $K$.
The geometric fundamental group of the product is the product of
the curve groups; this follows over $\mathbb C$ from the product
of topological groups and finite generation, and then by Riemann
existence and base-extension invariance.  At an interior coordinate
the inertia image is zero by the centre implication.  At a boundary
coordinate the image is evaluation at $\gamma_j$ followed by that
curve's cusp homomorphism.  This homomorphism is zero for
$\mathbb A^1$, whose geometric group is trivial, and is injective
for every other nonproper curve by the cusp proof in
Claim~B.102 (including $\mathbb G_m$).  Independent factors cannot
cancel in a direct product.  Claim~B.97 gives precisely the test
displayed above.  The argument is unchanged after the finite residue
extension needed to split a boundary point, since that extension is
unramified for the valuation.
\end{proof}
\paragraph{Check record.} checked only in this writing, 2026-11-18: GPT-6 Astra (OpenAI), ChatGPT Pro mode; internal effort value not exposed. Own detailed proof written and reread; primary source hypotheses matched. Through N3.7 read; targeted checks passed, as root message records


\begin{claim}{B.104}{Proved}{The excluded factor types can kill a later meridian}
For $\mathbb A^1$ a DVR at infinity has no centre and trivial
geometric inertia image.  The type $(0,0)$ can also cause failure in
a nontrivial tower: the $\mathbb G_m$-bundle
\[
 U=\mathcal O_{\mathbb P^1}(-1)^\times\longrightarrow\mathbb P^1
\]
is an elementary tower of types $(0,0),(0,2)$, and its geometric
fundamental group is trivial.  A divisorial valuation along the zero
section of $\mathcal O(-1)$ has nonzero inertia and no centre on
$U$, yet has trivial global image.
\end{claim}
\begin{proof}
The first assertion follows from
$\pi_1(\mathbb A^1_{\overline K})=1$ in characteristic zero.
The map sending a nonzero vector in a line in $K^2$ to that vector
identifies the second total space with $\mathbb A^2\setminus\{0\}$.
Over $\mathbb C$ it retracts onto $S^3$, so is simply connected;
Riemann existence gives the geometric assertion over every
characteristic-zero field.  The zero section is a smooth divisor in
the line bundle and supplies a DVR with inertia $T$.  It has no
centre on the complement.  A good relative compactification of the
bundle is $\mathbb P(\mathcal O\oplus\mathcal O(-1))$, with its
zero and infinity sections as boundary.  In topological terms the
boundary map from $\pi_2(\mathbb P^1)$ to the fibre $\pi_1$ is
multiplication by $-1$, the degree of the line bundle.  This also
explains why a fibre-injection argument cannot include projective-line
bases indiscriminately.
\end{proof}
\paragraph{Check record.} checked only in this writing, 2026-11-18: GPT-6 Astra (OpenAI), ChatGPT Pro mode; internal effort value not exposed. Own detailed proof written and reread; primary source hypotheses matched. Not yet reported


\begin{claim}{B.105}{Proved}{A global criterion using finite-cover residues}
Keep the SNC compactification and centre of Claim~B.100.  Let
$P_x$ be its punctured strictly local neighbourhood over
$\overline K$, with group $T^r$.  Define a subgroup
$\mathcal R_x\subset(\mathbb Q/\mathbb Z)^r$ as follows.
For every connected finite \emph{\'etale} cover
$W\to U_{\overline K}$ admitting a pointed lift
$P_x\to W$, for every $n>0$, and every class
$\alpha\in H^1(W,\mu_n)$, include its restriction to $P_x$.
The latter is the residue vector
$(a_1/n,\ldots,a_r/n)$ in the indicated group.  Then
\[
 I_v\to\pi_1(U)\text{ is trivial}
 \quad\Longleftrightarrow\quad
 \sum_i a_i\gamma_i\in n\Gamma
 \quad\text{for every such finite-cover residue vector}.
\]
This is valid for every smooth $U$, every elementary tower and every
factor type, and every valuation.  It uses cyclic cohomology on
finite covers, rather than an unspecified kernel of
$T^r\to\pi_1(U)$.
\end{claim}
\begin{proof}
Write $G=\pi_1(U_{\overline K})$ and $f:T^r\to G$ for the local
homomorphism.  A pointed lift to the cover defined by an open subgroup
$H\subset G$ exists exactly when $f(T^r)\subset H$; it is the
chosen component on which the local cover is trivial.  The residue
classes above are therefore precisely restrictions of finite-order
characters of these open subgroups.  They separate every nonidentity
element in the compact abelian subgroup $f(T^r)$.  To verify this,
let $a\ne1$ belong to that subgroup, choose a finite quotient
$q:G\to Q$ with $q(a)\ne1$, and put
$A=q(f(T^r))$ and $H=q^{-1}(A)$.  The finite group $A$ is abelian.
A cyclic character of $A$ detects $q(a)$; pulled back to $H$ it is
an element of $H^1(W,\mu_n)$ for the corresponding cover $W$.
It therefore belongs to the defining family and detects $a$.

It follows that the global inertia image is zero exactly when all
these characters vanish on the local image of $I_v$.  The dual
formula in Claim~B.100 makes this condition
$\sum(a_i/n)\gamma_i=0\bmod\Gamma$, which is the displayed
divisibility test.  All parts of the construction are geometric
finite-cover and Kummer-cohomology data, and the valuation enters
only through finitely many values in each test.  The residue vector
can equivalently be computed by extending a cyclic torsor to the
normalization and reading its Kummer orders on the branches.
The proof uses arbitrary finite quotients, so it is a profinite
statement, not a bounded-quotient experiment.

Here is a description of the same module in terms of finite
\'etale algebras and divisors, specifying how to calculate its
characters.  For each such pointed cover $W$, put
$F_W=\overline K(W)$ and
\[
 \mathcal K_n(W)=
 \{h\in F_W^\times:\operatorname{div}_W(h)
       \in n\operatorname{Div}(W)\}/F_W^{\times n}.
\]
The Kummer identification $H^1(W,\mu_n)=\mathcal K_n(W)$
holds because $W$ is regular and $n$ is invertible.  Indeed, a
torsor restricts to $Y^n=h$ at the generic point;
unramifiedness at codimension one makes each divisor order a
multiple of $n$.  Conversely these divisibilities make the
normalization in $Y^n=h$ unramified at codimension one: locally
divide $Y$ by the appropriate uniformizer power, leaving the
\'etale equation $Y^n=u$ for a unit.  Purity on the regular $W$
makes the normalization \'etale everywhere.  Generic restriction
is injective because a generic section of a finite \'etale torsor
over a normal scheme extends uniquely.  This proves the
identification, including torsors from nontrivial $n$-torsion
line bundles; it is not restricted to global units.

The pointed lift embeds $F_W$ in the fraction field of the chosen
strictly local branch.  For $h\in\mathcal K_n(W)$, factor its
pullback in that regular strictly local ring, allowing poles.
All irreducible factors except the boundary parameters $z_i$
have orders divisible by $n$.  Units have $n$th roots by
Hensel's lemma.  Its character vector is therefore
\[
 \left(\frac{\operatorname{ord}_{z_1}(h)}n,\ldots,
       \frac{\operatorname{ord}_{z_r}(h)}n\right)
       \pmod{\mathbb Z^r}.
\]
Thus $\mathcal R_x$ may be defined using these divisor orders
alone, ranging over finite \'etale algebras with the stated local
section.  This definition does not require a previously unknown
kernel of $T^r\to G$.  For any specified cover and $n$, the input
is the divisor group, the rational functions whose divisors are
divisible by $n$, and the chosen local branch.

For the original product compactification of Claim~B.103 the module
$\mathcal R_x$ is the full coordinate module on the non-affine-line
boundary factors.  For the blow-up in Claim~B.101 it is the
diagonal subgroup $\{(q,q):q\in\mathbb Q/\mathbb Z\}$.
Thus its exact test there is
$\gamma_1+\gamma_2\in\Gamma_{\rm div}$, not separate
divisibility of the two values.  These computations exhibit how a
noninjective boundary group is handled without mistaking its local
image for its global one.
\end{proof}
\paragraph{Check record.} checked only in this writing, 2026-11-18: GPT-6 Astra (OpenAI), ChatGPT Pro mode; internal effort value not exposed. Own detailed proof written and reread; primary source hypotheses matched. Not yet reported


\paragraph{Scope of the determination.}
Claim~B.105 is an exact finite-cover cohomological criterion, with a
proof that the indicated abelian tests suffice even when the ambient
fundamental group is nonabelian.  It is not a finite algorithm from
the list of factor types, nor a computed classification of
$\mathcal R_x$ for every nonproduct tower.  That limitation of the
finite-cover criterion is superseded by Claims~B.124--B.125:
for every tower in the Artin good-pair convention they identify
the global boundary kernel with the closure of its finite integral
relation lattice, calculate that lattice recursively from the
discrete bundle data, and give the exact Smith divisibility test
for every valuation rank and every factor type.  This is a finite
geometric determination; no uniform terminating algorithm from
arbitrary defining equations is asserted.  The proofs above settle
the universal local-injectivity assertion negatively, verify D3--D5
and E3, verify D7 for DVRs with every allowed factor type, give the
larger value-group range of the centre margin, and give nonzero
invisible inertia outside that range.  They do not infer a global
centre theorem from local injectivity at a crossing.

\paragraph{Checks and provenance.}
The displayed arguments were written and reread by GPT-6 Astra
(OpenAI), ChatGPT Pro mode, on 18 November 2026.  There is no claim
of human or different-model verification.  The finite code checks
only the stated integral matrices and finite quotients.  Temkin's
valuation notes, Kuhlmann's Chapter~9, Hoshi's published paper,
the supplied SGA~1 source, and Stacks Tag~0BMB were actually opened.
No priority claim is made for these formulations.

\subsection{Finite boundary relations for all elementary factor types}

In this subsection an elementary curve fibration means an open
subscheme of a smooth proper relative curve whose complement is
finite \'etale over the base.  This is the Artin good-pair convention
used in the surrounding manuscript: cote~149, scans~011, 075 and
079; the question is in scan~120, manuscript page~59.  Empty
boundary is allowed.  Thus all types $(g,n)$ occur, including
$(0,0)$, $(0,1)$, $(0,2)$ and $(1,0)$.  An arbitrary smooth
nonproper morphism with changing behaviour at infinity is not
silently substituted for such a fibration.

\begin{claim}{B.124}{Proved}{Abelian subgroups retain their full topology in every elementary tower}
Let $U_d\to\cdots\to U_0=\operatorname{Spec}K$ be an elementary
tower in the preceding sense, where $K$ has characteristic zero
and every fibre is geometrically connected.  Choose a complex
realization of a field of definition for the tower and put
$G_i=\pi_1^{\rm top}(U_i(\mathbb C))$.
For every stage there is an exact sequence of discrete groups
\[
 1\longrightarrow N_i\longrightarrow G_i
       \longrightarrow G_{i-1}\longrightarrow1.             \tag{N3.L1}
\]
The group $N_i$ is a finitely generated free group or a closed
orientable surface group when the fibre is hyperbolic.  For
fibre type $(0,2)$ or $(1,0)$ it is respectively a quotient of
$\mathbb Z$ or $\mathbb Z^2$ by the image of the topological
transgression; for types $(0,0)$ and $(0,1)$ it is trivial.
In particular the abelian quotients here may have torsion.

Every $G_i$ and each of its finite-index subgroups is residually
finite and good for cohomology with finite modules.  Completion
of (N3.L1) is exact, including injectivity at its left end.
Moreover, for every finitely generated abelian subgroup
$A\subset G_i$, the homomorphism
\[
                 \widehat A\longrightarrow\widehat G_i
                                                               \tag{N3.L2}
\]
is injective.  Thus the profinite topology induced on $A$ is
its full profinite topology.
\end{claim}

\begin{proof}
\medskip\noindent\textbf{The discrete kernels, including the exceptional types.}
A smooth proper complex curve family with a finite \'etale
boundary is a differentiable locally trivial bundle of pairs.
For completeness, lift coordinate vector fields on a small
ball in the base to vector fields on the compact total fibres,
choose them tangent to the finitely many moving marked points,
and integrate them.  Compactness supplies a common interval of
existence.  The resulting trivializations of the proper pair
restrict to trivializations of its complement.  The homotopy
sequence therefore gives
\[
 \pi_2(U_{i-1}(\mathbb C))\xrightarrow{\partial_i}
 \pi_1(F_i)\longrightarrow G_i\longrightarrow G_{i-1}
 \longrightarrow1 .
\]
The image of $\partial_i$ is central in $\pi_1(F_i)$.  One can
check this without an asphericity assumption on the base.
Pull the bundle back to a representative sphere and trivialize
it over the two hemispheres.  The clutching map is a loop of
self-diffeomorphisms of $F_i$ in the identity component.  Its
evaluation at a point gives the boundary loop.  For any loop
$a$ of $F_i$, evaluation of the clutching loop on all points
of $a$ gives a square homotopy whose boundary is the commutator
of these two loops.  Their classes commute.

A hyperbolic finite-type surface has centre-free fundamental
group.  For a free group this follows from its free action on
its Cayley tree; for a closed surface of genus at least two,
nonidentity elements act as hyperbolic isometries of the
universal covering disc, and their centralizers preserve their
axes and are cyclic.  Thus $\partial_i=0$ in these cases.
The other fibre groups are respectively $\mathbb Z,\mathbb
Z^2,1,1$, which gives exactly the asserted list of $N_i$.
In the abelian cases the transgression image is a subgroup of
$\mathbb Z^e$, $e\le2$, and has a finite integral basis.  It
must be divided out before taking a profinite completion.

\medskip\noindent\textbf{The elementary group-cohomology inputs.}
The goodness and finite-module cohomology finiteness of free
and positive-genus orientable surface groups were proved in
B.102 by effacement.  Explicitly, first kill the module action
by a finite cover.  Degree-one classes are killed by kernels
of finite homomorphisms.  On a compact orientable surface,
degree-two pullback multiplies the orientation class by the
cover degree, so a degree divisible by the module exponent
kills it; higher cohomology vanishes.  Free groups have no
cohomology above degree one.  The same assertions hold for all
finite-index subgroups.  Effacement on finite-index subgroups,
or equivalently the comparison of the two universal
cohomological $\delta$-functors using finite coinduced modules,
proves goodness.

Finitely generated abelian groups have the same properties.
For $\mathbb Z$ the two-term resolution with differential
$t-1$ gives the comparison with the procyclic group for every
finite module.  Iterating the Lyndon--Hochschild--Serre
comparison proves it for $\mathbb Z^e$.  A finite group equals
its completion; taking a product with it and using the same
spectral sequence handles the torsion part.  In every fixed
degree all cohomology groups are finite: use the finite-rank
resolutions of the free part and the finite bar terms of the
finite part.  Every subgroup of finite index is of the same
form.  These basic groups are residually finite; for the
surface and free cases this also follows from the separability
theorem used below.

\medskip\noindent\textbf{Completion of an extension.}
Suppose $1\to N\to G\to Q\to1$ is an extension in which $N$
is one of the groups just listed, and $Q$ and its finite-index
subgroups are good and residually finite.  Given a
finite-index characteristic subgroup $N_0\subset N$, put
$F=N/N_0$.  After replacing $Q$ by a finite-index subgroup
$Q_0$, its outer action on the finite group $F$ is trivial.
The centralizer of $F$ in the inverse image of $Q_0$ in
$G/N_0$ then surjects onto $Q_0$, with kernel $Z(F)$:
an element with inner conjugation action can be multiplied
by an element of $F$ to centralize it.  The resulting central
extension has a class in $H^2(Q_0,Z(F))$.

Goodness identifies this class with a continuous class on
$\widehat Q_0$.  A continuous cocycle with finite values
factors through a finite quotient; after normalizing it,
restriction to that quotient's kernel is zero.  The extension
therefore splits on a finite-index subgroup $Q_1\subset Q_0$.
The image of a splitting is a finite-index subgroup of
$G/N_0$ intersecting $F$ trivially.  Its normal core gives
a finite quotient of $G$ which is injective on $F$.
Characteristic subgroups of finite index are cofinal in $N$,
since $N$ is finitely generated.  Consequently the induced
topology on $N$ is full and
\[
 1\longrightarrow\widehat N\longrightarrow\widehat G
      \longrightarrow\widehat Q\longrightarrow1
\]
is exact.  The kernel assertion on the right is the usual
quotient property of completion: finite quotients killing
$N$ are precisely the finite quotients of $Q$.
Residual finiteness follows as well.  An element outside
$N$ is detected in a quotient of $Q$, and a nonidentity
element of $N$ is detected by the finite quotients just
constructed.

For a finite $G$-module $M$, the discrete and continuous
Lyndon--Hochschild--Serre spectral sequences have the same
$E_2$ pages, because $N$ and $Q$ are good and
$H^b(N,M)$ is finite for every $b$.  The first-quadrant
comparison proves that $G$ is good.  A finite-index subgroup
of $G$ is an extension of a finite-index subgroup of $Q$
by a finite-index subgroup of $N$, so the entire argument
also applies to it.  Induction proves these assertions for
all stages of the tower.

\medskip\noindent\textbf{The topology on abelian subgroups.}
Every finitely generated abelian subgroup of a finitely
generated abelian group has full induced topology, by Smith
normal form.  Every nontrivial abelian subgroup of a free
or closed hyperbolic surface group is cyclic, by the
centralizer argument above.  Its powers are separable:
Scott's Theorem~3.2, equivalently Theorem~3.3, says that
every finitely generated subgroup of a surface group is
closed in its profinite topology
\citep[Theorems~3.2--3.3, p.~561]{ROOT_Scott1978}.
We use that theorem with the correction to its hyperbolic
convexity proof in \citet[\S1, pp.~217--219]{ROOT_Scott1985}.
For a cyclic subgroup $A=\langle a\rangle$ and $m>0$,
separate the finitely many elements $a,\ldots,a^{m-1}$
from $\langle a^m\rangle$ and intersect the resulting
finite-index subgroups.  Its intersection with $A$ is
exactly $\langle a^m\rangle$.  Taking a normal core shows
that the induced topology is full.  This also covers the
finitely generated free groups, which are surface groups
of noncompact surfaces.

Now let $A\subset G$ be finitely generated abelian in an
extension as above.  Put $A_0=A\cap N$ and let $A_1$ be
its image in $Q$.  Both are finitely generated abelian.
Their completions form an exact sequence
\[
 0\longrightarrow\widehat A_0\longrightarrow\widehat A
        \longrightarrow\widehat A_1\longrightarrow0 ,
\]
as is seen directly by decomposing finitely generated
abelian groups into their primary and free parts.
By induction $\widehat A_1\to\widehat Q$ is injective.
Hence an element of the kernel of
$\widehat A\to\widehat G$ belongs to $\widehat A_0$.
The latter injects into $\widehat N$ by the basic case,
and $\widehat N\to\widehat G$ is injective by the extension
argument.  The kernel is zero.  Starting at $G_0=1$
proves (N3.L2) at every stage.

Finally, the tower and its pairs descend to a finitely
generated subfield of $K$, which embeds into $\mathbb C$.
Geometric \'etale groups are unchanged by algebraically
closed characteristic-zero base extension, and Riemann
existence identifies them with the completions just used
\citep[Remark~3.3(c), Theorem~3.4]{MilneEtale}.
Thus no embedding of the entire field $K$ into $\mathbb C$
is assumed.
\end{proof}
\paragraph{Check record.} Checked only in this writing, 18 November 2026: GPT-6 Astra (OpenAI), ChatGPT Pro mode, root complete proof and reading; arithmetic-notes reader full second proof reading with local comparison and finite Kummer joins incorporated. Separate effort setting not exposed. B125 arithmetic certificate: 23863 exact comparisons in23 cases.


\begin{claim}{B.125}{Proved}{The finite integral boundary test for every elementary tower}
Let $U/K$ be any elementary tower of B.124, put $L=K(U)$,
and let $v$ be any valuation of $L$ trivial on $K$, with
value group $\Gamma$.  Choose a smooth proper
compactification $\overline U$ with simple normal crossings
boundary.  At the centre of a prolongation of $v$, label the
boundary branches $z_1,\ldots,z_r$, and put
$\gamma_i=v(z_i)>0$.  If the centre is in $U$, take $r=0$.
Choose corresponding positive complex meridians
$m_1,\ldots,m_r$ after spreading out the labelled local
pair.  They commute in the discrete group
$G=\pi_1^{\rm top}(U(\mathbb C))$.  Define its integral
relation lattice
\[
 \Lambda_x=\left\{(a_i)\in\mathbb Z^r:
                m_1^{a_1}\cdots m_r^{a_r}=1\text{ in }G\right\}.
                                                               \tag{N3.L3}
\]
This lattice has a finite basis.  Choose a unimodular
integer matrix $S$ giving Smith coordinates
\[
 S\Lambda_x=d_1\mathbb Z e_1\oplus\cdots\oplus
                 d_s\mathbb Z e_s,\qquad d_j>0,
\]
and set $\eta_j=\sum_i S_{ji}\gamma_i$.
Then, in the entire arithmetic group defining $U$,
\[
 I_v\longrightarrow E_U^{\overline L}\text{ is trivial}
 \quad\Longleftrightarrow\quad
 \begin{cases}
   \eta_j\in d_j\Gamma,&1\le j\le s,\\
   \eta_j\in\Gamma_{\rm div},&s<j\le r .
 \end{cases}                                                \tag{N3.L4}
\]
Here $\Gamma_{\rm div}=\bigcap_{n\ge1}n\Gamma$.
The test holds in every dimension, for every value group
and rank, and for all elementary factor types.
The finite relation lattice is obtained recursively from
the discrete extension data (N3.L1): first compute relations
after projection to $G_{i-1}$, then compute the relations
of their lifts in $N_i$.  These last calculations are
integer relations in a cyclic group or Smith calculations
in a finitely generated abelian group.  The transgression
subgroup is included in the latter's presentation.

The local boundary homomorphism
$\widehat{\mathbb Z}(1)^r\to\pi_1(U_{\overline K})$
is injective exactly when $\Lambda_x=0$.  Thus neither
the factor types alone nor the mere existence of a smooth
normal crossings compactification implies local injectivity.
\end{claim}

\begin{proof}
\medskip\noindent\textbf{Compactification and local comparison.}
Such a compactification exists for every smooth separated
finite-type $U$ in characteristic zero.  The exact result
we use is Temkin's Theorem~1.1.15: a regular finite-type
scheme separated over a noetherian quasi-excellent
characteristic-zero base, together with an open subscheme
whose complement is an SNC divisor, has a regular proper
compactification retaining that SNC boundary
\citep[Theorem~1.1.15]{ROOT_Temkin2016}.
Apply it over $\operatorname{Spec}K$ with the open
subscheme equal to $U$ itself and empty initial divisor.
Regularity is smoothness over the perfect field $K$.
Properness gives the centre of every valuation, without
a restriction on rank.

We may pass to algebraically closed constants and to a
strictly local neighbourhood at that centre.  Constants
have value zero, and extending them preserves the value
group and does not alter the inertia group in the geometric
kernel; this is the unramified comparison in B.97--B.100.
The local Kummer group is
$T^r$, where $T=\widehat{\mathbb Z}(1)$, and the inertia
map is
\[
 \operatorname{Hom}(\Gamma,T)\longrightarrow T^r,\qquad
          \chi\longmapsto(\chi(\gamma_1),\ldots,\chi(\gamma_r)).
                                                               \tag{N3.L5}
\]
This is the all-rank formula proved in B.97 and B.100,
not an application of a discrete-valuation comparison.

The local homomorphism into the geometric group is the
completion of the homomorphism
$\mathbb Z^r\to G$, $e_i\mapsto m_i$.
Over $\mathbb C$ this follows by taking a small polydisc
with boundary coordinates $z_i$: its complement has
fundamental group $\mathbb Z^r$, and the positive loops
have precisely the Kummer characters $z_i^{1/n}$.
Riemann existence identifies these characters and their
maps into the global group.

For arbitrary $K$, spread out the tower, the proper pair,
the finite labelled branches and a neighbourhood in
their intersection stratum to a finitely generated field.
Along a connected smooth stratum the normal punctured
polydiscs form a locally trivial torus bundle; transporting
a base point changes the local map by conjugation.
The same statement for finite \'etale covers follows from
the local Kummer description of B.100.  Given any finite
Galois cover, choose an integer $n$ divisible by its
ramification indices along the finitely many labelled
branches.  Abhyankar's lemma and purity make its
normalization \'etale after adjoining $z_i^{1/n}$.
Restriction to the stratum of this root cover is a finite
locally constant cover with its $\mu_n^r$-action.
An \'etale path along the connected labelled stratum
identifies these actions up to conjugacy.  This includes
geometric generic or nonclosed points and geometric
closed points.  Applied to every finite Galois cover, it
gives exactly the equality of kernels of the local maps,
so a complex realization at such a closed point suffices.
Characteristic-zero base extension and Riemann existence,
in the forms of \citet[Remark~3.3(c), Theorem~3.4]{MilneEtale},
then give the required local-global comparison.  Choices
of paths conjugate the map, and a change of Tate orientation
multiplies all coordinates by one unit of
$\widehat{\mathbb Z}$; neither changes its kernel or the
vanishing test.

\medskip\noindent\textbf{Why no additional profinite relations appear.}
Let $A=\langle m_1,\ldots,m_r\rangle\subset G$.
It is a finitely generated abelian group and
$A=\mathbb Z^r/\Lambda_x$.  Completion of this abelian
presentation is exact.  B.124 says that
$\widehat A\to\widehat G$ is injective.  Consequently
\[
 \ker\bigl(T^r\longrightarrow\pi_1(U_{\overline K})\bigr)
       =\overline{\Lambda_x}\,(1).                         \tag{N3.L6}
\]
This proves the substantive comparison with finite integral
relations; residual finiteness of $G$ alone would not prove
it, because it would not ensure the full induced topology
on $A$.  In particular a nonzero closed relation in the
profinite boundary group cannot appear without a discrete
integral relation.

\medskip\noindent\textbf{Smith coordinates and valuation divisibility.}
In Smith coordinates the quotient in (N3.L6) is
\[
 \widehat A\simeq
       \prod_{j=1}^s\mathbb Z/d_j\mathbb Z
         \ \times\ \widehat{\mathbb Z}^{\,r-s}.
\]
The image of (N3.L5) in its $j$th coordinate is evaluation
at $\eta_j$, reduced modulo $d_j$ for a finite factor.
The finite Kummer pairing is perfect, so evaluation modulo
$d_j$ is identically zero exactly when
$\eta_j\in d_j\Gamma$.  For clarity, every character of
$\Gamma/d_j\Gamma$ is a character of the subgroup
\[
 \Gamma/d_j\Gamma\lhook\joinrel\longrightarrow
       (\Gamma\otimes\mathbb Q)/\Gamma,\qquad
 \gamma\bmod d_j\Gamma\longmapsto\gamma/d_j\bmod\Gamma .
\]
It extends to a character of
$(\Gamma\otimes\mathbb Q)/\Gamma$, since the target
$\mu_\infty$ is divisible and hence injective as an
abelian group.  On this subgroup its values remain in
$\mu_{d_j}$.  Thus restricting the full inertia group
does not lose any character needed for this finite test.
Evaluation in $T$ is zero exactly on $\Gamma_{\rm div}$
by B.97.  This proves (N3.L4).
Inertia maps trivially to the constant-field Galois group,
so its image in the full arithmetic group is zero exactly
when the geometric image just calculated is zero.

\medskip\noindent\textbf{Recursive calculation of the finite lattice.}
We spell out the finite group data used by (N3.L3), so the
answer is not a reference to an unknown profinite kernel.
For any commuting tuple $a_1,\ldots,a_r$ in $G_i$,
project it to $G_{i-1}$.  Recursively let the columns of
an integer matrix $B$ form a basis of the relation lattice
of this projected tuple.  All possible relations of the
original tuple are therefore of the form $Bv$.
For each column $B_j$ compute
\[
                       b_j=\prod_k a_k^{B_{kj}}\in N_i.
                                                               \tag{N3.L7}
\]
These elements commute.
If $N_i$ is free or a closed hyperbolic surface group,
they generate a cyclic subgroup, or are all the identity.
In the nontrivial case choose its generator $b$ and write
$b_j=b^{e_j}$.  The required lattice is
$B\ker((e_j):\mathbb Z^{\operatorname{rank}B}\to\mathbb Z)$.
If they are all the identity it is $B\mathbb Z^{\operatorname{rank}B}$.
If $N_i=\mathbb Z^e/T_i$ is abelian, represent the $b_j$
by the columns of an integer matrix $C$.  Then take
\[
        B\{v\in\mathbb Z^{\operatorname{rank}B}:Cv\in T_i\}.
                                                               \tag{N3.L8}
\]
Integer kernels and Smith normal form give a finite basis.
For $N_i=1$ there is no further relation to impose.
The recursion starts at $G_0=1$, whose relation lattice
is $\mathbb Z^r$.

The subgroup $T_i$ is the image of $\partial_i$ in
$\pi_1(F_i)=\mathbb Z^e$, $e\le2$; its finite presentation
is part of the discrete bundle data.  It can be specified
by the evaluated clutching loops on spheres in the base,
choosing finitely many whose images generate this subgroup.
Every subgroup of $\mathbb Z^e$ is finitely generated,
so such finite data always exist.  For a line bundle
complement this is exactly the subgroup of spherical
Chern numbers.  We do not assert a uniform terminating
algorithm extracting these topological data from arbitrary
equations over an unspecified field.  The theorem is an
exact geometric determination by finite integral data,
with an explicit finite lattice calculation once those
data are specified; it makes no effectiveness claim of
that additional kind.

\medskip\noindent\textbf{All types and explicit checks.}
For a product compactification, the missing-point
meridians of every factor other than $\mathbb A^1$ are
independent, and those of $\mathbb A^1$ are zero.
This gives B.103 directly from (N3.L4).
For the blow-up of B.101 the relation lattice is generated
by $(1,-1)$; its remaining free coordinate is
$\gamma_1+\gamma_2$, so the test is
$\gamma_1+\gamma_2\in\Gamma_{\rm div}$.

More generally let
$U=\mathcal O_{\mathbb P^1}(-d)^\times$, with $d\in\mathbb Z$.
This is a tower of types $(0,0),(0,2)$, with proper relative
pair $\mathbb P(\mathcal O\oplus\mathcal O(-d))$ and the
zero and infinity sections removed.  The sphere
transgression is multiplication by $-d$, its clutching
degree.  At the zero boundary the lattice is $d\mathbb Z$.
For $d\ne0$ the exact test is
$v(z)\in |d|\Gamma$; for $d=0$ it is
$v(z)\in\Gamma_{\rm div}$.  Thus the same factor types
give different answers, including the complete killing
of inertia when $|d|=1$.  At a DVR with $v(z)=1$, the
image has order $|d|$ when $d\ne0$, and is all of $T$
when $d=0$.  This includes B.104 and shows concretely why
the transgression cannot be omitted.

The exact integer changes of basis, finite-factor
annihilator tests and these line-bundle examples are
checked in \texttt{code/n3\_lattice\_certificate.py}, with
its JSON certificate and saved output.  The program checks
the stated finite computations; the passage to every
valuation group is the Kummer and completion proof above.
Combining this theorem with B.102 gives the strongest
centre statement proved here: for towers without
$(0,0),(0,1)$ factors and $\Gamma_{\rm div}=0$, trivial
inertia is equivalent to a centre on $U$; outside that
range the complete answer is (N3.L4), and B.99, B.101
and B.104 exhibit the failures of the proposed broader
centre or local-injectivity assertions.
\end{proof}
\paragraph{Check record.} Checked only in this writing, 18 November 2026: GPT-6 Astra (OpenAI), ChatGPT Pro mode, root complete proof and reading; arithmetic-notes reader full second proof reading with local comparison and finite Kummer joins incorporated. Separate effort setting not exposed. B125 arithmetic certificate: 23863 exact comparisons in23 cases.


% AI-produced by GPT-6 Astra (OpenAI), ChatGPT Pro mode, 18 November 2026.
% Root second reading and configuration targeted second reading recorded, 18 November 2026.
% No priority or external/human verification is claimed.
\subsection{Norms of actual point maps over the unbounded nilpotent tower}

Fix a prime $p\geq5$, a prolongation of the $p$-adic valuation to
$\overline{\mathbb Q}$, and write
\[
 K_p=\operatorname{Frac}(\mathbb Z_{(p)}^{\rm sh}),\quad
 I=G_{K_p},\quad N=\bigcap_{c\geq1}\gamma_{c+1}(I),\quad
 K=\overline{\mathbb Q}^{N}=K_{p,\rm nil}.
\]
The lower-central terms are closed.  The quotient $I/N$ is the
maximal pronilpotent quotient.  We keep a fixed complex embedding
$\overline{\mathbb Q}\hookrightarrow\mathbb C$ as well; it is
used only after the arithmetic Kummer statements have been proved.
The $p$-adic and complex absolute values are never identified.
For a field $E$ containing $K$ put
$\kappa_E:E^\times\to H^1(G_E,\mathbb Q_p(1))$.

\begin{claim}{B.106}{Proved}{The three arithmetic inputs at every finite extension}
For every finite extension $E/K$:
\begin{enumerate}
\item $\ker\kappa_E=\mu_\infty$.
\item If $b\in E^\times$, $\alpha\in\mathbb Q_p$, and
      $\kappa_E(b)=\alpha\kappa_E(p)$, then $\alpha\in\mathbb Q$.
\item For every finite extension $E'/E$, complex conjugation in the
      chosen complex embedding does not centralize $G_{E'}$.
\end{enumerate}
These assertions re-prove the inputs A1-B.57--B.58 used below; they
are not deductions from a reconstruction theorem.
\end{claim}
\begin{proof}
The inertia group has a normal wild pro-$p$ subgroup and tame quotient
$\prod_{\ell\ne p}\mathbb Z_\ell(1)$.  In a finite nilpotent
quotient its Sylow subgroups commute.  Consequently
\[
 Q=I/N=Q_{(p)}\times T_{p'},\qquad
 T_{p'}=\prod_{\ell\ne p}\mathbb Z_\ell(1).
\]
The entire tame quotient survives because it is abelian.  The
cyclotomic character factors through $Q$, and on $T_{p'}$ its image
is $\mu_{p-1}\subset\mathbb Z_p^\times$.  All roots of unity lie
in $K$, since their action is already abelian.

Suppose $\kappa_E(a)=0$.  Since $G_E$ fixes all roots of unity,
integral Kummer cohomology is a group of homomorphisms to
$\mathbb Z_p(1)$ and is torsion free.  Rational vanishing therefore
means that all compatible $p$-power roots of $a$ belong to $E$.
Choose a finite Galois number field $D/\mathbb Q$ containing $a$
and algebraic generators of $E/K$, and replace $E$ by $E'=KD$.
This field is stable under the local decomposition group.  Let
$M/\mathbb Q_p$ be the completion of $D$ at the chosen place, made
Galois by the choice of $D$.  Let $a_1,\ldots,a_s$ be the local
Galois conjugates of $a$, and put
\[
 F=K_pD(\mu_{p^\infty}),\qquad
 R=F(a_j^{1/p^m}:1\leq j\leq s,\ m\geq1).
\]
These are Galois extensions of $K_p$, with $R\subset E'$.
The group $D_R=\operatorname{Gal}(R/F)$ embeds in
$\mathbb Z_p(1)^s$ by Kummer theory and is torsion free.  The image
$C_R$ of $N$ in $\operatorname{Gal}(R/K_p)$ is finite, since
$E'/K$ is finite.

Choose $t\in T_{p'}$ with cyclotomic value a primitive
$(p-1)$st root $\zeta$, and a lift $\tau\in I$ of $(1,t)$.
It is central modulo $N$, so its commutator with $D_R$ lies in
$C_R\cap D_R=1$.  Put $J=I_M=I\cap G_D$ and
$J_0=J\cap\ker\chi_p$.  The group $J_0$ maps onto $D_R$.
Conjugation of Kummer roots gives
\[
 \kappa_{J_0}(\tau a)=\zeta\kappa_{J_0}(a).                 \tag{B.NORM.1}
\]
Restriction from $H^1(J,\mathbb Q_p(1))$ to
$H^1(J_0,\mathbb Q_p(1))$ is injective.  Indeed $J/J_0$ is
abelian with nontrivial scalar cyclotomic character, and a cocycle
on an abelian group with such a character is a coboundary: comparing
$c(gh)$ and $c(hg)$ gives
\[
 (\chi(g)-1)c(h)=(\chi(h)-1)c(g).
\]
Choosing $g$ with $\chi(g)\ne1$ proves the assertion.
Restriction from $G_M$ to $I_M$ is injective as well, since
$\mathbb Q_p(1)^{I_M}=0$.  Thus (B.NORM.1) holds in the rational
Kummer group of the finite local field $M$.

Applying the valuation gives $v(a)=\zeta v(a)$, hence $v(a)=0$.
Applying the principal-unit logarithm gives
$\log(\tau a)=\zeta\log(a)$.  The Baker--Brumer theorem says
that $p$-adic logarithms of algebraic units that are linearly
independent over $\mathbb Q$ remain independent over
$\overline{\mathbb Q}$; the complete primary proof opened here
is Rousseau (1969), the unnumbered theorem on printed pages
11-01--11-02.  If $\log(a)\ne0$, this theorem first makes the two
logarithms rationally dependent and then makes their ratio
$\zeta$ rational.  This is impossible for $p\geq5$.
Therefore $\log(a)=0$.  The logarithm has torsion kernel on the
unit group of the finite local field, so $a$ is a root of unity.
The converse is immediate.  This proves (1), with every finite
extension included.

For (2), make the same normal enlargement and write $N'=G_{E'}$,
which is normal in $I$ and finite index in $N$.  Conjugating the
assumed equality on $N'$ by $g\in I$ gives
$\kappa_{E'}(g(b)/b)=0$, because the Kummer class of the rational
number $p$ is invariant.  By (1), each ratio is a root of unity.
The finitely many local conjugates lie in a number field, with a
finite group of roots of unity.  One integer $m>0$ therefore gives
$b^m\in K_p$.

The class $\kappa_I(b^m)-m\alpha\kappa_I(p)$ restricts to zero
on $N'$.  But
$H^1(I/N',\mathbb Q_p(1))=0$.  To see this, its finite normal
subgroup $N/N'$ has zero rational first cohomology and fixes the
coefficient module, so inflation--restriction reduces to $Q$.
The central tame element with cyclotomic value $\zeta\ne1$ makes
every cocycle on $Q$ a coboundary by the preceding displayed identity.
Thus the class was zero already on $I$.  The element $b^m\in K_p$
lies in a finite unramified local extension $M_0/\mathbb Q_p$.
The local restriction to inertia is injective as above.  Applying
its valuation, normalized by $v(p)=1$, gives
$m\alpha=v(b^m)\in\mathbb Z$.  This proves (2).

For (3), first $N\ne1$: the splitting field of $X^p-p$ over
$K_p$ has totally ramified Galois group
$C_p\rtimes\mathbb F_p^\times$, with faithful scalar action,
whose nilpotent residual is nonzero.  The radical degree is $p$
by its value, and the cyclotomic tame degree is $p-1$, so this
calculation does give the full indicated finite quotient of $I$.
The group $N$ lies in wild pro-$p$ inertia, since the tame quotient
is abelian.  It has no torsion, as follows directly from Kummer
theory.  A nonidentity finite-order element would have a power of
prime order $q$; let $H$ be the cyclic subgroup that power generates
and $F=\overline{\mathbb Q}^{H}$.  Then $\mu_\infty\subset F$
and $[\overline{\mathbb Q}:F]=q$.  Kummer theory gives a generator
$\beta$ with $\beta^q=a\in F\setminus F^q$.
Choose $\gamma\in\overline{\mathbb Q}$ with $\gamma^q=\beta$.
Taking the degree-$q$ norm gives
\[
 \operatorname{N}(\gamma)^q
   =\operatorname{N}(\beta)=(-1)^{q-1}a.
\]
The sign is a $q$th power in $F$, since all roots of unity belong
to $F$, contradicting $a\notin F^q$.  Thus $N$ is infinite and
every finite-index subgroup $G_{E'}$ is nontrivial.

The field $E'$ has its nontrivial rank-one henselian valuation.  If
complex conjugation $c$ centralized $G_{E'}$, it would preserve
$E'$.  The transported valuation would be another rank-one henselian
valuation.  If inequivalent, the two are independent, and the
F.~K.~Schmidt theorem (Kuhlmann, Theorem~9.44) would make $E'$
separably closed, a contradiction.  Thus $c$ preserves the
valuation, and uniqueness of prolongation puts $c$ in the local
decomposition group.  Its image in the unramified quotient
$\widehat{\mathbb Z}$ is trivial, and then its image in the
torsion-free tame inertia quotient is trivial.  It would be an
element of order two in wild pro-$p$ inertia, impossible for the
odd prime $p$.  This proves (3).
\end{proof}
\paragraph{Check record.} checked only in this writing, 2026-11-18: GPT-6 Astra (OpenAI), ChatGPT Pro mode; internal effort value not exposed. Own detailed proof written and reread; primary source hypotheses matched. B8.1--2 and B8.4--5 read in full by root; B8.3 separately checked by configuration agent with global-continuation expansion


\begin{claim}{B.107}{Proved}{Rational residues and norm moduli on affine space}
Let $E/K$ be finite.  A class
$\xi\in H^1(C,\mathbb Z_p(1))$ on a smooth curve over $E$ whose
evaluation at every point over every finite extension is an actual
Kummer value has rational geometric boundary residues.

More generally, let $V\subset\mathbb A^n_E$ be the complement
of a hypersurface, and let
$\eta\in H^1(V,\mathbb Z_p(1))$ have that same actual-value
property.  After a finite extension splitting the irreducible
geometric boundary hypersurfaces, with equations $h_1,\ldots,h_s$,
there are $q_j\in\mathbb Q\cap\mathbb Z_p$ and an arithmetic
class $\beta$ such that
\[
 \eta=\beta+\sum_jq_j\kappa(h_j)
       \quad\hbox{in }H^1(V,\mathbb Q_p(1)).               \tag{B.NORM.2}
\]
If $b(t)$ is any actual value of $\eta$ at
$t\in V(\overline{\mathbb Q})$, then in the fixed complex
embedding
\[
 |b(t)|^2=C\prod_j|h_j(t)|^{2q_j},\qquad C>0,\quad
 C\in\overline{\mathbb Q}.                              \tag{B.NORM.3}
\]
The constant is independent of $t$; a finite extension of point
definition does not change it.
\end{claim}
\begin{proof}
Split the curve boundary and fix a point $q$.  Riemann--Roch gives
a rational function $g$ with sole pole at $q$ and with zeros at
every other boundary point.  Delete zero, infinity and the finitely
many branch values from $\mathbb P^1$.  On the inverse image the
map is finite \emph{\'etale}; corestrict $\xi$ to the resulting
punctured rational line.  Its residues at removed branch values are
zero.  Its residue at infinity is exactly $\operatorname{res}_q\xi$.
The norm formula for a discrete valuation multiplies residues by
residue-field degrees, not by ramification indices; after splitting
the points these degrees are one.  This follows first from
$v(Na)=\sum f_w w(a)$ at finite Kummer levels and then in the
inverse limit.  Hence the corestricted class is an arithmetic class
plus $\alpha\kappa(T)$, with
$\alpha=-\operatorname{res}_q\xi$.

Corestriction evaluates as the product of the field norms of actual
point values, so its evaluations are still actual values.  Choose
$t_0$ and $p^m t_0$ in the punctured line, with $m>0$; there are
only finitely many excluded choices.  Subtracting evaluations gives
an actual value $b$ with
$\kappa_E(b)=m\alpha\kappa_E(p)$.  Claim~B.106(2) makes
$\alpha$ rational.  This independently proves the boundary argument
of A1-B.60.

For $V$ the geometric unit group modulo constants is freely
generated by the $h_j$, and its Picard group is zero because its
coordinate ring is a localization of a polynomial UFD.  The Kummer
sequence and the inverse limit therefore give (B.NORM.2) initially
with coefficients in $\mathbb Z_p$.  At a smooth geometric point
of a chosen hypersurface away from the others, take a transverse
line defined over a finite extension.  The restricted class has
the actual-value property, and its residue is exactly the selected
coefficient, with intersection multiplicity one.  The curve result
makes every coefficient rational.  Subtracting the indicated Kummer
classes leaves an arithmetic class, either directly by the same
unit calculation or by the kernel part of Hochschild--Serre.

Choose one algebraic basepoint $t_0\in V$ and a positive integer
$m$ clearing all the $q_j$.  Comparing evaluations at $t$ and
$t_0$ and applying Claim~B.106(1) over a common finite extension
gives
\[
 \left(\frac{b(t)}{b(t_0)}\right)^m
       \prod_j\left(\frac{h_j(t_0)}{h_j(t)}\right)^{mq_j}
       \in\mu_\infty.
\]
Its complex absolute value is one.  Taking the positive $m$th root
proves (B.NORM.3), with
$C=|b(t_0)|^2\prod_j|h_j(t_0)|^{-2q_j}$.
This is positive algebraic because all the numbers in the formula
are algebraic and all the exponents are rational.  The formula
uses one fixed $t_0$; it establishes uniformity of the constant,
rather than choosing a separate constant at each point.
\end{proof}
\paragraph{Check record.} checked only in this writing, 2026-11-18: GPT-6 Astra (OpenAI), ChatGPT Pro mode; internal effort value not exposed. Own detailed proof written and reread; primary source hypotheses matched. B8.1--2 and B8.4--5 read in full by root; B8.3 separately checked by configuration agent with global-continuation expansion


\begin{claim}{B.108}{Proved}{Reconstructing a norm polynomial without assuming continuity}
Let $X/E$ be smooth geometrically connected affine of dimension
$n\geq1$, and let
$\phi:X_{\rm et}\to\mathbb A^1_{E,\rm et}$ be a pinned geometric
morphism over $E_{\rm et}$.  Write
$f:X(\overline{\mathbb Q})\to\overline{\mathbb Q}$ for its actual
point map.  Suppose it is nonconstant.  For any finite generically
\emph{\'etale} morphism $g:X\to\mathbb A^n_E$ of degree $d$,
the polynomials
\[
 P_t(Z)=\prod_{x\in g^{-1}(t)}(Z-f(x))                    \tag{B.NORM.4}
\]
on its geometric \emph{\'etale} fibres agree, on a nonempty
Zariski open set of algebraic $t$, with a monic polynomial
$P_g\in\overline{\mathbb Q}(t_1,\ldots,t_n)[Z]$ of degree $d$.
No continuity or analyticity of $f$ is assumed.
\end{claim}
\begin{proof}
Relative pinning gives an actual $G_E$-equivariant point map,
including $f(x)\in E(x)$ over every finite field of definition.
This is Carlson--Haine--Wolf, Proposition~3.6, whose hypotheses
are a perfect base field, an algebraic extension of it, and a
target topologically of finite type.  Pulling back the open
complement of any $a\in\overline{\mathbb Q}$ shows that
$S_a=f^{-1}(a)$ is a closed subset after a finite scalar extension.
It is proper when $f$ is nonconstant: if the generic point mapped
to a closed point, continuity and specialization would make the
entire point map constant.  Since $g$ is finite, $g(S_a)$ is a
proper closed subset of $\mathbb A^n$.

Fix $a$.  Choose a principal open $V_a\subset\mathbb A^n$
avoiding $g(S_a)$ and the branch locus, and corestrict the pullback
of $\kappa(Z-a)$ along
$g^{-1}(V_a)\to V_a$.  Its evaluation is the actual value
\[
 b_a(t)=\prod_{x\in g^{-1}(t)}(f(x)-a)=(-1)^dP_t(a).
\]
Indeed its factors lie in the residue fields of the fibre points;
the product is their field norm and belongs to the field of $t$.
Claim~B.107 gives
\[
 |P_t(a)|^2=R_a(t),\qquad
 R_a(t)=C_a\prod_j|h_{a,j}(t)|^{2q_{a,j}},                \tag{B.NORM.5}
\]
where $C_a>0$ is algebraic and uniform in $t$, and all $q_{a,j}$
are rational.  The right side is positive, real analytic, and
algebraic in the independent variables $t,\overline t$ off a
finite union of algebraic hypersurfaces.  Its logarithm is
pluriharmonic there.

We now recover the coefficients from finitely many such equations.
For a monic polynomial $P(Z)=\sum_{k=0}^d c_kZ^k$, $c_d=1$,
the squared modulus $|P(r+is)|^2$ is a polynomial in the real
variables $r,s$, of degree at most $2d$ in each.  Its values on
the grid $0\leq r,s\leq2d$ determine this polynomial by
two-variable Lagrange interpolation.  On changing variables to
$Z=r+is$, $W=r-is$, the coefficient of $Z^kW^d$ is exactly
$c_k$.  Thus each coefficient is a fixed algebraic linear
combination of the $(2d+1)^2$ squared moduli on that grid.
A translated and nonzero algebraically scaled grid has the same
property.  This is a coefficient reconstruction, not a continuity
assumption on the original function.

Choose such a grid $A$, and let $\Omega\subset\mathbb C^n$ be
the complement of the finitely many hypersurfaces needed for its
$R_a$'s.  It is connected.  Applying the interpolation functionals
to these real-analytic functions defines a monic polynomial
\[
 P(t,Z)=Z^d+c_{d-1}(t)Z^{d-1}+\cdots+c_0(t)              \tag{B.NORM.6}
\]
whose coefficients are real analytic on $\Omega$ and algebraic
over $\overline{\mathbb Q}(t,\overline t)$.  At every algebraic
$t\in\Omega$ they equal the coefficients of (B.NORM.4).
For every other algebraic $a$, density of algebraic points and
real-analytic continuation now give
$|P(t,a)|^2=R_a(t)$ wherever both sides are defined.
In particular $\log|P(t,a)|^2$ is pluriharmonic away from its
zeros.  The differential identity extends across the auxiliary
hypersurfaces wherever $P(t,a)\ne0$, because the expression in
the derivatives of $P$ is real analytic there.  Hence this
pluriharmonicity can be used at a fixed $t$ for every algebraic
$a$ except the finitely many roots, and then for every complex
$a$ by density.

Here is the analytic implication, including mixed branches.  For
each $j,k$ put
\[
 Q_{jk}(a)=\frac{P_{t_j\overline t_k}(t,a)}{P(t,a)}
          -\frac{P_{t_j}(t,a)P_{\overline t_k}(t,a)}{P(t,a)^2}.
\]
At a fixed $t$ these are rational holomorphic functions of $a$,
and pluriharmonicity reads
$Q_{jk}(a)+\overline{Q_{kj}(a)}=0$.  A holomorphic rational
function equal to an antiholomorphic one on an open set is
constant.  Since the polynomial is monic, all its differentiated
coefficients have lower $a$-degree, and $Q_{jk}(a)\to0$ as
$a\to\infty$.  Consequently every $Q_{jk}$ is zero.

\paragraph{Local algebraic separation of the root types.}
At a generic real-analytic germ, complexify $\bar t$ to independent
variables $w$. The coefficients of $P$ belong to a finite algebraic
extension of $\overline{\mathbf Q}(t,w)$. In a finite splitting
field the derivations $\partial_{t_j}$ and $\partial_{w_k}$ extend
uniquely, since the characteristic is zero. Decompose $P$ into its
distinct roots $\lambda_\nu$ with positive multiplicities $m_\nu$.
The double and simple partial-fraction poles of
\[
 \partial_{t_j}\partial_{w_k}\log P(t,w,a)=0
\]
give respectively
\[
 m_\nu(\partial_{t_j}\lambda_\nu)
        (\partial_{w_k}\lambda_\nu)=0,
 \qquad
 m_\nu\partial_{t_j}\partial_{w_k}\lambda_\nu=0.
\]
If one $t$-derivative of a root is nonzero, every $w$-derivative
is zero; if one $w$-derivative is nonzero, every $t$-derivative
is zero. This is an identity in a field, so no point-dependent
choice of type is involved. A root with both sets zero is constant.

A root killed by all $w$-derivations is algebraic over
$\overline{\mathbf Q}(t)$. Here is the elementary differential
algebra argument. Take its monic minimal polynomial over
$\overline{\mathbf Q}(t,w)$. Derivations commute with the
automorphisms of a separable normal closure, so all its conjugates
are also killed. Its coefficients, the elementary symmetric
functions in those conjugates, belong to the simultaneous constant
field of the $w$-derivations, namely $\overline{\mathbf Q}(t)$.
The same argument applies with $t$ and $w$ exchanged. A root killed
by both sets is algebraic over $\overline{\mathbf Q}$ and hence
belongs to that algebraically closed field.

We have therefore obtained monic algebraic coefficient germs
\[
 P(t,Z)=H_0(t,Z)A_0(\bar t,Z),
\]
where the roots of $H_0$ are holomorphic or constant, and those of
$A_0$ are antiholomorphic and nonconstant. Repeated roots stay with
their full multiplicities. A holomorphic nonconstant algebraic
germ cannot equal an antiholomorphic one: complexifying such an
equality makes all derivatives in both independent variable sets
zero. This observation gives uniqueness of the division into the
two root types, with constants assigned to $H_0$.

\paragraph{Global continuation, including collisions of opposite types.}
Remove from $\Omega$ a complex algebraic hypersurface containing
the poles and branch loci of the finite algebraic coefficient
extensions of $H_0(t,Z)$ and of the conjugate of
$A_0(\bar t,Z)$. Call the resulting connected complex Zariski
open set $\Omega'$. Each coefficient germ can now be continued
along every path in $\Omega'$. Their product continues to $P$,
by uniqueness of continuation of real-analytic germs. This does
not require the roots of $H_0$ to avoid the roots of $A_0$ on the
real diagonal; continuation of the algebraic coefficients of
the two factors is defined independently of those collisions.

After continuation around any loop, let the resulting monic
factors be $H_1(t,Z)$ and $A_1(\bar t,Z)$. They satisfy
\[
 H_1(t,Z)A_1(\bar t,Z)=H_0(t,Z)A_0(\bar t,Z).
\]
Continuation preserves the vanishing of the relevant derivations
and cannot turn a nonconstant algebraic germ into a constant one.
Uniqueness of the two root types, just proved, forces $H_1=H_0$
and $A_1=A_0$. Thus every coefficient of each factor is single
valued on $\Omega'$; the first is holomorphic there, and the
second becomes holomorphic after conjugating the coordinates.

A single-valued holomorphic function on $\Omega'$ algebraic over
$\mathbf C(t)$ belongs to $\mathbf C(t)$. To see this, normalize
the base in its irreducible function-field extension and remove
the remaining branch and pole locus. This gives a connected
finite \'etale cover. Its complex analytic space is connected:
an irreducible complex variety remains analytically connected
after deleting a proper algebraic subset. A globally defined
branch is a section of this covering and has open-and-closed
image, so the connected cover has degree one. The same reasoning
applies to the antiholomorphic coefficients after conjugation.

Finally a rational function in $\mathbf C(t)$ algebraic over
$\overline{\mathbf Q}(t)$ belongs to
$\overline{\mathbf Q}(t)$. One concrete descent argument is to
bound numerator and denominator degrees and choose sufficiently
many algebraic parameter points avoiding their poles and the
degenerate specializations of the algebraic equation. The values
are algebraic. Rational interpolation, normalized by one nonzero
coefficient, then recovers the function by a finite linear system
with algebraic entries. Normalize one denominator coefficient
known to be nonzero. A solution exists over $\mathbf C$, hence
also over $\overline{\mathbf Q}$ by Gaussian elimination.
Sufficiently many interpolation points determine the rational
function uniquely (cross-multiply and interpolate a polynomial
of twice the degree bound), so this algebraic solution represents
the original function. Consequently
\[
 \begin{gathered}
 P(t,Z)=H(t,Z)A(\overline t,Z),\\
 H\in\overline{\mathbf Q}(t)[Z],\qquad
 A\in\overline{\mathbf Q}(\bar t)[Z].
 \end{gathered}\tag{B.NORM.8}
\]

It remains to exclude a nonconstant $A$.  Choose a finite number
field $M$, stable under complex conjugation, containing the
coefficients of $H,A,g,X$ and all auxiliary hypersurfaces, and
work with $\sigma\in G_{EM}$.  Equivariance of the actual point
map gives $P_{\sigma t}=\sigma(P_t)$ at every algebraic parameter
where defined.  Cancelling the common monic factor $H(\sigma t,Z)$
in (B.NORM.8) gives
\[
 A(c\sigma t,Z)=A(\sigma c t,Z),                        \tag{B.NORM.9}
\]
where $c$ denotes coordinatewise complex conjugation.  If a
coefficient $a(w)$ of $A$ is nonconstant, put
$\tau=c\sigma c\sigma^{-1}$.  This automorphism fixes $M$, and
(B.NORM.9) says $a(\tau u)=a(u)$ on a nonempty Zariski open set
of algebraic $u$.  Since $\tau$ fixes the coefficients,
$\tau(a(u))=a(u)$.  A nonconstant rational function on affine
space has cofinite image in the algebraically closed field
$\overline{\mathbb Q}$ (restrict to a suitable line if necessary).
Thus $\tau$ fixes a cofinite subset, and then fixes every element
by subtraction of two elements in that subset.  Hence $\tau=1$.
This would hold for every $\sigma\in G_{EM}$, contradicting
Claim~B.106(3).  All coefficients of $A$ are therefore constant;
by the convention on its roots it has degree zero.  The polynomial
$P=H$ is rational in $t$, as required.
\end{proof}
\paragraph{Check record.} checked only in this writing, 2026-11-18: GPT-6 Astra (OpenAI), ChatGPT Pro mode; internal effort value not exposed. Own detailed proof written and reread; primary source hypotheses matched. B8.1--2 and B8.4--5 read in full by root; B8.3 separately checked by configuration agent with global-continuation expansion


\begin{claim}{B.109}{Proved}{Generic norms force one algebraic graph}
In the setting of Claim~B.108 there is a unique
$h\in E(X)$ such that $f(x)=h(x)$ at every algebraic point of
some nonempty Zariski open subset of $X$.  The rational function
is regular everywhere and agrees with $f$ at every algebraic
point.
\end{claim}
\begin{proof}
Choose one finite Noether normalization $g$.  Claim~B.108 gives
a monic rational equation $P_g(g(x),f(x))=0$ on a dense open.
After shrinking to a smooth affine $X_0$, all its coefficients are
regular.  The Zariski closure $Z$ of the graph of $f$ over $X_0$
is therefore contained in a scheme finite over $X_0$.
This proves algebraicity of the containing correspondence before
any assertion about the graph itself.  Let $Z_1,\ldots,Z_r$ be
the irreducible components dominating $X_0$, and let their degrees
over $X_0$ be $e_i\geq1$.  There is at least one such component,
since the graph has a point over every algebraic point of $X_0$.
The actual graph is dense in each $Z_i$ after the other components
have been removed: otherwise its closure would not have that
irreducible component.  Since $Z_i$ dominates $X_0$, these graph
points project to a dense subset of $X_0$.  Thus the coordinate
$y$ is nonconstant on each $Z_i$: a component with $y=a$ would
make $S_a=f^{-1}(a)$ dense in $X_0$, whereas $S_a$ is a proper
closed subset for a nonconstant pinned point map.

We choose a second finite normalization
$g':X_0\to\mathbb A^n$ with the following additional property:
each $(g',y):Z_i\to\mathbb A^n\times\mathbb A^1$ is birational
onto its image, and these images are pairwise distinct.  Here is a
generic-projection proof.  Embed $X_0$ in affine space.  Over the
function field $\overline{\mathbb Q}(y)$ the union of the generic
fibres of the $Z_i$ is a reduced union of $(n-1)$-dimensional
subvarieties of that affine space: a fibre of
$Z\subset X_0\times\mathbb A^1$ embeds in $X_0$ when $y$ is
fixed.  A generic linear projection to $\mathbb A^n$ is
generically one-to-one on that union.  The incidence proof counts
pairs of distinct points, of dimension $2(n-1)$, and the $n$
independent conditions that their images coincide; the resulting
double locus has dimension at most $n-2$, less than that of a
component.  The characteristic is zero, so generic degree one is
birational.  The same estimate separates different components.
For $n=1$ the generic fibres are finite reduced sets, and a generic
linear function separates them directly.  These dimension estimates
are on the fibres over the generic value of $y$.  The conditions
first give a nonempty open subset of coefficient space over
$\overline{\mathbb Q}(y)$.  It contains an open defined by
finitely many nonvanishing polynomials in the coefficients whose
coefficients are rational in $y$.  After clearing denominators
and expanding in $y$, failure of one such nonvanishing identically
in $y$ is the simultaneous vanishing of finitely many polynomials
over $\overline{\mathbb Q}$.  Its complement is a nonempty open
in coefficient space over $\overline{\mathbb Q}$.  Thus these
conditions can be imposed simultaneously with the nonempty open
condition of finite Noether normalization of $X_0$.  They have a
choice over a finite extension of $E$, as permitted throughout.

Write $d'=[\overline{\mathbb Q}(X_0):
\overline{\mathbb Q}(g')]$.  Birationality of $(g',y)$ says
\[
 \overline{\mathbb Q}(Z_i)=
 \overline{\mathbb Q}(g',y).
\]
The minimal polynomial of $y$ over
$\overline{\mathbb Q}(g')$ therefore has degree $e_i d'$.
These minimal polynomials are distinct for different components.
On the other hand Claim~B.108 supplies a monic norm polynomial
$P_{g'}$ of degree $d'$ vanishing on the actual graph, hence on
all its dominating closure components.  Their distinct minimal
polynomials divide it.  Thus
\[
                    d'\sum_i e_i\leq d',
\]
so $r=1$ and $e_1=1$.  The graph agrees with one rational function
$h$ off a proper closed subset.  Uniqueness on a dense open and
the $G_E$-equivariance of $f$ imply $\sigma h=h$ for every
$\sigma\in G_E$, hence $h\in E(X)$.

We now return to the original smooth affine $X$ and remove both
possible poles and exceptional point values; agreement on a dense
set alone would not do this.  Enlarge the exceptional subset to a
proper closed $B\subset X$ that contains $X\setminus X_0$ and
outside which agreement and regularity hold.  Fix any algebraic
point $x_0\in B$.  After a finite scalar extension take a finite
linear normalization $g':X\to\mathbb A^n$ of this original $X$
such that
$t_0=g'(x_0)$ is an \emph{\'etale} fibre and
\[
                g'^{-1}(t_0)\cap B=\{x_0\}.
\]
Such projections exist after enlarging the affine embedding by
polynomial functions of degree at most a fixed sufficiently large
integer.  These functions separate values and first jets at any
two distinct points.  Impose $g'(x_0)=0$.  At each other point
$x$, the conditions $g'(x)=0$ have codimension $n$ in the
parameter space of the $n$ functions; requiring also a singular
differential has one further independent determinantal condition,
because first jets can be prescribed.  The incidence over the
$n$-dimensional $X\setminus\{x_0\}$ therefore has image of
positive codimension.  The incidence over $B\setminus\{x_0\}$,
with just the value conditions, also has image of positive
codimension since $\dim B<n$.  At $x_0$ invertibility of the
differential is a nonempty open condition.  The highest-degree
forms can simultaneously be chosen with no common zero on the
projective boundary, which is the nonempty open condition making
the map finite.  Thus a parameter outside these finitely many
proper closed subsets gives exactly the asserted fibre.  This
also supplies a direct jet-and-incidence proof of the required
smooth-fibre Bertini step.

Choose the interpolation grid in Claim~B.108 to avoid the finitely
many actual values of $f$ on this fibre.  For each grid value $a$,
the closed set $g'(S_a)$ does not contain $t_0$, so the principal
open $V_a$ used in that proof can be chosen to contain $t_0$.
The reconstructed coefficients are then real analytic near $t_0$,
are rational by Claim~B.108, and have no pole at $t_0$.
Their polynomial is equal generically to the ordinary norm
characteristic polynomial of $h$, since $f=h$ off $B$.
Consequently $h$ satisfies a monic equation with coefficients
regular at $t_0$.  The local rings of $X$ are normal, so integrality
makes $h$ regular at every point of that fibre, including $x_0$.

At $t_0$ the polynomial equals the product of $(Z-f(x))$ by the
finite-grid reconstruction, and equals the product of $(Z-h(x))$
by specialization of the finite \emph{\'etale} norm.  Every factor
except the one at $x_0$ is already the same.  Cancelling them gives
$f(x_0)=h(x_0)$.  Since $x_0$ was arbitrary, this proves both
global regularity and equality at all geometric points.
\end{proof}
\paragraph{Check record.} checked only in this writing, 2026-11-18: GPT-6 Astra (OpenAI), ChatGPT Pro mode; internal effort value not exposed. Own detailed proof written and reread; primary source hypotheses matched. B8.1--2 and B8.4--5 read in full by root; B8.3 separately checked by configuration agent with global-continuation expansion


\begin{claim}{B.110}{Proved}{Pinned reconstruction over the nilpotent union for $p\geq5$}
For every $p\geq5$ and every finite extension $E/K_{p,\rm nil}$,
the functor
\[
 \operatorname{Sch}^{\rm ft}_E[\operatorname{UH}^{-1}]
       \longrightarrow\operatorname{RTop}^{\rm pin}_E,
 \qquad X\longmapsto X_{\rm et}
\]
is fully faithful.  Here the target uses geometric morphisms over
$E_{\rm et}$, with their relative structural isomorphisms, that
take closed points to closed points.  Equivalently, for finite-type
$E$-schemes with seminormal source, pinned relative geometric
morphisms of their \emph{\'etale} topoi come from unique morphisms
of $E$-schemes.
\end{claim}
\begin{proof}
Claims~B.108--B.109 construct an actual morphism
$X\to\mathbb A^1$ agreeing with any pinned point map when $X$
is smooth geometrically connected affine.  A constant point map is
the constant morphism; if needed a rational point ensures its value
lies in the base field.  For a smooth connected source that is not
geometrically connected, pass to a finite extension splitting its
geometric components, apply the construction to each, and descend
the unique resulting morphism using equality of its geometric
point values.  Every arithmetic assertion used above is stable
under those finite extensions by Claim~B.106.

We apply the exact reduction theorem of Carlson--Haine--Wolf
(2024), Theorem~4.18, implication (4.18.4)$\Rightarrow$(4.18.1).
It holds for \emph{any} field, without their later amenability
hypothesis.  It requires precisely an actual morphism agreeing on
geometric points for each affine regular connected source with a
rational point and target $\mathbb A^1$.  In characteristic zero
regular finite-type schemes are smooth, and absolute weak
normalization of a regular scheme is itself.  Thus the requirement
has been met.  Their proof uses h-descent and reconstructs ring
operations from geometric-point evaluations; in particular it
identifies the entire inverse-image functor, not merely the point
map.  Finally their Corollary~2.32 identifies the localization at
universal homeomorphisms with the stated seminormal-source
formulation in characteristic zero.  This closes the
arbitrary-source and pointwise-trivial-functor gaps left by
A1-B.59--B.61 for these fields.
\end{proof}
\paragraph{Check record.} checked only in this writing, 2026-11-18: GPT-6 Astra (OpenAI), ChatGPT Pro mode; internal effort value not exposed. Own detailed proof written and reread; primary source hypotheses matched. B8.1--2 and B8.4--5 read in full by root; B8.3 separately checked by configuration agent with global-continuation expansion


\paragraph{What this changes, and what it does not decide.}
The new step is the use of all finite norm polynomials and a generic
projection to eliminate a multivalued algebraic correspondence.
It handles the proper-cohomology component on an arbitrary source
by corestriction to a principal open in affine space, whose Picard
group is zero.  It does not assume that this component vanished
on the original curve.  The finite interpolation script checks
the coefficient reconstruction for sample degrees only; all
degrees and the profinite arithmetic claims are proved above.

The remaining B8 list is not empty: other mixed-characteristic
henselian fields not covered by the proved positive and negative
classes still require further work.
The maximal metabelian fields and every finite extension are
settled subsequently by Claim~B.136; B.138 gives a uniform
all-prime proof for the nilpotent unions, including $p=2,3$.
Those later proofs do not use the torsion-only kernel needed
in the present argument.  Here we recheck the two small-prime
counterexamples to that arithmetic input.
For $p=2$, the splitting field of $X^{2^m}-2$ over $K_2$ has
a $2$-group as Galois group: the cyclotomic and radical degrees
are both powers of two.  It is nilpotent, so every such root
belongs to $K_{2,\rm nil}$.  Thus the nontorsion element $2$
has zero rational $2$-Kummer class there.

For $p=3$, set $z=\sqrt{-3}$ and $a=(2+z)/(2-z)$.
The quadratic extension $F_0=K_3(z)=K_3(\mu_3)$ sends $a$
to $a^{-1}$.  The extension
$R_m=F_0(\mu_{3^m},b)$, $b^{3^m}=a$, is Galois over $K_3$
with a normal $3$-subgroup $P$ and quotient $C_2$.
Cauchy's theorem supplies an involution $s$; it inverts
$\mu_{3^m}$ and has $s(b)=\xi b^{-1}$.  Its square gives
$\xi^2=1$, hence $\xi=1$ since $\xi\in\mu_{3^m}$.
If $g\in P$ has $g(b)=\eta b$ and
$g(\zeta_{3^m})=\zeta_{3^m}^{u}$, direct substitution gives
$sgs^{-1}(b)=\eta b$ and
$sgs^{-1}(\zeta_{3^m})=\zeta_{3^m}^{u}$.
It follows that $s$ centralizes $P$ and
$\operatorname{Gal}(R_m/K_3)=P\times C_2$ is nilpotent.
Thus every $3$-power root of $a$ belongs to $K_{3,\rm nil}$.
The equality $a+a^{-1}=2/7$ proves that $a$ is not a root
of unity, since the sum of a root of unity and its inverse is
an algebraic integer.  These calculations re-prove the two
exceptions in A1-B.57.  They refute input~(1) at the two small
primes, not reconstruction itself.  No such field is promoted by
removing $p\geq5$ from the proof.

The current Carlson--Stix revision, arXiv:2410.22474v2 of
20 September 2026, was also opened.  Its Theorem~A/\allowbreak Theorem~4.1
assumes a sub-$p$-adic field.  Its new relative
Theorem~D/\allowbreak Theorem~6.12 assumes a geometrically unibranch
noetherian Nagata base whose component function fields are
sub-$p$-adic.  Neither directly supplies the conclusion above.
Indeed $K_{p,\rm nil}$ contains $\mu_\infty$, whereas the
algebraic constant field of a finitely generated extension of a
$p$-adic local field has only finitely many roots of unity.
Its point-theoretic open-or-constant Proposition~3.17 requires
Kummer faithfulness, which fails here on $\mu_\infty$.

\paragraph{Checks and provenance.}
These are new proofs in this continuation, written by GPT-6 Astra
(OpenAI), ChatGPT Pro mode, on 18 November 2026.  They make no
claim of priority or human verification.  The Kummer inputs were
re-proved from the local inertia and logarithm arguments; the
norm/interpolation and generic-projection steps are separate from
the earlier affine-line proof.  A second reading, and its exact
scope, must be recorded before these working statements are
treated as checked in the final register.

% AI-produced by GPT-6 Astra (OpenAI), ChatGPT Pro mode, 18 November 2026.
% Working identifiers. No priority or human/external verification is claimed.
\subsection{Bare normalizers, finite-cover detection, and the CM residue}

Throughout this subsection $U/k$ is a smooth geometrically connected
affine hyperbolic curve over a number field.  Write
\[
 F=\pi_1(U_{\overline k}),\qquad E=\pi_1(U),\qquad
 p:E\longrightarrow G_k,\qquad
 R=\bigcap_\ell\ker(F\longrightarrow F^{(\ell)}).
\]
Groups and subgroups are profinite and closed, and maps are continuous.
The notation $F^{(\ell)}$ means the maximal pro-$\ell$ quotient of
the \emph{original} geometric group.  The source question is
cote~14, scans~003, 007, 011 and 041.  An arithmetic normalizer has
image $G_L$ for an actual finite extension $L/k$; we retain this
field in every character statement.  It need not equal $k$.

\begin{claim}{B.111}{Proved}{The precise scope and numbering of the cyclotomic test}
Over a number field, a nontrivial closed pro-$\ell$ subgroup of
$F$ is contained in cusp inertia if and only if it is isomorphic
to $\mathbb Z_\ell$ and is essentially cyclotomically normalized.
Here the latter condition means that some subgroup of its arithmetic
normalizer has open image in $G_k$ and acts by $\chi_\ell\circ p$.
This is the opened published theorem of Nakamura as stated by Stix.

In Stix's 2012 version the relevant references are Definition~26,
Definition~29, Theorem~30 and Corollaries~31--32.  In the supplied
2008 source, version~3.05, the corresponding numbered statements
are Lemma~19, Theorem~20 and Corollaries~21--22; the definitions
are unnumbered prose.  The count 22--25 in the prompt is not the
count produced by this source's shared theorem counter.

Every original maximal pro-$\ell$ quotient, and every quotient
of it, kills $R$.  Moreover every cusp inertia group has trivial
intersection with $R$.
\end{claim}
\begin{proof}
Stix (2012), Theorem~30, has precisely the three equivalent
conditions stated above, under the hypothesis that the
$\ell$-adic cyclotomic character is non-Tate.  Definition~26 and
Lemma~28 express that hypothesis by the absence, after finite
extension, of a nonzero map
$\mathbb Q_\ell(1)\to V_\ell A$ for an abelian variety $A$.
For a number field choose a finite place of good reduction away
from $\ell$.  Its Frobenius has complex absolute values $q^{1/2}$
on $V_\ell A$ and $q$ on $\mathbb Q_\ell(1)$, proving the
required absence.  The same argument works after any finite
extension.  The theorem is therefore applicable to the actual
full group here.  Corollary~31 characterizes maximal full
procyclic cusp groups; Corollary~32 is the section formulation.
They do not replace the character hypothesis by bare normalization.

For the numbering, the environments sharing the \texttt{thm}
counter include conjectures and remarks.  We count every actual
beginning of such an environment after stripping comments.
The labels \texttt{lem:Ribet}, \texttt{thm:naka} and
\texttt{cor:anWF} occur at 19, 20 and 21, with the following
cuspidal-section corollary at 22.  The script
\texttt{stix\_numbering\_certificate.py} records the declarations,
all numbered environments and these matches.  The published
content was matched separately, rather than inferred from the
counter calculation.

The blindness of the specified quotients is immediate from the
definition of $R$.  To verify the cusp intersection assertion,
use a standard geometric surface presentation with $2g$ handle
generators and $r-1$ independent cusp generators.  If $r>1$, every
cusp generator has nonzero image in the torsion-free pro-$\ell$
abelianization.  If $r=1$, hyperbolicity and affineness give
$g\geq1$, and the cusp word is the inverse of
$\prod_{j=1}^g[x_j,y_j]$.  Its leading Magnus term is
$-\sum_j[X_j,Y_j]$, which is nonzero over $\mathbb Z_\ell$,
also for $\ell=2$.  In either case the cusp's $\mathbb Z_\ell$
component maps injectively into $F^{(\ell)}$: its nonzero image
is torsion free.  A nontrivial closed subgroup of full cusp
inertia has a nontrivial component at some prime, so cannot lie
in $R$.  Cusp disjointness and the geometric cusp normalizer
used below are Stix (2012), Proposition~16(1)--(3), corresponding
to Proposition~11 and Corollaries~12--13 of the 2008 source.
\end{proof}
\paragraph{Check record.} checked only in this writing, 2026-11-18: GPT-6 Astra (OpenAI), ChatGPT Pro mode; internal effort value not exposed. Own detailed proof written and reread; primary source hypotheses matched. All N2.1--7 read in full by root; finite-cover topology join requested and inserted


\begin{claim}{B.112}{Proved}{A bare character descends after a power and is detected on a cover}
Let $1\ne A\leq F$ be procyclic and let $N\leq E$ normalize
$A$, with $p(N)=G_L$.  For every $\ell$ with $A_\ell\ne1$
write $\theta_\ell:N\to\mathbb Z_\ell^\times$ for its
conjugation character.  After replacing $N$ by $AN$,
put $D=N\cap F$ and $d_\ell=\ell-1$ for odd $\ell$,
$d_2=1$.  There is a character
\[
 \lambda_\ell:G_L\longrightarrow\mathbb Z_\ell^\times,
 \qquad \lambda_\ell\circ p=\theta_\ell^{d_\ell}.
\]
There are a finite extension $L_1/L$ and a finite \'etale cover
$Y\to U_{L_1}$ in whose geometric homology the image of $A_\ell$
is a nonzero line, with character $\psi$, such that
$\lambda_\ell|_{G_{L_1}}=\psi^{d_\ell}$.
If $A$ is noncuspidal, this line maps injectively to
$V_\ell\operatorname{Jac}(\overline Y)$; in particular
$\psi$ has Frobenius weight one and Hodge--Tate weights in
$\{0,1\}$, with cyclotomic Hodge--Tate weight one.
\end{claim}
\begin{proof}
The replacement is legitimate because $A$ is abelian: $AN$ is
compact, normalizes $A$, and has the same arithmetic image.
It supplies the sometimes missing inclusion $A\leq D$.
The affine group $F$ is free profinite, so it is torsion free,
and every closed pro-$\ell$ subgroup has cohomological dimension
at most one and is free pro-$\ell$.  The cohomological facts and
freeness criterion are stated explicitly in Morales, arXiv
2211.12390v3, Section~8.1 and Theorem~8.5 (the latter cites
Ribes--Zalesskii, Theorem~7.7.4); the primary text was opened.

Choose an $\ell$-Sylow subgroup $S$ of $D$ containing the normal
subgroup $A_\ell$.  For each $s\in S$, the pro-$\ell$ group
$\overline{\langle A_\ell,s\rangle}$ is generated by two
elements, is metabelian, and is free pro-$\ell$.  It must be
procyclic: a free pro-$\ell$ group on two generators has
nonmetabelian finite $\ell$-group quotients, for example those
detected by a nonzero iterated commutator in its Magnus algebra.
Thus $S$ centralizes $A_\ell$.  The image of a Sylow subgroup
under a profinite quotient is a Sylow subgroup of the image.
Consequently $\theta_\ell(D)$ has trivial pro-$\ell$ part and
is contained in the finite prime-to-$\ell$ subgroup of
$\mathbb Z_\ell^\times$.  The displayed power kills it.
Factoring through the compact quotient $N/D=G_L$ gives the
continuous character $\lambda_\ell$.  No section is used.

Let $a$ be a generator of $A_\ell$ and choose an open
characteristic subgroup $V\triangleleft F$ with $a\notin V$.
Such subgroups are cofinal because $F$ is finitely generated.
Put $H=VA_\ell$.  It is open and normalized by $N$.
The quotient $H/V$ is a nontrivial finite cyclic $\ell$-group,
so the map $A_\ell\to H^{\rm ab}_\ell$ is nonzero.
The latter is a finite free $\mathbb Z_\ell$-module, as $H$
is the geometric group of a finite cover of an affine curve.
The map from $\mathbb Z_\ell$ is therefore injective.

Descend the finite cover associated to $H$ to a finite extension
and realize its arithmetic group as an open subgroup $E_Y\leq E$
with $E_Y\cap F=H$.  The group $N\cap E_Y$ has open arithmetic
image, denoted $G_{L_1}$.  Its geometric kernel is contained in
$H$ and acts trivially on $H^{\rm ab}$, so the action on the
detected line is the genuine arithmetic action of $G_{L_1}$.
This proves the character comparison, including its direction.

The boundary sequence for the smooth compactification is
\[
 0\longrightarrow W\longrightarrow H^{\rm ab}_\ell
   \longrightarrow T_\ell\operatorname{Jac}(\overline Y)
   \longrightarrow0,
\]
where $W$ is the cuspidal permutation lattice with its one
degree relation, tensored with $\mathbb Z_\ell(1)$.
If the line maps to zero in the last term, split the finitely
many cusps by another finite extension.  Then its character is
$\chi_\ell$.  Claim~B.111 applied to the \emph{actual}
$A_\ell\leq F$ puts it in one cusp inertia group.  Since
$A$ centralizes $A_\ell$, Stix's Proposition~16(3) puts all of
$A$ in that group.  Thus for noncuspidal $A$ the image in the
Jacobian is nonzero, hence injective over $\mathbb Q_\ell$.
The asserted weights are the good-reduction Frobenius theorem
and the Hodge--Tate theorem for the Tate module of an abelian
variety; Tate's statement is quoted in Serre (1989), Chapter~III,
Section~1.2.  This proves exactly where the weight-one alternative
enters the bare-normalizer problem.
\end{proof}
\paragraph{Check record.} checked only in this writing, 2026-11-18: GPT-6 Astra (OpenAI), ChatGPT Pro mode; internal effort value not exposed. Own detailed proof written and reread; primary source hypotheses matched. All N2.1--7 read in full by root; finite-cover topology join requested and inserted


\begin{claim}{B.113}{Proved}{The arithmetic exclusion of LM4 B.29, with its exact field}
In Claim~B.112 let $M/\mathbb Q$ be a finite Galois extension
containing $L$, and let $\mathcal P_M$ be the rational primes
unramified in $M$ whose Frobenius conjugacy class is that of a
complex conjugation in $\operatorname{Gal}(M/\mathbb Q)$.
Either $A$ is contained in a unique cusp inertia group or
\[
                   A_\ell=1\qquad(\ell\in\mathcal P_M).
\]
Thus cofinite prime support forces cuspidality.  If $L$ is
totally real, every such $A$ is cuspidal, without a support
condition.
\end{claim}
\begin{proof}
Suppose $A$ is noncuspidal and $A_\ell\ne1$.
Write $d=d_\ell$ and use the line $\psi$ of Claim~B.112.
Its power $\psi^d$ is the restriction of $\lambda=\lambda_\ell$.
It follows that $\lambda$ is unramified outside finitely many
places and Hodge--Tate.  For ramification this follows by taking
places unramified in $L_1/L$ and outside the finite exceptional
set of $\psi$.  Descent of the Hodge--Tate property through a
finite extension is Serre (1989), Appendix to Chapter~III,
Corollary~2, printed page~III--32.

Tate's local algebraicity theorem and global reciprocity express
$\lambda|_{G_M}$, on an open subgroup of the local units at
$\ell$, by a $\mathbb Q_\ell$-defined algebraic character of
$(\operatorname{Res}_{M/\mathbb Q}\mathbb G_m)_{\mathbb Q_\ell}$.
Its integral coefficients $n_\tau$ are indexed by embeddings
$\tau:M\hookrightarrow\overline{\mathbb Q}$.
It is trivial on a finite-index subgroup of global units.
The precise source inputs are Serre (1989), Chapter~III,
pages~III--7, III--9 and III--11, and Chapter~II,
Section~3.1, pages~II--29--30; these were opened.
For clarity, the required coefficient constraint follows by
applying each complex embedding to that unit identity and using
Dirichlet's unit lattice:
\[
             n_\tau+n_{c\tau}=2a
             \quad\hbox{for every embedding }\tau.
\]
The sum is constant because the logarithms of units span the
hyperplane defined by the product formula, with conjugate pairs
identified.  If $M$ is totally real, this directly gives
$n_\tau=a$ for every $\tau$.

For $\ell\in\mathcal P_M$, choose the prime of $M$ whose
Frobenius is the chosen $c$.  The decomposition group fixes the
coefficient vector because the algebraic character is defined
over $\mathbb Q_\ell$.  Hence $n_{c\tau}=n_\tau$, and the
display again gives one integer $n_\tau=a$ for all $\tau$.
Removing $\chi_\ell^a$ leaves a character trivial on an open
subgroup of each local unit group and unramified outside a finite
set.  It factors through a finite ray-class group.  We obtain
\[
                   \lambda|_{G_M}=\varepsilon\chi_\ell^a,
                   \qquad \varepsilon\text{ of finite order}.
\]
After a common finite extension the equality with $\psi^d$
therefore has two incompatible consequences: Frobenius weights
give $d=2a$, whereas Hodge--Tate weights give $a\in\{0,d\}$.
This proves the exclusion.  Chebotarev gives
$\mathcal P_M$ positive density, proving the cofinite-support
conclusion.  If $L$ is totally real, take a totally real Galois
closure $M$.  The unit argument then works at every $\ell$,
without a Frobenius condition.  Cusp uniqueness follows from
Stix's geometric cusp-intersection theorem.  This re-proves
LM4 B.29, including the preliminary enlargement of the normalizer.
\end{proof}
\paragraph{Check record.} checked only in this writing, 2026-11-18: GPT-6 Astra (OpenAI), ChatGPT Pro mode; internal effort value not exposed. Own detailed proof written and reread; primary source hypotheses matched. All N2.1--7 read in full by root; finite-cover topology join requested and inserted


\begin{claim}{B.114}{Proved}{The CM field obstruction and a necessary qualification over Q(i)}
If $A$ in Claim~B.112 is noncuspidal, its actual normalizer
field $L$ contains a CM subfield.  For each supported $\ell$
and each finite Galois $M/\mathbb Q$ containing $L$, some CM
subfield of $M$ embeds in $\mathbb Q_\ell$.  In particular,
if the normalizer image is the full $G_{\mathbb Q(i)}$, its
support consists only of primes split in $\mathbb Q(i)$.

For an arbitrary open image $G_L\subset G_{\mathbb Q(i)}$
one must use $L$ in this statement.  The character argument
alone does not exclude primes inert in $\mathbb Q(i)$.
For example $L=\mathbb Q(i,\sqrt{-2})$ admits a
$\mathbb Q_3$-valued geometric character of Frobenius weight
one although $3$ is inert in $\mathbb Q(i)$.
\end{claim}
\begin{proof}
We give the field-theoretic step of LM6 A.53 explicitly.
Let $G=\operatorname{Gal}(M/\mathbb Q)$ and
$H=\operatorname{Gal}(M/L)$.  The algebraic coefficient function
$n:G\to\mathbb Z$ of $\lambda_\ell|_{G_M}$ is right
$H$-invariant, since restriction from $L$ is pullback by the
norm from $M$ to $L$.  Subtract its mean, writing $n=a+b$.
The unit identity used in Claim~B.113, after every embedding,
gives
\[
 b(cg)=-b(g)
 \quad\hbox{for every complex conjugation }c\in G.
\]
The function $b$ is nonzero: if all exponents were equal, the
finite-times-cyclotomic conclusion and the incompatible weights
in that proof would apply without a restriction on $\ell$.

Let $J$ be the normal subgroup generated by products of two
complex conjugations.  It fixes $b$ on the left and, since $J$
is normal, also on the right.  Thus $b$ is right $JH$-invariant.
Modulo $J$ all complex conjugations have the same central image.
No complex conjugation belongs to $JH$: otherwise its left
action on $b$ would coincide, modulo $J$, with a right action
fixing $b$, contradicting $cb=-b\ne0$.
It follows that $M^{JH}\subset L$ is totally imaginary and
has the same nontrivial involution under every complex embedding;
the fixed field of that involution is totally real.  Its degree
over that fixed field is two, so it is a CM field.

Put $P=\operatorname{Stab}_G(b)$ for the left action.  It
contains $J$ and no complex conjugation.  The same argument
makes $M^P$ a CM field.  Since the character is
$\mathbb Q_\ell$-defined, the local decomposition group
$D_v$ fixes $b$, so $D_v\subset P$ and
$M^P\hookrightarrow M_v^{D_v}=\mathbb Q_\ell$.
This includes ramified primes.  For $L=\mathbb Q(i)$ take
$M=L$; its only CM subfield is itself, and its embedding into
$\mathbb Q_\ell$ is equivalent to splitting of $\ell$.

To verify the stated limitation construct the elliptic curve
\[
             A:y^2=x^3+4x^2+2x.
\]
Its discriminant is $512$ and $j=8000$.  The standard
degree-two map, here checked by substitution,
\[
 x'=x+4+2/x,\qquad y'=y(1-2/x^2),
\]
maps $A$ to $A':y'^2=x'^3-8x'^2+8x'$.
The latter is its quadratic twist by $-2$.
Writing $s=\sqrt{-2}$, the map
$(x',y')\mapsto(x'/(-2),y'/(-2s))$ is an isomorphism
$A'\to A$ over $\mathbb Q(s)$.  The resulting endomorphism
$\alpha$ of $A$ has differential $s$, since the two-isogeny
preserves the invariant differential.  Hence $\alpha^2=[-2]$:
endomorphisms of an elliptic curve in characteristic zero are
determined by their differentials.  This proves the claimed
CM action without relying on the value of $j$ as a catalogue.

The polynomial $T^2+2$ has two roots in $\mathbb Q_3$,
by Hensel lifting of its simple roots modulo three.
Thus $V_3A$ splits into two $\mathbb Q_3$-lines preserved
by $G_L$, because $\alpha$ is defined over $L$.  Each has
Frobenius weight one and Hodge--Tate weights in $\{0,1\}$.
Yet $T^2+1$ has no root modulo three, so $3$ is inert in
$\mathbb Q(i)$.  This is a counterexample to the proposed
\emph{character-theoretic inference}, not an asserted subgroup
$A\leq F$ solving the normalizer problem.  The symbolic and
finite-field calculations are in \texttt{n2\_certificates.py}.
LM6 A.53 itself correctly uses the actual normalizer field;
the qualification is necessary when applying it to the prompt.
\end{proof}
\paragraph{Check record.} checked only in this writing, 2026-11-18: GPT-6 Astra (OpenAI), ChatGPT Pro mode; internal effort value not exposed. Own detailed proof written and reread; primary source hypotheses matched. All N2.1--7 read in full by root; finite-cover topology join requested and inserted


\begin{claim}{B.115}{Proved}{The isogeny covers force the CM residue into a deeper kernel}
Let $U=E_0\setminus\{O\}$ for
$E_0:y^2=x^3-x$ over $\mathbb Q(i)$, and let
$1\ne C\leq R$ be procyclic with open arithmetic normalizer
image.  Put
\[
 P_n=\ker(F\longrightarrow F^{\rm ab}/nF^{\rm ab}),\qquad
 R_{\rm nil}(D)=\bigcap_{j\geq1}\gamma_j(D).
\]
Then
\[
 \begin{gathered}
 C\subset R_{\rm nil}(F')\subset F'',\\
 R_{\rm nil}(F')=
 \bigcap_{n\geq1,\,\ell\ {
m prime}}
 \ker(P_n\longrightarrow P_n^{(\ell)})
 \quad\hbox{inside }F'.
 \end{gathered}
\]
The inclusion $R_{\rm nil}(F')\subset R$ is strict.
Every weight-one detector of $C$ on a finite cover lies in
the Prym kernel of its compactified Jacobian map to $E_0$,
and such a detector has compactification genus at least two.
\end{claim}
\begin{proof}
We first re-prove the elliptic weighted argument used by LM6
A.35.  Let $D$ be the maximal pro-$\ell$ group of an elliptic
curve minus a nonempty finite divisor and let $B\simeq\mathbb
Z_\ell$ be arithmetically normalized with zero proper Tate
homology.  After a finite extension choose free pro-$\ell$
generators $x,y,c_1,\ldots,c_{r-1}$, with $x,y$ the handle
generators and the $c_i$ cuspidal.  Give $X,Y$ weight one and
the cusp variables $C_i$ weight two in the integral Magnus
embedding.  Its separation gives a first nonzero weighted term
of a generator $b$ of $B$, of degree $d\geq2$.

The filtration is preserved arithmetically.  Its generator
eigenvalues at good Frobenius have weights one on $X,Y$ and
two on the $C_i$.  Every eigenvalue on weighted degree $d$
therefore has complex absolute value $q^{d/2}$: its associated
graded module is a subquotient of the homogeneous tensor algebra
on those generators.  Inner conjugation is the identity on the
first nonzero term, so the normalizer character of $B$ occurs
on this term.  The finite-cover detector of Claim~B.112,
performed inside $D$, also realizes that same character in
curve homology, whose Frobenius weights are one or two.
Here is the required comparison with an actual curve cover.
An open subgroup of $D$ contains an open normal subgroup with
finite $\ell$-group quotient.  Its inverse image in the full
geometric group is therefore subnormal through a finite chain
with $\ell$-power indices.  Every finite $\ell$-group quotient
of that inverse image extends to a finite $\ell$-group quotient
of the original group after taking successive cores: at a
normal step of index $\ell$, the intersection of the conjugate
kernels has $\ell$-group quotient, and the remaining extension
also has $\ell$-power order.  Induction through the chain gives
the full induced pro-$\ell$ topology.  Consequently the open
subgroup of $D$ is precisely the maximal pro-$\ell$ group of
the corresponding finite \'etale curve cover, and its
abelianization is that curve's $\ell$-adic homology.
A common good Frobenius therefore forces $d=2$.

The degree-two primitive elements with zero linear handle term
are exactly
$a[X,Y]+\sum_i b_iC_i$.  To check this integrally, the
group-like identity gives zero coefficients of
$X\otimes X,Y\otimes Y$ in the reduced coproduct and
opposite coefficients of $X\otimes Y,Y\otimes X$.
The $C_i$ are already primitive in weighted degree two.
No division by two is used.  By the Weil pairing,
$[X,Y]$ has character $\det T_\ell E=\chi_\ell$, as do
all the rationalized cusp generators.  A nonzero torsion-free
leading term consequently forces the normalizer character
itself to be cyclotomic on the indicated open arithmetic group.

Now $C\subset R\subset F'$ lies in every $P_n$.
The associated cover is
$E_0\setminus E_0[n]\xrightarrow{[n]}E_0\setminus\{O\}$.
If $C$ has nonzero image in some $P_n^{(\ell)}$, its
$\ell$-component maps injectively to that procyclic image.
Its proper homology maps to zero under
$[n]:T_\ell E_0\to T_\ell E_0$, since its image in
$F^{(\ell)}$ is zero.  Multiplication by $n$ is injective
on this torsion-free module, even for $\ell\mid n$.
The proper homology is therefore zero.  The preceding weighted
argument makes the actual character on $C_\ell$ cyclotomic.
Claim~B.111 then puts $C_\ell$ in full cusp inertia, which
contradicts $C\subset R$.  This proves all the claimed
isogeny-cover vanishings without using the separate, deeper
pro-$\ell$ CM theorem of LM6 A.20.

For clarity, $R=R_{\rm nil}(F)$: the image of the intersection
of the nested closed lower-central subgroups in any finite
quotient is the intersection of their images, by compactness.
Its quotient therefore has exactly the finite nilpotent
quotients.  A finite nilpotent group is the product of its
Sylow subgroups, so these quotients are detected by the
maximal pro-prime quotients.  The same argument applies to
any profinite group $D$, without finite-generation assumptions.

For the group equality, restrictions of all the maps
$P_n\to P_n^{(\ell)}$ kill $R_{\rm nil}(F')$.
Conversely an element of $F'$ outside this subgroup survives
in a finite quotient $G$ of $F/R_{\rm nil}(F')$.  Its
derived subgroup $G'$ is nilpotent.  Take $n$ divisible by
the exponent of $G/G'$.  The image of $P_n$ in $G$ is
exactly $G'$: it is contained there by the choice of $n$
and contains it because $F'\subset P_n$.  A Sylow projection
of $G'$ detects the chosen element and factors through some
$P_n^{(\ell)}$.  This proves the reverse inclusion.
The containment in $F''$ follows because the abelian quotient
of $F'$ is pronilpotent.

Strictness is witnessed by the two-generated finite group
$G=\mathbb F_5^2\rtimes C_3$, where a generator acts by
\[
             T=\begin{pmatrix}0&-1\\1&-1\end{pmatrix}.
\]
It has $T^3=1$, $T-1$ invertible and no invariant line over
$\mathbb F_5$.  A vector and the generator generate $G$.
Thus $G'=\gamma_\infty(G)=\mathbb F_5^2\ne0$, while
$\gamma_\infty(G')=1$.  Compactness shows that the image
of an intersection of nested closed lower-central terms in
a finite quotient is the intersection of their images.
The image of $R$ is therefore $\mathbb F_5^2$, whereas
the deeper residual subgroup has zero image.  This is checked
by an actual finite-group certificate.

Finally for any finite cover $Y\to U$, the proper homology
of $C_\ell$ maps to zero in $T_\ell E_0$ because the
original pro-$\ell$ quotient kills it.  It belongs to the
kernel of the Jacobian map, the Prym part.  A compactification
of genus one maps to $E_0$ by a nonconstant isogeny up to
translation, which is injective on rational Tate modules.
Genus zero cannot map nonconstantly to $E_0$.  Hence a
nonzero weight-one residual detector requires genus at least
two.  No assertion that such a Prym line lifts to $C$ has
been made.
\end{proof}
\paragraph{Check record.} checked only in this writing, 2026-11-18: GPT-6 Astra (OpenAI), ChatGPT Pro mode; internal effort value not exposed. Own detailed proof written and reread; primary source hypotheses matched. All N2.1--7 read in full by root; finite-cover topology join requested and inserted


\begin{claim}{B.116}{Proved}{An explicit real descent and its fixed subgroup}
The CM curve $E_0/\mathbb Q(i)$ has the based real descent
\[
 E':y^2=x^3+4x,\qquad b=(0,0),\qquad
 (x,y)\longmapsto(u^2x,u^3y),\quad u=(1+i)/2.
\]
For $U'=E'\setminus\{O\}$ and its rational-point action,
complex conjugation $c\in G_{\mathbb Q}$ satisfies
\[
 F\rtimes\langle c\rangle
     \simeq\widehat{\mathbb Z}\amalg C_2,
 \qquad C_F(c)=1.
\]
The isomorphism over $\mathbb Q(i)$ transports this additional
involution to the geometric group in N2.  It lies outside
$G_{\mathbb Q(i)}$ and is not a real place of that field.
\end{claim}
\begin{proof}
The identity $u^4=-1/4$ verifies the displayed curve isomorphism,
which fixes the origin and $(0,0)$.  The cubic $x(x^2+4)$ has
one real root, so $E'(\mathbb R)$ has one connected component,
a circle.  Removing $O$ makes that fixed locus a real line.
Here is a concrete based topological computation.  The
automorphism $(x,y)\mapsto(-x,iy)$ makes the period lattice
stable under multiplication by $i$.  A rank-two lattice with
this property is a rank-one $\mathbb Z[i]$-module; since that
ring is principal, complex uniformization gives
$E'(\mathbb C)\simeq\mathbb C/(\mathbb Z+i\mathbb Z)$
after scaling.  Conjugation fixing the origin lifts to
$z\mapsto\alpha\overline z$, where
$\alpha\in\{1,-1,i,-i\}$.  The possibilities $\alpha=\pm1$
have two fixed circles and $\alpha=\pm i$ one fixed circle.
The real-component calculation above therefore forces the
latter case.  Changing the square-lattice coordinate by a
unit if necessary gives $\alpha=i$.

Thus in coordinates $(x,y)\in\mathbb R^2/\mathbb Z^2$ the
involution is $(x,y)\mapsto(y,x)$.  The unique nonzero real
two-torsion point, namely $b=(0,0)$ on the cubic, corresponds
to $(1/2,1/2)$ on the square torus.  The horizontal loop
through this point and its vertical image under $c$ form a
free basis $b_1,b_2$ of the torus punctured at the origin.
Indeed their complement is an open square containing the
removed point, whose punctured interior retracts to its
boundary.  Conjugation exchanges these two actual based
loops.  Hence the real topological extension is
\[
 \langle b_1,b_2,c\mid c^2=1,\ cb_1c=b_2,\ cb_2c=b_1\rangle
       \simeq\mathbb Z*C_2.
\]
This computation retains the rational-point involution itself,
not merely its outer class.

Riemann existence, SGA~1, Expos\'e~XII, Theorem~5.1, identifies
finite algebraic and topological covers and their equivariant
descent.  The free index-two kernel has its full profinite
topology: characteristic finite-index subgroups inside any
specified finite-index subgroup are stable under the involution
and yield finite quotients of the semidirect product.  Completion
therefore gives the asserted free profinite product.  Herfort--Ribes
(1985), Theorem~B (or its strengthened Theorem~B prime), places
the centralizer of the nonidentity element of the $C_2$ factor
inside that factor.  Its intersection with the geometric kernel
is trivial.  Their Theorem~A also says that every involution in
this product is conjugate to $c$, so every such involution has
trivial geometric centralizer.

The base change to $\mathbb Q(i)$ identifies the based geometric
groups and their $G_{\mathbb Q(i)}$ actions.  It does not put
$c$ inside this subgroup: complex conjugation sends $i$ to $-i$.
This is an external descent operation available for the
characteristic subgroup $R$, and not an involution supplied by
the normalizer field of an arbitrary $C$.
\end{proof}
\paragraph{Check record.} checked only in this writing, 2026-11-18: GPT-6 Astra (OpenAI), ChatGPT Pro mode; internal effort value not exposed. Own detailed proof written and reread; primary source hypotheses matched. All N2.1--7 read in full by root; finite-cover topology join requested and inserted


\begin{claim}{B.117}{Proved}{The LM4 real restrictions have a CM analogue for the whole kernel}
For the based real descent in Claim~B.116, every nontrivial
Sylow subgroup of $R$ is a free pro-prime group of rank at
least two.  There is no nontrivial closed abelian subgroup of
$R$ stable under this full $G_{\mathbb Q}$ action, and no
nontrivial closed pro-odd subgroup stable under that action.
Consequently
\[
 Z(R)=1,\qquad O_{2'}(R)=1,
\]
and a nontrivial $R$ has a nonabelian $2$-Sylow subgroup.
These conclusions do not assert that every subgroup with an
open normalizer image in $G_{\mathbb Q(i)}$ is stable under
this full action.
\end{claim}
\begin{proof}
The group $R$ is characteristic in $F$, so is stable under the
based $G_{\mathbb Q}$ action.  If $S$ is an $\ell$-Sylow
subgroup of $R$, profinite Sylow conjugacy gives
\[
       N_{R\rtimes G_{\mathbb Q}}(S)\twoheadrightarrow
                              G_{\mathbb Q}.
\]
Indeed any Galois translate of $S$ is conjugate to it by an
element of $R$.  This is surjectivity of a compact normalizer
projection; it is not a section or invariance of $S$ under
the original section.  The subgroup $S$ is free pro-$\ell$.
If it had rank one, Claim~B.113 with normalizer field
$\mathbb Q$ would make it cuspidal, contradicting
Claim~B.111.  Thus its rank is at least two.

For a closed abelian subgroup $A\leq R$, each Sylow is both
abelian and free pro-prime, hence either zero or $\mathbb Z_\ell$.
Therefore $A\simeq\prod_{\ell\in\Sigma}\mathbb Z_\ell$ is
procyclic, by choosing a generator in each component and using
the Chinese remainder theorem in every finite quotient.  If it
is $G_{\mathbb Q}$-stable, its semidirect product with that
Galois group normalizes it, and Claim~B.113 again contradicts
its inclusion in $R$.  This gives $A=1$ and, in particular,
$Z(R)=1$.

Let $H\leq R$ be pro-odd and $c$-stable.  For every $c$-stable
open normal $V\triangleleft H$, the fixed-point map
$C_H(c)\to C_{H/V}(c)$ is surjective.  If a coset is fixed,
a lift conjugates the complement $\langle c\rangle$ to
another complement in $V\rtimes\langle c\rangle$.
These are $2$-Sylow subgroups, since $V$ is pro-odd, and
Sylow conjugacy corrects the lift by an element of $V$.
Claim~B.116 therefore makes the induced involution on every
finite quotient $H/V$ fixed-point free.

A finite group with a fixed-point-free involution $\alpha$
is abelian and $\alpha$ is inversion.  Indeed
$h\mapsto h^{-1}\alpha(h)$ is injective (equal values give
a fixed quotient of the two elements), hence surjective.
Each such value is sent to its inverse by $\alpha$.
Surjectivity shows that inversion is an automorphism of the
whole group, forcing commutativity.  Taking the cofinal system
of $c$-stable $V$ proves that $H$ is abelian and inverted by
$c$.  If it is stable under the full $G_{\mathbb Q}$, the
previous paragraph gives $H=1$.  The largest closed normal
pro-odd subgroup $O_{2'}(R)$ is characteristic, so vanishes.
If $R$ itself were nontrivial and pro-odd, the same argument
would contradict it.  Its $2$-Sylow is therefore nontrivial,
and is nonabelian by the first paragraph.

This is a re-proof of the mechanisms of LM4 B.33, B.50 and
B.51 for this particular kernel.  It requires the actual real
descent action and the full totally real normalizer image at
the indicated steps; neither is inferred for an individual
open-normalized subgroup in the original CM question.
\end{proof}
\paragraph{Check record.} checked only in this writing, 2026-11-18: GPT-6 Astra (OpenAI), ChatGPT Pro mode; internal effort value not exposed. Own detailed proof written and reread; primary source hypotheses matched. All N2.1--7 read in full by root; finite-cover topology join requested and inserted



\noindent\emph{Statement B.118 and its produced argument are retained in Part II.}


\paragraph{Checks and provenance.}
The arguments and the corrections to the finite-cover and
normalizer-field inferences were written and reread by GPT-6
Astra (OpenAI), ChatGPT Pro mode, 18 November 2026.  The
published Stix text, its supplied 2008 source, Serre's specified
pages, the Herfort--Ribes original text, and the Morales source
were opened.  Finite calculations check the explicitly identified
matrix, isogeny and numbering assertions only.  No computation
is represented as proving the unresolved full profinite assertion,
and no priority, human, or different-model verification is claimed.

% AI-produced by GPT-6 Astra (OpenAI), ChatGPT Pro mode, 18 November 2026.
% No priority or human/external verification is claimed.
\subsection{The metabelian field: centralizers, affine automorphisms, and norm descent}

Fix a prime $p$ and the chosen valuation on $\overline{\mathbb Q}$.
Put
\[
 K_p=\operatorname{Frac}(\mathbb Z_{(p)}^{\rm sh}),\quad
 I=G_{K_p},\quad F_p=\overline{\mathbb Q}^{I'},\quad
 B_p=\overline{\mathbb Q}^{I''}.
\]
A prime denotes the closed derived subgroup.  The chosen local
absolute Galois group is denoted $D_p$.  Henselization and completion
identify $D_p$ with $G_{\mathbb Q_p}$ and $I$ with its inertia.
All field extensions occurring below are inside $\overline{\mathbb Q}$.

\begin{claim}{B.119}{Proved}{Wild inertia has trivial local centralizer, at every prime}
For a finite extension $M/\mathbb Q_p$, let $P_M$ be wild inertia.
Then
\[
 C_{G_M}(P_M)=1,\qquad C_{G_M}(P_M'')=1.
\]
Moreover $G_M$ is torsion free, including when $p=2$.
\end{claim}
\begin{proof}
The primary source for wild freeness is Iwasawa (1955),
Theorem~2(ii), printed page~463, together with the paragraph immediately
before it: the wild group is free pro-$p$ of countably infinite rank.
Kochloukova--Zalesskii (2011), Theorem~5.1, Case~I, gives the
procyclic centralizer and normalizer of a nonidentity element in a
finitely generated free pro-$p$ group and recalls that closed
subgroups of free pro-$p$ groups are free.  Both exact primary
passages were opened.  For the infinite-rank group here their needed
consequence follows by taking the closed subgroup generated by the
centralizing element and two noncommuting elements.  It is finitely
generated and free pro-$p$, and a nonzero element cannot centralize
two noncommuting elements in it.  Thus the centralizer in $P_M$ of
any nonabelian subgroup is trivial.

Write $T=M^{\rm tame}$, normalize $v_M(M^\times)=\mathbb Z$,
and put $e=v_M(p)$.  The residue field of $T$ is
$\overline{\mathbb F}_p$, its value group is
$\mathbb Z_{(p)}$, and it contains $\mu_p$, also for $p=2$.
Suppose $c\in C_{G_M}(P_M)\cap I_M$ has nontrivial tame image.
Its image in $I_M/P_M=\prod_{\ell\ne p}\mathbb Z_\ell(1)$
has infinite order in at least one component.  We may therefore
choose a prime-to-$p$ integer $m$ as large as necessary such that,
for a uniformizer $\pi_M$ and $z=\pi_M^{1/m}$,
\[
 c(z)=\eta z,\qquad \eta\in\mu_m,\quad
 \operatorname{ord}(\eta)>p-1,\quad
               1/m<ep/(p-1).
\]
Indeed use increasing powers of a prime at which the tame component
is nonzero.  Since prime-to-$p$ roots of unity reduce injectively,
the residue of $\eta$ does not lie in $\mathbb F_p^\times$.

Kummer theory identifies $T^\times/T^{\times p}$ with
$H^1(P_M,\mu_p)$.  Because $c$ centralizes $P_M$, its action
on this cohomology is multiplication by $\chi_p(c)\bmod p$:
it acts only on the coefficient module.  Let
$a\in\{1,\ldots,p-1\}$ represent that nonzero residue.
Applying this to $1+z$ makes
\[
                     r=\frac{1+\eta z}{(1+z)^a}
\]
a $p$th power in $T$.  But its leading term is
$(\eta-a)z$, with unit coefficient, so $v(r-1)=1/m$.
If $r=y^p$, the residue of $y$ is one, since in characteristic
$p$ the only $p$th root of one is one.  Put $t=y-1$.
In the binomial expansion of $(1+t)^p-1$, if
$v(t)<e/(p-1)$ the unique lowest term is $t^p$; if
$v(t)>e/(p-1)$ the lowest term is $pt$; and at equality
all terms have value at least $ep/(p-1)$.
Our strict bound $v(r-1)<ep/(p-1)$ consequently forces
\[
                  v(t)=1/(mp)\notin\mathbb Z_{(p)},
\]
a contradiction.  This includes $p=2$ without any change.
Thus $c$ has trivial tame image and belongs to $P_M$.
The centre of the nonabelian free group $P_M$ is trivial,
so $c=1$.

Let $C=C_{G_M}(P_M)$.  Since $P_M$ is normal in $G_M$,
so is $C$.  Hence $[C,I_M]\subset C\cap I_M=1$.
An element of $C$ therefore centralizes tame inertia as well.
If its residue exponent is $\nu\in\widehat{\mathbb Z}$,
conjugation on tame inertia is exponentiation by $q^\nu$,
where $q$ is the residue cardinality of $M$.
This action is faithful: for every $n\geq2$, $q$ has
exact order $n$ modulo $q^n-1$, which is prime to $p$.
Thus trivial action gives $\nu\in n\widehat{\mathbb Z}$
for every $n$, and $\nu=0$.  The element lies in $C\cap I_M$
and is one.  This proves the first centralizer equality by a
local Kummer argument, without invoking reconstruction of the tame
field.

The subgroup $P_M''$ is nonabelian.  For an explicit witness,
surject $P_M$ onto $\operatorname{UT}_9(\mathbb F_p)$,
using eight of its free generators for adjacent transvections.
With $x_{ij}=1+E_{ij}$, the identities
\[
 [[x_{12},x_{23}],[x_{34},x_{45}]]=x_{15},\qquad
 [[x_{56},x_{67}],[x_{78},x_{89}]]=x_{59},\qquad
                         [x_{15},x_{59}]=x_{19}\ne1
\]
hold over every $\mathbb F_p$.  Thus its second derived group
is already nonabelian.  Now $C_{G_M}(P_M'')$ is normal in
$G_M$, and its commutator with $P_M$ lies in
$C_{P_M}(P_M'')=1$.  It centralizes $P_M$, so the first
part makes it trivial.

For completeness, here is an elementary torsion-freeness argument,
including the small prime.  A torsion element has a power of prime
order $q$.  Its image in the unramified quotient
$\widehat{\mathbb Z}$ is zero, so it fixes the strict-unramified
field, which contains $K_p$.  Let $H$ be the cyclic group of order
$q$ it generates and let $A=\overline{\mathbb Q}^{H}$.
Then $\overline{\mathbb Q}/A$ is cyclic of degree $q$.
For odd $q$, $H$ fixes $\mu_q$ because its automorphism group
has order $q-1$.  Kummer theory writes this extension as
$A(\beta)$, $\beta^q=a\notin A^{\times q}$.
Choose $\gamma$ with $\gamma^q=\beta$ in the algebraically
closed field.  Taking norms gives
$N(\gamma)^q=N(\beta)=a$, a contradiction.
For $q=2$, the same argument works if $i\in A$, since
$N(\beta)=-a$ and $-1$ is a square.  If $p$ is odd then
$i\in K_p\subset A$.  If $p=2$ and $i\notin A$,
we have $\overline{\mathbb Q}=A(i)$.  A primitive cube root
$\zeta$ lies in $K_2\subset A$, and
\[
                  (\zeta^2)^2+\zeta^2=-1.
\]
Thus $w=\zeta^2+i\zeta$ has norm $-1$.  A square root
$v$ of $w$ exists in $A(i)$, giving $N(v)^2=-1$ in $A$,
contrary to $i\notin A$.  No prime-order element exists.
The same proof applies to the local group via the chosen
henselization--completion identification.
\end{proof}
\paragraph{Check record.} checked only in this writing, 2026-11-18: GPT-6 Astra (OpenAI), ChatGPT Pro mode; internal effort value not exposed. Own detailed proof written and reread; primary source hypotheses matched. Root full second reading of META.1--4 and META.K passed, including all small-prime and finite-grid joins


\begin{claim}{B.120}{Proved}{The metabelian coefficient centralizer is trivial}
For every finite extension $E/B_p$,
\[
                         C_{G_{\mathbb Q}}(G_E)=1.
\]
This re-proves the coefficient-centralizer assertion of AS-FF.40
without its dependence on a reconstruction theorem for the tame field.
\end{claim}
\begin{proof}
Enlarge $E$ to $E_1=B_pD$, where $D/\mathbb Q$ is a finite
Galois number field containing generators of $E/B_p$.
Then $H=G_{E_1}=I''\cap G_D$.  Let $M=D_v$ be the chosen
local completion.  Its wild group satisfies
\[
                       P_M''\subset H.
\]
Indeed $P_M\subset I\cap G_D$, and taking two derived
subgroups gives the inclusion.  Claim~B.119 shows that
$H$ is nontrivial, even nonabelian; hence $E_1$ is not
separably closed.

Suppose $\sigma\in G_{\mathbb Q}$ centralizes $G_E$.
It centralizes $H$, and therefore preserves its fixed field
$E_1$.  This field has the henselian rank-one valuation induced
by $v$.  Transport by $\sigma$ gives another such valuation.
The F. K. Schmidt theorem in Kuhlmann, Chapter~9, Theorem~9.44,
says that two independent nontrivial henselian valuations force
separable closedness.  Distinct rank-one valuations are
independent.  They must therefore be equivalent.  A henselian
valuation has a unique prolongation to the algebraic closure,
so $\sigma$ preserves the chosen valuation on
$\overline{\mathbb Q}$ and belongs to $D_p$.

Because $D$ is finite Galois, some positive power
$\sigma^r$ fixes it and belongs to $G_M\subset D_p$.
It centralizes $P_M''\subset H$, so Claim~B.119 gives
$\sigma^r=1$.  The torsion-freeness proved there gives
$\sigma=1$.  This proves the assertion for the original $E$,
not only for the enlargement.
\end{proof}
\paragraph{Check record.} checked only in this writing, 2026-11-18: GPT-6 Astra (OpenAI), ChatGPT Pro mode; internal effort value not exposed. Own detailed proof written and reread; primary source hypotheses matched. Root full second reading of META.1--4 and META.K passed, including all small-prime and finite-grid joins


\paragraph{The exact Kummer input.}
We use the following statement of B.56, and re-prove its arithmetic
content here.  For every finite $E/B_p$, with rational $p$-Kummer map
$\kappa_E$, put
\[
 W=\{a\in B_p^\times:a^m\in F_p^\times
                              \text{ for some }m\geq1\}.
\]
Then
\[
 \ker\kappa_E=W,\qquad W\cap(1-W)=F_p\setminus\{0,1\}.
                                                        \tag{META.K}
\]
To check the first assertion, enlarge to $E_1=B_pD$ with $D$
locally Galois and containing the element under consideration; put
$J=I\cap G_D$, $H=I''\cap J$,
$\Gamma=J/H$, and $C=(I'\cap J)/H$.
A Kummer class vanishing on $H$ inflates from
$H^1(\Gamma,\mathbb Q_p(1))$.
Both $C$ and $\Gamma/C$ are abelian, $C$ fixes the coefficients,
and $\chi_p$ has infinite image on $\Gamma/C$.
For an abelian group with nontrivial scalar action, the cocycle
identity
$(\chi(g)-1)c(h)=(\chi(h)-1)c(g)$ makes its first cohomology
zero.  Inflation--restriction therefore injects this cohomology
into $H^1(C,\mathbb Q_p(1))$.
Every $g\in I'$ centralizes $C$ and fixes the coefficients,
so fixes the inflated class.  Thus
$\kappa_J(g(a)/a)=0$.  All these algebraic numbers belong to
the finite local field $M=D_v$.
Restriction of $H^1(G_M,\mathbb Q_p(1))$ to $I_M=J$ is
injective, since the inertia invariants of the coefficient module
are zero.  Valuation and the logarithm on sufficiently small
principal units show that the local rational Kummer kernel is
exactly the finite group $\mu(M)$.
An integer killing it makes $a^m$ fixed by all of $I'$, so
$a^m\in F_p$.  Since $F_p$ contains all roots of unity,
this radical lies in its maximal abelian extension $B_p$.
Conversely all divisions of an element of $W$ are radicals of
an element of $F_p$, so lie in $B_p$ and give zero Kummer class.
Integral and rational vanishing agree here because $G_E$ fixes
$\mathbb Z_p(1)$ and homomorphisms into that torsion-free module
have no torsion.

For the second assertion, if $a,1-a\in W$, then for every
$g\in I'$ we have $g(a)=\zeta a$ and
$g(1-a)=\eta(1-a)$ with $\zeta,\eta\in\mu_\infty\subset F_p$.
Their sum equals one.  If $\zeta\ne\eta$, it expresses
$a=(1-\eta)/(\zeta-\eta)\in F_p$; if they are equal,
they both equal one.  Thus every $g$ fixes $a$ and $a\in F_p$.
The converse is immediate.  This proof includes every prime and
all finite extensions; no small-prime exception is hidden in (META.K).

\begin{claim}{B.121}{Proved}{Every pinned affine-line autoequivalence has an affine point map}
Let $E/B_p$ be finite and let $\phi$ be a pinned relative
autoequivalence of $\mathbb A^1_{E,\mathrm{et}}$ over $BG_E$.
Its actual equivariant point map is
\[
                         f(t)=at+b,
                 \qquad a\in E^\times,\quad b\in E.
\]
This is an independent re-proof of AS-FF.44, whose full earlier
proof was subsequently read and matched.  The assertion concerns
the actual point map.  It does not by itself identify every geometric morphism with that point map,
or prove reconstruction for arbitrary pinned morphisms.
\end{claim}
\begin{proof}
The actual bijection $f:\overline{\mathbb Q}\to
\overline{\mathbb Q}$ is supplied by Carlson--Haine--Wolf,
Proposition~3.6.  Fix distinct points $u,v$ and a third point
$x=u+c(v-u)$, with $c\in F_p\setminus\{0,1\}$.
Pass to a finite extension of $E$ containing these points and
their images.  The punctured line at $u$ is carried to the
punctured line at $f(u)$.  Its geometric $p$-Kummer module
has rank one, so the induced map multiplies it by a unit
$\alpha_u\in\mathbb Z_p^\times$.  Subtracting evaluations
at $x$ and $v$ cancels the arithmetic summand and gives
\[
 \kappa\!\left(\frac{f(x)-f(u)}{f(v)-f(u)}\right)
       =\alpha_u\kappa\!\left(\frac{x-u}{v-u}\right)=0.
\]
The puncture at $v$ gives the same assertion for one minus the
left argument.  Formula (META.K) puts that argument in $F_p$.
Applying the same reasoning to the inverse autoequivalence proves
that $f$ maps each $F_p$-affine line bijectively onto an
$F_p$-affine line.

Here is the affine-geometry step in the present setting.
Regard $\overline{\mathbb Q}$ as a vector space over $F_p$;
its dimension exceeds one, since $I'$ is nontrivial.
Normalize by
\[
 T(t)=\frac{f(t)-f(0)}{f(1)-f(0)}.
\]
It fixes zero and one and preserves lines in both directions.
It preserves affine planes, since the plane through three
noncollinear points is the union of lines joining points of two
intersecting lines.  Within a plane it preserves parallelism,
characterized by disjointness.  The parallelogram construction
therefore gives $T(x+y)=T(x)+T(y)$ for linearly independent
$x,y$.  For dependent nonzero $x,y$, choose $z$ independent
of their line and compute $T(x+y+z)$ in the two ways; the
case $x+y=0$ follows by the same construction or by comparing
$T(z)$ with $T(x+(z-x))$.  Thus $T$ is additive everywhere.
For each nonzero $x$ write $T(cx)=\tau_x(c)T(x)$.
Applying additivity to independent $x,y$ shows
$\tau_x=\tau_y=\tau_{x+y}$; an auxiliary independent vector
handles any pair of directions.  The common bijection
$\tau:F_p\to F_p$ is additive, multiplicative by evaluating
$T(cd x)$ in two ways, and fixes one.  It is a field
automorphism and $T$ is $\tau$-semilinear.

The puncture at zero, normalized by the evaluation at one,
also gives, with one fixed $\alpha\in\mathbb Z_p^\times$,
\[
                   \kappa(T(t))=\alpha\kappa(t)
                                  \qquad(t\ne0).
\]
The equality holds after every finite scalar extension needed for
the point.  Consequently
$T(uv)/(T(u)T(v))\in W$ for $u,v\ne0$.
Fix $u\ne0$ and $v\notin F_p$, and put
$A=T(uv)/T(u)$, $B=T(v)$.  Semilinearity applied to $v+c$
shows that
\[
                         \frac{A+s}{B+s}\in W
                         \qquad(s\in F_p)
                                                        \tag{META.R}
\]
whenever numerator and denominator are nonzero.
For $g\in I'$, the quotient of the $g$-translate of (META.R)
by itself is a root of unity.  Restrict $s$ to rational integers.
All values of
\[
                 \frac{(gA+s)(B+s)}{(gB+s)(A+s)}
\]
then belong to one fixed number field containing
$A,B,gA,gB$.  Its roots of unity form a finite set.
A rational function taking values in a finite set at infinitely
many arguments is constant, and its value at infinity is one.
If $A\ne B$, comparison of the two linear factors forces
$gA=A$ and $gB=B$.  For every $g\in I'$ this gives
$B\in F_p$, contrary to $v\notin F_p$ and
$T^{-1}(F_p)=F_p$.  Hence $A=B$ and $T(uv)=T(u)T(v)$.
When $v\in F_p$ the same equality is already semilinearity.
Thus $T$ is a field automorphism of $\overline{\mathbb Q}$,
so belongs to $G_{\mathbb Q}$.

Equivariance of $f$ makes $f(0)$ and $f(1)$ elements of $E$.
The normalized $T$ therefore commutes with $G_E$.
Claim~B.120 gives $T=1$, and the stated affine formula follows.
The finite-root-of-unity argument above is essential: absolute
value one alone would not force a rational function to be constant.
\end{proof}
\paragraph{Check record.} checked only in this writing, 2026-11-18: GPT-6 Astra (OpenAI), ChatGPT Pro mode; internal effort value not exposed. Own detailed proof written and reread; primary source hypotheses matched. Root full second reading of META.1--4 and META.K passed, including all small-prime and finite-grid joins


\begin{claim}{B.122}{Proved}{A finite Gaussian grid descends all actual norm coefficients by one derived level}
Let $E/B_p$ be finite, $X/E$ smooth geometrically integral affine,
and $\phi:X_{\mathrm{et}}\to\mathbb A^1_{E,\mathrm{et}}$
a pinned relative morphism with nonconstant actual point map $f$.
Fix a finite generically \'etale Noether projection
$g:X\to\mathbb A^n_E$ of degree $d$.
After a finite enlargement $E_1=B_pD$ there is a nonempty
principal open $V\subset\mathbb A^n_{E_1}$ such that for every
finite number field $D'\supset D$ and every
$t\in V(F_pD')$, the polynomial
\[
                P_t(Z)=\prod_{x\in g^{-1}(t)}(Z-f(x))
\]
has its coefficients in $F_pD'$.  The fibre is taken on the
\'etale locus, so has exactly $d$ points.
Equivalently, for arbitrary algebraic $t$ in that open,
its coefficients belong to $F_pD(t_1,\ldots,t_n)$.
For $X=\mathbb A^n$ and the identity projection this says
$f(t)\in F_pD(t_1,\ldots,t_n)$, with one $D$ independent of $t$.
\end{claim}
\begin{proof}
Take the Gaussian grid
\[
          \mathcal A=\{r+is:0\leq r,s\leq4d\}
                         \subset\mathbb Q(i)\subset F_p.
\]
For each $a\in\mathcal A$ the set of actual preimages of $a$
is the geometric-point set of a proper closed subset of $X$.
Finiteness of $g$ makes its image a proper closed subset of affine
space.  Remove these images and the branch locus, and choose a
nonempty principal open $V$ in the complement.
Corestrict the Kummer pullback of the coordinate minus $a$ from
$g^{-1}(V)$ to $V$.  On every rational point $t$ its evaluation
is the actual Kummer value of
\[
                     b_a(t)=P_t(a),
\]
up to the fixed harmless sign from the chosen norm convention.
The sign is a root of unity and cancels below.
This is the norm compatibility proved directly in Claim~B.107;
it uses the product over the finite \'etale fibre, not an
assumed continuity of $f$.

The coordinate ring of a principal open of affine space is a UFD.
After factoring its finitely many defining polynomials and making a
finite extension, its geometric first $p$-adic Kummer cohomology is
freely generated by the Kummer classes of those factors.  The
Hochschild--Serre sequence, or the Kummer sequence at each finite
level followed by inverse limits, consequently writes each norm class
as
\[
                   \xi_a=\beta_a+\sum_j\lambda_{aj}\kappa(h_j),
                  \qquad \lambda_{aj}\in\mathbb Q_p.
                                                        \tag{META.N}
\]
The finitely many $h_j$ and all scheme and projection coefficients
lie in one finite number field $D$.  No rationality of the
$\lambda_{aj}$ is asserted or needed.

Put $E'=B_pD'$ and $F=F_pD'$, and choose
$\sigma\in I'\cap G_{D'}=G_F$.  It normalizes $E'$ and fixes
$\mathbb Q_p(1)$, all $h_j$, and all points $t\in V(F)$.
After subtracting (META.N) at two points $t,t_0$, its evaluation
is $\sigma$-invariant.  Formula (META.K) gives
\[
 \frac{\sigma(b_a(t)/b_a(t_0))}{b_a(t)/b_a(t_0)}\in W.
                                                        \tag{META.C}
\]
We use the following elementary strengthening of this membership.
If $x$ is algebraic, $\sigma\in I'$, and
$u=\sigma(x)/x\in W$, then $u$ is a root of unity.
Every $I'$-translate of an element of $W$ differs from it by
a root of unity.  Moreover $\sigma$ has finite order on $x$,
so the product of finitely many consecutive translates of $u$
telescopes to one.  Hence a positive power of $u$ is a root of
unity, and so is $u$.  Applying this to (META.C) gives
\[
 |\sigma P_t(a)P_{t_0}(a)|^2
                   =|P_t(a)\sigma P_{t_0}(a)|^2
                              \qquad(a\in\mathcal A),
                                                        \tag{META.G}
\]
for the fixed complex embedding.  The notation $\sigma P_t$ means
coefficientwise Galois action; every grid point is fixed by $\sigma$.

Both sides of (META.G), as polynomials in the real and imaginary
parts of $a$, have degree at most $4d$ in each variable.
The rectangular grid determines such a polynomial, so they agree
for every complex $a$.  The two holomorphic monic polynomials
$\sigma P_t\,P_{t_0}$ and $P_t\,\sigma P_{t_0}$ have the
same absolute value everywhere.  Their zeros with multiplicity
therefore coincide, and their common leading coefficient makes
them equal.
Choose $t_1\in V(F)$ for which the actual output multisets at
$t_0$ and $t_1$ are disjoint.  This is possible: remove the
finitely many proper closed images under $g$ of the preimages of
the values in the first fibre; $F^n$ is Zariski dense even over
$\overline{\mathbb Q}$.  The resulting $P_{t_0},P_{t_1}$ are
coprime.  The polynomial identity then gives
$P_{t_0}\mid\sigma P_{t_0}$, hence equality by monicity and
degree.  It follows that $\sigma P_t=P_t$ for every $t\in V(F)$.
Since this holds for all $\sigma\in G_F$, the coefficients
belong to $F$.  All arguments remain valid for every $D'\supset D$,
proving the stated uniformity and the last reformulation.
\end{proof}
\paragraph{Check record.} checked only in this writing, 2026-11-18: GPT-6 Astra (OpenAI), ChatGPT Pro mode; internal effort value not exposed. Own detailed proof written and reread; primary source hypotheses matched. Root full second reading of META.1--4 and META.K passed, including all small-prime and finite-grid joins


\begin{claim}{B.123}{Proved}{Reconstruction at every bounded nilpotency level, including the small primes}
For every prime $p$ and every integer $c\geq1$, put
\[
               K_{p,c}=\overline{\mathbb Q}^{\gamma_{c+1}(I)}.
\]
Every finite extension of $K_{p,c}$ admits closed-point-preserving
reconstruction of finite-type schemes from their relative \'etale
topoi, after universal homeomorphisms are inverted.
In particular this holds over every finite extension of
$F_p=K_{p,1}$, at every prime.
This re-proves the positive assertions AS-FF.37 and AS-FF.39 by the
actual-point norm method; it does not assert the same for their
unbounded union when $p=2,3$.
\end{claim}
\begin{proof}
We first give the elementary cohomology input, since it also
explains why a finite bound on the nilpotency class matters.
If $\Gamma$ is a profinite nilpotent group of finite class and
$V=\mathbb Q_p(\chi)$ is a one-dimensional continuous module
with nontrivial $\chi$, then $H^1(\Gamma,V)=0$.
For an abelian group this is the cocycle identity used in (META.K).
For larger class let $Z$ be its centre.  If $\chi|_Z\ne1$,
then $V^Z=0$ and $H^1(Z,V)=0$, so inflation--restriction gives
the assertion.  If $\chi|_Z=1$, then
$H^1(Z,V)=\operatorname{Hom}(Z,\mathbb Q_p)$ with the scalar
$\chi$ action of $\Gamma/Z$: conjugation on $Z$ is trivial.
Its invariant part is zero.  Inflation--restriction reduces the
assertion to the smaller-class quotient $\Gamma/Z$ with the
same nontrivial character.  This proves the vanishing by induction.

Let $K=K_{p,c}$ and let $E/K$ be finite.  All roots of unity
belong to $K$, since $\gamma_{c+1}(I)\subset I'$.
For each element $a$ (or $b$ in the scalar test) under
consideration, enlarge $E$ to $KD$ with $D/\mathbb Q$ finite
Galois containing that element as well as generators of $E/K$,
and put $M=D_v$, $J=I\cap G_D=I_M$, and
$H=\gamma_{c+1}(I)\cap J$.  The quotient $\Gamma=J/H$
is nilpotent of class at most $c$.  Its cyclotomic character
has infinite, in particular nontrivial, image.
If $a\in E^\times$ has zero rational $p$-Kummer class, its
local class on $J$ restricts to zero on $H$, so inflates
from $H^1(\Gamma,\mathbb Q_p(1))=0$.
Restriction from $G_M$ to $J$ is injective because
$\mathbb Q_p(1)^J=0$.  The local valuation and principal-unit
logarithm make $a$ a root of unity.  Conversely every root of
unity has all its divisions in $K$.  We have thus proved, for
all finite $E/K$,
\[
                         \ker\kappa_E=\mu_\infty.
\]
For an equality $\kappa_E(b)=\alpha\kappa_E(p)$, use the
same argument on the difference of the two classes.  The local
equality follows, and the normalized local valuation gives
\[
                     \alpha=\frac{v_M(b)}{v_M(p)}\in\mathbb Q.
\]
No Baker--Brumer theorem or finite tame root of unity is needed
for these two conclusions at a bounded level.

There is also no complex conjugation centralizing $G_E$, or
any of its open subgroups.  To check nontriviality rather than
assume it, $H$ contains $\gamma_{c+1}(P_M)$.
This subgroup is nonabelian: in a quotient
$\operatorname{UT}_{2c+3}(\mathbb F_p)$ of the countably
infinite-rank free group $P_M$, its image contains
$1+E_{1,c+2}$ and $1+E_{c+2,2c+3}$, whose commutator is
$1+E_{1,2c+3}\ne1$.
Thus $KD$ is not separably closed.  If complex conjugation
centralized $H$, the same F. K. Schmidt argument as in
Claim~B.120 would put it inside $D_p$, contrary to the
torsion-freeness in Claim~B.119.  Repeating this after any
finite extension proves the required open-subgroup version.

These are exactly the three arithmetic inputs used in
Claims~B.107--B.110: the stable Kummer kernel is roots of
unity, proportionality to the actual Kummer value of $p$
forces a rational scalar, and complex conjugation does not
centralize any relevant open arithmetic group.
For completeness, their uses are as follows.  The curve norm
residue argument uses the scalar statement to make every boundary
residue rational.  The root-of-unity kernel then makes actual
norm absolute values finite products of algebraic functions to
rational powers, with one uniform positive algebraic constant.
Finite Gaussian interpolation recovers their monic coefficient
polynomials as real analytic algebraic functions.  The differential
identities and algebraic continuation separate holomorphic and
antiholomorphic factors; the last arithmetic input excludes the
latter.  The finite-projection degree argument makes the remaining
correspondence a regular graph and treats its exceptional points
on the original source.  None of these steps uses a property
specific to $K_{p,\mathrm{nil}}$ beyond those three inputs and
$K\subset\overline{\mathbb Q}$, both now verified here.
Finally Carlson--Haine--Wolf, Theorem~4.18, gives full pinned
reconstruction from that actual-point theorem, and Corollary~2.32
identifies the localization by universal homeomorphisms.
Those source hypotheses were checked in Claim~B.110 and apply
unchanged to the present characteristic-zero fields and every
finite extension.  This completes the asserted bounded-level
reconstruction, including $p=2,3$.
\end{proof}
\paragraph{Check record.} checked only in this writing, 2026-11-18: GPT-6 Astra (OpenAI), ChatGPT Pro mode; internal effort value not exposed. Own detailed proof written and reread; primary source hypotheses matched. Root full second reading of META.5 passed after explicitly including each tested a or b in the finite Galois field D


\begin{claim}{B.126}{Proved}{Two translated grids give full arithmetic equivariance of the norm data}
Under the hypotheses of Claim~B.122, one can choose $D$ and $V$
so that, with $F=F_pD$,
\[
            P_{\sigma t}(Z)=\sigma\bigl(P_t(Z)\bigr)
       \qquad\bigl(\sigma\in G_F,\ t\in V(\overline{\mathbb Q})\bigr).
\]
Here the action on a polynomial is coefficientwise.  In particular,
if the source is $\mathbb A^n_E$, its entire actual point map becomes
$G_{F_pD}$-equivariant for one finite $D$, after enlarging $D$
finitely.  This is an actual-point descent statement; it does not
assert descent of the geometric morphism of topoi.
\end{claim}
\begin{proof}
Choose $a_0\in B_p\setminus F_p$.  Such an element exists:
$I'$ contains the nontrivial derived subgroup of wild inertia,
is itself a closed subgroup of free wild pro-$p$ inertia,
and is therefore a nontrivial free pro-$p$ group.  Its
abelianization is nonzero, so $I''\ne I'$.
Use the two target grids
\[
 \mathcal A_0=\{r+is:0\leq r,s\leq4d\},\qquad
                  \mathcal A_1=a_0+\mathcal A_0.
\]
Include $a_0$ in $D$, together with the finitely many norm-class
and open-set coefficients for both grids.  The proof of
Claim~B.122 still applies: every $\sigma\in I'\cap G_D$
fixes every point of both grids.  The interpolation argument for
one of them already proves $P_t\in F[Z]$ for $t\in V(F)$.

Fix $t_0\in V(F)$.  In (META.N), evaluating at $t_0$ gives
\[
          \beta_a=\kappa(b_a(t_0))-
                                \sum_j\lambda_{aj}\kappa(h_j(t_0)).
\]
All displayed algebraic arguments now belong to $F$: this uses the
norm coefficient descent just proved.  Thus $\beta_a$ is
$G_F$-invariant as a cohomology class over $E_1=B_pD$.
The $h_j$ have coefficients in $D$, so the whole class $\xi_a$
is invariant.  Evaluating at arbitrary algebraic $t$ and transporting
by $\sigma\in G_F$ consequently gives, after one common finite
extension of $E_1$,
\[
            \frac{P_{\sigma t}(a)}{\sigma P_t(a)}\in W
                       \qquad(a\in\mathcal A_0\cup\mathcal A_1).
                                                        \tag{META.2G}
\]
This uses the stable kernel (META.K); it does not assume that $t$
is fixed by $\sigma$.

Put $R(Z)=P_{\sigma t}(Z)/\sigma P_t(Z)$, a quotient of two
monic degree-$d$ polynomials.  For each $\tau\in I'$, membership
in $W$ implies that $\tau(R(a))/R(a)$ is a root of unity.
On the first Gaussian grid, $\tau$ fixes $a$, so the same
absolute-square interpolation as in (META.G) gives $\tau R=R$.
Indeed cross multiplication involves degree-$2d$ monic
polynomials, and the $(4d+1)^2$ values determine their
absolute-square polynomial.  Equality of absolute values gives
equality of the monic products, even if common factors occur.
Thus
\[
                              R(Z)\in F_p(Z).
\]
For the second grid write $a=a_0+z$, with
$z\in\mathcal A_0$.  Applying $\tau$ to (META.2G) now gives
\[
       |(\tau R)(\tau a_0+z)|=|R(a_0+z)|.
\]
The identical interpolation argument applied to the rational
function $S(Z)=R(a_0+Z)$ gives $\tau S=S$ for every $\tau\in I'$.
Consequently $R(a_0+Z)\in F_p(Z)$ as well.  No invariance of
$a_0$ under $I'$ has been assumed.

These two rationality assertions force $R$ to be constant.
If it were nonconstant, its Laurent expansion at infinity,
where its value is one, would have the form
\[
 R(Z)=1+cZ^{-m}+c'Z^{-m-1}+O(Z^{-m-2}),
               \qquad m\geq1,\quad 0\ne c\in F_p,\quad c'\in F_p.
\]
The coefficient of $Z^{-m-1}$ in $R(a_0+Z)$ is
$c'-ma_0c$.  Its membership in $F_p$ would give $a_0\in F_p$,
a contradiction.  Hence $R=1$, proving full equivariance of the
norm-polynomial assignment.

For the assertion on the entire affine-space point map, take the
identity projection, so $d=1$.  Remove the finitely many closed
sets $S_a=f^{-1}(a)$ for the two grids and cover their complement
by finitely many nonempty principal opens.  Perform the preceding
construction on each such open and take a common finite $D$
containing the coefficients of every $S_a$ as well.
The actual function is $G_F$-equivariant on that complement.
Each $S_a$ is $G_F$-stable because its equations have coefficients
in $D$, and its value $a$ also lies in $D$.
Thus $f(\sigma x)=a=\sigma f(x)$ on $S_a$ too.
This proves the asserted global actual-point equivariance.
\end{proof}
\paragraph{Check record.} checked only in this writing, 2026-11-18: GPT-6 Astra (OpenAI), ChatGPT Pro mode; internal effort value not exposed. Own detailed proof written and reread; primary source hypotheses matched. Root full second reading of META.6 passed: invariant normalized beta, finite common extensions, both grids, monic absolute-square interpolation, Laurent contradiction, and finite-open/fibre patching


\begin{claim}{B.127}{Proved}{Actual evaluations detect every Kummer class on a principal affine open}
Let $E/B_p$ be finite and let
\[
             V=\operatorname{Spec}E[X_1,\ldots,X_n,1/h]
\]
be nonempty.  If $\xi\in H^1(V,\mathbb Q_p(1))$ has zero
evaluation at every $x\in V(L)$ for every finite extension
$L/E$, then $\xi=0$.  Consequently two pinned relative morphisms
from $V$ to $\mathbb G_m$ with the same actual point map pull
back the same rational $p$-Kummer class.

For a pinned morphism with affine-space source, choose the finite
field $D$ of Claim~B.126 and put $F=F_pD$.  For every
$a\in\overline{\mathbb Q}$, write
$U_a=\mathbb A^n\setminus f^{-1}(a)$ and
$\xi_a=\phi^*\kappa(Y-a)$ on its nonempty domain.  Then the
actual equivariance of Claim~B.126 also gives exact
cohomological equivariance
\[
                    \sigma\xi_a=\xi_{\sigma a}
                         \qquad(\sigma\in G_F),
                                                        \tag{META.EV}
\]
after the finite scalar extensions needed to define these classes.
No ambiguity from classes with identically zero actual evaluations
occurs in this assertion.
\end{claim}
\begin{proof}
The case $n=0$ is evaluation at the single rational point, so
assume $n>0$.  A finite scalar extension suffices to split the
geometric irreducible factors of $h$.  Restriction on rational
$p$-cohomology is injective: its composite with corestriction is
multiplication by the nonzero extension degree.  We may therefore
make this scalar extension.  Choose a finite Galois number field
$D_0$ containing all coefficients, these factors, and generators
of the original $E/B_p$, and replace $E$ by $E_0=B_pD_0$.
Let $h_1,\ldots,h_r$ be the distinct nonconstant geometric factors.

The localized polynomial ring is factorial, with unit group
$E_0^\times\times\mathbb Z^r$, generated modulo constants by
the $h_j$.  Its Picard group is zero.  The Kummer sequences at
each $p$-power level, followed by inverse limit and tensoring
with $\mathbb Q_p$, consequently give
\[
              \xi=\beta+\sum_{j=1}^r\lambda_j\kappa(h_j),
 \qquad \beta\in H^1(E_0,\mathbb Q_p(1)),\quad
                         \lambda_j\in\mathbb Q_p.
                                                        \tag{META.E1}
\]
There is no inverse-limit correction: $E_0$ contains every
$p$-power root of unity, and the transition maps on these
degree-zero groups are surjective.
Choose $t_0\in\mathbb Q^n$ avoiding all $h_j=0$;
$\mathbb Q^n$ is Zariski dense over $\overline{\mathbb Q}$,
by induction on the number of variables.  Evaluation at $t_0$
gives
\[
                \beta=-\sum_j\lambda_j\kappa(h_j(t_0)).
                                                        \tag{META.E2}
\]
Thus the possibly nonalgebraic arithmetic class in (META.E1)
has become an explicit finite $\mathbb Q_p$-linear combination
of actual Kummer classes.

Choose a nonidentity element $g\in I''\cap G_{D_0}$.
Such an element exists: for $M_0=(D_0)_v$, the nonabelian
group $P_{M_0}''$ constructed in Claim~B.119 is contained
in $I''\cap G_{D_0}$.  The element $g$ fixes $D_0$ and has
$\chi_p(g)=1$, since the cyclotomic character is abelian.
There is $b\in\overline{\mathbb Q}^\times$ such that
\[
                         g(b)/b\notin\mu_\infty.
                                                        \tag{META.E3}
\]
Indeed, choose an algebraic $a$ moved by $g$ and put $b=a+c$
with $c\in\mathbb Q$.  The ratios $(g(a)+c)/(a+c)$ all lie
in one number field, whose group of roots of unity is finite.
Each root of unity other than one can occur for at most one
$c$, and one cannot occur because $g(a)\ne a$.
Avoiding these finitely many $c$ and the zero denominator
proves (META.E3).

Let $D_1/\mathbb Q$ be finite Galois containing $D_0$ and $b$;
it also contains $g(b)$.  Put
\[
 M=(D_1)_v,\qquad J=I\cap G_{D_1}=I_M,\qquad
 H=I''\cap J=G_{B_pD_1}.
\]
The field $M/\mathbb Q_p$ is finite Galois.  In particular $g$
acts on $M$ and normalizes $J$ and $H$, although it need not fix
$b$.  For any $u\in\mathbb Q^n$ and every sufficiently large
integer $m$, the point
$t_m=t_0+p^m b u$ lies in $V(D_1)$.  The class
\[
        \theta_m=\sum_j\lambda_j
              \kappa_M\bigl(h_j(t_m)/h_j(t_0)\bigr)
               \in H^1(G_M,\mathbb Q_p(1))
\]
restricts to zero on $H$, by (META.E2) and the assumed
vanishing of all evaluations.  Its restriction to $J$
therefore inflates from $H^1(J/H,\mathbb Q_p(1))$.
Conjugation by $g$ acts trivially on $J/H$: for $j\in J$,
the commutator $[g,j]$ belongs to $I''$ because $g\in I''$,
and to $J$ because $J$ is normal in $I$.
It also acts trivially on the coefficient module, as
$\chi_p(g)=1$.  Hence $g\theta_m|_J=\theta_m|_J$.

Restriction from $H^1(G_M,\mathbb Q_p(1))$ to
$H^1(I_M,\mathbb Q_p(1))$ is injective: the inertia
invariants of $\mathbb Q_p(1)$ vanish, and
inflation--restriction applies.  We conclude $g\theta_m=\theta_m$.
Kummer functoriality and the local principal-unit logarithm
therefore give, for all sufficiently large $m$,
\[
 \sum_j\lambda_j\log
 \frac{h_j(t_0+p^m g(b)u)}{h_j(t_0)}
 =
 \sum_j\lambda_j\log
 \frac{h_j(t_0+p^m b u)}{h_j(t_0)} .
                                                        \tag{META.E4}
\]
All logarithms are the convergent logarithm of an element
near one in the single finite local field $M$.
The action is compatible with this logarithm because
$\log(gz)=g\log(z)$; its power-series coefficients are
rational, and $g$ fixes $\mathbb Q_p$ and all $\lambda_j$.
It also fixes the $h_j$ and $t_0$ by the choice of $D_0$.

For this fixed $u$ the power series
\[
 F_u(Z)=\sum_j\lambda_j
          \log\frac{h_j(t_0+Zu)}{h_j(t_0)}\in M[[Z]]
\]
converges near zero and has constant term zero.
Equation (META.E4) says that
$F_u(g(b)p^m)-F_u(bp^m)=0$ for infinitely many $m\to\infty$.
A convergent nonzero power series cannot have such zeros:
its first nonzero term dominates all higher terms sufficiently
close to zero.  Therefore
$F_u(g(b)Z)=F_u(bZ)$ as power series.
For every $k\ge1$, (META.E3) gives $g(b)^k-b^k\ne0$,
so comparison of the coefficient of $Z^k$ shows $F_u=0$.

Apply this for every $u\in\mathbb Q^n$ to the multivariable
convergent series
\[
 F(\mathbf Z)=\sum_j\lambda_j
             \log\frac{h_j(t_0+\mathbf Z)}{h_j(t_0)}.
\]
Its homogeneous part of each degree vanishes on
$\mathbb Q^n$, hence is the zero polynomial.  Thus $F=0$.
Differentiating its formal expansion gives the rational-form
identity
\[
                          \sum_j\lambda_j\frac{dh_j}{h_j}=0.
\]
The expansion map at $t_0$ is injective on rational functions
whose denominators are nonzero there.  At the generic point
of the prime divisor $h_j=0$, the residue of the displayed
form is $\lambda_j$: $h_j$ has order one and every other
$h_k$ is a unit.  Hence all $\lambda_j=0$, and (META.E2)
then gives $\beta=0$.  This proves the injectivity assertion.

The consequence for two pinned morphisms follows by subtracting
their pulled-back classes and using actual point evaluation.
For (META.EV), Claim~B.126 makes the actual evaluations of
$\sigma\xi_a$ and $\xi_{\sigma a}$ identical, on the same
open domain and over every finite field of point definition.
Restrict their difference to any nonempty principal open in
that domain and apply the proved assertion.  Restriction from
the domain to this open is injective on $H^1(-,\mathbb Q_p(1))$:
the source is normal, so the map on fundamental groups is
surjective, and the coefficient character is constant here.
The equality follows on the whole domain, proving (META.EV).
\end{proof}
\paragraph{Check record.} checked only in this writing, 2026-11-18: GPT-6 Astra (OpenAI), ChatGPT Pro mode; internal effort value not exposed. Own detailed proof written and reread; primary source hypotheses matched. Root full second reading of META.7 passed: finite scalar restriction/corestriction, unit decomposition and lim1, non-root-unit dilation, one Galois completion, trivial quotient/coefficient action, local restriction injection, power-series identity and differential residues


\begin{claim}{B.128}{Proved}{The norm data descend to one local arithmetic group}
Under the hypotheses of Claim~B.122, the finite number field
$D$ can be chosen so that the norm-polynomial assignment satisfies
\[
 P_{\sigma t}(Z)=\sigma(P_t(Z))
 \qquad
 \bigl(\sigma\in D_p\cap G_D,\ t\in V(\overline{\mathbb Q})\bigr).
                                                        \tag{META.LOC}
\]
Here $D/\mathbb Q$ is finite Galois, $M=D_v$ is its completion
at the chosen $p$-adic place, and $D_p\cap G_D\cong G_M$.
Consequently every pinned relative morphism
$\mathbb A^n_{E,\mathrm{et}}\to\mathbb A^1_{E,\mathrm{et}}$,
with $E/B_p$ finite, has actual point map equivariant under one
such open local absolute Galois group.  Its pulled-back rational
$p$-Kummer classes satisfy (META.EV) for this larger group as well.
This asserts neither descent of the whole geometric morphism nor
point-theoretic lifting of every normalized class to $M$.
\end{claim}
\begin{proof}
Use the two grids of Claim~B.126, including
$a_0\in B_p\setminus F_p$.  First choose a finite Galois number
field $D_0$ containing their points, the coefficients of $V$ and
all $h_j$ in the finitely many decompositions (META.N), and
generators of the given finite extension $E/B_p$.
Choose $t_0\in V(D_0)$.  For each grid point $a$, its norm
evaluation $b_a(t_0)=P_{t_0}(a)$ is an actual nonzero algebraic
number.  Enlarge $D_0$ to one finite Galois number field $D$
containing all these finitely many numbers.  Evaluation gives
\[
       \beta_a=\kappa(b_a(t_0))-
                            \sum_j\lambda_{aj}\kappa(h_j(t_0)).
                                                        \tag{META.LOC1}
\]
Every algebraic argument in (META.LOC1), and every coefficient
of $h_j$, belongs to $D$.  No rationality of
$\lambda_{aj}\in\mathbb Q_p$ is needed.

The field $B_p$ is stable under $D_p$: inertia is normal in the
local decomposition group, and $I''$ is characteristic in $I$.
Thus $D_p\cap G_D$ acts on $E_1=B_pD$ and on the cohomology
of $V_{E_1}$.  Each term in (META.LOC1) is fixed under this
action, and so is $\xi_a$.  This includes the full action on
$\mathbb Q_p(1)$: the Kummer map is Galois equivariant, so
the class of an element fixed by $\sigma$ is fixed even when
$\chi_p(\sigma)\ne1$; the scalar $\lambda_{aj}$ itself is
fixed as an element of $\mathbb Q_p$.

For $\sigma\in D_p\cap G_D$ and arbitrary algebraic $t$,
evaluation and transport now give
\[
                   \frac{P_{\sigma t}(a)}{\sigma P_t(a)}\in W
             \qquad(a\in\mathcal A_0\cup\mathcal A_1),
\]
over a common finite extension of $E_1$.  The stable kernel
(META.K) applies to this extension.  The rest of the two-grid
proof of Claim~B.126 applies verbatim to
$R=P_{\sigma t}/\sigma P_t$.  Its auxiliary automorphisms
$\tau\in I'$ act on $W$ independently of the chosen $\sigma$.
The first grid makes $R\in F_p(Z)$, the translated grid makes
$R(a_0+Z)\in F_p(Z)$, and the nonzero first Laurent coefficient
of a hypothetical nonconstant $R$ would imply $a_0\in F_p$.
Thus $R=1$, proving (META.LOC).

For an affine-space source use the identity projection and the
finite principal-open patching of Claim~B.126.  Include in
$D$ the equations of the finitely many excluded grid fibres
and their values.  These fibres are then stable under
$D_p\cap G_D$, and the same conclusion holds on them as on
their complement.  A constant point map is handled simply by
including its algebraic value in $D$.
Finally Claim~B.127 converts this actual equivariance into
the stated exact cohomological equivariance.  The fixed field
of $D_p\cap G_D$ in $\overline{\mathbb Q}$ is the chosen
henselization $D^h$, embedded in $M$; hence the point assertion
also says that actual values on $D^h$-points belong to $D^h$.
It does not turn the original topos morphism over $B_pD$ into
a morphism over $D^h$.
\end{proof}
\paragraph{Check record.} checked only in this writing, 2026-11-18: GPT-6 Astra (OpenAI), ChatGPT Pro mode; internal effort value not exposed. Own detailed proof written and reread; primary source hypotheses matched. Root full second reading of META.8 passed: finite actual base values, B\_p stability under D\_p, full Tate-coefficient transport, arbitrary-point kernel test, independent I-prime two-grid argument, finite patching, cohomological consequence, henselization fixed field


\begin{claim}{B.129}{Proved}{Universal projections give local equivariance on every smooth affine source}
Let $E/B_p$ be finite, let $X/E$ be smooth, geometrically
integral and affine, and let
$\phi:X_{\mathrm{et}}\to\mathbb A^1_{E,\mathrm{et}}$
be pinned and relative.  Its actual point map
$f:X(\overline{\mathbb Q})\to\overline{\mathbb Q}$
is equivariant under $D_p\cap G_D\cong G_{D_v}$ for one
finite Galois number field $D$ containing a field of definition
for $X$ and generators of $E/B_p$.
The choice of $D$ works for every point and every element of
this local group simultaneously.  It does not assert that
$\phi$ descends to the henselization $D^h$.
\end{claim}
\begin{proof}
A constant actual map is handled by adjoining its value to $D$,
and a zero-dimensional geometrically integral smooth source is
$\operatorname{Spec}E$.  We may assume $n=\dim X>0$ and that
$f$ is nonconstant.  Fix a closed embedding $X\subset\mathbb A^m$
over a finite field of definition.  Let $\overline X\subset
\mathbb P^m$ be its projective closure, and put
$X_\infty=\overline X\cap\{X_0=0\}$, of dimension at most $n-1$.
We first work with the full affine space of $n\times m$ matrices
$A$.  The condition
\[
                 \mathbb P(\ker A)\cap X_\infty=\varnothing
                                                        \tag{META.UP1}
\]
defines a nonempty open parameter locus.  Openness follows by
projecting the closed incidence correspondence with the projective
scheme $X_\infty$.  For a fixed point of $X_\infty$, the condition
that its vector be killed by $A$ imposes $n$ independent linear
conditions on $A$.  The incidence therefore has dimension at most
$nm+(n-1)-n=nm-1$, proving nonemptiness of the complement.
Choose a nonempty principal open $S$ inside this locus.

For $s=A\in S$, write $g_s(x)=Ax$.  The universal map
\[
 G:X\times S\longrightarrow\mathbb A^n\times S,
                    \qquad(x,s)\longmapsto(g_s(x),s)
                                                        \tag{META.UP2}
\]
is finite.  Indeed (META.UP1) makes
$[X_0:\mathbf X]\mapsto[X_0:A\mathbf X]$ a morphism
$\overline X\times S\to\mathbb P^n\times S$ with no
base points.  It is proper, and its inverse image of
$\mathbb A^n\times S$ is exactly $X\times S$.
Consequently (META.UP2) is proper; it is also affine, and a
proper affine finite-type morphism is finite.

This finite map is flat of one constant positive degree $d$.
Both its source and target are smooth of dimension $n+nm$.
At a point of the target, its regular local ring has the same
dimension as each local source ring above it: the residue-field
extensions are finite, so the dimension formula for these smooth
varieties gives equality.  Each source ring is Cohen--Macaulay.
A regular system of parameters in the target is therefore a
system of parameters, hence a regular sequence, in each source
ring.  The Koszul criterion for a finite module over a regular
local ring makes the resulting finite module free.  Explicitly,
the regular sequence makes its positive Koszul homology zero;
in particular $\operatorname{Tor}_1(-,k)=0$, and Nakayama's
lemma applied to a lifted minimal basis gives freeness.
Thus the map is finite flat.  Its rank is constant because
$\mathbb A^n\times S$ is connected.  It is positive because
each finite map $g_s$ has a nonempty $n$-dimensional source and
thus has full-dimensional closed image in $\mathbb A^n$.

Choose the two grids of Claim~B.126 for this degree $d$.
For each grid point $a$, let $C_a=f^{-1}(a)$, with its reduced
closed structure.  It is a proper closed subset of $X$.
After a finite scalar enlargement, the finitely many $C_a$ and
the universal projection data have a common finite field of
definition.  The sets $G(C_a\times S)$ are closed and proper:
$G$ is finite, and $\dim(C_a\times S)\le n+nm-1$.
The branch locus of (META.UP2) is closed and proper as well,
since the finite dominant map is generically separable in
characteristic zero.  Let $T^\circ\subset\mathbb A^n\times S$
be the complement of these finitely many closed sets.
It is nonempty.  Cover it by finitely many principal opens
$V_1,\ldots,V_q$ in $\mathbb A^{n+nm}$, each contained in
$T^\circ$.

Precompose $\phi$ with the scheme projection
$X\times S\to X$.  This gives a pinned morphism whose actual
point map is $(x,s)\mapsto f(x)$; no product of arbitrary
topos morphisms is being assumed.  Over each $V_j$ the map
$G^{-1}(V_j)\to V_j$ is finite \'etale of degree $d$ and
its fibres avoid every grid fibre.  Corestrict the pulled-back
grid Kummer classes.  Their actual evaluations are precisely
the norm-polynomial values
\[
             P_{s,t}(Z)=\prod_{x\in g_s^{-1}(t)}(Z-f(x)).
                                                        \tag{META.UP3}
\]
The proof of Claim~B.128 uses only a finite \'etale map
over a principal affine open, its actual norm evaluations,
and the Kummer decomposition on that principal open.
Those data are exactly present here; an extension of the map
to the whole ambient affine space is not used in that proof.
Apply its normalization and two-grid argument on every $V_j$.
There are finitely many opens and grid classes.  Thus one
finite Galois number field $D$ contains all their coefficients
and actual base values, together with equations for $X,S,C_a$
and the grid values themselves.  For
$\sigma\in D_p\cap G_D$ we obtain, simultaneously on all
of $T^\circ$,
\[
                    P_{\sigma s,\sigma t}
                                  =\sigma P_{s,t}.
                                                        \tag{META.UP4}
\]
This is the uniformity supplied by the product source: the
field $D$ is fixed before any individual point or projection
parameter is subsequently chosen.

We record the exact elementary projection fact needed at such
a point.  Fix $x\in X(\overline{\mathbb Q})$ and a proper
closed subset $C\subset X_{\overline{\mathbb Q}}$ not
containing $x$.  There is a nonempty open set of matrices
$s\in S$ for which the whole fibre $g_s^{-1}(g_s(x))$ is
reduced and avoids $C$.
For avoidance, the incidence
\[
       \{(y,A):y\in C,\ A(y-x)=0\}
\]
has dimension at most $nm+\dim C-n\le nm-1$, since $y\ne x$.
Its image closure is therefore a proper closed subset of
the matrix space.  For reducedness away from $x$, consider
the same incidence with $y\in X\setminus\{x\}$.
It is a vector bundle of rank $nm-n$ over this smooth variety,
and hence is smooth of dimension $nm$.
Its map to matrix space, restricted over $S$, is quasifinite
because every $g_s$ is finite.  If it is not dominant, exclude
the closure of its image.  If it is dominant, its generic
fibre is reduced and separable: it is a zero-dimensional
localization of a regular scheme over a characteristic-zero
field.  The non-étale locus therefore has image with proper
closure: it is a proper closed subset of the irreducible smooth
$nm$-dimensional incidence, so it and its image closure have
dimension at most $nm-1$.  Excluding that closure makes every remaining
point of the incidence fibre \'etale over parameter space.
This is exactly nonsingularity of the linear intersection
at every $y\ne x$, since its relative tangent space is
$\ker(A|_{T_yX})$.
At $x$ itself, the condition that $A|_{T_xX}$ be invertible
is nonempty and open.  Intersecting these finitely many
nonempty opens proves the projection fact.

Now fix $\sigma\in D_p\cap G_D$ and a point $x$.
If $f(x)$ belongs to a grid, equivariance already holds:
the corresponding $C_a$ and $a$ are defined over $D$ and
are fixed by $\sigma$.  The same observation excludes
$f(\sigma x)$ belonging to a grid when $f(x)$ does not.
Suppose, for a contradiction in the remaining case, that
\[
                b=\sigma^{-1}(f(\sigma x))\ne f(x).
\]
The proper closed set
$C=f^{-1}(b)\cup\bigcup_a C_a$ does not contain $x$.
The projection fact supplies a nonempty open subset of $S$
whose fibres through $x$ are reduced and avoid $C$.
It has a matrix $s$ with rational entries: rational affine
points are Zariski dense even over $\overline{\mathbb Q}$.
In particular $\sigma s=s$.
The point $(s,g_s(x))$ lies in $T^\circ$, so (META.UP4)
applies and $g_s(\sigma x)=\sigma g_s(x)$.
The value $f(\sigma x)$ is a root of
$P_{s,\sigma g_s(x)}=\sigma P_{s,g_s(x)}$.
Hence $b$ is a root of $P_{s,g_s(x)}$, and there is some
$y\in g_s^{-1}(g_s(x))$ with $f(y)=b$.
This contradicts the choice of a fibre avoiding $C$.
Thus $f(\sigma x)=\sigma f(x)$ for every $x$ and every
$\sigma$ in the one fixed local group, as asserted.
\end{proof}
\paragraph{Check record.} checked only in this writing, 2026-11-18: GPT-6 Astra (OpenAI), ChatGPT Pro mode; internal effort value not exposed. Own detailed proof written and reread; primary source hypotheses matched. Root full second reading of META.9 passed: universal projective/proper-affine finiteness, regular-sequence finite flatness, finite principal cover, one normalized field before point choices, incidence avoidance and reducedness, rational fixed parameters and multiset contradiction; explicit bad-incidence dimension join inserted


\paragraph{The intermediate descent obstruction and its later resolution.}
Claims~B.121--B.122 and B.126 give restrictions on actual point
maps and uniform $G_{F_pD}$-equivariance of every norm-polynomial
assignment.  For an affine-space source, this includes the entire
actual point map for one finite $D$.  By themselves they do not settle B8 for
$B_p$.  Claim~B.127 additionally proves that actual evaluations
detect all rational $p$-Kummer classes on principal affine opens;
the resulting exact cohomological equivariance has no remaining
ambiguity from a class invisible at every point.  Claim~B.128
strengthens these equivariance statements to one open local absolute
Galois group, after including finitely many actual base values in
the field of definition.  Claim~B.129 obtains one such group
for every smooth geometrically integral affine source, by a
universal family of finite projections.  In particular,
(META.N) can be made to
descend to $F_pD$ after normalizing by an actual value at $t_0$
and enlarging $D$ to contain the resulting finite algebraic data.
Nevertheless these intermediate conclusions do not identify its evaluation over $F_pD$ with
the actual Kummer class there: their difference may lie in
\[
 \ker\bigl(H^1(F_pD,\mathbb Q_p(1))\longrightarrow
                         H^1(B_pD,\mathbb Q_p(1))\bigr).
\]
This kernel contains the Kummer classes of $F_p^\times$ and is
not zero.  Thus even field-rationality and uniform arithmetic
equivariance of the norm coefficients are not silently promoted to
a point-theoretic lift of every descended class,
a descended morphism of topoi, or a regular morphism of schemes.
Claim~B.136 later resolves the full metabelian reconstruction
assertion by a separate norm argument: cyclic local quotients
and exact logarithmic derivatives recover the reduced norm,
and the two-grid test removes the common factor.
The displayed kernel remains a real obstruction to the
unjustified intermediate descent inference; the later proof
does not assume that kernel is zero.

\subsection{Complete pointed families with no free fixed norm vector}

\begin{claim}{B.130}{Proved}{A primitive-zero prime-power cover excludes every pointed free norm vector}
Let $C/k$ be a smooth proper geometrically connected curve of
genus at least two over a finite field of characteristic $p$.
Suppose $C_{\bar k}$ has a connected cyclic \'etale cover
$D\to C_{\bar k}$ of degree $\ell^n$, where $\ell\ne p$ is
a prime and $n\ge1$, whose primitive Jacobian factor has
geometric $p$-rank zero.  For every finite extension $k'/k$
and every $P\in C(k')$, the fixed norm module $M_P^F$ has
no vector with free deck orbit at every geometric normal level.
Here $F$ is the arithmetic Frobenius of $k'$ through $P$.

In particular, this answers B2 affirmatively for every pointing
and every finite field of definition of either of the following
geometric genus-two families in characteristic two:
\begin{enumerate}
\item the ordinary models
\[
 z^2+z=ax+\frac b x+\frac c{x+1},\qquad
 abc\ne0,\qquad (a+b+c)^3=abc;
\]
\item the two-rank-one models
\[
 z^2+z=x^3+ax+\frac b x,\qquad b\ne0,\qquad ba^{16}=1.
\]
\end{enumerate}
Geometric isomorphism to one of these models suffices; its
displayed coefficients need not belong to the initial field.
The covers in these two cases have degrees three and five,
respectively.  No assertion is made here for the remaining
genus-two loci.  The case of a degree-one cover is excluded.
\end{claim}

\begin{proof}
Let $G=\langle\delta\rangle$ be the cyclic group of order
$\ell^n$, let $T\subset G$ be its subgroup of order $\ell$,
and put $D_0=D/T$.  Both $D$ and $D_0$ are geometrically
normal covers of $C$.  The new rational character summands
for $D\to D_0$ are exactly the characters of exact order
$\ell^n$: a character is trivial on $T$ precisely when it
factors through $G/T$.  The averaging idempotent of $T$
therefore shows that
\[
 B=\bigl(\ker(J_D\longrightarrow J_{D_0})\bigr)^0
\]
is isogenous to the primitive factor in the hypothesis.
Thus $B$ has geometric $p$-rank zero.  Equivalently
$B(\bar k)[p^\infty]=0$; this is an assertion about
geometric points of its group schemes, not about their
possibly nontrivial connected parts.

Fix any $k'$ and $P$ as in the statement.  The curve cover,
its cyclic action, and the quotient all descend to a finite
extension of $k'$.  Choose a geometric point $P_D$ above
$P$ and enlarge that finite field further until $P_D$ and
all deck transformations are rational.  Call the resulting
field $L$.  Its image $P_0$ on $D_0$ is rational too.
This is a permitted enlargement of the chosen geometric
cover; the zero geometric $p$-rank hypothesis is unchanged.
In particular it holds for the actual pointed descent just
chosen, independently of its possible constant twist.

Pullback and norm for $\pi:D_L\to(D_0)_L$ satisfy
$\pi_*\pi^*=[\ell]$.  Since $\ell$ is prime to $p$,
the averaging idempotent on the geometric $p$-divisible
groups splits off the old factor.  Its complementary
\'etale part is that of $B$, hence is zero.  On finite
rational $p$-primary groups it follows that
\[
 J_D(L)[p^\infty]=
                    \pi^*J_{D_0}(L)[p^\infty].
                                                        \tag{NM.Z1}
\]
Equivalently, the pullback--norm isogeny gives
$\#J_D(L)/\#J_{D_0}(L)=\#B(L)$, which is prime to $p$.
The same conclusions hold over every extension of $L$.

Suppose $m\in M_P^F$ had free orbit at every normal level.
The chosen pointing identifies the displayed covers with
levels of its geometric pointed tower.  Since
$m$ is fixed by $F$, it is fixed by $F^{[L:k']}$.
Its coordinate $m_D$ on $D$ consequently belongs to
$J_D(L)[p^\infty]$.  By (NM.Z1), every element of the
nontrivial subgroup $T$ fixes $m_D$.  This contradicts
freeness at the normal level $D\to C$.  This is the
specific normal-step obstruction of B.9, with its transfer
from a finite constant extension spelled out; no
point-separation theorem is used in this deduction.

For the ordinary family, B.85 constructs a connected cyclic
three cover with primitive two-rank zero.  For the
two-rank-one family, B.86 constructs the analogous cyclic
five cover.  Their full cover and $p$-rank proofs have
been checked above.  A geometric model and its cover
descend after a finite constant extension, so the argument
applies to any finite-field form of either model and every
choice of $P$.  Both families are nonempty: $a=b=c=1$
satisfies the first equation, and $a=b=1$ satisfies the
second.  Since $n\ge1$ was essential to the nontrivial
subgroup $T$, no conclusion from a degree-one primitive
factor is being inserted into this argument.
\end{proof}
\paragraph{Check record.} Checked only in this writing, 18November2026: GPT-6 Astra (OpenAI), ChatGPT Pro mode. Root complete proof written and read; finite-checks reader separately checked the prime-power primitive factor, arbitrary pointing and finite-field transfer. Separate effort setting not exposed.


\subsection{Local analytic recovery over the metabelian field}

The arguments in this subsection retain the notation $B_p,F_p,W,I,D_p$
of the preceding subsection.  The word ``analytic'' always refers to
convergent power series over a finite extension of $\mathbb Q_p$.
It is not a continuity assertion about an initially arbitrary point map.

\begin{claim}{B.131}{Proved}{Pinned point maps are continuous for the henselian valuation}
Let $E/B_p$ be finite and let $\phi:X_{\mathrm{et}}\to
Y_{\mathrm{et}}$ be a pinned relative morphism between finite-type
$E$-schemes.  For every finite extension $E'/E$, its actual point
map $X(E')\to Y(E')$ is continuous for the topology induced by
the chosen $p$-adic valuation.
If this point map is equivariant under $G_M$ for a finite local
field $M$ with $M\cap\overline{\mathbb Q}\subset E$, its
restriction to points in any finite extension
$L/M$ that are algebraic over $\mathbb Q$ is continuous for the
subspace $p$-adic topology.  This statement alone does not extend
the map to points of the completed field $L$.
\end{claim}
\begin{proof}
For a field $k$, an \'etale-open subset of $Y(k)$ is an image
of $Y'(k)$ for an \'etale morphism $Y'\to Y$; these images form
a basis.  We may take $Y'\to Y$ quasicompact and separated.
For finite-type schemes, inverse image by a geometric morphism
of \'etale topoi carries such an object to an object represented
by a quasicompact separated \'etale scheme $X'\to X$
\citep[Recollection~4.13 and Corollary~A.9]{CarlsonHaineWolf2024}.
At a rational point $x$, relative actual-point evaluation
identifies the fibre of $X'$ with the fibre of $Y'$ at $f(x)$,
as \'etale schemes over $k$; this is precisely the relative
point identification of \citet[Proposition~3.6]{CarlsonHaineWolf2024}.
Thus
\[
 f^{-1}\bigl(\operatorname{image}(Y'(k)\to Y(k))\bigr)
       =\operatorname{image}(X'(k)\to X(k)).
\]
The point map is continuous for the \'etale-open systems.
The same argument applies after every finite scalar extension.
If the finite-type schemes are not separated, apply this
argument on affine charts: the underlying point map is
Zariski continuous, the inverse images of target charts
are open, and these charts form open covers for both
topologies under consideration.

Each finite $E'/B_p$ is henselian for the nontrivial chosen
valuation and is not separably closed.  To see the latter,
enlarge it to $B_pD$ with $D$ finite Galois; its absolute group
contains the nontrivial group $P_{D_v}''$ of Claim~B.119.
The exact opened theorem of Johnson, Tran, Walsberg and Ye says
that, on a non-separably-closed henselian field, the \'etale-open
system is the system induced by its valuation topology
\citep[Corollary~6.13, p.~4055; proof of Theorem~6.5]{AN_JTWY2024}.
All of these hypotheses have just been checked.  This proves
the first assertion.

For the last assertion put $L^{\mathrm{alg}}=L\cap
\overline{\mathbb Q}$ inside the chosen $p$-adic algebraic closure.
If $G_M$-equivariance holds, the actual value at a point defined
over $L^{\mathrm{alg}}$ is fixed by $G_L$ and hence belongs
to $L^{\mathrm{alg}}$.  The compositum of $E$ with
$L^{\mathrm{alg}}$ is finite over $E$ whenever $E$ contains
the henselian algebraic subfield of $M$, as will always be
arranged below.  Restrict the already continuous map on this
finite extension to $L^{\mathrm{alg}}$-points.  The common
absolute value gives exactly the subspace topologies in
question.  No extension to the completion has been used.
\end{proof}
\paragraph{Check record.} checked only in this writing, 2026-11-18: GPT-6 Astra (OpenAI), ChatGPT Pro mode; internal effort value not exposed. Own detailed proof written and reread; primary source hypotheses matched. Root full second reading of META.10--12 passed, including exact opened JTWY source and CHW hypotheses; two clarity joins on affine charts and equivariant scalar errors inserted


\begin{claim}{B.132}{Proved}{A cyclic local quotient makes every lifting error logarithm scalar}
Let $M/\mathbb Q_p$ be finite and let $N\triangleleft G_M$
be a closed normal subgroup fixing $\mu_{p^\infty}$.
Suppose $L/M$ is finite cyclic and $N$ maps onto
$\operatorname{Gal}(L/M)$.  If
$\theta\in H^1(G_L,\mathbb Q_p(1))$ restricts to zero on
$N\cap G_L$, then its local logarithm belongs to $M\subset L$.

In the setting $E_1=B_pD$ and $M=D_v$ of Claim~B.128,
put $N=G_{E_1}\subset G_M$, using henselian local comparison.
After enlarging $M$ by a finite extension inside the fixed
field of $N$, one can find such an $L/M$ of cyclic degree
$p^s$ for any prescribed $s\ge1$.  In particular one can
require $[L:M]\ge3$.
\end{claim}
\begin{proof}
Put $H=N\cap G_L$.  The coefficient module has trivial
$H$-action, and inflation--restriction makes $\theta$ inflate
from $H^1(G_L/H,\mathbb Q_p(1))$.  For $g\in N$, conjugation
on $G_L/H$ is trivial: $G_L$ is normal in $G_M$, and
$[g,G_L]\subset N\cap G_L$.  The action of $g$ on
$\mathbb Q_p(1)$ is also trivial.  Hence $g\theta=\theta$.
The local Kummer logarithm is Galois equivariant.  The action
just calculated therefore fixes its value in $L$ under the
image of $N$, which is the full group of $L/M$.  The value
lies in $M$.
Here the logarithm is the $\mathbb Q_p$-linear map obtained
from local Kummer theory and the principal-unit logarithm,
extended by $\log p=0$.  Its equivariance follows directly
from the rational logarithm power series.  Only finite
$\mathbb Q_p$-linear combinations of actual Kummer classes
are needed below, so the same argument may be read on those
combinations alone.

For the second assertion $N$ is normal in $G_M$, because
$I''$ is characteristic in the normal inertia group and
$G_M=D_p\cap G_D$.  It fixes every root of unity.  There is
$a\in E_1^\times$ whose integral $p$-Kummer class on $N$
is nonzero.  Indeed, take $c\in B_p\setminus F_p$, whose
existence was proved in Claim~B.126.  If $c\notin W$ use
$a=c$; otherwise $1-c\notin W$ by
$W\cap(1-W)=F_p\setminus\{0,1\}$ and use $a=1-c$.
The stable kernel (META.K) gives the required nonvanishing.
Because $N$ fixes the coefficients, this integral class is
a homomorphism to $\mathbb Z_p(1)$ with image $p^r\mathbb Z_p(1)$
for a finite $r\ge0$.  Its reduction modulo $p^r$ is zero,
so $a$ has a $p^r$-th root $\alpha$ in $E_1$.
The class of $\alpha$ has full image $\mathbb Z_p(1)$.

Enlarge $M$ to contain $\alpha$ and $\mu_{p^s}$.
These are algebraic numbers in $E_1$, so this enlargement
is inside the local fixed field of $N$; $N$ stays normal
in the resulting smaller local Galois group.  Put
$L=M(\alpha^{1/p^s})$.  Surjectivity of the Kummer character
of $\alpha$ on $N$, modulo $p^s$, shows both that $L/M$
has degree $p^s$ and that $N$ maps onto its Galois group.
Since $\mu_{p^s}\subset M$, this group is cyclic.
The construction is local but uses algebraic-number
generators, so all the fields have dense subfields of
elements algebraic over $\mathbb Q$ as required for point
evaluation.  For $p=2$ choose $s\ge2$; for odd $p$ choose
$s\ge1$.
\end{proof}
\paragraph{Check record.} checked only in this writing, 2026-11-18: GPT-6 Astra (OpenAI), ChatGPT Pro mode; internal effort value not exposed. Own detailed proof written and reread; primary source hypotheses matched. Root full second reading of META.10--12 passed, including exact opened JTWY source and CHW hypotheses; two clarity joins on affine charts and equivariant scalar errors inserted


\begin{claim}{B.133}{Proved}{Regular actual point maps on principal affine opens over the metabelian field}
Let $E/B_p$ be finite, let
$U=\operatorname{Spec}E[T_1,\ldots,T_n,1/h]$ be nonempty,
and let $\phi:U_{\mathrm{et}}\to\mathbb A^1_{E,\mathrm{et}}$
be pinned and relative.  Its actual geometric point map is
induced by a unique element of
$E[T_1,\ldots,T_n,1/h]$.
In particular a pinned affine-space morphism has a polynomial
actual point map, for every prime $p$.
This theorem concerns actual point maps on the stated sources;
it does not by itself identify the full inverse-image functor
or cover an arbitrary smooth nonrational source.
\end{claim}
\begin{proof}
The constant case and $n=0$ follow from relative point
evaluation.  Assume that the actual map $f$ is nonconstant
and $n>0$.

\medskip\noindent\textbf{One finite local model and its normalized classes.}
Choose $a_0\in B_p\setminus F_p$ and the two degree-one grids
\[
 \mathcal A_0=\{r+is:0\le r,s\le4\},\qquad
 \mathcal A_1=a_0+\mathcal A_0,\qquad
 \mathcal A=\mathcal A_0\cup\mathcal A_1.
\]
For $a\in\mathcal A$, the fibre $f^{-1}(a)$ is a proper
closed subset of $U$.  Choose a nonempty principal open
$V\subset U$ avoiding their union.  After a finite scalar
extension split the finitely many geometric prime divisors
in its complement and write their equations as $h_1,\ldots,h_r$.
Kummer theory on this factorial open gives
\[
       \xi_a=\phi^*\kappa(Y-a)
            =\beta_a+\sum_j\lambda_{aj}\kappa(h_j),
                      \qquad\lambda_{aj}\in\mathbb Q_p.
                                                        \tag{META.A1}
\]
Apply Claim~B.129 to $U$ and enlarge its finite field $D$
to contain these data and all grid points.  Normalize (META.A1)
at one actual point $t_*\in V(D)$ and include the finitely
many actual values $f(t_*)-a$ in $D$, as in Claim~B.128.
With $E_1=B_pD$ and $M=D_v$, the entire actual map is
$G_M$-equivariant.  The classes in (META.A1) lift explicitly
to $V_M$ with
\[
 \beta_a=\kappa(f(t_*)-a)-
                         \sum_j\lambda_{aj}\kappa(h_j(t_*)).
                                                        \tag{META.A2}
\]
Every algebraic argument and every polynomial coefficient
in (META.A2) belongs to $M$ and is algebraic over $\mathbb Q$.
The lift is not yet asserted to have the correct actual
evaluation over $M$.

Use Claim~B.132 to enlarge $M$ inside the fixed field
of $N=G_{E_1}$ and choose a cyclic extension $L/M$ of
degree $q=p^s\ge3$ onto whose Galois group $N$ surjects.
All already established equivariance and coefficient
conditions persist under this enlargement.  Write
$L^{\mathrm{alg}}=L\cap\overline{\mathbb Q}$.
For every $t\in V(L^{\mathrm{alg}})$, actual equivariance
puts $f(t)$ in $L^{\mathrm{alg}}$.
The difference between the local Kummer class of $f(t)-a$
and the evaluation of the lifted (META.A1) restricts to zero
on $N\cap G_L$.  Its logarithm therefore lies in $M$.
If $\rho:L\to L/M$ is the quotient as an $M$-vector space,
this says
\[
 \rho\bigl(\log(f(t)-a)\bigr)
       =\rho\left(\sum_j\lambda_{aj}\log h_j(t)\right)
                       \qquad(a\in\mathcal A).
                                                        \tag{META.A3}
\]
The constant terms of (META.A2) lie in $M$ and disappear
under $\rho$.  The logarithm is locally analytic away from
zero, with derivative $dz/z$.

\medskip\noindent\textbf{A primitive output exists.}
Let $L_0$ be the unique maximal proper subfield of the cyclic
$p$-power extension $L/M$.  Suppose that $f(t)\in L_0$ for
every $t\in V(L^{\mathrm{alg}})$.  Take a nonidentity
$\tau\in\operatorname{Gal}(L/L_0)$ and choose
$b\in L^{\mathrm{alg}}$ with $\tau(b)/b$ not a root of unity.
The argument of (META.E3), replacing a moved element by
that element plus a rational constant, proves this choice.
For $u\in\mathbb Q^n$ and all sufficiently large $m$, apply
(META.A3) at $t=t_*+p^m bu$ and at $\tau t$.
Actual equivariance gives $f(\tau t)=\tau f(t)=f(t)$.
Writing the unprojected error as
$e_a(t)=\log(f(t)-a)-\sum_j\lambda_{aj}\log h_j(t)$,
Galois equivariance and the coefficients in $M$ give
$e_a(\tau t)=\tau e_a(t)=e_a(t)$, since $e_a(t)\in M$.
Consequently, for each $a$,
\[
 \sum_j\lambda_{aj}\log
       \frac{h_j(t_*+p^m\tau(b)u)}{h_j(t_*)}
 =\sum_j\lambda_{aj}\log
       \frac{h_j(t_*+p^m bu)}{h_j(t_*)}.
\]
The convergent-series and multivariable argument in
Claim~B.127 now gives
$\sum_j\lambda_{aj}\,dh_j/h_j=0$, and its prime-divisor
residues give every $\lambda_{aj}=0$.
This applies to every centre in both grids.
By (META.A1)--(META.A2), for arbitrary algebraic $t\in V$
the quotient $(f(t)-a)/(f(t_*)-a)$ then belongs to $W$.
Apply the two-grid argument to
\[
                         \frac{Z-f(t)}{Z-f(t_*)}.
\]
The first grid makes this rational function belong to $F_p(Z)$,
the translated grid makes its translate by $a_0$ belong to
$F_p(Z)$, and its first nonzero Laurent coefficient would
force $a_0\in F_p$.  It is therefore one.  Thus $f$ is
constant on $V$, hence on $U$ since a fibre is closed.
This contradiction proves the existence of
$t_0\in V(L^{\mathrm{alg}})$ for which
$z_0=f(t_0)$ generates $L/M$.

\medskip\noindent\textbf{An analytic extension is constructed, rather than assumed.}
Use two distinct rational centres $a,b\in\mathcal A_0$.
The $M$-analytic map
\[
 \Psi:L\setminus\{a,b\}\longrightarrow (L/M)^2,
       \qquad z\longmapsto
          \bigl(\rho\log(z-a),\rho\log(z-b)\bigr)
\]
has injective derivative at $z_0$.
Indeed a derivative vector in its kernel satisfies
$\delta z=c(z_0-a)=d(z_0-b)$ with $c,d\in M$.
As $z_0\notin M$ and $a\ne b$, this forces $c=d=0$.
Choose $q$ of the $2(q-1)$ scalar $M$-coordinates of $\Psi$
whose Jacobian has rank $q$.  The analytic inverse function
theorem gives an inverse for this selected-coordinate map
on a neighbourhood of $z_0$.  This theorem follows here
from the usual power-series inverse construction after
shrinking until the nonlinear part is a strict contraction.

The selected coordinates of the right side of (META.A3)
are $M$-analytic functions of the $qn$ scalar coordinates of
$t\in L^n$ near $t_0$.  Claim~B.131 makes the actual
function continuous on the algebraic $L$-points.
After shrinking their neighbourhood it takes values in the
inverse-function neighbourhood of $z_0$.  Composing the
analytic right side with that inverse constructs an
$M$-analytic map $F$ on a full neighbourhood in $L^n$ that
agrees with $f$ on all algebraic points there.
These points are dense because a number field is dense in
its chosen completion.  Thus all equations (META.A3), for
every centre in the two grids, hold for $F$ throughout the
neighbourhood.  No unrestricted continuous extension to a
completion was invoked.

\medskip\noindent\textbf{Separating the embeddings by finite logarithmic tests.}
Regard the $M$-analytic germ as a germ on
$\operatorname{Res}_{L/M}\mathbb A^n$ and extend its
coefficients to $L$.  Since $L/M$ is Galois, the input
coordinates become $q$ independent blocks
$\mathbf u_0,\ldots,\mathbf u_{q-1}$, one for each
embedding of $L$ over $M$.  The output has conjugate
analytic germs $z_0(\mathbf u),\ldots,z_{q-1}(\mathbf u)$.
Their values at the centre are the distinct conjugates of
the primitive element $f(t_0)$; hence every difference
$z_i-z_j$, $i\ne j$, is a unit in the local analytic ring.
The condition that the logarithmic error lie in the scalar
subspace $M$ is, after this extension of coefficients,
\[
 \log(z_i-a)-\log(z_0-a)
            =H_a(\mathbf u_i)-H_a(\mathbf u_0),
 \qquad H_a(\mathbf u)=\sum_j\lambda_{aj}\log h_j(\mathbf u),
                                                        \tag{META.A4}
\]
for $i>0$ and $a\in\mathcal A$.
All logarithms here mean their local analytic germs about
the corresponding nonzero values; the identities include
their constants and may be differentiated.

Take derivatives $\partial$ and $\partial'$ in different
input blocks.  The right side of (META.A4) has zero mixed
derivative.  The resulting left side is
\[
 \frac{\partial\partial'z_i}{z_i-a}
 -\frac{\partial z_i\,\partial'z_i}{(z_i-a)^2}
 -\frac{\partial\partial'z_0}{z_0-a}
 +\frac{\partial z_0\,\partial'z_0}{(z_0-a)^2}=0.
                                                        \tag{META.A5}
\]
After multiplication by the two squared denominators its
numerator has degree at most three in $a$.
The four centres $0,1,2,3$ therefore make it the zero
polynomial.  Expansion at its distinct poles $z_i,z_0$
shows that both mixed second derivatives and both products
of first derivatives in (META.A5) are zero.
The ring of analytic germs is an integral domain.  Thus each
nonconstant $z_i$ depends on at most one input block.
None is constant: a constant germ would make $f$ constant
on a nonempty $p$-adic open set of algebraic points, which
is Zariski dense, contradicting the proper closed fibre
condition.  This reasoning applies to every conjugate
output component.

For a block $k$ different from $i$ and $0$, first
differentiation of (META.A4), at two distinct centres,
shows $\partial_k z_i=\partial_k z_0=0$, since $z_i-z_0$
is a unit.  Hence the one block on which $z_0$ depends
belongs to $\{i,0\}$ for every $i>0$.
There are at least two such indices because $q\ge3$,
so that block is $0$.  The block for $z_i$ is either $i$
or $0$.  If it were $0$, differentiating (META.A4) in
block $i$ would make every $H_a$ constant.  Differentiating
then in block $0$ at two centres would make both $z_i$
and $z_0$ constant, a contradiction.
Thus $z_i$ depends exactly on its matching block $i$.

It follows that
$\log(z_i-a)-H_a(\mathbf u_i)$ equals
$\log(z_0-a)-H_a(\mathbf u_0)$ with independent variables
on the two sides.  Both sides are constant germs.
In the original embedding this proves that $F$ is
$L$-analytic and that
\[
         \frac{dF}{F-a}
                  =\sum_j\lambda_{aj}\frac{dh_j}{h_j}
                    =:\omega_a\qquad(a\in\mathcal A).
                                                        \tag{META.A6}
\]
The forms $\omega_a$ are rational forms with coefficients
in $M$.  Choose two distinct centres $a,b$.
As $F$ is nonconstant in characteristic zero,
$\omega_a-\omega_b$ is nonzero.  For a coordinate $k$
where its coefficient is nonzero, (META.A6) gives
\[
           F=\frac{a\omega_{a,k}-b\omega_{b,k}}
                         {\omega_{a,k}-\omega_{b,k}}
                          =:R\in M(T_1,\ldots,T_n).
                                                        \tag{META.A7}
\]

\medskip\noindent\textbf{Algebraic coefficients and global recovery.}
The rational function $R$ has coefficients algebraic over
$\mathbb Q$.  To verify this without assuming continuity
in a different absolute value, write it as a reduced
numerator and denominator of minimal maximum total degree
$d$.  The space of pairs of degree at most $d$ defining
this rational function is one-dimensional: any second pair
is a polynomial multiple of the reduced pair, and a
nonconstant multiplier raises the maximum degree.
Evaluation at all algebraic points in the analytic
neighbourhood imposes precisely these linear equations,
because that set is Zariski dense.  Finitely many of the
equations therefore already cut out the one-dimensional
space.  Their coefficients and their actual values $f(t)$
are algebraic numbers.  Solving this finite system and
normalizing one nonzero coefficient gives a numerator and
denominator over $\overline{\mathbb Q}$.
Thus $R\in\overline{\mathbb Q}(\mathbf T)$.

Equation (META.A6) is an identity of rational forms,
since it holds as a germ on an analytic open set.
For each $a\in\mathcal A$, its residues at every affine
prime divisor show that $R-a$ has no zero or pole outside
the divisors $h_j=0$, and that
$\lambda_{aj}=\operatorname{ord}_{h_j}(R-a)\in\mathbb Z$.
Factoriality therefore gives
\[
                  R-a=c_a\prod_j h_j^{\lambda_{aj}},
                         \qquad c_a\in\overline{\mathbb Q}^\times.
\]
In particular $R-a$ is a unit on $V$.
After adjoining the finite algebraic coefficient set,
the difference $\xi_a-\kappa(R-a)$ is an arithmetic class.
At the algebraic point $t_0$ the actual functions agree,
so this arithmetic difference vanishes after a finite
extension.  Restriction on rational $p$-cohomology is
injective by corestriction, and hence the difference is zero.
Thus for every algebraic $t\in V$ and every grid centre,
\[
                            \frac{f(t)-a}{R(t)-a}\in W.
\]
Apply the two-grid argument once more, now to
$(Z-f(t))/(Z-R(t))$.  It gives $f(t)=R(t)$ on all of
$V(\overline{\mathbb Q})$, not merely on the local
analytic neighbourhood.

Finally $R$ has no pole on the original principal open
$U$.  Otherwise choose a smooth algebraic point on a
pole divisor in $U$ where the reduced numerator is nonzero,
and a transverse algebraic line through it not contained
in $U\setminus V$.  Along algebraic points on this line
with parameter $p^m\to0$, the values of $R$ tend to
infinity in the $p$-adic absolute value.  They equal $f$
on $V$, contradicting Claim~B.131 at the chosen point.
The factorial coordinate ring of $U$ now gives
$R\in\Gamma(U_{\overline{\mathbb Q}},\mathcal O)$.
At any remaining point of $U$, choose an algebraic line
approaching it through $V$.  Continuity of $f$ and of this
regular $R$ makes their values equal there as well.
All sequences used here lie in one finite extension of
$E$ after adjoining their algebraic data, exactly the
scope of Claim~B.131.

The actual map is $G_E$-equivariant.  Equality on all
geometric points therefore makes the regular function
$R$ invariant under $G_E$, so it descends to
$\Gamma(U,\mathcal O)=E[\mathbf T,1/h]$.
Uniqueness follows from geometric-point density on this
reduced affine scheme.  This proves the assertion.
\end{proof}
\paragraph{Check record.} checked only in this writing, 2026-11-18: GPT-6 Astra (OpenAI), ChatGPT Pro mode; internal effort value not exposed. Own detailed proof written and reread; primary source hypotheses matched. Root full second reading of META.10--12 passed, including exact opened JTWY source and CHW hypotheses; two clarity joins on affine charts and equivariant scalar errors inserted


\begin{claim}{B.134}{Proved}{A nonzero geometric Kummer class invisible at every point}
At $p=5$, let $E=B_5$ and choose $i\in E$ whose reduction
at the chosen place is $2\in\mathbb F_5$.
On the elliptic curve $A:y^2=x^3-x$ over $E$ there is a
nonzero class $\xi\in H^1(A,\mathbb Q_5(1))$ whose
evaluation is zero at every $A(E')$-point for every finite
extension $E'/E$.  The same holds after restriction to any
nonempty affine open of $A$.
Consequently Claim~B.127 does not extend to arbitrary
smooth sources.  This is not a counterexample to reconstruction.
\end{claim}
\begin{proof}
The equation has discriminant $64$, so it defines an elliptic
scheme over the valuation ring $R$ of $E$.  This valuation
ring is strictly henselian: it is henselian and its residue
field is $\overline{\mathbb F}_5$.
The automorphism $[i](x,y)=(-x,iy)$ extends to this model.
On the invariant differential $dx/(2y)$ it acts by $i$.
The endomorphism $\alpha=[2]+[i]$ therefore has differential
$2+i$, a unit since its reduction is $4$.
Translation by points of the smooth group scheme shows
that $\alpha$ is \'etale everywhere.  It is proper and
quasifinite, and hence finite \'etale.

Its generic degree is five.  After a complex embedding,
multiplication by $2+i$ on a Gaussian lattice has matrix
$\left(\begin{smallmatrix}2&-1\\1&2\end{smallmatrix}\right)$
and index five.  The same calculation shows that
$\ker\alpha^n$ is cyclic of order $5^n$.
Indeed write $(2+i)^n=a_n+ib_n$.
The multiplication matrix has determinant $5^n$ and
$\gcd(a_n,b_n)=1$: a common rational prime would have
to be $5$, but divisibility by $5=(2+i)(2-i)$ contradicts
coprimality of these two Gaussian primes.
Its Smith form is therefore $\operatorname{diag}(1,5^n)$.
The source of every generic cover $\alpha^n:A\to A$
is connected, so these cyclic torsors are nontrivial.

The kernels of the model maps are finite \'etale over
the strictly henselian ring $R$, and hence constant.
Their compatible inverse system is a free $\mathbb Z_5$-module
of rank one.  Fix compatible cyclic identifications and,
since $\mu_{5^\infty}\subset E$, identify this system
with the Tate-twisted coefficient system.
The covers $\alpha^n$ define an integral class whose
geometric character has image all of $\mathbb Z_5(1)$.
Its image $\xi$ in rational cohomology is nonzero.

For a finite extension $E'/E$, the valuation ring $R'$
is still strictly henselian with residue field
$\overline{\mathbb F}_5$.
Every point of $A(E')$ extends to an $R'$-point by
properness.  Its inverse image under every model cover
$\alpha^n$ is a nonempty finite \'etale $R'$-scheme,
hence a disjoint union of copies of $\operatorname{Spec}R'$.
Thus each evaluated torsor is trivial.
Compatible lifts can be chosen successively, since the
same argument makes $\alpha:A(E')\to A(E')$ surjective.
The integral evaluation, and hence its rational image,
is zero.  Restriction to a nonempty affine open remains
nonzero because its fundamental group surjects onto
that of the normal connected proper curve.
The evaluations there are unchanged.  This proves the claim.
\end{proof}
\paragraph{Check record.} checked only in this writing, 2026-11-18: GPT-6 Astra (OpenAI), ChatGPT Pro mode; internal effort value not exposed. Own detailed proof written and reread; primary source hypotheses matched. Root full reading of META13/B134 passed: finite etale Gaussian endomorphism towers, strictly Henselian fibres, compatible lifts, nonzero geometric character, and affine-open restriction


\paragraph{Scope after the local analytic argument.}
Claim~B.133 establishes actual-point reconstruction on
principal affine-space opens at $B_p$ and every finite extension,
uniformly in the prime.  By itself it does not replace the
arbitrary-source input in Carlson--Haine--Wolf, Theorem~4.18.
For a general finite Noether projection the actual data are a
monic polynomial of output values, not a single output value.
The logarithmic quotient map used above can lose rank when
conjugate output polynomials have common factors.
Claim~B.136 below resolves this remaining step: reduced
ratios on a constant-gcd stratum, algebraic root types and
a large cyclic quotient recover the norm polynomial;
normalization and two grids remove its common factor.
That later theorem proves full reconstruction at $B_p$
and every finite extension.  The present intermediate
argument is retained to distinguish precisely what the
scalar analytic calculation establishes on its own.

\paragraph{Sources and checks.}
The new continuity input was read in the published JEMS PDF:
Johnson--Tran--Walsberg--Ye, Corollary~6.13 and the proof of
Theorem~6.5, printed pp.~4054--4055.  The reduction from pinned
topos maps to actual continuous maps was proved above using
the already opened Carlson--Haine--Wolf statements.
The local group, analytic and finite-grid arguments were written
and reread in this writing; their second-reading status is recorded
separately.  No priority or external human verification is claimed.

\subsection{A finite logarithmic test for a monic polynomial}

\begin{claim}{B.135}{Proved}{Rational logarithmic derivatives at twice the degree determine the coefficients}
Let $C$ be a field of characteristic zero, put
$K=C(T_1,\ldots,T_n)$, and let $\mathscr F/K$ be a field
with the commuting derivations $\partial_1,\ldots,\partial_n$
extending the coordinate derivations of $K$.  Assume that the
common field of constants of $\mathscr F$ is exactly $C$.
Let $P(Z)\in\mathscr F[Z]$ be monic of degree $d\ge1$.
Suppose that there are at least $2d$ distinct elements $a\in C$
such that $P(a)\ne0$ and
\[
             \frac{\partial_jP(a)}{P(a)}\in K
                   \quad(1\le j\le n).                 \tag{B.135.1}
\]
Then $P\in K[Z]$.
The bound $2d$ is optimal in this generality.

In particular the assertion applies to the fraction field of
the ring of convergent power-series germs at a point of $L^n$,
where $L/\mathbb Q_p$ is finite and $C=L$.
The hypothesis concerns exact logarithmic derivatives before
passing to any quotient by scalar-valued errors.
\end{claim}
\begin{proof}
Call a nonzero element $u\in\mathscr F$ a logarithmic element
if $\partial_j u/u\in K$ for every $j$.
Declare two such elements equivalent when their ratio belongs
to $K^\times$.  Representatives of distinct equivalence classes
are linearly independent over $K$.  Indeed, suppose otherwise
and choose a relation of smallest positive length,
\[
                    u_1+\sum_{i=2}^m c_i u_i=0,
                    \qquad c_i\in K^\times.
\]
There is no relation of length one.  For each $j$, differentiate
this relation and subtract $(\partial_j u_1/u_1)$ times the
original one.  The coefficient of $u_1$ disappears.  Minimality
therefore gives, for every $i>1$,
\[
 \partial_jc_i+c_i\left(
             \frac{\partial_j u_i}{u_i}
             -\frac{\partial_j u_1}{u_1}\right)=0.
\]
Equivalently, all derivatives of $c_i u_i/u_1$ vanish.
The common-constant hypothesis puts this ratio in $C^\times$,
so $u_i/u_1\in K^\times$, a contradiction.

Let $\mathcal A$ be a set of $2d$ centres satisfying
(B.135.1), and put $u_a=P(a)$.
Partition these logarithmic elements into their equivalence
classes, and choose representatives $v_\chi$.
Interpolation at any $d+1$ centres shows that every coefficient
of $P$ belongs to the $K$-span of the $u_a$; this is permitted
because $2d\ge d+1$.  The independence just proved gives a
unique expansion
\[
                 P(Z)=\sum_\chi v_\chi P_\chi(Z),
                    \qquad P_\chi\in K[Z],\quad
                           \deg P_\chi\le d.           \tag{B.135.2}
\]
The leading coefficient is $1$.  Include $1$ among the
logarithmic elements in the independence argument.  It follows
that one of the represented classes is the trivial class
$K^\times$; choose its representative $v_0=1$.
Coefficient comparison in (B.135.2) then gives
\[
       P_0\text{ monic of degree }d,\qquad
                   \deg P_\chi\le d-1\quad(\chi\ne0).
\]
At a centre $a$, the value $u_a$ belongs to a single class.
Consequently every $P_\chi(a)$ except the coefficient of that
class is zero.  Since $P_0$ has at most $d$ roots, at least
$d$ of the $2d$ centres satisfy $P_0(a)\ne0$.
At each of these centres every $P_\chi$ with $\chi\ne0$
vanishes.  Its degree is at most $d-1$, so it is identically
zero.  Thus $P=P_0\in K[Z]$.

For the analytic application translate the centre to the origin.
The field of meromorphic germs embeds in the iterated formal
Laurent-series field
$L((T_1))\cdots((T_n))$.  Successively applying the coordinate
derivations to its Laurent expansions shows that an element
annihilated by all of them belongs to $L$: in characteristic
zero every nonconstant Laurent monomial has a nonzero derivative
in some variable.  Thus its common constant field is exactly
$L$, as required.  No arithmetic statement about the
coefficients of $P$ is used in this application.

The number of centres in this lemma cannot in general be
lowered to $2d-1$.  For $d\ge2$, choose monic $A(Z)\in C[Z]$
of degree $d$ and $B(Z)\in C[Z]$ of degree $d-1$, each with
distinct roots in $C$, and with disjoint root sets.
In the analytic case let $u=\exp(T_1)$ on a sufficiently small
neighbourhood of zero and put $P=A+uB$.
At each of the $d$ roots of $A$, the value is a nonzero
constant times $u$; at each of the $d-1$ roots of $B$, it is a
nonzero constant.  Their logarithmic derivatives are rational,
but $P$ is not rational because $u$ is not rational.
If this exponential were rational in all $n$ variables, its
vanishing partial derivatives in the other variables would
put it in $C(T_1)$.  For completeness, if $R\in C(T_1)$ satisfied $R'=R$ and
$R\ne0$, every finite zero or pole of $R$ would give a pole
of $R'/R=1$.  Thus $R$ would be constant, a contradiction.
For $d=1$, one centre likewise does not suffice: use $P=Z+u$
and the centre zero.
\end{proof}
\paragraph{Check record.} Checked only in this writing, 18November2026: GPT-6 Astra (OpenAI), ChatGPT Pro mode; root full proof written and read, finite-checks reader full second reading PASS including sharpness and analytic constant field. Separate effort setting not exposed.


\subsection{Norm recovery and reconstruction over the metabelian field}

\begin{claim}{B.136}{Proved}{Full pinned reconstruction over every finite extension of the metabelian inertia field}
Let $p$ be any prime, let $I$ be the chosen local inertia
subgroup of $G_{\mathbb Q}$, and put
$B_p=\overline{\mathbb Q}^{I''}$.
For every finite extension $E/B_p$, the \'etale-topos functor
on finite-type $E$-schemes, localized at universal
homeomorphisms, is fully faithful for pinned relative
geometric morphisms over $\operatorname{Spec}(E)_{\mathrm{et}}$.
The relative structural isomorphism is retained.

More explicitly, for every smooth geometrically integral
affine $X/E$, every such pinned morphism
$X_{\mathrm{et}}\to\mathbb A^1_{E,\mathrm{et}}$ has the
actual geometric point map of a unique regular function
on $X$.  The statement is uniform in $p$, including $2$ and $3$.
\end{claim}
\begin{proof}
We first prove the actual-point assertion.  The zero-dimensional
and constant cases follow from relative point evaluation.
Assume $n=\dim X>0$ and that the point map $f$ is nonconstant.
Every fibre $S_a=f^{-1}(a)$ is a proper closed subset after
the finite scalar extension needed to define $a$.
All finite scalar extensions made in this proof remain
finite over $B_p$, so the stable kernel (META.K) applies
throughout.

\medskip\noindent\textbf{A normalized norm system over one finite local field.}
Take any finite Noether normalization
$g:X\to\mathbb A^n$ of degree $d$.
Choose $a_0\in B_p\setminus F_p$ and put
\[
 \mathcal A_0=\{r+is:0\le r,s\le4d\},\qquad
 \mathcal A_1=a_0+\mathcal A_0,\qquad
                  \mathcal A=\mathcal A_0\cup\mathcal A_1.
\]
Choose a nonempty principal open $V\subset\mathbb A^n$
avoiding the branch locus and all the closed images $g(S_a)$.
Over it the map is finite \'etale of degree $d$.
Corestriction of the pulled-back classes of $Y-a$ gives
classes $\xi_a\in H^1(V,\mathbb Q_p(1))$ with actual
evaluations
\[
 \kappa(P_t(a)),\qquad
               P_t(Z)=\prod_{x\in g^{-1}(t)}(Z-f(x)).
                                                        \tag{META.NR1}
\]
Here $P_t$ is monic of degree $d$ and none of these grid
values is zero.  More precisely, corestriction of the
class of $Y-a$ evaluates as
$\kappa((-1)^dP_t(a))$; the factor $(-1)^d$ is torsion
and has zero rational $p$-adic Kummer class, giving the
displayed formula.
The Kummer calculation on the factorial principal open,
after splitting its prime boundary factors $h_j$, gives
\[
                  \xi_a=\beta_a+
                         \sum_j\lambda_{aj}\kappa(h_j),
                           \qquad\lambda_{aj}\in\mathbb Q_p.
                                                        \tag{META.NR2}
\]
Normalize at an actual point as in Claim~B.128.
That claim supplies a finite Galois number field $D$,
$E_1=B_pD$ and $M=D_v$ such that all coefficients,
grid points and normalizing actual values belong to $D$,
the norm assignment is $G_M$-equivariant, and each
$\beta_a$ is an explicit finite linear combination of
Kummer classes of elements of $D^\times$.

Choose $t_*\in V(D)$ and enlarge $D$ to contain all points
of the finite \'etale fibre above $t_*$.
Thus the cover has $d$ analytic inverse branches over $M$
on a sufficiently small ball about $t_*$.
The inverse function construction can be made uniformly
convergent on that ball after every finite extension of
$M$: its coefficients are in $M$ and the same contraction
bound applies.  All its values at algebraic parameter
points are algebraic, since they belong to a finite
algebraic fibre.  The $d$ branches exhaust that fibre.

Put $N=G_{E_1}\triangleleft G_M$ and use Claim~B.132
to enlarge $M$ inside the fixed field of $N$ and choose
a cyclic extension $L/M$ of degree
\[
                            q=p^s>d+1
                                                        \tag{META.NR3}
\]
onto whose Galois group $N$ surjects.
All the previous coefficient and equivariance assertions
persist.  Let $\tau$ generate $\operatorname{Gal}(L/M)$.
Write $L^{\mathrm{alg}}=L\cap\overline{\mathbb Q}$.
For algebraic $t$ in the chosen $L$-ball, equivariance
puts every coefficient of $P_t$ in $L^{\mathrm{alg}}$.
These coefficients are continuous functions of $t$.
Indeed the analytic inverse branches of $g$ are continuous,
the actual $f$ is continuous on points over the finite
extension $E_1L^{\mathrm{alg}}/E_1$ by Claim~B.131,
and the coefficients in (META.NR1) are symmetric polynomial
expressions in their finitely many values.
Their values lie in $L$, so they have its subspace topology.

For $a\in\mathcal A$, set locally
$H_a(t)=\sum_j\lambda_{aj}\log h_j(t)$.
The evaluation of the normalized lift (META.NR2) over $L$
differs from the actual Kummer class of $P_t(a)$ by a class
vanishing on $N\cap G_L$.  Claim~B.132 therefore gives
\[
                       \log P_t(a)-H_a(t)\in M.
                                                        \tag{META.NR4}
\]
The arithmetic normalizing logarithms already belong to
$M$.  Taking $\tau^i$ of (META.NR4) and subtracting it,
with $1\le i<q$, gives
\[
 \log\frac{\tau^iP_t(a)}{P_t(a)}
                     =H_a(\tau^i t)-H_a(t).
                                                        \tag{META.NR5}
\]
No point-theoretic lifting of (META.NR2) over $M$ is
asserted in deriving this equality.

\medskip\noindent\textbf{Analytic reduced ratios on a constant-degree stratum.}
For every algebraic parameter in the ball define
\[
 c_i(t)=\deg\gcd(P_t,\tau^iP_t),\qquad
          r_i(t,Z)=\frac{\tau^iP_t(Z)}{P_t(Z)}.
\]
Choose an algebraic parameter $t_0$ where
$\sum_{i=1}^{q-1}c_i(t)$ is minimal.
Such a value exists because this sum takes only finitely
many nonnegative integer values.
The condition that the gcd degree be at least a given
integer is the vanishing of subresultant minors, and hence
is closed in the polynomial coefficients.
Continuity of those coefficients permits shrinking around
$t_0$ so that $c_i(t)\le c_i(t_0)$ for all $i$.
Minimality of the sum then makes every $c_i(t)$ constant
on the algebraic points of this neighbourhood.
Put $e_i=d-c_i(t_0)$.
The reduced numerator and denominator of $r_i$ are monic
of the same degree $e_i$.

For each grid centre, (META.NR5), continuity, and the
local inverse relation between logarithm and exponential
give an analytic function equal to $r_i(t,a)$ near $t_0$:
\[
 r_i(t_0,a)\exp\bigl(
 H_a(\tau^it)-H_a(t)
       -H_a(\tau^it_0)+H_a(t_0)\bigr).
                                                        \tag{META.NR6}
\]
The neighbourhood is chosen so that each ratio to its
base value lies in a logarithm-injective principal-unit
neighbourhood.  Thus (META.NR6) recovers the actual value,
not only its logarithm.

If $e_i=0$, the ratio is one.  If $e_i>0$, interpolate
its monic reduced numerator and denominator from any
$2e_i$ distinct grid centres.
The interpolation matrix is invertible at $t_0$.
Indeed a homogeneous solution $(\delta A,\delta B)$ has
degrees at most $e_i-1$ and makes
$\delta A\,B-A\,\delta B$ vanish at $2e_i$ centres.
This polynomial has degree at most $2e_i-1$, so it is
zero.  Coprimality of the monic reduced pair $A,B$ then
forces $\delta A=\delta B=0$.
Cramer's rule applied to (META.NR6) gives analytic
coefficient functions for the reduced pair near $t_0$.
They coincide with the actual reduced pair at every
algebraic parameter there by uniqueness.
This constructs analytic $r_i$ without asserting that
the coefficients of $P_t$ were analytic.

Use scalar coordinates for
$\operatorname{Res}_{L/M}\mathbb A^n$ centred at $t_0$.
Choose the $M$-basis of $L$ in
$L\cap\overline{\mathbb Q}$, using powers of the
algebraic-number Kummer generator in Claim~B.132.
Algebraic scalar-coordinate points are then dense in
the scalar-coordinate balls and give actual algebraic
parameters in $L^n$.
These are $qn$ coordinates over $M$, even when
$t_0\notin M^n$.
After extending coefficients to $L$, they split into
independent blocks $\mathbf u_0,\ldots,\mathbf u_{q-1}$,
centred at the conjugates of $t_0$.
Let $\mathscr F$ be the field of meromorphic germs in
these blocks over $L$.
The semilinear action of $\tau$ permutes the blocks,
because the scalar coordinates before extension were
over $M$.  On actual scalar-coordinate points it acts
on the values in $L$.
Consequently the identities among actual polynomial
ratios give identities of analytic germs
\[
      \tau^j r_i=\frac{r_{i+j}}{r_j},\qquad r_0=1,
                    \quad i,j\in\mathbb Z/q\mathbb Z.
                                                        \tag{META.NR7}
\]
The logarithmic derivative form of (META.NR5) is
\[
             d\log r_i(a)
                  =dH_a(\mathbf u_i)-dH_a(\mathbf u_0).
                                                        \tag{META.NR8}
\]
These formulas concern germs in the scalar coordinates
at the single chosen point; they do not require its
conjugates to lie in that same small unsplit ball.

\medskip\noindent\textbf{Divisor types and positivity from a large cyclic quotient.}
Extend the commuting coordinate derivations uniquely to
an algebraic closure of $\mathscr F$; in characteristic
zero these extensions exist and still commute.
Factor $r_i$ there as a product of distinct linear factors
$(Z-\alpha)$ with nonzero integral multiplicities,
positive for zeros and negative for poles.
No analytic choice of a root is made.

For derivatives $\partial,\partial'$ in different blocks,
(META.NR8) gives $\partial\partial'\log r_i(a)=0$.
As a rational function of $a$, this expression has
denominator degree at most $4d$ and numerator degree
at most $4d-1$: the reduced numerator and denominator
have degree $e_i\le d$, and their leading coefficients
are constant one.
The $4d+1$ integer centres $0,\ldots,4d$ make that
numerator identically zero.
Partial fractions at a root $\alpha$ then give
\[
              \partial\partial'\alpha=0,
                        \qquad\partial\alpha\,\partial'\alpha=0.
                                                        \tag{META.NR9}
\]
Thus each nonconstant root is annihilated by all
derivations except those of a single block.
We call this its type.  If all its derivatives vanish,
it is algebraic over $L$: differentiating its monic
minimal polynomial shows that every coefficient of that
polynomial is a common constant, hence belongs to $L$.
Here the common constants of meromorphic germs are $L$,
as proved in B.135.
The same minimal-polynomial argument shows that a root
of type $j$ is algebraic over the field of meromorphic
germs in block $j$ alone.
For this last assertion, the elements of $\mathscr F$
annihilated by the other block derivations are exactly
that field: specialize the other variables at a nearby
constant point where a denominator is not identically
zero, and use the derivative identities to identify the
resulting meromorphic functions.
Nonconstant roots of different types are distinct,
since equality would annihilate every derivative.

First derivatives of (META.NR8) in blocks other than
$0,i$ vanish as well.  The same rational-function
interpolation and partial-fraction argument shows that
the nonconstant roots of $r_i$ have only types $0$ and $i$.
Let $D_i=\operatorname{div}_Z(r_i)$ be its finite zero-pole
divisor.  Equation (META.NR7) implies
\[
                         D_i-D_j=\tau^j D_{i-j}.
\]
For distinct nonzero $i,j$, the right side has only types
$i,j$ and constant roots.  Hence the type-zero part of
$D_i$ is the same for every $i\ne0$.
Write it as $-B$.
Taking $j=-i$ in the cocycle, or using
$\tau^i D_{-i}=-D_i$, identifies the type-$i$ part as
$\tau^iB$.  Thus
\[
                       D_i=\tau^iB-B+C_i,
                                                        \tag{META.NR10}
\]
where $C_i$ is supported on roots algebraic over $L$.
The divisor $B$ is invariant under the absolute Galois
group of $\mathscr F$: root type is defined by the
extended derivations, which commute with that group.

There are no constant roots in any $D_i$.
Suppose $c$ is such a zero or pole.
Specialize once at an algebraic parameter where the
analytic reduced ratio retains this zero or pole.
The ratio has algebraic-number coefficients there,
so $c$ is itself algebraic over $\mathbb Q$.
At every algebraic parameter in a sufficiently small
analytic open, a constant pole makes $c$ a root of $P_t$.
A constant zero makes it a root of $\tau^iP_t$; applying
one fixed lift of $\tau^{-i}$ gives a fixed algebraic
root of $P_t$ instead.
This would put that analytic open in the closed image
$g(S_c)$, or in the image for that fixed conjugate of $c$.
Such an image is proper because $g$ is finite and
$\dim S_c<n$.
Algebraic points of a nonempty $p$-adic open are Zariski
dense, so this is impossible.  Therefore $C_i=0$.

The divisor $B$ is effective.
If it had coefficient $-m<0$ at a nonconstant root
$\alpha$ of type zero, (META.NR10) would make
$\tau^j\alpha$ a pole of $r_j$ of multiplicity $m$
for every $1\le j<q$.
These roots have distinct types and hence are distinct.
The least common multiple of the reduced denominators
of the $r_j$, taken over $\mathscr F$, would have degree
at least $(q-1)m>d$.
But this is impossible on specialization to actual
parameters: each such reduced denominator divides the
actual monic degree-$d$ polynomial $P_t$.
Here specialization can be made without losing lcm degree.
Compute the finitely many gcds and lcms by subresultants
over the meromorphic-germ field and retain the nonzero
minors and denominator coefficients that specify their
degrees.  Avoiding their zeros and poles gives a nonempty
analytic open on which these same computations specialize.
A nonzero analytic germ over a finite extension of $M$
cannot vanish on an open set of the scalar $M$-coordinates;
the one-variable identity theorem, applied successively,
proves this assertion.  Algebraic scalar-coordinate points
are dense, so one can choose an actual algebraic parameter
in that open.  Its polynomial $P_t$ would be divisible
by an lcm of degree greater than $d$, a contradiction.
This is an analytic-open specialization argument; no
algebraic open was assumed before rationality was proved.

Let $Q(\mathbf u_0,Z)$ be the monic polynomial with
root divisor $B$, of degree $e$.
It is defined over $\mathscr F$ by the Galois invariance
of $B$; its coefficients are annihilated by the other
block derivations, and thus depend only on block zero.
Moreover it is exactly the reduced monic denominator of
$r_1$.  Indeed $B$ is effective, its roots have type zero,
and the roots of $\tau B$ have type one, so there is no
cancellation in (META.NR10).
The denominator was already constructed analytically by
interpolation.  Thus $Q$ has analytic coefficients over
$L$, and $0\le e\le d$.
The normalization at infinity is one, so (META.NR10) gives
the exact identity
\[
                     r_i=\frac{\tau^iQ}{Q}.
                                                        \tag{META.NR11}
\]
This conclusion also permits $e=0$, in which case $Q=1$;
no nonconstant-factor assumption is inserted here.

\medskip\noindent\textbf{Rational coefficients and removal of the invisible factor.}
Differentiate (META.NR11) and (META.NR8) in block zero.
Since $\tau^iQ$ depends on block $i$, for every grid centre
we obtain
\[
                   d\log Q(t,a)=dH_a(t)
                       =\sum_j\lambda_{aj}\,dh_j/h_j.
                                                        \tag{META.NR12}
\]
If $e>0$, B.135 applies with constant field $L$, the
meromorphic-germ field in block zero, and at least $2e$
of our centres.  Its hypotheses concern exact derivatives,
and (META.NR12) supplies exactly those hypotheses.
It yields $Q\in L(t)[Z]$.
The same conclusion is immediate when $e=0$.

Every coefficient of $Q$ is in fact a rational function
over $\overline{\mathbb Q}$.
On the analytic constant-degree neighbourhood, $Q_t$ is
the reduced denominator of an actual ratio and therefore
divides $P_t$.
At algebraic parameters its roots, and hence its
coefficients, are algebraic numbers.
Apply finite rational interpolation to each coefficient
function in $L(t)$, exactly as in Claim~B.133:
use its reduced numerator and denominator of minimal
maximum degree, whose normalized coefficient vector is
the unique solution of finitely many linear equations
at algebraic points in the analytic neighbourhood.
This proves $Q\in\overline{\mathbb Q}(t)[Z]$.

Taking prime-divisor residues in the rational identity
(META.NR12) shows that every $\lambda_{aj}$ is an integer
and that $Q(t,a)$ has no zero or pole on $V$.
Thus
\[
                  Q(t,a)=c_a\prod_jh_j(t)^{\lambda_{aj}}
\]
with an algebraic constant $c_a\ne0$.
Interpolation at $e+1$ centres also shows that every
coefficient of the monic $Q$ is regular on $V$.
After a finite extension, the cohomological difference
$\xi_a-\kappa(Q(t,a))$ is consequently arithmetic.

Choose an algebraic parameter $t_\dagger$ in the local
constant-degree neighbourhood and put
\[
                   C(Z)=P_{t_\dagger}(Z)/Q_{t_\dagger}(Z).
\]
This is a monic polynomial over $\overline{\mathbb Q}$
of degree $d-e$ because $Q_{t_\dagger}$ divides that
actual norm polynomial.
Its value at every grid centre is nonzero.
Evaluating the arithmetic difference there gives
\[
                    \xi_a=\kappa(Q(t,a))+\kappa(C(a)).
                                                        \tag{META.NR13}
\]
This equality holds after one finite scalar extension
containing the finitely many algebraic data; restriction
and corestriction justify arithmetic equality if a further
point field was used.
At an arbitrary algebraic parameter $t\in V$, (META.NR13)
and the stable kernel imply
\[
                        \frac{P_t(a)}{Q_t(a)C(a)}\in W
                                  \qquad(a\in\mathcal A).
\]
Both numerator and denominator of
$P_t(Z)/(Q_t(Z)C(Z))$ are monic of degree $d$.
The two-grid argument of Claim~B.126 therefore makes
that rational function one.  We have proved
\[
                           P_t(Z)=Q_t(Z)C(Z)
                    \quad\hbox{for every algebraic }t\in V.
                                                        \tag{META.NR14}
\]
If $C$ had positive degree, one of its algebraic roots
would be an output value on every finite fibre over $V$.
The proper closed image $g(S_c)$ would then contain $V$,
the same dimension contradiction used above.
Thus $C=1$, $e=d$, and every actual norm polynomial is
the monic polynomial $Q\in\overline{\mathbb Q}(t)[Z]$,
with coefficients regular on the chosen $V$.
In particular the alternative $e=0$ would force $d=0$
and is excluded automatically in this nonconstant case.

\medskip\noindent\textbf{One algebraic graph and its extension across every point.}
The preceding argument applies to every finite Noether
projection after any finite scalar extension.
We give the graph and generic-degree argument directly;
its earlier occurrence in Claim~B.109 is a comparison,
not a premise about the present field.
For one normalization, shrink to a smooth affine
$X_0\subset X$ on which the recovered coefficients of
$P_g(g(x),Z)$ are regular and every actual point satisfies
$P_g(g(x),f(x))=0$.
This monic equation puts the Zariski closure of the
actual graph in a scheme finite over $X_0$.
Remove the closed images of its nondominating components.
At least one component dominates, since the graph has
a point above every algebraic point of $X_0$.
Let its dominating irreducible components have degrees
$e_1,\ldots,e_r$.
The output coordinate is nonconstant on each component:
after removing the other components the actual graph
is dense in that component, or it would not be a
component of the graph closure.  Its projection is
therefore dense in $X_0$, and a constant output
would make a proper closed fibre dense in $X_0$.
Choose a generic finite projection $g':X_0\to\mathbb A^n$
such that $(g',Y)$ is birational on each component and
has distinct images on different components.
Over the generic value of $Y$, these are
$(n-1)$-dimensional subvarieties embedded in $X_0$;
the incidence of pairs with the same generic linear
projection to $\mathbb A^n$ has dimension at most $n-2$.
For $n=1$ it is separation of a finite set.
The conditions first give a nonempty open of projection
coefficient space over $\overline{\mathbb Q}(Y)$.
It contains an open defined by finitely many nonvanishing
polynomials in those coefficients, whose coefficients
are rational functions of $Y$.
Clear their denominators and expand in $Y$.
For any one such polynomial, vanishing identically in
$Y$ is the simultaneous vanishing of its finitely many
coefficient polynomials over $\overline{\mathbb Q}$.
The complement is a nonempty open over that field.
Intersect these finitely many opens with the nonempty
open condition of finite Noether normalization of $X_0$.
An algebraic point in this intersection gives the
required projection after a finite scalar extension.

If the degree of $g'$ is $d'$, the distinct minimal
polynomials of the output coordinate on these components
have degrees $d'e_i$.
All divide the actual norm polynomial of degree $d'$,
because it vanishes on the dense actual graph.
Hence $\sum_i e_i\le1$.
There is exactly one component and its degree is one.
The actual function therefore agrees with a rational
function $h\in\overline{\mathbb Q}(X)$ on a nonempty
Zariski open.  Its uniqueness and the original $G_E$
equivariance put $h$ in $E(X)$.

Continuity now removes all poles and exceptional point
values.  If $h$ had a pole divisor on the smooth $X$,
choose an algebraic smooth point of that divisor away
from its other exceptional loci and a transverse local
algebraic curve through it.
On a finite field of definition, \'etale local coordinates
and Hensel's lemma give algebraic points of that curve
approaching the chosen point through the dense agreement
open.  The pole makes their $h$-values unbounded, contrary
to Claim~B.131 for the actual map there.
Thus $h$ has no codimension-one pole.
A smooth affine scheme is normal, and the intersection
of its height-one valuation rings inside its function
field is its coordinate ring; therefore $h$ is regular.
At any remaining algebraic point choose an \'etale
coordinate neighbourhood and an algebraic line in its
coordinate space whose local lift is not contained in
the exceptional closed subset.
The same Henselian local inverse gives algebraic points
in the agreement open tending to that point.
Continuity of $f$ and of the regular $h$ gives equality
at the point.  All this takes place after one finite
extension for each chosen algebraic curve or local chart,
exactly the scope of Claim~B.131.

We have proved the required actual-point theorem for
every smooth geometrically integral affine source and
every finite extension of $B_p$.
A general smooth affine source is handled after a finite
extension splitting its finitely many geometric connected
components; smooth connected components are integral.
The resulting regular functions descend by their unique
actual values and finite Galois descent.

Finally apply Carlson--Haine--Wolf, Theorem~4.18,
in the implication from its actual-point condition
(4.18.4) to full pinned reconstruction, and
Corollary~2.32 for the localization by universal
homeomorphisms \citep{CarlsonHaineWolf2024}.
Their actual-point hypothesis asks for a morphism with
the prescribed geometric-point values on every affine
regular connected source with a rational point, with
target $\mathbb A^1$; it has no amenability hypothesis.
In characteristic zero a regular finite-type scheme is
smooth and its absolute weak normalization is itself.
The actual-point condition has therefore been proved
for every required source.  These exact source statements
were opened directly; the occurrence of the same
reduction in Claim~B.110 is not an additional premise.
The same argument starts with any finite $E/B_p$.
This proves the full statement, without a descent
assumption on the original geometric morphism.
\end{proof}
\paragraph{Check record.} checked only in this writing, 2026-11-18: GPT-6 Astra (OpenAI), ChatGPT Pro mode; internal effort value not exposed. Own detailed proof written and reread; primary source hypotheses matched. Root full reading of every paragraph passed: normalized norm classes; constant-gcd interpolation; algebraic root derivations and divisor cocycle; analytic-open lcm specialization; B135; constant-factor removal; graph degree; transverse-curve continuity; exact CHW reduction. Kummer sign clarification and field-input ledger inserted.


\paragraph{Field-dependent inputs used by the norm proof.}
The transfer to another algebraic mixed-characteristic
field must verify the following inputs; the argument
above does not identify that field with $B_p$.
First, finite extensions must be non-separably-closed
henselian fields with the chosen valuation, so that
Claim~B.131 supplies continuity of every actual point
map on their valued points.
Second, for a finite set of norm classes and grid centres,
normalization at one actual parameter must yield one
finite local field $M$ with normal $N=G_{E_1}\subset G_M$,
trivial cyclotomic character on $N$, a $G_M$-equivariant
actual norm assignment, and arithmetic terms $\beta_a$
that are explicit linear combinations of Kummer classes
of elements of $M$.
Third, after enlarging $M$ within the fixed field of $N$,
there must be cyclic extensions $L/M$ of arbitrarily
large $p$-power degree onto whose Galois groups $N$
surjects.  A nonzero Kummer class on $N$ with
$\mu_{p^\infty}$ in its fixed field supplies these
extensions by the proof of Claim~B.132.
Fourth, a local Kummer class over $L$ vanishing on
$N\cap G_L$ must have logarithm in $M$; the normality,
surjectivity, and trivial coefficient action just listed
prove precisely this assertion in Claim~B.132.
Finally, for every degree $d$ the stable Kummer kernels
over all finite scalar extensions must satisfy the
two-grid test: a quotient of monic degree-$d$ polynomials
whose values at the prescribed two grids lie in that
kernel is identically one.
That test is used both in obtaining the normalized
equivariance and in (META.NR14).
The principal-open Kummer decomposition, reduced-ratio
interpolation, root-type and lcm arguments, B.135,
finite-projection graph argument, and final CHW reduction
then use no other field-specific property.

\paragraph{Scope and provenance.}
This theorem addresses the full metabelian field and all
its finite extensions.  Other mixed-characteristic fields
listed in B8, including any remaining unbounded lower-central
case not covered elsewhere, are separate task rows.
The proof does not claim that point evaluations detect every
cohomology class on an arbitrary source; Claim~B.134
explicitly disproves that shortcut.
The norm argument uses principal affine bases, exact local
logarithmic errors, finite rational interpolation, and the
already independently read B.135.
No priority or external human verification is claimed.

\subsection{A uniform reconstruction argument for the unbounded nilpotent tower}

\begin{claim}{B.137}{Proved}{A tame signed Kummer relation and two finite grids}
Let $p$ be any prime, let $I=G_{K_p}$ be the fixed local inertia
group, and put
\[
 N=\bigcap_{c\ge1}\gamma_{c+1}(I),\qquad
 K=\overline{\mathbb Q}^{N}=K_{p,\mathrm{nil}}.
\]
There are $\tau\in I$, $b\in K$, and
$\epsilon\in\{1,-1\}$ with $\tau b\ne b$ and
$\tau i=i$, such that for every finite extension $E/K$
and every $a\in E^\times$,
\[
 \kappa_E(a)=0\quad\Longrightarrow\quad
                   \tau(a)/a^\epsilon\in\mu_\infty.
                                                        \tag{B.137.1}
\]
Uniformly in $p$, one may choose an odd prime
$\ell\nmid p(p-1)$, take $b=p^{1/\ell}$ in its
positive real embedding, and choose $\epsilon=1$.
At $p=2$ one may in particular take $\epsilon=1$, $b=2^{1/3}$ with
its positive real embedding, and $\tau b=\omega b$
for a primitive cube root $\omega$.
At $p=3$ one may take $\epsilon=-1$, $b=\sqrt3$,
and $\tau b=-b$.

For every integer $d\ge1$, set
\[
 \mathcal G_0=\{r+is:0\le r,s\le4d\},\qquad
                    \mathcal G_1=b+\mathcal G_0.
\]
If $P,Q\in\overline{\mathbb Q}[Z]$ are monic of degree
$d$, have nonzero values at these centres, and all the
classes $\kappa(P(a)/Q(a))$ vanish after finite extensions
of $K$ containing their data, then $P=Q$.
For the uniform choice, $\kappa_K(1-p^{1/\ell})\ne0$.
In particular $\kappa_K(1-2^{1/3})\ne0$ at $p=2$, and
$\kappa_K(3)\ne0$ at $p=3$; these classes remain nonzero
after every finite extension of $K$.
\end{claim}
\begin{proof}
Using the tame/wild inertia sequence of B.53, we
rederive the all-prime quotient structure appearing in B.57:
\[
                     I/N=Q_{(p)}\times T_{p'},\qquad
             T_{p'}=\prod_{\ell\ne p}\mathbb Z_\ell(1).
                                                        \tag{B.137.2}
\]
Indeed a finite inertia quotient has a normal wild
$p$-subgroup and cyclic prime-to-$p$ quotient; in a
nilpotent quotient its Sylow subgroups commute.  All
tame cyclic quotients survive, giving the displayed
prime-to-$p$ factor in the inverse limit.  The action
on roots of unity is abelian, so $K$ contains
$\mu_\infty$ and $N$ fixes every Tate coefficient.
The $p$-adic cyclotomic character on $T_{p'}$ has
finite image of order dividing $p-1$; it is
trivial at $p=2$ and has image $\{1,-1\}$ at $p=3$.

For the uniform choice, take an odd prime
$\ell\nmid p(p-1)$.  Such primes exist by Euclid's
argument.  The extension $K_p(p^{1/\ell})/K_p$
is cyclic tame of degree $\ell$: $\mu_\ell\subset K_p$
and the root has valuation $1/\ell$.
Choose an element of the pro-$\ell$ factor of
$T_{p'}$ with nontrivial image in this cyclic
quotient, and lift $(1,t)$ to $\tau\in I$.
Its $p$-cyclotomic value is one, since the
prime-to-$p$ part of that character has order
dividing $p-1$, while a pro-$\ell$ group cannot
map nontrivially to the pro-$p$ part.
The element $\tau$ therefore fixes all roots
of unity, including $i$, and sends
$b=p^{1/\ell}$ to $\omega b$ for a nontrivial
$\ell$-th root $\omega$.
This gives the claimed all-prime choice with
$\epsilon=1$.

For the displayed concrete alternatives, at $p=2$
choose $t\in T_{p'}$ acting nontrivially
on the tame cyclic extension $K_p(2^{1/3})/K_p$.
This extension has degree three: the valuation of
$2^{1/3}$ is $1/3$, and $\mu_3\subset K_p$.
At $p=3$, choose $t$ with cyclotomic value $-1$.
It acts by $-1$ on $\sqrt{-3}$ through its action
on $\mu_3$, and hence on $\sqrt3$, since
$i\in K_p$ is fixed.  The corresponding tame
quadratic extension is contained in $K$.
In each case lift the central element $(1,t)$ of
(B.137.2) to $\tau\in I$.
Its cyclotomic value is exactly the asserted
$\epsilon$; $\tau$ fixes $i$ in both cases.
No finite-order lift in $I$ is asserted or needed.

Suppose now that $\kappa_E(a)=0$.
Because $E$ contains all roots of unity,
$H^1(G_E,\mathbb Z_p(1))$ is a torsion-free group
of homomorphisms.  Rational vanishing is therefore
integral vanishing, so compatible $p$-power roots
of $a$ belong to $E$.
Take a finite Galois number field $D$ containing
$a$ and algebraic generators of $E/K$.
The field $E'=KD$ is stable under the local
decomposition group and contains the same roots
for every local conjugate $a_1,\ldots,a_m$ of $a$.
Write $M=D_v$ and
\[
 F=K_pD(\mu_{p^\infty}),\qquad
 R=F(a_j^{1/p^n}:1\le j\le m,\ n\ge1).
\]
Both are Galois over $K_p$, and $R\subset E'$.
The group $D_R=\operatorname{Gal}(R/F)$ embeds
by its Kummer characters into $\mathbb Z_p(1)^m$,
so it is torsion free.  The image $C_R$ of $N$
in $\operatorname{Gal}(R/K_p)$ is finite because
$E'/K$ is finite.
Centrality of $\tau$ modulo $N$ gives
\[
                         [\tau,D_R]\subset C_R\cap D_R=1.
                                                        \tag{B.137.3}
\]
Here $D_R$ is normal, so the commutator belongs
to both displayed groups.

Put $J=I_M=I\cap G_D$ and $J_0=J\cap\ker\chi_p$.
The latter group surjects onto $D_R$.
The natural coefficient action and (B.137.3) give
\[
                  \kappa_{J_0}(\tau a)
                            =\epsilon\kappa_{J_0}(a).
\]
Restriction from $H^1(J,\mathbb Q_p(1))$ to
$H^1(J_0,\mathbb Q_p(1))$ is injective.
In fact $J/J_0$ is abelian with nontrivial
cyclotomic character; for a cocycle $c$ on that
quotient, comparison of $c(gh)$ and $c(hg)$ gives
$(\chi(g)-1)c(h)=(\chi(h)-1)c(g)$, making $c$
a coboundary as soon as one $\chi(g)\ne1$.
Restriction from $G_M$ to $I_M$ is also injective,
since $\mathbb Q_p(1)^{I_M}=0$.
The equality thus holds in $H^1(G_M,\mathbb Q_p(1))$.
The rational Kummer kernel of the finite local
field $M$ consists exactly of its roots of unity:
valuation detects the uniformizer part, and the
principal-unit logarithm detects the remaining
part up to finite torsion.
This proves (B.137.1), including its finite-extension
quantifier.  The argument is the signed form of
the tame-central calculation in B.106; its proof
here uses neither a torsion-only kernel over $K$
nor a Baker--Brumer theorem at any prime.

Let $R_0(Z)=P(Z)/Q(Z)$ in the grid assertion.
For $a\in\mathcal G_0$, (B.137.1) and $\tau a=a$
give
\[
             S(a)\in\mu_\infty,\qquad
             S(Z)=\frac{\tau R_0(Z)}{R_0(Z)^\epsilon}.
\]
Write $S=A/B$ with monic numerator and denominator
of degree at most $2d$ (before cancelling if necessary).
Their degrees are equal and their leading coefficients
are one.  In the chosen complex embedding the polynomial
\[
 A(r+is)\overline A(r-is)-B(r+is)\overline B(r-is)
\]
has degree at most $4d$ separately in $r,s$ and
vanishes on the full $(4d+1)$ by $(4d+1)$ integer
grid.  Successive one-variable interpolation makes
it identically zero.  The invertible linear change
$z=r+is$, $w=r-is$ then gives
$A(z)\overline A(w)=B(z)\overline B(w)$ identically.
Fixing one $w$ where both denominators are nonzero
shows that $A$ and $B$ are proportional; monicity
makes them equal.  We have proved
\[
                           \tau R_0=R_0^\epsilon.
                                                        \tag{B.137.4}
\]

Apply the same argument on the second grid to
$R_0(b+Z)$.  It gives
$\tau R_0(\tau b+Z)=R_0(b+Z)^\epsilon$.
Together with (B.137.4), this says
$R_0(\tau b+Z)=R_0(b+Z)$.
A rational function in characteristic zero with
a nonzero additive period is constant: a finite
pole would have infinitely many translates, and
for a polynomial of positive degree the leading
term of its finite difference is nonzero.
The value at infinity of the original monic
equal-degree ratio is one.  Hence $R_0=1$.
There are finitely many grid values, so any finite
extensions used in their hypotheses may first be
replaced by one common finite compositum.

For the uniform choice, a zero class for $1-b$
would make $(1-\omega b)/(1-b)$ a root of unity.
In the positive real embedding of $b$,
\[
 |1-\omega b|^2-|1-b|^2
                  =2b(1-\operatorname{Re}\omega)>0,
\]
which excludes this possibility.
For the explicit alternative at $p=3$, a zero
class for $3$ would give $9\in\mu_\infty$ by
(B.137.1), which its valuation excludes.
At $p=2$, a zero class for $1-b$ would make
$(1-\omega b)/(1-b)$ a root of unity.
But $b>0$ in the chosen complex embedding and
\[
 |1-\omega b|^2=1+b+b^2>(1-b)^2=|1-b|^2.
\]
The ratio therefore has absolute value greater
than one.  Restriction to a finite extension is
injective in rational cohomology by restriction
and corestriction, or directly because a nonzero
homomorphism to $\mathbb Z_p$ cannot vanish on
a finite-index subgroup.  This proves the final
assertion.
\end{proof}
\paragraph{Check record.} Checked only in this writing, 18November2026: GPT-6 Astra (OpenAI), ChatGPT Pro mode. Root complete proof written and reread; finite-checks and arithmetic-notes readers passed the complete small-prime proof and the all-prime extension. Final all-prime clause rereading is recorded separately in CHECK\_EVENTS.tsv. Separate effort setting not exposed.


\begin{claim}{B.138}{Proved}{Uniform pinned reconstruction at every nilpotent inertia union}
For every prime $p$, let $K=K_{p,\mathrm{nil}}$ be as
in B.137.  For every finite extension $E/K$, the functor
from finite-type $E$-schemes localized at universal
homeomorphisms to their small \'etale topoi is fully
faithful for closed-point-preserving relative geometric
morphisms over $\operatorname{Spec}(E)_{\mathrm{et}}$,
with the relative structural isomorphism retained.
This settles this specified lower-central union,
and every finite extension of it, by one all-prime
argument without a torsion-only Kummer-kernel hypothesis.
\end{claim}
\begin{proof}
We verify the field-dependent inputs of B.136 and give
the transfer of its norm argument explicitly.
Write $N=G_K$.  This group is normal in the chosen
local decomposition group, because the lower-central
terms of its normal inertia subgroup are characteristic.
It fixes all roots of unity.  The nonzero classes of
B.137 show that $N$, and each of its finite-index
subgroups, have a nonzero homomorphism to $\mathbb Z_p$.
In particular no finite extension of $K$ is separably
closed.  All these fields are henselian for the chosen
mixed-characteristic valuation.

For every finite $E'/E$, actual point maps of pinned
relative topos morphisms are continuous on $E'$-points
for this valuation.  Here is the comparison being used.
Carlson--Haine--Wolf, Proposition~3.6, identifies the
inverse image of a point fibre with the actual relative
point fibre; the recollection in \S4.13 and Corollary~A.9 imply
that inverse images of quasi-compact separated \'etale
scheme objects are such scheme objects.  The inverse
image on $E'$-points of an \'etale image is therefore
an \'etale image.  Johnson--Tran--Walsberg--Ye,
Corollary~6.13, identifies this \'etale-open topology
with the valuation topology for a non-separably-closed
henselian valued field.  These hypotheses have just
been checked for every $E'$.  Affine charts give the
same assertion for all finite-type sources and targets.
This is the proof of B.131 with its field assumptions
verified here, not an application of its metabelian
statement to a different field
\citep{CarlsonHaineWolf2024,AN_JTWY2024}.

Let $X/E$ be smooth geometrically integral affine and
let $f$ be the actual point map of a pinned relative
morphism to the affine line.  The constant case is
immediate, so assume that $f$ is nonconstant.
Take any finite Noether normalization
$g:X\to\mathbb A^n$ of degree $d$ and the two grids
of B.137 for this $d$.  Each inverse image $S_a=f^{-1}(a)$
is proper closed.  Remove the branch locus and the
finitely many proper closed images $g(S_a)$, and choose
a nonempty principal open $V$ in their complement.
On it the finite cover is \'etale.  Put
\[
             P_t(Z)=\prod_{x\in g^{-1}(t)}(Z-f(x)).
                                                        \tag{B.138.1}
\]
Corestriction of the pulled-back rational Kummer classes
of $Y-a$ gives $\xi_a\in H^1(V,\mathbb Q_p(1))$
whose evaluation is $\kappa(P_t(a))$; the factor
$(-1)^d$ has zero rational Kummer class.

After a finite scalar extension, split the finitely
many prime boundary equations $h_j$ of $V$.
The principal-open Kummer calculation used in B.122
gives
\[
                  \xi_a=\beta_a+
                       \sum_j\lambda_{aj}\kappa(h_j),
                  \qquad\lambda_{aj}\in\mathbb Q_p.
                                                        \tag{B.138.2}
\]
That calculation depends on the geometric unit group
of this factorial open, and not on the particular
algebraic base field.  In detail its geometric units
modulo constants are freely generated by the $h_j$,
its Picard group is zero, and the Kummer exact sequences
and their inverse limits identify geometric first
cohomology with the corresponding finite-rank Tate
module.  Subtracting these unit classes leaves the
arithmetic class $\beta_a$.

Choose an algebraic parameter $t_*$ and include in a
finite Galois number field $D$ all scheme and boundary
coefficients, the grids, $t_*$, its finite fibre, the
values $P_{t_*}(a)$, and algebraic generators of $E/K$.
Replace the base by $E_1=KD$ and put $M=D_v$.
Evaluation at $t_*$ gives the exact normalization
\[
       \beta_a=\kappa(P_{t_*}(a))-
                    \sum_j\lambda_{aj}\kappa(h_j(t_*))
                       \quad\hbox{over }E_1.
                                                        \tag{B.138.3}
\]
Its right side is the restriction of an explicit class
over $M$.  Thus every datum in (B.138.2), including
its arithmetic term with the natural Tate coefficient
action, is fixed by $G_M$.
The field $E_1$ is stable under that local group and
$N_1=G_{E_1}$ is normal in $G_M$.

For $\sigma\in G_M$, compare evaluations of (B.138.2)
at $t$ and at $\sigma t$ over one common finite
extension of $E_1$ containing all point data.
Natural Kummer transport and (B.138.3) give
\[
           \kappa\!\left(
             \frac{P_{\sigma t}(a)}{(\sigma P_t)(a)}
                        \right)=0
                   \quad\hbox{for every grid centre }a.
\]
The centres are fixed by $\sigma$ because they belong
to $D$.  The finite-grid theorem B.137 therefore gives
\[
                         P_{\sigma t}=\sigma P_t.
                                                        \tag{B.138.4}
\]
The single field $M$ was chosen before $t$ and
$\sigma$.  This proves the required uniform norm
equivariance without assuming any descent of the
original topos morphism.

The explicit class in B.137 restricts nontrivially
to $N_1$.  Its integral image is $p^r\mathbb Z_p$ for
some $r\ge0$.  Adjoining a $p^r$-th root of its
argument inside $E_1$ makes the corresponding image
all of $\mathbb Z_p$; this is the integral Kummer
exact sequence, not a division of an arbitrary
cohomology class.  Enlarge $M$ by adjoining this
algebraic argument $\alpha\in E_1$ and the finitely
many elements $\mu_{p^s}\subset E_1$ for any chosen $s$.
The algebraic-number part of the enlarged local
field is still contained in the henselian field
$E_1$; the complete local field itself is not
being identified with a subfield of $E_1$.
Then
\[
                     L=M(\alpha^{1/p^s})
\]
is cyclic of degree $q=p^s$, and $N_1$ surjects onto
$\operatorname{Gal}(L/M)$.  We take $q>d+1$.
The group $N_1$ has not changed during this
enlargement, and (B.138.4) persists under restriction
to the smaller local Galois group.

The scalar logarithmic-error assertion also holds.
If $\theta\in H^1(G_L,\mathbb Q_p(1))$ restricts
to zero on $H=N_1\cap G_L$, it inflates from
$G_L/H$.  Every $n\in N_1$ acts trivially on this
quotient by conjugation, since $[n,G_L]\subset H$,
and acts trivially on its coefficients.
Consequently $n\theta=\theta$.
The natural local logarithmic map sends $\theta$
to an element of $L$ fixed by the surjective image
of $N_1$ in $\operatorname{Gal}(L/M)$, so that
element belongs to $M$.  This is precisely the
inflation argument of B.132 with its hypotheses
now verified for $N_1$.

We describe how these inputs recover the actual norm.
The algebraic-number field $M\cap\overline{\mathbb Q}$
is henselian and contained in $E_1$.  Its finite
extension obtained by adjoining $\alpha^{1/p^s}$
has completion $L$; the henselization--completion
comparison identifies it with $L\cap\overline{\mathbb Q}$.
Consequently $E_1(L\cap\overline{\mathbb Q})/E_1$
is finite.  This is the field on whose points
the continuity comparison is used below.
On an $L$-ball where $g$ has all its inverse branches,
(B.138.4) puts the coefficients of $P_t$ in $L$ at
algebraic parameters.  They are continuous by the
point-map continuity proved above.  The scalar-error
argument gives
\[
       \log P_t(a)-\sum_j\lambda_{aj}\log h_j(t)\in M.
                                                        \tag{B.138.5}
\]
Choose a parameter minimizing the sum of the gcd
degrees of $P_t$ with its $q-1$ nonidentity conjugates.
Continuity makes those degrees constant locally.
The exact logarithms of their ratios, obtained by
subtracting conjugates of (B.138.5), recover analytic
values by the exponential on a small neighbourhood.
Interpolation at twice the reduced degree makes
each reduced numerator and denominator analytic.
Invertibility is the monic coprime calculation
$\deg(\delta A B-A\delta B)\le2e-1$ in B.136.
Thus this step makes ratios analytic; it never
extends an arbitrary continuous point map to a
completion.

Over the split restriction-of-scalars coordinates
there are $q$ independent input blocks.
The mixed-derivative and $4d+1$-centre interpolation
argument of B.136 classifies each nonconstant root
of a reduced ratio by its unique input block.
The ratio cocycle makes its divisor
$\tau^i B-B+C_i$, where $C_i$ has constant roots.
Such a constant root would put a nonempty analytic
open in the proper closed image of a single $S_a$,
which is impossible.  A negative coefficient in
$B$ would give $q-1>d$ distinct conjugate poles
among the reduced denominators.  Their lcm cannot
divide the degree-$d$ actual $P_t$ after choosing
an algebraic parameter off the finitely many
nonzero analytic subresultants.  Hence $B$ is
effective.  These are the field-independent
divisor and specialization arguments proved in
full in B.136; their only numerical requirement
is the cyclic degree $q>d+1$ established here.

The monic polynomial $Q$ of divisor $B$ is the
analytic reduced denominator.  It depends only
on the identity block.  Differentiating the
ratio identity in that block gives
$d\log Q(t,a)=\sum_j\lambda_{aj}dh_j/h_j$ exactly.
B.135 makes its coefficients rational over $L$;
finite rational interpolation of their algebraic
values makes them rational over $\overline{\mathbb Q}$.
Divisor residues make all $\lambda_{aj}$ integral
and $Q(t,a)$ units on $V$.
Normalizing the remaining arithmetic class at
one algebraic parameter gives a fixed monic
polynomial $C$ with
\[
       \xi_a=\kappa(Q(t,a)) + \kappa(C(a)).
\]
At every algebraic parameter, B.137 now replaces
the metabelian two-grid step of B.136 and gives
$P_t(Z)=Q_t(Z)C(Z)$.
A positive-degree $C$ again puts all of $V$ in
one proper closed image $g(S_c)$.
Thus $C=1$, and every actual norm polynomial
has the recovered rational coefficients.

This applies to every finite normalization after
every finite scalar extension.  The graph-degree
argument of B.109, recalled in B.136, therefore
gives one rational function $h$ agreeing with $f$
on a nonempty open.  More precisely, if the
dominating graph components have degrees $e_i$,
a generic normalization of degree $d'$ separating
their output-coordinate generic fibres makes their
distinct minimal polynomials have total degree
$d'\sum e_i$.  They all divide the actual norm
polynomial of degree $d'$, so $\sum e_i=1$.
Continuity on finite valued extensions excludes
every pole along an algebraic transverse curve
and identifies the values at all remaining
points, exactly as in B.136.  Normality of smooth
$X$ makes $h$ regular.  Equality at every
geometric point makes $h$ invariant under the
original $G_E$, so it descends from $E_1$ to $E$.

We have proved the actual-point condition for
every required smooth geometrically integral
affine source.  Splitting geometric components
over a finite extension and using this same
uniqueness descends the result for every smooth
connected affine source.  Carlson--Haine--Wolf,
Theorem~4.18, implication (4.18.4) to (4.18.1),
now supplies full pinned reconstruction, and
Corollary~2.32 gives the stated localization
formulation in characteristic zero
\citep{CarlsonHaineWolf2024}.
All global algebraic scalar extensions used above are finite over
the fixed $K$, so the proof establishes the
assertion for every original finite $E/K$.
The nontrivial small-prime kernels in B.57
remain true; the signed finite-grid relation,
not their disappearance, resolves reconstruction.
\end{proof}
\paragraph{Check record.} Checked only in this writing, 18November2026: GPT-6 Astra (OpenAI), ChatGPT Pro mode. Root complete proof written and reread; finite-checks and arithmetic-notes readers passed the complete small-prime proof and the all-prime extension. Final all-prime clause rereading is recorded separately in CHECK\_EVENTS.tsv. Separate effort setting not exposed.


\subsection{Relative inertia fields in positive finite transcendence degree}

\begin{claim}{B.139}{Proved}{The algebraic closure of a finitely generated local subfield and its completion}
Fix a prime $p$ and an embedding
$D\subset\overline{\mathbb Q}_p$, where $D/\mathbb Q$
is finitely generated.  Write $M_0$ for its closure and
$\Omega$ for its algebraic closure inside
$\overline{\mathbb Q}_p$.
Then $M_0/\mathbb Q_p$ is finite, the chosen henselization
is $D^h=M_0\cap\Omega$, and completion induces an
equivalence between finite separable extensions of $D^h$
and those of $M_0$.
Restriction identifies $G_{M_0}$ with the decomposition
subgroup of $G_D$ for the chosen valuation.
Under this identification, $I_{M_0}$ is the absolute
Galois group of the chosen strict henselization field
of $D$ inside $\Omega$.

If $D_1/D$ is a finite extension in $\Omega$ and $M_1$
is its closure, then the same statements hold with
$D_1,M_1$ in place of $D,M_0$.
For every finite extension $L/M_1$ in
$\overline{\mathbb Q}_p$, its algebraic part
$L_{\mathrm{alg}}=L\cap\Omega$ is finite over $D_1^h$
and has completion $L$; the field $L$ admits an
$M_1$-basis lying in $L_{\mathrm{alg}}$.
There is an embedding $\Omega\hookrightarrow\mathbb C$
extending any prescribed complex embedding of
$\overline{\mathbb Q}\subset\Omega$.
\end{claim}
\begin{proof}
Choose finitely many generators of $D/\mathbb Q$.
Each is algebraic over $\mathbb Q_p$, so together they
belong to a finite extension $M'/\mathbb Q_p$.
The closure $M_0$ lies in $M'$ and contains
$\mathbb Q_p$, because $\mathbb Q$ is dense in
$\mathbb Q_p$.  It is therefore a finite extension
of $\mathbb Q_p$.
In particular the restricted value group of $D$
is discrete and its residue field is finite.

We use two exact opened results.
Temkin, Theorem~1.9.5(ii), says that a real-valued
field is henselian if and only if it is relatively
separably closed in its completion; his construction
in \S1.9.7 and Lemma~1.9.8 identify the henselization
with that relative separable closure
\citep{AN_TemkinValuations}.
Thus $D^h=M_0\cap\Omega$ in characteristic zero.
The Stacks Project, Lemma~15.115.2, Tag~09EJ,
says that every finite separable extension of the
fraction field of the completion of a DVR descends
to a finite separable extension of its fraction field
\citep{AN_Stacks09EJ}.
Applied to the valuation ring of $D^h$, this gives
essential surjectivity of completion on finite
separable extensions.

Here is the remaining full-faithfulness argument.
For a finite extension $F/D^h$, henselianity gives
one valuation and the completion is $FM_0$.
The same henselian criterion, applied to all finite
extensions, says that $F$ is relatively algebraically
closed in its completion.  Also $F$ and $M_0$ are
linearly disjoint over $D^h$: for a finite Galois
closure this follows because $D^h$ is relatively
separably closed in $M_0$, and the general finite
case follows by passage to that Galois closure.
Consequently $[FM_0:M_0]=[F:D^h]$.
A $D^h$-embedding $F_1\to F_2$ preserves the unique
valuation and extends continuously to the completions.
Conversely, an $M_0$-embedding of the completions
sends $F_1$ to elements algebraic over $D^h$ inside
the completion of $F_2$; those elements belong to
$F_2$ by its relative algebraic closedness.
The two constructions are inverse.
Finite products give the finite-etale-algebra version.

For a finite $L/M_0$, choose a descended finite
$F/D^h$ and the embedding supplied by its completion
isomorphism.  Its elements are algebraic over $D$,
so $F\subset\Omega$.  Since $F$ is relatively
algebraically closed in $L$ and $L/F$ is the
completion, every element of $L\cap\Omega$ belongs
to $F$.  Hence $F=L\cap\Omega$.
This also proves that an $F/D^h$-basis is an
$M_0$-basis of $L$ and lies in $\Omega$.
The same proof applies to $D_1$, which is still
finitely generated over $\mathbb Q$ and has finite
local closure $M_1$.

The equivalence of finite Galois categories identifies
their absolute Galois groups by restriction.
The valuation decomposition subgroup of $G_D$ fixes
$D^h$ by the defining universal property of
henselization.  Its inertia subgroup fixes the
strict henselization, since finite unramified
extensions correspond to finite residue-field
extensions.  This proves the claimed identifications.

Finally $\Omega/\overline{\mathbb Q}$ has finite
transcendence degree.  Extend the prescribed embedding
of $\overline{\mathbb Q}$ by choosing algebraically
independent complex images of a transcendence basis.
Extend through the algebraic extension into the
algebraically closed field $\mathbb C$.
No compatibility between this complex embedding and
the chosen $p$-adic topology is asserted or required.
\end{proof}
\paragraph{Check record.} checked only in this writing, 2026-11-18: GPT-6 Astra (OpenAI), ChatGPT Pro mode; internal effort value not exposed. Own detailed proof written and reread; primary source hypotheses matched. Root full reading of B139 passed; basis wording correction inserted; exact Temkin/Stacks primary passages independently opened and read


\begin{claim}{B.140}{Proved}{Reconstruction for every finite derived length and the relative nilpotent union}
Keep the data of Claim~B.139 and put
$I=I_{M_0}\subset G_D$ using its restriction embedding.
Write $I^{(0)}=I$ and
$I^{(r+1)}=\overline{[I^{(r)},I^{(r)}]}$.
Define subfields of $\Omega$ by
\[
 F_{D,r}=\Omega^{I^{(r)}}\quad(r\ge0),\qquad
 K_{D,\mathrm{nil}}
   =\Omega^{\bigcap_{c\ge1}\gamma_{c+1}(I)}.
\]
For every prime $p$, every such finitely generated
$D\subset\overline{\mathbb Q}_p$, and every finite
extension $E$ of any $F_{D,r}$ with finite $r$, or of
$K_{D,\mathrm{nil}}$,
inside $\Omega$, the functor
\[
 \operatorname{Sch}^{\mathrm{ft}}_E[\operatorname{UH}^{-1}]
       \longrightarrow \operatorname{RTop}^{\mathrm{pin}}_E
\]
is fully faithful.  The target consists of geometric
morphisms over $\operatorname{Spec}(E)_{\mathrm{et}}$
that preserve closed points, with their relative
structural isomorphisms retained.
The same conclusion holds for any subfield $F\subset E$
for which $E/F$ is algebraic, and for every finite
extension of such an $F$.

These classes occur in every finite transcendence
degree over $\mathbb Q$.  For $r\ge1$ the derived
fields contain all roots of unity and have unbounded
wild ramification index, as does the nilpotent union.
Thus the positive assertion includes fields of
positive transcendence degree beyond the bounded-wild
criterion and beyond sub-$p$-adic fields.
It does not classify all mixed-characteristic
henselian fields of finite transcendence degree.
\end{claim}
\begin{proof}
All fields in the statement are algebraic over $D$.
Write $K_D^{sh}=F_{D,0}=\Omega^I$ for the relative
strict henselization field.  It is henselian
with residue field $\overline{\mathbb F}_p$, and
its algebraic extensions retain henselianity and
that residue field.  For $r\ge1$ the displayed
kernels are characteristic in $I$, so they are
normal in $G_{M_0}$.  They fix every root of unity
because the cyclotomic action is abelian.
We verify the arithmetic inputs of the norm proof,
including their stability under finite extensions,
before transporting its geometric part.

\medskip\noindent\textbf{The stable kernel at any finite derived length.}
Fix $r\ge2$, and set $K=F_{D,r}$ and
$A=F_{D,r-1}$ in this paragraph.  Define
\[
 W_{D,r}=\{a\in F_{D,r}^\times:a^m\in A^\times
                                  \text{ for some }m\ge1\}.
\]
For every finite $E/F_{D,r}$ we claim
\[
 \ker\bigl(E^\times\longrightarrow
                        H^1(G_E,\mathbb Q_p(1))\bigr)=W_{D,r},
 \qquad W_{D,r}\cap(1-W_{D,r})=A\setminus\{0,1\}.
                                                        \tag{META.RL1}
\]
For $r=2$ this is the relative metabelian kernel.
For larger $r$ the quotient need not be metabelian;
the following finite-orbit argument handles the
additional cohomological term.
Suppose that $a\in E^\times$ has zero class.
Choose a finite Galois extension $D_1/D$ containing
$a$ and generators of $E/F_{D,r}$, and replace $E$ by
$E_1=F_{D,r}D_1$.
Let $M_1$ be the closure of $D_1$ and put
\[
 J=I\cap G_{D_1}=I_{M_1},\quad H=I^{(r)}\cap J,
 \quad\Gamma=J/H,\quad C=(I^{(r-1)}\cap J)/H.
\]
The local class of $a$ in $H^1(J,\mathbb Q_p(1))$
inflates from $\Gamma$, because it vanishes on $H$.
The group $C$ is abelian and fixes the coefficient
module $V=\mathbb Q_p(1)$.
Every $g\in I^{(r-1)}$ normalizes $J$, because
$D_1/D$ is Galois, and acts trivially on $C$, on
$\Gamma/C$, and on $V$.
For the first assertion use
$[I^{(r-1)},I^{(r-1)}]\subset I^{(r)}$;
for the second use normality of $I^{(r-1)}$ in $I$.
Inflation--restriction gives the exact sequence
\[
 0\longrightarrow H^1(\Gamma/C,V)
 \longrightarrow H^1(\Gamma,V)
 \longrightarrow H^1(C,V)^{\Gamma/C}.
                                                        \tag{META.RL1a}
\]
The action of $g$ is the identity on both outer
terms.  Therefore $(g-1)^2=0$ on the middle term.
Let $\theta$ be the inflated Kummer class of $a$.
The element $a$ has a finite $g$-orbit in $D_1$,
so $g^m(a)=a$ for some positive integer $m$.
Naturality of Kummer theory gives $g^m\theta=\theta$.
On this vector the unipotence equation gives
\[
       0=(g^m-1)\theta=m(g-1)\theta.
\]
Since the coefficient field is $\mathbb Q_p$,
we obtain $g\theta=\theta$.
No finite-order assertion about the action on the
whole cohomology group is used.
Thus $\kappa_J(g(a)/a)=0$ for every
$g\in I^{(r-1)}$.

All these elements belong to the one finite local
field $M_1$: the automorphisms in $G_{M_0}$ preserve
$D_1$ and its valuation closure.
Restriction from $H^1(G_{M_1},\mathbb Q_p(1))$ to
$H^1(I_{M_1},\mathbb Q_p(1))$ is injective because
the coefficient module has no inertia invariants.
Valuation and the principal-unit logarithm show
that the rational Kummer kernel of $M_1$ is its
finite group of roots of unity.
One integer annihilating that group makes $a^m$
fixed by every $g\in I^{(r-1)}$, so $a^m\in A$.
Since $A$ contains all roots of unity, its radical
extensions are abelian and lie in $F_{D,r}$.
This proves that $a\in W_{D,r}$.
Conversely all $p$-power divisions of an element
of $W_{D,r}$ are radicals of an element of $A$ and
belong to $F_{D,r}$, so their Kummer class vanishes.
Rational and integral vanishing agree since $G_E$
acts trivially on $\mathbb Z_p(1)$ and continuous
homomorphisms into $\mathbb Z_p$ are torsion free.

If $a,1-a\in W_{D,r}$, then for any $g\in I^{(r-1)}$
there are roots of unity $\zeta,\eta\in A$ with
$g(a)=\zeta a$ and $g(1-a)=\eta(1-a)$.
If $\zeta\ne\eta$, their sum expresses
$a=(1-\eta)/(\zeta-\eta)\in A$.
If they are equal, their sum forces both to be one.
Either some unequal pair already puts $a$ in $A$,
or every $g$ fixes $a$.  This proves the second
formula in (META.RL1); its converse is immediate.

There exists $b\in F_{D,r}\setminus A$.
Indeed the wild inertia group of the finite local
field $M_0$ is free pro-$p$ of infinite rank by
Iwasawa, Theorem~2(ii) \citep{AN_Iwasawa1955}.
For each $h\ge1$ it surjects onto
$\operatorname{UT}_{2^h+1}(\mathbb F_p)$ by
sending free generators to adjacent elementary
transvections.  The exact commutator identity
\[
 [1+E_{ab},1+E_{bc}]=1+E_{ac}\qquad(a<b<c)
\]
shows, by a balanced induction, that its $h$-th
derived subgroup contains $1+E_{1,2^h+1}\ne1$.
Hence $I^{(h)}$ is nontrivial for every $h$.
For $h\ge1$ it is contained in the pro-$p$
wild inertia, since $I/P_{M_0}$ is abelian.
Every nontrivial pro-$p$ group has nonzero
abelianization: take a nontrivial finite $p$-group
quotient and then a quotient of order $p$.
Thus $I^{(r)}$ is proper in $I^{(r-1)}$, proving
the asserted strict field extension.
For any chosen $b\notin A$, at least one of
$b$ and $1-b$ is outside $W_{D,r}$, by (META.RL1).
It has a nonzero Kummer class, which remains
nonzero after every finite extension by rational
restriction and corestriction.

For degree $d$, use the Gaussian integer grid
$\mathcal G_0=\{u+iv:0\le u,v\le4d\}$ and
$\mathcal G_1=b+\mathcal G_0$.
If a monic equal-degree ratio $R=P/Q$ takes values
in the stable kernel at these grids, (META.RL1)
implies $g(R(a))/R(a)\in\mu_\infty$ for every
$g\in I^{(r-1)}$.
Such a $g$ fixes the first grid.
In the fixed embedding $\Omega\hookrightarrow\mathbb C$,
write the numerator and denominator of $gR/R$
as monic polynomials $U,V$ of equal degree at
most $2d$ before cancellation.
The polynomial
$U(u+iv)\overline U(u-iv)-V(u+iv)\overline V(u-iv)$
has degree at most $4d$ in each variable and
vanishes on the full grid.  It is zero by
successive interpolation.  The invertible
substitution $z=u+iv,w=u-iv$ gives
$U(z)\overline U(w)=V(z)\overline V(w)$.
Fixing an allowed $w$ and using monicity gives
$U=V$, and therefore $gR=R$.
On the translated grid the same calculation gives
$gR(gb+Z)=R(b+Z)$.
Choose $g$ moving $b$.
The identity $gR=R$ then gives a nonzero additive
period for $R$.  Its finite pole set, or the
leading term of a polynomial finite difference,
forces it to be constant; monicity makes it one.
This proves the required two-grid test for the
relative finite-derived-length field.

\medskip\noindent\textbf{The relative lower-central union.}
Now set $N=\bigcap_c\gamma_{c+1}(I)$ and
$K=\Omega^N$.
Finite inertia quotients have a normal $p$-subgroup
and a cyclic prime-to-$p$ quotient.
Their nilpotent quotients have commuting Sylow
subgroups, and all tame cyclic quotients survive.
Consequently
\[
       I/N=Q_{(p)}\times\prod_{\ell\ne p}\mathbb Z_\ell(1).
                                                        \tag{META.RL2}
\]
Let $e_0=e(M_0/\mathbb Q_p)$ and choose an odd
prime $\ell\nmid e_0p(p-1)$.
The extension obtained by adjoining
$b=p^{1/\ell}$ to the relative strict henselization
field is cyclic tame of degree $\ell$:
its value $1/\ell$ is absent from the value group
$e_0^{-1}\mathbb Z$, while all $\ell$-th roots
of unity are already present.
Choose a central element of (META.RL2) in the
pro-$\ell$ tame factor with nontrivial image on
this radical, and lift it to $\tau\in I$.
Its $p$-cyclotomic value is one, since
$\ell\nmid p(p-1)$; inertia also fixes all
prime-to-$p$ roots of unity.  Hence $\tau$ fixes
every root of unity and sends $b$ to $\omega b$
for some $\omega\ne1$.

For completeness, the stable-kernel consequence
of its centrality has exactly the same local proof
in this relative field.
If $a$ has zero Kummer class over a finite $E/K$,
its compatible $p$-power roots lie in $E$.
Choose finite Galois $D_1/D$ containing $a$ and
generators of $E/K$.
Let $F$ be the compositum of $D_1$, the relative
strict henselization field, and $\mu_{p^\infty}$;
let $R/F$ be generated by all $p$-power roots of
the finitely many local conjugates of $a$.
Then $R\subset KD_1$, both $F$ and $R$ are
Galois over the relative strict henselization
field, and $\operatorname{Gal}(R/F)$ is a
torsion-free subgroup of $\mathbb Z_p(1)^m$.
The image of $N$ in $\operatorname{Gal}(R/K_D^{sh})$
is finite because $KD_1/K$ is finite.
Centrality of $\tau$ modulo $N$ makes its
commutator with $\operatorname{Gal}(R/F)$ both
finite and an element of that torsion-free
group, hence trivial.
Writing $J=I_{M_1}$ and $J_0=J\cap\ker\chi_p$,
the group $J_0$ surjects onto this radical group.
Therefore $\kappa_{J_0}(\tau a)=\kappa_{J_0}(a)$.
Restriction to $J_0$ is injective on
$H^1(J,\mathbb Q_p(1))$: its inflation kernel
is the first cohomology of the abelian group
$J/J_0$ with nontrivial scalar character, and
the cocycle identity
$(\chi(g)-1)c(h)=(\chi(h)-1)c(g)$ makes that
cohomology zero;
restriction from $G_{M_1}$ to $J$ is also injective.
The finite-local-field Kummer kernel now gives
\[
             \kappa_E(a)=0\quad\Longrightarrow\quad
                              \tau a/a\in\mu_\infty.
                                                        \tag{META.RL3}
\]

Choose the complex embedding of $\Omega$ so that
the rational radical $b=p^{1/\ell}$ is positive.
Then
\[
 |1-\omega b|^2-|1-b|^2
                    =2b(1-\operatorname{Re}\omega)>0.
\]
Thus (META.RL3) proves that $1-b$ has nonzero
Kummer class over $K$ and every finite extension.
For the grids $\mathcal G_0$ and $b+\mathcal G_0$,
the same absolute-square interpolation first gives
$\tau R=R$ and then $R(\tau b+Z)=R(b+Z)$.
The nonzero period gives $R=1$.
This verifies both arithmetic inputs uniformly
in the relative nilpotent case, including $p=2,3$.

\medskip\noindent\textbf{Verification of the norm proof over $\Omega$.}
For $K=F_{D,r}$ with $r\ge2$, or for the
nilpotent union, let $E/K$ be finite.
For finitely many coefficients of a smooth affine
source, a finite projection, its principal-open
boundary, grid centres and actual normalizing
values, choose one finite Galois $D_1/D$
containing all those data and generators of $E/K$.
Its completion $M_1$ is finite over $\mathbb Q_p$
by Claim~B.139.
Put $E_1=KD_1$ and $N_1=G_{E_1}$.
Characteristicity of the original inertia kernel
makes $N_1$ normal in $G_{M_1}$, and it fixes
$\mathbb Q_p(1)$.
The explicit nonzero classes just constructed
remain nonzero on $N_1$.  Thus every finite
extension of $K$ is non-separably-closed.
The continuity proof of Claim~B.131 applies:
inverse images of quasi-compact separated etale
objects and relative point fibres give etale-open
continuity, and the opened Johnson--Tran--Walsberg--Ye,
Corollary~6.13, identifies that topology with the
valuation topology on each such finite extension
\citep{CarlsonHaineWolf2024,AN_JTWY2024}.

We spell out the arithmetic normalization, so that
no descent of a topos morphism is inserted.
For the monic norm polynomial $P_t$ on a principal
open $V$ of affine space, pulled-back and corestricted
Kummer classes have the decomposition
\[
 \xi_a=\beta_a+\sum_j\lambda_{aj}\kappa(h_j),
                     \qquad\lambda_{aj}\in\mathbb Q_p,
                                                        \tag{META.RL4}
\]
because the geometric unit group modulo constants
is free on the prime boundary equations $h_j$
and the geometric Picard group is zero.
Evaluation at a fixed $t_*$, all of whose data
were included in $D_1$, identifies $\beta_a$ with
the finite linear combination
\[
       \kappa(P_{t_*}(a))-
                     \sum_j\lambda_{aj}\kappa(h_j(t_*)).
\]
It is the restriction of a class over $M_1$.
Consequently (META.RL4) is $G_{M_1}$-invariant,
with its natural Tate coefficient action.
Comparison at $t$ and $\sigma t$ puts every
grid value of $P_{\sigma t}/\sigma P_t$ in the
stable kernel, over one common finite extension
of $E_1$.
The respective two-grid test proved above gives
$P_{\sigma t}=\sigma P_t$ for all algebraic
parameters and all $\sigma\in G_{M_1}$.

A nonzero integral Kummer class on $N_1$ has
image $p^r\mathbb Z_p$.
Its argument has a $p^r$-th root $\alpha$ in
$E_1$, since its reduction modulo $p^r$ is zero.
The class of $\alpha$ has image all of
$\mathbb Z_p$.
Enlarge the finite local $M_1$ inside the local
fixed field of $N_1$ by $\alpha$ and
$\mu_{p^s}$.  The extension
$L=M_1(\alpha^{1/p^s})$ is cyclic of degree
$q=p^s$ and receives a surjection from $N_1$.
Take $q>d+1$ for the projection degree $d$.
Claim~B.139 shows that its algebraic part
$L\cap\Omega$ is finite over $M_1\cap\Omega$
and that $L$ has an $M_1$-basis in $L\cap\Omega$.
The field $M_1\cap\Omega$ is contained in $E_1$:
the finitely many algebraic elements used to enlarge
$M_1$ belong to $E_1$, and the chosen henselization
of their field belongs to the henselian field $E_1$.
In particular $E_1(L\cap\Omega)/E_1$ is finite,
and its algebraic scalar coordinates are dense
in the local analytic coordinates.

If a local class in $H^1(G_L,\mathbb Q_p(1))$
vanishes on $N_1\cap G_L$, it inflates from
the quotient.  Conjugation by $N_1$ acts trivially
on that quotient and on the coefficients, because
$N_1\triangleleft G_{M_1}$, $L/M_1$ is Galois,
and $\chi_p(N_1)=1$.
Its local logarithm is therefore fixed by
$\operatorname{Gal}(L/M_1)$ and belongs to $M_1$.
This proves the scalar-error input of the norm
argument for these fields.

The remaining proof in Claim~B.136 now works
over $\Omega$ in place of $\overline{\mathbb Q}$.
To make that transfer explicit, the norm
coefficients at parameters in $L\cap\Omega$
belong to that field by the proved local
equivariance and are continuous by the finite
extension continuity just checked.
Minimizing the finitely many conjugate gcd
degrees and interpolating the exact logarithms
of reduced monic ratios makes those ratios
analytic without assuming the norm was analytic.
In restriction-of-scalars coordinates the
mixed-derivative argument classifies their
nonconstant roots by input block.
Their cocycle has divisor $\tau^jB-B+C_j$.
A constant root belongs to $\Omega$ after
one specialization at an actual $\Omega$-point;
it is excluded by the proper closed actual fibre.
A negative coefficient of $B$ would give
$q-1>d$ distinct denominator roots.
Nonzero analytic subresultants stay nonzero
at some scalar-coordinate $\Omega$-point,
by density and the analytic identity theorem,
contradicting divisibility by the degree-$d$
actual norm.  Thus $B$ is effective and its
monic polynomial $Q$ is the already analytic
reduced denominator, depending on one block.

The exact equations
$d\log Q(t,a)=\sum_j\lambda_{aj}dh_j/h_j$
and B.135 put $Q$ in $L(t)[Z]$.
Its values at actual parameters have coefficients
in the algebraically closed field $\Omega$;
finite rational interpolation therefore puts
$Q$ in $\Omega(t)[Z]$.
Divisor residues make all $\lambda_{aj}$
integral and $Q(t,a)$ a unit on $V$.
One actual normalization then gives a fixed
monic $C\in\Omega[Z]$ with
$\xi_a=\kappa(Q(t,a))+\kappa(C(a))$.
The two-grid test proves $P_t=Q_tC$ everywhere
on $V$.  A positive-degree $C$ would put
all of $V$ in the proper closed image of
one actual fibre, so $C=1$.
All of these arguments use algebraicity over
$D$, not algebraicity over $\mathbb Q$.
The complex embedding is used only for the
finite-grid absolute-square identities; the
analytic calculations use the $p$-adic topology.

We have recovered the norm polynomial for
every finite Noether projection.
The self-contained graph-degree proof in
Claim~B.136, over $\Omega$, gives a
unique rational function agreeing with the
actual map on a dense open.
Its generic projection conditions descend
from $\Omega(Y)$ to a nonempty open in
coefficient space over $\Omega$ by clearing
denominators and expanding in $Y$.
All coefficients used lie in some finite
extension of $E$, since $\Omega/E$ is algebraic.
Transverse algebraic curves and valuation
continuity exclude codimension-one poles
and identify every exceptional point value.
Normality makes the rational function regular;
its original $G_E$-equivariance descends it to $E$.
Splitting geometric connected components over
a finite extension handles every smooth affine
source required by Carlson--Haine--Wolf,
Theorem~4.18(4.18.4).  That theorem and
Corollary~2.32 give the asserted full pinned
reconstruction and localization
\citep{CarlsonHaineWolf2024}.

\medskip\noindent\textbf{Algebraic descent and the size of the new family.}
Suppose $F\subset E$ and $E/F$ is algebraic.
Their algebraic closures can be identified with
the same $\Omega$.
A relative pinned map over $F$ base changes
to one over $E$: for finite scalar extensions
this is slicing at the finite etale constant
object.  These slice inverse-image functors
are compatible under further finite extension.
Every finitely presented etale object, its
morphisms and every finite covering family
over $E$ descend to one finite subextension
$E_i/F$.  The compatible functors on these
finite levels therefore define the inverse-image
functor on the limit etale site; they preserve
finite limits and covering families and so
define the required geometric morphism.
Closed-point preservation persists, since
every coordinate algebraic over $E$ is
algebraic over $F$.
For a smooth affine source over $F$, the
actual-point theorem over $E$ gives a regular
function after base change to $E$.
Viewed over $\Omega$, it has the original
$G_F$-equivariant actual point values.
Two regular functions with the same values
on all geometric points are equal on a
reduced finite-type scheme.  Hence this
function is $G_F$-invariant and descends
to $F$ by the coefficient form of faithfully
flat descent.  CHW's actual-point reduction
therefore gives reconstruction over $F$.
For a finite $F'/F$, the compositum $EF'/E$
is finite and is covered by what we proved;
applying the same descent gives the claim
for $F'$.

The cases $r=0,1$ of the finite-derived statement
follow by this algebraic descent from $F_{D,2}$.
For each integer $s\ge0$, choose algebraically
independent $a_1,\ldots,a_s\in\mathbb Q_p$
and put $D=\mathbb Q(a_1,\ldots,a_s)$.
Such a finite tuple exists because the
algebraic closure of a finitely generated
field over $\mathbb Q$ is countable whereas
$\mathbb Q_p$ is uncountable.
Its closure is $\mathbb Q_p$.
The stated fields are algebraic over $D$, so
have transcendence degree $s$.
Those with derived index at least one, and
the nilpotent union, contain $\mu_{p^\infty}$, which gives
unbounded $p$-parts of ramification indices,
and indeed contain $\mu_\infty$.
A sub-$\ell$-adic field for any prime $\ell$
embeds into a finitely generated extension
of a finite $\ell$-adic field; the algebraic
constant field of such an extension is
finite over that local field and has only
finitely many roots of unity.
Already $F_{D,0}$ contains infinitely many
prime-to-$p$ roots of unity.  Thus every stated
$F_{D,r}$ and the nilpotent union are not
sub-$\ell$-adic.
For $s>0$ they are also outside the earlier
number-field-algebraic $B_p$ and nilpotent
union statements.
This proves the specified enlargement of
the positive classes without turning it
into a classification of the remaining
mixed-characteristic fields.
Finally, the union of all the finite derived
fields is $\Omega$: every finite quotient of
$I$ is an extension of a finite $p$-group by
a cyclic prime-to-$p$ group, hence solvable.
The algebraically closed field $\Omega$ fails
reconstruction.  A nonidentity automorphism of
$\overline{\mathbb Q}$ extends to $\Omega$
after fixing a transcendence basis and gives
a semilinear pinned automorphism of the affine
line's etale topos over the terminal base
$\operatorname{Spec}(\Omega)_{\mathrm{et}}$.
Its point function fixes every rational number
but moves some algebraic number.
No regular polynomial over $\Omega$ has this
point function, since agreement with $T$ on
infinitely many rational points forces that
polynomial to equal $T$.
Thus the finite-derived conclusion has not
been extended through that infinite union.
\end{proof}
\paragraph{Check record.} checked only in this writing, 2026-11-18: GPT-6 Astra (OpenAI), ChatGPT Pro mode; internal effort value not exposed. Own detailed proof written and reread; primary source hypotheses matched. Root full reading of every paragraph of B140 passed, including the remainder after RL1, finite-derived filtration, relative tame grids, all norm-input transfers, finite-site algebraic descent, examples, and the infinite-union boundary; exact Temkin/Stacks passages independently opened and read


\paragraph{Source and check scope.}
The new comparison inputs were opened on
18 November 2026: Temkin's author PDF,
Theorem~1.9.5(ii), \S1.9.7 and Lemma~1.9.8,
printed p.~25; and the entire statement and
proof of Stacks Tag~09EJ.
The proof explicitly replaces finite number
fields by finite extensions of $D$ with finite
local completions, and explicitly replaces
algebraic numbers by elements of $\Omega$.
The complex and $p$-adic embeddings are separate.
Second-reading events and status are recorded
separately; no priority, human verification,
or verification by a different model is claimed.

\subsection{A uniform genus-two criterion for the prescribed point}

\begin{claim}{B.141}{Proved}{Canonical-point classes and the unique potentially bad twist}
Let $C/k$ be a smooth proper geometrically connected curve of genus
two over $k=\mathbf F_q$ of characteristic $p$, let $P\in C(k)$, and put
\[
 J=J_C,\qquad \delta=[K_C-2P]\in J(k).
\]
Let $\ell\ne p$ be an odd prime for which
$f=\operatorname{ord}_{\ell}(p)$ is even.
For every nonzero additive character
$\chi:J(k)\twoheadrightarrow\mathbf Z/\ell$ with
$\chi(\delta)\ne0$, the associated constant cyclic
$P$-split Lang cover $D_\chi\to C$ has
\[
                  p\nmid \frac{\#J_{D_\chi}(k)}{\#J(k)}.
                                                            \tag{B.141.1}
\]
In particular, if $\delta\notin\ell J(k)$, this prescribed-point
class-number witness exists already over $k$.
If $r=\dim_{\mathbf F_\ell}J(k)/\ell J(k)$, this condition gives
at least $\ell^{r-1}$ distinct pointed covers of degree $\ell$,
where an identification of the deck group with
$\mathbf Z/\ell$ is not retained.

More precisely, fix any nonzero $\chi$, put $a=\chi(\delta)$,
and index its constant twists by $t\in\mathbf Z/\ell$ so that
a closed point of degree $m$ acquires the added Frobenius
translation $mt$.  Among these $\ell$ twists, the only
possibly bad twist is
\[
                              2t+a=0.                       \tag{B.141.2}
\]
This twist need not be bad.  In particular, the number of bad
twists is at most one, regardless of the index of
$\langle p\rangle$ in $(\mathbf Z/\ell)^\times$.

If $\chi(\delta)\ne0$, the same cover after base change to
every extension $k_n/k$ of degree $n$ prime to $\ell$ is
again a $P$-split class-number witness.
The sufficient condition here is the nonvanishing of the class
of $\delta$ in $J(k)/\ell J(k)$ at the stated field; nonzero
$\ell$-primary torsion by itself is not this condition.

A nonempty explicit family is
\[
 C_p:\ y^2=x^6+1,\qquad P=(0,1),\qquad
          p\ge5,\quad p\equiv2\text{ or }5\pmod9 .
                                                            \tag{B.141.10}
\]
Here one may take $\ell=3$.
For $p=5$, the character specified in the proof gives a cover
of genus four with
\[
  \#J_{C_5}(\mathbf F_5)=36,\qquad
  \#J_D(\mathbf F_5)=432,\qquad
  \frac{\#J_D(\mathbf F_5)}{\#J_{C_5}(\mathbf F_5)}=12.
                                                            \tag{B.141.11}
\]
\end{claim}
\begin{proof}
We give the pointed construction and the complete character calculation.
The Lang map $F_q-1:J\to J$ is a separable isogeny with
constant kernel $J(k)$.
Factoring its source by $\ker\chi$ gives a separable isogeny
with constant cyclic kernel of order $\ell$.
Its pullback under $j_P:R\mapsto[R-P]$ has a chosen rational
point over $P$, and all of that fibre splits.
The pullback is geometrically connected: the Abel--Jacobi map
induces the isomorphism on prime-to-$p$ geometric abelian
fundamental groups, so the nonzero cyclic quotient remains
nonzero on the curve.  Equivalently, pullback along $j_P$
identifies the degree-zero $\ell$-torsion of the dual Jacobian
with that of $J$, by its principal polarization.
This is the same Lang and Kummer construction used in B.2;
no point-section separation theorem is used.

Choose a primitive $\ell$-th root of unity $\zeta$ in characteristic
zero.  For $i\in(\mathbf Z/\ell)^\times$ write
\[
       S_i=\sum_{R\in C(k)}\zeta^{\,i\chi([R-P])},
                     \qquad a=\chi(\delta).
\]
Expanding the unramified character Euler product by effective
divisors gives
\[
                        L_i(T)=1+S_iT+q\zeta^{ia}T^2.
                                                            \tag{B.141.3}
\]
Here is a direct Riemann--Roch verification, including the last
coefficient.  Tensoring by $\mathcal O(-dP)$ identifies
$\operatorname{Pic}^d(C)(k)$ with $J(k)$.
Every degree-two class other than the canonical class has one
effective divisor over $k$, while the canonical class has
$q+1$ of them: the dimensions of sections are respectively
one and two.  Summing the nontrivial character over all of
$J(k)$ gives zero, leaving the coefficient
$q\zeta^{i\chi([K_C-2P])}$.
For $d\ge3$, Riemann--Roch gives the same section dimension
$d-1$ in every class, so the number
$(q^{d-1}-1)/(q-1)$ of effective divisors in each class is
constant and the character sum is zero.
The degree-one coefficient is the stated sum over $C(k)$.
The rational point $P$ ensures these Picard classes are represented
by line bundles over $k$, so no descent obstruction enters this count.

For the twist indexed by $t$, the corresponding factor is
\[
              L_{i,t}(T)=
       1+\zeta^{it}S_iT+q\zeta^{i(a+2t)}T^2.
                                                            \tag{B.141.4}
\]
Indeed, on a closed point of degree $m$, the character value
is multiplied by $\zeta^{imt}$.
The regular-representation factorization of each closed-point
Euler factor gives
\[
 \frac{Z(D_t,T)}{Z(C,T)}
       =\prod_{i=1}^{\ell-1}L_{i,t}(T),\qquad
 \frac{\#J_{D_t}(k)}{\#J(k)}
       =\prod_{i=1}^{\ell-1}L_{i,t}(1).
                                                            \tag{B.141.5}
\]
The second product is an integer, since it is a rational algebraic
integer, and it equals the positive displayed ratio.
Thus the class-number criterion used here is precisely the
nonvanishing of every factor modulo $p$.

The genus-two canonical pencil is the hyperelliptic pencil,
with an involution $\iota$ defined over $k$ satisfying
$R+\iota R\sim K_C$.
For completeness, the pencil is base-point-free because a
degree-one divisor cannot have two independent sections on a
positive-genus curve; Riemann--Roch then gives
$h^0(K_C-R)=1<h^0(K_C)=2$.
It defines a degree-two map to the projective line.
That map is separable, including in characteristic two:
a purely inseparable degree-two extension of a rational
function field over the perfect field $k$ is rational.
Its nontrivial involution is therefore defined over $k$,
and permutes $C(k)$.
Applying this permutation to the character sum gives the exact identity
\[
                       S_{-i}=\zeta^{-ia}S_i.                \tag{B.141.6}
\]

Reduce $\mathbf Z[\zeta]$ into an algebraic closure of
$\mathbf F_p$.
The image $\bar\zeta$ still has order $\ell$.
Since $f$ is even and $(\mathbf Z/\ell)^\times$ is cyclic,
\[
                             p^{f/2}\equiv-1\pmod\ell.
                                                            \tag{B.141.7}
\]
The coefficients of every $S_i$ are rational integers.
Consequently, if the reduction of $L_{i,t}(1)$ is zero,
then raising that equality to the $p^{f/2}$-th power shows
that the reduction of $L_{-i,t}(1)$ is also zero.
Writing these two equalities out using $q=0$ in the residue field,
\[
        \overline S_i=-\bar\zeta^{-it},
                     \qquad
        \overline S_{-i}=-\bar\zeta^{it}.
\]
Together with (B.141.6), they imply
$\bar\zeta^{\,i(a+2t)}=1$.
Both $i$ and $2$ are invertible modulo $\ell$.
It follows that $a+2t=0$.
This proves (B.141.2) simultaneously for all primitive characters.
The $P$-split twist is $t=0$, so $a\ne0$ excludes every vanishing
factor in (B.141.5), proving (B.141.1).

The dual of the finite vector space $J(k)/\ell J(k)$ separates
its nonzero elements.  If $\delta$ is nonzero in that quotient,
there are $(\ell-1)\ell^{r-1}$ linear functionals nonzero on it.
Each is surjective.  The nonzero scalar multiples have the same
kernel and give the same pointed cover, and different kernels give
different geometric cyclic quotients.
Dividing by $\ell-1$ proves the count of covers.

Finally let $k_n/k$ have degree $n$ prime to $\ell$.
The character of the base-changed Lang cover is
\[
       \chi_n=\chi\circ\operatorname{Nm}_{k_n/k}:J(k_n)
                              \longrightarrow\mathbf Z/\ell.
                                                            \tag{B.141.8}
\]
This follows either from closed-point Artin symbols or from
$F_q^n-1=(F_q-1)(1+F_q+\cdots+F_q^{n-1})$ in the Lang construction.
On a class rational over $k$, the norm is multiplication by $n$.
Thus $\chi_n$ is nonzero, and for the extended canonical-point class,
\[
                            \chi_n(\delta)=n\chi(\delta)\ne0.
\]
The same proof over $k_n$, whose characteristic is still $p$,
gives the claimed persistence for this very cover.
In particular the point stays split by base change.
If instead $\delta\in\ell J(k)$, that divisibility persists under
every field extension.  Therefore this sufficient test cannot
first become applicable after enlarging the constants.

\medskip\noindent\textbf{The explicit whole family and its finite certificate.}
For the primes in (B.141.10), the polynomial $x^6+1$ is squarefree.
Its smooth proper double cover of the projective line has genus
two and two rational points $\infty_+,\infty_-$ above infinity.
The canonical class is their sum.
There is a degree-two morphism
\[
 \pi:C_p\longrightarrow E_p,\qquad
 E_p:\ v^2=u^3+1,\qquad (u,v)=(x^2,y).
                                                            \tag{B.141.12}
\]
Both points at infinity map to the origin $O$ of $E_p$.
The point $R=(0,1)=\pi(P)$ has exact order three:
the tangent line $v=1$ meets the cubic at $R$ with multiplicity
three, so its divisor gives $3R=O$, and $R\ne O$.
Consequently pushforward of degree-zero divisors gives
\[
                         \pi_*(\delta)=-2R=R.
                                                            \tag{B.141.13}
\]
Cubing permutes $\mathbf F_p$ because $p\equiv2\pmod3$.
The usual quadratic-character count, followed by that permutation,
therefore gives
\[
 \#E_p(\mathbf F_p)
   =p+1+\sum_{u\in\mathbf F_p}
                  \left(\frac{u^3+1}{p}\right)
   =p+1 .
\]
The assumptions modulo nine give $v_3(p+1)=1$.
In a finite abelian group whose order has three-adic valuation one,
a nonzero point of order three is not divisible by three.
Thus $R\notin3E_p(\mathbf F_p)$.
If $\delta$ belonged to $3J_{C_p}(\mathbf F_p)$, applying the norm
$\pi_*$ would contradict (B.141.13).
This proves $\delta\notin3J_{C_p}(\mathbf F_p)$ throughout the
family.  Since $\operatorname{ord}_3(p)=2$, the criterion already
proved applies to each member, and its persistence applies over
every extension of degree prime to three.
This family conclusion uses the displayed algebraic argument,
not a finite search.

At $p=5$, let $T=(4,0)\in E_5(\mathbf F_5)$.
The six-element elliptic group is generated by $R$ of order
three and $T$ of order two.
Choose $\chi_E(R)=1$, $\chi_E(T)=0$, and put
$\chi=\chi_E\circ\pi_*:J_{C_5}(\mathbf F_5)\to\mathbf Z/3$.
Equation (B.141.13) gives $\chi(\delta)=1$, so this character is
surjective.  Its values on the pointed curve classes are
\[
\begin{array}{c|rrrrrr}
 Q&(0,1)&(0,4)&(2,0)&(3,0)&\infty_+&\infty_-\\ \hline
 \chi([Q-P])&0&1&2&2&2&2 .
\end{array}                                                  \tag{B.141.14}
\]
For a primitive cube root $\zeta$, the two nontrivial character
factors are therefore
\[
  L_1(T)=1+3\zeta^2T+5\zeta T^2,\qquad
  L_2(T)=1+3\zeta T+5\zeta^2T^2.
\]
Their product is
\[
             1-3T+4T^2-15T^3+25T^4 ,
                                                            \tag{B.141.15}
\]
whose value at one is twelve.
Direct enumeration gives
$\#C_5(\mathbf F_5)=6$ and $\#C_5(\mathbf F_{25})=46$.
There is also a short independent check of the latter count.
The sixth powers in $\mathbf F_{25}^{\times}$ form
$\boldsymbol\mu_4=\mathbf F_5^\times$, with six preimages
per value.  The values of $1+\boldsymbol\mu_4$ are
$2,3,4,0$; every nonzero element of $\mathbf F_5$ is a
square in $\mathbf F_{25}$.
Thus the affine count is $2+6(2+2+2+1)=44$, and the two
points at infinity give forty-six.
The first two Newton identities and the reciprocal genus-two
Weil polynomial give
$\#J_{C_5}(\mathbf F_5)=(6^2+46)/2-5=36$.
Together with (B.141.5), this proves (B.141.11).

The complete finite certificate is
\path{finite_checks/code/b141_canonical_example.py}, with its
retained JSON and output.
It enumerates both curve point sets, uses
$\mathbf F_{25}=\mathbf F_5[u]/(u^2-2)$, retains and checks the
whole elliptic addition and character tables, verifies
$R\notin3E_5(\mathbf F_5)$ and all six pointed character values,
and computes (B.141.15) both by quadratic-integer multiplication
and by an integer resultant.
The cover itself is the explicit Lang quotient based at $P$
for the fully specified character $\chi$.
\end{proof}
\paragraph{Check record.} GPT-6 Astra Pro, 2026-11-18: authored full proof and exact program; root and configuration each read the complete criterion and separately the full family addition. Root reran the certificate; configuration independently derived the F25 point count. Effort not separately exposed.


\paragraph{Exact reduction of the remaining B1 locus.}
This proves B1 for every genus-two marked curve satisfying
$[K_C-2P]\notin\ell J_C(k)$ for at least one odd prime
$\ell\ne p$ with even $\operatorname{ord}_{\ell}(p)$.
It is a condition on the entire marked curve and its given field,
and includes every prime-to-$\ell$ constant extension of each
certified pointed cover.
A counterexample to B1 in genus two must therefore satisfy
\[
 [K_C-2P]\in\ell J_C(k)
 \quad\text{for every odd }\ell\ne p
          \text{ with even }\operatorname{ord}_{\ell}(p).
                                                            \tag{B.141.9}
\]
Only primes dividing $\#J_C(k)$ need be checked in this necessary
condition.  It leaves every Weierstrass marking, for which the
canonical-point class is zero, as well as other possible markings.
It proves neither the existence of a witness nor its impossibility
on that remaining locus, and does not settle the full threshold
$p>(2g-2)^2$.

The preceding B.28 bounds the number of bad twists by
$(\ell-1)/f$ in genus two; the bound here is one whenever $f$ is
even, and identifies the only potentially bad twist relative to
the prescribed point.  B.23 has canonical-point class zero at its
displayed marking and does not supply this sufficient family.
The explicit characteristic-two examples in B.21--B.27 and A.50
do not contain the family (B.141.10).
The proof above is a direct character calculation; no priority
or claim of a complete literature classification is made.

\paragraph{Check record.}
Written and checked by GPT-6 Astra Pro, finite-checks reader,
18 November 2026; the separate effort setting is not exposed.
The complete B.23 and B.28--B.29 texts were reread to check the
scope of the predecessor statements.
The adjacent B.21--B.22 statements were also reread when checking
that the explicit family is outside those finite examples.
The root and configuration readers independently passed the
complete criterion proof before the explicit-family addition;
the configuration reader also passed the entire addition and
the finite certificate, and supplied the independent
$\mathbf F_{25}$ count above.
These are recorded as separate check events.
This statement was prepared after checkpoint04 and was not
included in its frozen PDF.

\subsection{Normal algebraic extensions and the remaining subgroup boundary}

\begin{claim}{B.142}{Proved}{An abelian normal quotient supplies the arithmetic inputs for reconstruction}
Keep $D$, $M_0$, $\Omega$ and
$I=I_{M_0}\subset G_D$ as in Claim~B.139.
Let $H$ be a nontrivial closed normal subgroup of
$G_{M_0}$, and let $B$ be a closed normal subgroup
of $I$ such that
\[
 H\subsetneq B\subset I\cap\ker\chi_p,
                         \qquad [B,B]\subset H.
                                                        \tag{META.NG1}
\]
Put $K=\Omega^H$ and $A=\Omega^B$.
Then every finite extension of $K$ satisfies full
closed-point-preserving relative etale reconstruction
for finite-type schemes localized at universal
homeomorphisms, with the structural isomorphism
to the base retained.

The arithmetic assertion needed here is the following
one-sided description of the stable rational Kummer
kernel: for every finite $E/K$ and every $a\in E^\times$,
\[
 \kappa_E(a)=0\quad\Longrightarrow\quad
  g(a)/a\in\mu_\infty\ (g\in B),
 \quad a^m\in A\text{ for some }m\ge1.
                                                        \tag{META.NG2}
\]
No equality between this kernel and all radicals
of $A$ is asserted.
\end{claim}
\begin{proof}
The inertia group fixes all prime-to-$p$ roots
of unity, and $B$ also fixes all $p$-power roots
by its cyclotomic condition.  Hence
$\mu_\infty\subset A\subset K$.
All these fields are algebraic extensions of
the henselian field $D^h$, with the fixed
mixed-characteristic valuation.

Suppose that $a\in E^\times$ has zero rational
Kummer class.  Choose a finite Galois extension
$D_1/D$ inside $\Omega$ containing $a$ and
generators of $E/K$.  Let $M_1$ be its finite
local closure, and replace $E$ by $KD_1$.
Set
\[
 J=I\cap G_{D_1}=I_{M_1},\quad H_1=H\cap J,
 \quad\Gamma=J/H_1,\quad C=(B\cap J)/H_1.
\]
The local class of $a$ in $H^1(J,V)$, where
$V=\mathbb Q_p(1)$, vanishes on $H_1$ and
inflates from a unique $\theta\in H^1(\Gamma,V)$.
The group $C$ is abelian and acts trivially on $V$.
For any $g\in B$, Galoisness of $D_1/D$ makes
$g$ normalize $J$; normality of $H$ and $B$
makes the resulting action preserve $\Gamma$
and $C$.
It acts trivially on $C$, because $[B,B]\subset H$,
and on $\Gamma/C$, because $B$ is normal in $I$.
It acts trivially on $V$ by $\chi_p(g)=1$.

Inflation--restriction gives
\[
 0\longrightarrow H^1(\Gamma/C,V)
 \longrightarrow H^1(\Gamma,V)
 \longrightarrow H^1(C,V)^{\Gamma/C}.
                                                        \tag{META.NG3}
\]
The action of $g$ is the identity on both outer
terms.  Thus $(g-1)^2=0$ on the middle term.
Since $a\in D_1$ has a finite $g$-orbit, there
is $s\ge1$ with $g^s(a)=a$.
Naturality of Kummer theory and injectivity of
inflation give $g^s\theta=\theta$.
The unipotence identity now yields
\[
               0=(g^s-1)\theta=s(g-1)\theta,
\]
so $g\theta=\theta$ in characteristic zero.
This uses a finite orbit only for the specific
Kummer vector, not for all of $H^1(\Gamma,V)$.
We obtain $\kappa_J(g(a)/a)=0$.

The elements $a$ and $g(a)$ belong to the one
finite local field $M_1$, since $g$ preserves
$D_1$ and its valuation closure.
Restriction from $H^1(G_{M_1},V)$ to
$H^1(I_{M_1},V)$ is injective: the coefficient
module has no inertia invariants.
The rational Kummer kernel of $M_1$ is its
finite group of roots of unity.
Indeed valuation detects the uniformizer part,
and the principal-unit logarithm detects an
open subgroup of units modulo its finite
torsion; the residue unit group is finite.
Hence $g(a)/a\in\mu(M_1)$ for every $g\in B$.
One integer annihilating $\mu(M_1)$ makes
$a^m$ fixed by all of $B$, proving (META.NG2).
The argument included every finite $E/K$.

Choose $b\in K\setminus A$, which exists by
$H\subsetneq B$.
At least one of $b$ and $1-b$ has nonzero
Kummer class over $K$.
Otherwise (META.NG2) would give, for every $g\in B$,
\[
        g(b)=\zeta_g b,\qquad
        g(1-b)=\eta_g(1-b),
                 \quad\zeta_g,\eta_g\in\mu_\infty\subset A.
\]
If some pair is unequal, its sum expresses
$b=(1-\eta_g)/(\zeta_g-\eta_g)\in A$.
If every pair is equal, the sum makes each
pair one and all $g$ fix $b$, again a contradiction.
This nonzero class remains nonzero over every
finite extension by rational restriction and
corestriction.  In particular all the fields
used below are non-separably-closed.

For a degree $d\ge1$ use
$\mathcal G_0=\{u+iv:0\le u,v\le4d\}$ and
$\mathcal G_1=b+\mathcal G_0$.
Suppose a quotient $R=P/Q$ of monic degree-$d$
polynomials over $\Omega$ has nonzero grid
values and all those values are zero Kummer
classes after finite extensions of $K$.
For every $g\in B$, (META.NG2) gives
$g(R(a))/R(a)\in\mu_\infty$.
Since $g$ fixes the first grid, write
$gR/R=U/V$ with monic polynomials of equal
degree at most $2d$ before cancellation.
Under an embedding $\Omega\hookrightarrow\mathbb C$,
the polynomial
\[
 U(u+iv)\overline U(u-iv)-V(u+iv)\overline V(u-iv)
\]
vanishes on the whole grid and has degree at
most $4d$ separately in $u,v$.
Successive interpolation makes it zero.
The invertible change $z=u+iv,w=u-iv$ gives
$U(z)\overline U(w)=V(z)\overline V(w)$.
Fixing an allowed $w$ shows proportionality,
and monicity gives $gR=R$.
Repeating this on $b+\mathcal G_0$ gives
$gR(gb+Z)=R(b+Z)$.
Choose $g$ moving $b$.
Then $R$ has the nonzero additive period $gb-b$.
A finite pole would have infinitely many
translates; a nonconstant polynomial has a
nonzero leading finite difference.
Thus $R$ is constant, and its value at infinity
is one.  We have proved the uniform two-grid test.

We now verify the remaining inputs of the
norm proof, rather than treating (META.NG2)
alone as reconstruction.
For a smooth geometrically integral affine
source and a pinned map to $\mathbb A^1$, fix
any finite Noether projection and the two
grids just constructed.
As in (META.RL4), choose a principal open
avoiding the branch locus and the proper
closed images of the grid fibres.
The geometric unit calculation, whose Picard
group is zero, gives norm classes
\[
                 \xi_a=\beta_a+
                       \sum_j\lambda_{aj}\kappa(h_j),
                  \qquad\lambda_{aj}\in\mathbb Q_p.
\]
Include all scheme and boundary coefficients,
grid centres, one algebraic normalizing parameter
and its actual norm values in a finite Galois
$D_1/D$.  Evaluation at that parameter identifies
each $\beta_a$ with an explicit finite linear
combination of Kummer classes from $D_1$.
With $E_1=KD_1$ and $M_1=\overline{D_1}$
in the local topology, the group
$N_1=G_{E_1}=H\cap G_{D_1}$ is normal in
$G_{M_1}$ and fixes $\mathbb Q_p(1)$.
All normalized classes are $G_{M_1}$-invariant.
Comparison of their evaluations at $t$ and
$\sigma t$ puts every grid value of
$P_{\sigma t}/\sigma P_t$ in the stable kernel.
The two-grid test gives
$P_{\sigma t}=\sigma P_t$ uniformly in
$\sigma\in G_{M_1}$ and the parameter.

The nonzero Kummer class constructed above
restricts nontrivially to $N_1$.
Dividing its integral image $p^r\mathbb Z_p$
by taking an actual $p^r$-th root in $E_1$
gives a Kummer homomorphism onto $\mathbb Z_p$.
After adjoining its algebraic argument and
$\mu_{p^s}$ to the finite local field, it
produces a cyclic $L/M_1$ of any prescribed
degree $q=p^s$, with surjective image of $N_1$.
These enlargements lie in the local fixed
field of $N_1$, so normality persists.
Claim~B.139 puts $M_1\cap\Omega$ in
the henselian $E_1$, identifies
$L\cap\Omega$ as a finite extension of it,
and supplies an $M_1$-basis of $L$ in that
algebraic part.

If a class over $L$ vanishes on
$N_1\cap G_L$, it inflates from the quotient.
The group $N_1$ acts trivially on that quotient
and on the coefficient module.
Its local logarithm is therefore in $M_1$,
because $N_1$ surjects onto $\operatorname{Gal}(L/M_1)$.
Actual point maps are valuation-continuous
over all finite extensions: the qcs etale-site
and relative-point comparison of CHW gives
etale-open continuity, and JTWY Corollary~6.13
applies to these non-separably-closed henselian
fields \citep{CarlsonHaineWolf2024,AN_JTWY2024}.

All inputs of the norm proof in
Claim~B.136 have now been verified over
$\Omega$, with the local comparison of
Claim~B.139.
To specify the use, take $q>d+1$.
Continuity and minimum conjugate gcd degree
give a constant-degree neighbourhood on the
dense algebraic parameter set.
Exact local logarithms and $2e$-centre
interpolation recover the reduced ratios.
Their extended coordinate derivations and
cocycle give the root divisor
$\tau^jB_0-B_0+C_j$.
Constant roots are algebraic over $D$ by
one actual specialization and are excluded
by proper closed actual fibres.
A negative coefficient of $B_0$ would give
more than $d$ distinct denominator roots;
specialization off the nonzero analytic
subresultants contradicts the actual degree.
The resulting effective divisor gives the
analytic monic denominator $Q$ in one block.
The exact logarithmic equations satisfy
B.135, so $Q\in L(t)[Z]$; interpolation of
its values in $\Omega$ gives $Q\in\Omega(t)[Z]$.
Residues are integral, and one actual
normalization gives an arithmetic constant
factor $C$ with $P_t=Q_tC$ by the proved
two-grid test.
Proper closed fibres force $C=1$.
This recovers all finite-projection norms.
The fully proved graph-degree argument in
Claim~B.136 makes their graph a single
rational function.  Valuation continuity
on algebraic transverse curves removes
poles and exceptional point values.
It is a regular function over the original
base by Galois descent.
Geometric components are handled by one
finite extension and the same uniqueness.
Finally CHW Theorem~4.18 and Corollary~2.32
give full reconstruction in the stated
relative category \citep{CarlsonHaineWolf2024}.
Every step starts with any finite $E/K$,
as required.
\end{proof}
\paragraph{Check record.} checked only in this writing, 2026-11-18: GPT-6 Astra (OpenAI), ChatGPT Pro mode; internal effort value not exposed. Own detailed proof written and reread; primary source hypotheses matched. Root full 425-line reading of B142--143 passed: B normal in inertia suffices for the quotient action, H normal in the local group gives all norm inputs, one local roots-of-unity group bounds all g, both grids and nonzero class, full norm transfer, no radical converse, nonabelian normal-subgroup lemma, minimal derived index and H containing B-prime nontrivial, correct core-in-open and algebraic descent


\begin{claim}{B.143}{Proved}{Every proper Galois local algebraic extension reconstructs}
Keep the data of Claim~B.139.
Let $1\ne N\triangleleft G_{M_0}$ be any
closed normal subgroup.
Then $F=\Omega^N$, every finite extension of
$F$ in $\Omega$, and every algebraic subfield
of any of those fields satisfy the reconstruction
property of Claim~B.142.
Equivalently, every proper Galois algebraic
extension of $D^h=M_0\cap\Omega$ reconstructs,
with the same finite-extension and algebraic
descent assertions.

More generally let $D^h\subset F\subset\Omega$.
If for some finite local extension $M_1/M_0$
the closed subgroup
\[
 \operatorname{Core}_{G_{M_1}}(G_F\cap G_{M_1})
   =\bigcap_{g\in G_{M_1}}g(G_F\cap G_{M_1})g^{-1}
                                                        \tag{META.NG4}
\]
is nontrivial, then $F$ and every finite
extension of $F$ reconstruct.
Thus within these algebraic local subfields
the remaining cases have trivial normal
core in every open local absolute Galois
group.
\end{claim}
\begin{proof}
We first show that every nontrivial closed
normal subgroup of a finite-local-field
absolute Galois group is nonabelian.
Let $P=P_{M_0}$ be its wild inertia and
suppose $N$ is abelian.
Then $A=N\cap P$ is abelian and normal in $P$.
If $a\in A$ is nonidentity, choose two
noncommuting elements $s,t\in P$ and let
$P_0=\overline{\langle a,s,t\rangle}$.
Wild inertia is free pro-$p$ of infinite
rank, and its finitely generated closed
subgroup $P_0$ is free pro-$p$.
The nontrivial abelian normal subgroup
$A\cap P_0$ is procyclic, because it lies
in the procyclic centralizer of $a$ in
$P_0$.
The normalizer of a nonidentity procyclic
subgroup in a finitely generated free
pro-$p$ group is procyclic as well.
Thus normality would make $P_0$ procyclic,
contrary to the choice of $s,t$.
These exact free pro-$p$ statements are
Iwasawa Theorem~2(ii) and Kochloukova--Zalesskii
Theorem~5.1, Case~I, opened and used with
the same finite-generator reduction in
Claim~B.119
\citep{AN_Iwasawa1955,AN_KZ2011}.
We obtain $A=1$.
Normality of both $N$ and $P$ gives
$[N,P]\subset N\cap P=1$.
The local centralizer theorem
$C_{G_{M_0}}(P)=1$, proved at every prime
in Claim~B.119, now gives $N=1$.
This contradiction proves the assertion.
The argument applies to every finite
extension of $M_0$ as well.

Return to a nontrivial normal $N$ and put
$B=N'=\overline{[N,N]}$.
It is nontrivial by what we just proved,
and it is normal in $G_{M_0}$.
The quotient $G_{M_0}/I$ and the cyclotomic
image are abelian, so
\[
                         B\subset I\cap\ker\chi_p.
\]
We do not assert that $B$ is pro-$p$.
Every finite quotient of $I$ is an extension
of a finite $p$-group by a cyclic prime-to-$p$
group, hence solvable.
It follows that $\bigcap_r I^{(r)}=1$:
a nonidentity element has a nonidentity
image in some finite solvable quotient,
whose derived series terminates.
There is consequently a least integer
$r\ge1$ with $B\not\subset I^{(r)}$.
Minimality gives $B\subset I^{(r-1)}$.
Set
\[
                              H=B\cap I^{(r)}.
\]
This is a closed normal subgroup of $G_{M_0}$,
proper in $B$.
The commutator inclusion
$[B,B]\subset [I^{(r-1)},I^{(r-1)}]\subset I^{(r)}$
makes $B/H$ abelian and nontrivial.
Moreover $H\ne1$:
the nontrivial normal subgroup $B$ is
nonabelian by the first paragraph, so
$1\ne B'\subset H$.
Every hypothesis of Claim~B.142 has
therefore been proved.
It reconstructs $K=\Omega^H$ and all its
finite extensions.
Since $H\subset N$, the field $K$ contains
$F=\Omega^N$, and $K/F$ is algebraic.
The explicit algebraic-descent argument
in Claim~B.140 now proves reconstruction
over $F$.
For any finite $F'/F$, the compositum
$KF'/K$ is finite; reconstruction there
descends to $F'$.
The same argument gives the stated
algebraic subfield assertion.

For (META.NG4), write its nontrivial core
as $C$.  It is a closed normal subgroup
of $G_{M_1}$.
By Claim~B.139 one may choose a finite
extension $D_1/D$ inside $\Omega$ with
closure $M_1$; its henselization is
$M_1\cap\Omega$.
Apply the first assertion with this base
and $C$.
The reconstructed proper Galois field
$\Omega^C$ contains $F$, since $C\subset G_F$.
Its extension of $F$ is algebraic.
Descent, and descent after every finite
compositum, prove the assertion.

Finally every henselian valued subfield
$F\subset\overline{\mathbb Q}_p$ of
finite transcendence degree contains the
henselization of a finitely generated
$D\subset F$ containing a transcendence
basis.  Then $F/D^h$ is algebraic and
the criterion just proved applies whenever
its indicated core is nontrivial.
No assertion for core-free subgroups is
hidden in this last reduction.
\end{proof}
\paragraph{Check record.} checked only in this writing, 2026-11-18: GPT-6 Astra (OpenAI), ChatGPT Pro mode; internal effort value not exposed. Own detailed proof written and reread; primary source hypotheses matched. Root full 425-line reading of B142--143 passed: B normal in inertia suffices for the quotient action, H normal in the local group gives all norm inputs, one local roots-of-unity group bounds all g, both grids and nonzero class, full norm transfer, no radical converse, nonabelian normal-subgroup lemma, minimal derived index and H containing B-prime nontrivial, correct core-in-open and algebraic descent


\paragraph{Remaining scope.}
The normality argument is a new sufficient
criterion within the algebraic closure of
a finitely generated subfield of
$\overline{\mathbb Q}_p$.
It does not replace a closed subgroup by
its normal closure or its normal core
without checking that the core is nontrivial.
Nor does it embed arbitrary henselian
mixed-characteristic fields into
$\overline{\mathbb Q}_p$.
In particular fields with more general
residue fields or value groups, and the
remaining center-free core-free subgroup
cases, require separate arguments.
No priority or external human verification
is claimed.

\subsection{Explicit core-free local fields with a structural-cell obstruction}

\begin{claim}{B.144}{Proved}{Procyclic wild fixed fields fail reconstruction in every finite transcendence degree}
Keep the relative local data $D,M_0,\Omega$ of
Claim~B.139, and put $P=P_{M_0}$.
Choose $s\ge1$, with $s\ge2$ if $p=2$, such that
$1+p^s\in\chi_p(P)$, and choose
$\tau\in P$ with $\chi_p(\tau)=1+p^s$.
Then
\[
 C=\overline{\langle\tau\rangle}\simeq\mathbb Z_p,
 \qquad F_\tau=\Omega^C                                      \tag{META.PC1}
\]
is a henselian mixed-characteristic field of
transcendence degree $\operatorname{trdeg}_{\mathbb Q}D$,
with residue field $\overline{\mathbb F}_p$.
Neither $F_\tau$ nor any of its finite extensions
satisfies closed-point-preserving relative etale
reconstruction with the structural isomorphism retained.
For every finite $M_1/M_0$ in the fixed local algebraic closure,
\[
       \operatorname{Core}_{G_{M_1}}(C\cap G_{M_1})=1.
                                                               \tag{META.PC2}
\]
For $M_0=\mathbb Q_p$ one can take
$\chi_p(\tau)=1+p$ if $p$ is odd and
$\chi_2(\tau)=5$ if $p=2$.
Consequently this supplies negative mixed-characteristic
families in every finite transcendence degree, outside
the positive normal-core criterion of Claim~B.143.
\end{claim}
\begin{proof}
We first justify the choices at every prime.
The $p$-power cyclotomic extension of $\mathbb Q_p$
is totally ramified and has infinite pro-$p$ part.
A finite local extension $M_0/\mathbb Q_p$ intersects
it in a finite extension, so the cyclotomic image of
$G_{M_0}$ is open in the cyclotomic image of
$G_{\mathbb Q_p}$.
For the inertia assertion specifically,
$I_{M_0}=I_{\mathbb Q_p}\cap G_{M_0}$ is an open
subgroup of $I_{\mathbb Q_p}$; its cyclotomic image
has finite index in the image of $I_{\mathbb Q_p}$.
It therefore contains an open subgroup of $1+p\mathbb Z_p$
(of $1+4\mathbb Z_2$ when $p=2$).
The quotient $I_{M_0}/P$ is pro-prime-to-$p$.
Hence the pro-$p$ part of that image comes from $P$:
its image in $\chi_p(I_{M_0})/\chi_p(P)$ is both pro-$p$
and pro-prime-to-$p$, and is therefore trivial.
This proves the existence of $s$ and $\tau$.
For $M_0=\mathbb Q_p$ the cyclotomic description gives
the two stated explicit image values.

The continuous homomorphism $\mathbb Z_p\to C$,
$z\mapsto\tau^z$, is surjective by construction.
Its kernel is a closed subgroup of $\mathbb Z_p$.
Every nonzero such subgroup is open, which would make
$C$ finite.
But $\chi_p(\tau)=1+p^s$ has infinite order, as is
seen by its nonzero $p$-adic logarithm.
The kernel is therefore zero and $C\simeq\mathbb Z_p$.
This avoids any possible torsion qualification at $p=2$.

Since $C\subset P\subset I_{M_0}$, the field $F_\tau$
contains the strict henselization of $D$ for the chosen
valuation.  It is algebraic over $D^h$, hence henselian.
The residue field of the strict henselization is
$\overline{\mathbb F}_p$; every residue element in the
local algebraic closure is already algebraic over
$\mathbb F_p$, so the residue field remains exactly
$\overline{\mathbb F}_p$.
The extension $F_\tau/D$ is algebraic and therefore
has the stated transcendence degree.
It is not separably closed, since its absolute
Galois group is the nontrivial group $C$.

We give the retained-structural-cell obstruction
in full, including its actual point function.
This is the necessary center condition already
appearing in the proof of B.50, now applied to an
explicit local core-free field family.
Let $F$ be any characteristic-zero field with algebraic
closure $\Omega$ and a nonidentity
$c\in Z(G_F)$.
Write $X=\mathbb A^1_F$ and
$p_*:X_{\mathrm{et}}\to BG_F$ for its structural
geometric morphism.
For a continuous $G_F$-set $S$, the action of $c^{-1}$
is a $G_F$-equivariant natural automorphism of $S$;
centrality is exactly what ensures equivariance.
Applying $p^*$ gives a natural automorphism of $p^*$.
Let $\alpha_c:p_*\xrightarrow{\sim}p_*$ be its
adjoint mate.
Thus, using the right-adjoint convention of
CHW Notation~2.2, the pair
\[
                  (\operatorname{id}_{X_{\mathrm{et}}},\alpha_c)
                                                               \tag{META.PC3}
\]
is a well-typed relative geometric automorphism:
its retained structural cell has source
$p_*\operatorname{id}$ and target $p_*$
\citep[Notation~2.2]{CarlsonHaineWolf2024}.
Its underlying topological map is the identity,
so it preserves closed points and is pinned.

For an actual point $a\in\Omega=X(\Omega)$,
the relative point obtained by composing with
(META.PC3) is the point $c(a)$.
Here is an explicit check of the direction of the cell.
For any etale $V\to X$, applying $c^{-1}$ to the
coordinates gives a compatible bijection
\[
                    V_{c(a)}(\Omega)\longrightarrow V_a(\Omega).
\]
These bijections identify the two geometric fibre
functors.  On an etale object pulled back from $F$
the bijection is the action of $c^{-1}$ on its
$G_F$-set of geometric points.
It is consequently the prescribed mate of $\alpha_c$
on the base, and gives the required relative
identification, including the structural cell.
The actual point function of (META.PC3) is therefore
\[
                              a\longmapsto c(a).                \tag{META.PC4}
\]

This relative morphism cannot become the etale
topos morphism of an $F$-scheme morphism
$h:\mathbb A^1_F\to\mathbb A^1_F$, even up to a
relative $2$-isomorphism.
Indeed composing such an isomorphism with every
relative $\Omega$-point would give the same actual
point by CHW Proposition~3.6: its hypotheses hold
because $F$ is perfect, $\Omega/F$ is algebraic,
and $X$ is finite type \citep[Proposition~3.6]{CarlsonHaineWolf2024}.
We would have $h(a)=c(a)$ for all $a\in\Omega$.
The morphism $h$ is a polynomial in $F[T]$.
Since $c$ fixes the infinite subset $\mathbb Q\subset F$,
that polynomial agrees with $T$ at infinitely many
arguments and is therefore $T$.
But $c$ moves some element of $\Omega$, a contradiction.
Nor can localization at universal homeomorphisms
remove the contradiction: if that localized
reconstruction property held, the characteristic-zero
comparison in CHW Corollary~2.32 would realize
this relative morphism by an ordinary polynomial
on the smooth seminormal source $X$
\citep[Corollary~2.32]{CarlsonHaineWolf2024}.
This proves failure in the exact relative category,
including its relative $2$-isomorphisms.

The necessary group condition is in fact slimness:
every open subgroup of $G_F$ must have trivial center.
We justify the finite-extension transfer used here.
Suppose reconstruction holds over $F$, and let
$E/F$ be finite.
An $E$-scheme is an $F$-scheme together with its
map to the finite etale object $\operatorname{Spec}E$.
On etale topoi this is the slice construction for
the corresponding finite locally constant object
of $BG_F$; a map into that slice is equivalently
a map over $BG_F$ preserving the chosen section
of the pulled-back finite object.
Given a pinned relative morphism over $E$, forget
to $F$ and reconstruct it there, using seminormal
finite-type sources as permitted by Corollary~2.32.
The resulting $F$-scheme map preserves the chosen
finite-object section: the original relative
$E$-cell supplies this equality, and the
reconstructing relative $F$-isomorphism preserves
it upon pullback.
Thus the scheme map is an $E$-scheme map.
The universal property of the slice identifies
its relative $E$-morphism with the given one;
there is no further choice of a structural cell
once that section is specified.
Faithfulness descends in the same way, or follows
by forgetting to $F$.
This proves reconstruction over every finite $E/F$.
If an open subgroup $G_E\subset G_F$ had a
nonidentity central element, (META.PC3)--(META.PC4)
over $E$ would contradict that reconstruction.

Apply this argument to $F=F_\tau$ and $c=\tau$.
Every finite extension $F'/F_\tau$ has absolute
Galois group an open subgroup of $C$, hence
$p^m\mathbb Z_p$ for some $m\ge0$.
It again has a nonidentity central element,
so the same construction proves failure over $F'$.

Finally $C\cap G_{M_1}$ is an open subgroup of
$C$ for every finite $M_1/M_0$.
Its core in $G_{M_1}$ is abelian and normal.
The group-theoretic argument in Claim~B.143
proves that every nontrivial closed normal subgroup
of a finite-local-field absolute Galois group is
nonabelian.  The core must consequently be trivial,
which proves (META.PC2).
The same conclusion holds for the absolute Galois
subgroup of every finite $F'/F_\tau$.

For any $n\ge0$, choose $n$ algebraically independent
elements of $\mathbb Q_p$ over $\mathbb Q$ and let
$D$ be the field they generate over $\mathbb Q$.
Its closure is $\mathbb Q_p$, so the construction
applies with the explicit cyclotomic values stated
above and with transcendence degree $n$.
\end{proof}
\paragraph{Check record.} checked only in this writing, 2026-11-18: GPT-6 Astra (OpenAI), ChatGPT Pro mode; internal effort value not exposed. Own detailed proof written and reread; primary source hypotheses matched. Root full original theorem reading and full rereading of final joins passed: mate the inverse of c, actual a-to-c(a), relative 2-isomorphism exclusion, seminormal UH contraposition, p2 and finite-extension/core assertions, explicit open-inertia cyclotomic argument and finite-etale slice transfer proving necessary slimness; status Proved


\paragraph{Literature boundary.}
Asayama's Theorem~1.4 and Corollary~1.5 produce
Kummer-faithful algebraic extensions of number fields
with abelian absolute Galois group.
His Remark~1.6 records the gap in the formerly
asserted general implication from Kummer-faithfulness
to slimness \citep{AN_Asayama2026}.
This distinguishes Kummer-faithfulness from the
center obstruction proved above.
No henselian property of Asayama's examples is
used or asserted here, and no Kummer-faithfulness
claim is made for the explicit fields (META.PC1).
The remaining local subgroup gate is now restricted
to slim groups whose cores in every open
local absolute Galois group are trivial.
No priority or external human verification is claimed.

\subsection{A prescribed-point criterion from an automorphism orbit}

\begin{claim}{B.145}{Proved}{Long automorphism orbits force pointed class-number witnesses}
Let $C/k$ be smooth, proper and geometrically connected of genus
$g\ge2$ over $k=\mathbf F_q$ of characteristic $p$.
Let $P\in C(k)$, let $\ell\ne p$ be an odd prime, and let
$\chi:J_C(k)\twoheadrightarrow\mathbf Z/\ell$.
Suppose there is $\tau\in\operatorname{Aut}_k(C)$ and
$u\in(\mathbf Z/\ell)^\times$ such that
\[
 \chi\circ\tau_*=u\chi,\qquad
 u\in\langle p\rangle,\qquad
 m:=\operatorname{ord}_{\ell}(u)>g-1 .
                                                            \tag{B.145.1}
\]
Put $a=\chi([\tau P-P])$.
Index the constant twists of the associated pointed Lang cover by
$t\in\mathbf Z/\ell$, so that a closed point of degree $d$
acquires the added Frobenius translation $dt$.
The only possibly bad twist, meaning one whose relative class
number is divisible by $p$, is
\[
                              t=\frac{a}{u-1}.
                                                            \tag{B.145.2}
\]
In particular, if $a\ne0$, the specified $P$-split cover, at
$t=0$, is a class-number witness already over $k$.
That same cover remains a witness over every constant extension
of degree prime to $\ell$.
There is no restriction that the order of $\tau$ be prime to $p$.

An explicit genus-three plane-quartic example is
\[
 C:\ Y^3Z-YZ^3=X^4,\qquad
 k=\mathbf F_{27}=\mathbf F_3[t]/(t^3-t-1),\qquad
 P=(x,y)=(1,t),\qquad \ell=7 .
                                                            \tag{B.145.3}
\]
The automorphism is $\tau(x,y)=(x,y+1)$.
The character specified below has $u=2=3^2\bmod7$ and $a=1$.
Its constant cyclic cover has genus fifteen and
\[
\begin{split}
 \#J_C(k)&=18928,\\
 \frac{\#J_D(k)}{\#J_C(k)}&=427439613252172921
                                  \equiv1\pmod3,\\
 \#J_D(k)&=8090576999637129048688 .
\end{split}                                                   \tag{B.145.4}
\]
At this particular field and marking the least admissible
witness prime is seven.
All twenty-four points of $C(k)\setminus C(\mathbf F_3)$
are certified by this character construction.
The four original $\mathbf F_3$ pointings are not certified
by the nonzero-shift hypothesis, but the separate finite twist
calculation also gives them a degree-seven witness over $k$,
with relative class number $121740744925904896$.
For each original pointing the first field admitting an
odd-prime class-number witness has degree three over
$\mathbf F_3$, and seven is the least prime at that field.
For the displayed seven twists over $k$, base change to
$\mathbf F_{27^n}$ has relative class number prime to three
exactly when $4\nmid n$ for twist one, and exactly when
$28\nmid n$ for every other twist.
At the stated field, the split pointed models are twist one
for the four original markings and twists zero, four and six
for the twenty-four new markings.  Twists two, three and five
have no split $k$-rational marking; their assertion is only
the relative class-number unit test.
\end{claim}
\begin{proof}
\medskip\noindent\textbf{The general character polynomial and its degree modulo $p$.}
We use the connected constant Lang cover attached to $\chi$,
with its chosen rational point over $P$, as constructed in B.2.
Fix a primitive $\ell$-th root of unity $\zeta$.
For $i\ne0\bmod\ell$ let
\[
 A_d(i)=\sum_{\substack{D\ge0\\ \deg D=d}}
              \zeta^{\,i\chi([D-dP])},\qquad
                 L_i(T)=\sum_{d\ge0}A_d(i)T^d .
                                                            \tag{B.145.5}
\]
The sum is over effective divisors defined over $k$.
Its Euler product is the unramified character factor of the cover.
We recall the exact Riemann--Roch reduction needed here.

Identify $\operatorname{Pic}^d(C)(k)$ with $J_C(k)$ using $dP$.
The number of effective divisors in a line-bundle class $M$
is $(q^{h^0(M)}-1)/(q-1)$.
For $d>2g-2$, that number is the same in every class, so
nontrivial character orthogonality gives $A_d(i)=0$.
For $g\le d\le2g-2$, Riemann--Roch and the same orthogonality
give the exact identity
\[
 A_d(i)=q^{d+1-g}\zeta^{\,i c}
                      A_{2g-2-d}(-i),\qquad
             c=\chi([K_C-(2g-2)P]).
                                                            \tag{B.145.6}
\]
Indeed replace $M$ by $K_C\otimes M^{-1}$.
Its section dimension differs from $h^0(M)$ by $g-1-d$.
The contribution of the constant term in the numerator
$q^{h^0(M)}-1$ vanishes after summing the character.
Factoring out $q^{d+1-g}$ leaves the effective-divisor count
for $K_C\otimes M^{-1}$, with character multiplier
$\zeta^{ic}$.
The rational point ensures these line-bundle classes and their
linear systems are defined over $k$.
Consequently
\[
                   \deg(\overline L_i)\le g-1,\qquad
                            \overline L_i(0)=1
                                                            \tag{B.145.7}
\]
after reduction at any prime above $p$.
The polynomial is therefore not identically zero.

\medskip\noindent\textbf{The affine action on twists.}
The automorphism $\tau$ permutes the effective divisors of
each degree.  For any such divisor,
\[
 [\tau D-dP]=\tau_*[D-dP]+d[\tau P-P].
\]
Applying $\chi$ and summing gives
\[
                  A_d(iu)=\zeta^{-ida}A_d(i),\qquad
                         L_{iu}(T)=L_i(\zeta^{-ia}T).
                                                            \tag{B.145.8}
\]
Choose $j\ge1$ with $p^j=u\bmod\ell$.
If the $i$-th character factor of twist $t$ vanishes modulo $p$,
then $\overline L_i(\bar\zeta^{it})=0$.
Raising this equality to the $p^j$-th power and using
(B.145.8) gives
\[
              \overline L_i(\bar\zeta^{\,i(ut-a)})=0.
                                                            \tag{B.145.9}
\]
Here coefficients are sums of rational-integer multiples of
roots of unity, so the same power changes the character index
from $i$ to $iu$.
The image $\bar\zeta$ has exact order $\ell$, because
$\ell\ne p$.

The affine map $t\mapsto ut-a$ has the unique fixed point
$a/(u-1)$, and every other orbit has length $m$.
Indeed, after translating by that fixed point, it is
multiplication by $u$ on $\mathbf F_\ell$.
The assumption $m>g-1\ge1$ ensures $u\ne1$.
For any nonfixed $t$, iteration of (B.145.9) would give
$m$ distinct roots of the nonzero polynomial
$\overline L_i$, contradicting (B.145.7).
This excludes a vanishing factor for every primitive index $i$,
with the same unique exceptional twist for all of them.

As in the direct Euler-factor calculation of B.141, the product
of the nontrivial character factors at one is exactly
$\#J_{D_t}(k)/\#J_C(k)$.
It is an integer, and it is divisible by $p$ precisely when one
of those reduced factors vanishes.
This proves (B.145.2), including the assertion for the
prescribed $P$-split twist.

For a degree-$n$ constant extension with $\ell\nmid n$,
the same cover has character
$\chi_n=\chi\circ\operatorname{Nm}_{k_n/k}$.
Norm commutes with $\tau_*$, so the eigencharacter relation
still holds with $u$.
Moreover
$\chi_n([\tau P-P])=n\chi([\tau P-P])=na\ne0$.
The character is nonzero since its restriction to $J_C(k)$
is $n\chi$.
The already proved argument over $k_n$ proves persistence.

\medskip\noindent\textbf{A nonempty genus-three example.}
The three partial derivatives of
$Y^3Z-YZ^3-X^4$ in characteristic three are
$-X^3,-Z^3,Y^3$; they have no common projective zero.
Thus (B.145.3) is a smooth plane quartic of genus three.
Its unique point at infinity is $O=(0:1:0)$.
The line $Z=0$ cuts out $4O$, so $K_C\sim4O$.
There is a degree-two morphism
\[
 \pi:C\longrightarrow E,\qquad
 E:\ v^2=u^3-u,\qquad (u,v)=(y,x^2),
                                                            \tag{B.145.10}
\]
which sends $O$ to the elliptic origin and intertwines $\tau$
with $\tau_E(u,v)=(u+1,v)$.

The point $R=(t,1)$ belongs to $E(k)$.
The doubling law gives
$2R=(t+1,1)$, $4R=(t+2,1)$ and $8R=R$.
Hence $R$ has exact order seven.
Over $\mathbf F_3$, the elliptic curve has four points and
Frobenius trace zero.  Its two Frobenius roots have sum zero,
so their cubes also have sum zero and $\#E(k)=28$.
All four points of $E[2]$ are already rational over
$\mathbf F_3$.
It follows that
\[
                    E(k)=E[2]\oplus\langle R\rangle .
\]
Define $\chi_E:E(k)\to\mathbf Z/7$ by
$\chi_E(R)=1$ and $\chi_E(E[2])=0$.
Since $\tau_E(R)=2R$ and $\tau_E$ permutes $E[2]$,
$\chi_E\circ\tau_E=2\chi_E$.
Set $\chi=\chi_E\circ\pi_*:J_C(k)\to\mathbf Z/7$.
For $P=(1,t)$, one has $\pi(P)=R$ and
\[
                 \chi([\tau P-P])=\chi_E(2R-R)=1.
\]
Thus $\chi$ is surjective, its eigenvalue is $u=2=3^2\bmod7$,
and $m=3>2=g-1$.
The cover is a witness by the general theorem.
This verifies the criterion's hypotheses geometrically and
arithmetically, before any finite class-number computation.

\medskip\noindent\textbf{The complete finite arithmetic certificate.}
The program \path{finite_checks/code/b145_quartic_example.py}
uses the displayed cubic field and
$\mathbf F_{729}=k[s]/(s^2+1)$.
Its elements are encoded by $a+27b$ with $a,b\in k$;
the $27$-power Frobenius sends $s$ to $-s$.
It retains every degree-one point and every degree-two
closed-point orbit used below, together with their elliptic
norms and character values.
The elliptic addition and character tables are also retained
and checked on every pair.

The exact curve counts over $\mathbf F_3,\mathbf F_9,k$ and
$\mathbf F_{729}$ are respectively $4,4,28,676$.
For the first three, Newton's identities give the
$\mathbf F_3$ Frobenius polynomial
\[
                         (X^2-3)^2(X^2+3).
\]
Cubing its roots gives $(X^2-27)^2(X^2+27)$ over $k$,
and hence $\#J_C(k)=(-26)^2\cdot28=18928$.
For the character based at $P$, the degree-one and
degree-two effective-divisor occupancy vectors are
\[
\begin{split}
 (n_0,\ldots,n_6)&=(8,8,0,8,0,0,4),\\
 (m_0,\ldots,m_6)&=(116,116,116,116,116,34,116).
\end{split}                                                   \tag{B.145.11}
\]
The second vector includes the $28\cdot29/2=406$ unordered
pairs of rational points, with repetition, and all
$(676-28)/2=324$ closed points of degree two.
The character of a degree-two closed point is computed by
the elliptic norm of its two geometric points minus $2R$.
Thus these are the complete coefficient sums in (B.145.5),
not point-count proxies for the degree-two coefficient.

One has $c=\chi([4O-4P])=-4=3\bmod7$.
Writing
\[
       S=8+8\zeta+8\zeta^3+4\zeta^6,
\]
the resulting full character factor is
\[
 L_1(T)=1+ST-82\zeta^5T^2+
             27\zeta^3 S(\zeta^{-1})T^3+729\zeta^3T^4 .
                                                            \tag{B.145.12}
\]
Here the degree-two vector reduces to $-82\zeta^5$
using $1+\zeta+\cdots+\zeta^6=0$, and the two high-degree
coefficients follow from (B.145.6).
At one, it is represented by
\[
 A(Z)=225+8Z+216Z^2+953Z^3+108Z^4-82Z^5+4Z^6 .
\]
The complete integer norm is
\[
 \operatorname{Res}_Z(1+Z+\cdots+Z^6,A(Z))
                         =427439613252172921 .
                                                            \tag{B.145.13}
\]
The program computes this norm both as an integer resultant
and as the determinant of multiplication on
$\mathbf Z[Z]/\Phi_7$; it also retains the full relative
zeta numerator of degree twenty-four.
As a separate unit check, $\operatorname{ord}_7(3)=6$,
so $\Phi_7$ is irreducible modulo three, whereas
$A(Z)\bmod3=2Z+2Z^3+2Z^5+Z^6$ is nonzero and is not
a multiple of $\Phi_7$.
This proves (B.145.4).
Since $5\nmid18928$ and five is the only smaller odd prime
different from three, no admissible smaller prime can give
a constant pointed cyclic quotient at this field.
Seven is therefore the least witness prime at this field.

For clarity about the marking, $C(\mathbf F_3)$ consists of
$O$ and $(0,j)$ for $j\in\mathbf F_3$.
Their images under $\pi$ lie in $E[2]$, so their
$\chi([\tau P-P])$ values are zero.
The complete $k$-point table gives nonzero
$\chi_E(\pi(P))$ at every other point: there are eight
points at each of the three values $1,2,4$.
These are exactly the other twenty-four points.
For each, the automorphism shift is
$(2-1)\chi_E(\pi(P))\ne0$, proving the assertion for
all of those new pointings.
The three sets are permuted by $\tau$, and the pointed
character factors depend only on the displayed value,
so their relative class number is the same integer.

\medskip\noindent\textbf{The remaining twist and the original pointings.}
The criterion leaves twist one undecided in this example,
since $a/(u-1)=1$.
The complete finite calculation settles it too.
The exact norms for all twists are
\[
\begin{array}{c|r}
 t & \#J_{D_t}(k)/\#J_C(k)\\ \hline
 0,4,6&427439613252172921\\
 1&121740744925904896\\
 2,3,5&56515205339274361 .
\end{array}                                                  \tag{B.145.14}
\]
Every displayed integer is a three-adic unit.
For an original $\mathbf F_3$ point $P_0$, one has
$\chi_E(\pi(P_0))=0$, whereas $\chi_E(\pi(P))=1$.
Changing the pointing from $P$ to $P_0$ therefore adds one
to every degree-one exponent and $d$ to every degree-$d$
exponent.  Its split model is exactly twist one.
This proves the witness assertion for all four original
pointings, by the finite calculation rather than the
nonzero-shift criterion.

There is also an exact test through every constant extension
for each of these seven covers.
Work now in characteristic three and put
\[
 U=\bar\zeta+\bar\zeta^2+\bar\zeta^4,\qquad
                    \eta=1+2U .
\]
The quadratic-period identity $U^2+U+2=0$ gives
$\eta^2=-1$.
Reducing (B.145.12) and substituting $\bar\zeta T$ gives
\[
 \overline L_{1,1}(T)=1+\eta T-T^2=(1-\eta T)^2,
 \qquad
 \overline L_{1,t}(T)
                  =(1-\eta\bar\zeta^{\,t-1}T)^2 .
                                                            \tag{B.145.15}
\]
All other primitive character factors are its conjugates.
The nonzero reduced Frobenius roots consequently have
order four for $t=1$, and order twenty-eight for $t\ne1$,
since $\eta$ and a nontrivial seventh root have coprime
orders four and seven.
Under a degree-$n$ extension, every Frobenius root is
replaced by its $n$-th power.
Roots that reduce to zero cannot contribute a factor
divisible by three in the evaluation at one.
For the other roots, such divisibility occurs exactly
when their $n$-th power is one.
Hence
\[
 3\nmid\frac{\#J_{D_t}(\mathbf F_{27^n})}
                 {\#J_C(\mathbf F_{27^n})}
 \quad\Longleftrightarrow\quad
 \begin{cases}
  4\nmid n,&t=1,\\
  28\nmid n,&t\ne1 .
 \end{cases}                                                \tag{B.145.16}
\]
This is a proof for all extension degrees, not extrapolation
from a finite list of extensions.
The program separately checks the coefficient identities
in (B.145.15) and the exact orders four and twenty-eight
in $\mathbf F_3[Z]/\Phi_7$.

Finally the base Frobenius polynomial already computed gives
\[
 \#J_C(\mathbf F_3)=16,\qquad
 \#J_C(\mathbf F_9)=256,\qquad
 \#J_C(\mathbf F_{27})=16\cdot7\cdot13^2 .
                                                            \tag{B.145.17}
\]
A constant cyclic pointed cover of odd prime degree requires
a quotient of the rational Jacobian of that order.
The first two orders exclude all such covers.
At degree three the twist-one witness exists and the third
order makes seven the least admissible prime there.
This proves the claimed field-first minimum for every
original $\mathbf F_3$ marking.
\end{proof}
\paragraph{Check record.} GPT-6 Astra Pro, 2026-11-18: full proof and exact finite certificate written and checked. Root and configuration each completed full initial proof readings and separate complete final-addon readings, including the expanded program. Root independently executed the final certificate at05:43-05:44UTC, all PASS. Configuration independently calculated the exceptional twist norm and exact unit-root orders; its code/output reading does not claim a complete rerun. Effort not separately exposed.


\paragraph{Exact scope and predecessor comparison.}
The general statement proves B1 for the entire marked-curve
family satisfying (B.145.1) with nonzero automorphism shift.
It includes B.141 by taking the hyperelliptic involution in
genus two, and the genus-three example proves that its scope
is larger.  It imposes a real automorphism and eigencharacter
condition; it does not manufacture such an automorphism for
an arbitrary curve, and it does not settle the unrestricted
threshold $p>(2g-2)^2$.
The finite quartic certificate treats new
$\mathbf F_{27}$ pointings and, by the separate twist calculation,
the first-field minimum for all four original $\mathbf F_3$
pointings of this quartic.  It gives the exact constant-extension
unit test for each of the seven displayed covers.
It does not determine the least prime when arbitrary extension
degrees are allowed, the least witness among all covers over
every extension, or the full B6 census.

\paragraph{Check record.}
Written and checked by GPT-6 Astra Pro, finite-checks reader,
18 November 2026; the separate effort setting is not exposed.
All finite inputs and complete outputs are retained with the
program and JSON certificate.
The infinite family is proved by the automorphism and
Riemann--Roch argument, not inferred from the finite example.
Independent second-reading events are recorded separately.

\subsection{Discrete mixed-characteristic bases with finitely generated residue field}

\begin{claim}{B.146}{Proved}{Completion, Kummer kernels and the deep-unit subspace beyond finite local fields}
Let $D$ be finitely generated over $\mathbb Q$, with
a discrete valuation of mixed characteristic $(0,p)$
whose residue field $k$ is finitely generated over
$\mathbb F_p$.
Let $M$ be its completion, choose an algebraic closure
$\overline M$, and let $\Omega$ be the algebraic closure
of $D$ inside $\overline M$.
The following assertions hold, also after every finite
extension of the valued field $D$.

First, $D^h=M\cap\Omega$, and completion gives an
equivalence on finite separable extensions of $D^h$
and $M$.  Thus $G_M$ identifies with the chosen
decomposition subgroup of $G_D$, with its inertia
subgroup $I_M$.
For finite $L/M$, its algebraic part $L\cap\Omega$
is finite over $D^h$, dense in $L$, and $L$ has an
$M$-basis in that algebraic part.
The field $\Omega$ admits a complex embedding.

Second, the rational $p$-Kummer kernel of $M$ is
exactly the finite group $\mu(M)$.
There is a canonical injective $\mathbb Q_p$-linear map
\[
         \iota_M:M\lhook\joinrel\longrightarrow
                      H^1(G_M,\mathbb Q_p(1)),                 \tag{META.DV1}
\]
characterized on sufficiently small $x$ by
$\iota_M(x)=\kappa_M(\exp x)$.
It is compatible with finite extensions and their
Galois actions.
No assertion that (META.DV1) exhausts the cohomology
is made.

Finally the image of $I_M$ in $G_{\mathbb Q_p}$ is
$I_{L_0}$ for some finite extension $L_0/\mathbb Q_p$.
In particular the cyclotomic character on $I_M$ is
nontrivial, restriction
$H^1(G_M,\mathbb Q_p(1))\to H^1(I_M,\mathbb Q_p(1))$
is injective, and every successive closed derived
quotient of $I_M$ is nontrivial.
\end{claim}
\begin{proof}
Write $v_M(M^\times)=\mathbb Z$ and $e=v_M(p)>0$.
The completion is a complete discrete valuation
field with the same value group and residue field.
The real-valued henselian comparison in Temkin,
Theorem~1.9.5(ii), \S1.9.7 and Lemma~1.9.8,
identifies the henselization with the relative
separable closure in the completion
\citep{AN_TemkinValuations}.
In characteristic zero this gives $D^h=M\cap\Omega$.
The proof of Stacks Tag~09EJ, Lemma~15.115.2,
works for an arbitrary DVR and supplies descent
of every finite separable extension of $M$
\citep{AN_Stacks09EJ}.
Here is also the full-faithfulness comparison.
A henselian field is relatively separably closed
in its completion.  Since $D^h$ is relatively
algebraically closed in $M$, it is linearly disjoint
from $M$ in each finite Galois extension; hence
completion preserves finite extension degrees.
Every finite extension of $D^h$ is henselian,
and its field embeddings preserve the uniquely
extended valuation, so extend uniquely to embeddings
of its completion.
Conversely an embedding of those completions carries
the elements algebraic over $D^h$ into the relative
algebraic part of the target, which is exactly the
original finite henselian field.
This proves full faithfulness and the claimed
equivalence.  It identifies the Galois and inertia
groups by preserving the residue extension.
It also proves the assertion about $L\cap\Omega$;
any basis before completion remains a basis afterwards.
In particular the algebraic part is dense.
Finally an embedding of the algebraic closure of
$\mathbb Q$ in $\Omega$ into $\mathbb C$ extends
across a finite transcendence basis and then across
its algebraic closure, giving the asserted embedding.
These arguments used neither finiteness of the
residue field nor local compactness.

We next prove the Kummer statements, including the
distinction between deep units and all units.
The algebraic closure $k_0$ of $\mathbb F_p$ in $k$
is finite.  Indeed if $t$ is a finite transcendence
basis, the degrees of all finite constant subfields
are bounded by $[k:\mathbb F_p(t)]$.
Moreover
\[
                       \bigcap_{n\ge0}k^{p^n}=k_0.             \tag{META.DV2}
\]
To see this, take a normal proper model of $k/k_0$.
The divisor of an element in every $k^{p^n}$ has
all its integer coefficients divisible by every
$p^n$, and is therefore zero.
Normality makes that rational function regular
everywhere, and properness makes it algebraic over
$k_0$.  It consequently belongs to $k_0$.
The reverse inclusion holds because the finite
field $k_0$ is perfect.

Put $U_m=1+\mathfrak m_M^m$.
For an integer $m>e/(p-1)$ the convergent power
series for logarithm and exponential give inverse
group isomorphisms
\[
             U_m\xrightarrow{\log}\mathfrak m_M^m,
                 \qquad\log(u^p)=p\log u.                    \tag{META.DV3}
\]
For completeness, in the logarithm the terms
$x^j/j$ tend to zero for $v_M(x)>0$; for the
exponential, $v_M(j!)\le ej/(p-1)$ gives convergence
when $v_M(x)>e/(p-1)$.
The first term has strictly smaller valuation
than all subsequent terms on that smaller ball;
the formal inverse identities then converge and
give (META.DV3).
Repeated $p$-th powers carry all of $U_1$ into
$U_m$ after a uniformly bounded number $N$ of
iterations.  In fact, if $v_M(u-1)=a\ge1$, the
binomial expansion gives
\[
                  v_M(u^p-1)\ge\min(pa,a+e)>a.
\]
Thus $N=m$ is already a sufficient bound.

The topology on $U_m$ induced by the $p$-power
quotients of $M^\times$ is its full $p$-adic topology.
If $b^{p^n}\in U_m$, valuation gives $v_M(b)=0$,
and reduction gives $\overline b^{p^n}=1$.
In characteristic $p$ this forces $\overline b=1$,
so $b\in U_1$.
For $n\ge N$ we obtain the explicit inclusions
\[
 U_m^{p^n}\ \subset\ U_m\cap(M^\times)^{p^n}
                 \ \subset\ U_m^{p^{\,n-N}}.                 \tag{META.DV4}
\]
Completeness and (META.DV3) now give an injection
\[
 U_m=\varprojlim_n U_m/U_m^{p^n}
       \lhook\joinrel\longrightarrow
             \varprojlim_n M^\times/(M^\times)^{p^n}.
\]
The inverse-limit Kummer sequence identifies the
right side with $H^1(G_M,\mathbb Z_p(1))$.
Its possible $\varprojlim^1$ term is zero:
the groups of $p^n$-th roots of unity in $M$ are
finite, so their inverse system is Mittag--Leffler.
Continuous cocycles of a profinite group in
$\mathbb Q_p(1)$ have bounded image; multiplication
by a power of $p$ puts them in $\mathbb Z_p(1)$.
The same observation for coboundaries proves
\[
 H^1(G_M,\mathbb Z_p(1))\otimes_{\mathbb Z_p}\mathbb Q_p
                         =H^1(G_M,\mathbb Q_p(1)).
\]
After tensoring the injection with $\mathbb Q_p$,
(META.DV3) identifies its domain with $M$.
This is (META.DV1); explicitly, for any $x\in M$
and sufficiently large $r$,
\[
                \iota_M(x)=p^{-r}\kappa_M(\exp(p^rx)).
\]
The formula proves independence of $r$, as well
as extension compatibility and Galois equivariance.

Suppose $\kappa_M(a)=0$.
Some $p^j$ multiple of the integral Kummer class
is zero, so $a^{p^j}\in(M^\times)^{p^n}$ for every $n$.
The value group gives $v_M(a)=0$, and reduction
together with (META.DV2) puts $\overline a$ in
$k_0^\times$.
Hensel's lemma lifts it to a root of unity $\zeta$
of order prime to $p$.
Then $u=a/\zeta\in U_1$ and $\kappa_M(u)=0$.
Raise it to $p^N$ to enter $U_m$.
The injection just proved gives $\log(u^{p^N})=0$,
and (META.DV3) gives $u^{p^N}=1$.
Thus $a$ is a root of unity; the converse follows
by tensoring the torsion Kummer class with $\mathbb Q_p$.
There are only finitely many roots of unity in $M$:
the prime-to-$p$ ones inject into the finite group
$k_0^\times$, while a primitive $p^n$-th root has
normalized $p$-adic value
$v_p(\zeta_{p^n}-1)=1/(p^{n-1}(p-1))$.
The latter follows from the Eisenstein polynomial
$\Phi_{p^n}(1+T)$, and those denominators must divide
$e$.  This also includes $p=2$.
All these assertions apply to finite extensions,
whose residue fields remain finitely generated.

It remains to justify the useful finite-local
quotient of inertia without asserting wild freeness.
Embed $\overline{\mathbb Q}_p$ in $\overline M$
over the canonical copy $\mathbb Q_p\subset M$,
and put
\[
                  T=M^{\mathrm{nr}}\cap\overline{\mathbb Q}_p.
\]
Finite unramified constant extensions stay etale
after base change to the valuation ring of $M$,
so $\mathbb Q_p^{\mathrm{nr}}\subset T$.
Every finite $L'/\mathbb Q_p$ inside $T$ has
$e(L'/\mathbb Q_p)\mid e$, since $M^{\mathrm{nr}}$
has the same value group as $M$.
Also
\[
 [L'\mathbb Q_p^{\mathrm{nr}}:\mathbb Q_p^{\mathrm{nr}}]
                                      =e(L'/\mathbb Q_p).
\]
For this equality, take the maximal unramified
subextension of $L'/\mathbb Q_p$ and an Eisenstein
polynomial for a uniformizer of the remaining
totally ramified extension.
After unramified base change the polynomial remains
Eisenstein, giving exactly the displayed degree.
The union $T/\mathbb Q_p^{\mathrm{nr}}$ thus has
bounded finite subextension degree, at most $e$,
and is finite.  Hence $T=L_0\mathbb Q_p^{\mathrm{nr}}$
for one finite $L_0/\mathbb Q_p$.
The image of $I_M$ in $G_{\mathbb Q_p}$ is closed
and has fixed field $T$, so is exactly $I_{L_0}$.
The local cyclotomic character has infinite image
there.  Thus $\mathbb Q_p(1)^{I_M}=0$, and
inflation--restriction proves the asserted injection.

Finally $I_{L_0}$ has strictly decreasing closed
derived series.  Its wild inertia is free pro-$p$
of infinite rank by Iwasawa Theorem~2(ii)
\citep{AN_Iwasawa1955}.
For every $r$, mapping finitely many free generators
to the adjacent elementary matrices of
$\operatorname{UT}_{2^r+1}(\mathbb F_p)$ shows that
its $r$-th derived subgroup is nontrivial:
balanced commutators use
$[1+E_{ab},1+E_{bc}]=1+E_{ac}$ and terminate at
$1+E_{1,2^r+1}$.
Hence every derived subgroup of $I_{L_0}$ is
nontrivial.  From the first commutator onward it
is pro-$p$, and a nontrivial pro-$p$ group has
a nontrivial abelian quotient, by taking any
nontrivial finite $p$-group quotient and then
a quotient of order $p$.
The first derived quotient is already nontrivial
by tame inertia.
Surjectivity $I_M\to I_{L_0}$ carries closed
derived subgroups onto closed derived subgroups,
so proves strictness for $I_M$ as well.
\end{proof}
\paragraph{Check record.} checked only in this writing, 2026-11-18: GPT-6 Astra (OpenAI), ChatGPT Pro mode; internal effort value not exposed. Own detailed proof written and reread; primary source hypotheses matched. Root full reading of B146 and all B147 lines255--782 passed: perfect core, deep-unit topology and injection, finite Kummer kernel, finite-local inertia quotient, strict derived groups, imperfect residue, nil radical relation, split fibre before cyclic q, scalar differences and exact ratio sign, analytic density without compactness, full norm/graph/descent and rank-one strict-scope separation


\begin{claim}{B.147}{Proved}{Reconstruction at derived and nilpotent inertia fields with arbitrary finite residue transcendence}
Keep $D,M,\Omega,k$ of Claim~B.146, and put
$I=I_M\subset G_D$.
For the closed derived series $I^{(0)}=I$ and
$I^{(r+1)}=\overline{[I^{(r)},I^{(r)}]}$, and the
closed lower central series $\gamma_c(I)$, set
\[
 F_r=\Omega^{I^{(r)}}\quad(r\ge0),\qquad
 K_{\mathrm{nil}}=\Omega^{\,\bigcap_{c\ge1}\gamma_{c+1}(I)}.
                                                               \tag{META.DV5}
\]
Every finite extension of every $F_r$ and of
$K_{\mathrm{nil}}$ satisfies full closed-point-preserving
relative etale reconstruction for finite-type schemes
localized at universal homeomorphisms, with the
structural isomorphism retained.
The property also descends to any subfield over which
one of these reconstructed fields is algebraic.

In particular, for the Gauss $p$-valuation on
$D=\mathbb Q(t_1,\ldots,t_s)$, this gives positive
families with residue transcendence degree $s$ for
every $s\ge1$ and every prime $p$.
No assertion for all proper normal subgroups of
$G_M$ is included.
\end{claim}
\begin{proof}
We first establish the arithmetic inputs, and then
give the precise change in the analytic norm proof.
The tame/wild inertia theorem for an arbitrary DVR,
Stacks Tag~0BUA, Lemma~58.13.5, gives
\[
             1\longrightarrow P\longrightarrow I
               \longrightarrow \prod_{\ell\ne p}\mathbb Z_\ell(1)
               \longrightarrow1,                            \tag{META.DV6}
\]
with $P$ pro-$p$ \citep{AN_Stacks0BUA}.
The source does not require a perfect residue field.
All the prime-to-$p$ tame orders occur here by
adjoining roots of a uniformizer to the strict
henselization; units have roots of those orders by
Hensel's lemma.  The strict henselian residue field
is $k^{\mathrm{sep}}$, which need not be perfect.
Nothing below replaces it by an algebraic closure.

\medskip\noindent\textbf{The derived fields.}
Let $r\ge2$, put $K=F_r$ and $A=F_{r-1}$, and
let $E/K$ be finite.
If $a\in E^\times$ has zero rational Kummer class,
choose a finite Galois $D_1/D$ containing $a$ and
generators of $E/K$, and replace $E$ by $KD_1$.
Its completion $M_1$ is finite over $M$ and again
has finitely generated residue field.
Set
\[
 J=I\cap G_{D_1}=I_{M_1},\quad
 H=I^{(r)}\cap J,\quad \Gamma=J/H,\quad
 C=(I^{(r-1)}\cap J)/H,\quad V=\mathbb Q_p(1).
\]
The Kummer class on $J$ vanishes on $H$, which fixes
$V$, so inflates from a unique element of
$H^1(\Gamma,V)$.
For $g\in I^{(r-1)}$, Galoisness of $D_1/D$ makes
$g$ normalize $J$.
It acts trivially on $C$, on $\Gamma/C$ and on $V$.
Inflation--restriction gives a two-step filtration
of $H^1(\Gamma,V)$ on whose pieces $g$ is the identity:
\[
 0\longrightarrow H^1(\Gamma/C,V)
  \longrightarrow H^1(\Gamma,V)
  \longrightarrow H^1(C,V)^{\Gamma/C}.
\]
Thus $(g-1)^2=0$.
The particular Kummer vector has finite $g$-orbit,
because $a\in D_1$ has one.  If $g^n$ fixes it,
the equality $g^n-1=n(g-1)$ on that vector makes
$g$ fix it as well.
Therefore $\kappa_J(g(a)/a)=0$.
The inertia restriction injection and the finite
root-of-unity Kummer kernel of Claim~B.146 give
\[
       g(a)/a\in\mu(M_1)\quad(g\in I^{(r-1)}),\qquad
                         a^m\in A\text{ for some }m\ge1.       \tag{META.DV7}
\]
One exponent works for all $g$, since $g$ preserves
the one field $M_1$ and $\mu(M_1)$ is finite.
This proof includes every finite $E/K$.

By Claim~B.146 the derived inclusions are strict,
so choose $b\in K\setminus A$.
Both fields contain all roots of unity, since the
cyclotomic action of inertia is abelian.
If the classes of $b$ and $1-b$ both vanished,
(META.DV7) would give
$g(b)=\zeta_gb$ and $g(1-b)=\eta_g(1-b)$, with
$\zeta_g,\eta_g\in\mu_\infty\subset A$.
If some $\zeta_g\ne\eta_g$, adding expresses $b$
as an element of $A$; if all are equal, adding
makes them one, so again every $g$ fixes $b$.
This contradiction supplies a nonzero Kummer class
over $K$, surviving every finite extension by
rational restriction and corestriction.

For projection degree $d$ use the grids
\[
       \mathcal G_0=\{u+iv:0\le u,v\le4d\},
                         \qquad \mathcal G_1=b+\mathcal G_0.
\]
If a monic equal-degree rational function $R=P/Q$
has zero Kummer classes at all the grid values,
(META.DV7) makes $gR(a)/R(a)$ a root of unity at
every first-grid centre.
In a complex embedding of $\Omega$, the difference
of the squared absolute values of its monic
numerator and denominator has degree at most $4d$
separately in $u,v$.
Successive interpolation makes it zero.
Writing $z=u+iv,w=u-iv$ gives the product identity
$U(z)\overline U(w)=V(z)\overline V(w)$; monicity
then gives $gR=R$.
The translated grid similarly gives
$gR(gb+Z)=R(b+Z)$.
Choose $g$ moving $b$.
A rational function with the nonzero additive
period $gb-b$ is constant, by its finite pole set
or its leading polynomial finite difference.
Its value at infinity is one, so $R=1$.
This proves the required uniform two-grid test.

\medskip\noindent\textbf{The nilpotent union.}
Put $N=\bigcap_c\gamma_{c+1}(I)$ and $K=\Omega^N$.
Every finite inertia quotient is an extension of
a finite $p$-group by a cyclic prime-to-$p$ group.
In a nilpotent quotient all Sylow subgroups commute,
and all tame cyclic quotients survive.
Consequently
\[
       I/N=Q_{(p)}\times\prod_{\ell\ne p}\mathbb Z_\ell(1),
                                                               \tag{META.DV8}
\]
where $Q_{(p)}$ is pro-$p$.
Let $e=v_M(p)$ for the integer normalization.
Choose an odd prime $\ell\nmid ep(p-1)$, and put
$b=p^{1/\ell}$.
The strict henselization has value group
$e^{-1}\mathbb Z$ in the normalization $v(p)=1$;
thus adjoining $b$ gives a tame cyclic extension
of degree $\ell$.
Choose a central element in the pro-$\ell$ factor
of (META.DV8) moving $b$, and lift it to $\tau\in I$.
Its $p$-cyclotomic value is one: a pro-$\ell$ group
maps neither to a nontrivial pro-$p$ subgroup of
$\mathbb Z_p^\times$ nor to its torsion subgroup
of order dividing $p-1$.
Inertia fixes all other roots of unity, so $\tau$
fixes $\mu_\infty$ and $\tau b=\omega b\ne b$.

Here is the stable-kernel argument over this
complete DVR.
If $\kappa_E(a)=0$ for a finite $E/K$, then
compatible $p$-power roots of $a$ belong to $E$:
$G_E$ fixes $\mathbb Z_p(1)$, so its integral
$H^1$ is a torsion-free group of homomorphisms.
Choose finite Galois $D_1/D$ containing the data.
Let $F$ be the compositum of $D_1$, the relative
strict henselization, and $\mu_{p^\infty}$;
let $R/F$ be generated by all $p$-power roots of
the finitely many inertia conjugates of $a$.
Both fields are Galois over the strict henselization,
$R\subset KD_1$, and
$D_R=\operatorname{Gal}(R/F)$ is a torsion-free
closed subgroup of $\mathbb Z_p(1)^m$.
The image of $N$ in the Galois group of $R$
over the strict henselization is finite, since
$KD_1/K$ is finite.
Centrality of $\tau$ modulo $N$ puts
$[\tau,D_R]$ both in that finite image and in
$D_R$, hence makes it trivial.

For $J=I_{M_1}$ and $J_0=J\cap\ker\chi_p$,
the group $J_0$ surjects onto $D_R$.
The commutator equality, including the trivial
coefficient action of $\tau$, therefore gives
$\kappa_{J_0}(\tau a)=\kappa_{J_0}(a)$.
Restriction from $H^1(J,V)$ to $H^1(J_0,V)$ is
injective: $J/J_0$ is abelian with nontrivial
cyclotomic character, and comparing a cocycle
at $xy$ and $yx$ gives
$(\chi(x)-1)c(y)=(\chi(y)-1)c(x)$, so it is a
coboundary.
Claim~B.146 also gives inertia restriction
injectivity and the Kummer kernel in $M_1$.
We conclude
\[
             \kappa_E(a)=0\quad\Longrightarrow\quad
                                  \tau(a)/a\in\mu_\infty.      \tag{META.DV9}
\]
Choose the complex embedding with $b>0$.
Since
$|1-\omega b|^2-|1-b|^2=2b(1-\operatorname{Re}\omega)>0$,
the class of $1-b$ is nonzero and remains so over
every finite extension.
For $\mathcal G_0,b+\mathcal G_0$, the preceding
absolute-square interpolation applied to $\tau$
first gives $\tau R=R$ and then
$R(\tau b+Z)=R(b+Z)$.
The nonzero period gives $R=1$.
Thus both arithmetic inputs hold at every prime,
including $2$ and $3$.

\medskip\noindent\textbf{Local normalization and cyclic extensions.}
For either field $K$ just treated, start with an
arbitrary finite $E/K$, a smooth geometrically
integral affine source $X/E$, and a pinned relative
map to $\mathbb A^1_E$.
If its actual point map is nonconstant, each
actual fibre is a proper closed subset.
Choose a finite Noether projection of degree $d$
and the two grids above.
On a principal open $V$ of affine space avoiding
the branch locus and all the grid-fibre images,
the pulled-back and corestricted Kummer classes are
\[
          \xi_a=\beta_a+\sum_j\lambda_{aj}\kappa(h_j),
                      \qquad\lambda_{aj}\in\mathbb Q_p.        \tag{META.DV10}
\]
The geometric Picard group is zero, and the
geometric units modulo constants are freely
generated by the prime boundary equations $h_j$,
which proves this decomposition by Kummer theory
and inflation--restriction.
The actual evaluation is $\kappa(P_t(a))$ for
the monic degree-$d$ norm polynomial $P_t$;
the factor $(-1)^d$ from corestricting $Y-a$
is torsion and disappears in rational Kummer theory.

Choose one finite Galois $D_1/D$ containing
all coefficients, grid centres, a normalizing
parameter $t_*$ and its actual norm values,
all coordinates of the finitely many points in
one good projection fibre, and generators of $E/K$.
Let $M_1$ be its completion and $E_1=KD_1$.
The normality of the inertia kernel gives
$N_1=G_{E_1}\triangleleft G_{M_1}$ and
$\chi_p(N_1)=1$.
Evaluation at $t_*$ expresses
$\beta_a=\kappa(P_{t_*}(a))-\sum_j\lambda_{aj}
\kappa(h_j(t_*))$ as the restriction of a class
over $M_1$.
Thus (META.DV10), including its Tate coefficient
action, is $G_{M_1}$-invariant.
Comparison at $t$ and $\sigma t$ puts all grid
values of $P_{\sigma t}/\sigma P_t$ in the stable
kernel over a common finite extension.
The proved two-grid test yields
$P_{\sigma t}=\sigma P_t$ uniformly.

The nonzero Kummer class restricts nontrivially
to $N_1$.
Its integral image is $p^u\mathbb Z_p$; taking
an actual $p^u$-th root of its argument in $E_1$
makes that image all of $\mathbb Z_p$.
Adjoin this argument and $\mu_{p^s}$ to $M_1$
inside the local fixed field of $N_1$.
The resulting cyclic extension $L/M_1$ generated
by a $p^s$-th root has degree $q=p^s>d+1$ and
receives a surjection from $N_1$.
Normality persists under these finite enlargements.
Claim~B.146 supplies an $M_1$-basis of $L$
in $L\cap\Omega$, and
$E_1(L\cap\Omega)/E_1$ is finite:
the henselian $E_1$ contains $M_1\cap\Omega$
after the indicated algebraic enlargements.
All scalar algebraic parameter sets used below
are consequently dense.

Actual maps are valuation-continuous over every
finite extension of $E_1$.
Indeed CHW Proposition~3.6 and the representable
etale inverse-image comparison in \S4.13 and
Corollary~A.9 give etale-open continuity, and
JTWY Corollary~6.13 identifies that topology
with the valuation topology for these henselian
non-separably-closed fields
\citep{CarlsonHaineWolf2024,AN_JTWY2024}.
The nonzero Kummer class established above
verifies non-separable-closedness at every finite
extension.  Affine charts handle nonseparated
finite-type schemes.

\medskip\noindent\textbf{The scalar-error step using deep units only.}
The chosen good etale projection fibre already
splits over $M_1$, before the cyclic $L/M_1$
was chosen.  Work near it with the $d$ analytic
inverse branches; no further base enlargement
is used to split this fibre.
For actual parameters in $L\cap\Omega$, local
norm equivariance puts the coefficients of $P_t$
in $L$; continuity gives continuity of those
coefficients on the dense algebraic parameter set.
At such a parameter put
\[
 \theta_t=\kappa_L(P_t(a))-\beta_a|_{G_L}
                  -\sum_j\lambda_{aj}\kappa_L(h_j(t)).
\]
It vanishes on $N_1\cap G_L$.
By inflation--restriction it inflates from
$G_L/(N_1\cap G_L)$.
Conjugation by $N_1$ is the identity on this
quotient and on $V$, since $N_1$ is normal in
$G_{M_1}$, $L/M_1$ is Galois, and $\chi_p(N_1)=1$.
Hence every $\theta_t$ is $N_1$-invariant.

Choose an actual centre $t_0$ and shrink the ball
so that, for every grid centre and every $h_j$,
the ratios $P_t(a)/P_{t_0}(a)$ and
$h_j(t)/h_j(t_0)$ lie in a deep unit group of $L$.
This uses only their actual continuity, not an
analytic extension of $P_t$.
Claim~B.146 now gives the exact identity
\[
 \theta_t-\theta_{t_0}
   =\iota_L\left(
      \log\frac{P_t(a)}{P_{t_0}(a)}
        -\sum_j\lambda_{aj}\log\frac{h_j(t)}{h_j(t_0)}
             \right).                                        \tag{META.DV11}
\]
Injectivity and Galois equivariance of $\iota_L$
make the scalar in parentheses fixed by $N_1$.
Since $N_1$ surjects onto $\operatorname{Gal}(L/M_1)$,
that scalar lies in $M_1$.
This is the exact replacement for a logarithm on
all of $H^1(G_L,V)$; residue-unit classes need not
belong to the deep-unit subspace.

Write $\tau$ for a generator of $\operatorname{Gal}(L/M_1)$
and $r_i(t,Z)=\tau^iP_t(Z)/P_t(Z)$ at actual points.
Subtract the $\tau^i$-transform of (META.DV11).
With the signs in the order transform minus original,
the result is
\[
 \log\frac{r_i(t,a)}{r_i(t_0,a)}
 =\sum_j\lambda_{aj}
   \left(
    \log\frac{h_j(\tau^it)}{h_j(\tau^it_0)}
       -\log\frac{h_j(t)}{h_j(t_0)}
   \right).                                                   \tag{META.DV12}
\]
After a further shrink, all finite linear combinations
are in the exponential convergence ball.
Exponentiating therefore recovers each ratio value
as $r_i(t_0,a)$ times an analytic function.
No global branch of logarithm on arbitrary residue
units has been chosen.
Differentiating gives exactly the matching-block
logarithmic equations of the norm proof in
Claim~B.136.

\medskip\noindent\textbf{Analytic recovery and the algebraic graph.}
We explain why that proof applies over the present
complete DVR without a local-compactness hypothesis.
Minimize the finite integer sum of conjugate gcd
degrees on the actual parameter ball.
At a minimizing point the relevant nonzero
subresultants remain nonzero by continuity;
the minimum then makes every gcd degree locally
constant on the dense actual set.
Recenter (META.DV11)--(META.DV12) at this
minimizing actual point.
For each reduced monic degree-$e_i$ ratio, its
values at $2e_i$ centres determine its numerator
and denominator analytically by (META.DV12).
The interpolation kernel is zero because
$\delta A B-A\delta B$ has degree at most $2e_i-1$
and has $2e_i$ roots; coprimality and monicity
then make both variations zero.

Use scalar coordinates for $\operatorname{Res}_{L/M_1}$,
with an algebraic $M_1$-basis and centre $t_0$.
After extending coefficients to $L$, the $q$
embedding blocks are permuted by the semilinear
cyclic action even if $t_0\notin M_1$.
The extended derivations of algebraic meromorphic
germs and the $4d+1$-centre rational identity
classify each nonconstant root by a unique block.
The rational identity has denominator degree at
most $4d$ and numerator degree at most $4d-1$,
so the centre count still proves it identically.
The divisor cocycle gives
$D_i=\tau^i B_0-B_0+C_i$.
A constant root is algebraic over $D$ after one
actual specialization, and its occurrence on an
open set contradicts proper closed actual fibres.
Thus $C_i=0$.
A negative coefficient of $B_0$ would give
$q-1>d$ distinct poles among the reduced ratios.
Their denominator least common multiple divides
the actual degree-$d$ polynomial at actual points,
which is impossible after a degree-preserving
specialization.

That specialization uses a nonempty analytic open,
not an algebraic open.
There are only finitely many nonzero minors and
subresultants to avoid.
A nonzero convergent power-series germ over a
complete nonarchimedean characteristic-zero field
cannot vanish on a scalar open ball: the
one-variable assertion follows by the first
nonzero term near an accumulation point, and
successive variables prove the multivariable
assertion.
Thus their complements are open and dense; a
finite intersection is nonempty and meets the
dense algebraic scalar parameter set.
Neither Baire's theorem nor local compactness
is used.

The effective type-zero divisor is already the
analytic monic reduced denominator $Q$ of $r_1$.
The exact derivatives obtained from (META.DV12)
give
$d\log Q(t,a)=\sum_j\lambda_{aj}\,dh_j/h_j$.
The constants of the meromorphic-germ differential
field are $L$, as follows term by term in a
power series.
Claim~B.135 therefore gives $Q\in L(t)[Z]$.
Its values at actual algebraic parameters are
algebraic over $D$, since it divides $P_t$ there;
finite rational interpolation gives
$Q\in\Omega(t)[Z]$.
Residues of its logarithmic derivatives make the
exponents integral.
Normalizing once more at an actual parameter
and using the two grids gives $P_t=Q_tC$ for a
single constant monic polynomial $C$.
Any root of a positive-degree $C$ would force
every projection fibre to meet one proper closed
actual fibre, contrary to finite-image dimension.
Hence $C=1$, and all projection norms have been
reconstructed.

The graph argument of Claim~B.136 is algebraic
over $\Omega$ and uses no additional local-field
property: one recovered norm bounds graph closure
finitely over a dense source open; a generic finite
linear projection separates its dominating
components and is birational on the generic
output fibres.
If their degrees over the source are $e_i$, the
distinct minimal output polynomials have degrees
$e_i d'$ and their product divides the recovered
degree-$d'$ norm.  Thus $\sum_i e_i\le1$, forcing
a unique component of degree one.
The finite coefficient conditions descend from
$\Omega(Y)$ to a nonempty open set over $\Omega$
by clearing denominators and expanding in $Y$.
Thus the actual map agrees densely with a rational
function, which descends to the original field
by equivariance.
The valuation continuity already proved excludes
poles using algebraic transverse curves, and
excludes exceptional point values by approximation
in an etale chart.
Normality of the smooth source makes the function
regular.
Geometric components descend after a finite
extension by the same uniqueness.
CHW Theorem~4.18 and Corollary~2.32 give the full
stated reconstruction property
\citep{CarlsonHaineWolf2024}.

\medskip\noindent\textbf{Descent and strict additional scope.}
The proof has included every finite extension of
$F_r$ for $r\ge2$ and of $K_{\mathrm{nil}}$.
For an algebraic subfield, base change a pinned map
along the algebraic extension using the compatible
finite constant-object slices.
Every finitely presented etale object, morphism
and finite covering family descends to one finite
scalar extension, so these compatible slice functors
give the required inverse-image functor on the limit
site, preserving finite limits and covers.
Closed points remain closed since their residue
coordinates remain algebraic.
Reconstruct there and descend the unique resulting
scheme morphism by its equivariant actual point map.
This is the finite-site algebraic descent already
proved in Claim~B.140; it uses no finite-local-field
assumption.
Applying it to $F_0,F_1\subset F_2$ proves their
assertions, including every finite extension by
taking its finite compositum with $F_2$.

For $D=\mathbb Q(t_1,\ldots,t_s)$ the Gauss
valuation has group $\mathbb Z$ and residue
$\mathbb F_p(t_1,\ldots,t_s)$.
All fields in (META.DV5) are algebraic over $D$,
so have transcendence degree $s$ over $\mathbb Q$;
their residue fields are algebraic over that
residue field and hence have transcendence degree
$s$ over $\mathbb F_p$.
For $s>0$ they are not isomorphic to any henselian
subfield of $\overline{\mathbb Q}_p$ carrying its
induced $p$-adic valuation.
Indeed both valuations would be nontrivial of
rank one on the same non-separably-closed field.
If inequivalent they are independent, contradicting
the two-henselian-valuations theorem of F.\ K.\ Schmidt,
in the opened formulation Kuhlmann Theorem~9.44
\citep{AN_Kuhlmann9}; if equivalent their residue
transcendence degrees would agree, whereas every
residue field inside $\overline{\mathbb Q}_p$ is
algebraic over $\mathbb F_p$.
Thus this is a strictly additional field family.

The argument deliberately makes no assertion that
wild inertia for these residue fields is free,
or that all its nontrivial normal subgroups are
nonabelian.
Consequently it does not extend the all-proper-normal
conclusion of Claim~B.143 to the present bases.
That normal-subgroup boundary, other algebraic
subgroups, and more general nondiscrete valuations
remain separate questions.
No priority or external human verification is claimed.
\end{proof}
\paragraph{Check record.} checked only in this writing, 2026-11-18: GPT-6 Astra (OpenAI), ChatGPT Pro mode; internal effort value not exposed. Own detailed proof written and reread; primary source hypotheses matched. Root full reading of B146 and all B147 lines255--782 passed: perfect core, deep-unit topology and injection, finite Kummer kernel, finite-local inertia quotient, strict derived groups, imperfect residue, nil radical relation, split fibre before cyclic q, scalar differences and exact ratio sign, analytic density without compactness, full norm/graph/descent and rank-one strict-scope separation


\subsection*{Targeted primary-source audit for Statements B.43--B.55}

This audit records passages actually opened on 18 November 2026 by
GPT-6 Astra Pro in Pro mode; a separate effort setting is not exposed.
It verifies the external inputs specified below and the comparison
diagrams used in B.49 and B.54.  It does not claim an independent
re-proof of every statement in the interval B.43--B.55.  In particular,
the construction of compatible logarithmic special-fibre isomorphisms
in B.48 and the extension from covering categories to the whole small
\'etale site in B.49 are separate arguments of this paper, rather than
assertions attributed to Lepage.

\paragraph{Hoshi--Tsujimura: slimness, not an unstated reconstruction theorem.}
The revised March 2023 author PDF was read at Definition~4.3,
Lemma~4.4 with its proof, and Corollary~4.5 with its proof, on
printed pp.~11--12.  Definition~4.3 defines slimness by trivial centres
of all open subgroups and states its equivalent formulation using
centralizers of normal open subgroups.  Corollary~4.5(1) applies when
the field is finitely generated and transcendental over some subfield.
Part~(3) applies when the field embeds in the fraction field of a
Noetherian local domain of mixed characteristic.  Thus every finitely
generated characteristic-zero field occurring in B.43 has the stated
slimness: use (1) for positive transcendence degree and a localization
of the ring of integers at a nonzero prime in (3) for a number field.
Taking the normal core of an arbitrary open subgroup gives triviality
of its centralizer in the whole group.  The same source consequently
supplies exactly the centralizer property of $G_{\mathbb Q}$ used in
B.47.  This does not by itself assert triviality of a centralizer in
the much larger group $\operatorname{Aut}(\overline K)$; B.47 proves
that additional assertion by its formal-arc argument.
\citep[Definition~4.3, Lemma~4.4, Corollary~4.5(1),(3)]{HoshiTsujimura2023}

\paragraph{Aschenbrenner--van den Dries--van der Hoeven and continuity.}
The opened author-hosted book PDF has creation and modification metadata
7 January 2025.  The following numbers and printed pagination refer
to that online version; their pagination in the 2017 printed volume
has not been independently compared.  Corollary~3.3.22, on printed
p.~124, characterizes algebraic closedness of an equicharacteristic-zero
valued field by henselianity, algebraically closed residue field, and
divisible value group.  The definition immediately before
Proposition~3.3.25, on p.~125, specifies the universal property of a
henselization with uniqueness of the valued embedding.  That
proposition constructs it; Corollary~3.3.26 says it is algebraic and
immediate.  These statements justify the precisely stated uses in
B.47.  Lemma~3.4.2 and its full proof, on p.~135, say that a valuation
is henselian if and only if both its coarsening and the induced
valuation on the residue field are henselian.  No finite-rank or
discreteness hypothesis is imposed in this lemma.  The finite-rank
argument in B.55 is used to obtain a nontrivial rank-one coarsening,
not to satisfy an additional hypothesis of the source.
\citep[Corollary~3.3.22, Proposition~3.3.25, Corollary~3.3.26,
Lemma~3.4.2]{ADH2017}

Khassa--Khanduja's publisher PDF was read at Lemma~2.1 and its entire
proof, printed p.~159.  It states continuity of every automorphism
of a henselian valued field that is not separably closed, for the
valuation topology; the rank is arbitrary.  Its proof transports
the henselian valuation by the automorphism and invokes the
dependence theorem for two henselian valuations on a field that is
not separably closed.  The stronger conclusion used in B.47,
preservation of the normalized discrete valuation, is a separate
elementary consequence: equal rank-one valuation topologies yield
the same valuation ring, and an order automorphism of the discrete
value group is the identity.  It is not attributed to the source
for arbitrary higher-rank valuations.
\citep[Lemma~2.1]{KhassaKhanduja2010}

\paragraph{Lepage: the full arithmetic covering categories.}
The publisher PDF was read at Theorem~1.6, printed p.~1437,
Theorem~1.8 and its whole proof, pp.~1438--1439, and
Corollary~1.9, p.~1439.  The precise distinction is as follows.
Theorem~1.6 assumes an fs log-regular scheme $X$ and a prime set
$L$ consisting of primes invertible on $X$.  Restriction from $X$
to its log-trivial open $X_{\mathrm{tr}}$ gives an equivalence
between the corresponding finite Kummer--\'etale and finite
\'etale covering categories of pro-$L$ type.
Theorem~1.8 assumes a Noetherian henselian local base $B$, with
closed point $s$, and a connected fs log scheme $X$ whose
underlying scheme is proper over $B$.  Pullback gives
\[
 \operatorname{KCov}(X)\xrightarrow{\sim}
 \operatorname{KCov}(X_s).
\]
Here $X_s$ carries the pulled-back log structure.  Neither geometric
base change to an algebraic closure of the residue field nor passage
to the geometric log fundamental group occurs in this conclusion.
The proof retains the arithmetic extension over the ordinary
fundamental group of $B$ and passes through its strict henselization.
In contrast, the convenient combined Corollary~1.9 explicitly assumes
a complete DVR.  The use over a merely henselian DVR in B.49 is
therefore correctly based on Theorems~1.6 and 1.8 themselves.
\citep[Theorems~1.6, 1.8, Corollary~1.9]{Lepage2013}

For a henselian DVR $R$ of residue characteristic zero, every positive
integer is a unit of $R$.  The vertical-and-marked log structure of
the stable model is fs and log regular: smooth charts and the node
chart $xy=u\pi^d$ are toroidal.  Taking $L$ to be all primes in
Theorem~1.6 retains every finite covering.  The special fibre is
geometrically connected and the underlying stable model is proper.
Thus all hypotheses for both comparison functors hold.

The following formal explanation spells out the functoriality needed
in B.49.  Write $j_X:U_X\hookrightarrow X$ for the log-trivial open
and $i_X:X_s\to X$ for the strict closed-fibre immersion.  Put
\[
 \mathcal E_X=i_X^*\,(j_X^*)^{-1}:
 \operatorname{FEt}(U_X)\xrightarrow{\sim}
 \operatorname{KCov}(X_s),
\]
where a quasi-inverse is supplied with its unit and counit.
Both original functors are actual pullback functors.  For a morphism
$q:X'\to X$ of the marked log models, the pullback squares commute
up to their canonical pullback isomorphisms.  A prescribed comparison
on the closed fibre consequently lifts uniquely to a comparison of
extended covers, once the closed-fibre identification is specified:
this is precisely full faithfulness of $i_X^*$.  Associativity and
all cocycle identities follow because their two sides have the same
closed-fibre restriction and this restriction is faithful.  The same
argument applies to a finite diagram at once.  In particular it
applies to mark-forgetting contractions whenever the special-fibre
isomorphisms commute with their logarithmic maps, as proved separately
in B.48.  It also applies to finite valued field extensions and the
Galois descent diagrams; one does not choose unrelated outer group
isomorphisms for the different marked sets.

For the arithmetic assertion, apply the same two equivalences also
to the log trait $B=\operatorname{Spec}R$ itself, which is proper
over itself.  They give
\[
 \operatorname{FEt}(K)\simeq\operatorname{KCov}(B)
 \simeq\operatorname{KCov}(s).
\]
The pullback square for the structure morphism $X\to B$ commutes
with both comparison functors.  An isomorphism of log special fibres
\emph{over this fixed log point} therefore induces an equivalence
over $\operatorname{FEt}(K)$, or over $BG_K$ with the corresponding
structural 2-isomorphism.  An isomorphism only after geometric
base change would not supply this assertion.  Identifying the actual
cusp residue groups $G_{K(a)}$ still uses B.48's local marked charts
and the equality of actual fields $K(F(a))=K(a)$; it is not a
consequence of an unlabelled geometric comparison alone.

\paragraph{AHLS: the precise non-fs approximation.}
The source inspected is arXiv:2411.13662v1; its title-page margin
dates this version 20 November 2024.  A previous bibliography
annotation saying 22 November should be corrected.  The targeted
statements and their proof passages were read in the version fixed
by that identifier.  Definition~4.2.1, printed p.~53, defines a
finite Kummer--\'etale map to have an \emph{integral} underlying
scheme map.  Remark~4.2.2 expressly warns that it need not be finite.
Proposition~4.2.6, p.~55, is an equivalence of categories of these
covers under saturated base change for an affine charted inverse
system of qcqs saturated log schemes.  Definition~3.4.4, p.~40,
requires compatible strict \'etale chart covers arising from direct
systems of saturated prelog rings.  It is stronger data than just
an uncharted inverse system of the underlying schemes.
\citep[Definition~3.4.4, Definition~4.2.1, Remark~4.2.2,
Proposition~4.2.6]{LA_AHLS2024}

Lemma~5.3.2 and the paragraph after it, pp.~65--66, apply to a
microbial valuation ring of residue characteristic zero.  They
provide a filtered union of local excellent regular subalgebras
$\Lambda_i$, each containing the chosen pseudouniformizer $\pi$,
with local maps to $R$ and strict normal-crossings divisor
$V(\pi)_{\mathrm{red}}$.  They can be henselian when $R$ is.
Its characteristic-zero proof uses resolution over $\mathbb Q$,
so the subalgebras can be $\mathbb Q$-algebras.  With
$M_i=\Lambda_i\cap\Lambda_i[1/\pi]^\times$, the paragraph after
the lemma identifies the full valuative prelog ring as their
filtered colimit.  It expressly notes that these $M_i$ need not be
fs, although their associated log structures are fs and log regular.
The regularity implication for smooth log schemes over these bases
is also stated there, with its Kato Theorem~8.2 citation.
\citep[Lemma~5.3.2 and the following paragraph]{LA_AHLS2024}

The cofinal refinement in B.54 supplies the extra finite chart data
that this application requires.  For every nonzero node parameter
$a\in\mathfrak m_R$, choose $N$ with $\pi^N/a\in R$ and
include both $a$ and $\pi^N/a$ at a stage.  Then $a\in M_i$,
so the node equation has a log chart over the chosen base divisor,
not over an unrelated boundary.  Example~2.4.5, p.~15, proves that
\[
 P\longrightarrow P[e_1,e_2]/(e_1+e_2=a)
\]
is smooth, vertical, and saturated for every saturated $P$ and
$a\in P$, with saturated target.  This applies at every refined
stage.  Marked smooth charts use the additional $\mathbb N$ factor.
The saturated curve maps have ordinary underlying base changes;
the finite strict \'etale chart cover and its transition maps are
descended together.  This verifies the affine charted hypothesis
for the proper models.  Because the ring maps are local inclusions,
the maps on their residue fields are injective and their colimit is
$\kappa$.  Reduction of the same chart system thus verifies that
hypothesis for the log closed fibres as well.  Lepage's two
equivalences apply at the Noetherian stages.  AHLS continuity on the
two log terms and ordinary finite-\'etale continuity on the open
term give the asserted limiting equivalence.
\citep[Example~2.4.5, Definition~3.4.4,
Proposition~4.2.6]{LA_AHLS2024}

Proposition~5.3.1, pp.~65--67, identifies the log structure of a
smooth \emph{vertical} log scheme over the microbial valuation
base with the one induced by its generic fibre, under the stated
discrete, algebraically closed, or equicharacteristic-zero
hypotheses.  It applies first to the vertical structure in B.54;
the horizontal marked sections are then added by the separate
smooth charts.  A marked curve is not itself vertical.  This
distinction is respected by the proof under audit.
\citep[Proposition~5.3.1]{LA_AHLS2024}

Finally Proposition~4.4.3, pp.~60--62, makes the canonical
valuative log map associated with a finite tame extension of a
henselian valued field Kummer--\'etale.  Its underlying map is
integral, so it is a finite Kummer--\'etale cover in
Definition~4.2.1.  Corollary~4.4.9, p.~64, identifies the log
fundamental group of the canonical log valuation ring with the
maximal tame quotient of $G_K$.  In residue characteristic zero
the residue characteristic exponent is one, so all finite separable
extensions are tame in the definition immediately before
Proposition~4.4.3.  These statements therefore retain all of $G_K$
in B.54, without asserting module-finiteness of the integral closure.
\citep[Proposition~4.4.3, Corollary~4.4.9]{LA_AHLS2024}

\paragraph{Hassett and the non-Noetherian base.}
The author's version revised 14 March 2002 was opened at the
definition of the moduli problem and Theorem~2.1, printed
pp.~3--4.  Its initial definition uses Noetherian bases and weights
$0<a_i\leq1$ with $2g-2+\sum a_i>0$.  The theorem supplies a
smooth proper Deligne--Mumford moduli stack over $\mathbb Z$ with
projective coarse scheme and a universal curve.  For $g=0$ and all
weights one, with at least three marks, every stable labelled
geometric curve has trivial automorphisms: its labelled tree and
all special points are fixed, and each component has at least
three of them.  Thus this stack is the fine projective scheme.
Extension of its finite-presentation moduli interpretation to an
arbitrary base is the finite-presentation descent argument in B.54,
not an assertion that the source initially formulates its moduli
problem over arbitrary bases.  The ordinary valuative criterion
for that proper scheme gives the model over the given valuation
ring, and uniqueness gives its compatibility with valued base
extension.  No potential-stable-reduction theorem over a different
fraction field is substituted here.
\citep[Theorem~2.1]{LA_Hassett2003}

\paragraph{Outcome and boundaries of this check.}
The specified primary inputs support their stated uses in
B.43, B.47, B.49, B.54, and B.55.  Their descendants B.50--B.51
use these through the cited internal assertions.  No source-scope
gap was found in the all-prime arithmetic comparison.  The source
date correction for AHLS is bibliographic.  The fixed-log-point,
marked-chart, and coherence requirements described above are
essential parts of the application and should remain explicit in
the proofs.  No priority claim is made.

\section{Consequences and limits}
\subsection{The manuscript's statements and relative class numbers}

\begin{claim}{C.1}{Proved}{The two curve assertions in cote 149, scan 36}
Over every field finitely generated over its prime field, the
point-section map for a smooth hyperbolic curve is injective:
use the full geometric group in characteristic zero, and its
prime-to-$p$ quotient in characteristic $p$. For a proper curve
of genus at least two over a finite field, the covering condition
of scan 36 forces $P=Q$.
\end{claim}
\begin{proof}
In positive characteristic this is exactly
\statementref{A.1} and \statementref{A.33}. Their dependencies
are the same-prime compactification lemma, the cuspidal criterion
with its orientation and Kummer evaluation, the relative
configuration exactness, and the diagonal boundary transport;
each was independently proved in Part A. The covering condition
is equivalent to section conjugacy by \statementref{A.31}, and
the main theorem then gives equality. Characteristic zero is the
all-prime Kummer argument in \statementref{A.30} and
\statementref{A.33}.

Thus the sentence in scan 36, \emph{«en déduire $P=Q$»}, holds in
the specified curve scope. The phrase
\emph{«sans hyp. de caractéristique»} requires retaining the
finite-generation and hyperbolicity hypotheses. It does not make
the corresponding abelianized assertion true in positive
characteristic, and it does not include proper elliptic curves.
\end{proof}
\paragraph{Check record.} Also checked in this writing, 2026-11-18: GPT-6 Astra (OpenAI, ChatGPT Pro mode), root; separate effort setting not exposed. Read scope of field extensions, semiabelian Kummer collision versus full sections, independence limits, and virtual-puncture limitations. C.3 wording is corrected to avoid asserting formal independence not proved by its example.


\begin{claim}{C.2}{False as stated}{The characteristic-$p$ prime-to-$p$ extension of Théorème 3 in all dimensions}
The asserted characteristic-$p$ extension of cote 149, scan 32,
Théorème 3, for every geometrically connected finite-type scheme
embedded in a semiabelian variety, is false, already in dimension
one. Over $\mathbf F_2$, take
\[
 E:\quad y^2+xy=x^3+1,\qquad O=\infty,\quad T=(0,1).
\]
The two points have identical prime-to-two section classes in
the whole arithmetic extension. For an arbitrary semiabelian
$G/K$ over a finitely generated characteristic-$p$ field, the
exact replacement is
\[
 \ker\bigl(G(K)\to\hbox{prime-to-}p\hbox{ section classes}\bigr)
       =G(K)[p^\infty]
\]
after using the identity point to identify this section set with
abelian Galois cohomology. A subscheme embedded in $G$ has
injective point sections whenever no difference of distinct
rational points lies in $G(K)[p^\infty]$.
\end{claim}
\begin{proof}
The explicit checks in \statementref{A.24} verify that $E$ is
smooth and $T$ has exact order two. Its geometric fundamental
group is abelian, its prime-to-two part is the product of its
prime-to-two Tate modules, and the difference of the two point
sections is the Kummer class of $T$, which is zero. This is a
conjugacy assertion in the whole group, not a finite-level test.

For a semiabelian variety the same abelian fundamental-group
description and \statementref{A.30} give the exact kernel.
Functoriality of fundamental groups for an embedding $X\to G$
shows that a collision on $X$ maps to a point difference in this
kernel. Its avoidance is therefore a sufficient condition on
$X$. We do not assert the converse for $X$: hyperbolic curves
may still separate points nonabelianly when their generalized
Jacobian Kummer map fails to do so, as the negative control
\statementref{A.25} and the main theorem show. Products
$E\times\mathbf G_m^{d-1}$ give the same counterexample in every
positive dimension.

This matches the warning on scan 34 about translation by
$p$-torsion and the failure of the abelianized map (**) on scan
35. Removing a point from an elliptic curve can restore the
nonabelian hyperbolic assertion, as proved in Part A; it does
not invalidate the warning about (**).
\end{proof}
\paragraph{Check record.} Also checked in this writing, 2026-11-18: GPT-6 Astra (OpenAI, ChatGPT Pro mode), root; separate effort setting not exposed. Read scope of field extensions, semiabelian Kummer collision versus full sections, independence limits, and virtual-puncture limitations. C.3 wording is corrected to avoid asserting formal independence not proved by its example.


\paragraph{Readings and the limits of the transcript check.}
We checked the supplied second-pass transcriptions, not absent
images of the original scans. The alternatives
\emph{à l'aise / l'air} on scan 35 do not affect a hypothesis.
The margin's uncertain reading involving an \emph{anneau} is
mathematically consistent with spreading to a finite-type model
and specializing at a closed point, as in \statementref{A.32};
the proof does not depend on selecting that reading. The
alternative \emph{ou / en} in the preceding sentence also changes
no mathematical quantifier used here. In scan 36, a cover $Y'$
in the displayed covering condition is a cover of $Y$, as its
chosen point and generalized Jacobian require; the transcription's
$X$ is the noted variable slip. We interpret the torsion there
as $p$-primary, rather than as necessarily killed by $p$.
No statement in this paper depends on resolving a doubtful word.

\begin{claim}{C.3}{Proved}{The precise class-number consequences and a sharp limitation}
Point separation gives the isolating pointed cover of
\statementref{B.19}; it gives no numerical value for the
$p$-part of its relative class number. There is no universal
assertion that every nontrivial pointed prime-to-$p$ cover has
$p$-unit relative class number, or that every such cover has
$p$-divisible relative class number. Both behaviours occur for
the same curve, point, field, and prime degree in our computation.

For $C_\alpha/\mathbf F_{64}$ of \statementref{B.21},
$\alpha$ with bit code $58$, and $P=\infty$, the degree-seven
characters $(1,0)$ and $(1,1)$ in the recorded basis give
respectively
\[
 h(D_{(1,0)})/h(C)=46713473689,\qquad
 h(D_{(1,1)})/h(C)=55586549824.
\]
The first is odd and the second is even. Nevertheless the
point-section map of $C$ is injective over this field and every
constant extension.
\end{claim}
\begin{proof}
The exact occupancy vectors for these two characters are
$(9,12,17,12,12,17,12)$ and $(11,18,14,8,8,14,18)$.
The degree-two twisted polynomials, including leading coefficient
64 and the infinity pointing, are supplied in
\statementref{B.21}. Their full integer resultants give exactly
the two displayed integers; both are retained in the computed
JSON, and neither is a conjectural value. The main theorem and
its checked inputs apply to this same curve.

A hypothetical collision would force the relative divisibility
in the applicable class-number criterion. Its nonexistence
removes that premise; contraposition does not turn it into
existence of a cover with nondivisible relative class number.
Our explicit constructions prove such existence for the
fourteen requested $\mathbf F_4$ pointed cases and for the
supersingular test curve at primes 191 and 229, as shown in
Part B. Those are additional calculations, not formal
consequences of the point-separation theorem. The universal
quantifier in B1 therefore remains a separate question.
\end{proof}
\paragraph{Check record.} Also checked in this writing, 2026-11-18: GPT-6 Astra (OpenAI, ChatGPT Pro mode), root; separate effort setting not exposed. Read scope of field extensions, semiabelian Kummer collision versus full sections, independence limits, and virtual-puncture limitations. C.3 wording is corrected to avoid asserting formal independence not proved by its example.


\begin{claim}{C.4}{Proved}{What transfers to the other research line}
The characteristic-zero construction of a virtual puncture
does not prove the surjectivity of the section map for the
tripod over $\mathbf Q$, and does not construct a birational
lift of an arbitrary section. The present finite-field
injectivity theorem supplies no such assertion. Likewise a
non-point-theoretic full section over a finite field obtained
from a norm-module class in Part B is not a counterexample
over a field finitely generated over $\mathbf Q$.
\end{claim}
\begin{proof}
The pullback construction in the next subsection supplies an
extension for every section, without giving a rational point or
a splitting of its virtual cusp. Recognizing it as the punctured
group of an actual point is an additional condition. The
characteristic-zero separation criterion compares two already
geometric point sections; it has no arbitrary section as a
second input. The finite-field norm-module realization changes
both the characteristic and the arithmetic Galois group, so no
base-change map transfers it to $\mathbf Q$. These are exact
limitations of the maps constructed, not assertions that the
characteristic-zero surjectivity or birational-lifting questions
are false. They address the separate LM6 line referred to in
the supplied intake.
\end{proof}
\paragraph{Check record.} Also checked in this writing, 2026-11-18: GPT-6 Astra (OpenAI, ChatGPT Pro mode), root; separate effort setting not exposed. Read scope of field extensions, semiabelian Kummer collision versus full sections, independence limits, and virtual-puncture limitations. C.3 wording is corrected to avoid asserting formal independence not proved by its example.




\subsection{The precise separation consequences for the published theorems}

In this subsection $\Pi_Z^\Sigma$ always retains the entire arithmetic
quotient $G_k$. For a proper curve it is the geometrically pro-$\Sigma$
\'etale group; for an open curve it is the geometrically pro-$\Sigma$
tame group. The set $\Sigma$ may contain $p$, but must contain some
$r\ne p$. The relative comparison used in this scope is
\statementref{A.13}, not a claim that the full pro-$\Sigma$
configuration fibre sequence is exact.

\paragraph{The exceptional set in the source.}
Sa\"idi--Tamagawa 2017, Definitions~1.1, 1.2 and 1.4, first takes
the exceptional locus of the map from points of the universal
pro-cover to \emph{commensurability classes} of decomposition
subgroups, then takes its image $E_Z$ in the closed points of the
compactification. Its ``almost injective in the strong sense'' means
finite exceptional locus for a map of sets. In contrast,
Definition~2.9 calls a proper curve almost $\Sigma$-separated when
the residue fields of its exceptional points lie in a proper
algebraic subfield of $\bar k$. Proposition~2.12 assumes
$J_Z$-largeness and proves this latter support statement; it does
not assert that $E_Z$ is finite. The earlier cofinite result is
ST2009 Proposition~1.8(vi), with almost injectivity of the
closed-point decomposition map in Corollary~1.11(iii).
These are distinct assertions.\citep[Definitions~1.1, 1.2, 1.4,
2.6, 2.9; Proposition~2.12]{LA_ST2017}
\citep[Proposition~1.8(vi), Corollary~1.11]{LA_ST2009}

\begin{claim}{C.5}{Proved}{C1, C5}
Let $X/k_X$ and $Y/k_Y$ be proper smooth geometrically connected
curves of genus at least two over finite fields. Let $\Sigma_X,
\Sigma_Y$ each contain a prime different from the respective
characteristic, and suppose
\[
 \ker\!\left(G_{k_X}\longrightarrow
       \Aut\!\left(\prod_{s\notin\Sigma_X}T_sJ_X\right)\right)
 \ne1,
 \qquad
 \ker\!\left(G_{k_Y}\longrightarrow
       \Aut\!\left(\prod_{s\notin\Sigma_Y}T_sJ_Y\right)\right)
 \ne1.
\]
At the characteristic prime $T_p$ denotes the \'etale Tate module.
For every isomorphism $\alpha:\Pi_X^{\Sigma_X}\simeq
\Pi_Y^{\Sigma_Y}$, the two prime sets, characteristics and field
cardinalities coincide under the group-theoretic identifications,
and $\alpha$ preserves Frobenius. Moreover:
\begin{enumerate}
\item $E_X=E_Y=\varnothing$, and the same is true on every finite
connected \'etale covering arising from an open subgroup;
\item the decomposition-group construction gives a unique bijection
$\phi:X^{\cl}\simeq Y^{\cl}$, preserving residue-field degrees and
carrying $D_x$ to $D_{\phi(x)}$ up to conjugacy at every point;
\item every corresponding pair of finite marked sets $S,T$ may be
used in the cuspidally abelian reconstruction of ST2017
Corollary~3.3; the Green-compatibility assertion in Proposition~3.12
holds at all points and at every corresponding finite covering;
\item $\alpha$ is induced by the unique isomorphism of curves with
their corresponding pro-coverings in ST2017 Theorem~4.22.
\end{enumerate}
The fourth assertion was already the global conclusion of that
theorem. The improvement is the removal of the exceptional-point
restrictions in the preceding reconstruction statements.
\end{claim}
\begin{proof}
The displayed kernel conditions are exactly $J_X$-largeness and
$J_Y$-largeness in Definition~2.6. Proposition~1.7 recovers the
geometric subgroup and the prime sets. Lemma~2.7 gives largeness
relative to the constant fields, and Proposition~2.5 gives equality
of the characteristics and cardinalities and preservation of
Frobenius. Proposition~2.12 gives almost $\Sigma$-separation.
\statementref{A.23}, using the arbitrary-prime-set comparison
\statementref{A.13} and the retained boundary, upgrades it to
$E=\varnothing$. By ST2017 Proposition~2.11(i), an \'etale
$\Sigma$-cover has exceptional set equal to the inverse image of
the original exceptional set, so emptiness holds at every such
level, including constant extensions. No largeness hypothesis on
the new Jacobian is silently added.

Theorem~1.9 reconstructs the set of decomposition subgroups from
the Frobenius group. Emptiness makes the association with points
injective, hence gives $\phi$ on all closed points. The arithmetic
image of $D_x$ is $G_{k(x)}$, which proves the degree assertion.
Corollary~3.3 and Proposition~3.12 now have all their original
hypotheses with $X^{\cl}\setminus E_X=X^{\cl}$ and similarly
for $Y$ and for every covering. Substitution proves precisely the
marked-point and Green-compatibility assertions. Finally the
hypotheses of Theorem~4.22 are exactly those displayed here,
so its curve and pro-covering isomorphism applies.
\citep[Proposition~1.7, Theorem~1.9, Proposition~2.5,
Lemma~2.7, Propositions~2.11--2.12, Corollary~3.3,
Proposition~3.12, Theorem~4.22]{LA_ST2017}
\end{proof}
\paragraph{Check record.} Also checked in this writing, 2026-11-18: GPT-6 Astra Pro; Pro mode; separate effort setting not exposed; ANABELIAN\_A1\_SUCCESSOR/configuration. Also checked in this writing: Read complete carried proof and all numbered ST inputs. Add A35 for the repaired imported central/abelian/fixed-r reconstruction. A23 gives emptiness; A13 is explicitly invoked in the carried proof for its full Sigma scope. Etale-cover equality avoids imposing largeness on every new Jacobian. Published4.22 geometric conclusion is already known.


The $J$-largeness assumption includes every cofinite prime set and
some infinite-complement sets described in ST2017 Remark~2.8.
The present argument does not replace this condition by an arbitrary
nonempty $\Sigma$. In particular, the singleton-prime collisions
of \statementref{A.26} prevent such a replacement. Nor does
$E=\varnothing$ remove the omitted-prime torsion
$J(k)\{\Sigma'\}$ or the quotient by
$k^\times\{\Sigma'\}$ in the pseudo-function reconstruction
(Definition/Remark~4.1 and Theorem~4.2). The finite-index step in
the proof of Theorem~4.22 uses Remark~4.7's separate control of
complementary-prime torsion along a constant-field tower. Point
separation alone does not provide that control.

\begin{claim}{C.6}{Proved}{C1, C5}
Let $U,V$ be hyperbolic curves over finite fields, with an
isomorphism of geometrically pro-$\Sigma$ tame groups. Assume,
as in ST2017 Theorem~4.23, that there are corresponding open
subgroups defining covers $U',V'$ whose smooth compactifications
$X',Y'$ have genus at least two, with $\Sigma$ respectively
$J_{X'}$-large and $J_{Y'}$-large. Then $E_U=E_V=\varnothing$,
all closed-point decomposition labels are unique, and the unique
scheme isomorphism with its pro-coverings in Theorem~4.23 has
no exceptional-point restriction. The matched-open-subgroup
hypothesis is part of this assertion.
\end{claim}
\begin{proof}
By C.5, $E_{X'}=E_{Y'}=\varnothing$. If two interior points
of the universal pro-cover of $U$ had commensurate decomposition
groups, restrict to their common open arithmetic subgroup and to
the given open subgroup for $U'$. After a finite constant
extension they give distinct rational points of $U'$ with equal
section subgroups. They remain distinct on $X'$, and projection
to $\Pi_{X'}^\Sigma$ contradicts its all-extension separation.
Two distinct upstairs points above the same geometric point are
already excluded by ST2017 Proposition~1.3. Cusp decomposition
groups are recognized by their nontrivial inertia, and their
commensurability separation is the cusp assertion of ST2009
Proposition~1.8(i),(ii), recalled in ST2017 Theorem~1.8. Thus
no interior or cuspidal exception remains. The same argument
applies to $V$. Theorem~4.23 then applies with its original
matched-cover hypotheses. This proves both the new separation
clause and the already published geometric conclusion.
\citep[Proposition~1.3, Theorems~1.8 and 4.23]{LA_ST2017}
\end{proof}
\paragraph{Check record.} Also checked in this writing, 2026-11-18: GPT-6 Astra Pro; Pro mode; separate effort setting not exposed; ANABELIAN\_A1\_SUCCESSOR/configuration. Also checked in this writing: Read complete proof. Commensurate interior groups restrict to a common finite arithmetic level; Proposition1.3 excludes two universal-cover points above the same geometric point, and C5 excludes the remaining distinct compactified labels. Cusps are separated by the source inertia assertion. Matched-cover largeness hypotheses retained.


\begin{claim}{C.7}{Proved}{C2, C5}
Let $U,V$ be hyperbolic curves over finite fields, and let
$\Sigma_U^\dagger$ and $\Sigma_V^\dagger$ consist of every prime
different from their respective characteristics. Every isomorphism
of their geometrically pro-$\Sigma$ tame groups arises from a unique
isomorphism of curves and corresponding pro-coverings. All their
closed-point decomposition maps are injective, including after
every constant extension and every open-subgroup covering. For
proper $X,Y$, the restrictions to complements of exceptional sets
in ST2009 Corollary~2.3, Theorem~3.6 and Proposition~3.8 can be
replaced by all points.
\end{claim}
\begin{proof}
For $\Sigma=p'$, \statementref{A.1} gives proper point
separation after every finite extension and
\statementref{A.22} gives $E_X=\varnothing$. If $p$ is also
retained, a collision would survive in the prime-to-$p$ quotient,
so it is excluded as well. At every proper \'etale level the
curve still has genus at least two; the same proof applies there.
For a nonproper hyperbolic curve choose a corresponding finite
tame cover whose compactification has genus at least two, as in
the proof of ST2009 Corollary~3.10, and use the argument of
C.6. The cusp separation is unchanged.

The exact published Isom statements are ST2009 Theorem~3.9
for proper curves and Corollary~3.10 for arbitrary hyperbolic
curves. They permit either all primes or all primes except $p$.
Their hypotheses are met here. In the proper reconstruction,
Corollary~2.3 now applies to arbitrary finite $S,T$;
$\mathcal O_{E_X}=K_X$ and $\mathcal O_{E_Y}=K_Y$ in
Theorem~3.6; and Proposition~3.8 has no excluded points at any
finite covering. The proof of Theorem~3.9 is the formal
combination of Theorem~3.7 with Proposition~3.8, and the proof
of Corollary~3.10 descends compatible proper isomorphisms from
the cofinal family of such covers. These are the precise places
where the separation result removes restrictions.
\citep[Corollary~2.3, Theorems~3.6--3.7, Proposition~3.8,
Theorem~3.9, Corollary~3.10]{LA_ST2009}

This simplification does not set $J_X(k_X)[p^\infty]$ equal
to zero. The $p$-power correction to functions in Theorem~3.6(ii)
therefore remains necessary. Similarly, the reconstruction of
additive structure is still a separate part of the proof; an
injective point-label map alone is not an Isom theorem.
\end{proof}
\paragraph{Check record.} Also checked in this writing, 2026-11-18: GPT-6 Astra Pro; Pro mode; separate effort setting not exposed; ANABELIAN\_A1\_SUCCESSOR/configuration. Also checked in this writing: Read complete proof; directly uses A1 point separation, A22 equivalence with exceptional-set emptiness, C6 affine descent and A35 imported quotient repairs. Read journal3.6, 3.8, 3.9 and3.10 proof loci. The p-primary correction in reconstruction of functions remains. No full-group Isom theorem is applied to only a relative-r group.


\paragraph{What becomes of the proposed configuration route.}
ST2009 arXiv Remark~3.13, journal Remark~3.14, would first obtain
an isomorphism of the \emph{full} prime-to-$p$ punctured groups
from the proposed full fibre exactness, then apply the published
Isom theorem to infer equality of the punctures. The full
exactness premise is false by Part~A. Replacing the group by its
cuspidally pro-$r$ quotient does not permit that last application
of the full-group Isom theorem. What does work is the distinct
chain
\[
 \text{relative exactness and diagonal transport}
 \Longrightarrow
 \text{the cuspidal criterion of Part A}
 \Longrightarrow E_X=\varnothing.
\]
This chain is independent of the published Isom conclusion.
Combining it with the explicitly retained cuspidally abelian
and fixed-$r$ reconstruction inputs gives C.7 by the source's
function-reconstruction argument. It is a corrected route to
point separation followed by the known Isom proof, rather than
a restoration of the false full fibre premise.

\begin{claim}{C.8}{Proved}{C5}
For proper hyperbolic curves over finite fields and
$\Sigma^\dagger=p'$, the $\Sigma$-separation assumption in
Mochizuki's original Definition~1.18, Proposition~2.6(ii),(iii)
and Corollary~3.11 is automatic. Thus a Frobenius-preserving
isomorphism is point-theoretic, arbitrary finite corresponding
marked sets are allowed in Theorem~2.5(i), the indicated proper
and cuspidally abelian cusp decomposition groups are
commensurably terminal, and the isomorphism is totally globally
Green-compatible. The quotient repairs of A.35
are used at the fibre-identification inputs.
\end{claim}
\begin{proof}
Definition~1.18(i) asks for closed-point decomposition
injectivity on every connected finite \'etale covering arising
from an open subgroup. C.7 proves exactly this requirement.
Decomposition-group recognition and this injectivity change
quasi-point-theoreticity into point-theoreticity under
Definition~1.18(ii), as in Remark~1.18.2(ii). Theorem~2.5(i)
assumes point-theoreticity and gives the marked-set assertion.
Proposition~2.6(ii),(iii) now have their separation assumption
satisfied and give the two commensurable terminality assertions.
Finally Corollary~3.11 assumes finite fields,
$\Sigma^\dagger=p'$ and separation; all three conditions hold.
Its fixed-prime fibre uses are in the repaired scope of A.35.
These are applications of the stated source results with their
actual assumptions, not an assertion that every conclusion of
the entire paper has been reverified.
\citep[Definition~1.18, Remark~1.18.2, Theorem~2.5(i),
Proposition~2.6(ii),(iii), Corollary~3.11]{LA_Mochizuki2007}
\end{proof}
\paragraph{Check record.} Also checked in this writing, 2026-11-18: GPT-6 Astra Pro; Pro mode; separate effort setting not exposed; ANABELIAN\_A1\_SUCCESSOR/configuration. Also checked in this writing: Read complete proof and all listed original Mochizuki statements. All-cover separation in C7 exactly discharges Definition1.18. A35 repairs the specified fibre inputs. All-prime hypotheses in function reconstruction and all other source assumptions remain.


For reference, the original/LaTeX numbering pairs in the last
claim are: Definition~1.18/\allowbreak Definition~1.5;
Theorem~2.5/\allowbreak Theorem~2.1; Proposition~2.6/\allowbreak Proposition~2.3;
Corollary~3.11/\allowbreak Corollary~3.1. The source's
Theorem~2.5(ii) assumes $\Sigma$ is \emph{all} primes in order
to recover principal divisors exactly. Separation does not remove
that hypothesis. Its original Theorem~3.10 (LaTeX Theorem~3.1)
already permits arbitrary matching rational points; the new
fact makes the matching point unique, rather than supplying a
previously absent arbitrary-point existence assertion.

\paragraph{Bounds on the consequence claim.}
Wakabayashi's reconstruction theorems use a single invertible
prime. Their pro-$l$ point labels cannot be declared injective
from a theorem about the larger prime-to-$p$ group; Part~A gives
actual singleton-prime collisions. The virtual-curve construction
in Sun 2025 already applies to arbitrary sections, and Sun 2026's
specified open-injection criterion has its own geometricity
hypotheses. Point separation neither constructs a nongeometric
section nor realizes its virtual extension as a punctured
algebraic curve. Finally no arbitrary-homomorphism anabelian
theorem follows merely by replacing exceptional sets with the
empty set in the checked Isom statements. These limits are part
of the scope of C1, C2 and C5.

\subsection{The separation-hypothesis census in the remaining named papers}

\begin{claim}{C.9}{Proved}{Exact removal of the named separation conditions}
In the original Mochizuki 2007 text, the explicit separation or
point-theoreticity conditions relevant to C5 are the loci in the
following table. For finite fields with $\Sigma^\dagger=p'$,
point-theoreticity and separation at these loci follow from
\statementref{A.1} and its all-cover corollaries, as applied in
C.7--C.8. Other hypotheses in the same statements remain.
Wakabayashi 2016 and the two Sun preprints have no additional
unproved finite-field prime-to-$p$ separation condition whose
removal strengthens their stated theorems. The distinctions in
the table are part of the conclusion.
\end{claim}
\begin{center}\small
\begin{longtable}{p{0.25\linewidth}p{0.33\linewidth}p{0.34\linewidth}}
\toprule
Locus&Condition or use&Consequence and retained restriction\\*\midrule
\endhead
Mochizuki Definition~1.18; Remark~1.18.2(i)--(iii)&
Separation on every finite open-subgroup cover; passage from
quasi-point-theoretic to point-theoretic; possible failure of
maximality or terminality of point decomposition groups.&
Separation and unique point labels hold at every such level.
For a Frobenius-preserving isomorphism the matching point is
unique. The terminality conclusion is Proposition~2.6 below.\\
Proposition~2.3(ii)&
Proves separation when $\Sigma$ is all primes.&
Its finite-field point-separation conclusion now also holds for
$\Sigma=p'$. Its assertion about detecting principal divisors
in (i) is a different statement and is not extended by this
substitution.\\
Definition~2.4; Theorem~2.5(i)&
Point-theoreticity supplies matching point sets used in Green
compatibility and marked cuspidally abelian reconstruction.&
All finite corresponding sets are allowed, functorially, with
the same cusp and proper-quotient data. The auxiliary definition
of Green compatibility is unchanged.\\
Theorem~2.5(ii)&
Adds $\Sigma=\mathrm{Primes}$ to reconstruct principal
divisors and the stated push-forward of multiplicative groups.&
Point-theoreticity is automatic in the finite-field setting,
but the all-prime hypothesis is retained. Separation alone does
not remove complementary-primary torsion.\\
Proposition~2.6(ii),(iii)&
Assumes $\Sigma$-separation for commensurable terminality in
the proper group and in the cuspidally abelian punctured group.&
Both terminality conclusions hold without an additional
separation assumption in the finite-field prime-to-$p$ cases.\\
Remark~2.6.1; Corollary~2.7&
The remark compares an already given punctured-group
isomorphism with the reconstructed abelian one. The corollary
assumes all primes, point-theoreticity, and total local principal
Green compatibility.&
The point-label condition is automatic as above. Existence of
the punctured-group isomorphism in the remark is not deduced
from separation. In the finite-field all-prime corollary,
Corollary~3.11 supplies Green compatibility as in C.8;
the Isom conclusion was already part of the source's program.\\
Corollary~2.8; Proposition~2.10(iii)&
The former is a point-theoretic affine Isom theorem with all
primes; the latter supplies point-theoreticity for quasi-Belyi
curves over nonarchimedean local fields.&
The finite-field all-prime premise of the former is automatic;
the prime-to-$p$ version is obtained by the separately checked
ST2009 argument, C.7. The local-field result has no new
finite-field consequence from R1.\\
Theorem~3.10; Corollary~3.11&
The theorem assumes matching rational point decomposition
groups, not separation. The corollary assumes finite fields,
$\Sigma^\dagger=p'$, and separation.&
The matching point in the theorem becomes unique; arbitrary
matching-point existence was already present. The corollary's
separation condition disappears and total global Green
compatibility follows.\\
Wakabayashi Theorems~3.9,4.4; Corollary~3.10 and
Remark~3.10.1&
One invertible prime $l$; matching decomposition groups in the
punctured theorem; the remark explicitly allows several
matching closed points.&
The source does not assume a mixed-prime separation condition.
R1 cannot make its single-prime point correspondence unique;
the collisions in Part~A prohibit that inference.\\
Sun 2025, Definition~3.1, Proposition~3.7, Corollary~3.8,
Theorems~A--C&
Virtual pullbacks are defined for arbitrary sections in
characteristic zero. Geometric sections identify their actual
fibres. Theorems~A and C use number fields or mixed-characteristic
local fields. Theorem~B allows any characteristic-zero field, with
its stated nonempty-boundary and cyclotomic hypotheses.&
There is no conditional finite-field separation clause to
remove. Point separation does not prove that a virtual curve
has an algebraic realization or that its section is geometric.\\
Sun 2026, Lemma~5.6; Theorem~5.7 and Corollary~5.8&
Lemma~5.6(i) uses the already known injectivity of points into
full sections over an arithmetic field; (ii) uses an additional
open-immersion geometricity theorem. The later results impose
specific virtual covering and cusp conditions.&
No such condition is supplied merely by R1. The source's
arithmetic fields are number fields or mixed-characteristic
local fields, and its groups are full \'etale groups.
None of these assertions becomes a finite-field prime-to-$p$
section-surjectivity theorem.\\
\bottomrule
\end{longtable}
\end{center}
\begin{proof}
We searched the complete opened texts for the source's
separation, point-theoreticity, point-injectivity and geometricity
terms, then read the surrounding numbered statements and their
proof uses. In Mochizuki's original text, the explicit
$\Sigma$-separation occurrences are Definition~1.18,
Remark~1.18.2, Proposition~2.3(ii), Proposition~2.6(ii),(iii)
and Corollary~3.11. The point-theoreticity occurrences add
Definition~2.4, Theorem~2.5, Remark~2.6.1, Corollaries~2.7--2.8
and Proposition~2.10(iii). Theorem~3.10 requires corresponding
point groups, but not separation. This produces the displayed
finite list; no unexamined whole-paper assertion is substituted
for it.

C.7 proves separation on every required cover, C.8 proves
the resulting point-theoreticity and terminality, and A.35
and A.41 supply the repaired quotient fibre inputs. Substituting
these established assertions into the displayed hypotheses gives
the consequences in the first eight rows. In particular, the
proof of Theorem~2.5(i) restricts the reconstructed cuspidally
abelian configuration extension to the matched point groups;
the repaired abelian fibre is exactly what this restriction
requires. The all-prime requirement in (ii) is separate and
has not been removed. The proof of Corollary~3.11 uses the
fixed-$l$ reconstruction for each $l\in\Sigma^\dagger$
and recognizes noncuspidal decomposition groups there. A.41
repairs that input, while all-cover separation makes each of
the matching labels unique. This is precisely total global
Green compatibility.

Wakabayashi's Remark~3.10.1 expressly distinguishes existence
of a matching point from uniqueness. Its surrounding theorems
remain in the single-prime category described in A.37, so
the larger prime-to-$p$ result does not alter their hypotheses.
Sun 2025 constructs the virtual pullback before making any
assumption that its section is a point. Sun 2026 defines an
arithmetic field to mean a number field or a mixed-characteristic
local field; its Lemma~5.6(i) invokes the injectivity portion
of the section conjecture in that already known full-group
setting. Its part (ii) and later geometricity criteria have
additional inputs, not a missing point-separation premise.
These readings prove the last three rows and complete the
stated separation-condition census. They also explain why no
new arbitrary-homomorphism theorem or section-surjectivity
theorem has been inferred here.
\citep[Definition~1.18, \S2, Theorem~3.10,
Corollary~3.11]{LA_Mochizuki2007}
\citep[Theorems~3.9,4.4, Remark~3.10.1]{LA_Wakabayashi2016}
\citep[Definition~3.1, Proposition~3.7, Corollary~3.8,
Theorems~A--C]{LA_Sun2025}
\citep[\S2, Lemma~5.6, Theorem~5.7, Corollary~5.8]{LA_Sun2026}
\end{proof}
\paragraph{Check record.} Also checked in this writing, 2026-11-18: GPT-6 Astra Pro; Pro mode; separate effort setting not exposed; ANABELIAN\_A1\_SUCCESSOR/configuration. Also checked in this writing: Read complete statement, table and proof; searched full opened source texts for separation/point-theoreticity terms and read the listed statements. Corrected the Sun2025 row to distinguish TheoremB characteristic-zero scope from A/C arithmetic-field scope. A26 is the explicit singleton-collision comparison; A37 and A41 supply the source-scope and repair census.



\subsection{Virtual punctures and the characteristic-zero fibre construction}

\begin{claim}{C.10}{Proved}{C4}
Let $k$ have characteristic zero, let $C/k$ be hyperbolic, and consider
the projection $C^{(2)}=(C\times_k C)\setminus\Delta_C\to C$.
Work in a specified quotient category in which there is an exact
sequence
\[
 1\longrightarrow F\longrightarrow E\xrightarrow{q}\Pi_C
 \longrightarrow1.
\]
For every section $s:G_k\to\Pi_C$ the profinite pullback
\[
 \Pi_s=E\times_{\Pi_C,s}G_k
\]
is an extension $1\to F\to\Pi_s\to G_k\to1$. If $s=s_P$ is the
section of $P\in C(k)$, the natural fibre map identifies $\Pi_s$ with
the corresponding quotient of $\Pi_{C\setminus P}$, whenever the
assumed exact fibration statement includes that fibre identification.
\end{claim}
\begin{proof}
The kernel of the pullback projection is $F\times\{1\}$. Surjectivity
follows from that of $q$, proving exactness. For a geometric section
the map from the arithmetic group of the actual fibre has the specified
geometric fibre map and is the identity on $G_k$. If those geometric
groups agree, the extension isomorphism criterion proves that the
arithmetic groups agree. These arguments are valid for profinite groups
and continuous maps, without requiring the section to come from a point.
\end{proof}
\paragraph{Check record.} Also checked in this writing, 2026-11-18: GPT-6 Astra Pro; Pro mode; separate effort setting not exposed; ANABELIAN\_A1\_SUCCESSOR/configuration. Also checked in this writing: Read complete proof; exactness follows by profinite pullback. Actual-fibre identification is conditional on the specified exact fibre comparison. The construction alone supplies no section of the new extension.


\paragraph{The quotient qualification is necessary.}
An exact sequence in a quotient cannot recover information already
discarded by that quotient. For example, killing the entire diagonal
inertia gives the product group with its exact projection sequence;
its fibre is the unpunctured geometric group. Consequently the phrase
``the punctured group in every exact quotient'' means the corresponding
punctured quotient and its retained boundary data. It cannot literally
mean the full punctured fundamental group in a quotient that kills
the new cusp. This is a limitation of that wording, not a failure of
the pullback construction.

\begin{claim}{C.11}{Proved}{C4}
Assume the chosen configuration quotient also retains a diagonal
boundary decomposition group $D_\Delta$ mapping surjectively to
$\Pi_C$, with kernel the surviving inertia $I_\Delta$. Then
\[
 D_s=D_\Delta\times_{\Pi_C,s}G_k\subset\Pi_s
\]
is the distinguished virtual cusp extension. The pair $(\Pi_s,D_s)$
depends only on the geometric conjugacy class of $s$, up to the
corresponding inner ambiguity. For $s_P$ it is the actual punctured
fibre group with its new cusp in the chosen quotient.
\end{claim}
\begin{proof}
Suppose $s'=\Inn(\gamma)s$ for $\gamma\in\Delta_C$. The diagonal
decomposition map has a geometric lift $d$ of $\gamma$; conjugation
by $d$ identifies the inverse images of the two section subgroups and
their intersections with $D_\Delta$. Thus it transports both the
extension and the cusp. If the second proper projection is retained,
it changes that projection by $\Inn(\gamma)$, which may be removed
by a geometric inner automorphism in the target fibre because its
map to the other proper factor is surjective. This leaves the cusp
conjugacy class unchanged. For $s_P$, a local parameter at $P$ gives
the ordinary boundary decomposition description, and the exact fibre
comparison gives the asserted identification.
\end{proof}
\paragraph{Check record.} Also checked in this writing, 2026-11-18: GPT-6 Astra Pro; Pro mode; separate effort setting not exposed; ANABELIAN\_A1\_SUCCESSOR/configuration. Also checked in this writing: Read complete proof. A geometric lift in retained diagonal boundary gives simultaneous transport of extension and cusp; optional second proper projection correction uses geometric fibre surjectivity. Boundary lift and actual cusp comparison are additional hypotheses, not consequences of bare exactness.


The diagonal lift is extra geometric information. Exactness of
$E\to\Pi_C$ alone does not specify which lift preserves a cusp.
Also $\Pi_s$ is an extension, not automatically a split extension;
the construction has not supplied a section of $\Pi_s\to G_k$.
In particular the pullback of the diagonal cusp extension along a
nongeometric section is a virtual boundary datum, whose realizability
as a rational puncture has not been proved by writing the fibre product.

\begin{claim}{C.12}{Proved}{A14, C4}
The construction in C.10 is already Sun's virtual fundamental group.
The two opened preprints do not assert that every virtual puncture is
the puncture at a rational point.
\end{claim}
\begin{proof}
Sun 2025 Definition~3.1 starts with a smooth morphism of connected separated normal
finite-type schemes over a characteristic-zero field, inducing a
surjection on fundamental groups, and a Galois section of the base,
and defines
the virtual group by this pullback. Corollary~3.2 records its extension
and the independence of the geometric kernel from the section.
For a hyperbolic-curve family, Proposition~3.7 gives the fibre comparison,
and Corollary~3.8 identifies the virtual group with the actual fibre
when the section is geometric. This is a direct predecessor, rather
than a new construction of this paper.

Sun 2026 develops additional reconstruction and geometricity criteria.
Its introduction's Theorem~D corresponds to Corollary~5.8 in the body
(the introduction calls the reference a lemma). It concerns a number
field or a mixed-characteristic local field, virtual curves arising
from curves of types $(0,3)$ and $(1,3)$, and an open injection of their
virtual groups respecting their geometric kernels. Under those
hypotheses it forces the geometric index to be two and both sections
to be geometric. These restrictive hypotheses are not available for
an arbitrary section over a finitely generated extension of $\Q$.
\citep[Definition~3.1, Corollary~3.2, Proposition~3.7 and
Corollary~3.8]{LA_Sun2025}
\citep[Theorem~D and Corollary~5.8]{LA_Sun2026}
\end{proof}
\paragraph{Check record.} Also checked in this writing, 2026-11-18: GPT-6 Astra Pro; Pro mode; separate effort setting not exposed; ANABELIAN\_A1\_SUCCESSOR/configuration. Also checked in this writing: Read complete carried proof and the cited Sun2025 definitions/proofs and Sun2026 Lemma5.6/Theorem5.7/Corollary5.8. Corrected smooth surjective morphism to smooth morphism inducing a surjection on pi1. Sun2026 introduction mislabels bodyCorollary5.8 as a lemma; stated type/open-injection hypotheses retained.


\begin{claim}{C.13}{Proved}{C4}
Suppose $k$ is finitely generated over $\Q$, $C/k$ is hyperbolic,
and the base quotient retains the full arithmetic fundamental group,
or at least the point-separating Kummer maps used below. A cusp-preserving
isomorphism over the identity of this base quotient between the
punctures at $P,Q\in C(k)$ forces $P=Q$.

More generally, in the full base quotient, if the virtual puncture
$(\Pi_s,D_s)$ is isomorphic over $\Pi_C$, preserving the distinguished
cusp, to the actual puncture at $P$, then $s$ and $s_P$ are geometrically
conjugate. Conversely a geometric section has that puncture by C.10.
\end{claim}
\begin{proof}
Projecting the cusp decomposition groups to $\Pi_C$ identifies the
point sections up to geometric conjugacy. If the smooth proper
compactification $\bar C$ has positive genus, abelianization gives
zero Kummer class for $[P-Q]\in J_{\bar C}(k)$. The
Mordell--Weil--Lang--N\'eron finite-generation theorem applies to
$J_{\bar C}(k)$. For a finitely generated
abelian group $A$, the map to all its finite-quotient completions is
injective: $\bigcap_{n\ge1}nA=0$. The Kummer exact sequences therefore
make the all-prime Kummer map injective on $J_{\bar C}(k)$. Hence $[P-Q]=0$.
A nonconstant function with divisor $P-Q$ would have degree one and
identify $\bar C$ with $\mathbf P^1$, impossible for positive genus; thus
$P=Q$.

If $\bar C$ has genus zero, extend $k$ finitely so that two boundary
points $a,b$ are rational. The existence of $P$ identifies $\bar C$
with $\mathbf P^1$. A degree-one coordinate with zero at $a$ and pole
at $b$ gives an injective map $C\to\mathbf G_m$ on rational points.
The full Kummer map on $k^\times$ is injective. To see this, let $K$
be the algebraic closure of $\Q$ in $k$, a number field. If an element
of $k^\times$ has an $n$th root in $k$ for every $n$, all its divisorial
valuations on a normal projective $K$-model are zero, so it lies in
$K^\times$. Its roots also lie in $K$, because they are algebraic over
$K$. Their valuations at every finite prime are zero; the finitely
generated unit group of $K$ has trivial intersection of all its
$n$th-power subgroups. Thus the element is one. Applying this to the
ratio of the coordinate values at $P,Q$ again proves $P=Q$.

For the virtual statement, the image of its distinguished cusp in
$\Pi_C$ is $s(G_k)$ by definition. An over-$\Pi_C$, cusp-preserving
comparison with the actual puncture makes this subgroup conjugate
to $s_P(G_k)$. Their projections to $G_k$ identify the corresponding
sections. Conversely C.10 and C.11 give the comparison when $s$
is geometric.
\end{proof}
\paragraph{Check record.} Also checked in this writing, 2026-11-18: GPT-6 Astra Pro; Pro mode; separate effort setting not exposed; ANABELIAN\_A1\_SUCCESSOR/configuration. Also checked in this writing: Read complete proof. Positive-genus recognition uses the all-prime Jacobian Kummer injection; genus-zero recognition uses a split degree-one coordinate and the valuation/unit proof that the intersection of all powers is trivial. Virtual converse is exactly projection of its retained cusp group; no existence or geometricity is presumed.


This gives a recognition equivalence for the open characteristic-zero
question; it neither constructs a nongeometric section nor proves
that none exists. The finite-field proof of AG.40 used finiteness of
$J_C(k)$ to obtain a principal multiple of $P-Q$, constant-field point
counts, and the weight calculation eliminating nonconstant proper
Kummer classes. Those steps cannot simply be reused over $k$.
In a singleton-$l$ quotient, even the abelian Kummer map on $J_C(k)$
has the prime-to-$l$ torsion subgroup in its kernel. Thus exactness of
a chosen configuration quotient is not by itself a proof of point
separation for its base quotient. The full-group case, including
affine curves, is covered by the genus-zero and positive-genus
arguments above.\citep[\S1.2]{LA_CiperianiStix2015}


\part{Arguments whose verification or completion remains open}
\section*{Part A: retained statements and produced arguments}
\begin{claim}{A.28}{Open}{The remaining arbitrary-prime-set classification}

The assertions above prove emptiness for $J_C$-large sets (once
the stated transport theorem is inserted), infiniteness for
singletons, and a finite-or-infinite dichotomy. They do not decide
the alternative for every arbitrary nonempty prime set that is
neither a singleton nor $J_C$-large. That additional classification
asked in A12 remains unproved in this audit.

\end{claim}

\paragraph{State of the unresolved assertion.}

The remaining obstruction is exact: the Kummer comparison confines a
possible exceptional pair to a fibre over complementary-prime torsion,
but when the complementary Tate representation is faithful its kernel
gives no proper field supporting those pairs. The single-prime central
arithmetic subgroup proof cannot be applied to a multi-prime completion:
its automorphism group need not have a pro-\(\Sigma\) kernel on a finite
Frattini quotient. Neither assertion is silently assumed. The next
approach would be to classify the kernel of the full outer Frobenius
action for these prime sets, or to establish sparse support of the
possible point-pair differences along a constant-field tower.




\paragraph{Current proved reductions and remaining scope.} A.53--A.61 add the all-cover Frobenius-order criterion, its retained-characteristic-prime extension, precise limitations of finite-prime and bounded-generator methods, an explicit cofinal alternating homology tower over the fixed finite field, an actual first arithmetic order jump, and the precise orbit/multiplicity thresholds for a second factor. The remaining universal task is to determine whether the omitted-prime orders on this tower are unbounded and, when the action is faithful, to decide point separation.

\paragraph{Reason for placement in Part II.} The stated universal assertion remains unproved; the accompanying text records its exact current gate.
\section*{Part B: retained statements and produced arguments}
\begin{claim}{B.5}{Open}{The remaining prescribed-point existence assertion}
For every prescribed $(C,P)$ as in B1, prove or refute the
existence of a prime $\ell$ for which at least one first-field
test in B.4 succeeds. This remains undecided here, including
the universal assertion for pairs $(p,g)$ with
$p>(2g-2)^2$. B.3 gives the corrected numerical sufficient
conditions; B.28--B.29 prove unpointed existence and quantify
the prescribed-point obstruction. The R8 example has the
explicit least-prime witness proved in B.26--B.27, and B.30
gives a plane-quartic instance. None of these removes the
arbitrary-curve and prescribed-point quantifiers.
\end{claim}



\paragraph{Current proved reductions and remaining scope.} B.89 proves that any prescribed finite set of admissible odd primes can fail permanently for suitable pointed curves of each fixed genus and characteristic. Thus no witness-prime bound depending only on these two integers can settle the unrestricted question, including its displayed large-characteristic range. B.141 removes the entire genus-two marked-curve locus with a nonzero canonical-point class modulo some admissible odd prime of even p-order; its explicit congruence family includes characteristics above that threshold. B.145 proves the broader automorphism-eigencharacter criterion in every genus, including a genus-three quartic family with all constant-extension tests.

\paragraph{Reason for placement in Part II.} The stated universal assertion remains unproved; the accompanying text records its exact current gate.

\begin{claim}{B.12}{Open}{The universal free-orbit assertion of B2}
It remains undetermined here whether $M_P^F$ has no nonzero
all-level free vector for every genus-two curve in characteristic
two over every finite field of definition, and whether an
all-level free vector exists for any proper hyperbolic curve in
any characteristic. The general fixed classes and their full
section realizations have been determined above, but these results
do not decide nonemptiness of the inverse limit in
\statementref{B.8} for all such curves.
\end{claim}
\paragraph{The exact remaining step.}
To prove the universal nonexistence statement one must produce,
for each pointed curve, a stable normal level $H$ with
$S_H=\varnothing$. A witness as in \statementref{B.9} suffices,
but point separation supplies no such witness: it excludes classes
represented compatibly by point differences, whereas the
surjection in \statementref{B.7} produces many sections which
are not point sections. To refute the statement one must instead
prove $S_H\ne\varnothing$ for every stable normal level of one
fully specified pointed curve, not merely compute regular vectors
at a finite list of levels. Compactness would then give a genuine
all-level counterexample. Either sort of all-level free vector,
if it exists, automatically has a non-point-theoretic section
realization by \statementref{B.7}; none can come from a pair of
distinct points after the main point-separation theorem.



\paragraph{Current proved reductions and remaining scope.} B.130 removes two complete nonempty geometric genus-two families from this gate: the ordinary cubic locus of B.85 and the two-rank-one quintic locus of B.86. Every finite-field form, finite extension and rational pointing in those families has no all-level free fixed norm vector. The other genus-two loci and the general all-level existence assertion remain undecided.

\paragraph{Reason for placement in Part II.} The stated universal assertion remains unproved; the accompanying text records its exact current gate.

\begin{claim}{B.24}{Open}{The literal degree-one-permitted primitive-cover assertion}
For every proper curve of genus at least two over
$\overline{\mathbf F}_p$, does there exist a connected cyclic
\emph{\'etale} cover of odd degree prime to $p$, allowing degree one,
whose primitive part has $p$-rank zero, with the convention
$P_1=J_C$?
\end{claim}
The preceding counterexample does not settle this statement.
The identity cover proves it for every curve of $p$-rank zero,
including that example. For positive $p$-rank, degree one fails.
The exact Frobenius-orbit formulation above extends to arbitrary
genus with character spaces of dimension $g-1$: one needs
nilpotence of each orbit composite, not just a noninvertible
arrow. The conductor bound proved here uses scalar Frobenius
and the one-dimensional genus-two character spaces, so it does
not bound the positive-$p$-rank cases. We tried transferring the
scalar obstruction, using a single nonordinary character, and
using the identity convention; each fails for this stated reason.
The smallest unresolved scope in this work is positive-two-rank
genus-two curves in characteristic two. No such curve is claimed
to be a counterexample.




\paragraph{Current proved reductions and remaining scope.} B.85--B.86 prove primitive-zero covers on the two stated genus-two loci. B.87 shows that no fixed finite list of odd cyclic degrees suffices for all ordinary genus-two curves. It does not provide a single curve avoiding every odd degree, and the degree-one-allowed universal assertion remains open.

\paragraph{Reason for placement in Part II.} The stated universal assertion remains unproved; the accompanying text records its exact current gate.

\begin{claim}{B.32}{Open}{The unbounded quartic classification in B6}
The least witness over arbitrary constant extensions for
\emph{every} smooth plane quartic over $\mathbf F_3$ and
$\mathbf F_5$ remains undecided here. B.30 is an exact
least-prime and first-field result for one pointed quartic;
B.31 excludes four actual covers at all later fields.
The independent program \path{code/pointed_quartic_tests.py}
also computes the point counts and Weil polynomials of
$X^4+Y^4+Z^4=0$ over $\mathbf F_3$, at $(1:1:1)$,
and $X^4+Y^4-Z^4=0$ over $\mathbf F_5$, at $(1:0:1)$.
Their base-field Jacobian orders are $64$ and $256$,
respectively, so they have no base-field odd-prime Lang
quotient. This is not an all-extension exclusion.

The previous return's large base-field census programs were
not attached and have not been rerun here. We do not treat
their numerical tables as newly verified certificates.
B.4 reduces each fixed-prime all-extension question to a
finite calculation, but gives no terminating bound on the
prime search. It also remains undecided whether some curve
has a nonabelian witness of the stated type and has no
class-number witness over any extension. Failure of finitely
many characters, primes, or fields cannot prove that assertion.
\end{claim}



\paragraph{Current proved reductions and remaining scope.} B.145 settles the quartic $Y^3Z-YZ^3=X^4$ over $\mathbf F_3$ at its earliest witness field, for all four original rational markings: the exact field-first minimum is degree three and prime seven. The seven displayed covers have exact class-number unit tests through every constant extension. This does not determine least primes over arbitrary fields, all-cover minima at each extension, or the full quartic census.

\paragraph{Reason for placement in Part II.} The stated universal assertion remains unproved; the accompanying text records its exact current gate.

\begin{claim}{B.52}{Open}{Exact remaining scope of the henselian-field question B8}
The classification demanded by B8 remains open for general
finite-transcendence-degree fields over $\mathbb Q$ with a
henselian mixed-characteristic valuation outside the following
proved positive and negative classes. All positive assertions
below concern full closed-point-preserving relative etale
reconstruction of finite-type schemes, with the structural
isomorphism retained and universal homeomorphisms inverted.
They include every finite extension of the stated field.

For every prime $p$, let
$K_p=\operatorname{Frac}(\mathbb Z_{(p)}^{sh})$ and $I=G_{K_p}$.
The fields
\[
 B_p=\overline{\mathbb Q}^{I''},\qquad
 K_{p,\mathrm{nil}}=
   \overline{\mathbb Q}^{\,\bigcap_{c\ge1}\gamma_{c+1}(I)}
\]
are positive by B.136 and B.138, respectively, including $p=2,3$.
Every fixed bounded-nilpotency field is also covered by B.123.

More generally, let $D\subset\overline{\mathbb Q}_p$ be
finitely generated over $\mathbb Q$, let $M$ be its finite
local completion, and let $\Omega$ be its algebraic closure
inside $\overline{\mathbb Q}_p$.
B.139--B.140 prove the corresponding finite-derived and
lower-central-union assertions in every finite transcendence
degree. B.142--B.143 prove positivity for
$\Omega^N$ for every nontrivial closed normal subgroup
$N\triangleleft G_M$, and for every algebraic subfield
over which such a positive field is algebraic.
A nontrivial normal core of the field's absolute Galois
subgroup in some open local Galois group is consequently
a sufficient criterion, with the relative local base
chosen as in B.139. No normal closure is substituted
for that core.

There is a further positive class with no finite-residue
restriction. Let $D/\mathbb Q$ be finitely generated with
a discrete mixed-characteristic valuation whose residue
field is finitely generated over $\mathbb F_p$.
Let $M$ be its completion, let $\Omega$ be the algebraic
closure of $D$ in $\overline M$, and let $I=I_M\subset G_D$.
B.146--B.147 prove positivity for
\[
 \Omega^{I^{(r)}}\quad(r\ge0),\qquad
 \Omega^{\,\bigcap_{c\ge1}\gamma_{c+1}(I)},
\]
and algebraic descent as above.
The Gauss valuations on $\mathbb Q(t_1,\ldots,t_s)$
give strictly additional families with residue
transcendence degree $s$ for every $s>0$.
The all-proper-normal-subgroup conclusion of B.143
is not asserted for these more general complete bases.

On the negative side, B.144 proves that positivity
requires a slim absolute Galois group: every open
subgroup must have trivial centre.
It gives explicit henselian fields of every finite
transcendence degree, and all their finite extensions,
whose absolute Galois groups are wild procyclic
$\mathbb Z_p$ groups. Their cores in every open
finite-local Galois group are trivial.
Reconstruction is also false for every field in the
stated finite-transcendence-degree setting admitting
a nontrivial henselian valuation of residue
characteristic zero, by B.49 and B.54--B.55.
This includes a mixed-characteristic valuation with
such a henselian coarsening.
An algebraically closed field is negative by B.53;
in particular the union of all finite derived fields
in the finite-local setting is distinguished from
the positive lower-central union.

Thus the remaining finite-local-subfield cases are
slim and have trivial normal core in every open
local group, apart from any further positive cases
already covered by the stated criteria.
Outside that setting, other algebraic Galois
subgroups, residue-field configurations and
nondiscrete valuations remain in the quantifier.
Neither unbounded wild ramification nor a
non-torsion Kummer kernel alone decides it.
\end{claim}
\begin{proof}[Scope verification and unresolved step]
B.136 proves the arbitrary-source norm-polynomial step
at every prime. B.137--B.138 verify the uniform two-grid
and local normalization inputs for the entire
lower-central union, including $2,3$.
Their final Carlson--Haine--Wolf reduction reconstructs
the full relative inverse-image functor.

B.139 proves the finite-extension and completion
comparison for the relative finitely generated base.
B.140 checks the finite-derived and nilpotent stable
Kummer kernels, two-grid rigidity, arbitrary-source
norm recovery and algebraic descent in that setting.
B.142 supplies the more general abelian normal
quotient criterion. B.143 proves that every
nontrivial normal subgroup of a finite-local
absolute Galois group is nonabelian; a least strict
derived step supplies the criterion. Algebraic
descent gives the asserted normal-core consequence.
All of these claims include finite extensions.

B.146 treats a complete discrete valuation field
with finitely generated, possibly imperfect residue
field. Its additive logarithmic map injects the
complete field into rational Kummer cohomology;
it is not asserted to be onto.
B.147 uses ratios in deep units to place the exact
scalar error in this subspace. It verifies every
remaining norm and graph input without assuming
local compactness or freeness of the general wild
inertia group. Its final algebraic descent and
two-henselian-valuations argument establish the
claimed additional scope.

For the negative procyclic fields, B.144 constructs
a nontrivial central absolute Galois element and
uses its inverse-image natural transformation to
produce an actual-point map on the affine line
which no polynomial realizes. Its finite-etale
slice argument shows that positivity would pass
to every finite extension, proving the slimness
necessity and the stated negative family.
The residue-zero conclusions use the arithmetic
covering-category comparison and valuation
approximation of B.49 and B.54--B.55; those
residue-characteristic hypotheses remain essential.

For a general remaining mixed-characteristic field,
the norm proof still requires a group controlling
normalized norm equivariance, arbitrarily large
cyclic quotients with the stated normal-subgroup
surjectivity, the scalar logarithmic comparison,
and a stable finite-grid rigidity test.
These inputs have now been verified for all
positive classes above, but not for every field
in the remaining quantifier.
The displayed negative constructions do not
decide those residual cases either.
The full classification therefore remains Open.
\end{proof}
\paragraph{Current proved reductions and remaining scope.} B.136--B.147 remove the stated finite-derived and lower-central families in every finite transcendence degree, all proper Galois algebraic extensions of the specified relative local bases, and all fields with a nontrivial core in some open local group. Explicit wild procyclic fixed fields and every finite extension are negative; slimness is necessary. B.146--B.147 add derived and nilpotent inertia fields over complete mixed-characteristic discrete valuation bases with finitely generated residue fields of arbitrary transcendence degree. The remaining gate is the general mixed-characteristic field classification outside the stated positive and negative classes. B.134 remains a counterexample to general point-evaluation detection of Kummer classes, and is not a reconstruction counterexample.

\paragraph{Reason for placement in Part II.} The stated universal assertion remains unproved; the accompanying text records its exact current gate.

\begin{claim}{B.67}{Open}{The exact unsettled quantifier in B5}
For every prime $p$ and every proper smooth connected curve
$C/\overline{\mathbf F}_p$ of genus at least two, there exists
a connected finite \emph{\'etale} cover $D\to C$ whose geometric
Galois closure has order prime to $p$, and a nonconstant morphism
$D\to H$ to a positive-genus supersingular curve $H$.
This assertion remains undecided here for arbitrary $C$.
\end{claim}

We proved the assertion for every odd-characteristic hyperelliptic
source, including characteristic three, and for the further tame
Galois classes in B.38. In characteristic two we proved the
specified cyclic-three and uniformly triple-ramified cases in
B.37. These are statements about actual morphisms of curves and
prime-to-$p$ Galois closures, not only about nonordinary Jacobians.

We tried the corrected torsion-position construction on an elliptic
quotient, then the low-genus quotient and repeated-branch-index
constructions. The first requires such a quotient or an already
constructed correspondence; the latter require the explicit
Galois and index hypotheses in B.38. For a general tame map,
an unramified point in a branch fibre prevents ramification
cancellation, and tameness alone does not control the prime
divisors of the Galois closure. The general tame Belyi theorem
therefore does not finish this step. Nonordinary covering theorems
do not furnish a supersingular target. These failed reductions
are not counterexamples to the displayed assertion.

The first unhandled genus is two in characteristic two for sources
outside B.37; in odd characteristic the first unhandled genus is
three for nonhyperelliptic sources outside B.38. A next needed
step is a prime-to-$p$ covering construction that creates the
uniformly divisible branch fibres used in B.36--B.38 for one of
these remaining classes, or a genuine obstruction to every such
supersingular correspondence. No finite search is asserted to
exclude all covers of a remaining source.




\paragraph{Current proved reductions and remaining scope.} B.85--B.86 give actual supersingular correspondences on the ordinary cubic and two-rank-one quintic loci, with the entire cover and map written explicitly. The other parameters and higher-genus cases remain open; the bounded exclusions B.83--B.84 are not universal nonexistence results.

\paragraph{Reason for placement in Part II.} The stated universal assertion remains unproved; the accompanying text records its exact current gate.

\begin{claim}{B.118}{Open}{The remaining full CM and general bare-normalizer assertions}
The existence or exclusion of a nontrivial closed procyclic
$C\leq R$ with open arithmetic normalizer image in
$G_{\mathbb Q(i)}$ remains unproved here.  The same general
bare-normalizer assertion for an affine hyperbolic curve over
a number field remains unproved.  The preceding claims prove
the restrictions on any candidate, not its existence or
nonexistence.
\end{claim}
\begin{proof}[Recorded attempt and remaining implication]
A CM candidate lies in $R_{\rm nil}(F')$, dies in every
maximal pro-prime quotient of every multiplication-isogeny
cover, and is detected only by a weight-one line on some
genus-at-least-two covering Jacobian's Prym part.  Its actual
normalizer field contains a CM subfield and satisfies the
local embedding and prime-support conditions above.  A
cyclotomic detector would instead force cuspidality and
contradict the common-kernel condition.

Varying finite covers does not by itself eliminate the
weight-one line: Claim~B.112 proves that there is a detector,
but not that one can choose it in the boundary submodule.
The finite-cover argument would settle the assertion if every
such candidate admitted a cover in which its detected line
has zero proper homology.  That additional assertion is not
proved, and the isogeny covers cannot establish it, because
Claim~B.115 proves that they kill the entire candidate.
Passing from a Prym quotient to an elliptic curve also fails
to produce the needed punctured map unless the preimage of
the removed point is contained in the existing boundary.
The proper Jacobian morphism alone gives no such containment.

The arithmetic-core approach can be carried out exactly, but
it stops at a compatible homology statement.  Start with the
finite-cover detector $H$ of Claim~B.112.  Choose a
characteristic open subgroup $W$ of $F$ contained in $H$.
Then $C_\ell\cap W$ is open in $C_\ell$, and its image
in $W^{\rm ab}_\ell\otimes\mathbb Q_\ell$ is nonzero:
the map to $H^{\rm ab}_\ell$ sends its generator to a
nonzero integer multiple of the original detected vector.
Use the rational-point section for the real model $U'/\mathbb Q$.
Since $W$ is characteristic, $W\rtimes G_{\mathbb Q}$ is
an open subgroup of $\pi_1(U')$ and gives an actual finite
cover $Y/\mathbb Q$ with a chosen rational point over $b$.
The intersection of the old normalizer with this arithmetic
cover group has open image.  After taking a finite Galois
extension $L_1/\mathbb Q$ inside that image, the detected
line is $G_{L_1}$-stable.  Its finitely many translates span
a $G_{\mathbb Q}$-stable subrepresentation $V$ of the
homology of $Y$, and restriction of $V$ to $G_{L_1}$ is a
sum of one-dimensional geometric characters of weight one.
Their span has zero intersection with boundary homology:
at a common good Frobenius the former eigenvalues have
absolute value $q^{1/2}$ and the latter absolute value $q$.
Thus $V$ injects into the proper Jacobian homology.  All its
lines map to zero in $V_\ell E'$, since the map $Y\to U'$
is defined over $\mathbb Q$ and the original line comes
from $C\subset R$.  The entire span is a Prym representation.

This arithmetic core does not create a $c$-stable procyclic
subgroup.  In fact no one-dimensional geometric character
of weight one can be defined over a number field with a
real embedding: the coefficient-function proof of
Claim~B.114, applied to that character itself, would put a
CM subfield in a field that embeds in $\mathbb R$.
If conjugation fixed a $G_{L_1}$-character occurring in $V$,
it would preserve its isotypic summand.  The involution is
diagonalizable there over $\mathbb Q_\ell$, so one of its
eigenlines would give such a character over
$L_1^{\langle c\rangle}$, a field with a real embedding.
Consequently conjugation pairs distinct character summands.
Their sums contain fixed vectors, but those vectors have no
justified lifts to fixed geometric elements.  The explicit
model of Claim~B.116 already demonstrates this obstruction:
$c$ exchanges $b_1,b_2$, the nonzero homology vector
$[b_1]+[b_2]$ is fixed, and nevertheless $C_F(c)=1$.
Passing to an arithmetic core therefore supplies additional
weight-one Prym data, not a contradiction with the fixed-
subgroup calculation.

The external real involution gives actual new restrictions
on $R$ and on full-action-stable subgroups, but a normalizer
with image $G_L$ for a totally imaginary $L$ need contain
no such involution.  No Sylow subgroup or candidate is
silently assumed invariant under $c$.
Finally an open normalizer image does not give an arithmetic
section of that normalizer.  We do not invoke the unrefereed
local section assertion used in LM6 A.61--62, nor substitute
an infinite-index normalizer for the fundamental group of
a finite-type source in an open Hom theorem.  The missing
all-cover weight-one exclusion or a compatible full profinite
construction of a candidate is retained as the exact remaining
task, rather than reported as a solved reduction.
\end{proof}
\paragraph{Reason for placement in Part II.} The stated universal assertion remains unproved; the accompanying text records its exact current gate.

\bibliographystyle{plainnat}
\bibliography{ANABELIAN_A1_SUCCESSOR_references}
\end{document}
